\documentclass[dvipdfmx,a4paper]{jsarticle}

%---------------- 数式・記号 ----------------
\usepackage{amsmath,amssymb}
\usepackage{amsthm}
\usepackage{mathtools}

%---------------- 図・箇条書き ----------------
\usepackage{graphicx}     % \resizebox（TikZ 図の全幅化）
\usepackage{longtable}     % 付録 B の主要記号一覧
\usepackage{enumitem}     % enumerate[label=...]/description[style=nextline]
\usepackage{tikz}
\usetikzlibrary{fit,backgrounds,arrows.meta,calc,positioning}
%---------------- 引用 ----------------
\usepackage[backend=biber]{biblatex}
\addbibresource{references.bib}
%---------------- 定理様 環境（§3 基準・全セクション共通）----------------
\newtheoremstyle{breakstyle}
  {10pt}       % 上の空白
  {10pt}       % 下の空白
  {\normalfont}
  {}
  {\bfseries}
  {.}
  {\newline}  % 見出しの後で改行
  {}

\theoremstyle{breakstyle}
\theoremstyle{definition}
\newtheorem{definition}{定義}[section]
\newtheorem{theorem}{定理}[section]
\newtheorem{assumptioninner}{仮定}[section]
\renewcommand{\theassumptioninner}{A\thesection.\arabic{assumptioninner}}
\newenvironment{assumption}
  {\begin{assumptioninner}}
  {\end{assumptioninner}}
\newtheorem{proposition}{命題}[section]
\newtheorem{lemma}{補題}[section]
\newtheorem{observation}{観察記録}[section]
\newtheorem{remark}{備考}[section]
\newtheorem{problem}{問題}

%---------------- メタ情報 ----------------
\title{自律認知システム理論：倫理的協調知性の数理的基盤\\（Phase 1：決定論的定式化）}
\author{NF}
\date{}
% PDF の作成者欄を固定する（hyperref を使わないので dvipdfmx の special で指定する）
\AtBeginDvi{\special{pdf:docinfo<</Author(NF)>>}}

\begin{document}
\maketitle

%================ 要旨 ================
\begin{abstract}
本論文は、
自己と環境を区別し、
入力を受けて行動を選択し、
内部状態を遷移させ、
内部モデル状態を更新しながら存続を図る系を、
共通の形式で記述する対象として自律認知システム
（ACS：Autonomous Cognitive System）を定式化する。
本理論は著者が研究する倫理的協調知性理論
（ECIT：Ethical Coordinative Intelligence Theory）の基盤として、
倫理的協調を可行性保存の観点から形式的に検討するための
前段階の形式枠組みを与えることを目的とする。
本論文の形式的内容は
Phase~1（決定論的・離散時間で、状態・行動・内部モデル状態空間にコンパクト性を仮定する枠組み）
に限定され、
単一ACSの内部動学を中心に扱う。
組織ACSについては、
単一ACS枠組みを社会的組織へ構造的に適用する
限定的な記述として扱う。
ACSを7タプル $\Sigma=(X,U,E,f,\Pi,M,\mathrm{Update})$ として定義し、
持続評価関数 $\Pi$ から派生する
可行集合 $K:=\{x\in X\mid \Pi(x)\geq 0\}$ を可行性記述の中心的対象とする。
長期的な可行性の維持に関わる制約軸として、
自己維持（SM）・環境利用（EU）・可行応答（FR）の三原理を導入し、
それぞれの時刻別充足述語と論理的地位を区別して整理する。
FRは、必要な局面で可行な応答が得られること、
およびその必要応答ポテンシャルの保持・獲得・回復・再構成に関わる制約軸である。
これら三原理を原理水準の制約軸とし、
これとは別の機構水準の記述として、
可行性に関わる内部評価と応答の形成を担う
三機構情報処理アーキテクチャ
（TIPA：Triadic Information-Processing Architecture）、
その作動を支える物理的基盤層（PSL）および
実装部位層（ISL）からなる内部機構構造を与える。
あわせて、一段の作動記録が次の時刻へ閉じる標準作動形式と時刻間の接続、
および外部対象との関係性の最小構造と関係品質を定式化する。
さらに、可行行動が存在しない閉塞と、
構造的個体としての同一性を成立させていた組織関係が
同一個体として回復しない形で失われた崩壊とを区別する。
また、ACSが参加する集団・組織と
組織自身のACS性を区別し、
組織自身について整形式なACS構造が実現する場合を、
個体ACSの参加の有無によらず組織ACSとして記述し、先行節の定義を組織水準へ適用する。
ACS理論は、異なる自律主体の行動を、特定の文化的善悪判断に先立って、
可行性上の要請、内部状態、外部環境、関係および内部機構の作動との
構造的な関係として記述するための共通の形式を目指すものであり、
本論文はその決定論的・構造的な基礎を与える。
本論文は記述理論かつ構造理論の立場をとり、
人間、生物、AI、ロボット、組織等への
実証的妥当性を本論文の形式的結果から直接には主張しない。
決定論的枠組みのもとでの補助的詳細化はPhase~1.$\alpha$へ、
確率的拡張、連続時間化、因果的関係動学、
分岐を許す時間的自己同一性、
および社会動態の定量化はPhase~2へ、
ECITとの正式な接続はPhase~3へ委譲する。
\end{abstract}
\bigskip
\noindent\textbf{キーワード}：自律認知システム、可行性、Viability Theory、持続評価関数、三原理（自己維持・環境利用・可行応答）、可行性保存制約、TIPA、組織ACS、閉塞、崩壊。


%================ Section 1 導入 ================
%======================================================================
\section{導入}
\label{sec:intro}

本論文は、自律的に振る舞う認知システムが長期にわたり存続するための
構造的条件を、数理的な枠組みのもとで記述することを目的とする。
本論文では、自己と環境を区別し、入力を受けて行動を選択し、
内部状態と内部モデルを更新しながら存続を図る系を、
自律認知システム（ACS：Autonomous Cognitive System）として抽象化する。
人間、自律的AIシステム、および後続節で構造的に適用される
組織的システムは、この抽象化の適用先に含まれうるが、
各領域における実証的な妥当性は本論文の主張に含めない。
本論文は、ACSの存続を支える条件を、
規範的な命令としてではなく、可行性に関する構造的な制約軸として定式化する。

本論文が扱う範囲は Phase~1、
すなわち決定論的・離散時間・コンパクト空間・単一ACSの内部動学を
中心とする枠組みに限定される。
確率的拡張、連続時間化、社会動態の本格的展開は後続の段階に委譲する。
以下ではまず背景と動機を述べ（\S\ref{subsec:background}）、
本論文の立場とスコープ、および構造・機能概念の分類規則を定め
（\S\ref{subsec:scope}）、
主要概念を予告したうえで（\S\ref{subsec:preview}）、
本論文の貢献（\S\ref{subsec:contributions}）、
関連研究（\S\ref{subsec:related}）、
全体構成（\S\ref{subsec:roadmap}）を順に述べる。

%----------------------------------------------------------------------
\subsection{背景と動機}
\label{subsec:background}

社会のなかで生きる主体の行動は、
熟慮による選択として現れることもあれば、
熟慮的な選択過程を経ない局所的・反応的・自動的な行動として
現れることもある。
こうした行動は、多くの場合、
その主体が属する社会や文化の規範に照らして評価される。
しかし規範は社会や文化ごとに異なり、
同じ行動が一方では妥当とされ、他方では非難されることがある。
規範による評価だけに依拠すると、
異なる規範のもとにある主体どうしが、
互いの行動を共有可能な形で理解する手がかりを得にくい。
人間と自律的AIシステムのように、
成り立ちの異なる主体が同じ社会のなかで行動する場合には、
この問題はいっそう明確になる。

本論文の関心は、
主体の行動がどのような構造的条件と結びついて成立するのかを、
特定の社会や文化の善悪の判断を前提としない形式で記述することにある。
本論文では、行動の成立に関わる可行性上の要請、
内部状態と内部モデル、外部環境、および内部機構の作動の組を、
その行動の構造的源泉と呼ぶ。
構造的源泉は、行動の単一の原因を指す語ではない。
この記述のために本論文は、
人間、自律的AIシステムその他の自律的に振る舞う系を、
共通の抽象構造をもつ自律認知システムとして定義する。
この共通構造のもとでは、
主体の行動を善悪の基準によって判定するのではなく、
どの可行性上の要請が関わり、
どの内部機構の作動を経てその行動が成立したかとして記述できる。
熟慮的な評価を経る応答と、
局所的な保護応答や自動化された処理による応答も、
同じ構造のもとで区別して記述される（\S\ref{sec:tipa}）。
また、内部モデルや評価過程が関与する行動については、
主体が利用可能な情報にもとづいて形成した評価と、
形式モデル上の可行性とが一致することを要求しない。
このため、誤認や評価の誤差を伴う行動も、
同じ構造的記述の対象となる。
本論文は、このような共有可能な記述言語の
決定論的・構造的な基礎を与えることを目指す。
ここでいう記述は、行動の予測や原因の一意な特定を意味しない。
その射程の限定は \S\ref{subsec:scope} で述べる。

本論文が提示するACS理論
（自律認知システム理論／Autonomous Cognitive System Theory）は、
著者が研究を進める倫理的協調知性理論
（ECIT：Ethical Coordinative Intelligence Theory）の
基本構造理論として位置づけられる。
ECITは、人間とAIを含む複数の主体が、
対立を固定化させずに協調を続けるための枠組みをどう与えるか、
という問いに関わる。
この問いを論じるためには、
協調の当事者となる主体それぞれの行動を、
特定の規範を前提とせずに共有可能な形で記述する基盤が要る。
ACS理論は、その記述の基盤を与える理論として構想されている。
すなわちACS理論は、
三原理である自己維持・環境利用・可行応答、
自己と環境の構造的結合、
および外部対象との関係構造を記述する枠組みをECITに提供する。
これらがECITの協調概念とどのように対応するかは、
Phase~3 で検討する。

このとき本論文が採る基本的着想は、
倫理を善悪の判断としてではなく、
系が可行であり続けるための制約として記述することである。
すなわち、「こうあるべき」という規範的要請を外部から導入するのではなく、
長期にわたり可行性を保つ系がどのような構造を備えるかを記述し、
その記述のなかに倫理の礎となる条件が
構造的に現れるかどうかを問う。
行動の構造的源泉を記述するという本論文の関心は、
この着想と結びついている。
主体の行動を、可行性に関する要請と、
それに応答する内部機構の作動との関係として記述することが、
特定の善悪の判断を前提としない記述言語の中身となるからである。
この立場のもとで、三原理は価値的な命令ではなく、
長期可行性に関わる相異なる構造的制約軸として扱われる。

ECITはACS理論の上位に位置づけられる枠組みであるが、
その本格的な展開は後続の段階の課題とし、
ACS理論との接続は Phase~3 に位置づける。
本論文では、ECITの概要を動機づけの水準で非形式的に述べるにとどめ、
その規範的・社会適用的側面を Phase~1 の主張に混入させない。
ECITの具体的な要素やそれらの機構への対応づけ、
社会的応用に関する主張は、
本論文の形式的帰結としてではなく、
Phase~3 における今後の研究方向として位置づけられる
（\S\ref{subsec:phase_outlook}）。

%----------------------------------------------------------------------
\subsection{本論文の立場とスコープ}
\label{subsec:scope}

本論文は、倫理的判断そのものを与えるものではない。
本論文が扱うのは、ACSが長期にわたり可行性を保つうえで、
どのような構造的条件が要請されるかである。
この意味において、本論文は記述理論であり、規範理論ではない。
また構造理論であり、価値理論ではない。
本論文が記述するのは「ACSはいかにあるべきか」ではなく、
「長期にわたり可行であるACSはいかなる構造を備えるか」である。
倫理に相当する条件は、外部から課される規範としてではなく、
可行性保存上の制約として記述される。
ECITが扱う規範的・社会適用的な側面は、
本論文の主張には含めず、後続の段階へ委譲する。

\S\ref{subsec:background} で述べた行動の記述について、
その射程を次のように限定する。
第一に、記述は主体の行動を予測することでも、
ある行動の原因を一意に特定することでもない。
本論文の枠組みは行動選択の一般法則を与えず、
一つの行動に複数の構造的要請と内部機構が関与しうることを認める。
第二に、記述の言語が特定の文化の善悪の判断を前提としないことと、
個別の主体へ適用する際のモデル化が
観察者の選択から独立であることとは区別される。
可行集合や必要環境要因をどのように設定するかには、
モデル化の選択が介在する。
第三に、本論文の共通構造が適用されるのは、
\S\ref{sec:acs_def} の定義と構成要件を満たす系であり、
個別の領域における経験的な当てはまりは
本論文の主張に含めない（\S\ref{sec:limitations_and_outlook}）。

長期可行性に関する本論文の形式的な解析は、
可行集合 $K$ と、必要環境充足に相対する
$E_{\mathrm{req}}$ ロバスト可行性核 $\mathrm{Viab}_f(K)$
（定義~\ref{def:viab_core}）を中心とする。
可行性核に留まる応答を選び続ける軌道についての結論
（命題~\ref{prop:kernel_maintenance}）は、
必要環境要因の充足と行動の選択に関する前件のもとでの条件付きの結果であり、
本論文は、それらの前件が実在のACSにおいて維持されることを主張しない
（備考~\ref{rem:kernel_response_scope}）。
また本論文は、可行集合への所属と、機能的終端または崩壊の不在との同値を
一般的な帰結として置かない
（備考~\ref{rem:principle_time_domain_scope}、備考~\ref{rem:sm_not_collapse}）。

本論文の形式的内容は Phase~1 に限定される。
Phase~1 は決定論的・離散時間・コンパクト空間を前提とし、
単一ACSの内部動学を中心に扱う。
この枠組みのもとでは、確率的な記述や期待値的な評価は用いず、
有限窓における決定論的な集計量のみを用いる。
Phase~1 の前提と各対象の型を変えずに、
規則・閾値・操作的基準を具体化する補助的詳細化は Phase~1.$\alpha$ とし
（\S\ref{subsec:phase1a}）、
確率的拡張・連続時間化・社会動態の定量化は Phase~2 の課題とし、
ECITとの本格的な接続は Phase~3 の課題とする。

なお、本論文で扱う組織ACSは、
社会動態の一般理論ではなく、
組織自身がACSの定義を満たす場合に、
単一ACSの枠組みを社会的組織へ構造的に適用する記述として位置づけられる。
したがって、社会的相互作用の確率的変動、制度形成、
集団内の長期的進化、ネットワーク構造の動的変化を
一般的に記述することは、Phase~1 の主張には含めない。
社会は、ACSその他の主体が帰属・所属・活動する場として扱う
（定義~\ref{def:org_terms}）。

本論文の各主張は、
明示された仮定のもとで（under stated assumptions）、
本枠組みの範囲内で（within present framework）
成立するものとして述べられる。

\subsubsection{本論文が対象とするACSと「必要」の意味}
\label{subsubsec:real_world_scope}

ACSは、物理的制約のもとで作動する実在のシステムを記述するための
概念である。
生物学的に実現されるACSはもとより、
計算機上で実現されるACSも、
有限の計算資源、有限の応答時間、有限の記憶容量、
供給とエネルギーの制約、
および構成要素の損傷・劣化可能性から自由ではない。
本論文は、この物理的制約の被拘束性をACSの副次的性質ではなく、
その内部機構を論じる際の前提として扱う。
したがって本論文は、物理的に実装されない仮想的なACSを
別の型として立てず、ACSを一律に扱う。

本論文が想定する環境は、
時間変化、部分観測性、有限資源、有限応答時間、
損傷・劣化可能性を伴い、
他主体との相互作用を含みうる環境である。
ここで、他主体との相互作用を「含みうる」とすることは、
Phase~1 において多主体動学を一般的に形式化することを意味しない。
単一ACSから見た他主体の作用は、
必要に応じて外部環境入力の一部として取り扱われ、
多ACS間相互作用そのものの動学的展開は後続段階へ委譲する。

一方、完全に固定された環境、完全情報が与えられる環境、
実質的に無制約な資源をもつ環境その他の退化した設定は、
本論文の考察対象に含めない。
この除外の一部は、後続節において形式的な仮定として与えられる。
すなわち、環境列が時間的に一定でないことは
\S\ref{sec:three_principles} の非定常環境の仮定
（仮定~\ref{ass:nonstationary}）として、
外部環境状態の差異が可行行動集合に一切影響しない退化的な系の除外は
\S\ref{sec:acs_def} の可行行動集合に対する環境入力の非退化性の仮定
（仮定~\ref{ass:env_viable_nondegeneracy}）として、
それぞれ導入される。
上に挙げたその他の退化した設定（情報および資源に関するもの）の除外は、
Phase~1 では考察対象の射程を定める宣言にとどめ、
独立の形式的仮定としては置かない。
なお、現在可行な行動の少なくとも一つについて、
必要環境要因がすべて充足される外部環境状態のうちに
次状態を $K$ に置かないものがあること
（可行な応答の弱点、定義~\ref{def:nonresponse_env_set}）は、
対象クラスを限定する条件ではなく、
\S\ref{sec:three_principles} の応答多様性縮退と長期可行性の条件付き解析において
命題の前提としてだけ用いられる。

以上より、本論文で以下「必要」と記す場合、
別途より強い論理的必然性が証明されている場合を除き、
上記の環境において長期的可行性の維持を指向するACSに対する
構成的必要性を意味する。
上記の退化した設定を用いてこの構成的必要性を否定することは、
本論文の射程外とする。

この射程の限定は、後続節で個別に証明される
定義的必然性、条件付き必然性または数学的必要性の強さを
弱めるものではない。
とくに、\S\ref{sec:three_principles} において
明示的な仮定のもとで証明される
各命題・補題の論理的地位は、
各結果自身の仮定、定義域および限定によって定まり、
本項の射程限定によって置き換えられない。

\subsubsection{構造・機能概念の三分類}
\label{subsubsec:abc_classification}

本論文は、ある機能がACSに要求されることと、
その機能を本論文がどのような形式で表現するかと、
その表現を解析するために導入する理論量とを、
それぞれ別の事柄として扱う。
これらを区別せずに読むと、
本論文が採用する機構分解や内部表現が、
実在ACSに一対一で備わる部位として要求されているかのように
受け取られうる。
そこで次の三分類を導入する。

\begin{definition}[構造・機能概念の三分類]
\label{def:abc_classification}

本論文では、ACSの構造、機能、機構形式および時刻依存量について、
それらの理論上の地位を区別するため、次の三分類を用いる。

\begin{description}

\item[分類A：構成的必要機能]
\S\ref{subsubsec:real_world_scope} で述べた環境において
長期的な可行性の維持を指向するACSに対して、
構成的に要求される機能をいう。
分類Aは機能的必要性を表し、
当該機能が独立した物理部位、計算モジュール、
状態変数または明示的内部表現として
実装されることを要求しない。
同じ機能的要請を満たす機能的に等価な実現を許す。
分類Aへの指定それ自体は、
その必要性があらゆる数学的環境について
証明済みであることを意味しない。
その必要性の射程は \S\ref{subsubsec:real_world_scope} に従う。

\item[分類B：標準機能形式]
分類Aの構成的必要機能その他のACS内部機能を、
本論文の形式体系において分析・比較・記述可能な責務へ分解するために
採用する標準的な機構、機能分解、内部表現、
接続構造または作動形式をいう。
分類Bは単なる例示ではなく、本論文が採用する標準機能形式である。
一方、生物学的または人工的実装において、
各分類B要素が一対一の独立部位または独立モジュールとして
存在することは要求しない。
また、分類Aの機能が機能的に等価な別の実装構成によって
担持される可能性を排除しない。

\item[分類C：記述・診断・比較概念]
ACSの状態、作動水準、時刻間変化、能力、偏り、健全性、
外部対象との関係、および失敗状態への近接その他を
記述、診断、比較または追跡するために
本論文が導入する理論上の構成物をいう。
これには、スカラー量およびベクトル量のほか、
記述子、カテゴリー、状態帯、集合、経路集合その他の
分析上の構成物を含む。
分類Cの概念は、実在ACSが同一の記号、尺度または内部変数として
保持または計算することを要求しない。
分類Cへの指定は、
その概念がACSの7タプルへ新たな構成要素または
独立状態座標を追加することを意味しない。
また本論文は、別途定義・仮定・命題による接続がない限り、
分類Cの特定値、カテゴリーへの所属または状態帯への所属のみから
三原理の時刻別充足述語の真偽、閉塞または崩壊を一般的な帰結として置かない。

\end{description}
\end{definition}

\begin{remark}[三分類の適用規則]
\label{rem:abc_classification_rules}

定義~\ref{def:abc_classification} の三分類は、
以下の規則に従って適用する。

\begin{enumerate}
  \item 分類は、名称の印象ではなく、
        当該の定義・備考等が直接導入する理論対象に対して付与する。

  \item 原則として、一つの理論対象には一つの分類を付与する。
        一つの記述が分類Aの必要機能と分類Bの標準実現形式を
        同時に含む場合には、
        必要機能と実現形式を別の理論対象として分離する。

  \item 分類Aは機能的必要性を表す。
        したがって本論文は、分類Aの機能が必要であることから、
        特定の分類B機構、物理部位または計算モジュールが
        一対一に存在しなければならないことを一般的な帰結として置かない。

  \item 分類Bは本論文が採用する標準機能形式を表す。
        分類Bであることは単なる任意例を意味しないが、
        対応する機能の物理的・生物学的・計算的実装を
        一意に指定するものでもない。

  \item 分類Cは記述・診断・比較のための概念であり、
        ACSの7タプルへ新たな構成要素を追加するものではない。
        実装が分類Cと同型の内部量を実際に保持する場合でも、
        その事実だけで分類を変更しない。

  \item 三分類は、理論対象が本論文内で担う
        機能的・表現的・記述的地位を区別する軸である。
        これに対し、\S\ref{sec:three_principles} で用いる
        定義的必然性、条件付き必然性、構造的要請、
        ならびに個々の定理・命題の証明強度は、
        論理的地位・主張強度を表す別の軸である。
        したがって、分類Aであること自体から
        数学的必要性または定理としての証明済み性は導かれず、
        逆に定理として証明された対象の論理的地位が
        分類A/B/Cによって弱められることもない。
        仮定の導入地位および強さも、三分類とは別の軸である。

  \item 本論文は、分類Cの概念の端点、閾値、特定値、
        カテゴリーへの所属または状態帯・経路集合への所属から、
        別途明示された接続条件なしに
        閉塞 $A_t=\emptyset$ または
        崩壊 $\mathrm{Col}_t=1$ を一般的な帰結として置かない。

  \item 本分類は、主として \S\ref{sec:tipa} 以降で導入される
        機能、機構形式、補助的内部表現、作動形式、
        時刻依存量、関係記述および失敗状態の分析概念に適用する。
        \S\ref{sec:acs_def} のACS 7タプルを構成する基礎集合・写像そのもの、
        \S\ref{sec:three_principles} の三原理、
        および \S\ref{sec:acs_def} における
        閉塞・機能的終端・崩壊そのものの基礎定義は、
        本分類より上位の基礎定義・構造条件または対象現象として扱い、
        原則として分類A/B/Cの表記対象外とする。
        一方、\S\ref{sec:rel_min} の関係記述子・関係カテゴリー等、
        および \S\ref{sec:collapse_cond} の記述的状態帯・経路集合等は、
        その直接の理論的役割に応じて原則として分類Cの対象とする。
\end{enumerate}

後続節では、分類対象の初出となる正式な定義または備考の表題末尾に、
原則として「（分類A）」「（分類B）」「（分類C）」のいずれかを付す。
同一対象を後続箇所で参照するたびに分類表記を反復する必要はない。
\end{remark}

\subsubsection{形式表現とACS内部認識の区別}
\label{subsec:description_levels_global}

本論文で導入する数式、指標、分類、状態帯、経路集合その他の形式対象は、
理論上の記述・分析のための表現である。
本論文は、その存在だけから、
ACS自身が同一の記号・尺度で当該量の真値を内部に保持し、
完全に認識し、または直接計算することを一般的な帰結として置かない。

ACS内部での評価、診断、応答形成または更新に必要な情報と機能は、
各節で明示される機能要件および標準機能形式に従う。
したがって、
論文上の形式表現が整形式に定義されること、
ACS内部で当該対象を測定・認識・利用できること、
および外部観測者が当該対象を測定できることを区別する。

想定する読者は、Viability Theory、凸解析、動的システム論、
Clarke劣微分の基礎を持つ専門家である。
本論文では、これらの基礎的概念については逐一の導入を行わず、
必要な箇所で参照する。

\subsubsection{記号と添字の規約}
\label{subsec:notation}

本論文の記号では、対象を識別する添字を下付きに置き、
種類、個体、成分、窓、時刻の順に並べる
（例えば $d_{\mathrm{eff},i,t}$、$\bar Q_{i\leftarrow j,W_Q,t}$）。
上付きには、演算と、同じ基本記号で表される対象の変種を区別する固定したラベル
（$E^{\mathrm{bg}}$ 等）だけを置く。
括弧内には、写像、対応または演算の入力を置く。
単一個体だけを扱う文脈では個体の添字を省略する。
規約の全体は付録~\ref{app:notation} に示す。

%----------------------------------------------------------------------
\subsection{主要概念の予告}
\label{subsec:preview}

本論文では、ACSを次の7タプルとして定義する。
\[
\Sigma=(X,U,E,f,\Pi,M,\mathrm{Update}).
\]
ここで $X$ は内部状態空間、$U$ は行動空間、
$E$ は必要環境空間、$f$ は状態遷移写像、
$\Pi$ は持続評価関数、$M$ は内部モデル構造、
$\mathrm{Update}$ は内部モデル更新作用素である。
各構成要素の厳密な定義と型は
\S\ref{sec:acs_def} で与える。
持続評価関数 $\Pi$ から派生する
可行集合
\[
K:=\{x\in X\mid\Pi(x)\ge0\}
\]
と、その時刻ごとの可行行動集合は、
本論文で定める可行性に関わるACSの状態と行動を特徴づける中心的な対象となる。
ACSの7タプルおよびその基礎集合・写像は、
定義~\ref{def:abc_classification} の三分類より上位の
基礎定義として扱う。

ACSの長期可行性に関わる制約軸として、
三原理を導入する。
自己維持（SM）は現在状態が可行集合に属すること、
環境利用（EU）は一段階の可行性に関して臨界となる必要環境要因が充足されること、
可行応答（FR）は臨界必要要因が充足される局面で
次状態を可行集合に置く応答が得られることに対応する。
外部環境状態の差異に対する自己と環境の構造的結合は、
三原理のいずれでもなく EU の前提として扱い、
外部対象との関係品質 $Q$ は \S\ref{sec:rel_min} の記述量とする。
これらの厳密な定式化と、
それぞれが長期可行性に対して持つ論理的地位は
\S\ref{sec:three_principles} で扱う。
三原理は、三分類による機能形式の区別より上位の
制約軸として位置づける。

これらの原理水準の条件と、
ACSにおいてそれらに関わる差異を作動へ接続する
具体的な機能構造とを区別するため、
\S\ref{sec:tipa} では、
情報統合機構M1、適応的評価機構M2、応答行動機構M3からなる
TIPA（Triadic Information-Processing Architecture）、
物理的基盤層PSL、
実装部位層ISL（Implementation Substrate Layer）を定式化する。
ここで、本論文が対象とするACSに構成的に要求される機能は
必要に応じて分類Aとして明示し、
TIPA・PSL・ISLにおける具体的な機構分解、
内部表現、接続構造および作動形式は、
それらを本論文の形式体系内で分析するための
分類Bの標準機能形式として区別する。
この区別により、必要機能の存在を、
特定の生物学的部位または人工的モジュールの
一対一の存在と同一視しない。

また、\S\ref{sec:extended_dynamics} では、
\S\ref{sec:tipa} で定義された同一時刻構造を再定義せず、
一段の作動記録が次の時刻へ閉じる標準作動形式を与え、
内部モデル、機能的記憶、偏向性、実効的発展自由度、
予測精度と実効的時間地平、基盤・実装部位の健全度その他について、
時刻 $t$ から $t+1$ への決定論的な変化を記述する。
このうち、作動水準、変化量、偏り、能力、健全性等を
記述・診断・比較するために導入される理論量は、
主として分類Cとして位置づける。
分類Cの量は、実在ACSが同じ数値表現を内部に保持することを
要求するものではない。

以上により本論文では、
「ACSに構成的に必要な機能」、
「本論文が採用する標準機能形式」、
「理論上の記述・診断・比較量」を区別しながら、
同一の形式的枠組みの内部で接続する。

%----------------------------------------------------------------------
\subsection{本論文の貢献}
\label{subsec:contributions}

本論文の貢献は、
自律的システム、知能、倫理、協調を
概念的に関連づけることにあるのではなく、
これらに関わる構造を単一の形式的枠組みの内部で
接続する点にある。
すなわち、可行性の保存、内部機構、
時刻間の内部動学、関係性の記述、
失敗状態の区別を、
決定論的・離散時間の共通語彙のもとで定式化する。
Phase~1 では確率的記述・連続時間動学・
多主体ネットワーク動態を導入せず、
この限定によって理論の最小骨格を明確化する。
具体的な貢献は次の六点である。

第一の貢献は、ACSの形式的定義である。
人間・自律的AIシステム・組織的システムとして
記述されうる対象を共通の形式で扱うために、
ACSを、単なるエージェント、制御器、
あるいは報酬最大化システムとしてではなく、
\S\ref{subsec:preview} に示した7タプルとして定式化した
（\S\ref{sec:acs_def}）。
この定式化のもとで、
自己維持は外部から与えられた目的の達成としてではなく、
可行集合 $K$ とそこから導かれる可行行動集合を通じて表現される。
これにより、
状態が可行集合を外れること、
可行行動が消失すること、
および構造的個体としての同一性を成立させていた組織関係が
同一個体として回復しない形で失われることを
区別して扱うための基礎が与えられる。
また、外部環境状態を外生背景成分と行動応答成分に分け、
行動応答成分が当該ACSの行動に応じて遷移する形を与える
（定義~\ref{def:env_realization}）。
また、内部モデル $M$ とその更新作用素 $\mathrm{Update}$ を
7タプルの内部に明示することで、
ACSを単なる状態遷移系としてではなく、
自己の状態と環境に関する内部的な処理・更新を伴う
自律認知システムとして記述するための形式的基盤を与える。

第二の貢献は、三原理の定式化とその位置づけである。
長期可行性に関わる三つの制約軸として、
自己維持・SM（Self-Maintenance）、
環境利用・EU（Environmental Utilization）、
可行応答・FR（Feasible Response）
を定義した（\S\ref{sec:three_principles}）。
それぞれの時刻別充足述語は、
SMが現在状態の可行性を、
EUが一段階の可行性に関して臨界となる必要環境要因の充足を、
FRが臨界必要要因の充足される局面で次状態を $K$ に置く応答の存在を問う。
SMは可行集合の定義と同値な条件であり、
EUの時刻別充足述語が偽であることは閉塞を伴い、
閉塞は臨界必要要因の充足の有無によって二分される。
外部環境状態の差異に対する構造的環境結合条件は、
三原理のいずれでもなくEUの前提として位置づける。
三原理の時刻別充足述語の間については、
示した条件付き含意と、本論文が主張しない含意とを区別して整理した
（\S\ref{subsec:principle_linkage}）。
あわせて、三原理に関係する状態を有限時間窓で集約する
三原理時間統合評価条件（T-STAB）を定め、
応答多様性の縮退と可行な応答の弱点が持続する場合に頑健応答不能が持続すること、
およびこれに強仮定を加えると有限時間で閉塞に至ることを示した
（\S\ref{subsec:sd_necessity}）。
以上はいずれも価値判断としてではなく、
本枠組み内の構造的条件として扱われる。
多主体系における均衡の扱いは後続の段階の課題として残る。

第三の貢献は、
三原理を抽象的条件に留めず、
本論文が対象とするACSに構成的に要求される機能と、
それらを本論文内で分析可能にする標準機能形式とを
区別したうえで、内部機構の形式的構成を与えた点である
（\S\ref{sec:tipa}）。
本論文は、情報統合、可行性感応評価、応答形成、
基盤維持、入出力接続、実装担体その他の機能を、
必要に応じて分類Aの構成的必要機能として区別する。
そのうえで、これらを分析可能な責務へ分解する
分類Bの標準機能形式として、
情報統合機構M1、適応的評価機構M2、応答行動機構M3からなるTIPA、
物理的基盤層PSL、実装部位層ISLを定式化する。
また、識別、Priority、要請、候補形成、機能的記憶、
高速処理、Automaticityその他の内部表現・作動形式についても、
それらを現実の独立部位と同一視せず、
本論文における標準機能形式として位置づける。
これにより、原理水準の条件と実装可能な内部構造を切り離さず、
かつ特定の生物学的部位や人工的設計へ還元しない形で、
ACSの内部機構を価値中立的に記述する基礎を与える。
さらに、分類A・B・Cの区別を導入することで、
必要機能、標準実現形式、解析上の記述量を
存在論的に同一視しない読み方を明示した。

第四の貢献は、
単一ACSの時刻間の接続を、
同一時刻の機構構造と区別して整理した点である
（\S\ref{sec:extended_dynamics}）。
\S\ref{sec:extended_dynamics} は、
\S\ref{sec:tipa} のTIPA・PSL・ISLの
同一時刻における構造・責務を再定義しない。
原始的な機能の標準写像と読み出し規則を固定し、
環境実現を与えたとき、
一段の作動記録がただ一つ定まり次の時刻へ閉じること、
およびその記録の因果性を示した。
そのうえで、学習と忘却、
偏向性とその程度である偏向強度、
実効的発展自由度、
予測精度係数と実効的時間地平、
PSL・ISLの健全度を、
状態・内部モデル・外部環境状態または作動記録の有限区間からの読み出しとして時刻 $t$ から $t+1$ へ接続する。
同時に、学習による変化と改善、
「内部量の特定値」と「閉塞・崩壊」を
相互に同一視しない。
これにより\S\ref{sec:extended_dynamics}は、
特定の学習則や収束則を仮定するのではなく、
\S\ref{sec:tipa}の同一時刻構造を
後続の関係性・失敗状態解析へ接続する
時刻間のインターフェースを与える。

第五の貢献は、
ACSと外部対象との関係性を、
補助的な社会的記述としてではなく、
関係継続中に記録される可行性関連量の変化を記述する形式として
導入した点である
（\S\ref{sec:rel_min}）。
本論文では関係性の最小構造を定式化し、
観測、依存、支援、協力、共生、支配、管理、
操作、搾取、競争、回避、対立などを
価値中立的なカテゴリーとして整理する。
この関係記述は、外部対象との接続を
単なる有無として扱うのではなく、
関係継続中の有限窓において記録される
可行性、応答可能性、時間的余裕その他の関連量の変化を
記述するための形式的骨格を与える。
これらの差分は、外部対象の単独因果効果を意味しない。
また、評価可能な関係について、
記述者が指定する評価規則に相対する関係品質 $Q$ を導入する。
$Q$ は構造的環境結合条件やEUの条件ではなく、
ACS内部の評価でもない。
一方でPhase~1では、
関係ネットワークの一般動学や
社会的相互作用の確率過程までを主張しない。

第六の貢献は、
失敗状態の区別と組織への構造的適用である。
可行行動が存在しない閉塞
$A_t=\emptyset$
と、
構造的個体としての同一性を成立させていた組織関係が
本枠組みに相対して回復しない形で失われた崩壊
$\mathrm{Col}_t=1$
とを別個の対象として定式化し（\S\ref{sec:acs_def}）、
両者の関係、外的回復可能性、軌道記録と到達経路を整理する
（\S\ref{sec:collapse_cond}）。
これにより、
「現時点で可行な行動が存在しないこと」と、
「同一個体としての構造的自己が回復不能となること」を
同一の失敗概念へ還元しない。
さらに、単一ACSの枠組みを、
単なるACS集団と区別された組織ACSへ適用する
（\S\ref{sec:org_acs}）。
組織自身がACSの定義を満たす場合には、
先行節の閉塞・崩壊・三原理の定義を組織水準へそのまま適用する。
参加する個体ACSの数は組織ACSの成立の条件に含めず、組織ACSの区分の記録として扱う。
構成員変化を含む組織・社会動態の一般化と定量化は Phase~2 の課題とする。

以上の成果の理論的限界と未解決問題は
\S\ref{sec:limitations_and_outlook} で開示され、
Phase~1.$\alpha$（決定論的枠組みのもとでの補助的詳細化）への委譲事項と、
Phase~2（確率的拡張・連続時間化・組織／社会動態）
および Phase~3（ECITとの接続）への展望が整理される。
Phase~1 の目的は、
経験的妥当性や確率的予測を主張することではなく、
倫理を可行性保存制約として記述する立場のもとで、
ACS理論の最小構造を明示し、
後続の段階への接続基盤を与えることにある。
各主張の強さは、
対応する定義・仮定・証明の強さを超えないように制御される。

%----------------------------------------------------------------------
\subsection{関連研究}
\label{subsec:related}

本論文の数理的基盤は、
いくつかの既存の枠組みに依拠する。

可行集合および可行核の概念は、
Viability Theory の枠組み
\cite{Aubin1991,AubinFrankowska1990} に基づく。
ただし、Viability Theory は
可行集合と可行核の一般理論を与えるものであり、
本論文が扱う内部機構の構成、
関係性記述、
閉塞と崩壊の区別といった記述は、
その主題に含まれない。
本論文の位置づけは、
可行性の語彙とこれらの構造記述を
単一の形式的枠組みの内部で接続する点にある。

持続評価関数の劣微分構造は、
非平滑解析における
Clarke の一般化勾配・一般化劣微分の枠組み
\cite{Clarke1990} に依拠して記述される。
ただし、本論文は非平滑解析そのものを拡張するものではなく、
持続評価関数の境界構造を記述するために
既存の数学的道具を用いるにとどまる。

標準的なAIエージェント論
\cite{RussellNorvig2020} は、
知覚に基づいて行動を選択する主体を扱う。
本論文はこの視点と接続するが、
ACSを知覚から行動への写像としてのみ扱うのではなく、
可行性評価、内部モデルの更新、
閉塞、崩壊を備えた構造として扱う点で立場を異にする。

同様に、強化学習 \cite{SuttonBarto2018} は
報酬最大化を中心概念として行動選択を扱うが、
本論文は報酬を第一原理としない。
ACSの存続基準は、
何らかの量の最大化としてではなく、
持続評価関数の符号条件
$\Pi(x_t)\ge0$、
すなわち可行性の保存として与えられる。
報酬に基づく評価・選択は、
TIPAにおける行動比較の一つの実装形態として
読みうるが（\S\ref{sec:tipa}）、
この対応の形式化は本論文の主張には含めない。
また、内部状態の設定点からの距離を駆動とし、
駆動の減少を報酬として扱う恒常性強化学習がある
\cite{KeramatiGutkin2014}。
本論文の自己維持（SM）の時刻別充足述語は、
設定点への接近ではなく閾値に対する符号条件 $\Pi(x_t)\ge0$ であり、
報酬を定めない。

また、自己維持を必須変数の維持として捉える発想は、
Ashby の古典的議論 \cite{Ashby1952} とも接続する。
本論文はこの自己維持の問題意識を共有する一方で、
可行集合、可行行動、内部モデル、
TIPA・PSL・ISL、関係性、
閉塞・崩壊を同一形式体系内で区別して記述することを目指す。

Ashby はさらに、調節の問題を必要多様性の法則として論じた
\cite{Ashby1956}。
\S\ref{sec:three_principles} の応答多様度
（定義~\ref{def:mu_eff}）は、
$E_{\mathrm{req}}(x)$ のうち次状態を $K$ に置く環境の範囲が等しい行動を
一つに数えた同値類の数であり、
外乱の多様性と調節の多様性を比べる法則を述べるものではない。

生命的な自律性の研究では、
系が自己の可行性の条件に対してどのように自身を調節するかが論じられてきた。
Di Paolo は、可行性の条件に関して自己を調節する性質として適応性を論じた
\cite{DiPaolo2005}。
Barandiaran と Egbert は、最小の代謝系のモデルで
可行な領域と不安定な領域を区別し、
可行な領域へ戻すのに必要な環境の変化を用いて、
規範への従属を定義した
\cite{BarandiaranEgbert2014}。
Barandiaran らは、行為主体の定義に
個体性、相互作用の非対称性、規範性の三条件を置いた
\cite{BarandiaranDiPaoloRohde2009}。
これらの三条件は本論文の三原理とは異なる。
本論文は、可行性条件と自律系の行動・調節との関係を扱う点で、
これらの研究と接点をもつ。
ただし本論文は、規範がエージェントに内在することを主張せず、
可行性に関わる構造的制約を
自己維持（SM）、環境利用（EU）、可行応答（FR）の
三原理として分けて定式化する。
また、Maturana と Varela は
生命系の自律性を自己産出（オートポイエーシス）の概念で論じた
\cite{MaturanaVarela1980}。
本論文の構造的自己（定義~\ref{def:s_struct}）は
7タプルとして定め、自己産出を条件に含めない。
Montévil と Mossio は、
生物の組織を、過程に作用する制約が互いに依存し合い互いを維持する閉包として特徴づけ、
これを相互作用する生物系の境界を定める手がかりとしても用いた
\cite{MontevilMossio2015}。
本論文の崩壊の条件 (SC-1)（定義~\ref{def:collapse}）は、
構造的同一性を成立させる組織関係の喪失を述べるが、
その組織関係を制約の閉包として特定することは定義に含めない。

自己と環境の境界、内部状態、行動、存続を一つの説明の体系で扱う研究として、
自由エネルギー原理がある。
Friston は、マルコフブランケットによって内部状態と外部状態が分けられる確率的な力学系では、
内部状態が自由エネルギーと呼ばれる量を最小化するように見え、
環境をモデル化してそれに作用することで自己の構造と機能を保つ
能動的な推論を行っているように見えると論じた
\cite{Friston2013}。
本論文も、自己--環境境界（定義~\ref{def:b_se}）、
内部モデルとその更新（定義~\ref{def:internal_model}、定義~\ref{def:update_operator}）、
行動、および存続を一つの形式で扱う。
ただし本論文の形式は決定論的であり、
時刻別の可行性を持続評価関数の符号条件 $\Pi(x_t)\ge0$ によって表し、
何らかの量の最小化を原理として置かない。
また、Kolchinsky と Wolpert は、
系が環境について持つ情報のうち、
系の存続の度合いを表す可行性関数の値を保つのに因果的に必要なものを意味情報と定義した
\cite{KolchinskyWolpert2018}。
本論文の必要環境要因（定義~\ref{def:environment_space}）は、
存続に必要な外部環境の条件を、
状態と外部環境状態の組に対する充足判定写像として与えるものであり、
系が環境について保持する情報を定義しない。
Ruffini は、決定論的で閉じた世界と調節器の組について、
アルゴリズム的な複雑さによる調節の性能指標と、
調節器が世界について持つ情報との関係を分析した
\cite{Ruffini2026}。
本論文は内部モデル $M$ を7タプルの成分として置くが（定義~\ref{def:internal_model}）、
この性能指標を採用せず、内部モデルの必要性を調節の性能から導かない。
López-Díaz と Gershenson は、
生物の行為者性が、異なる時間尺度で減衰し再生される制約に依存するとして、
自律性、目標指向性、行為者性、開放性の四つの条件を構造的な階層として整理した
\cite{LopezDiazGershenson2026}。
本論文の Phase~1 は、決定論的な離散時刻の標準形に限られ、
このような複数の時間尺度にわたる組織の動学を定式化しない。

本論文の機構記述のうち、
内部の予測と後の時刻の量との照合は、
予測符号化と接点をもつ。
Rao と Ballard のモデルでは、
上位の領野から下位の領野への結合が下位の活動の予測を運び、
下位から上位への結合が予測と実際の活動との残差を運ぶ
\cite{RaoBallard1999}。
\S\ref{sec:tipa} の M1-c は、
予測整合性誤差（定義~\ref{def:m1_temporal_features}）と、
保持記録を通じた予測外の判定（定義~\ref{def:prediction_record_check}）
を記述する。
本論文は、この分野からの機構的導出や
一対一対応を主張するものではなく、
また実証的検証を与えるものでもない。
これらの既存理論との関係を
「ACS理論から導出される」
あるいは
「ACS理論と数学的に同値である」
とは主張しない。
関連分野との対応関係の精密化と検証は、
後続の段階の課題とする。

自動的処理と制御的処理の区別は
Schneider と Shiffrin によって論じられた
\cite{SchneiderShiffrin1977}。
Logan の事例理論は、
格納された事例の想起に依拠する処理を自動的な処理とみなす
\cite{Logan1988}。
これらは、\S\ref{sec:tipa} の Automaticity
（備考~\ref{rem:automaticity_deliberation}）と、
その時刻間の形成・変容
（\S\ref{sec:extended_dynamics} の備考~\ref{rem:s5_fast_change}）
に近い対象を扱う。
本論文は Automaticity を独立した機構として導入せず、
自動的作動と熟慮的作動の間に普遍的な二値境界を置かない。

内部機構の共通の形式を与える試みとして、認知アーキテクチャの研究がある。
Laird、Lebiere、Rosenbloom は、
ACT-R、Sigma、Soar の三つの認知アーキテクチャを踏まえ、
人間に類する心の構造と処理、記憶と内容、学習、知覚と運動について、
共通の計算的な枠組みを提案した
\cite{LairdLebiereRosenbloom2017}。
\S\ref{sec:tipa} の TIPA・PSL・ISL も内部機構の標準形式を与えるが、
可行性に関わる評価を機構の記述に含め、
構成的必要機能（分類A）、標準機能形式（分類B）、記述・診断・比較概念（分類C）を区別する
（定義~\ref{def:abc_classification}）点で目的が異なる。
本論文は、認知アーキテクチャとしての網羅性や実装上の妥当性を主張しない。
Zhang らは、自律的なエージェントを、
定められた同一性と権限の範囲の中で課題の目標と自らの可行性の条件をともに保ち、
世界と自己のモデルを用いて継続的に行動し、
フィードバックを通じて制御の仕組みと長期的な能力を調整する系として定め、
一つの境界と、作動・可行性・反省・学習の四つのループからなる構成を提案した
\cite{ZhangEtAl2026}。
本論文は、TIPA・PSL・ISL とこれらのループとの間に一対一の対応を置かない。
また本論文は、エージェントの構成法を与えるものではなく、
可行性、内部機構、および閉塞と崩壊の区別を、定義と命題として記述する。
Ruffini と Castaldo は、
モデル化・目的関数・計画の三つの機能的役割を備え、
外乱のもとで自らのパターンを保つ調節をもつ系として、
エージェントをアルゴリズム情報理論と熱力学の観点から論じた
\cite{RuffiniCastaldo2026}。
本論文は、TIPA の三つの機構をこれらの役割と一対一に対応させない。

Brooks の階層型の制御では、
下位の層は上位の層が加わっても機能し続け、
上位の層は下位の層の出力を抑制してその役割を包摂しうる
\cite{Brooks1986}。
\S\ref{sec:tipa} の局所保護応答も、
TIPA の出力を用いずに生成される
（備考~\ref{rem:psl4_concurrency}）。
ただし、局所保護応答の発動は M3 の行動指令によって取り消されず、
その局所応答成分は実効行動に保持される
（定義~\ref{def:effective_action_resolve}）。
この点で、上位の層による出力の抑制を許す構成とは異なる。

強化学習の安全性の研究では、
仕様から合成した遮蔽機構（shield）が、
学習器に安全な行動の一覧を与える構成と、
学習器が選んだ行動が仕様に反する場合にだけそれを置き換える構成とが
提案されている \cite{AlshiekhEtAl2018}。
\S\ref{sec:tipa} の実効行動解決は、
TIPA の行動指令とは別の経路で生じる局所応答成分を実効行動に保持する点で
後者と比べられる。
ただし本論文は、仕様の合成や安全性を確保する手法を与えない。

De Jaegher と Di Paolo は、
相互作用の過程そのものが一種の自律性をとりうると論じた
\cite{DeJaegherDiPaolo2007}。
これに対し、\S\ref{sec:rel_min} の双方向合成
（定義~\ref{def:bidirect_relation}）は二つの単方向記録の順序対であり、
相互作用を新しいACSとして構成しない。
複数のACSが参加する組織を一つの系として扱う場合は、
\S\ref{sec:org_acs} の組織水準のACS構造の問題となる。

Castelfranchi は、
他の行為者の心についての信念に基づく弱い社会的行為と、
目標に基づく強い社会的行為を区別し、
目標の委任と採用を協力の構成要素として分析した
\cite{Castelfranchi1998}。
\S\ref{sec:rel_min} の関係カテゴリー（定義~\ref{def:category_set}）も
依存・協力・競争などの名称を含むが、
これらは関係形式を記述する名称であり、
成立条件として行為者の信念や目標の帰属を定めない。
また、関係品質は記述者が定める評価規則に相対する量であり（定義~\ref{def:Q_rule}）、
関係の継続中の可行性関連量の差分を、
その外部対象による因果効果とはしない。

Beer は、組織を管理サイバネティクスの立場から論じた
\cite{Beer1972}。
\S\ref{sec:org_acs} は、
組織にACSが参加することと組織そのものがACSであることを区別し、
組織水準の構造について先行節の定義を適用できる条件を記述する。
特定の統治形態や管理の方式は扱わない。
Masumori は、過程が互いの存続に寄与する関係を、
系の可行性への介入による効果によって定め、
自律性、行為者性、および自己を維持する単位どうしの整合（alignment）を
この関係によって論じた
\cite{Masumori2026}。
本論文の関係記述は、そのような介入による寄与を定義するものではなく、
関係の継続中の可行性関連量の差分を因果効果としない（\S\ref{sec:rel_min}）。
協調的な行動が有効な選択肢として残る条件と、
有効な選択肢が協調的な行動に絞られる条件の分析は、
問題~\ref{prob:capability_paradox} および問題~\ref{prob:ecit_full} に残される。

外乱を伴う系の可行性については、
一方の主体が他方の主体のすべての作用に対して状態を制約集合に保てる
初期状態の最大集合が差別化核（discriminating kernel）として論じられ、
その離散時間版と近似の算法が与えられている
\cite{CardaliaguetQuincampoixSaintPierre1994}。
定義~\ref{def:viab_core} の条件 $(\ast)$ は、
これと同じ $\exists u\,\forall e$ の型の量化条件を用いるものであり、
本論文はこの型の量化を新たに導入するものではない。
ただし、状態空間、外乱の範囲の定め方、および戦略や情報の構造の前提は
文献ごとに異なりうるので、本論文は差別化核との定義の同一性を主張しない。
本論文の構成上の特徴は、外乱の範囲を状態ごとの必要環境充足集合
$E_{\mathrm{req}}(x)$（定義~\ref{def:E_req}）に取り、
その範囲を必要環境要因 $\eta\in E$ と構成要件 CR-1 によって定める点にある。
外部環境状態空間 $E_{\mathrm{ext}}$ のすべての状態に対する頑健な核、
および楽観的な核との包含関係は命題~\ref{prop:kernel_standard_comparison} に示す。
Yamaki と Churiki は、
有限の格子環境で有限の地平の方策を評価する計算モデルを用いて、
生物の認知に類する方策選択を、
可行性の制約のもとで制御可能な将来を広げることとして捉えうるかを検討した
\cite{YamakiChuriki2026}。
本論文の三原理は方策を選ぶ目的関数ではなく（\S\ref{sec:three_principles}）、
FR の時刻別充足述語は将来の制御可能性の大きさを条件に含めない（定義~\ref{def:FR}）。
本論文は、制御可能な将来の拡大から、
可行性の維持を一般的な帰結として置かない。
また、安全制御の分野では、安全集合を関数の上位集合として与え、
その前方不変性を制御によって保つ制御障壁関数の枠組みが用いられている
\cite{AmesEtAl2019}。
可行集合 $K$ と可行行動集合 $A_t$ はこれと同じ形の集合を
一段階の可行性について扱うが、
本論文は安全集合の設計法や制御則の構成を与えない。

可行性理論は回復力の定義にも用いられている。
Martin は、撹乱を受けた後の状態から所与の時間地平までに制約集合へ戻るための
最小の費用を回復の費用とし、その逆数によって回復力を定めた
\cite{Martin2004}。
そこでは、制約を保ち続けられる状態、
所与の時間地平までに有限の費用で回復できる状態、
および回復できない状態が、
回復の費用がゼロ、有限、無限であることによって区別される。
本論文の閉塞（定義~\ref{def:blockage_sec2}）は、
実現した外部環境状態のもとで一段階の可行行動の集合が空であることであり、
崩壊（定義~\ref{def:collapse}）は、
構造的同一性を成立させる組織関係の喪失と、
実現した外生背景成分の列のもとでの通常の自律的更新と
行動応答成分の遷移 $g_O$ による継続、および
許容外的作用（定義~\ref{def:admissible_external_action}）を加えた別の継続に相対する
同一個体としての非回復性によって定まる。
本論文は、回復の費用が無限であることと崩壊とを同一視しない。

可行性の概念を人工的なエージェントに用いる研究もある。
Tsiokos は、資源台帳による実行可能性の制限、頑健な可行性核、
通信路容量としての実行可能なエンパワーメント、経験的なパッケージング自己写像の
四つの検査可能な構成要素によって行為者を定め、
その可行性核を後続台の意味論のもとでの最大不動点として計算した
\cite{Tsiokos2026}。
Marin と Chaudhary は、Aubin の可行性理論に基づいて、
自律的なAIエージェントの実行時の統治の枠組みを提案した
\cite{MarinChaudhary2026}。
本論文の $E_{\mathrm{req}}$ ロバスト可行性核
（定義~\ref{def:viab_core}）は、
$E_{\mathrm{req}}(x)$ のすべての環境に対して自身にとどまる行動が存在する
最大の部分集合として定め、これらの可行性核との同一性を主張しない。
また本論文は、行為者性の指標や統治の手続きを与えるものではなく、
ACSの内部機構と、閉塞・整形式実現の不成立・崩壊の区別
（定義~\ref{def:blockage_sec2}、定義~\ref{def:wf_realization}、
定義~\ref{def:collapse}）を記述する。

Dalrymple らは、AIが外界に及ぼす影響のモデル、
許容される結果の形式的な仕様、
仕様が満たされることの監査可能な証拠を与える検証の三つによって、
AIシステムの安全性について定量的な裏づけを得ることを目指す枠組みを提案した
\cite{Dalrymple2024}。
本論文はAIの安全性を確保する手法を与えるものではなく、
本論文の記述語彙をこうした仕様の記述に用いることは
本論文の主張に含めない。

以上の関連理論は、
本論文の全体を単独で包含する競合理論としてではなく、
可行性、非平滑境界、
エージェント行動、学習、自己維持、
自律性と調節、組織の閉包、自由エネルギー原理と意味情報、
予測、自動的処理、認知アーキテクチャ、階層型の制御、相互作用、社会的行為、組織、
外乱のもとでの可行性核、安全集合と行動の遮蔽、可行性に基づく回復力、恒常性に基づく強化学習、
AIエージェントの可行性と安全性等の
個別側面について、接続する既存基盤または比較の対象として参照される。
本論文の目的は、
これらの既存理論を置き換えることではなく、
それぞれが扱う一部の概念との接点・相違点を明示しつつ、
ACSの可行性、内部機構、時刻間動学、
関係性、失敗状態を
一つの決定論的形式枠組みへ接続することにある。

%----------------------------------------------------------------------
\subsection{本論文の構成}
\label{subsec:roadmap}

本論文の構成は次のとおりである。

\begin{itemize}

  \item \S\ref{sec:acs_def} では、
    ACSを7タプルとして形式的に定義し、
    持続評価関数、可行集合、可行行動集合、
    内部モデル、更新作用素、認知的処理能力の記述方針、
    構造的自己その他の基本語彙を確立する。

  \item \S\ref{sec:three_principles} では、
    三原理SM・EU・FRの時刻別充足述語を定義し、
    それぞれの論理的位置づけ、
    閉塞の二分、時刻別充足述語の間の条件付き含意と含意関係の限定、
    応答多様性縮退と長期可行性の条件付き接続、
    および三原理を有限時間窓上で統合する
    三原理時間統合評価条件（T-STAB）を扱う。

  \item \S\ref{sec:tipa} では、
    原理水準の三原理を再定義せず、
    本論文が対象とするACSに必要な機能を
    分類Aとして必要に応じて明示し、
    それらを本論文内で分析可能にする
    分類Bの標準機能形式として、
    TIPA、PSL、ISLからなる内部機構構造を定式化する。
    あわせて、識別、評価、応答形成、
    機能的記憶、高速処理、Automaticity、
    作動水準等の責務と境界を整理する。

  \item \S\ref{sec:extended_dynamics} では、
    \S\ref{sec:tipa} の同一時刻構造を再定義せず、
    一段の作動記録が次の時刻へ閉じる標準作動形式を与え、
    単一ACSにおける内部モデル、
    機能的記憶、各種記述量、
    PSL・ISLの健全度等の
    決定論的な時刻間変化を整理する。
    また、分類Cの記述・診断・比較量が
    ACSの追加構成要素ではないことを明示する。

  \item \S\ref{sec:rel_min} では、
    ACSが外部対象と取り結ぶ関係性の最小構造を、
    関係継続中に記録される可行性関連量の変化を記述する形式として扱う。
    外部対象の分類、単方向・双方向の関係記述、
    評価規則に相対する関係品質、
    関係カテゴリーおよび方向性を、
    Phase~1 の範囲で価値中立的に整理する。

  \item \S\ref{sec:collapse_cond} では、
    \S\ref{sec:acs_def} で定義した閉塞と崩壊を区別し、
    可行行動集合の消失と
    構造的自己の回復不能性を
    異なる失敗状態として扱う。
    また、外的回復可能性と、それらへの到達経路を
    Phase~1 の決定論的枠組みの範囲で整理する。

  \item \S\ref{sec:org_acs} では、
    単一ACSの枠組みを
    単なるACS集団と区別された組織ACSへ適用し、
    組織水準のACS構造の成立のもとで
    先行節の定義の組織水準への適用を扱う。

  \item 最後に \S\ref{sec:limitations_and_outlook} で、
    Phase~1 の理論的限界、Phase~1.$\alpha$ への委譲事項、未解決問題を開示し、
    Phase~2・Phase~3への展望と結論を述べる。

  \item 付録~\ref{app:notation} に記号添字規約を、
    付録~\ref{app:symbols} に各節で新たに用いる主要記号の一覧を、
    付録~\ref{app:index} に定義・仮定・命題等の一覧を置く。

\end{itemize}


%================ Section 2 ACSの形式的定義 ================
%======================================================================
\section{自律認知システムの形式的定義}
\label{sec:acs_def}

本節では、\S\ref{sec:intro} で予告した自律認知システム（ACS）の概念を、
後続節が共通に参照できる単一ACSの形式的基礎として定義する。

\S\ref{sec:intro} で述べた理論的位置づけ、Phase~1 の全体的スコープ、
構造・機能概念の三分類の一般規則、および論文全体の動機は本節では再論せず、
ACSの基礎型、構成要件、整形式な実現、可行性、原始的な機能、内部モデルおよび
更新構造の定義に必要な範囲のみを確認する。

本節は、ACSの発生起源、実装史または進化過程を扱わず、
ある対象が本論文の意味でACSとして記述されるための抽象的構造を与える。
以下では一貫して「主張の強さ $\leq$ 証明の強さ」という方針を採用する。
本節は閉塞、機能的終端および崩壊を定義するが（定義~\ref{def:blockage_sec2}、定義~\ref{def:functional_termination}、定義~\ref{def:collapse}）、
本節単独で主張しない事項は、本節末尾の「本節が主張しない事項と Phase 委譲」に列挙する。

%----------------------------------------------------------------------
\subsection{本節の記述範囲と共通規約}
\label{subsec:s2_scope}
\label{subsec:sec2_intro}

本小節では、本節全体を貫く \S\ref{sec:intro} との責務分離、
規約・構成要件・仮定の地位の区別、
および Phase~1 における責務境界を整理する。
本節で新たに用いる主要記号の一覧は付録~\ref{app:symbols_s2} に示す。

\subsubsection{導入節との責務分離}
\label{subsubsec:s2_s1_boundary}

論文全体の動機・理論的位置づけ・Phase 区分・三分類の一般規則は
\S\ref{sec:intro} に委ね、本節では形式的定義に必要な内容のみを扱う。
ここでの「定義」は、対象クラスの記述に用いる型と条件を固定する作業を指す。
一般的な倫理的評価、対象実装の経験的同定、既存理論に対する優越性の主張は
本節の結論に含めない。
\S\ref{sec:intro} の分類A・B・Cは、
7タプルそれ自体の成分数を変更する規則ではない（備考~\ref{rem:abc_classification_rules}）。

\subsubsection{規約・構成要件・仮定の地位の区別}
\label{subsubsec:s2_requirement_status}

本節では、記述上の取り決めである規約、ACSとして成立するための構成要件、
および個別命題の成立範囲を限定する仮定を、互いに異なる地位のものとして扱う。

\begin{description}[style=nextline]
  \item[規約（Convention）]
  論文が形式対象をどのように書き分けるかを指定する取り決めである。
  型分離規約 TSC（\S\ref{subsec:wellformedness}）および
  記述水準の規約（\S\ref{subsec:acs_7tuple}）が
  これに当たる。規約はACS自身の機能的能力ではない。
  \item[構成要件（Constitutive Requirement; CR）]
  本論文がACSと呼ぶ対象の成立条件を指定する。
  必要環境整合性 CR-1（\S\ref{subsec:environment_space}）
  および原始的な機能と基盤機能の具備 CR-2
  （\S\ref{subsec:three_primitive_functions}）が
  これに当たる。CR-1・CR-2 は7タプルから機械的に検査できる条件ではなく、
  その実現に対する意味論的条件であり、通常の数学的仮定の件数には数えない。
  \item[仮定（Assumption）]
  定義された対象について追加的な結果を論じる際に用いる条件であり、
  各仮定には導入地位（A・S・O）と強さ（弱・中・強）を付す。
  本節の仮定は \S\ref{subsec:assumption_list_sec2} に一覧する。
\end{description}

したがって、TSC をACS自身の機能的能力として解釈せず、
CR-1・CR-2 を数学的な追加仮定として数えない。
定義~\ref{def:acs} は、規約 TSC のもとで CR-1 および CR-2 を満たす実現を備える
構造をACSとして扱う。

\subsubsection{Phase~1 における責務境界と本節の構成}
\label{subsubsec:s2_phase1_boundary}

本節は決定論的・離散時間・コンパクト内部状態空間の単一ACS構造を対象とする。
時間添字は $t\in\mathbb N_0$、状態更新は $t\mapsto t+1$ とする。
外部環境状態は外生背景成分と行動応答成分に分解し、
外生背景成分の列は決定論的な解析において与えられる実現列として扱い、
その生成法則や発生頻度を本節では規定しない。
行動応答成分は、環境の側の遷移写像に従って当該ACSの行動に応答しうる
（定義~\ref{def:env_realization}）。
確率的拡張、連続時間化、非コンパクト空間への一般化、
必要環境要因の経験的同定手続き、自己--環境境界の動的再編成、
および社会的組織への構造的適用は、それぞれ本節の形式定義とは別の問題であり、
後続節または後続 Phase の課題とする。

本節の構成は次のとおりである。
\S\ref{subsec:wellformedness} で外部背景と自己--環境型分離および規約 TSC を、
\S\ref{subsec:primitive_functions} で原始的な機能と基盤機能を定義し、
\S\ref{subsec:acs_7tuple} でACSの7タプルと各成分の基本型を定める。
続いて \S\ref{subsec:self_formal_definition} で構造的自己・整形式実現・
時間的自己同一性と機能的終端・崩壊を、\S\ref{subsec:environment_space} で必要環境空間と CR-1 を、
\S\ref{subsec:state_space_basic_variables} で状態遷移・行動・外部環境入力を、
\S\ref{subsec:viability_function_and_blockage} で持続評価・可行性・
一段階可行環境を扱う。
\S\ref{subsec:three_primitive_functions} で基礎的な機能、CR-2 および原始的な機能の標準写像形式を、
\S\ref{subsec:internal_model_M} で内部モデル構造 $M$ を、
\S\ref{subsec:update_operator} で更新作用素 $\mathrm{Update}$ と時間平均内部モデル変化量を、
\S\ref{subsec:cognitive_capacity_policy} で認知的処理能力の記述方針を述べ、
\S\ref{subsec:assumption_list_sec2} で規約・構成要件・仮定を整理する。

%----------------------------------------------------------------------
\subsection{外部背景と自己--環境型分離}
\label{subsec:wellformedness}

本小節では、ACS定義に先立ち、内部状態と外部環境状態を
異なる型として扱うための型分離規約を定める。
本論文における「外部」（external）は普遍的な世界分類ではなく、
対象とするACS個体の境界に相対的な概念である。
したがって、複数のACS個体を区別する必要がある場合には、
個体 $\Sigma_i$ に相対的な外部環境状態空間を
$E_{\mathrm{ext},i}$ と書くことができる。
単一ACSのみを扱う本節では、混同の生じない範囲で
個体添字を省略して $E_{\mathrm{ext}}$ と書く。

$E_{\mathrm{ext}}$ は、当該ACSの状態遷移、感知その他の
外部入力記述に用いうる外部環境状態候補の型である。
その元を $e\in E_{\mathrm{ext}}$ とし、
時刻 $t$ の実現外部環境状態を
$e_t\in E_{\mathrm{ext}}$
と書く。
$E_{\mathrm{ext}}$ はACS 7タプルの独立成分ではなく、
7タプルが相対的に定義される外部背景型である。
Phase~1 では $E_{\mathrm{ext}}$ を位相空間として扱い、
状態遷移写像 $f$ の連続性を述べるための基盤とする。
$E_{\mathrm{ext}}$ のコンパクト性は一般には仮定せず、
これを要する命題ではその箇所で追加仮定として明示する。

本論文では、ACSの定義を支える基礎条件を、規約と構成要件の二種に区別する。
規約は、本論文が記述のために採用する取り決めであり、
当該ACSが満たすべき条件ではない。
以下で定める型分離規約 TSC（Type Separation Convention）は、
内部状態型と外部環境状態型を集合表現上でも分離するために採用する規約である。
構成要件系列 CR（Constitutive Requirement）は、
当該7タプル構造体が本論文におけるACSとして成立するために
満たされねばならない条件を表す。
本節では、必要環境整合性を CR-1（\S\ref{subsec:environment_space}）、
原始的な機能と基盤機能の具備を CR-2（\S\ref{subsec:three_primitive_functions}）とする。
いずれも、後続命題の成立範囲を限定する通常の仮定とは区別し、
仮定一覧には含めない。
構成要件の番号は本節における初出順である。

\subsubsection*{型分離規約 TSC（内部状態型と外部環境状態型の分離）}
\label{conv:TSC}

ACS個体 $\Sigma_i$ の内部状態空間 $X$ と、
その個体に相対的な外部環境状態空間 $E_{\mathrm{ext}}$ は、
本論文の集合表現上では異なる型として扱う。
Phase~1 では、この型分離を明示するためタグ付きの互いに素な表現を採用し、
\[
  X\cap E_{\mathrm{ext}}=\emptyset
\]
と書く。

ここでこの式は、内部状態を表す数学的対象と外部環境状態を表す
数学的対象を無標識の集合表現上で混同しないための構成上の規約である。
したがって、物理世界においてACS内部を構成する物質と外部環境を構成する
物質が存在論的に異なる種類であることを主張するものではない。
また、同一の数値列その他の表現が内部状態と外部環境状態の双方に
現れうる場合でも、型タグを付した対象としては区別される。

\begin{remark}[TSCの性格]
\label{rem:TSC_character}
$X\cap E_{\mathrm{ext}}\neq\emptyset$ となる無標識の集合表現を採用した場合でも、
内部状態引数と外部環境状態引数を異なるsortとして指定すれば、
後続の状態遷移写像の型自体は定義可能である。
したがって、TSC は写像の数学的整形式性そのものに不可欠な条件ではなく、
自己と環境の型を集合表現上でも明示的に分離するために本論文が採用する規約である。
タグ付きの互いに素な表現はその実装手段であり、
TSC が要求するのは型分離そのものである。
Phase~2 において自己--環境境界の流動化または再編成を扱う場合には、
本規約の扱いを改めて明示する。
\end{remark}

\begin{remark}[外部環境状態空間の個体相対性]
\label{rem:Eext_relative}
$E_{\mathrm{ext}}$ は対象ACSの境界に相対的であるため、
一般には異なるACS個体 $\Sigma_i,\Sigma_j$ が同一の
外部環境状態空間をもつことを要求しない。
例えば、$\Sigma_i$ にとって他個体 $\Sigma_j$ は
$\Sigma_i$ の外部環境の一部として現れうる一方、
$\Sigma_j$ 自身の内部状態は $\Sigma_j$ の状態空間に属する。
Phase~1 の単一ACS定義では、異なる個体に相対的な
$E_{\mathrm{ext}}$ 間の変換写像または共通母空間を新たに要求しない。
本論文では、これらの対応構造を導入せずに定式化を与える。
個体別の外部環境状態空間を相互に対応づける構造を要する定式化は
Phase~2 の課題とする。
\end{remark}

\begin{definition}[外部環境状態空間の分解]
\label{def:Eext_decomposition}
外部環境状態空間 $E_{\mathrm{ext}}$ を、二つの位相空間の積
\[
  E_{\mathrm{ext}}
  =
  E^{\mathrm{bg}}\times O,
  \qquad
  O
  :=
  \prod_{a\in\mathcal A_O} O_a
\]
として表し、積位相を与える。
その元を $e=(b,o)$、時刻 $t$ の実現外部環境状態を
$e_t=(b_t,o_t)$ と書く。
各成分は次のとおりである。
\begin{enumerate}[label=(\roman*)]
  \item $E^{\mathrm{bg}}$ を\emph{外生背景成分}の空間と呼ぶ。
  その元 $b$ は、当該ACSの行動に応答しない外部側の条件を表す。
  \item $\mathcal A_O$ は有限集合であり、時刻によらず固定する。
  その元 $a$ を\emph{枠}と呼ぶ。
  \item 各 $O_a$ は位相空間であり、区別された元 $\bot_a\in O_a$ を含む。
  $\bot_a$ を枠 $a$ の\emph{不在状態}と呼ぶ。
  $O$ を\emph{行動応答成分}の空間と呼び、
  $o\in O$ の第 $a$ 成分を $o_a\in O_a$ と書く。
\end{enumerate}
$\mathcal A_O=\emptyset$ のとき $O$ は一点集合であり、
$E^{\mathrm{bg}}$ が一点集合であることも排除しない。
\end{definition}

任意の外部環境状態空間は、$\mathcal A_O=\emptyset$、
$E^{\mathrm{bg}}=E_{\mathrm{ext}}$ とすることで
定義~\ref{def:Eext_decomposition} の形に表せる。
したがって、この分解は対象とするACSの範囲を狭めない。
行動応答成分を置くかどうか、どの成分を枠に割り当てるかは、
対象系の記述における選択である。
規約 TSC は分解の後も維持され、
\[
  X\cap\bigl(E^{\mathrm{bg}}\times O\bigr)=\emptyset
\]
である。

\begin{remark}[枠と行動応答成分の読み]
\label{rem:Eext_decomposition_reading}
定義~\ref{def:Eext_decomposition} について、次の読みを定める。
\begin{enumerate}[label=(\alph*)]
  \item 枠 $a$ は個物そのものではなく、
  外部環境のうち当該ACSの行動に応答しうる成分を記述するための位置である。
  個物が現れることや消えることは、
  $o_a$ が $\bot_a$ と $\bot_a$ 以外の値の間を移ることとして表す。
  同一の枠を時刻をまたいで再び用いて、異なる個物を表してよい。
  本論文は、枠の状態の列から個物の同一性を追跡しない。
  \item 行動応答成分は個物に限らない。
  当該ACSの行動に応答しうる外部側の量であれば、
  蓄積量（例えば当該ACSの排出物が周囲に蓄積した量）や
  局所の場（例えば当該ACSの発熱によって変わる周囲の温度）も
  枠に割り当てうる。
  \item $E^{\mathrm{bg}}$ と $O$ の区別は、当該ACSの行動に応答しうるか否かによる
  記述上の区別であり、外部世界の存在論的な分割ではない。
  \item 「対象」「目的対象」の語は、当該ACSの内部での識別に関わる語として
  \S\ref{sec:tipa} で用いる。
  枠の添字集合 $\mathcal A_O$ は理論水準の添字であり、
  ACS内部の対象表現と同一視しない。
  \S\ref{sec:rel_min} の外部対象集合 $\mathcal Y_{-i}$ とも同一視しない。
\end{enumerate}
記号について次を注意する。
外生背景成分の空間 $E^{\mathrm{bg}}$ は、
必要環境要因の集合 $E$（\S\ref{subsec:environment_space}）とは無関係である。
$O$ は \S\ref{sec:org_acs} の組織を表す記号 $\mathcal O$ とは無関係である。
不在状態 $\bot_a$ は、\S\ref{sec:tipa} で局所保護応答がないことを表す記号 $\bot$ とは別の記号である。
\end{remark}

%----------------------------------------------------------------------
\subsection{原始的な機能と基盤機能}
\label{subsec:primitive_functions}

本小節では、7タプルによるACSの定義に先立ち、
本論文がACSの機能を記述するときに用いる抽象的な機能区分と、
それを支える基盤機能を定める。
ここでは両者の意味だけを与え、
基礎的な機能、構成要件 CR-2 および型をもつ標準写像形式は
\S\ref{subsec:three_primitive_functions} で与える。

\begin{definition}[原始的な機能（分類A）]
\label{def:primitive_functions}
本論文は、情報を得てから応答を生成するまでのACSの機能的責務を、
次の四つの抽象的な機能区分に分けて記述し、
各区分を\emph{原始的な機能}と呼ぶ。
\begin{description}[style=nextline]
  \item[感知]
  自己の内部状態および外部環境状態から情報を得る機能的責務。
  \item[識別]
  得られた情報を区別し、後続の評価および行動に利用可能な
  内部表現を形成する機能的責務。
  \item[評価]
  内部表現、自己の現在の状態、応答の候補その他の評価の対象を、
  自己の状態に関わる内部基準に照らして非等価に扱う機能的責務。
  評価の対象としてこれらのすべてを扱うことは要求しない。
  \item[行動]
  利用可能な内部表現、評価または現在の状態の少なくともいずれかに応じて
  応答を生成する機能的責務。
  これらのすべてにもとづくことは要求しない。
\end{description}

ここでいう「原始的」とは、本論文の機能記述において
他の概念から定義せず、それ以上の区分に分けない単位として
採用することを指す。
生物学的・進化的な起源や古さ、または物理的な分解不能性を意味しない。

原始的な機能は、個体が有する具体的な機能
（基礎的な機能、定義~\ref{def:basic_functions}）の上位にある
抽象的な区分であり、特定の物理部位、計算モジュール、写像
または実装の分割を指さない。
一つの区分の責務を複数の機能が共同して担うこと、
および一つの機能が複数の区分の責務に関与することを排除しない。
四区分の列挙の順序は情報依存の典型を示すにとどまり、
時刻ごとの逐次実行または固定的な処理系列を意味しない。

本定義における「自己」は、ACSとして記述される対象系自身を指す
非形式的な用法であり、その形式的な対応物は
\S\ref{subsec:self_formal_definition} で定める。
評価における内部基準は、ACSが自ら担う選別の基準であり、
\S\ref{subsec:viability_function_and_blockage} で定義する
持続評価関数 $\Pi$ とは異なる。
本論文は、ACSが $\Pi$ を内部で保持または計算することを要求しない
（\S\ref{subsec:description_levels_global}）。
\end{definition}

\begin{definition}[基盤機能（分類A）]
\label{def:base_functions}
原始的な機能の各区分の責務を担う機能の活動を支え、
その担い手を持続させる機能を\emph{基盤機能}と呼ぶ。
基盤機能には、物質およびエネルギーの供給と輸送、代謝、修復、
基盤状態の調整その他が含まれうる。
基盤機能は原始的な機能の区分には含めず、
原始的な機能と同じ意味での「それ以上分けない単位」としても扱わない。
基盤機能は情報の取得から応答の生成までの責務を担うものではなく、
その責務を担う機能が作動できる条件を維持するものである。
一つの機能が、ある側面では原始的な機能の区分の責務に関与し、
別の側面では基盤機能として働くことを排除しない。
\end{definition}

身体の調整のうち、内部状態を感知し、設定された基準との差を識別・評価し、
内部へ作用して状態を調整する制御の働きは、
原始的な機能の四区分の責務が内部状態に向けて働く場合として扱いうる。
本論文はこの場合のために区分を加えない。
これに対し、物質の輸送や代謝のように機能の担い手を持続させる働きは、
基盤機能として扱う。
一つの過程が両方の働きに関与する場合、本論文はその過程を
いずれか一方に割り当てることを求めず、
同一の過程について、制御の責務を担う側面と
担い手を持続させる側面とを区別して記述しうるものとする
（定義~\ref{def:base_functions}）。
各側面をどの区分の責務または基盤機能として記述するかは
対象系の記述における選択であり、
本論文はその一般的な判定基準を Phase~1 では与えない。

\begin{remark}[原始的な機能と三原理の記述水準の区別]
\label{rem:primitive_vs_principles}
原始的な機能（定義~\ref{def:primitive_functions}）と
基盤機能（定義~\ref{def:base_functions}）は、機能の水準の概念である。
これに対し、\S\ref{sec:three_principles} の
自己維持（SM）・環境利用（EU）・可行応答（FR）の三原理は、
存続中のACSに対する制約軸であり、
本論文は原理そのものについて成立・不成立を述べない。
真偽をもつのは、同節で定める各原理の時刻別充足述語である。
本論文は、原始的な機能のいずれかの区分を特定の原理に一対一に対応させず、
原理ごとの専用の機能も置かない。
また本論文は、原始的な機能の各区分に関して個体が有する基礎的な機能
（定義~\ref{def:basic_functions}）または基盤機能が作動していることから
時刻別充足述語が真であることを、
時刻別充足述語が偽であることから特定の基礎的な機能または基盤機能の停止を、
いずれも一般的な帰結として置かない。
原始的な機能と、\S\ref{sec:three_principles} の三原理に関連する構造的機能
（定義~\ref{def:principle_related_functions}）との対応は、同節で扱う。
\end{remark}

%----------------------------------------------------------------------
\subsection{ACS 7タプルと基本型}
\label{subsec:acs_7tuple}

本小節では、外部環境状態空間 $E_{\mathrm{ext}}$ に相対的な
ACS個体の基礎形式を7タプルとして定義する。
7タプルとは、ACSを記述する7個の基本成分を
順序つきの組として並べたものである。
本小節では各成分の基本型を定め、
構造的自己の意味は \S\ref{subsec:self_formal_definition}、
必要環境空間 $E$ の存続必要性の意味論は
\S\ref{subsec:environment_space}、
原始的な機能に関する構成要件は
\S\ref{subsec:three_primitive_functions} でそれぞれ与える。
原始的な機能そのものは \S\ref{subsec:primitive_functions} で定義した。

\begin{definition}[自律認知システム]
\label{def:acs}
外部環境状態空間 $E_{\mathrm{ext}}$ に相対的に、
型分離規約 TSC（\S\ref{subsec:wellformedness}）のもとで、
構成要件 CR-1（必要環境整合性、\S\ref{subsec:environment_space}）および
構成要件 CR-2（原始的な機能と基盤機能の具備、\S\ref{subsec:three_primitive_functions}）を
満たす実現を備える以下の7タプル構造体を、\emph{自律認知システム}
（ACS: Autonomous Cognitive System）と定義する：
\[
  \Sigma
  =
  \bigl(
    X,\,
    U,\,
    E,\,
    f,\,
    \Pi,\,
    M,\,
    \mathrm{Update}
  \bigr).
\]

\medskip
\noindent\textbf{(1) 状態空間 $X$：}
$X$ はACSがとりうる内部状態の空間であり、
時刻 $t$ の内部状態を $x_t\in X$ と書く。
Phase~1 では $X\subseteq\mathbb{R}^n$ とし、
$X$ をコンパクト集合として扱う
（仮定~\ref{ass:compact_spaces}）。
$X$ はACSを構成する物質そのものではなく、
当該ACSが同一個体として実現されている間に取りうる
内部状態を形式的に表す状態空間である。

\medskip
\noindent\textbf{(2) 行動空間 $U$：}
$U$ はACSが選択しうる実効行動の空間であり、
時刻 $t$ の実効行動を $u_t\in U$ と書く。
Phase~1 では $U\subseteq\mathbb{R}^m$ とし、
$U$ をコンパクト集合として扱う
（仮定~\ref{ass:compact_spaces}）。
$U$ の内部には、外部環境へ作用する行動、内的認知状態に
作用する行動、および選択的待機を区別するラベル構造を
導入することがある（備考~\ref{rem:action_space_structure}）。
後続の機構記述において複数の内部指令源または局所応答が
区別される場合でも、状態遷移写像 $f$ に入力される $u_t$ は、
それらが当該実装の規則に従って解決された後の実効行動を表す。

\medskip
\noindent\textbf{(3) 必要環境空間 $E$：}
$E$ は、当該ACSが同一個体として存続するために不可欠な
外部環境要因を表す集合である。
ここでいう存続は、数学的対象としての7タプルが存在することではなく、
同一ACS個体としての整形式な実現が継続することを指す
（\S\ref{subsec:self_formal_definition}）。
Phase~1 では各必要環境要因を、内部状態 $x\in X$ と
外部環境状態 $e\in E_{\mathrm{ext}}$ の組に対する
二値の充足判定写像として表し、
\[
  E
  \subseteq
  \{0,1\}^{X\times E_{\mathrm{ext}}},
  \qquad
  E\neq\emptyset
\]
とする。
したがって任意の $\eta\in E$ について
\[
  \eta:
  X\times E_{\mathrm{ext}}
  \longrightarrow
  \{0,1\}
\]
である。
$\eta(x,e)=1$ は当該内部状態と外部環境状態の組において
その必要環境要因が充足されていることを、
$\eta(x,e)=0$ は非充足であることを表す。
Phase~1 では $\eta$ に連続性その他の正則性を課さない。
これを要する命題ではその箇所で追加仮定として明示する。
$E$ の元は外部環境状態そのものではないため、
実現外部環境状態は $e_t\in E$ ではなく
$e_t\in E_{\mathrm{ext}}$ と書く。
$E$ の各元がなぜ「必要」と呼ばれるかを拘束する構成要件 CR-1、
必要要因の外部依存性と共同充足、および要因分解の非一意性は
\S\ref{subsec:environment_space} で定める。

\medskip
\noindent\textbf{(4) 状態遷移写像 $f$：}
状態遷移写像は、内部状態、実効行動および実現外部環境状態から
次の内部状態を与える写像
\[
  f:
  X\times U\times E_{\mathrm{ext}}
  \longrightarrow
  X
\]
である。
したがって状態遷移は
\[
  x_{t+1}
  =
  f(x_t,u_t,e_t),
  \qquad
  x_t\in X,\quad
  u_t\in U,\quad
  e_t\in E_{\mathrm{ext}}
\]
と表される。
Phase~1 では $f$ を連続写像として扱う
（仮定~\ref{ass:f_continuity}）。
7タプル成分 $f$ は、必要環境空間上への制限としてではなく、
外部環境状態空間 $E_{\mathrm{ext}}$ 全体の上で定義する。
これにより、存続に必須ではないが状態遷移へ影響する外部入力、
中立的外部入力、および有害な外部入力も同一の型の下で記述できる。
$f$ の自律的機能過程と指向的過程への分解は
後続小節で扱う（仮定~\ref{ass:f_decomposition}）。

\medskip
\noindent\textbf{(5) 持続評価関数 $\Pi$：}
$\Pi$ はACSの内部状態の可行性を評価する実数値関数であり、
その型は
\[
  \Pi:X\longrightarrow\mathbb{R}
\]
である。
本論文では基本的に $\Pi(x_t)\geq0$ であるとき、
状態 $x_t$ は可行であると解釈する。
可行集合 $K$ は7タプルの独立成分ではなく、
\[
  K
  =
  \{\,x\in X\mid \Pi(x)\geq0\,\}
\]
として $\Pi$ から導出される
（定義~\ref{def:K}）。

\medskip
\noindent\textbf{(6) 内部モデル構造 $M$：}
$M$ はACSが外部環境、自己状態および行動可能性に関する情報を処理し、
応答形成に利用する内部モデル構造である。
本論文では
\[
  M=(M_1,M_2,M_3)
\]
とし、$M_1$ を情報統合機構、$M_2$ を適応的評価機構、
$M_3$ を応答行動機構とする三機構情報処理アーキテクチャ
（TIPA: Triadic Information-Processing Architecture）として扱う。
内部モデル状態の空間を $\mathcal{M}$ と書き、
時刻 $t$ の内部モデル状態を $M_t\in\mathcal{M}$ とする。
$M$ は機構構造を表す記号、$\mathcal{M}$ はその状態がとる空間、
$M_t$ は時刻 $t$ における内部モデル状態であり、これらを区別する。
内部モデルの最小構造は \S\ref{subsec:internal_model_M}、
TIPAの標準機能形式の詳細は \S\ref{sec:tipa} で扱う。

\medskip
\noindent\textbf{(7) 内部モデル更新作用素 $\mathrm{Update}$：}
$\mathrm{Update}$ は、実現内部状態および実現外部環境状態に応じて
内部モデル状態 $M_t\in\mathcal{M}$ を更新する作用素であり、
7タプルにおける包括的型は
\[
  \mathrm{Update}:
  \mathcal{M}\times X\times E_{\mathrm{ext}}
  \longrightarrow
  \mathcal{M}
\]
である。
構成要素別の詳細な更新型は
\S\ref{subsec:update_operator}（定義~\ref{def:update_operator}）で与え、
抽象型と詳細型の関係は備考~\ref{rem:update_type_two_layers} で整理する。

\medskip
以上の七成分により、ACSは、内部状態可能性、行動可能性、
外部世界に対する不可欠な存続依存性、状態遷移、可行性評価、
内部モデル構造および内部モデル更新を明示的に区別して記述される。
複数のACS個体を区別する必要がある場合には、個体添字を用いて
\[
  \Sigma_i
  =
  \bigl(
    X_i,\,
    U_i,\,
    E_i,\,
    f_i,\,
    \Pi_i,\,
    M_i,\,
    \mathrm{Update}_i
  \bigr)
\]
と表記し、その個体に相対的な外部環境状態空間を
$E_{\mathrm{ext},i}$ とする。
ただし本節では単一ACSの形式的定義を主対象とするため、
成分添字および $E_{\mathrm{ext}}$ の個体添字は
必要な場合を除き省略する。
\end{definition}

\subsubsection*{記述水準の規約（抽象構造型・個体・時刻別の値）}
\label{conv:structure_levels}

本論文では、ACSに関わる記号を次の三つの水準で書き分ける。
本規約は記述上の取り決めであり、ACS自身の機能的能力でも、
後続命題の成立範囲を限定する仮定でもない。

\begin{description}[style=nextline]
  \item[抽象構造型]
  定義~\ref{def:acs} が定める7タプル構造の型である。
  抽象構造型を集合として扱うことは要求しない。
  したがって個体 $\Sigma_i$ を抽象構造型の元とは書かず、
  「$\Sigma_i$ は抽象構造型の実現である」と書く。
  \item[個体 $\Sigma_i$]
  抽象構造型の一つの実現として識別されるACS個体の7タプル
  $\Sigma_i=(X_i,U_i,E_i,f_i,\Pi_i,M_i,\mathrm{Update}_i)$ である。
  Phase~1 では、個体の7タプル $\Sigma_i$ を時刻に依存しない形式的構造体として扱う
  （定義~\ref{def:s_struct}）。
  時刻とともに変わりうるのは、状態 $x_{i,t}\in X_i$、内部モデル状態その他の
  時刻別の値である。
  とくに内部モデル構造 $M_i$ の機構構成は固定し、
  各機構の状態 $M_{c,t}$、機能的記憶内容その他の時刻別の値の変化を、
  内部モデル更新および \S\ref{sec:extended_dynamics} の時刻間変化として記述する。
  機構構成そのものの時間変化は扱わない。
  \item[時刻別の値]
  個体の実現（状態列・環境列）に依存する量は、
  構造を入力とする関数形にせず、個体添字と時刻添字を下付きに置く。
  時刻に依存しない写像または対応と、その時刻評価値とを二段で書く。
  例えば可行行動集合については、対応 $A_i(x,e)$ と時刻評価値
  \[
    A_{i,t}:=A_i(x_{i,t},e_{i,t})
  \]
  を区別する（定義~\ref{def:at}）。
  時刻評価値の記号は、本文で定義したものに限って用いる。
\end{description}

内部モデルについては、機構構造 $M$（7タプルの成分）、
その状態空間 $\mathcal M$、時刻 $t$ の内部モデル状態 $M_t\in\mathcal M$ を区別する
（\S\ref{subsec:internal_model_M}）。
TIPA の第 $c$ 機構 $M_c$ と、その時刻 $t$ の状態 $M_{c,t}\in\mathcal M_c$ も、
時刻添字の有無によって型が異なる。
単一ACSのみを扱う文脈では個体添字を省略する。
個体の7タプルそのものが時刻とともに変化する場合の記述は
Phase~2 の課題とする（\S\ref{subsec:limitations_list} の L16）。

\subsubsection{基本型と各成分に関する補足}

\begin{remark}[可行集合 $K$ の位置づけ]
\label{rem:K_position}
可行集合 $K$ は7タプルの独立成分として保持されない。
$K$ は持続評価関数 $\Pi$ と本論文の可行性判定規約から
$K = \{\,x\in X\mid \Pi(x)\geq0\,\}$
として導出される派生的対象である。
$\Pi$ の具体的構成、$K$ の形式的定義、およびこれらの関係は
\S\ref{subsec:viability_function_and_blockage} で与える。
したがって $K$ はACS 7タプルの独立成分ではない。
\end{remark}

\begin{remark}[行動空間 $U$ の構造]
\label{rem:action_space_structure}
$U$ はコンパクト集合として与えられる
（仮定~\ref{ass:compact_spaces}）。
本論文では、$U$ の内部に
$U_{\mathrm{ext}}\subseteq U$（外部環境へ直接作用する行動）、
$U_{\mathrm{int}}\subseteq U$（内的認知状態を意図的に変化させる行動）、
$u_{\emptyset}\in U$（選択的待機）というラベル構造を識別する。
$U_{\mathrm{ext}}$ と $U_{\mathrm{int}}$ は $U$ の部分集合、
$u_{\emptyset}$ は $U$ の要素として識別されるラベル構造であり、
$U$ のコンパクト性を変更しない。
$u_{\emptyset}$ は行動の不在を表す記号ではなく、$U$ の正規の元であり、
他の元と同じく状態遷移写像 $f$ の行動引数となる。
位相構造への影響は Phase~2 の課題とする。
$u_{\emptyset}$ の選択時における自律的機能過程との関係は
後続の状態遷移分解で定める
（仮定~\ref{ass:f_decomposition}）。
行動が行動応答成分へ及ぼす作用は、ラベル $U_{\mathrm{ext}}$ によってではなく、
作用差（定義~\ref{def:action_difference}）によって記述する。
外部による強制は当該ACSの行動ではなく、
外部環境状態 $e_t$ を通じて $f$ に入る状態変化として表す。
同じ時刻に、当該ACSの行動 $u_t$ と外部による強制は共存しうる。
行動の区分は \S\ref{sec:tipa} の定義~\ref{def:action_classes} で、
目的の基準が当該ACSの内部の担体から読み出されることは
同節の命題~\ref{prop:purpose_self_attribution} で扱う。
\end{remark}

\begin{remark}[Updateの型宣言の二層構造]
\label{rem:update_type_two_layers}
7タプルにおける $\mathrm{Update}$ の型宣言
\[
  \mathcal{M}\times X\times E_{\mathrm{ext}}
  \longrightarrow
  \mathcal{M}
\]
は包括的な抽象型である。
構成要素別の詳細型
（定義~\ref{def:update_operator}、
\S\ref{subsec:update_operator}）は、
実装上利用可能な追加引数を含む写像族として定義されうる。
したがって抽象型と詳細型は直接同一の型である必要はなく、
抽象化レベルの差として整合的に扱う（under stated assumptions）。
詳細型に現れる追加引数は、同一ステップまたは先行処理により得られる
中間値・補助情報であり、それだけを理由として7タプルの独立成分を
追加しない。
\end{remark}

\subsubsection{7タプルの非冗長性}

\begin{remark}[7タプルの非冗長性]
\label{rem:tuple_nonredundancy}
本論文は、7タプル
$\Sigma=(X,U,E,f,\Pi,M,\mathrm{Update})$ を
Phase~1 のACS定義における基礎成分として採用する。
これらは、本論文が採用する表現体系の下で、それぞれ
内部状態可能性、行動可能性、外部世界に対する存続必要性、
状態遷移、可行性評価、内部モデル構造、内部モデル更新という
異なる構造的役割を担う。

特に、必要環境空間 $E$ は実現外部環境状態の型
$E_{\mathrm{ext}}$ そのものでも、その部分集合でもなく、
当該ACSが外部世界に対してもつ存続上の不可欠な依存性を
必要環境要因の集合として表す原始的構成情報である。
Phase~1 では $E$ を $X,f,\Pi$ その他の成分から
一意に導出される派生量として定義しない。
本論文は、$f$ が記述する環境差のうち何が同一ACSの存続に
構成的に不可欠であるかを $f$ から一意に定める規則を与えず、
また $\Pi$ から必要環境要因の分解粒度を一意に定める規則も与えない。
必要環境性の意味論的拘束と分解の非一意性は
\S\ref{subsec:environment_space} で扱う。

したがって、内部状態、行動、存続必要環境、状態遷移、可行性評価、
内部モデル、内部モデル更新という区別を保存する本論文の表現要件の下では、
いずれかの成分を削除すると、少なくとも一つの構造的役割が
明示的には保持されなくなるという意味で七成分を非冗長とみなす。

ただし本論文は、任意の数学的再符号化に対する絶対的最小性を主張しない。
例えば、拡張状態空間 $Y:=X\times\mathcal{M}$ と拡張遷移写像を用いて、
内部モデル状態と更新作用素を状態遷移側へ吸収する再符号化は
形式的には可能である。
その場合、本論文が明示的に保存する内部状態と内部モデル状態の区別、
および状態遷移と内部モデル更新の区別は失われる。
本備考が述べるのは、これらの役割分離を保存する
現在の形式体系の範囲内での非冗長性である
（within present framework）。
\end{remark}

%----------------------------------------------------------------------
\subsection{構造的自己と時間的自己同一性}
\label{subsec:self_formal_definition}

本小節では、ACS個体の「自己」を形式的に定義する。
ここでいう自己とは、主観的な意識や記憶の同一性ではなく、
ACSが同一の個体として同定される構造的な基盤を指す。
本論文では自己を三つの水準に分けて整理する。
構造的自己（個体そのもの）、動的自己状態（時刻ごとの状態）、
自己--環境境界（自己と環境の区別）である。
あわせて、これら三水準とは別に、
当該個体が各時刻において整形式に実現されているか否かを表す
実現水準を導入する。
この実現水準は、必要環境空間 $E$ の各要因がなぜ「必要」と呼ばれるかを
\S\ref{subsec:environment_space} で拘束する際に用いる語彙を与える。
本小節の末尾では、これとは別の対象として、
原始的な機能の各区分に関して構造的個体が有する基礎的な機能の活動と、
その活動を支える基盤機能にもとづく機能的終端と、
構造的個体としての同一性を成立させる組織関係にもとづく崩壊を定義する。
また、対象となる情報を自己側と環境側のいずれに帰属するものとして
識別したかを表す内部出力を、
\S\ref{sec:tipa} の識別成分として参照する。

本小節で新たに用いる記号は、構造的自己
$S_{\mathrm{struct}}(\Sigma_i)$、動的自己状態
$S_{\mathrm{dyn},t}$、自己--環境境界 $B_{SE}$、
整形式実現指標 $\mathrm{WF}_{i,t}$、
機能的終端判定時刻域 $\mathbb{T}_{\mathrm{FT},i}$、
機能的終端状態指標 $\mathrm{FTS}_{i,t}$、機能的終端指標 $\mathrm{FT}_{i,t}$、
崩壊状態指標 $\mathrm{CS}_{i,t}$、崩壊指標 $\mathrm{Col}_{i,t}$
および許容外的作用集合 $\Upsilon_{\mathrm{adm},i}$ である。
下付き $\mathrm{struct}$ は structural（構造的）、
$\mathrm{dyn}$ は dynamic（動的）に対応し、
$B_{SE}$ の $SE$ は self--environment（自己--環境）、
$\mathrm{WF}$ は well-formed（整形式）を表す。
あわせて、認知的な境界帰属識別の結果を表す成分として
\S\ref{sec:tipa} で定義される $\xi_{\mathrm{bnd},t}$ を参照する。
下付き $\mathrm{bnd}$ は boundary（境界）に対応する。

\subsubsection{構造的自己と整形式実現}

\begin{definition}[構造的自己]
\label{def:s_struct}
ACS個体 $\Sigma_i$ の\emph{構造的自己}を
\[
  S_{\mathrm{struct}}(\Sigma_i) := \Sigma_i
\]
として定義する。
ここで $\Sigma_i$ は定義~\ref{def:acs} で定義された7タプル
$\Sigma_i = ( X_i , U_i, E_i, f_i, \Pi_i, M_i, \mathrm{Update}_i)$
であり、構造的自己 $S_{\mathrm{struct}}(\Sigma_i)$ は
時刻に依存しない形式的構造体である。
個体添字 $i$ は対象とするACS個体を区別するための識別子であって、
冪を表さない。

構造的自己は、ACS個体 $\Sigma_i$ が同一の個体として
同定される基盤を与える。
ただし、数学的対象としての7タプルが存在することと、
その7タプルが記述する構造が各時刻において整形式に実現され、
かつ時刻をまたいで同一のACS個体として継続していることとは
区別される。
各時刻の実現を表す指標は定義~\ref{def:wf_realization} で導入し、
時刻間の同一個体性は定義~\ref{def:time_identity} で整理する。
両者と数学的対象としての存在との関係は
備考~\ref{rem:math_vs_realization} で述べる。
個体添字 $i$ を省略した $\Sigma$ および
$S_{\mathrm{struct}}(\Sigma)$ は、文脈上ただ一つの
ACS個体のみを扱う場合の略記であり、個体添字付きの
表記と同一の対象を指す。
\end{definition}

\begin{definition}[整形式実現指標（分類C）]
\label{def:wf_realization}
$\Sigma_i$ が定義~\ref{def:acs} の意味でACSであるとき、
各離散時刻 $t$ に対し、
$\Sigma_i$ に対応する構造が時刻 $t$ において
整形式（well-formed）に実現されているか否かを表す指標
\[
  \mathrm{WF}_{i,t}\in\{0,1\}
\]
を導入する。

$\mathrm{WF}_{i,t}=1$ は、時刻 $t$ において、
$\Sigma_i$ の七成分が当該実現においてその型のとおりに
与えられていることを表す。
$\mathrm{WF}_{i,t}=0$ はこれが成立しないことを表す。

$\mathrm{WF}_{i,t}$ は各時刻について個別に定まる指標であり、
異なる時刻の実現が同一のACS個体に属するか否かを
それ自体では述べない。
時刻間の同一個体性は、$\mathrm{WF}$ とは別に
定義~\ref{def:time_identity} の予備条件によって整理する。

$\mathrm{WF}_{i,t}$ は7タプルの成分ではなく、
実現の成否を記述するための時刻依存の指標である。
Phase~1 では $\mathrm{WF}_{i,t}$ の値を有限観測から
判定する操作的手続きを与えず、その定式化は Phase~2 の課題とする。
\end{definition}

\begin{remark}[数学的水準と実現水準の分離]
\label{rem:math_vs_realization}
定義~\ref{def:s_struct} の
$S_{\mathrm{struct}}(\Sigma_i)=\Sigma_i$ は
時刻に依存しない数学的対象であり、
その存在は $\mathrm{WF}_{i,t}$ の値に依存しない。
$\mathrm{WF}_{i,t}=0$ となる場合でも、
7タプル $\Sigma_i$ が数学的対象として消滅するのではない。
失われるのは、その7タプルが記述する構造が同一のACS個体として
実現され続けている状態である。

本論文において「存続」というとき、それは数学的対象としての
$\Sigma_i$ が存在し続けることではなく、
各時刻において $\mathrm{WF}_{i,t}=1$ が保たれ、
かつ時間的自己同一性の予備条件（定義~\ref{def:time_identity}）が
成立し続けることを指す。

この区別は、構成要件と実現指標の役割の違いに対応する。
構成要件 CR-1（\S\ref{subsec:environment_space}）および
CR-2（\S\ref{subsec:three_primitive_functions}）は、
個々の時刻における診断値ではなく、
当該7タプル構造体を本論文におけるACSのクラスへ
属させるための条件である。
これに対し $\mathrm{WF}_{i,t}$ は、
そのようにして定まったACS個体の実現が
時刻 $t$ において成立しているか否かを記録する時刻依存の指標である。

したがって、必要環境空間 $E$ の各要因が
「当該ACSの存続に必要である」という CR-1 の要求は、
$E$ を含む数学的構造 $\Sigma_i$ の存在についての条件ではなく、
$\mathrm{WF}_{i,t}$ の推移についての条件として述べられる。
この水準分離により、$E$ の意味論的拘束が
$E$ 自身を含む対象の存在を前提とするという循環は生じない。
\end{remark}

\subsubsection{動的自己と自己--環境境界}

\begin{definition}[動的自己状態]
\label{def:s_dyn}
ACS個体 $\Sigma_i$ の時刻 $t$ における
\emph{動的自己状態}を
\[
  S_{\mathrm{dyn},t} := (x_t,\ A_t)
\]
として定義する。
$x_t\in X$ は時刻 $t$ の状態ベクトル、
$A_t$ は可行行動集合（定義~\ref{def:at}）である。
$A_t = \emptyset$ は $S_{\mathrm{dyn},t}$ の消滅ではなく、
可行行動が存在しないこと（閉塞、定義~\ref{def:blockage_sec2}）を
意味する（within present framework）。
\end{definition}

\begin{definition}[自己--環境境界]
\label{def:b_se}
規約 TSC のもとで、自己--環境境界を表す構造的分離
$B_{SE}$ を
\[
  B_{SE}:=(X,E_{\mathrm{ext}}),
  \qquad
  X\cap E_{\mathrm{ext}}=\emptyset
\]
と定義する。
実現状態については、任意の時刻 $t$ において
\[
  x_t\in X,
  \qquad
  e_t\in E_{\mathrm{ext}}
\]
とする。
$B_{SE}$ は Phase~1 では固定された構造的分離であり、
内部状態の型 $X$ と外部環境状態の型 $E_{\mathrm{ext}}$ を
集合表現上で区別する。
この分離は、物理的・空間的境界または自己と非自己の
存在論的に完備した分割を主張するものではない。

必要環境空間 $E$（\S\ref{subsec:environment_space}）は
$B_{SE}$ が表す型分離とは別の対象である。
$B_{SE}$ は内部状態型と外部環境状態型の区別を表すのに対し、
$E$ は外部環境状態と内部状態の組に対する充足判定写像の集合として、
当該ACSが外部世界に対してもつ存続上の依存性を表す。
したがって $E$ は $E_{\mathrm{ext}}$ の部分集合ではなく、
$B_{SE}$ の第二成分と同一視されない。
\end{definition}

\begin{remark}[境界帰属識別の参照（分類B）]
\label{rem:xi_bnd_reference}
自己と環境の区別に関する認知的な識別結果は、
\S\ref{sec:tipa} において M1-b の識別状態
$\xi_t$ の全体の行の境界帰属成分
\[
  \xi_{\mathrm{bnd},t}
  \in
  \Xi_{\mathrm{bnd}}
  =
  \{\mathsf{self},\mathsf{env},\mathsf{und}\}
\]
として定義される。
本節ではこれを再定義せず、参照するにとどめる。
値 $\mathsf{self}$ は対象となる情報が自己側に帰属すると
識別されたこと、$\mathsf{env}$ は環境側に帰属すると
識別されたこと、$\mathsf{und}$ はいずれとも確定されなかったことを表す。

$\xi_{\mathrm{bnd},t}$ の値は識別の内容を表し、
識別機能が正しく作動していることを表さない。
M1-b の作動水準および実装部位の健全性は、
\S\ref{sec:tipa} の機能軸および
\S\ref{sec:extended_dynamics} の健全度指標により
別途記述される。
本節はそれらを再定義しない。
\end{remark}

\begin{proposition}[$\xi_{\mathrm{bnd},t}$ と $S_{\mathrm{struct}}$ の独立性]
\label{prop:xi_se_independence}
規約 TSC のもとで、構造的自己
$S_{\mathrm{struct}}(\Sigma_i)$ の存在は
$\xi_{\mathrm{bnd},t}$ の値に依存しない
（within present framework）。
\end{proposition}

\begin{proof}
$S_{\mathrm{struct}}(\Sigma_i) = \Sigma_i$ は7タプルとして
時刻に依存しない数学的対象である。
$\xi_{\mathrm{bnd},t}$ は M1-b の時刻 $t$ における出力であり、
$S_{\mathrm{struct}}(\Sigma_i)$ の定義に登場しない。
よって $\xi_{\mathrm{bnd},t}$ のいずれの値も
$S_{\mathrm{struct}}(\Sigma_i)$ の存在への論理的含意を持たない。
\end{proof}

\begin{remark}[認知的識別の失敗と構造の独立性]
\label{rem:xi_se_interpretation}
M1-b の識別機能が正しく作動しない場合、
ACSは自己と環境の識別を信頼可能に利用できない状態を生じうる。
本節では、この状態が可行性維持、閉塞または整形式実現に対して
及ぼす影響を一般的な帰結として置かず、本論文はその記述を Phase~1 では与えない。
ただし、$K$ と可行行動集合は内部モデルの状態によらず定まるので
（定義~\ref{def:at}）、識別の結果が状態の可行性に関わるのは、
実効行動の選択（\S\ref{sec:tipa}）を通じた範囲に限られる。

M1-b 自体が外部から書き換えられた場合、
出力値が $\mathsf{self}$、$\mathsf{env}$ または $\mathsf{und}$ の
いずれであるかにかかわらず、
構造的分離 $B_{SE}$ に対する誤分類を生じる状態が理論上生じうる。
この構造的分離 $B_{SE}$ と M1-b による自己・環境分類結果との
不整合は、Phase~1 では操作的判定基準が未定義の
中間状態として記録し、M 整合性検証機構の導入を
Phase~2 の課題とする。

なお、命題~\ref{prop:xi_se_independence} が述べるのは
$S_{\mathrm{struct}}(\Sigma_i)$ の存在についての独立性であり、
$\mathrm{WF}_{i,t}$ の値が
$\xi_{\mathrm{bnd},t}$ から独立であることを述べるものではない。
\end{remark}

\subsubsection{後続実現の非分岐性}

\begin{assumption}[個体同一性関係の非分岐性]
\label{ass:no_branching}
ACS個体 $\Sigma_i$ の時間的自己同一性を追跡する際、
時刻 $t$ におけるACS個体実現に対して、
同一個体の継続として対応づけられる時刻 $t+1$ の
ACS個体実現は高々一つである。

本仮定は、状態遷移写像
\[
  f:X\times U\times E_{\mathrm{ext}}\to X
\]
の一価性とは別の仮定である。
$f$ の一価性は、固定された入力
$(x_t,u_t,e_t)$ に対する次状態値 $x_{t+1}$ を
一意に定めるが、同一の次状態値をもつ複数の個体実現、
またはコピー個体の生成を排除しない。

\medskip
\noindent
導入地位：A（公理的）。強さ：強。
これは、同一ACS個体の継続として対応づけられる
後続の個体実現が複数に分岐する場合を、
Phase~1 の対象から除外する仮定である。
本仮定はコピー個体の生成可能性そのものを数学的に
否定するものではなく、コピーまたは複数の後続実現が
生じた場合の個体同一性を Phase~1 では定義しないことを表す。

同一個体性の分岐を許す場合には、時間的自己同一性関係、
構造的自己、および後続個体の対応関係を再定式化する必要がある。
その定式化は
\S\ref{sec:limitations_and_outlook} の
未解決問題として扱う。
\end{assumption}

\subsubsection{時間的自己同一性の予備条件}

\begin{definition}[時間的自己同一性の予備条件]
\label{def:time_identity}
本定義は、後掲の崩壊定義（定義~\ref{def:collapse}）、
および \S\ref{sec:extended_dynamics} における時間地平
$\widehat{\tau}_{t}$ の導入を先取りする予備的整理である。
本節では、時間的自己同一性の完全な存在・一意性・
到達性を主張しない。

ACS個体 $\Sigma_i$ の時刻 $t$ と時刻 $t+1$ の
動的自己状態が\emph{同一のACS個体に属する}ための
予備条件を、以下の三条件が同時に成立することとして
整理する（under stated assumptions）。

\medskip
\noindent\textbf{(TI-1) 構造的連続性の必要条件：}
構造的自己を担う個体実現が、時刻 $t$ から $t+1$ にかけて
同一個体として継続するための必要条件として、
\[
  \mathrm{WF}_{i,t}=1
  \quad\text{かつ}\quad
  \mathrm{WF}_{i,t+1}=1
\]
が成立することを要請する。

本条件は、両時刻において整形式な実現が成立していることを表す。
同一個体としての継続そのものは $\mathrm{WF}$ が単独で表すものではなく、
(TI-1) から (TI-3) を合わせて予備条件として整理される。
機能的終端・崩壊・閉塞との区別は備考~\ref{rem:time_identity_collapse} で述べる。

\medskip
\noindent\textbf{(TI-2) 動的自己状態の遷移整合性：}
時刻 $t$ の動的自己状態
\[
  S_{\mathrm{dyn},t}=(x_t,A_t)
\]
から時刻 $t+1$ の動的自己状態
\[
  S_{\mathrm{dyn},t+1}=(x_{t+1},A_{t+1})
\]
への継続が、同一のACS動学に従って定まることを要請する。
具体的には、
\[
  x_{t+1}=f(x_t,u_t,e_t),
  \qquad
  x_t\in X,\quad
  u_t\in U,\quad
  e_t\in E_{\mathrm{ext}}
\]
によって次状態 $x_{t+1}$ が生成され、さらに時刻 $t+1$ の
実現外部環境状態 $e_{t+1}\in E_{\mathrm{ext}}$ のもとで
\[
  A_{t+1}
  =
  \left\{
    u\in U
    \;\middle|\;
    f(x_{t+1},u,e_{t+1})\in K
  \right\}
\]
として可行行動集合 $A_{t+1}$ が定まることとする。

したがって、
\[
  (x_t,A_t)
  \longrightarrow
  (x_{t+1},A_{t+1})
\]
は、状態遷移写像 $f$ および可行行動集合の定義
（定義~\ref{def:state_transition}、定義~\ref{def:at}）
と整合する同一の状態遷移系上で記述される。

本条件は、$A_{t+1}$ が非空であることを要求しない。
したがって
\[
  A_{t+1}=\emptyset
\]
であっても、これは時刻 $t+1$ における閉塞
（定義~\ref{def:blockage_sec2}）を意味するにとどまる。
本論文は、これから動的自己状態
\[
  S_{\mathrm{dyn},t+1}
  =
  (x_{t+1},\emptyset)
\]
そのものの消滅、または
$\mathrm{WF}_{i,t+1}=0$
を一般的な帰結として置かない。

\medskip
\noindent\textbf{(TI-3) 個体同一性関係の非分岐性：}
時刻 $t$ のACS個体実現に対して、同一個体の継続として
対応づけられる時刻 $t+1$ のACS個体実現が高々一つであること
（仮定~\ref{ass:no_branching}）。
本条件は、状態遷移写像 $f$ の決定論性とは別に、
同一ACS個体から複数の後続個体が生じる場合を
Phase~1 の対象から除外する個体同一性上の仮定である。

\medskip
コピー自己を生じる分岐は、状態遷移写像 $f$ の一価性によってではなく、
仮定~\ref{ass:no_branching} により Phase~1 の対象から明示的に除外する。
非分岐性の緩和、および分岐を許す場合の時間的自己同一性の
再定式化は、\S\ref{sec:limitations_and_outlook} の
未解決問題~\ref{prob:branching_identity} として扱う。
本論文において自己は $S_{\mathrm{struct}}(\Sigma_i)$
として定義され、認識内容（$M_1$ の記憶内容）は
構造的存続の判定に影響しない。
主観的自己・意識・記憶の同一性は Phase~1 の記述対象に
含まれない（under stated assumptions）。
\end{definition}

\subsubsection{機能的終端と崩壊}
\label{subsubsec:collapse_sec2}

本項では、ACS個体の終端を二段に分けて定義する。
第一は\emph{機能的終端}であり、
原始的な機能（感知・識別・評価・行動）の各区分に関して構造的個体が有する
基礎的な機能の活動と、その活動を支える基盤機能のすべてが停止し、
そのいずれも同一個体として再び成立しない状態をいう。
第二は\emph{崩壊}（構造崩壊）であり、
構造的個体としての同一性を成立させていた組織関係が、
同一個体として回復しない形で失われた状態をいう。
本論文で単に「崩壊」というときは後者を指す。
いずれも、終端へ至る原因や到達の条件ではなく、
ある時刻の個体がすでにその状態にあると判定するための状態条件である。
本節はこれらを定義するが、
機能的終端または崩壊への到達、その十分条件または到達時刻を主張しない。
閉塞および記述的状態帯との関係、外的回復可能性、
および到達経路の整理は \S\ref{sec:collapse_cond} で扱う。
定義の順序は、機能的終端状態を判定する時刻域が崩壊状態指標を用いて定まるため、
崩壊を先に置く。
これは二つの状態が成立する時間的な順序を表すものではない。

機能的終端と崩壊に含まれる非回復性は、
考えうるあらゆる物理作用に対する不可能性ではなく、
本枠組みが考慮対象とする外的作用の範囲に相対して定める。
そのため、各ACS個体 $\Sigma_i$ に対し、
外的回復可能性の評価で考慮対象とする外的作用の集合を
$\Upsilon_{\mathrm{adm},i}$ と書き、その元を
$\upsilon\in\Upsilon_{\mathrm{adm},i}$ と書く。
本節では記号と型のみを宣言する。
$\Upsilon_{\mathrm{adm},i}$ は集合であり、本節ではその位相・測度・代数構造を定めず、
非空であることも仮定しない。
その内容は \S\ref{sec:collapse_cond}（定義~\ref{def:admissible_external_action}）で定める。
$\upsilon$ はACS自身が行動空間から選択する $u\in U$ とは役割の異なる変数であり、
状態遷移写像 $f$ の引数へ暗黙に代入しない。
(SC-2) および (FT-2) に現れる行動応答成分の遷移写像 $g_O$ の型は、定義~\ref{def:env_realization} で定める。
記号 $\Upsilon_{\mathrm{adm},i}$ は、\S\ref{sec:tipa} の M1-b の識別状態の値域 $\Xi$ とは
無関係である。

\begin{definition}[崩壊状態指標と崩壊指標]
\label{def:collapse}
ACS個体 $\Sigma_i$ と時刻 $t\in\mathbb{N}_0$ について、
次の (SC-1) および (SC-2) がともに成立するとき
\[
  \mathrm{CS}_{i,t}=1,
\]
それ以外のとき $\mathrm{CS}_{i,t}=0$ と定め、
$\mathrm{CS}_{i,t}\in\{0,1\}$ を\emph{崩壊状態指標}と呼ぶ。

\begin{description}[style=nextline]
  \item[(SC-1) 構造的同一性を成立させる組織関係の喪失]
  時刻 $t$ において、$\Sigma_i$ を構成していた構成要素の間にあって、
  構造的個体としての $\Sigma_i$ の同一性を成立させていた
  組織・結合関係が失われている。
  構成要素そのものが残存しているか否かは問わない。
  \item[(SC-2) 本枠組みに相対する同一個体としての非回復性]
  時刻 $t$ 以後について、実現した外生背景成分の列 $(b_s)_{s\geq t}$ のもとでの
  $\Sigma_i$ 自身の通常の自律的更新と行動応答成分の遷移 $g_O$ による継続、
  およびこれに許容外的作用 $\upsilon\in\Upsilon_{\mathrm{adm},i}$ を加えた別の継続の範囲内では、
  (SC-1) の組織・結合関係が同一個体 $\Sigma_i$ として再び成立する継続を、
  本枠組みは認めない。
\end{description}

(SC-1) および (SC-2) の判定は、対象系ごとに当該判定に先立って固定された判定規則に従う。
この判定規則は (SC-2) における同一個体性の判定基準を含み、
判定対象の経過または結果に応じて事後的に変更しない。
本論文は、この判定規則の一般形式を Phase~1 では与えない。

さらに、$\mathrm{CS}_{i,t}=1$ となる時刻 $t$ 以後については、
本枠組みは $\Sigma_i$ に同一個体として帰属する実現を認めない
（以下、\emph{帰属の規定}と呼ぶ）。
したがって、その時刻以後の各時刻 $s$ について、
$\Sigma_i$ に帰属する機能的活動および基盤機能は成立しない。
また、時刻 $s$ には定義~\ref{def:wf_realization} の意味での $\Sigma_i$ の実現が存在しないから、
$\mathrm{WF}_{i,s}=0$ である。

$\Sigma_i$ の\emph{崩壊指標} $\mathrm{Col}_{i,t}\in\{0,1\}$ を
\[
  \mathrm{Col}_{i,t}=1
  \quad:\Longleftrightarrow\quad
  \exists t'\in\{0,1,\ldots,t\}:\ \mathrm{CS}_{i,t'}=1
\]
と定め、それ以外のとき $\mathrm{Col}_{i,t}=0$ とする。
$\mathrm{Col}_{i,t}=1$ は、時刻 $t$ までに崩壊状態が成立したことを表す履歴である。
$\mathrm{Col}_{i,t}=1$ であるとき、
時刻 $t$ までに $\Sigma_i$ の\emph{崩壊}が成立済みであるという。
これは、時刻 $t$ に $\Sigma_i$ に帰属する実現が存在し、
その実現が崩壊した状態にあることを表すものではない。
\end{definition}

\begin{definition}[機能的終端状態指標と機能的終端指標]
\label{def:functional_termination}
ACS個体 $\Sigma_i$ について、
\[
  \mathbb{T}_{\mathrm{FT},i}
  :=
  \bigl\{
    t\in\mathbb{N}_0
    \bigm|
    \forall s\in\mathbb{N}_0\ (s<t \Rightarrow \mathrm{CS}_{i,s}=0)
  \bigr\}
\]
を\emph{機能的終端判定時刻域}と呼ぶ。
$\mathbb{T}_{\mathrm{FT},i}$ は、崩壊状態が一度も成立していない時刻と、
崩壊状態が初めて成立する時刻からなる。

各 $t\in\mathbb{T}_{\mathrm{FT},i}$ について、
次の (FT-1) および (FT-2) がともに成立するとき
\[
  \mathrm{FTS}_{i,t}=1,
\]
それ以外のとき $\mathrm{FTS}_{i,t}=0$ と定め、
$\mathrm{FTS}_{i,t}\in\{0,1\}$ を\emph{機能的終端状態指標}と呼ぶ。
$t\notin\mathbb{T}_{\mathrm{FT},i}$ については $\mathrm{FTS}_{i,t}$ を定めない。

\begin{description}[style=nextline]
  \item[(FT-1) 全面的機能停止]
  時刻 $t$ において、次の (a) および (b) がともに成立する。
  \begin{enumerate}[label=(\alph*)]
    \item 原始的な機能（感知・識別・評価・行動、定義~\ref{def:primitive_functions}）の
    各区分に関して構造的個体 $\Sigma_i$ が有する基礎的な機能
    （定義~\ref{def:basic_functions}）について、
    その機能的活動がすべて停止している。
    これは時刻 $t$ における機能的活動の全面的な不成立を指す状態条件であり、
    それ以前に活動が成立していたことを要求しない。
    また、選択的待機 $u_{\emptyset}$ の選択その他の
    待機または静止ではない（備考~\ref{rem:primitive_activity_cessation}）。
    \item (a) の基礎的な機能の活動を支える基盤機能
    （定義~\ref{def:base_functions}）が、
    すべて停止している。
  \end{enumerate}
  \item[(FT-2) 本枠組みに相対する同一個体としての非回復性]
  時刻 $t$ 以後について、実現した外生背景成分の列 $(b_s)_{s\geq t}$ のもとでの
  $\Sigma_i$ 自身の通常の自律的更新と行動応答成分の遷移 $g_O$ による継続、
  およびこれに許容外的作用 $\upsilon\in\Upsilon_{\mathrm{adm},i}$ を加えた別の継続の範囲内では、
  (a) の基礎的な機能のいずれかの機能的活動、
  または (b) の基盤機能のいずれかが、
  同一個体 $\Sigma_i$ として再び成立する継続を、
  本枠組みは認めない。
\end{description}

$\Sigma_i$ の\emph{機能的終端指標} $\mathrm{FT}_{i,t}\in\{0,1\}$ を、
すべての $t\in\mathbb{N}_0$ について
\[
  \mathrm{FT}_{i,t}=1
  \quad:\Longleftrightarrow\quad
  \exists t'\in\mathbb{T}_{\mathrm{FT},i}\cap\{0,1,\ldots,t\}:\ \mathrm{FTS}_{i,t'}=1
\]
と定め、それ以外のとき $\mathrm{FT}_{i,t}=0$ とする。
$\mathrm{FT}_{i,t}=1$ は、時刻 $t$ までに機能的終端状態が成立したことを表す履歴であり、
時刻 $t$ に $\Sigma_i$ が機能的終端状態にあることを表すものではない。
$\mathrm{FT}_{i,t}=1$ かつ $\mathrm{Col}_{i,t}=0$ であるとき、
$\Sigma_i$ は時刻 $t$ において\emph{機能的終端後}にあるという。

(FT-1) における機能的活動および基盤機能の判定、
および (FT-2) の判定は、対象系ごとに当該判定に先立って固定された判定規則に従う。
この判定規則は (FT-2) における同一個体性の判定基準を含み、
定義~\ref{def:collapse} の帰属の規定と整合するように固定され、
判定対象の経過または結果に応じて事後的に変更しない。
本論文は、この判定規則の一般形式を Phase~1 では与えない。
\end{definition}

(SC-2) と (FT-2) が量化する継続は、
定義~\ref{def:env_realization} と備考~\ref{rem:exogenous_environment} に従い、
外生的に与えられる背景成分の列を実現したものに固定したうえで、
当該ACSの自律的更新と行動応答成分の遷移によって生成される継続、
およびそれに許容外的作用を加えた継続である。
背景成分の列として起こりえた他の列を量化の対象としない。

\begin{proposition}[実現した軌道上の再成立]
\label{prop:realized_reestablishment}
時刻 $t$ について、実現した軌道の上のある時刻 $s>t$ に、
(SC-1) の組織・結合関係が同一個体 $\Sigma_i$ として再び成立していると
定義~\ref{def:collapse} の判定規則により判定されるならば、$\mathrm{CS}_{i,t}=0$ である。
同様に、$t\in\mathbb{T}_{\mathrm{FT},i}$ について、実現した軌道の上のある時刻 $s>t$ に、
(FT-1) の (a) の基礎的な機能のいずれかの機能的活動、
または (b) の基盤機能のいずれかが同一個体 $\Sigma_i$ として再び成立していると
定義~\ref{def:functional_termination} の判定規則により判定されるならば、
$\mathrm{FTS}_{i,t}=0$ である。
\end{proposition}

\begin{proof}
実現した軌道の時刻 $t$ 以後の部分は、実現した背景成分の列のもとで、
$\Sigma_i$ 自身の自律的更新と $g_O$ による遷移によって生成される
（備考~\ref{rem:exogenous_environment}）。
したがってそれは (SC-2) および (FT-2) が量化する継続の一つである。
その継続の上で再成立が判定されるから、本枠組みは再成立する継続を認めることになり、
(SC-2) は成立しない。よって $\mathrm{CS}_{i,t}=0$ である。
(FT-2) についても同様であり、$\mathrm{FTS}_{i,t}=0$ である。
\end{proof}

命題~\ref{prop:realized_reestablishment} の逆、
すなわち実現した軌道の上で再成立しないことから $\mathrm{CS}_{i,t}=1$ を導くことは置かない。
(SC-2) と (FT-2) の量化には、許容外的作用を加えた別の継続も含まれるからである。
実現した救助は軌道の上で記述し、許容外的作用は回復可能性の判定にだけ用いる。

状態指標 $\mathrm{CS}_{i,t}$・$\mathrm{FTS}_{i,t}$ はその時刻に各状態が成立したか否かを表し、
指標 $\mathrm{Col}_{i,t}$・$\mathrm{FT}_{i,t}$ は、
それが時刻 $t$ までに一度でも成立したことを保持する履歴指標である。
各時刻の状態指標が与えられれば、対応する履歴指標は一意に定まる。
これらの指標と時刻域 $\mathbb{T}_{\mathrm{FT},i}$ は7タプルの成分ではなく、
当該ACSがこれらを内部に保持または計算することを要求しない
（\S\ref{subsec:description_levels_global}）。

\begin{proposition}[機能的終端指標と崩壊指標の基本性質]
\label{prop:collapse_basic}
定義~\ref{def:collapse} および定義~\ref{def:functional_termination} のもとで、
次が成り立つ。
\begin{description}[style=nextline]
  \item[(C-2) 吸収性]
  $\mathrm{Col}_{i,t}=1$ ならば、すべての $s>t$ について $\mathrm{Col}_{i,s}=1$ である。
  同様に、$\mathrm{FT}_{i,t}=1$ ならば、すべての $s>t$ について $\mathrm{FT}_{i,s}=1$ である。
  \item[(C-3) 崩壊状態成立時刻の存在]
  $\mathrm{Col}_{i,t}=1$ ならば、ある $t'\leq t$ が存在して、
  時刻 $t'$ において (SC-1) および (SC-2) が成立している。
\end{description}
\end{proposition}

\begin{proof}
(C-3)：$\mathrm{Col}_{i,t}=1$ ならば、定義から $\mathrm{CS}_{i,t'}=1$ となる $t'\leq t$ が
存在し、$\mathrm{CS}_{i,t'}=1$ は (SC-1) および (SC-2) の成立を意味する。
(C-2)：この $t'$ はすべての $s>t$ について $t'\leq s$ を満たすから、
定義から $\mathrm{Col}_{i,s}=1$ である。
$\mathrm{FT}$ についても、$t'\in\mathbb{T}_{\mathrm{FT},i}$ かつ $\mathrm{FTS}_{i,t'}=1$ となる
$t'\leq t$ をとれば同様である。
\end{proof}

吸収性 (C-2) は履歴指標についての性質である。
$\mathrm{FT}_{i,s}=1$ は、時刻 $s$ に機能的終端状態にあることではなく、
時刻 $s$ までに機能的終端状態が成立したことを表す。

\begin{proposition}[機能的終端と崩壊の順序および崩壊後の整形式実現]
\label{prop:termination_order}
定義~\ref{def:collapse} および定義~\ref{def:functional_termination} のもとで、
$\mathrm{Col}_{i,t}=1$ とし、
\[
  t_{\mathrm{C}} := \min\{t'\in\mathbb{N}_0 \mid \mathrm{CS}_{i,t'}=1\}
\]
とおく（$\mathrm{Col}_{i,t}=1$ よりこの集合は空でなく、$t_{\mathrm{C}}\leq t$ である）。
このとき次が成り立つ。
\begin{enumerate}[label=(\roman*)]
  \item $t_{\mathrm{C}}\in\mathbb{T}_{\mathrm{FT},i}$ であり、$\mathrm{FTS}_{i,t_{\mathrm{C}}}=1$ である。
  \item $\mathrm{FT}_{i,t}=1$ である。
  すなわち、崩壊が成立している時刻 $t$ までに機能的終端状態は成立済みであり、
  機能的終端状態が初めて成立する時刻は $t_{\mathrm{C}}$ より後にならない。
  \item $t>t_{\mathrm{C}}$ ならば $t\notin\mathbb{T}_{\mathrm{FT},i}$ である。
  すなわち、崩壊が初めて成立した時刻より後に、機能的終端状態が新たに判定されることはない。
  \item $\mathrm{WF}_{i,t}=0$ である。
\end{enumerate}
\end{proposition}

\begin{proof}
(i)：$t_{\mathrm{C}}$ の最小性から、すべての $s<t_{\mathrm{C}}$ について $\mathrm{CS}_{i,s}=0$ であり、
$t_{\mathrm{C}}\in\mathbb{T}_{\mathrm{FT},i}$ である。
$\mathrm{CS}_{i,t_{\mathrm{C}}}=1$ であるから、定義~\ref{def:collapse} の帰属の規定により、
時刻 $t_{\mathrm{C}}$ において $\Sigma_i$ に帰属する機能的活動および基盤機能は成立せず、
(FT-1) の (a) および (b) が成立する。
また、時刻 $t_{\mathrm{C}}$ 以後に (a) の基礎的な機能のいずれかの機能的活動、
または (b) の基盤機能のいずれかが同一個体 $\Sigma_i$ として再び成立する継続は、
$\Sigma_i$ に同一個体として帰属する実現を含む。
本枠組みは時刻 $t_{\mathrm{C}}$ 以後にそのような実現を認めないから、(FT-2) が成立する。
よって $\mathrm{FTS}_{i,t_{\mathrm{C}}}=1$ である。

(ii)：(i) と $t_{\mathrm{C}}\leq t$ より、定義から $\mathrm{FT}_{i,t}=1$ である。
機能的終端状態が初めて成立する時刻は、(i) により $t_{\mathrm{C}}$ 以下である。

(iii)：$t>t_{\mathrm{C}}$ ならば $s=t_{\mathrm{C}}\in\{0,1,\ldots,t-1\}$ について
$\mathrm{CS}_{i,s}=1$ であるから、$t\notin\mathbb{T}_{\mathrm{FT},i}$ である。

(iv)：帰属の規定と定義~\ref{def:wf_realization} により、
$t_{\mathrm{C}}$ 以後の各時刻 $s$ で $\mathrm{WF}_{i,s}=0$ であり、$t\geq t_{\mathrm{C}}$ である。
\end{proof}

崩壊が初めて成立する時刻 $t_{\mathrm{C}}$ における $\mathrm{FTS}_{i,t_{\mathrm{C}}}$ は、
崩壊後の実現について状態を評価した値ではなく、
同一個体の終端の境界時刻における判定値である。
その値は、帰属の規定により (FT-1) および (FT-2) が成立することから定まる。
この境界の規定により、機能的終端と崩壊が同じ時刻に初めて成立する場合を、
離散時刻の上で表すことができる。

\begin{remark}[機能的終端・崩壊と他の状態の区別]
\label{rem:termination_hierarchy}
本論文は、次の状態を互いに区別する。
\begin{description}[style=nextline]
  \item[構造の破損・欠損]
  ACS個体を構成する部位、結合、情報構造または機能的構造の一部が、
  損傷・欠落・切断・劣化している状態。
  同一個体としての組織関係がなお維持されている限り、崩壊ではない。
  \item[部分的な機能不全・機能喪失]
  原始的な機能の各区分に対応する基礎的な機能その他の機能の一部が低下・停止し、
  またはその一部が回復しない形で失われている状態。
  原始的な機能のいずれかの区分に対応する基礎的な機能の一部が
  回復しない形で失われても、
  また一つ以上の区分に対応する基礎的な機能が
  回復しない形で全面的に失われても、
  他の区分の基礎的な機能または基盤機能の活動が残る限り、
  それだけでは機能的終端でも崩壊でもない。
  そのような喪失は $\Sigma_i$ のACS個体としての同定を取り消さない
  （構成要件 CR-2 の直後の説明）。
  \item[機能的終端]
  定義~\ref{def:functional_termination} の状態。ACS個体としての活動の終端であり、
  構造的個体そのものの崩壊ではない。
  \item[崩壊]
  定義~\ref{def:collapse} の状態。構造的個体としての同一性を成立させていた
  組織関係の、同一個体として回復しない喪失である。
  離散、分解、散逸、結合関係の断絶、構造的同一性を担う情報組織の消失などは
  その例となりうるが、定義そのものではない。
  構成要素が残存していても、それらの間の組織関係が回復しない形で失われていれば
  (SC-1)・(SC-2) は成立しうる。
  一方、修理または再組立により同一個体として組織関係が再び成立しうる分解・破損は、
  (SC-2) を満たさない。
\end{description}

構造の破損・欠損は、機能の低下を経て機能的終端に至る経路、
または崩壊に至る経路の要因となりうる。
ただし本論文は、構造の破損・欠損から機能的終端または崩壊を一般的な帰結として置かず、
部分的な機能喪失から機能的終端または崩壊を一般的な帰結として置かない。
冗長な構造、残存機能、代替経路、補償、修復または機能の縮退によって、
損傷後もACSとしての活動が続きうるためである。
また本論文は、機能的終端から同じ時刻の崩壊を一般的な帰結として置かない。
本論文が採用する時間関係は、命題~\ref{prop:termination_order} の
「崩壊が成立している時刻までに機能的終端状態は成立済みである」ことと、
「崩壊が初めて成立した時刻より後に機能的終端状態が新たに判定されることはない」ことだけである。

(FT-2) の非回復性は、(FT-1)(a) の原始的な機能の各区分に関する
基礎的な機能の活動と、(FT-1)(b) の基盤機能の両方について要求する。
したがって、基礎的な機能の活動または基盤機能のいずれか一つでも
同一個体として再び成立しうる時刻は、機能的終端状態の時刻ではない。

本論文は、閉塞 $A_{i,t}=\emptyset$（定義~\ref{def:blockage_sec2}）、
整形式実現の不成立 $\mathrm{WF}_{i,t}=0$、
および \S\ref{sec:three_principles} の三原理の時刻別充足述語が偽であることからは、
機能的終端も崩壊も一般的な帰結として置かない。
また、ある $\eta\in E$ について $x_{i,t}\notin D_\eta$
（定義~\ref{def:eta_satisfiable_domain}）であることは、
当該要因が自己側の状態のために充足不能であるという状態側の診断であり、
(FT-1)(b) の基盤機能の停止そのものとは同一視しない。

自己に関する記憶、自己認識の様式、その他の認識内容の喪失または
回復しない変化は、構造的自己の終端とは区別される。
定義~\ref{def:time_identity} のとおり、
認識内容は構造的存続の判定に影響しない。
したがって本論文は、これらの喪失から機能的終端または崩壊を一般的な帰結として置かない。
\end{remark}

\begin{remark}[機能的活動の全面的不成立と待機・静止の区別]
\label{rem:primitive_activity_cessation}
(FT-1)(a) は、原始的な機能の各区分に関して構造的個体が有する
基礎的な機能の活動の一部がなお成立している時刻を、
機能的終端から除くための条件である。
感知・識別・評価・行動は抽象的な機能区分であり、判定の対象は
定義~\ref{def:three_primitives} の標準写像形式そのものではなく、
構成要件 CR-2 が要請する責務を担う個体の基礎的な機能である。
したがって、機能的に等価な別の実現についても同じ責務の活動として判定する。

選択的待機 $u_{\emptyset}$（備考~\ref{rem:action_space_structure}）は、
行動の機能的責務による選択の結果であるから、
その選択が行われている時刻には行動の機能的活動は不成立ではない。
感知、識別または評価の活動が続く休止・静止も同様である。
また、自律的機能過程 $f_{\mathrm{auto}}$（仮定~\ref{ass:f_decomposition}）による状態変化、
および化学反応・温度変化・慣性運動・残留電流その他の物理的な変化の継続は、
それ自体では基礎的な機能の活動ではない。
したがって、(FT-1)(a) の成立時刻にも状態 $x_{i,t}$ は変化しうる。
ただし、$f_{\mathrm{auto}}$ に表される変化のうち基盤機能
（定義~\ref{def:base_functions}）の作動にあたるものは、
(FT-1)(b) の判定の対象である。
例えば、基礎的な機能の活動が停止していても、
物質の輸送のような基盤機能がなお作動している時刻には、
(FT-1)(b) は成立しない。
(FT-1) は、物理的な活動が一切存在しないことを要求しない。
本論文は、外部から観測される出力の不変、顕在的な運動の欠如、
または入力の変化の欠如のみから、(FT-1)(a) の成立を一般的な帰結として置かない。
(FT-1)(a) の判定規則は、当該の機能的活動が実際に担われているか否かを判定するものとする。

(FT-1)(b) は、(a) の活動を支える基盤機能のすべての停止を求める条件であり、
構成要素が物理的に残存していることを妨げない。
本論文は (FT-1) の (a) と (b) の相互の独立性・排反性を要求しない。
\end{remark}

\begin{remark}[機能的終端・崩壊と三原理]
\label{rem:collapse_not_principle_failure}
(FT-1) は原始的な機能の各区分に関して構造的個体が有する
基礎的な機能の活動と基盤機能に関する条件であり、
(SC-1) は構造的同一性を成立させる組織関係に関する条件である。
いずれも機能または構造的同一性に関する条件であり、
\S\ref{sec:three_principles} の三原理の時刻別充足述語の否定ではない。
機能的終端は、原始的な機能の各区分に関して構造的個体が有する
基礎的な機能の活動と、それを支える基盤機能のすべてが停止し、
そのいずれも同一個体として回復しない状態であり、
崩壊は、構造的個体そのものの終端である。
一方、三原理は存続中のACSに対する制約軸であり、
真偽をもつのはその時刻別充足述語である
（備考~\ref{rem:primitive_vs_principles}）。
したがって本論文は、機能的終端も崩壊も時刻別充足述語が偽であることとして定義せず、
機能的終端状態が成立した時刻以後の時刻別充足述語を偽と評価するのでもなく、
\S\ref{sec:three_principles} の時刻別充足述語の判定を
$\mathrm{FT}_{i,t}=0$ の時刻に適用する。
命題~\ref{prop:termination_order}(ii) により、そのような時刻では $\mathrm{Col}_{i,t}=0$ でもある。
ある量が機能的終端の後も形式上計算できることと、
それが時刻別充足述語の判定として意味をもつこととは区別する。
\end{remark}

\begin{remark}[機能的終端・崩壊と整形式実現指標の関係]
\label{rem:collapse_wf_relation}
整形式実現指標 $\mathrm{WF}_{i,t}$（定義~\ref{def:wf_realization}）は、
時刻 $t$ に $\Sigma_i$ の7成分が同一個体の実現において型のとおりに与えられているか否かを表し、
機能的活動が成立しているか否かを表さない。
したがって本論文は、$\mathrm{FT}_{i,t}$ と $\mathrm{WF}_{i,t}$ の間に
いずれの向きの含意も仮定しない。
機能的終端後で $\mathrm{WF}_{i,t}=1$ である時刻、
すなわち機能的活動は回復しない形で失われているが
7成分は型のとおりに与えられている時刻を排除しない。
機能的終端の後に整形式実現が失われることを述べる場合には、
構成要件 CR-1 の (CR-1a)(ii) の前件、すなわち
ある要因 $\eta\in E$ とある時刻 $t_0\geq1$ について
$\mathrm{WF}_{i,t_0-1}=1$、$\eta(x_{t_0-1},e_{t_0-1})=1$ であり、
すべての $s\geq t_0$ で $\eta(x_s,e_s)=0$ となることを、
別途示す必要がある。
本論文は、この前件を (FT-1)・(FT-2) からの一般的な帰結として置かない。

これに対し崩壊については、命題~\ref{prop:termination_order}(iv) により
$\mathrm{Col}_{i,t}=1$ ならば $\mathrm{WF}_{i,t}=0$ である。
この帰結は、崩壊状態の時刻以後に $\Sigma_i$ に帰属する実現を認めないという
本枠組みの規約（帰属の規定）と、整形式実現指標の定義とから従うものであり、
対象系の物理的過程についての主張ではない。
逆に本論文は、$\mathrm{WF}_{i,t}=0$ から $\mathrm{Col}_{i,t}=1$ も $\mathrm{FT}_{i,t}=1$ も一般的な帰結として置かない。
一時的に整形式実現を失っているが、機能的終端後にも崩壊にもない時刻を排除しない。
\end{remark}

\begin{remark}[機能的終端後と崩壊後の記録]
\label{rem:post_collapse_record}
機能的終端後の時刻にも、個体 $\Sigma_i$ の識別と記録は維持されうる。
状態 $x_{i,t}$ や外部環境状態 $e_{i,t}$ は、
構造的個体が崩壊していない限り記録の対象となりうる。
一方で、機能的終端後の実現は、現に作動するACSの実現としては扱わない。
$\mathrm{Col}_{i,t}=1$ となった後は、
本枠組みは $\Sigma_i$ に同一個体として帰属する実現を認めず、
機能的活動を帰属させる構造的個体そのものが終端している。
崩壊後の時刻を「機能的終端状態が続いている」時刻としては扱わない
（命題~\ref{prop:termination_order}(iii)）。
構成要件 CR-2 は、時刻ごとの診断値ではなく、
$\Sigma_i$ をACS個体として同定する条件である
（備考~\ref{rem:math_vs_realization}）。
本論文は、機能的終端後または崩壊後の実現を、時刻ごとに改めて CR-2 を満たす
現に作動するACSとして再判定することはしない。
また、部分的な機能喪失、一つ以上の区分の全面的な機能喪失、
機能的終端または崩壊によって、
それ以前に $\Sigma_i$ がACS個体として同定された履歴
および個体識別が遡って取り消されることはない。

\S\ref{sec:collapse_cond} のACS記録有効時刻域は $\mathrm{Col}_{i,t}=0$ の時刻からなり、
機能的終端後の時刻を含みうる。
そこでの記録は構造的個体としての追跡の記録であり、
現に作動するACSとしての機能の記述は $\mathrm{FT}_{i,t}=0$ の時刻に限る。
状態帯および到達経路を機能的終端と崩壊の二段に対してどのように記述するかは、
同節で扱う。
\end{remark}

%----------------------------------------------------------------------
\subsection{必要環境空間 $E$ と必要環境充足}
\label{subsec:environment_space}

本小節では、7タプルの第三成分である必要環境空間 $E$ の意味論を定める。
$E$ の基本型は定義~\ref{def:acs}(3) で
\[
  E\subseteq\{0,1\}^{X\times E_{\mathrm{ext}}},
  \qquad
  E\neq\emptyset
\]
として宣言済みであり、本小節では、その各元がなぜ「必要」と
呼ばれるかを拘束する構成要件 CR-1 を与える。
あわせて、必要要因がすべて充足される外部環境状態の集合
$E_{\mathrm{req}}(x)$ を派生的に定義し、
一時的な非充足と持続的な欠損の区別、
要因分解の非一意性、および必要性の同定可能性について整理する。

CR-1 は、\S\ref{subsec:self_formal_definition} で導入した
整形式実現指標 $\mathrm{WF}_{i,t}$
（定義~\ref{def:wf_realization}）を用いて記述する。
崩壊指標 $\mathrm{Col}_{i,t}$（定義~\ref{def:collapse}）は用いない。
これにより、$E$ の意味論的拘束が
崩壊定義に依存することを避ける。

本小節で新たに用いる記号は、必要環境要因 $\eta$、
要因別充足状態域 $D_\eta$、
必要環境充足集合 $E_{\mathrm{req}}(x)$ である。
下付き $\mathrm{req}$ は required（要求される）に対応する。
$E_{\mathrm{req}}(x)$ は $E_{\mathrm{ext}}$ の部分集合であり、
必要環境空間 $E$ そのものとは異なる対象である。

\begin{definition}[必要環境要因と必要環境空間]
\label{def:environment_space}
定義~\ref{def:acs}(3) のもとで、$E$ の各元
\[
  \eta:X\times E_{\mathrm{ext}}\to\{0,1\}
\]
を、ACS個体 $\Sigma_i$ の\emph{必要環境要因}と呼び、
$E$ を $\Sigma_i$ の\emph{必要環境空間}と呼ぶ。

$\eta(x,e)=1$ は、内部状態 $x$ と外部環境状態 $e$ の組において
当該要因が充足されていることを表し、
$\eta(x,e)=0$ は非充足であることを表す。
$\eta$ が内部状態 $x$ を引数に取るのは、
同一の外部環境状態のもとでも、内部の蓄え、活動水準、
または必要環境を利用する自己側の機能の状態
（備考~\ref{rem:env_utilization_function}）に応じて
充足の成否が異なりうるためである。
ただし、$\eta$ が外部環境状態に依存しない内部状態のみの条件となることは、
構成要件 CR-1 の (CR-1c) によって排除される。

$E$ の元は集合 $\{0,1\}^{X\times E_{\mathrm{ext}}}$ の元であるから、
同一の充足判定を与える二つの要因は同一視される。
Phase~1 では $E$ の濃度に制限を課さない。
\end{definition}

\begin{remark}[必要環境要因は特定の供給手段ではない]
\label{rem:factor_functional_class}
$E$ の元は、特定の物質、特定の供給経路または特定の装置ではなく、
それらを充足しうる経路を束ねた充足判定として与える。
例えば、代謝に要するエネルギーおよび栄養素の取得可能性、
または持続的な電力供給の可能性は、それぞれ一つの $\eta$ として
表されうるのに対し、その充足を実現する個別の供給手段は
$E$ の元ではない。

この扱いは、代替可能性を $\eta$ の内部に吸収するためである。
特定の供給手段を $E$ の元とすると、代替手段が存在する場合に
当該元は構成要件 CR-1 を満たさなくなる。
どの粒度で機能的な束ね方を行うかは実装依存であり、
その扱いは備考~\ref{rem:factor_nonuniqueness} で述べる。
\end{remark}

\begin{remark}[必要環境要因の標準的粗分類候補]
\label{rem:eta_coarse_classes}
本論文では、必要環境要因 $\eta\in E$ を粗く分類するための
標準的候補として、次の四つのラベルからなる集合
\[
  \mathsf{Cls}_\eta
  :=
  \{\,
    \mathrm{pow},\
    \mathrm{mat},\
    \mathrm{sink},\
    \mathrm{cond}
  \,\}
\]
を挙げる。
各ラベルの指す範囲は次のとおりである。

\begin{description}[style=nextline]
  \item[$\mathrm{pow}$（動力・利用可能自由エネルギー関係）]
  当該ACSの作動に要する動力または利用可能な自由エネルギーの
  取得可能性に関する必要性。
  \item[$\mathrm{mat}$（構造・物質資源関係）]
  構造の維持および補修・交換に要する物質資源の
  取得可能性に関する必要性。
  \item[$\mathrm{sink}$（散逸・排出関係）]
  熱、老廃物その他の副産物を外部へ散逸または排出しうることに
  関する必要性。
  \item[$\mathrm{cond}$（外部作動・境界条件関係）]
  $\mathrm{pow}$、$\mathrm{mat}$ または $\mathrm{sink}$ に
  還元されない外部側の条件、すなわち当該ACSが作動しうる
  物理的・環境的な境界状態、制約、必要な信号の受領可能性、
  対象への接近可能性、および相互作用の可能性が
  保たれていることに関する必要性。
\end{description}

本備考は分類候補の提示であり、定義、構成要件または仮定では
ない。
本論文は、この四分類が任意のACSに対して網羅的または排他的で
あることを証明していない。
したがって次の限定を付す。

\begin{enumerate}[label=(\alph*)]
  \item $\mathsf{Cls}_\eta$ による分類は $E$ の集合論的分割
  ではない。
  \item 同一の具体的な外部環境対象が複数のラベルに関わる
  必要性を充足しうる。
  \item 分類の対象は具体的な供給手段ではなく機能的必要性で
  ある（備考~\ref{rem:factor_functional_class}）。
  \item 各ACSが四つのラベルすべてに対応する要因をもつことを
  要求しない。
  \item 非常時に要する必要性を独立した第五のラベルとして
  立てない。
  \item 非常時の必要性についても、四つのラベルのいずれかに
  よって記述できる場合には、対応する要因が内部状態 $x$ に
  応じて条件付きで発動するものとして扱いうる。
  $\eta$ が $x$ を引数に取ること
  （定義~\ref{def:environment_space}）がこれを可能にする。
  本論文は、あらゆる非常時の必要性が四つのラベルによって
  記述されることを主張しない。
  \item より細かな分類は実装依存であり、Phase~1 では固定しない。
\end{enumerate}

概念的には
$\mathsf{Cls}_\eta\times\{\text{通常発動},\text{条件付き発動}\}$
という理解も可能であるが、本論文はこの積集合を新たな
正式な型として導入しない。

なお、外部環境の変化に対してACS側がどのように応答しうるかは、
$\eta$ が記述する外部依存とは理論階層の異なる事項であり、
\S\ref{subsec:update_operator} で別に扱う。
\end{remark}

\begin{definition}[要因別充足状態域（分類C）]
\label{def:eta_satisfiable_domain}
各必要環境要因 $\eta\in E$ に対し、
\[
  D_\eta
  :=
  \bigl\{\,
    x\in X
    \;\bigm|\;
    \exists e\in E_{\mathrm{ext}}:\ \eta(x,e)=1
  \,\bigr\}
  \subseteq X
\]
を $\eta$ の\emph{要因別充足状態域}と呼ぶ。
$x\in D_\eta$ は、内部状態 $x$ において、
ある外部環境状態のもとで当該要因が充足されうることを表し、
$x\notin D_\eta$ は、いかなる外部環境状態のもとでも
当該要因が充足されないことを表す。

$D_\eta$ は $\eta$、$X$ および $E_{\mathrm{ext}}$ から定まる派生集合であり、
7タプルの独立成分ではない。
定義式は状態遷移写像 $f$ を用いない。
記号 $D_\eta$ は、\S\ref{sec:tipa} の $M_1$ 出力成分 $D_{M1,t}$ とは無関係である。
\S\ref{sec:three_principles} の必要環境充足状態域 $D_{\mathrm{req}}$ との関係は、
同節の定義~\ref{def:req_domain} で述べる。
\end{definition}

\begin{remark}[必要環境利用機能]
\label{rem:env_utilization_function}
要因 $\eta$ と外部環境状態空間 $E_{\mathrm{ext}}$ を固定したもとでは、
内部状態 $x$ が $D_\eta$ に属するか否かは、
実現外部環境状態 $e_t$ の値には依存せず、内部状態 $x$ の性質として定まる。
ただし $D_\eta$ 自体は $E_{\mathrm{ext}}$ の取り方に相対的である。
本論文は、内部状態を $D_\eta$ に属させる自己側の条件の成立に関与する
一つまたは複数の既存機能の関与・協働を総称して、
当該要因に関する\emph{必要環境利用機能}と呼ぶ。
ここでいう既存機能には、原始的な機能の各区分に対応する基礎的な機能
（\S\ref{subsec:three_primitive_functions}）、
内部モデル構造とその更新
（\S\ref{subsec:internal_model_M}、\S\ref{subsec:update_operator}）、
および \S\ref{sec:tipa} の TIPA・PSL・ISL の標準機能形式が担う機能が含まれうる。
例えば、呼吸可能な大気の取得可能性を一つの必要環境要因とする場合、
呼吸・循環に関わる自己側の諸機能の協働がこれに当たりうる。
\S\ref{sec:three_principles} の環境利用（EU）との関係では、
必要環境利用機能は、EU の原理応答項目に関与する
既存機能（単一の機能である場合を含む）の協働の総称として位置づけられる。

本備考は総称を与えるものであり、
新たな写像、変数、7タプルの成分、機構、構成要件または仮定を導入しない。
必要環境利用機能を、既存の機能群と並ぶ独立の機能群としては扱わない。
また、構成的必要機能の候補としての記録にとどめ、
分類A表記を付与しない。
その構成的必要性の検討は Phase~1.$\alpha$ の課題とする（\S\ref{subsec:phase1a}）。

以下の限定を付す。

\begin{enumerate}[label=(\alph*)]
  \item ある $\eta\in E$ について $x\notin D_\eta$ ならば、
  $E_{\mathrm{req}}(x)=\emptyset$ である（定義~\ref{def:E_req}）。
  すなわち、当該要因に関する自己側の条件が成立しない内部状態では、
  外部環境状態のいかんにかかわらず必要環境充足は成立しない。
  \item 逆に、本論文は、すべての $\eta\in E$ について $x\in D_\eta$ であることから、
  $E_{\mathrm{req}}(x)\neq\emptyset$ を一般的な帰結として置かない。
  要因ごとの充足可能性から共同充足可能性への含意も仮定しない。
  \item 本論文は、$x\notin D_\eta$ それ自体から、
  $x\notin K$、閉塞 $A_t=\emptyset$（定義~\ref{def:blockage_sec2}）、
  $\mathrm{WF}_{i,t}=0$（定義~\ref{def:wf_realization}）、
  機能的終端（定義~\ref{def:functional_termination}）
  または崩壊（定義~\ref{def:collapse}）を一般的な帰結として置かない。
  これらとの関係は $\Pi$、$f$ および構成要件 CR-1 に依存する。
  \item 必要環境利用機能は、備考~\ref{rem:env_change_response} の
  環境変化応答機能とは区別される。
  前者は所与の必要環境を自己側で利用しうる状態にあるか否かに関わり、
  後者は外部環境の変化に対して現在または将来の応答を変化させうる
  横断的機能に関わる。
\end{enumerate}
\end{remark}

\subsubsection*{構成要件 CR-1（必要環境整合性）}
\label{req:CR1}

7タプル構造体 $\Sigma_i$ が本論文におけるACSであるためには、
その必要環境空間 $E$ について次の (CR-1a)、(CR-1b) および (CR-1c) が
成立することを要求する。

\medskip
\noindent\textbf{(CR-1a) 存続必要性：}
以下では、\emph{状態遷移列}を
\[
  (x_s,u_s,e_s)_{s\geq0},
  \qquad
  x_{s+1}=f(x_s,u_s,e_s),
  \quad
  u_s\in U,
  \quad
  e_s\in E_{\mathrm{ext}}
\]
を満たす列（定義~\ref{def:state_transition}）と呼ぶ。
状態遷移列は7タプルに従う数学的な列であり、
その定義は、各時刻において整形式な実現が成立していることを要件として含まない。
各時刻の実現の成否は $\mathrm{WF}_{i,s}$ が別に記述する。

各 $\eta\in E$ について、次の (i) および (ii) が成立することを要求する。

\begin{enumerate}
  \item[(i)] \textbf{欠損の実現可能性：}
  ある状態遷移列とある時刻 $t_0\geq1$ が存在して
  \[
    \mathrm{WF}_{i,t_0-1}=1,
    \qquad
    \eta(x_{t_0-1},e_{t_0-1})=1
  \]
  かつ
  \[
    \eta(x_s,e_s)=0
    \qquad
    \text{（すべての } s\geq t_0 \text{ について）}
  \]
  が成立する。

  \item[(ii)] \textbf{欠損時の存続不能性：}
  (i) の条件を満たす任意の状態遷移列と時刻 $t_0$ の組について、
  ある有限時刻 $T$（$T\geq t_0$）が存在して
  \[
    \mathrm{WF}_{i,s}=0
    \qquad
    \text{（すべての } s\geq T \text{ について）}
  \]
  が成立する。
\end{enumerate}

\medskip
\noindent\textbf{(CR-1b) 共同充足可能性：}
\[
  \exists x\in X:\quad
  \bigl\{
    e\in E_{\mathrm{ext}}
    \;\bigm|\;
    \forall \eta\in E,\ \eta(x,e)=1
  \bigr\}
  \neq\emptyset
\]
が成立する。

\medskip
\noindent\textbf{(CR-1c) 外部依存性（外部環境状態による非充足の実現可能性）：}
各 $\eta\in E$ について、ある状態遷移列とある時刻 $t_0\geq1$ が存在して
\[
  \mathrm{WF}_{i,t_0-1}=1,
  \qquad
  \eta(x_{t_0-1},e_{t_0-1})=1,
  \qquad
  x_{t_0}\in D_\eta
\]
かつ
\[
  \eta(x_s,e_s)=0
  \qquad
  \text{（すべての } s\geq t_0 \text{ について）}
\]
が成立する（$D_\eta$ は定義~\ref{def:eta_satisfiable_domain}）。

\medskip
(CR-1a)(i) は、$\eta$ の持続的な非充足が状態遷移列として
実現しうることを要求する。
この条件がなければ、たとえば
$\eta(x,e)=1$ が恒真である写像を $E$ に含めても
(ii) の前件が一度も成立せず、(ii) が空虚に成立してしまう。
(i) はこれを排除し、(ii) とあわせて、存続に不要な要因が $E$ に含まれないことを要求する。
時刻 $t_0-1$ において
$\mathrm{WF}_{i,t_0-1}=1$ および
$\eta(x_{t_0-1},e_{t_0-1})=1$ を課すのは、
非充足の開始時刻を $t_0$ として特定し、
かつ開始直前に整形式な実現が成立していた列のみを
証拠として用いるためである。
これにより (ii) において $T=t_0$ を取ることが可能となり、
非充足の開始と同時に実現が失われる場合と、
遅れて失われる場合とをともに表現できる。

(CR-1a)(ii) は、そのような非充足が持続する場合に、
有限時刻以降の実現が回復しないことを要求する。

(CR-1b) は、必要要因集合の全体が共同で充足不能となる
退化的構成を排除する。
これは (CR-1a)(i) が要求する各要因ごとの非空虚性とは別の条件であり、
本論文は両者の間にいずれの向きの含意も仮定しない。

(CR-1c) は、各要因について、持続的な非充足の開始時点の内部状態を固定したとき、
外部環境状態の違いによって当該要因の充足・非充足が分かれる
状態遷移列が少なくとも一つ存在することを要求する。
すなわち、非充足の開始時刻 $t_0$ において内部状態は $D_\eta$ に属し、
ある外部環境状態のもとでは当該要因が充足されうるにもかかわらず、
実現外部環境状態 $e_{t_0}$ のもとでは充足されない。
本条件は、非充足の原因が外部環境状態にあることを同定するものではない。
時刻 $t_0-1$ から $t_0$ にかけて内部状態も変化しうるからである。
この条件がなければ、ある $\rho:X\to\{0,1\}$ により
$\eta(x,e)=\rho(x)$ と書ける、外部環境状態に依存しない写像も、
$\rho$ の持続的な非充足が整形式実現の喪失を導く限り (CR-1a) を満たし、
内部状態のみの条件が必要環境要因として $E$ に含まれうる。
そのような $\eta$ では $D_\eta=\{x\in X\mid\rho(x)=1\}$ であり、
$\eta(x_{t_0},e_{t_0})=0$ は $x_{t_0}\notin D_\eta$ を意味するから、
(CR-1c) は成立しない。

(CR-1c) の証拠列においては
$x_{t_0}\in D_\eta$ かつ $\eta(x_{t_0},e_{t_0})=0$ であるから、
ある $e'\in E_{\mathrm{ext}}$ について $\eta(x_{t_0},e')=1$ であり、
\[
  \exists x\in X,\ \exists e,e'\in E_{\mathrm{ext}}:\quad
  \eta(x,e)\neq\eta(x,e')
\]
が成立する。
すなわち (CR-1c) は、充足判定が外部環境状態に依存することを含意する。
また (CR-1c) は (CR-1a)(i) の存在要求を含意する。
(CR-1a)(i) を独立に残すのは、(CR-1a)(ii) の量化範囲を、
開始時点の内部状態が $D_\eta$ に属する非充足に限らず、
$D_\eta$ に属さない内部状態で開始される非充足を含む
持続的な非充足の全体とするためである。

(CR-1c) は、$\eta$ が内部状態に依存することを排除しない
（定義~\ref{def:environment_space}、備考~\ref{rem:env_utilization_function}）。
また本条件は、外部環境状態への依存の程度、および上記の意味の非充足が
開始されうる内部状態の範囲を定量化しない。
外部依存が一部の内部状態においてのみ現れる要因をどこまで
必要環境要因として認めるか、およびその程度の定量化は
Phase~2 の課題とする。

CR-1 は、当該7タプル構造体を本論文におけるACSのクラスへ
属させるための条件であり、
個々の時刻における診断値でも、後続命題の成立範囲を限定する
仮定でもない（備考~\ref{rem:math_vs_realization}）。
CR-1 の充足を有限観測から判定する手続きは Phase~1 では与えず、
その定式化は Phase~2 の課題とする
（備考~\ref{rem:necessity_identification}）。

\begin{remark}[(CR-1a)(i) の外部背景相対性]
\label{rem:cr1_relativity}
(CR-1a)(i) は、$\eta$ の持続的非充足が状態遷移列として
実現しうることを要求するため、
外部環境状態空間 $E_{\mathrm{ext}}$ の取り方に相対的である。
$E_{\mathrm{ext}}$ に $\eta$ の非充足を伴う環境状態が
含まれない場合、当該 $\eta$ は (i) を満たさず、
その記述の下では $E$ の元とならない。
(CR-1c) も同様に $E_{\mathrm{ext}}$ の取り方に相対的であり、
内部状態が $D_\eta$ に属する時点で当該要因を非充足とする
外部環境状態を含む状態遷移列が実現しえない場合、
当該 $\eta$ は (CR-1c) を満たさない。

これは、ACSの7タプルが $E_{\mathrm{ext}}$ に相対的に定義される
という本論文の枠組み（\S\ref{subsec:wellformedness}）と整合する。
必要環境要因の同定が外部背景の取り方にどのように依存するかの
検討は Phase~2 の課題とする。
\end{remark}

\begin{remark}[(CR-1a)(ii) の量化の読み方]
\label{rem:cr1_quantification}
(CR-1a)(ii) の前件は、$\eta$ の非充足が時刻 $t_0$ 以降のすべての時刻に
わたって成立することを述べる。
$\eta$ は内部状態 $x_s$ と外部環境状態 $e_s$ の双方を引数に取るため、
この前件は次の二つの場合をともに排除している。
第一に、内部の蓄えその他により $\eta(x_s,e_s)=1$ が回復する場合。
第二に、行動 $u_s$ の選択を通じて後続の
$(x_{s'},e_{s'})$ が変化し $\eta(x_{s'},e_{s'})=1$ が回復する場合。
したがって、内部補償および許容行動による回復可能性は
前件の側に吸収されており、
(CR-1a)(ii) はこれらによって回復しない非充足のみを対象とする。

後件は、ある有限時刻以降すべての時刻で
$\mathrm{WF}_{i,s}=0$ となることを述べる。
時刻 $T$ の値は要因ごとに、また状態遷移列ごとに異なりうる。
本条件は $T$ の一様な上界を与えず、
$T=t_0$ となる場合、すなわち非充足の開始時点で
ただちに実現が失われる場合も排除しない。
\end{remark}

\begin{remark}[必要条件であって十分条件ではない]
\label{rem:cr1_one_way}
CR-1 が述べるのは、必要環境要因の回復しない非充足から
整形式実現の非回復への一方向の含意である。
必要環境要因の充足が整形式実現の継続に十分であること、すなわち
\[
  \forall \eta\in E:\ \eta(x_t,e_t)=1
  \;\Longrightarrow\;
  \mathrm{WF}_{i,t+1}=1
\]
を、本論文は主張しない。
必要環境要因がすべて充足されている場合であっても、
内部構造の損壊、行動の選択、外部からの作用その他により
整形式実現が失われる場合を排除しない。

したがって $E$ は、当該ACSの存続に関する必要条件の側を記述する
構成成分であり、存続の十分条件を与えるものではない。
\end{remark}

\begin{remark}[一時的な非充足と持続的な欠損]
\label{rem:cr1_timing}
ある時刻において
$\eta(x_t,e_t)=0$
が成立することから、本論文は、それ自体では
$\mathrm{WF}_{i,t}=0$ も
$\mathrm{WF}_{i,t+1}=0$ も一般的な帰結として置かない。
(CR-1a)(ii) の前件は時刻 $t_0$ 以降の持続を要求しており、
単発または断続的な非充足はこれを満たさない。

内部の蓄え、回復行動、および外部環境状態の変化によって
非充足が解消される場合、当該状態遷移列は (CR-1a)(ii) の前件を
満たさないため、後件についての帰結も生じない。
一時的な非充足がどの程度まで許容されるか、
すなわち回復余裕の定量化は Phase~1 では扱わず、
Phase~2 の課題とする（\S\ref{subsec:limitations_list} の L9）。
\end{remark}

\begin{definition}[必要環境充足集合（分類C）]
\label{def:E_req}
状態 $x\in X$ に対し、
$E$ のすべての必要環境要因が充足される外部環境状態の集合を
\[
  E_{\mathrm{req}}(x)
  :=
  \bigl\{
    e\in E_{\mathrm{ext}}
    \;\bigm|\;
    \forall \eta\in E,\ \eta(x,e)=1
  \bigr\}
  \subseteq
  E_{\mathrm{ext}}
\]
と定める。
これにより (CR-1b) は
\[
  \exists x\in X:\ E_{\mathrm{req}}(x)\neq\emptyset
\]
と書ける。

$E_{\mathrm{req}}$ は $X$ から $E_{\mathrm{ext}}$ への対応であり、
$E$、$X$ および $E_{\mathrm{ext}}$ から定まる派生的対象である。
定義式は状態遷移写像 $f$ を用いない。
ただし、$E$ の各元がなぜ必要と呼ばれるかを拘束する構成要件 CR-1 自体は、
状態遷移列と $\mathrm{WF}$ の推移を参照するため $f$ に依存する。
$E_{\mathrm{req}}$ は7タプルの独立成分ではない。
\end{definition}

\begin{remark}[$E_{\mathrm{req}}$ の性質と他の環境集合との区別]
\label{rem:E_req_properties}
定義~\ref{def:acs}(3) のとおり、Phase~1 では $\eta$ に
連続性その他の正則性を課さない。
したがって $E_{\mathrm{req}}(x)$ は一般に $E_{\mathrm{ext}}$ の
閉部分集合でもコンパクト部分集合でもない。
例えば $E_{\mathrm{ext}}=[0,1]$ に通常の位相を入れ、$E=\{\eta\}$、
$\eta(x,e)=1\iff e<1/2$ とすれば、$E_{\mathrm{req}}(x)=[0,1/2)$ は閉でなく、コンパクトでもない
（型の水準の例）。
これらの性質を要する命題では、その箇所で追加仮定として明示する。

$E_{\mathrm{req}}(x)$ は、必要環境要因の充足のみによって定まる集合であり、
当該環境状態のもとで可行な行動が存在するか否かを述べない。
一段階可行性に関する環境集合は
\S\ref{subsec:viability_function_and_blockage} で別に定義し、
両者を区別する。
すなわち、必要環境要因がすべて充足されている外部環境状態のもとでも
可行行動が存在しない場合、および
必要環境要因が充足されていない外部環境状態のもとでも
一段階の可行行動が存在する場合を、いずれも排除しない。
\end{remark}

\begin{remark}[要因分解の非一意性]
\label{rem:factor_nonuniqueness}
$E$ は、必要環境要因への分解の粒度まで一意に定まるとは限らない。
例えば $\eta_1,\eta_2\in E$ に対して連言
\[
  \eta_\wedge(x,e)
  :=
  \eta_1(x,e)\wedge\eta_2(x,e)
\]
を作り
\[
  E'
  :=
  \bigl(E\setminus\{\eta_1,\eta_2\}\bigr)\cup\{\eta_\wedge\}
\]
とすると、$E'$ は $E$ と同一の必要環境充足集合
$E_{\mathrm{req}}(x)$
を与え、CR-1 についても $E$ と同じ地位をもつ場合がある。

したがって本論文は、$E$ が要因分解の粒度まで一意であることを
要求しない。
Phase~1 では、対象とするACS個体について一つの必要環境要因分解 $E$ を
固定し、同一の必要環境充足集合を与える別の粗視化または細分化表現の
存在を排除しない。
分解の一意性を与える追加条件の検討は Phase~2 の課題とする。
\end{remark}

\begin{remark}[必要性の存立と同定可能性]
\label{rem:necessity_identification}
本論文は、必要環境要因が当該ACS個体について定まることを
観測者の知識から独立に扱う。
すなわち、ある $\eta$ が CR-1 の意味で必要であることは、
それが観測または記述されているか否かに依存しない。

一方で本論文は、未知のACS個体について $E$ を有限観測から
同定する手続きを与えない。
異なる外部環境状態のもとでの応答の観測、または内部モデルによる
予測は、その手がかりとなりうるが、
本論文はこれらから $E$ を復元する規則を提示しない。
必要環境要因の経験的同定、および同定誤差の扱いは
Phase~2 の課題とする。
この意味で、必要性の存立に関する記述と、
必要性の同定可能性に関する記述とを区別する。
\end{remark}

%----------------------------------------------------------------------
\subsection{状態遷移・行動・外部環境入力}
\label{subsec:state_space_basic_variables}

本小節では、ACSの離散時間動学を記述する。
定義~\ref{def:acs}(4) で型を宣言した状態遷移写像 $f$ による
一ステップの状態発展を定め、あわせて動学の記述に用いる
基本仮定を導入する。
これらの仮定は、\S\ref{subsec:acs_7tuple} の基本型宣言から
前方参照されているものを含む。
また、行動選択を経ない自律的成分と行動選択に媒介される
指向的成分への分解、および外部環境入力に対する
動学的非退化性を扱う。

必要環境空間 $E$ に関する意味論は
\S\ref{subsec:environment_space} で与えた。
本小節が扱う状態遷移写像は $E_{\mathrm{ext}}$ 全体の上で
定義されており、必要環境要因が充足されているか否かに
かかわらず、同一の型のもとで状態発展を記述する。

本小節で新たに用いる記号は、自律的機能過程
$f_{\mathrm{auto}}$、指向的過程 $f_{\mathrm{dir}}$、
結合演算 $\oplus$、および一ステップ差分 $\Delta Y_t$ である。
下付き $\mathrm{auto}$ は autonomous（自律的）、
$\mathrm{dir}$ は directional（指向的）に対応する。
ここでいう自律性は、状態遷移成分が行動選択に依存しないことを指し、
\S\ref{sec:tipa} で扱う自動性（automaticity）とは別の概念である。
後者は機構の作動形態に関する概念であり、
本小節の $f_{\mathrm{auto}}$ と対応関係をもたない。
記号 $\oplus$ は二つの状態値を結合する二項演算であり、
群または環の加法を表すものではない。

\subsubsection{空間のコンパクト性と連続性}

\begin{assumption}[$X$・$U$ のコンパクト性]
\label{ass:compact_spaces}
状態空間 $X\subseteq\mathbb{R}^n$ および
行動空間 $U\subseteq\mathbb{R}^m$ はコンパクト
（閉かつ有界）である。

\medskip
\noindent
導入地位：S（構造的）。強さ：強。
本仮定は、$X$ および $U$ 上の閉部分集合のコンパクト性、
連続実数値関数の極値存在、および
可行行動集合 $A_t$ のコンパクト性
（命題~\ref{prop:At_compact}）等を扱うための
構造的基盤を与える。
コンパクト性・閉集合・連続写像に関する標準的事実は、
一般位相の通常の結果に従う\cite{Munkres2000}。
本仮定は外部環境状態空間 $E_{\mathrm{ext}}$ には課さない。
Phase~2 では非コンパクト空間または局所コンパクト空間への
緩和を検討する。
\end{assumption}

\begin{assumption}[$f$ の連続性]
\label{ass:f_continuity}
状態遷移写像
\[
  f:X\times U\times E_{\mathrm{ext}}\to X
\]
は連続である。

\medskip
\noindent
導入地位：S（構造的）。強さ：中。
本仮定は、状態遷移に関する閉性と極限操作の扱いを
安定させるための構造的仮定である。
特に、可行集合 $K$ の閉性とあわせて、
可行行動集合 $A_t$ の閉性を導くために用いられる
（命題~\ref{prop:At_compact}）。
連続性は $E_{\mathrm{ext}}$ の位相
（\S\ref{subsec:wellformedness}）に関して述べられる。
必要環境要因 $\eta$ には連続性その他の正則性を課していないため
（定義~\ref{def:acs}(3)）、本仮定は
$E_{\mathrm{req}}(x)$ の開性・閉性その他の位相的性質を規定しない。
\end{assumption}

\begin{assumption}[$E_{\mathrm{ext}}$ の時刻固定性]
\label{ass:Eext_time_fixed}
Phase~1 において、外部環境状態空間 $E_{\mathrm{ext},t}$ は
時刻によらず同一の位相空間 $E_{\mathrm{ext}}$ として扱う。

\medskip
\noindent
導入地位：S（構造的）。強さ：中。
これは Phase~1 の決定論的・離散時間枠組みを単純化するための
構造的仮定である。
本仮定は、各時刻の実現外部環境状態 $e_t$ が時刻により
変化することを排除しない。
時間依存する外部環境状態空間への拡張は Phase~2 の課題とする。
\end{assumption}

\subsubsection{状態遷移写像}

\begin{definition}[状態遷移]
\label{def:state_transition}
ACS個体 $\Sigma_i$ の状態は離散時間で以下のように
発展する：
\[
  x_{t+1} = f(x_t, u_t, e_t).
\]
$x_t\in X$ は時刻 $t$ の状態ベクトル、
$u_t\in U$ は実効行動、
$e_t\in E_{\mathrm{ext}}$ は実現外部環境状態である。
この離散時間状態遷移は、離散時間制御系の標準的記法と
整合する\cite{Bertsekas2017}。

時刻に依存する実数値またはベクトル値の量 $Y_t$ について、
その一ステップ差分を
\[
  \Delta Y_t:=Y_{t+1}-Y_t
\]
と定義する。
この記法は、減算が定義されない集合値・分類値・関係には
適用しない。
状態空間 $X$ 自体は集合であり、時間差分の対象とはしない。

離散時間定式化は、Aubin の Viability Theory における
可行集合・可行行動という発想と概念的に整合する\cite{Aubin1991}。
ここでいう整合とは、数学的同値性や包含関係ではなく、
可行性保存制約を中心に状態遷移を記述するという点での
概念的矛盾の不在を意味する。
連続時間への拡張は Phase~2 の課題とする。
\end{definition}

\begin{remark}[状態遷移写像の読みの規約]
\label{rem:f_reading}
本論文では、$u\in U$ をモデル上の行動入力とし、
$f(x,u,e)$ を、内部状態 $x$、行動入力 $u$ および外部環境状態 $e$ のもとで
定まる次状態と読む。
この読みのもとでは、行動の試行、行動を実行できなかった場合、
および実行の成功の度合いを、$f$ の記述の中で互いに区別しない。
Phase~1 では、実行の成否または実行の度合いを表す独立の変数を導入しない。
それらの差異をモデルに含める場合には、状態 $x$、外部環境状態 $e$、
行動入力 $u$ の取り方、または $f$ の定義へ吸収する。
固定された $(x,u,e)$ に対しては、写像 $f$ が一つの次状態を与える。
定義~\ref{def:acs}(2) の実効行動 $u_t$ も、現実に実行が成功したことを意味せず、
内部の指令が解決された後に $f$ へ入力されるモデル上の行動入力をいう。
行動入力の実行可能性を $U$ とは別に表す構造は Phase~1 では導入しない。
この規約によって表現されない区別は、
\S\ref{sec:limitations_and_outlook} で理論的限界として扱う。
\end{remark}

\begin{definition}[行動応答成分の遷移写像と環境実現]
\label{def:env_realization}
定義~\ref{def:Eext_decomposition} の分解のもとで、全域・一価の写像
\[
  g_O:O\times X\times U\times E^{\mathrm{bg}}\longrightarrow O
\]
を\emph{行動応答成分の遷移写像}と呼び、
行動応答成分の実現値は
\[
  o_{t+1}=g_O(o_t,x_t,u_t,b_t)
\]
に従って発展するものとする。
$g_O(o,x,u,b)$ の第 $a$ 成分を $g_O(o,x,u,b)_a$ と書く。
外生背景成分の列 $(b_t)_{t\geq0}$ は外生的に与えられる列とする。
四つ組 $\bigl(E^{\mathrm{bg}},O,g_O,(b_t)_{t\geq0}\bigr)$ を、
当該ACSの\emph{環境実現}と呼ぶ。

$g_O$ はACS 7タプルの成分ではなく、
解析に用いる環境の側の構造である。
ACSが $g_O$ を内部に保持すること、または計算することは要求しない。
本論文は $g_O$ に連続性その他の正則性を課さない。
\end{definition}

定義~\ref{def:state_transition} と定義~\ref{def:env_realization} から、
初期値 $(x_0,o_0)$、背景成分の列 $(b_t)_{t\geq0}$ および実効行動の列
$(u_t)_{t\geq0}$ が与えられれば、
\[
  x_{t+1}=f\bigl(x_t,u_t,(b_t,o_t)\bigr),
  \qquad
  o_{t+1}=g_O(o_t,x_t,u_t,b_t)
\]
によって $(x_t,o_t)_{t\geq0}$ が一意に定まる。
$\mathcal A_O=\emptyset$ のとき $g_O$ は一点集合への定値写像であり、
実現外部環境状態の列は背景成分の列だけで定まる。
これは外部環境状態の列全体を外生的に与える場合に一致する。

\begin{remark}[実現外部環境状態の生成について]
\label{rem:exogenous_environment}
定義~\ref{def:state_transition} は、実現外部環境状態の列
$(e_t)_{t\geq0}$ の生成規則を7タプルの内部に置かない。
Phase~1 では、外生的に与えるのは背景成分の列 $(b_t)_{t\geq0}$ だけとし、
行動応答成分は定義~\ref{def:env_realization} の $g_O$ に従って
当該ACSの行動に応答しうるものとする。
したがって、ACSの軌道は7タプル $\Sigma$ だけからは定まらず、
$\Sigma$ と環境実現に相対して定まる。
この相対性は、外部環境状態の列全体を外生的に与える場合にも同じである。

背景成分の列の生成則、
他のACS個体を自身の遷移写像をもつ個体として同時に扱う結合動学、
および確率的な環境は Phase~2 の課題とする（\S\ref{subsec:limitations_list} の L2）。
他のACS個体の振る舞いを行動応答成分として $g_O$ に含める記述は、
環境実現の側のモデルであり、その個体自身の遷移写像を与えるものではない。
当該ACSと外部対象との関係構造は \S\ref{sec:rel_min} で扱う。
\S\ref{sec:rel_min} は関係の記述構造を与えるものであり、
外部環境状態の遷移則を与えるものではない。
\end{remark}

\begin{definition}[作用差（分類C）]
\label{def:action_difference}
$(o,x,b)\in O\times X\times E^{\mathrm{bg}}$、行動 $u\in U$ および枠 $a\in\mathcal A_O$ について、
\[
  g_O(o,x,u,b)_a\neq g_O(o,x,u_{\emptyset},b)_a
\]
が成り立つとき、行動 $u$ は $(o,x,b)$ において枠 $a$ に\emph{作用差}をもつという。
ここで $u_{\emptyset}$ は選択的待機である（備考~\ref{rem:action_space_structure}）。
\end{definition}

定義から直ちに、選択的待機 $u_{\emptyset}$ はどの $(o,x,b)$ においても
どの枠にも作用差をもたない。
また $\mathcal A_O=\emptyset$ のとき、作用差は定義されない。

\begin{remark}[作用差の射程]
\label{rem:action_difference_scope}
作用差は、待機と比べた一段階の値の違いの有無だけを表す。
次の限定を付す。
\begin{enumerate}[label=(\alph*)]
  \item 作用差は、成分 $O_a$ の取り方（どの区別を値として表すか）に相対的である。
  同じ物理的過程でも、$O_a$ の値が区別しない違いは作用差として現れない。
  したがって、作用差がないことは、物理的な接触や影響がないことを意味しない。
  \item 作用差は、比較の基準を $u_{\emptyset}$ に固定した量である。
  他の行動と比べた違いは、本定義の対象ではない。
  \item 作用差の大きさ（距離）は定めない。
  \item 作用差は、行動が結果を引き起こしたという因果的寄与の判定ではない。
  結果が他の原因によっても同時に生じる場合（過剰決定）や、
  他の原因が先に結果を生じさせる場合（先取り）の区別は与えない
  （\S\ref{sec:limitations_and_outlook}）。
  \item 本論文は、$u\in U_{\mathrm{ext}}$ であることと
  作用差をもつことの間に一般的な含意を置かない。
  距離、身体条件その他の事情により、外部へ向けた行動が作用差をもたない場合を排除しない。
\end{enumerate}
\end{remark}

\begin{remark}[自己の側への取り込みの表現]
\label{rem:intake_joint_transition}
外部の物や量が当該ACSの自己の側へ移ること（取り込み）は、
本論文では、型をまたいで同一の数学的対象が移動することとしては表さず、
同一の一段階における二つの遷移の対として表す。
すなわち、ある枠 $a$ について $o_{t,a}\neq\bot_a$ かつ
$g_O(o_t,x_t,u_t,b_t)_a\neq o_{t,a}$（値が $\bot_a$ となる場合を含む）であることと、
$f(x_t,u_t,e_t)$ が $o_{t,a}$ に依存しうることの組として表す。
規約 TSC は維持され、物質の同一性を $X$ と $O$ の型をまたいで追跡しない。
必要環境要因 $\eta$ の観点からは、
取り込みの前には $\eta(x_t,e_t)=1$ が外部側の状態を通じて成り立ち、
取り込みの後には同じ要因の充足が内部状態 $x_{t+1}$ を通じて成り立つ、
という読みが可能である（備考~\ref{rem:env_utilization_function}）。
本備考は新たな写像、変数、構成要件または仮定を導入しない。
取り込みの型の分類、および外部と内部に共通する資源量の収支は
Phase~1 では扱わない。
\end{remark}

\subsubsection{状態遷移の二成分分解}

\begin{assumption}[状態遷移の二成分分解]
\label{ass:f_decomposition}
Phase~1 で扱う状態遷移写像 $f$ は、自律的機能過程
$f_{\mathrm{auto}}$ と指向的過程 $f_{\mathrm{dir}}$、
および連続な二項演算 $\oplus$ により、
\[
  f(x_t,u_t,e_t)
  =
  \oplus\bigl(
    f_{\mathrm{auto}}(x_t,e_t),
    f_{\mathrm{dir}}(x_t,u_t,e_t)
  \bigr)
\]
と表される。

ここで
\[
  f_{\mathrm{auto}}:X\times E_{\mathrm{ext}}\to X,
\]
\[
  f_{\mathrm{dir}}:X\times U\times E_{\mathrm{ext}}\to X,
\]
\[
  \oplus:X\times X\to X
\]
であり、$\oplus$ は連続である。

$f_{\mathrm{auto}}$ は行動選択に依存せず、
内部状態と外部環境状態によって定まる自律的成分を表し、
$f_{\mathrm{dir}}$ は実効行動 $u_t$ を引数として受け、
これに依存しうる指向的成分を表す。
本仮定は、$f_{\mathrm{dir}}$ が実際に $u_t$ に依存すること、
すなわち $u\neq u'$ に対して
$f_{\mathrm{dir}}(x,u,e)\neq f_{\mathrm{dir}}(x,u',e)$ となる組が
存在することを要求しない。

さらに、自律的機能過程は定数写像ではないとする。
すなわち
\[
  \exists (x,e),(x',e')\in X\times E_{\mathrm{ext}}:
  \quad
  f_{\mathrm{auto}}(x,e)\neq f_{\mathrm{auto}}(x',e')
\]
が成立する。

選択的待機 $u_{\emptyset}$（備考~\ref{rem:action_space_structure}）
については、任意の $x\in X$ および
$e\in E_{\mathrm{ext}}$ に対し、
\[
  \oplus\bigl(
    f_{\mathrm{auto}}(x,e),
    f_{\mathrm{dir}}(x,u_{\emptyset},e)
  \bigr)
  =
  f_{\mathrm{auto}}(x,e)
\]
が成立する。

\medskip
\noindent
導入地位：S（構造的）。強さ：強。
一般の連続写像 $f$ がこの二成分分解をもつとは限らないため、
本分解は Phase~1 の構造的仮定として明示する。
これは、行動選択を経ない自律的成分と、
行動選択に媒介される指向的成分を数学的記述上で分離し、
選択的待機
\[
  f(x,u_{\emptyset},e)
  =
  f_{\mathrm{auto}}(x,e)
\]
を記述するための仮定である。

ただし本仮定は、閉塞状態において
$f_{\mathrm{auto}}$ に対応する実装上の自律的機能過程が
実際に作動し続けること、または
実効的時間地平 $\widehat{\tau}_{t}>0$ が成立することを
規定しない（備考~\ref{rem:blockage_auto}）。
また本仮定は、$f_{\mathrm{auto}}$ の作動時に
$\mathrm{WF}_{i,t}=1$ が保たれることを要求しない。
$f_{\mathrm{auto}}$、$f_{\mathrm{dir}}$、$\oplus$ の具体的
実装形式は Phase~2 以降の課題とする。
\end{assumption}

\begin{remark}[$\oplus$ と選択的待機の代数的地位]
\label{rem:oplus_neutrality}
仮定~\ref{ass:f_decomposition} は、$\oplus$ に対して
結合律・交換律・大域的単位元の存在等の一般的な代数構造を
要求しない。
選択的待機 $u_{\emptyset}$ について必要なのは、
各 $(x,e)\in X\times E_{\mathrm{ext}}$ に対して
\[
  \oplus\bigl(
    f_{\mathrm{auto}}(x,e),
    f_{\mathrm{dir}}(x,u_{\emptyset},e)
  \bigr)
  =
  f_{\mathrm{auto}}(x,e)
\]
が成立することのみであり、
$f_{\mathrm{dir}}(x,u_{\emptyset},e)$ が
$\oplus$ の大域的単位元であることを要求しない。
したがって、$u_{\emptyset}$ の中立性は
$f_{\mathrm{auto}}$ の値に対する局所的な条件として扱う。
\end{remark}

\subsubsection{外部環境入力の非退化性}

\begin{assumption}[外部環境入力に対する動学的非退化性]
\label{ass:env_input_nondegeneracy}
状態遷移写像 $f$ は外部環境入力に対して非退化である。
すなわち
\[
  \exists x\in X,\
  \exists u\in U,\
  \exists e,e'\in E_{\mathrm{ext}},\
  e\neq e':
  \quad
  f(x,u,e)\neq f(x,u,e')
\]
が成立する。

\medskip
\noindent
導入地位：S（構造的）。強さ：弱。
本仮定は、外部環境状態の差異が状態遷移に一切影響しない
退化的な系を Phase~1 の対象から除くためのものである。
本仮定は、環境差が可行性または可行行動集合に影響することを
述べるものではない。
可行行動集合の水準における非退化性は
\S\ref{subsec:viability_function_and_blockage} で扱う。
\end{assumption}

\begin{remark}[非退化条件の水準の区別]
\label{rem:nondegeneracy_levels}
本論文では、外部環境入力に関する非退化性を三つの水準で
区別する。
第一は本小節の仮定~\ref{ass:env_input_nondegeneracy} が
述べる状態遷移の水準であり、
環境差が次状態に影響しうることを要求する。
第二は \S\ref{subsec:viability_function_and_blockage} で扱う
可行行動集合の水準であり、
環境差が可行行動の集合に影響しうることを要求する。
第三は \S\ref{subsec:three_primitive_functions} で扱う
感知の水準であり、
環境差が感知出力に影響しうることを要求する。

可行行動集合の水準における非退化性は
仮定~\ref{ass:env_input_nondegeneracy} を含意する
（命題~\ref{prop:nondegeneracy_implication}）。
逆は一般には成立しない
（備考~\ref{rem:nondegeneracy_converse}）。
感知の水準については、本論文は、別途接続条件を置かない限り、
他の二つの水準との間に相互含意を仮定しない。

いずれの水準の非退化性も、必要環境空間 $E$ の
構成要件 CR-1（\S\ref{subsec:environment_space}）とは
別の条件である。
本論文はこれらの非退化性から必要環境要因の存在を一般的な帰結として置かず、
また CR-1 からこれらの非退化性を一般的な帰結として置かない。
\end{remark}

%----------------------------------------------------------------------
\subsection{持続評価・可行性・一段階環境可行性}
\label{subsec:viability_function_and_blockage}

本小節では、状態の可行性を評価する持続評価関数 $\Pi$ を定義し、
そこから可行集合 $K$ と可行行動集合を導く。
これらの可行性概念は、制御系における制御不変集合の発想と
接続する\cite{Blanchini1999}。
さらに、可行行動が一つも存在しない状態を閉塞として定義し、
一段階可行性を許す外部環境状態の集合を派生的に定める。

崩壊は閉塞とは区別される事象であり、その定義は
\S\ref{subsec:self_formal_definition}（定義~\ref{def:collapse}）で与えた。
本小節では、本論文が閉塞から整形式実現の不成立
（$\mathrm{WF}_{i,t}=0$、定義~\ref{def:wf_realization}）も
崩壊も一般的な帰結として置かないことを明示する。

本小節の可行性概念は、必要環境要因の充足
（\S\ref{subsec:environment_space}）とは別の対象である。
必要環境充足集合 $E_{\mathrm{req}}(x)$ が必要環境要因の充足のみで
定まるのに対し、本小節で定める一段階可行環境対応は
可行行動の存在によって定まる。
両者の関係は備考~\ref{rem:fixed_environment_and_state_dependent_viability}
で整理する。

本小節で新たに用いる記号は、持続評価関数 $\Pi$、閾値ベクトル
$\theta$、可行集合 $K$、可行行動対応 $A(x,e)$、
時刻 $t$ の可行行動集合 $A_t$、および一段階可行環境対応
$\mathcal{E}_{\mathrm{viab}}^{(1)}(x)$ である。
上付き $(1)$ は一段階を表す識別子であって冪ではない。
記号 $\rightrightarrows$ は集合値対応を表す。

\subsubsection{持続評価関数と可行集合}

\begin{definition}[持続評価関数]
\label{def:Pi}
閾値ベクトル $\theta = (\theta_1, \ldots, \theta_n)
\in\mathbb{R}^n$ を所与とする。
\emph{持続評価関数} $\Pi : X \to\mathbb{R}$ を
以下で定義する：
\[
  \Pi(x) := \min_{k\in\{1,\ldots,n\}}
  \bigl(x_k - \theta_k\bigr),
  \qquad
  x\in X.
\]
$x_k$ は状態ベクトル $x$ の第 $k$ 成分、
$\theta_k$ はその次元の可行閾値である。
すなわち $\Pi(x)$ は、各次元の余裕
$x_k - \theta_k$ のうち最小のものを与える
（最弱次元による評価；備考~\ref{rem:min_type_pi}）。
時刻 $t$ の状態への適用は
$\Pi(x_t)=\min_k(x_{k,t}-\theta_k)$ と書く。
Phase~1 では各 $\theta_k$ を時間不変の定数として扱う。
動的閾値への拡張は Phase~2 以降の課題とする。
\end{definition}

\begin{remark}[$\min$ 型 $\Pi$ の採用理由]
\label{rem:min_type_pi}
各次元の閾値 $\theta_k$ が、その次元の最低可行水準、
すなわちこれを下回ると他次元の状態にかかわらず可行性が
損なわれる水準として定義される場合、
可行性判定には各次元の条件
\[
  x_k\geq\theta_k
  \qquad (k=1,\ldots,n)
\]
が同時に要求される。
この非補償的な閾値規約のもとでは、
\[
  \Pi(x) = \min_{k\in\{1,\ldots,n\}} (x_k-\theta_k)
\]
は、最も余裕の小さい次元によって可行性を評価する形式となり、
各次元の最低可行水準を直接的に表す。
加重線形型や一部のノルム型のように次元間の補償を許す形式は、
本小節で採用する各 $\theta_k$ の非補償的な意味とは整合しない。
したがって Phase~1 では、この閾値規約を直接表現する
基礎形式として $\min$ 型の $\Pi$ を採用する。
ただし、同一の可行集合 $K$ を $0$ 以上の水準集合として与える
他の連続関数表現の存在を排除するものではない
（命題~\ref{prop:theta_from_pi}；under stated assumptions）。
\end{remark}

\begin{definition}[可行集合]
\label{def:K}
$\Pi$（定義~\ref{def:Pi}）から、\emph{可行集合} $K$ を
以下の派生定義として構成する：
\[
  K := \{x\in X \mid \Pi(x)\geq 0\}.
\]
$K$ は $X$ の閉部分集合である（$\Pi$ の連続性と
$\{x\in\mathbb{R}\mid x\geq 0\}$ の閉性による）。
\end{definition}

\begin{proposition}[可行性の等価性と $K$ の導出可能性]
\label{prop:theta_from_pi}
(i)\ $\Pi$ の定義（定義~\ref{def:Pi}）のもとで、
\[
  \forall x\in X:
  \quad
  \Pi(x)\geq 0 \iff x\in K .
\]
(ii)\ $K\neq\emptyset$ および $K$ の閉性のもとで、
$\Pi$ の定義形式に依らず $K$ は適切な連続関数
$\Pi : X \to\mathbb{R}$ の $0$ 以上の水準集合として
表現できる。
したがって、(ii) は空でない閉集合 $K$ が連続実数値関数の $0$ 以上の水準集合として
表現できることを示す。
ただし、当該関数が定義~\ref{def:Pi} の $\min$ 型で表せることまでは主張しない。
\end{proposition}

\begin{proof}
(i)：任意の $x\in X$ について
$\Pi(x) = \min_k(x_k - \theta_k)\geq 0
\iff \forall k,\ x_k\geq\theta_k
\iff x\in K$。

\smallskip
(ii)：$K\neq\emptyset$ のもとで、距離関数
$\operatorname{dist}(x, K)$ は $X$ 上で実数値をとり
well-defined である。
さらに $K\neq\emptyset$ より
$x\mapsto\operatorname{dist}(x, K)$ は $X$ 上で $1$-Lipschitz であり、
したがって連続である。
ゆえに距離型関数
$\Pi_d(x) := -\operatorname{dist}(x, K)$ も $X$ 上連続である。
$K$ の閉性より $\operatorname{dist}(x, K)=0 \iff x\in K$ であり、
$\Pi_d(x)\leq 0$ とあわせて
$K = \{x \mid \Pi_d(x)\geq 0\}$ が成立する。
\end{proof}

\subsubsection{可行行動集合と閉塞}

\begin{definition}[可行行動対応と可行行動集合]
\label{def:at}
状態 $x\in X$ と外部環境状態 $e\in E_{\mathrm{ext}}$ に対し、
\emph{可行行動対応}
\[
  A:X\times E_{\mathrm{ext}}\rightrightarrows U
\]
を
\[
  A(x,e) := \{u\in U \mid f(x, u, e)\in K\}
\]
によって定義する。
時刻 $t$ における\emph{可行行動集合}を
\[
  A_t := A(x_t,e_t)
\]
と書く。
$A(x,e)$ は時刻に依存しない対応であり、
$A_t$ はその時刻 $t$ の実現値である。
個体を明示する場合には、個体 $\Sigma_i$ の成分を用いて、対応
\[
  A_i:X_i\times E_{\mathrm{ext},i}\rightrightarrows U_i,
  \qquad
  A_i(x,e):=\{u\in U_i\mid \Pi_i\bigl(f_i(x,u,e)\bigr)\geq0\}
\]
と、その時刻評価値
\[
  A_{i,t}:=A_i(x_{i,t},e_{i,t})
\]
を書く（\S\ref{subsec:acs_7tuple} の記述水準の規約）。
$\Pi_i(f_i(x,u,e))\geq0$ は、$\Pi_i$ から定義~\ref{def:K} と同じ規則で定まる
可行集合に $f_i(x,u,e)$ が属することと同値である。
単一ACSのみを扱う文脈では $A_i=A$、$A_{i,t}=A_t$ である。

$A(x,e)$ は、状態 $x$ と外部環境状態 $e$ のもとでの
一ステップ遷移後の次状態を $K$ に属させる行動の全体である。
特に $x\in K$ の場合には、次状態を $K$ 内に留める行動の全体を表す。
$A$ を集合値対応として読む場合は、
集合値解析の枠組みと接続する\cite{AubinFrankowska1990}。
\end{definition}

\begin{remark}[可行行動集合の読み]
\label{rem:at_reading}
定義~\ref{def:at} の可行行動対応は、本備考によって変更しない。
本論文では、$A(x,e)$ を、状態 $x$ と外部環境状態 $e$ のもとで
次状態を可行集合 $K$ に属させる行動入力の全体、
すなわち当該の $(x,e)$ における\emph{構造的応答ポテンシャルの集合}と読む。
この読みは記述上の呼称である。
この読みも定義~\ref{def:at} も、
$A(x,e)$ の元が実際に選択されること、
および複数時刻にわたって可行性が維持されることを要件として含まない。
選択的待機 $u_{\emptyset}$（備考~\ref{rem:action_space_structure}）も $U$ の元であり、
$f(x,u_{\emptyset},e)\in K$ であれば $u_{\emptyset}\in A(x,e)$ である。
行動の選択と可行行動集合との関係は \S\ref{sec:three_principles} で扱う。
\end{remark}

\begin{proposition}[可行行動集合の閉性とコンパクト性]
\label{prop:At_compact}
仮定~\ref{ass:compact_spaces}（$U$ のコンパクト性）および
仮定~\ref{ass:f_continuity}（$f$ の連続性）のもとで、任意の
$x\in X$ および $e\in E_{\mathrm{ext}}$ に対し、可行行動対応の値
$A(x,e)$ は $U$ の閉部分集合であり、したがってコンパクトである
（$A(x,e)=\emptyset$ の場合を含む）。
\end{proposition}

\begin{proof}
写像 $u\mapsto f(x,u,e)$ は $f$ の連続性より $U$ 上連続であり、
$K$ は $X$ の閉部分集合である（定義~\ref{def:K}）。
したがって
\[
  A(x,e)=\{u\in U \mid f(x,u,e)\in K\}
\]
は、閉集合 $K$ の連続写像による逆像として $U$ において閉である。
$U$ はコンパクトであるから、その閉部分集合である $A(x,e)$ は
コンパクトである。
空集合はコンパクトであるから、$A(x,e)=\emptyset$ の場合も
結論は成立する。
\end{proof}

\begin{definition}[閉塞]
\label{def:blockage_sec2}
時刻 $t$ において可行行動集合が空であること
\[
  A_{i,t} = \emptyset
\]
が成立するとき、ACS個体 $\Sigma_i$ は\emph{閉塞}状態に
あるという。
閉塞においては、いかなる $u\in U$ を選択しても、
一ステップ後の次状態を $K$ に属させることができない。
すなわち
\[
  \forall u\in U,\qquad f(x_t,u,e_t)\notin K
\]
が成立する。
\end{definition}

\begin{remark}[閉塞と整形式実現の不成立・崩壊の区別]
\label{rem:blockage_not_collapse_sec2}
閉塞
$A_{i,t}=\emptyset$
は、崩壊とは区別される事象である。
本論文では「崩壊」の語を、構造的個体としての同一性を成立させていた組織関係が
同一個体として回復しない形で失われた状態が
時刻 $t$ までに成立したこと、すなわち
$\mathrm{Col}_{i,t}=1$
が成立する場合に限定して用いる
（定義~\ref{def:collapse}）。
これに対して閉塞は、現在状態および実現外部環境状態のもとで、
次状態を可行集合 $K$ に属させる行動が存在しないことを表す。
したがって本論文は、閉塞 $A_{i,t}=\emptyset$ から
崩壊 $\mathrm{Col}_{i,t}=1$ を一般的な帰結として置かない。

また本論文は、閉塞から当該時刻の整形式実現の不成立
$\mathrm{WF}_{i,t}=0$ も一般的な帰結として置かない（定義~\ref{def:wf_realization}）。
閉塞が成立する場合であっても、同一ACS個体としての
整形式な実現が保たれている限り
$\mathrm{WF}_{i,t}=1$ でありうる
（備考~\ref{rem:time_identity_collapse}）。
閉塞と崩壊の区別、および両者に至る経路の整理は
\S\ref{sec:collapse_cond} で扱う。
\end{remark}

\begin{remark}[時間的自己同一性と機能的終端・崩壊・閉塞の区別]
\label{rem:time_identity_collapse}
予備条件 (TI-1) は、時間的自己同一性の比較が
両時刻における整形式な実現を前提とすることを述べる。
ただし本論文は、$\mathrm{WF}_{i,t}=0$ であることだけから
機能的終端または崩壊を一般的な帰結として置かず、
閉塞 $A_{i,t}=\emptyset$ からも機能的終端または崩壊を一般的な帰結として置かない。
また、閉塞が成立していても $\mathrm{WF}_{i,t}=1$ である時刻を排除せず、
その場合には予備条件 (TI-1) は成立しうる。
機能的終端・崩壊と整形式実現との関係は
備考~\ref{rem:collapse_wf_relation} に従う。
\end{remark}

閉塞・整形式実現の不成立・機能的終端・崩壊の論理関係を図~\ref{fig:s2_termination} に示す。

\begin{figure}[tbp]
\centering
\begin{tikzpicture}[
  font=\footnotesize,
  box/.style={draw, rounded corners=2pt, align=center, inner sep=3pt},
  cond/.style={font=\scriptsize, inner sep=1pt},
  note/.style={font=\scriptsize, align=center, inner sep=1pt},
  plabel/.style={font=\small, anchor=west, inner sep=0pt},
  imp/.style={-{Latex[length=1.8mm]}, thick},
  bar/.style={line width=3pt},
  tick/.style={densely dotted, gray},
]
\node[plabel] at (0,0.9) {(a)\ 同じ時刻 $t$ における含意};
\node[box] (col) at (4.2,0) {崩壊\\$\mathrm{Col}_{i,t}=1$};
\node[box] (ft)  at (1.6,-1.9) {機能的終端\\$\mathrm{FT}_{i,t}=1$};
\node[box] (wf)  at (6.8,-1.9) {整形式実現の不成立\\$\mathrm{WF}_{i,t}=0$};
\node[box] (bl)  at (10.4,-0.95) {閉塞\\$A_{i,t}=\emptyset$};
\draw[imp] (col) -- node[cond, above left]{命題~\ref{prop:termination_order}(ii)} (ft);
\draw[imp] (col) -- node[cond, above right]{命題~\ref{prop:termination_order}(iv)} (wf);
\node[note, anchor=north] at (bl.south) {他の三つとの間に\\線を引かない};

\begin{scope}[yshift=-4.3cm]
\node[plabel] at (0,0.7) {(b)\ 崩壊が成立する場合の時刻別指標（$\mathrm{Col}_{i,t}=1$ となる $t$ があるとき）};
\draw[-{Latex[length=1.6mm]}] (2.4,-3.7) -- (12.3,-3.7);
\node[cond, anchor=north east] at (12.3,-3.75) {時刻};
\node[cond, anchor=north] at (2.4,-3.75) {$0$};
\coordinate (sF) at (6.0,0);
\coordinate (tC) at (9.0,0);
\draw[tick] (sF |- 0,-0.1) -- (sF |- 0,-3.7);
\draw[tick] (tC |- 0,-0.1) -- (tC |- 0,-3.7);
\node[note, anchor=north] at (sF |- 0,-3.8) {機能的終端状態が\\初めて成立する時刻};
\node[note, anchor=north] at (tC |- 0,-3.8) {$t_{\mathrm{C}}$：崩壊状態が\\初めて成立する時刻};
\node[cond] at (7.5,-4.05) {$\leq$};
\node[anchor=east] at (2.3,-0.4) {$\mathrm{FT}_{i,t}$};
\draw[bar, gray!45] (2.4,-0.4) -- (sF |- 0,-0.4);
\draw[bar] (sF |- 0,-0.4) -- (12.1,-0.4);
\node[cond, above] at (4.2,-0.35) {$0$};
\node[cond, above] at (10.6,-0.35) {$1$（以後 $1$）};
\node[anchor=east] at (2.3,-1.1) {$\mathrm{Col}_{i,t}$};
\draw[bar, gray!45] (2.4,-1.1) -- (tC |- 0,-1.1);
\draw[bar] (tC |- 0,-1.1) -- (12.1,-1.1);
\node[cond, above] at (5.7,-1.05) {$0$};
\node[cond, above] at (10.6,-1.05) {$1$（以後 $1$）};
\node[anchor=east] at (2.3,-1.8) {$\mathrm{WF}_{i,t}$};
\draw[bar] (tC |- 0,-1.8) -- (12.1,-1.8);
\node[cond, above] at (10.6,-1.75) {$0$};
\node[cond, above] at (5.7,-1.75) {（この図は値を定めない）};
\node[anchor=east, align=right, font=\scriptsize] at (2.3,-2.5) {原理判定\\時刻域};
\draw[|-] (2.4,-2.5) -- (sF |- 0,-2.5);
\draw[fill=white, line width=0.45pt] (sF |- 0,-2.5) circle (1.4pt);
\node[cond, above] at (4.2,-2.45) {$\mathrm{FT}_{i,t}=0$ の時刻};
\node[anchor=east] at (2.3,-3.2) {$\mathbb{T}_{\mathrm{FT},i}$};
\draw[|-|] (2.4,-3.2) -- (tC |- 0,-3.2);
\node[cond, above] at (5.7,-3.15) {$t_{\mathrm{C}}$ 以前の時刻};
\node[cond, above] at (10.5,-3.15) {機能的終端状態を判定しない};
\end{scope}
\end{tikzpicture}
\caption{閉塞・整形式実現の不成立・機能的終端・崩壊の論理関係。
(a) 矢印は同じ時刻 $t$ における含意を表す。
線のない組については、本図は含意関係を表示しない（備考~\ref{rem:blockage_not_collapse_sec2}・備考~\ref{rem:termination_hierarchy}・備考~\ref{rem:collapse_wf_relation}）。
(b) 図は機能的終端状態が $t_{\mathrm{C}}$ より前に初めて成立する場合を描くが、両者が同じ時刻に初めて成立する場合を含む（命題~\ref{prop:termination_order}(i)(ii)）。
$\mathrm{FT}_{i,t}$ と $\mathrm{Col}_{i,t}$ は履歴指標であり、一度 $1$ になると以後 $1$ である（命題~\ref{prop:collapse_basic}）。
$t_{\mathrm{C}}$ より後の時刻には機能的終端状態を判定しない（命題~\ref{prop:termination_order}(iii)）。
崩壊後の横棒は履歴指標および帰属の規定による値の表示であり、同一個体としての実現軌道の継続を表さない（定義~\ref{def:collapse}）。
閉塞は時刻の上の段階として置かない。}
\label{fig:s2_termination}
\end{figure}

\begin{remark}[閉塞と状態遷移の記述可能性]
\label{rem:blockage_auto}
閉塞
$A_{i,t}=\emptyset$
は、いかなる行動を選択しても一ステップ後の次状態を
可行集合 $K$ に属させることができないことを表す。

一方、閉塞の成立は、状態遷移写像 $f$ またはその自律的成分
$f_{\mathrm{auto}}$ が数学的に未定義になることを意味しない。
したがって、任意の $u\in U$ に対する
$f(x_t,u,e_t)$
は、状態遷移写像の定義域上で引き続き定義される。
ただし、その次状態は $K$ に属さない。

$\mathrm{WF}_{i,t}=1$ が保たれ、状態遷移に必要な
実現条件が保たれている場合には、閉塞後の状態変化も
$x_{t+1}=f(x_t,u_t,e_t)$
によって記述できる。

ただし本論文は、閉塞状態において $f_{\mathrm{auto}}$ に対応する
実装上の自律的機能過程が実際に作動し続けること、
および実効的時間地平
$\widehat{\tau}_{t}>0$
が成立することを、閉塞の定義または
仮定~\ref{ass:f_decomposition} のみからの一般的な帰結として置かない。
本論文はこれらの成立条件を与えない。
なお $\widehat{\tau}_{t}>0$ は予測精度係数が正であることと同値であり
（\S\ref{sec:extended_dynamics} 定義~\ref{def:tau_hat_section5}）、
その値は記述者が指定する読み出し規則に依存する。
\end{remark}

\begin{remark}[閉塞の意味の範囲]
\label{rem:blockage_scope_sec2}
閉塞 $A_{i,t}=\emptyset$ は、現在の状態と実現外部環境状態のもとで、
次の一段の状態を $K$ に置く行動が存在しないことだけを表す。
\begin{enumerate}[label=(\alph*)]
  \item $x_t\in K$ の時刻の閉塞を\emph{維持の局面における閉塞}と読み、
  次の時刻に $K$ を外れることがどの行動によっても避けられないことを表す。
  $x_t\notin K$ の時刻の閉塞を\emph{回復の局面における閉塞}と読み、
  一段では $K$ に戻れないことを表す。
  これらは読みの呼称であり、定義と記号を追加しない
  （\S\ref{sec:three_principles} の維持の局面と回復の局面の区別に対応する）。
  \item 本論文は、回復の局面における閉塞から、有限段で $K$ に戻れないことを一般的な帰結として置かない。
  数段を要する回復の途中にある個体は、回復が完了する直前の段に至るまで、
  各時刻に回復の局面における閉塞にありうる。
  \item 可行行動集合と閉塞は、採用した離散時間の一段遷移仕様、
  すなわち時間解像度と、それに対応して定める行動空間 $U$ および状態遷移写像 $f$ に相対して定まる。
  本論文は、異なる一段遷移仕様による形式化のあいだで閉塞の判定が保存されることを仮定しない。
  \item 本論文は、次のいずれの成立からも $A_{i,t}=\emptyset$ を一般的な帰結として置かない。
  物理的に移動できないこと、目的や課題が進まないこと、
  数段先に可行性の喪失が避けられないこと、
  当該ACSが内部で有効な行動候補を見いだせないこと、
  当該ACSが打つ手がないと評価すること。
  これらと閉塞が同じ時刻に成り立つことは排除しない。
\end{enumerate}
\end{remark}

\subsubsection{一段階環境可行性}

\begin{definition}[一段階可行環境対応（分類C）]
\label{def:one_step_viable_environment}
各状態 $x\in X$ に対し、一段階可行環境対応
\[
  \mathcal{E}_{\mathrm{viab}}^{(1)}:X\rightrightarrows E_{\mathrm{ext}}
\]
を
\[
  \mathcal{E}_{\mathrm{viab}}^{(1)}(x)
  :=\{\,e\in E_{\mathrm{ext}}\mid A(x,e)\neq\emptyset\,\}
\]
によって定義する。
すなわち $\mathcal{E}_{\mathrm{viab}}^{(1)}(x)$ は、状態 $x$ に
おいて少なくとも一つの一段階可行行動を残す外部環境状態の
集合である。

この対応は $f$、$K$、$U$ および $E_{\mathrm{ext}}$ から導かれる
派生構造であり、ACSの7タプルへ新たな独立成分を追加するもの
ではない。
\end{definition}

\begin{remark}[一段階可行性と必要環境充足の区別]
\label{rem:fixed_environment_and_state_dependent_viability}
Phase~1 の一つの分析内では $E_{\mathrm{ext}}$ を固定する
（仮定~\ref{ass:Eext_time_fixed}）。
一方、$\mathcal{E}_{\mathrm{viab}}^{(1)}(x)$ は
内部状態 $x$ に依存して変化しうる。
したがって、内部状態の変化によって対応可能な環境範囲が
変化する事態は、固定された $E_{\mathrm{ext}}$ の内部における
$\mathcal{E}_{\mathrm{viab}}^{(1)}(x)$ の変化として表される。

$\mathcal{E}_{\mathrm{viab}}^{(1)}(x)$ は、各要素が個別に必要で
あることを表す集合ではなく、一段階可行性を許す外部環境状態の
集合である。
必要環境充足集合 $E_{\mathrm{req}}(x)$
（定義~\ref{def:E_req}）とは定まり方が異なり、
本論文は両者の間にいずれの向きの包含関係も要求しない。
すなわち、必要環境要因がすべて充足されている外部環境状態の
もとで可行行動が存在しない場合、および必要環境要因が
充足されていない外部環境状態のもとで一段階の可行行動が
存在する場合を、いずれも排除しない。

また本論文は、一段階可行性それ自体から、複数時点にわたる可行性、
周期的な環境要請または将来の環境状態への到達可能性を
一般的な帰結として置かない。
周期的または遅延を伴う環境要請の記述には、
複数時点にわたる環境列および状態軌道の別途定式化を要する。

実現外部環境状態 $e_t$ について、
\[
  e_t\notin\mathcal{E}_{\mathrm{viab}}^{(1)}(x_t)
  \quad\Longleftrightarrow\quad
  A(x_t,e_t)=\emptyset
\]
である。
このとき閉塞が成立する（定義~\ref{def:blockage_sec2}）が、
崩壊 $\mathrm{Col}_{i,t}=1$（定義~\ref{def:collapse}）とは区別される。
\end{remark}

\begin{remark}[実現外部環境状態と内部推定の区別についての予告]
\label{rem:actual_environment_and_internal_estimate}
本小節に現れる $e_t\in E_{\mathrm{ext}}$ は、システム記述上の
実現外部環境状態である。
これは、内部モデルが形成する環境推定とは区別される対象であり、
両者の一致を仮定しない。
内部環境推定の型、および実現外部環境状態との関係は
\S\ref{subsec:internal_model_M} で定める。
本小節では $E_{\mathrm{ext}}$ 上の距離、擬距離または推定誤差
評価写像を導入しないため、両者の乖離を定量化しない。
\end{remark}

\begin{assumption}[可行行動集合に対する環境入力の非退化性]
\label{ass:env_viable_nondegeneracy}
外部環境状態の差異は可行行動集合に影響しうる。
すなわち
\[
  \exists x\in X,\
  \exists e,e'\in E_{\mathrm{ext}},\
  e\neq e':
  \quad
  A(x,e)\neq A(x,e')
\]
が成立する。

\medskip
\noindent
導入地位：S（構造的）。強さ：弱。
本仮定は、外部環境状態の差異が可行行動の集合に一切影響しない
退化的な系を Phase~1 の対象から除くためのものである。
本仮定は、当該の差異が必要環境要因の充足に由来することを
述べるものではない。
\end{assumption}

\begin{proposition}[非退化性の水準間の含意]
\label{prop:nondegeneracy_implication}
仮定~\ref{ass:env_viable_nondegeneracy} は
仮定~\ref{ass:env_input_nondegeneracy}
（状態遷移の水準における非退化性）を含意する。
\end{proposition}

\begin{proof}
仮定~\ref{ass:env_viable_nondegeneracy} により、
ある $x\in X$ と $e\neq e'$ について
$A(x,e)\neq A(x,e')$ である。
一般性を失うことなく
$u\in A(x,e)\setminus A(x,e')$ を取る。
定義~\ref{def:at} より $f(x,u,e)\in K$ かつ
$f(x,u,e')\notin K$ である。
$K$ への所属が異なるから
$f(x,u,e)\neq f(x,u,e')$ であり、
これは仮定~\ref{ass:env_input_nondegeneracy} の
条件を満たす組を与える。
\end{proof}

\begin{remark}[逆向きの含意について]
\label{rem:nondegeneracy_converse}
命題~\ref{prop:nondegeneracy_implication} の逆は一般には成立しない。
$f(x,u,e)\neq f(x,u,e')$ であっても、
両者がともに $K$ に属するか、ともに $K$ に属さない場合には
$u$ の可行性は変化せず、
$A(x,e)=A(x,e')$ となりうるためである。
したがって本論文は、状態遷移の水準における非退化性から
可行行動集合の水準における非退化性を一般的な帰結として置かない。
感知の水準における非退化性との関係は
\S\ref{subsec:three_primitive_functions} で整理する。
本論文は、別途接続条件を置かない限り、
感知の水準と、状態遷移の水準および可行行動集合の水準との間に
相互含意を仮定しない。
\end{remark}

\begin{assumption}[可行実現可能性]
\label{ass:viable_realizability}
必要環境要因が充足された外部環境状態のもとで、
可行性を保つ一ステップ遷移が少なくとも一つ存在する。
すなわち
\[
  \exists x\in K,\
  \exists e\in E_{\mathrm{req}}(x),\
  \exists u\in U:
  \quad
  f(x,u,e)\in K
\]
が成立する。

\medskip
\noindent
導入地位：S（構造的）。強さ：中。
本仮定は、いかなる状態および必要環境充足状態においても
可行性を保つ遷移が存在しない退化的な系を
Phase~1 の対象から除く。
本仮定は、状態遷移列上の各時刻において
$e_t\in E_{\mathrm{req}}(x_t)$ が成立することを要求しない。
\end{assumption}

\begin{remark}[可行実現可能性から従う事項]
\label{rem:viable_realizability_consequences}
仮定~\ref{ass:viable_realizability} から、次が従う。

第一に、$\exists x\in K$ であるから
$K\neq\emptyset$
である。
これにより、命題~\ref{prop:theta_from_pi}(ii) が仮説として
置く $K\neq\emptyset$ は、本仮定のもとで充足される。

第二に、当該の $x$ と $e$ について $u\in A(x,e)$ であるから
$A(x,e)\neq\emptyset$ であり、
\[
  e\in\mathcal{E}_{\mathrm{viab}}^{(1)}(x)
  \cap
  E_{\mathrm{req}}(x)
\]
である。
したがって、必要環境充足集合と一段階可行環境対応の
共通部分が空でない状態が少なくとも一つ存在する。

第三に、構成要件 CR-1 の (CR-1b)（共同充足可能性、
\S\ref{subsec:environment_space}）が要求する
$E_{\mathrm{req}}(x)\neq\emptyset$ なる $x$ の存在は、
本仮定から $x\in K$ の範囲で従う。
ただし本論文は、本仮定から CR-1 の全体が従うこと、
および CR-1 から本仮定が従うことを一般的な帰結として置かない。
CR-1 は当該7タプル構造体をACSのクラスへ属させる構成要件であり、
本仮定は後続命題の成立範囲を限定する仮定である。
\end{remark}

%----------------------------------------------------------------------
\subsection{原始的な機能と基盤機能の具備および標準写像形式}
\label{subsec:three_primitive_functions}

本小節では、\S\ref{subsec:primitive_functions} で定めた原始的な機能に関して、
7タプル構造の実現が有する\emph{基礎的な機能}を定義し、
その具備を要請する構成要件 CR-2 を与える。
そのうえで、原始的な機能を本論文の形式体系で分析するための
標準写像形式 $P_{\mathrm{sens}}$（感知）、$P_{\mathrm{disc}}$（識別）、
$P_{\mathrm{val}}$（評価）、$P_{\mathrm{act}}$（行動）と、
それらの非退化条件および作動上の依存関係を述べる。

構成要件 CR-2 は、定義~\ref{def:acs} が前方参照している
二つの構成要件のうちの一方である。
他方の CR-1（必要環境整合性）は
\S\ref{subsec:environment_space} で与えた。
両者はいずれも当該7タプル構造体を本論文におけるACSのクラスへ
属させるための条件であるが、対象が異なる。
CR-1 は必要環境空間 $E$ の各元が存続に必要であることを拘束し、
CR-2 は原始的な機能の各区分に対応する基礎的な機能の具備とその機能的非退化性、
およびそれらの活動を支える基盤機能の具備を要請する。
本論文は一方から他方を一般的な帰結として置かない。

本小節で新たに用いる記号は、感知情報の空間 $\mathcal{I}$、
内部表現の空間 $\mathcal{R}_{\mathrm{int}}$ および
評価値の集合 $\mathcal{V}$ である。
下付き $\mathrm{int}$ は internal（内部）に対応する。

\begin{definition}[基礎的な機能]
\label{def:basic_functions}
定義~\ref{def:acs} の7タプル構造 $\Sigma$ の実現が、
原始的な機能（定義~\ref{def:primitive_functions}）の
各区分に関して有する機能を、
その実現の\emph{基礎的な機能}と呼ぶ。
ACS個体 $\Sigma_i$ については、$\Sigma_i$ の基礎的な機能という。
基礎的な機能は機能的責務に関する記述であり、
特定の物理部位、計算モジュール、標準写像
または一対一の実装分割を意味しない。
一つの区分の責務を複数の基礎的な機能が共同して担うこと、
および一つの基礎的な機能が複数の区分の責務に関与することを排除しない。
\end{definition}

本定義の対象を「ACS個体」ではなく「7タプル構造の実現」とするのは、
次の構成要件 CR-2 が、ACSであるための条件として本定義を用いるためである。

\subsubsection*{構成要件 CR-2（原始的な機能と基盤機能の具備）}
\label{req:CR2}

ACSであるための構成要件として、7タプル $\Sigma$ の実現が
次の (CR-2a)〜(CR-2c) を満たすことを要請する。
\begin{description}[style=nextline]
  \item[(CR-2a) 基礎的な機能の具備]
  原始的な機能の四区分（感知・識別・評価・行動）のそれぞれについて、
  基礎的な機能（定義~\ref{def:basic_functions}）を具備する。
  \item[(CR-2b) 機能的非退化性]
  (CR-2a) の基礎的な機能は、各区分について次の意味で退化していない。
  \begin{enumerate}[label=(\roman*)]
    \item 感知：外部環境状態を固定したとき、異なる内部状態から
    異なる情報が得られる場合が少なくとも一つ存在し、
    かつ内部状態を固定したとき、異なる外部環境状態から
    異なる情報が得られる場合が少なくとも一つ存在する。
    \item 識別：少なくとも一つの内部状態において、
    外部環境状態の違いが、得られた情報を経て異なる内部表現に至る場合が存在する。
    \item 評価：評価の対象となる入力の違いを非等価に扱う場合が
    少なくとも一つ存在する。
    さらに、可行集合 $K$（定義~\ref{def:K}）について
    $K\neq\emptyset$ かつ $X\setminus K\neq\emptyset$ であるときは、
    $K$ に属する状態と属さない状態の違いに対して
    感知・識別を通じて評価が異なる場合が少なくとも一つ存在する（可行性感応性）。
    標準写像形式では、同じ内部モデル状態・外部環境状態のもとで比較する。
    \item 行動：利用可能な内部表現、評価または現在の状態の
    少なくともいずれかの違いに応じて、生成される応答が異なる場合が
    少なくとも一つ存在する。
    これらのすべての違いに応答が依存することは要求しない。
  \end{enumerate}
  \item[(CR-2c) 基盤機能の具備]
  (CR-2a) の基礎的な機能の活動を支える基盤機能
  （定義~\ref{def:base_functions}）を具備する。
\end{description}

(CR-2a) および (CR-2c) の具備は、当該の機能が実現の構成として備わっていることを
指し、各時刻においてその機能が作動していることを要求しない。
(CR-2b) も、各時刻の作動の有無ではなく、具備された機能が
退化していないことに関する条件である。
時刻ごとの作動の停止は、機能的終端の判定
（定義~\ref{def:functional_termination}）で扱う。

(CR-2b) は分類Aの水準で述べた機能的条件であり、
特定の写像の性質として与えたものではない。
定義~\ref{def:three_primitives} の非退化条件は、
本論文が標準写像形式を用いて記述する場合に、
(CR-2b) の各項を写像の性質として表したものである。
ある実現の基礎的な機能と標準写像形式との対応づけは
対象系の記述における選択であり、本論文はその一般的な判定基準を
Phase~1 では与えない。

ここで CR-2 は、7タプルに追加される独立成分を導入するもの
ではなく、7タプル $\Sigma$ がACSとして適格であるための
条件である。
定義~\ref{def:three_primitives} は、原始的な機能の機能的責務を
本論文の形式体系において分析するための標準写像形式
$P_{\mathrm{sens}},P_{\mathrm{disc}},P_{\mathrm{val}},P_{\mathrm{act}}$ を与える。
CR-2 は、当該実現がこの標準写像形式そのものを
部品として備えることを要求せず、
同じ機能的責務を満たす機能的に等価な実現を許す。
$P_{\mathrm{sens}},P_{\mathrm{disc}},P_{\mathrm{val}},P_{\mathrm{act}}$ は7タプル
$\Sigma=(X,U,E,f,\Pi,M,\mathrm{Update})$ の成分でも、
そこから一意に導出される対象でもない。

原始的な機能のいずれかの区分について基礎的な機能を構成として欠く系は、
本論文におけるACSの定義を充足しない。

構成要件 CR-2 は、時刻ごとの診断ではなく、
$\Sigma_i$ をACS個体として同定する条件である
（備考~\ref{rem:post_collapse_record}）。
したがって、ACS個体として同定された $\Sigma_i$ が、
その後の損傷その他により、原始的な機能の一つ以上の区分について
基礎的な機能を回復しない形で全面的に失った場合にも、
そのことのみによって $\Sigma_i$ のACS個体としての同定は
取り消されない。
本論文は、そのような時刻に改めて CR-2 を判定して
$\Sigma_i$ をACSでないとすることはせず、
それ以前の同定を遡って取り消すこともしない。
このとき $\Sigma_i$ の機能の状態は、
機能的終端（定義~\ref{def:functional_termination}）の成否、
および \S\ref{sec:collapse_cond} の記述的状態帯によって記述する。
一つ以上の区分の全面的な喪失は、それだけでは機能的終端でも崩壊でもない
（備考~\ref{rem:termination_hierarchy}）。

構成要件 CR-2 は、後続命題の成立範囲を制限する仮定ではなく、
ACSという対象を定義するための構成要件である。
したがって、本論文では CR-2 を通常の仮定一覧には含めず、
構成要件 CR-1 および規約 TSC と同様に、
ACS定義を支える基礎条件として扱う。

\begin{remark}[原始的な機能の構成的必要性（分類A）]
\label{rem:three_primitives_necessity}
\S\ref{sec:intro} で述べた環境において長期的な可行性の維持を
指向するACSには、原始的な機能（定義~\ref{def:primitive_functions}）の
各区分について、次の機能的責務が構成的に要求される。
第一に、自己の内部状態および外部環境状態から情報を得る責務（感知）である。
第二に、得られた情報を区別し、内部表現を形成する責務（識別）である。
第三に、内部表現、自己の現在の状態、応答の候補その他の評価の対象を、
自己の状態に関わる内部基準に照らして非等価に扱う責務（評価）である。
第四に、利用可能な内部表現、評価または現在の状態の少なくともいずれかに応じて
応答を生成する責務（行動）である。
評価と行動のいずれについても、列挙した対象または情報のすべてを
用いることは要求しない
（構成要件 CR-2 の (CR-2b)、定義~\ref{def:three_primitives}）。

評価の責務を区分として置くのは、
感知・識別・行動の三つの責務だけでは、
応答の生成を自己の状態に関わる基準へ結びつける責務が含まれず、
長期的な可行性の維持を「指向する」という対象規定を
担う責務が欠けるためである。
評価の区分は機能の水準の概念であり、
\S\ref{sec:three_principles} の三原理に関連する
構造的機能のうち評価に関わるものと整合するが、
三原理の時刻別充足述語の真偽を表すものではない
（備考~\ref{rem:primitive_vs_principles}）。
ただし、評価の内部基準は持続評価関数 $\Pi$ ではなく、
本論文はACSが $\Pi$ を内部で保持または計算することを要求しない。

時刻をまたいで内部表現を保持する機能は、
原始的な機能の区分とは区別し、内部モデル構造
（\S\ref{subsec:internal_model_M}）および
\S\ref{sec:tipa} で扱う。
保持する機能は、基礎的な機能の活動を支える限りで基盤機能
（定義~\ref{def:base_functions}）に含まれうる。
この場合、機能的終端の条件 (FT-1)(b) の対象となる
（定義~\ref{def:functional_termination}）。
内部モデル更新は区分として加えるものではなく、
対象系においていずれかの区分の責務を担う場合には、
その区分に関する基礎的な機能として扱いうる。

これらの機能的要請は、それぞれが独立した物理部位、
計算モジュールまたは明示的な写像として実装されることを要求しない。
同じ機能的要請を満たす機能的に等価な実現を許す。
また、本備考は、これらの責務の必要性が
あらゆる数学的環境について証明済みであることを主張しない。
その必要性の射程は \S\ref{sec:intro} の規定に従う。

以下の定義~\ref{def:three_primitives} は、
この機能的要請を本論文の形式体系において
分析・比較・記述するために採用する標準写像形式である。
構成要件 CR-2 が要求するのは、
原始的な機能の各区分に対応する基礎的な機能の具備、
その機能的非退化性 (CR-2b) および基盤機能の具備であり、
特定の写像分解の実装を要求するものではない。
同定義の非退化条件は、(CR-2b) を標準写像形式の上で表したものである。
\end{remark}

\begin{definition}[原始的な機能の標準写像形式（分類B）]
\label{def:three_primitives}
非空集合 $\mathcal I,\mathcal R_{\mathrm{int}},\mathcal V$ をそれぞれ
感知情報、内部表現、評価値の型とする。以下の全域一価写像を固定する。
\[
\begin{aligned}
P_{\mathrm{sens}}&:X\times E_{\mathrm{ext}}\to\mathcal I,\\
P_{\mathrm{disc}}&:\mathcal M\times\mathcal I\to\mathcal R_{\mathrm{int}},\\
P_{\mathrm{val}}&:\mathcal M\times\mathcal R_{\mathrm{int}}\to\mathcal V,\\
P_{\mathrm{act}}&:\mathcal M\times\mathcal R_{\mathrm{int}}\times\mathcal V\to U.
\end{aligned}
\]
これらは原始的な機能を記述する標準写像であり、機能区分そのものでも
7タプルの追加成分でもない。7タプルだけから一意に導出されるとはしない。
$\mathcal I,\mathcal R_{\mathrm{int}},\mathcal V$ には、本定義では位相・順序を課さず、
写像の連続性も要求しない。評価値は候補別の評価情報を含む束であってよい。
真の状態・環境への直接アクセスを内部機構に要求せず、同時刻の入力は
\[
 I_t=P_{\mathrm{sens}}(x_t,e_t),\quad
 r_t=P_{\mathrm{disc}}(M_t,I_t),\quad
 v_t=P_{\mathrm{val}}(M_t,r_t),\quad
 u_t=P_{\mathrm{act}}(M_t,r_t,v_t)
\]
として接続する。内部表現は利用可能情報を保持する成分を含めてよく、
すべての引数への実質的依存を要求しない。

標準写像を用いる記述では、CR-2b の非退化性を次の条件で表す。
\begin{enumerate}
\item 感知の内部状態差への感応：
$\exists e\in E_{\mathrm{ext}},\ x,x'\in X:
P_{\mathrm{sens}}(x,e)\ne P_{\mathrm{sens}}(x',e)$。
\item 感知の外部環境差への感応：
$\exists x\in X,\ e,e'\in E_{\mathrm{ext}}:
P_{\mathrm{sens}}(x,e)\ne P_{\mathrm{sens}}(x,e')$。
\item 識別：ある $M\in\mathcal M,x\in X,e,e'\in E_{\mathrm{ext}}$ について
\[
P_{\mathrm{disc}}(M,P_{\mathrm{sens}}(x,e))
\ne P_{\mathrm{disc}}(M,P_{\mathrm{sens}}(x,e')).
\]
\item 評価：$P_{\mathrm{val}}$ はその定義域上で非定数である。
さらに $\emptyset\ne K\ne X$ のとき、ある
$M\in\mathcal M,e\in E_{\mathrm{ext}},x\in K,x'\in X\setminus K$ について
\[
\begin{split}
P_{\mathrm{val}}(M,P_{\mathrm{disc}}(M,P_{\mathrm{sens}}(x,e)))\\
\ne P_{\mathrm{val}}(M,P_{\mathrm{disc}}(M,P_{\mathrm{sens}}(x',e))).
\end{split}
\]
\item 行動：$P_{\mathrm{act}}$ はその定義域上で非定数である。
\end{enumerate}
感知の条件は $|X|\ge2$、識別の条件は $|\mathcal R_{\mathrm{int}}|\ge2$、
評価の条件は $|\mathcal V|\ge2$、行動の条件は $|U|\ge2$ を含意する。
$\mathcal M$ が単点でも、情報の差による非退化な作動は許される。
可行性感応性は同じ $M,e$ のもとで比較する存在条件であり、
すべての状態差の識別、評価の正確性、$\Pi$ との単調対応、可行行動の選択を要求しない。
$K=\emptyset$ または $K=X$ では可行性感応性の追加条件を課さない。
$\emptyset\ne K\ne X$ の場合には、この追加条件から評価の非定数性が従う。
本論文は、評価・行動の非定数性だけから、実行時の入力像上での非定数性を一般的な帰結として置かない。

$U$ の選択的待機 $u_\emptyset$ および指向的内部・外部行動の区別は
備考~\ref{rem:action_space_structure} に従う。
自律的機能過程 $f_{\mathrm{auto}}$ は選択対象ではない。
\end{definition}

\begin{proposition}[評価の可行性感応性が要求を課す範囲]
\label{prop:feasibility_sensitivity_scope}
仮定~\ref{ass:env_viable_nondegeneracy} のもとで、
\[
  K\neq\emptyset,
  \qquad
  X\setminus K\neq\emptyset
\]
が成り立つ。
したがって、この仮定のもとでは、
定義~\ref{def:three_primitives} の評価の可行性感応性は
「何も要求しない」場合に該当せず、要求を課す条件となる。
\end{proposition}

\begin{proof}
仮定~\ref{ass:env_viable_nondegeneracy} により、ある $x\in X$ と
$e\neq e'$ について $A(x,e)\neq A(x,e')$ である。
したがって、一方に属し他方に属さない $u\in U$ が存在する。
一般性を失わず $u\in A(x,e)\setminus A(x,e')$ とすると、
可行行動対応の定義（定義~\ref{def:at}）により
$f(x,u,e)\in K$ かつ $f(x,u,e')\notin K$ である
（命題~\ref{prop:nondegeneracy_implication} の証明と同じ推論）。
前者から $K\neq\emptyset$ である。
$f:X\times U\times E_{\mathrm{ext}}\to X$ であるから
$f(x,u,e')\in X\setminus K$ であり、$X\setminus K\neq\emptyset$ である。
\end{proof}

構成要件 CR-2 は、本節の仮定に依存せずACSのクラスを定める条件である。
命題~\ref{prop:feasibility_sensitivity_scope} は CR-2 の内容を変更せず、
上記の仮定を用いる範囲では、評価について
可行性感応性まで要求されることを述べるにとどまる。

\begin{remark}[感知の非退化性と他の水準との関係]
\label{rem:sensing_nondegeneracy_relation}
定義~\ref{def:three_primitives} の感知の外部環境状態入力に関する
非退化性は、外部環境入力に関する非退化性の第三の水準である
（備考~\ref{rem:nondegeneracy_levels}）。
第一の水準は状態遷移に関するもの
（仮定~\ref{ass:env_input_nondegeneracy}）、
第二の水準は可行行動集合に関するもの
（仮定~\ref{ass:env_viable_nondegeneracy}）である。

本論文は、別途接続条件を置かない限り、
感知の水準と他の二つの水準との間に相互含意を仮定しない。
感知の標準写像形式 $P_{\mathrm{sens}}$ と状態遷移写像 $f$ の間に接続条件を置いていない
ためであり、これは両者が独立であることを証明したという
主張ではない。

また、感知の非退化性は必要環境空間 $E$ の構成要件 CR-1
（\S\ref{subsec:environment_space}）とも別の条件である。
ある外部環境状態差が感知されることと、
その差が必要環境要因の充足に関わることとは、
本論文では区別される。
\end{remark}

\begin{remark}[原始的な機能の標準写像形式の作動依存]
\label{rem:P_dependency}
定義~\ref{def:three_primitives} の合成は同時刻の情報生成上の接続規約であり、
CR-2 の機能具備からその特定の写像分解を導くものではない。
可行性感応性はこの接続を通じた差異の伝達を要求する。
$P_{\mathrm{act}}$ は $r_t,v_t$ のすべての成分を利用する必要はない。
局所保護応答を含む標準形式では、局所情報を $r_t$ の成分として利用可能にし、
詳細評価の結果を用いない応答も表現できる。
この合成規約は物理的な逐次実行や遅延の下限を要求しない。
\end{remark}

\begin{remark}[作動依存と形式的存在の区別]
\label{rem:functional_vs_existential}
備考~\ref{rem:P_dependency} は機能的作動依存を述べる
ものであり、形式的存在の依存ではない。
写像としての $P_{\mathrm{disc}} : \mathcal M\times\mathcal{I}\to\mathcal{R}_{\mathrm{int}}$
は、$P_{\mathrm{sens}}$ の存在とは独立に型として定義可能である。
$P_{\mathrm{val}}$ と $P_{\mathrm{act}}$ についても同様である。
本備考が主張するのは、ACS 内部で各機能が意味のある入力を
得て作動するための依存順序であって、写像の数学的存在性の
順序ではない。
\end{remark}

\begin{remark}[情報依存と物理的時間順序の区別]
\label{rem:primitive_no_temporal_order}
同時刻の合成順序は、入力がどの出力に由来するかを指定する。
各写像の実装を物理的な逐次モジュールへ固定しない。
$M_t$ は先行時刻の情報を保持しうるが、同時刻の $r_t,v_t$ が
あらかじめ $M_t$ に含まれることを要求しない。
実効行動の効果は $x_{t+1}=f(x_t,u_t,e_t)$ を介して後続の感知に反映されうる。
$f$ の行動引数への非退化な依存は一般には要求しない
（仮定~\ref{ass:f_decomposition}）。
\end{remark}

%----------------------------------------------------------------------
\subsection{内部モデル構造 $M$}
\label{subsec:internal_model_M}

本小節では、ACSが環境・自己状態・行動可能性を処理するための
内部モデル構造 $M$ を定義する。
本論文では $M$ を情報統合・適応的評価・応答行動の三機構に
分解して扱い、この三成分構造を
\emph{三機構情報処理アーキテクチャ}
（TIPA: Triadic Information-Processing Architecture）
と呼ぶ。
TIPA は $M_1,M_2,M_3$ の機構構成および相互接続を表す総称であり、
7タプルの新たな独立成分ではない。
また、TIPA の三機構は、原始的な機能という抽象的な機能区分とも、
定義~\ref{def:three_primitives} の標準写像形式とも区別される。
三機構と原始的な機能の区分の間に一対一の対応を置かない。

各機構の動作原理と相互接続の詳細は
\S\ref{sec:tipa} で述べ、本小節では7タプルの一成分
としての $M$ の型と基本性質を与える。
あわせて、\S\ref{subsec:viability_function_and_blockage} から
委譲された内部環境推定の型、および
\S\ref{subsec:three_primitive_functions} から委譲された
内部表現の符号化の扱いを本小節で定める。

本小節で新たに用いる記号は、内部モデル状態の空間
$\mathcal{M}$（花文字エム）、その直積成分
$\mathcal{M}_1,\mathcal{M}_2,\mathcal{M}_3$、
直積距離 $d_{\mathcal{M}}$、
TIPA成分別変化量 $\delta_{M,c,t}$、
TIPA成分別変化ベクトル
$\boldsymbol{\delta}_{M,t}$、
最大集約変化量 $\delta_t$、
および内部環境推定 $\hat{e}_t$ である。

\subsubsection{内部モデルの型と状態空間}

\begin{assumption}[$\mathcal{M}$ のコンパクト性と成分距離の尺度規約]
\label{ass:M_compact}

各 $c=1,2,3$ について、非空コンパクト距離空間
$(\mathcal M_c^{\mathrm L},d_c^{\mathrm L})$ と $(\mathcal M_c^{\mathrm W},d_c^{\mathrm W})$ を与え、
\[
 \mathcal M_c=\mathcal M_c^{\mathrm L}\times\mathcal M_c^{\mathrm W},\qquad
 d_c((l,w),(l',w'))=\max\{d_c^{\mathrm L}(l,l'),d_c^{\mathrm W}(w,w')\}
\]
とする。上付き $L,W$ は学習成分・作動成分の識別ラベルであり、冪ではない。
各因子には距離位相、直積には積位相を入れる。
各因子が非単点なら距離を直径 $1$ に正規化し、単点なら距離は $0$ とする。
有限直積はコンパクトであり、上の max 距離は積位相を与える。
いずれかの因子が非単点なら $d_c$ の直径は $1$、両方単点なら $0$ である。
この分解は状態担体の指定であり、学習成分と作動成分の更新の独立性、
時間尺度の分離、既存の任意の計量との等長性を主張しない。

内部モデル状態空間の各成分
$\mathcal{M}_c$（$c=1,2,3$）は、非空なコンパクト距離空間
\[
  (\mathcal{M}_c,d_c),
  \qquad
  \mathcal{M}_c\neq\emptyset
\]
として与えられる。
非空性は、各時刻において内部モデル状態
$M_{c,t}\in\mathcal{M}_c$ が存在するために要する。

Phase~1 では、三成分間で距離尺度の任意な選択が
最大集約距離に影響することを避けるため、
各成分距離 $d_c$ は、非退化な
$\mathcal{M}_c$ に対して
\[
  \operatorname{diam}_{d_c}(\mathcal{M}_c)=1
  \qquad
  \text{if }|\mathcal{M}_c|\geq2
\]
を満たすよう正規化された距離として採用する。

$\mathcal{M}_c$ が一点集合である場合には
\[
  d_c\equiv0
\]
とする。

この尺度規約は、三機構の機能的重要性または
実装上のコストが等しいことを仮定するものではなく、
各成分距離を共通の相対尺度上で比較可能にするための
計量上の規約である。

\medskip
\noindent
導入地位：S（構造的）。強さ：強。
各成分のコンパクト性から、有限直積
\[
  \mathcal{M}
  =
  \mathcal{M}_1
  \times
  \mathcal{M}_2
  \times
  \mathcal{M}_3
\]
のコンパクト性が従う。
また、各 $\mathcal{M}_c$ の直径が有限であることにより、
本仮定の尺度規約による成分距離の正規化が可能となる。
この直積構造は、機構間の情報依存
\[
  M_1\to M_2\to M_3
\]
を否定しない。
Phase~2 では局所コンパクト化等の緩和を検討する。
\end{assumption}

\begin{definition}[内部モデル状態の学習成分と作動成分]
\label{def:model_learning_working}
$\mathcal M^{\mathrm L}:=\prod_{c=1}^3\mathcal M_c^{\mathrm L}$、
$\mathcal M^{\mathrm W}:=\prod_{c=1}^3\mathcal M_c^{\mathrm W}$ とし、
$M_t^{\mathrm L}=(M_{c,t}^{\mathrm L})_{c=1}^3$、$M_t^{\mathrm W}=(M_{c,t}^{\mathrm W})_{c=1}^3$ を
$M_t$ の座標射影とする。
学習成分は持続的な内容・学習・忘却を、作動成分は後続時刻へ持ち越す
短期履歴、前時刻の評価・行動の記録等を担う。
どの記録をどの機構成分に保持するかは、更新写像の依存制約に従って指定する。
$\mathcal M^{\mathrm W}$ からの読み出しが実際の過去の値と一致すること、
任意の履歴を符号化できることは、この直積構造からは従わない。
$\mathcal M^{\mathrm L}$ と $\mathcal M^{\mathrm W}$ は $\mathcal M$ の座標の整理であり、
7タプルに追加する状態ではない。
\end{definition}


\begin{definition}[内部モデル構造]
\label{def:internal_model}

内部モデル構造を
\[
  M=(M_1,M_2,M_3)
\]
とする。
各成分の機能的役割の概要は次のとおりである。

$M_1$ は情報統合機構であり、感知情報の統合・境界管理・
時間的連続性の維持を担う。
$M_2$ は適応的評価機構であり、三原理に関連する
機構層代理評価と Priority 警告等の生成を担う。
$M_3$ は応答行動機構であり、行動決定に必要な情報整理、
候補構成・評価結果の利用および TIPA 行動指令の決定を担う。

$M_1$・$M_2$・$M_3$ は機能的に並行して作動しうる
（仮定~\ref{ass:functional_concurrency}、
\S\ref{sec:tipa}）。
この機能的並行性は、三機構が逐次的にのみ作動するのではなく、
必要な上流出力が下流機構の入力として利用可能となる
$M_1\to M_2\to M_3$ の生成依存関係と両立することを表す。
定義~\ref{def:update_operator} の更新式に現れる
\[
  M_1\to M_2\to M_3
\]
の順序は情報依存関係を表すものであり、
三機構が実時間上で逐次的にしか作動しないことを意味しない。

本小節は、三機構が同一時刻に相互作用することを述べない。
機構間の同時刻における情報授受の可否、およびその時刻規約は
\S\ref{sec:tipa} で扱う。

各成分の詳細な定義・動作原理・相互接続構造は
\S\ref{sec:tipa}
（TIPA：三機構情報処理アーキテクチャ）で与える。
\end{definition}

\begin{definition}[内部モデル状態空間]
\label{def:model_parameter_space}

仮定~\ref{ass:M_compact} のもとで、
内部モデル構造 $M$ に対応する内部モデル状態空間
$\mathcal{M}$ を
\[
  \mathcal{M}
  =
  \mathcal{M}_1
  \times
  \mathcal{M}_2
  \times
  \mathcal{M}_3
\]
と定義する。

各 $\mathcal{M}_c$ は、仮定~\ref{ass:M_compact} で
与えられた非空なコンパクト距離空間
\[
  (\mathcal{M}_c,d_c),
  \qquad
  c=1,2,3
\]
である。

さらに、各成分距離 $d_c$ について、
仮定~\ref{ass:M_compact} の尺度規約に従い、
$|\mathcal{M}_c|\geq2$ の場合には
\[
  \operatorname{diam}_{d_c}(\mathcal{M}_c)=1
\]
とし、$\mathcal{M}_c$ が単点集合である場合には
\[
  d_c\equiv0
\]
とする。
この正規化は各成分距離の乗法的尺度を共通の相対範囲へ
揃えるための規約であり、
三機構の機能的重要性、実装コスト、または物理的変化量が
等しいことを意味しない。

任意の
\[
  m=(m_1,m_2,m_3),
  \qquad
  m'=(m'_1,m'_2,m'_3)
  \in\mathcal{M}
\]
に対し、直積距離 $d_{\mathcal{M}}$ を
\[
  d_{\mathcal{M}}(m,m')
  :=
  \max\left\{
    d_1(m_1,m'_1),
    d_2(m_2,m'_2),
    d_3(m_3,m'_3)
  \right\}
\]
と定義する。

時刻 $t$ の内部モデル状態を
\[
  M_t
  =
  (M_{1,t},M_{2,t},M_{3,t})
  \in\mathcal{M}
\]
と記し、
\[
  M_{c,t}\in\mathcal{M}_c,
  \qquad
  c=1,2,3
\]
とする。

仮定~\ref{ass:M_compact} と有限直積の性質により
$\mathcal{M}$ はコンパクトである。

なお、直積距離 $d_{\mathcal{M}}$ は、
$\mathcal{M}$ の距離空間構造およびコンパクト性を
記述するための計量であり、
TIPA 三機構の変化内容を余すところなく表す量としては
解釈しない。
各機構の動的変化は、
定義~\ref{def:tipa_component_change} の
TIPA成分別変化ベクトル
$\boldsymbol{\delta}_{M,t}$ により別途記述する。
\end{definition}

\begin{remark}[$M$ の独立成分としての保持]
\label{rem:M_independence}

本論文では、状態遷移写像 $f$ から
内部モデル構造 $M$ が一意に復元可能であることを仮定しない。

内部モデル構造
$M=(M_1,M_2,M_3)$
は、内部モデル状態空間
\[
  \mathcal{M}
  =
  \mathcal{M}_1
  \times
  \mathcal{M}_2
  \times
  \mathcal{M}_3
\]
の三成分分解に機構的解釈を与える。

行動の標準写像形式
\[
  P_{\mathrm{act}}:\mathcal M\times\mathcal R_{\mathrm{int}}\times\mathcal V\to U
\]
および更新作用素
\[
  \mathrm{Update}:
  \mathcal{M}\times X\times E_{\mathrm{ext}}
  \to\mathcal{M}
\]
が直接作用する対象は、
機構構造 $M$ そのものではなく、
その状態空間 $\mathcal{M}$ の元 $M_t$ である。

本論文では、状態遷移 $f$ と内部モデル構造 $M$ を
異なる構造的役割として明示的に分離する表現方針のもとで、
$M$ を7タプルの独立成分として保持する。

本記述は、$f$ から $M$ が導出不能であることの
一般定理を主張するものではない。

同一の $f$ に対応しうる内部モデル構造 $M$ の
一意性または非一意性に関する形式的検討は、
TIPA の具体的写像構造が確定する
\S\ref{sec:tipa} 以降、
および Phase~2 の課題とする。
\end{remark}

\subsubsection{内部環境推定}

\begin{definition}[内部環境推定（分類B）]
\label{def:internal_env_estimate}

時刻 $t$ において内部モデルが形成する外部環境状態の推定を
\[
  \hat{e}_t\in E_{\mathrm{ext}}
\]
と書き、\emph{内部環境推定}と呼ぶ。
上付きハットは推定量であることを表す識別子である。

$\hat{e}_t$ は実現外部環境状態 $e_t$ と同一の空間
$E_{\mathrm{ext}}$ に値をとるが、本論文は両者の一致を仮定しない。
$\hat{e}_t$ は7タプルの成分ではなく、
内部モデル状態 $M_t$ を含む利用可能な機構入力に基づき、
$M_1$ の処理によって生成される派生量である。
本論文は $\hat{e}_t$ が $M_t$ のみから定まることを要求しない。
$\hat{e}_t$ を生成する具体的な写像は $M_1$ に置き、
その形式は \S\ref{sec:tipa} で与える。
\end{definition}

\begin{remark}[実現外部環境状態と内部環境推定の乖離]
\label{rem:estimate_gap}

定義~\ref{def:internal_env_estimate} は、
$e_t$ と $\hat{e}_t$ の関係について
一致、近接またはいかなる誤差限界も要求しない。
この区別は、実現外部環境状態に基づくシステム水準の記述と、
内部モデルが生成する環境推定に基づく機構水準の記述とを
同一視しないためのものである。

Phase~1 では $E_{\mathrm{ext}}$ 上の距離、擬距離または推定誤差
評価写像を導入しないため、両者の乖離を定量化しない。
推定誤差の定量化、およびそれが可行性維持に及ぼす影響の評価は
\S\ref{sec:tipa} および Phase~2 の課題とする（\S\ref{subsec:limitations_list} の L6）。

本論文の状態遷移
$x_{t+1}=f(x_t,u_t,e_t)$
（定義~\ref{def:state_transition}）に現れるのは
実現外部環境状態 $e_t$ であり、$\hat{e}_t$ ではない。
\end{remark}

\begin{remark}[内部モデル状態と内部表現の保持]
\label{rem:M_encodes_representation}
同時刻の内部表現と評価値は定義~\ref{def:three_primitives} の中間出力である。
その一部を後続時刻で再利用する場合、保持内容は
定義~\ref{def:model_learning_working} の学習成分または作動成分から読み出す。
読み出し規則は対象の記述ごとに固定し、過去の真値の復元、任意の履歴の保持、
単射・全射・連続な符号化を要求しない。
各成分への書き込みは定義~\ref{def:update_operator} の型に従う。
包括的型に十分な引数があることだけでは、制限された各成分更新による
任意の記録の実現可能性は保証されない。
\end{remark}

\subsubsection{機構成分とTIPA統合}

\begin{definition}[TIPA成分別変化量と最大集約距離（分類C）]
\label{def:tipa_component_change}

時刻 $t$ および $t+1$ の内部モデル状態を
\[
  M_t
  =
  (M_{1,t},M_{2,t},M_{3,t}),
  \qquad
  M_{t+1}
  =
  (M_{1,t+1},M_{2,t+1},M_{3,t+1})
  \in\mathcal{M}
\]
とする。

TIPA の各機構における一ステップ変化量を、
各成分距離 $d_c$ を用いて
\[
  \delta_{M,c,t}
  :=
  d_c(M_{c,t+1},M_{c,t}),
  \qquad
  c=1,2,3
\]
と定義する。

これら三成分をまとめた
\emph{TIPA成分別変化ベクトル}を
\[
  \boldsymbol{\delta}_{M,t}
  :=
  \bigl(
    \delta_{M,1,t},
    \delta_{M,2,t},
    \delta_{M,3,t}
  \bigr)
  \in\mathbb{R}_{\geq0}^{3}
\]
と定義する。

$\boldsymbol{\delta}_{M,t}$ は、
$M_1$・$M_2$・$M_3$ の変化を相互に集約せずに記述する
一次的な動的記述量である。
これは本論文が分析のために導入する記述量であり、
当該ACSがこの量を内部に保持することを要求しない。

各成分距離 $d_c$ は、
仮定~\ref{ass:M_compact} の尺度規約により
共通の相対尺度へ正規化されているため、
\[
  0\leq\delta_{M,c,t}\leq1,
  \qquad
  c=1,2,3
\]
である。

したがって
$\boldsymbol{\delta}_{M,t}$ は、
各TIPA機構がそれぞれの状態空間全体に対して
どの程度相対的に変化したかを成分別に記述する。
この正規化は、
三機構の機能的重要性または実装上のコストが
等しいことを意味しない。

一方、内部モデル状態空間 $\mathcal{M}$ 上の
直積距離による最大集約値を
\[
\begin{aligned}
  \delta_t
  &:=
  d_{\mathcal{M}}(M_{t+1},M_t)\\
  &=
  \max\left\{
    \delta_{M,1,t},
    \delta_{M,2,t},
    \delta_{M,3,t}
  \right\}\\
  &=
  \|\boldsymbol{\delta}_{M,t}\|_{\infty}
\end{aligned}
\]
とする。

したがって $\delta_t$ は、
内部モデル変化の全成分情報そのものではなく、
三機構のうち最大の一ステップ変化を表す
二次的な集約量である。

特に、
\[
  \delta_t=0
  \iff
  \boldsymbol{\delta}_{M,t}=(0,0,0)
\]
であるが、異なる成分別変化ベクトルが
同一の $\delta_t$ を与えることがある。

したがって Phase~1 では、
TIPA の動的変化を分析する際には
$\boldsymbol{\delta}_{M,t}$ を一次的記述として用い、
$d_{\mathcal{M}}$ および $\delta_t$ は、
直積空間の距離、コンパクト性、および
最大成分変化の集約を必要とする場合に用いる
（within present framework）。
\end{definition}

\begin{observation}[$M_1/M_2/M_3$ の機能的分離方針]
\label{obs:tipa_integration}

本論文では、TIPA の三機構
$M_1,M_2,M_3$ を、それぞれ異なる機能的役割を担う
構成要素として分離して記述する。

第一に、$M_1$ と $M_2$ の分離は、
情報統合・状態推定を担う機構と、
その情報に基づいて可行性および警告を評価する機構とを
区別して記述するために保持される。
この区別は、\S\ref{sec:tipa} で定義される
Priority~3 警告を含む評価構造の記述上の基盤となる。

第二に、$M_2$ と $M_3$ の分離は、
評価を担う機構と、行動空間 $U$ への応答生成を担う機構とを
異なる機能的役割として記述するために保持される。
両者の具体的写像の定義域・値域および相互接続は
\S\ref{sec:tipa} で与える。

本観察は、二機構の統合が数学的に不可能であること、
または任意の統合が必ず機能低下を生じさせることを
主張するものではない。
三機構を分離して保持することは、
Phase~1 における TIPA の記述上の構成方針である。

機構統合を許す再定式化およびその機能的等価性の検討は、
Phase~1.$\alpha$ 以降の課題とする
（within present framework）。
\end{observation}

\begin{remark}[$M$ と $\mathrm{Update}$ の相互関係]
\label{rem:M_update_relation}

内部モデル構造 $M$ と、その状態を更新する規則
$\mathrm{Update}$（\S\ref{subsec:update_operator}）は、
7タプルにおいて異なる構造的役割をもつ独立成分として
保持される。

時刻 $t$ の内部モデル状態は
$M_t\in\mathcal{M}$
であり、行動の標準写像形式
\[
  P_{\mathrm{act}}:\mathcal M\times\mathcal R_{\mathrm{int}}\times\mathcal V\to U
\]
は $M_t$、同時刻の内部表現 $r_t$ と評価値 $v_t$ を入力として行動を生成する。
したがって、内部モデル構造 $M$ は、
その状態空間 $\mathcal{M}$ と内部モデル状態 $M_t$ を介して、
ACSの行動生成を支える機構的基盤として位置づけられる。

$M$ と $\mathrm{Update}$ の基本的な形式関係は
\S\ref{subsec:update_operator} で定める。
これらを TIPA の標準機構形式へ対応づける記述は
\S\ref{sec:tipa} で扱い、
時刻間の変化に関する事項は
\S\ref{sec:extended_dynamics} で扱う。
\end{remark}

%----------------------------------------------------------------------
\subsection{内部モデル更新作用素 $\mathrm{Update}$}
\label{subsec:update_operator}

本小節では、内部モデル状態を更新する作用素
$\mathrm{Update}$ を構成要素別に定義し、
更新に伴って実際に生じた内部モデル変化の時間平均を導入する。
これは実現した変化の記述量であり、認知的処理能力を表す量ではない
（\S\ref{subsec:cognitive_capacity_policy}）。

一ステップの内部モデル変化量そのもの、すなわち
TIPA成分別変化量 $\delta_{M,c,t}$、
TIPA成分別変化ベクトル $\boldsymbol{\delta}_{M,t}$、
および最大集約変化量 $\delta_t$ は
\S\ref{subsec:internal_model_M}
（定義~\ref{def:tipa_component_change}）で定義済みである。
本小節ではこれらを再定義せず、
$\mathrm{Update}$ によって生じた一ステップ変化量として参照する。

本小節で新たに用いる記号は、
最大成分時間平均更新速度 $\lambda_{T,t}$、
成分別時間平均更新速度 $\lambda_{M,c,T,t}$ および
$\boldsymbol{\lambda}_{M,T,t}$ である。

$T\in\mathbb{N}$, $T\ge1$ は評価に用いる有限時間窓の長さを表す。
時間窓の開始時刻が $0$ である場合には、
$\lambda_{T,0}$、
$\lambda_{M,c,T,0}$、
$\boldsymbol{\lambda}_{M,T,0}$ をそれぞれ
$\lambda_T$、
$\lambda_{M,c,T}$、
$\boldsymbol{\lambda}_{M,T}$ と略記する。

\subsubsection{更新作用素の定義}

\begin{definition}[Update作用素の構造]
\label{def:update_operator}

\textbf{(1) 包括的写像（型宣言の詳細は
備考~\ref{rem:update_type_two_layers} 参照）：}
\[
  \mathrm{Update} :
  \mathcal{M}\times X\times E_{\mathrm{ext}}
  \longrightarrow
  \mathcal{M}.
\]
本小節の包括的型では、$\mathrm{Update}$ の入力を
$(M_t,x_t,e_t)$ に限定する。
同時刻の $r_t,v_t,u_t$ は固定した標準写像のもとで $(M_t,x_t,e_t)$ から
定まるので、その参照のために包括的型の独立引数を追加する必要はない。
ただし、各成分更新がその情報を参照できるかは下記の詳細型に制約される。
遷移後に新しく取得される観測は、次時刻の入力として扱い、先取りしない。

\medskip
\noindent\textbf{(2) 構成要素別更新
（仮定~\ref{ass:M_compact} の直積分解のもとで）：}
\[
\begin{aligned}
  \mathrm{Update}_1
  &: \mathcal{M}_1\times X\times E_{\mathrm{ext}}
  \to\mathcal{M}_1,\\
  \mathrm{Update}_2
  &: \mathcal{M}_2\times X\times E_{\mathrm{ext}}
  \times\mathcal{M}_1
  \to\mathcal{M}_2,\\
  \mathrm{Update}_3
  &: \mathcal{M}_3\times X\times E_{\mathrm{ext}}
  \times\mathcal{M}_1\times\mathcal{M}_2
  \to\mathcal{M}_3,
\end{aligned}
\]
\[
\begin{aligned}
  M_{1,t+1}
  &=
  \mathrm{Update}_1(M_{1,t},x_t,e_t),\\
  M_{2,t+1}
  &=
  \mathrm{Update}_2(
    M_{2,t},x_t,e_t,M_{1,t+1}
  ),\\
  M_{3,t+1}
  &=
  \mathrm{Update}_3(
    M_{3,t},x_t,e_t,
    M_{1,t+1},M_{2,t+1}
  ).
\end{aligned}
\]

$\mathrm{Update}_2$ が $M_{1,t+1}$ に、
$\mathrm{Update}_3$ が $M_{1,t+1}$ および
$M_{2,t+1}$ に依存することは、
TIPA における情報依存関係
\[
  M_1\to M_2\to M_3
\]
を反映する
（\S\ref{sec:tipa}）。
この順序は同一ステップ内の情報依存関係を表すものであり、
TIPA 三機構が実時間上で逐次的にのみ作動することを
意味しない。

ここで
\[
  M_t
  =
  (M_{1,t},M_{2,t},M_{3,t})
  \in\mathcal{M}
\]
に対し、7タプルに含まれる包括的更新作用素を
\[
  \mathrm{Update}(M_t,x_t,e_t)
  :=
  \bigl(
    M_{1,t+1},
    M_{2,t+1},
    M_{3,t+1}
  \bigr)
\]
と定める。

したがって、$\mathrm{Update}_2$ および
$\mathrm{Update}_3$ に現れる追加引数は、
同一ステップ内の先行する構成要素更新によって生成される
中間値であり、包括的写像への独立した外部入力ではない。

\medskip
\noindent\textbf{(3) 一ステップ更新量の参照：}

$\mathrm{Update}$ によって生じた一ステップの内部モデル変化は、
定義~\ref{def:tipa_component_change} の
$\delta_{M,c,t}$、$\boldsymbol{\delta}_{M,t}$ および
$\delta_t$ によって記述する。
本小節はこれらを再定義しない。
$\boldsymbol{\delta}_{M,t}$ を一次的な変化記述として用い、
$\delta_t$ は三機構のうち最大の一ステップ変化を表す
二次的な集約指標として用いる。
したがって $\delta_t$ は、内部モデル全体の成分別変化パターンを
余すところなく表す量ではない。
\end{definition}

\begin{remark}[$\mathrm{Update}$ の環境引数について]
\label{rem:update_env_argument}
$\mathrm{Update}$ の環境引数は実現外部環境状態
$e_t\in E_{\mathrm{ext}}$ であり、
内部環境推定 $\hat{e}_t$
（定義~\ref{def:internal_env_estimate}）ではない。

これは、$\mathrm{Update}$ をシステム記述の水準で定める
Phase~1 の抽象化である。
機構水準において、各機構が実際に参照する入力が
実現外部環境状態であるか内部環境推定であるかは
\S\ref{sec:tipa} で扱う。
本小節の包括的型だけでは観測経由性は課されない。
\S\ref{subsec:s5_observation} ではこれを追加仮定として明示する。

また、内部表現の空間
$\mathcal{R}_{\mathrm{int}}$
（定義~\ref{def:three_primitives}）は
$\mathrm{Update}$ の独立引数に含まれない。
その理由は備考~\ref{rem:M_encodes_representation} に述べたとおりであり、
内部表現から内部モデル更新への具体的な情報経路は
\S\ref{sec:tipa} で扱う。
\end{remark}

\begin{definition}[Update整合系列]
\label{def:update_consistent_sequence}

系列
\[
  (M_t,x_t,e_t)_{t\geq0},
  \qquad
  M_t\in\mathcal{M},\quad
  x_t\in X,\quad
  e_t\in E_{\mathrm{ext}}
\]
が、各時刻 $t$ において
\[
  M_{t+1}
  =
  \mathrm{Update}(M_t,x_t,e_t)
\]
を満たすとき、これを\emph{Update整合系列}と呼ぶ。

本小節で内部モデルの変化量および更新に関する指標を述べる際には、
特に断らない限り Update整合系列を対象とする。
Update整合性の定義は、各時刻において
$\mathrm{WF}_{i,t}=1$ が成立することを要件として含まない
（定義~\ref{def:wf_realization}）。
\end{definition}

\subsubsection{軌道上の学習停止}

\begin{definition}[系列上の学習停止]
\label{def:update_stop_on_trajectory}
定義~\ref{def:update_consistent_sequence} の Update整合系列に沿い、
\[
\delta^{\mathrm L}_{c,t}:=d_c^{\mathrm L}(M^{\mathrm L}_{c,t+1},M^{\mathrm L}_{c,t}),\qquad
\delta^{\mathrm L}_t:=\max_{c=1,2,3}\delta^{\mathrm L}_{c,t}
\]
を学習成分の一段変化量（分類C）とする。
すべての $t\in\mathbb N_0$ で $\delta^{\mathrm L}_t=0$ のとき、この系列上で
\emph{学習停止}が成立するという。これは、すべての $c,t$ で
$M^{\mathrm L}_{c,t+1}=M^{\mathrm L}_{c,t}$ と同値である。
作動成分の変化は排除せず、$\delta_t=0$ と同一視しない。
$0\le\delta^{\mathrm L}_t\le\delta_t\le1$ であり、全状態の不変は学習停止を含意するが、
逆は一般には成り立たない。
写像全体で学習射影が恒等、すなわち
$\operatorname{pr}_L\mathrm{Update}(M,x,e)=\operatorname{pr}_L M$
ならば任意の整合系列で学習停止となる。
本論文は、一系列上の学習停止から写像全体のこの性質を一般的な帰結として置かない。
ここで $\operatorname{pr}_L:\mathcal M\to\mathcal M^{\mathrm L}$ は座標射影である。
\end{definition}

\subsubsection{時間平均内部モデル変化量}

\begin{definition}[時間平均内部モデル変化量（分類C）]
\label{def:update_quality}

以下の
$\lambda_{T,t}$、
$\lambda_{M,c,T,t}$ および
$\boldsymbol{\lambda}_{M,T,t}$ は、対象ACS $\Sigma$ について固定された
一つのUpdate整合系列
（定義~\ref{def:update_consistent_sequence}）に沿って定める。
7タプル $\Sigma$ は、初期状態 $x_0$、初期内部モデル状態 $M_0$、
環境実現（定義~\ref{def:env_realization}）および行動列を
一意に定めない（備考~\ref{rem:exogenous_environment}）。
したがってこれらの指標は $\Sigma$ のみからは定まらない。
本論文では評価系列への依存性を記法上省略する。
括弧内の $\Sigma$ は対象ACSを識別するものであり、
7タプルのみからこれらの値が一意に定まることを
意味しない。

\medskip
\noindent\textbf{最大成分時間平均更新速度と成分別時間平均更新速度：}

開始時刻 $t$、時間窓長 $T$ に対する
最大成分時間平均更新速度を
\[
  \lambda_{T,t}
  =
  \frac{1}{T}
  \sum_{s=t}^{t+T-1}\delta_s
\]
とする。

ここで $\delta_s$ は、時刻 $s$ における
TIPA 三機構の成分別変化量
$\boldsymbol{\delta}_{M,s}$ の最大値である
（定義~\ref{def:tipa_component_change}）。
したがって $\lambda_{T,t}$ は、
三機構すべての変化量を加算した総変化量ではなく、
各時刻における最大成分変化の時間平均を表す。

同じ時間窓について、成分別時間平均更新速度を
\[
  \lambda_{M,c,T,t}
  :=
  \frac{1}{T}
  \sum_{s=t}^{t+T-1}
  \delta_{M,c,s},
  \qquad
  c=1,2,3
\]
と定義し、
\[
  \boldsymbol{\lambda}_{M,T,t}
  :=
  \bigl(
    \lambda_{M,1,T,t},
    \lambda_{M,2,T,t},
    \lambda_{M,3,T,t}
  \bigr)
\]
を TIPA 三機構の成分別時間平均更新速度とする。

$\boldsymbol{\lambda}_{M,T,t}$ は成分別変化を集約せずに
記述する一次的指標であり、
$\lambda_{T,t}$ は最大成分変化を集約した
スカラー指標として区別して用いる。
これらは本論文が分析のために導入する記述量であり、
当該ACSがこれらを内部に保持することを要求しない。

開始時刻が $0$ の場合には、
\[
  \lambda_T
  :=
  \lambda_{T,0},
\]
\[
  \lambda_{M,c,T}
  :=
  \lambda_{M,c,T,0},
  \qquad
  \boldsymbol{\lambda}_{M,T}
  :=
  \boldsymbol{\lambda}_{M,T,0}
\]
と略記する。

\medskip
これらは、時間窓において実際に生じた内部モデル変化の記述量であり、
当該ACSが原理的に達成しうる最大の更新速度その他の容量を表さない
（\S\ref{subsec:cognitive_capacity_policy}）。
具体的な実装系または観測データに対する
実証的測定・推定方法は Phase~2 以降の課題とする。
\end{definition}

\begin{remark}[環境変化応答機能；構成的必要機能候補]
\label{rem:env_change_response}

本小節までに、原始的な機能の標準写像形式
$P_{\mathrm{sens}},P_{\mathrm{disc}},P_{\mathrm{val}},P_{\mathrm{act}}$
（定義~\ref{def:three_primitives}）、内部モデル構造 $M$
（定義~\ref{def:internal_model}）および更新作用素
$\mathrm{Update}$（定義~\ref{def:update_operator}）が
\S\ref{sec:acs_def} 内で出そろった。
本備考では、これらを横断して成立しうる機能を
\emph{環境変化応答機能}として記録する。

すなわち、外部環境状態の変化に対して、ACSが利用可能な感知、
内部表現形成、評価、予測、内部モデル更新および行動形成の
一部または複数を用い、現在または将来の応答を変化させうる
横断的機能をいう。
本備考は新たな写像、変数または応答時間量を定義しない。

この機能は、必要環境要因 $\eta\in E$
（\S\ref{subsec:environment_space}）とは理論階層が異なる。
$\eta$ が記述するのは、ACSが存続するために、
内部状態と外部環境状態の組において
何が満たされていなければならないかであり、
その外部環境状態への依存は構成要件 CR-1 の (CR-1c) によって拘束される。
所与の必要環境を自己側で利用しうるか否か
（必要環境利用機能、備考~\ref{rem:env_utilization_function}）は、
$\eta$ の内部状態への依存として表されうる。
これに対し本備考が記述するのは、外部環境が変化したときに
ACS側が利用可能な機能をどのように用いて応答しうるかである。
したがって本機能は、備考~\ref{rem:eta_coarse_classes} の
標準的粗分類候補 $\mathsf{Cls}_\eta$ に加えるべき第五のラベル
ではなく、これと同一の一覧に並べない。

以下の限定を付す。

\begin{enumerate}[label=(\alph*)]
  \item 本備考は、処理の固定系列を要求しない。
  感知から行動に至る各段階を毎時刻すべて経ることを
  要求せず、変化の検出から局所急性異常への保護応答に至る急性の応答と、
  識別・評価を経て必要に応じた内部モデル更新を行い後続の応答を
  変更する適応的な応答の双方を包含しうる。
  それぞれに対応しうる標準機能形式は
  \S\ref{sec:tipa} で扱う。

  \item 本備考は
  \[
    e\neq e'
    \;\Longrightarrow\;
    \mathrm{Update}(M,x,e)\neq\mathrm{Update}(M,x,e')
  \]
  を要求しない。
  外部環境状態が変化しても内部モデル更新を要しない場合があり、
  また現在の行動の変更のみによって応答が実現される場合もある。
  この扱いは、$f$ の行動引数に関する非退化性を一般には
  要求しない本論文の方針
  （仮定~\ref{ass:f_decomposition}、
  \S\ref{subsec:state_space_basic_variables}）、および
  恒等的な更新を排除しない本小節の扱い
  （定義~\ref{def:update_stop_on_trajectory}）と整合する。

  \item 本論文は、本備考から、整形式実現の維持
  （$\mathrm{WF}_{i,t}=1$、定義~\ref{def:wf_realization}）、
  可行性 $\Pi(x_t)\geq0$、非閉塞 $A_t\neq\emptyset$
  （定義~\ref{def:blockage_sec2}）、または非崩壊
  （定義~\ref{def:collapse}）を一般的な帰結として置かない。
  すなわち、応答しうる機能の存在から存続の成功を一般的な帰結として置かない。

  \item 外部からの変化が応答に先立ってACSの構造を
  回復不能に損なう場合、当該ACSが存続しないことをもって
  応答しうる機能の欠如とは判定しない。

  \item 本論文は、環境変化の後に応答すべき一律の期限を
  Phase~1 では課さない。
  応答の余地が存在する時間尺度は、変化の種別および実装に
  依存する。

  \item 予測されていなかった変化を含みうるが、
  本論文は、環境変化応答機能に一般に適用する予測外の判定条件および許容誤差を
  Phase~1 で形式化していないため、この意味での「予測外の変化」を形式的に定義しない
  （備考~\ref{rem:estimate_gap}）。
  \S\ref{sec:tipa} の予測整合性誤差と二種類の予測外は、M1 の出力と保持記録に相対して定める量である
  （定義~\ref{def:m1_temporal_features}、定義~\ref{def:prediction_record_check}）。
\end{enumerate}

本備考は構成的必要機能の候補としての記録であり、
現段階では次のいずれも行わない。
第一に、$\mathsf{Cls}_\eta$ への追加。
第二に、構成要件 CR-2（\S\ref{subsec:three_primitive_functions}）
への原始的な機能の新たな区分としての追加。
第三に、新たな構成要件の導入。
第四に、分類A表記の付与。

これらの判断には、存続に関わる環境変化の形式的定義、
応答を要しない環境差との区別、応答の余地が存在しない変化の
除外方法、内部モデル更新を伴わない応答の扱い、
即時の保護応答と適応的応答に共通する抽象条件、
恒等的な更新をもつ実現を排除する条件、および
応答しうる機能と応答の成功との分離を要する。
これらの検討は Phase~1.$\alpha$ または Phase~2 の課題とする
（\S\ref{subsec:phase1a}、within present framework）。
\end{remark}

%----------------------------------------------------------------------
\subsection{認知的処理能力の記述方針}
\label{subsec:cognitive_capacity_policy}

本小節では、ACSの認知的処理能力について Phase~1 で採用する記述方針を述べる。
本論文は Phase~1 において、認知的処理能力を形式的に定義せず、
単一の量としても導入しない。

\begin{remark}[認知的処理能力の記述方針]
\label{rem:cognitive_capacity_policy}
本論文は、認知的処理能力を単一のスカラー量として定義しない。
その形式化にあたっては、少なくとも次の三つを区別する必要がある。

\begin{description}[style=nextline]
  \item[量的容量]
  一時刻あたりに保持・統合・更新の対象としうる情報または表現の規模。
  \item[時間応答性]
  新しい情報、誤差または環境の変化を内部表現や評価へ反映するまでの速さ。
  離散時間の本枠組みでは、所定の統合・訂正の基準に到達するまでの
  段数（潜時）が形式化の候補となる。
  時間応答性は、一時刻あたりの処理の規模とは別の量である。
  \item[統合の質]
  新しい情報と既存の内部表現を、対象の状態や評価の対象との整合を保ちながら
  統合・更新できる程度。
  状態推定と評価の正確性、および誤差の検出と訂正を含みうる。
\end{description}

これらに関して、本論文は次の区別を置く。
第一に、容量、すなわちある条件のもとで達成しうる範囲と、
実現量、すなわち実際に生じた変化とを区別する。
定義~\ref{def:update_quality} の時間平均内部モデル変化量は後者の記述量であり、
認知的処理能力を表さない。
第二に、能力とその使い方の偏向とを区別する。
原始的な機能の責務の遂行（そのための資源・評価・注意その他の配分または利用を含む）に現れる偏向性の程度は、
\S\ref{sec:extended_dynamics} の偏向性（定義~\ref{def:bias_tendency}）として記述し、
その程度を偏向強度 $\beta_{\iota,t}$（定義~\ref{def:beta_accumulated}）で表して、
認知的処理能力には含めない。
第三に、三つの区分を一つの実数へ集約する一般式も、
区分の間の全順序も置かない。
第四に、認知的処理能力を、\S\ref{sec:tipa} の機能軸が表す作動の程度、
および \S\ref{sec:extended_dynamics} の健全度が表す基盤の状態と
同一視しない。

本備考は新たな記号、写像、構成要件または仮定を導入しない。
三つの区分の形式化と測定は Phase~1.$\alpha$ または Phase~2 の課題とする
（\S\ref{subsec:limitations_list} の L3、\S\ref{subsec:phase1a}）。
\end{remark}

%----------------------------------------------------------------------
\subsection{整形式性・構成要件・仮定の整理}
\label{subsec:assumption_list_sec2}

本小節では、\S\ref{sec:acs_def} で導入した規約、構成要件および
仮定を一覧として整理する。
本小節はこれらを新たに定義する場所ではない。
各対象の正式な定義は、それぞれの意味論的所有先である小節に
置かれており、本小節は所在と地位を再掲するにとどめる。

本論文は、ACS定義を支える条件を次の三種に区別する。
第一に規約であり、これは本論文が記述のために採用する
取り決めであって、当該ACSが満たすべき条件ではない。
第二に構成要件であり、これは当該7タプル構造体の実現が、
本論文におけるACSとして認められるために満たすべき
意味論的条件である。
第三に仮定であり、これは後続命題の成立範囲を限定するために
導入される条件である。
規約と構成要件は仮定一覧に含めない。

\subsubsection*{規約および構成要件}

独立した項目として本小節で再掲する規約は二件である。

\begin{description}[style=nextline]
  \item[記述水準の規約（\S\ref{subsec:acs_7tuple}）]
  抽象構造型、その実現である個体 $\Sigma_i$、
  および時刻別の値を書き分け、個体の実現に依存する量を
  時刻に依存しない写像とその時刻評価値の二段で書く規約。
  \item[型分離規約 TSC（\S\ref{subsec:wellformedness}）]
  内部状態型 $X$ と外部環境状態型 $E_{\mathrm{ext}}$ を
  集合表現上でも分離するために採用する規約。
  タグ付きの互いに素な表現はその実装手段であり、
  TSC が要求するのは型分離そのものである。
\end{description}

構成要件は二件である。

\begin{description}[style=nextline]
  \item[構成要件 CR-1（必要環境整合性、
  \S\ref{subsec:environment_space}）]
  必要環境空間 $E$ の各元が当該ACSの存続に真に必要であること
  (CR-1a)、必要要因の全体が同時に充足されうること
  (CR-1b)、および各要因について、開始時点の内部状態を固定したとき
  外部環境状態の違いによって充足・非充足が分かれる持続的な非充足が
  実現しうること (CR-1c)。
  \item[構成要件 CR-2（原始的な機能と基盤機能の具備、
  \S\ref{subsec:three_primitive_functions}）]
  感知・識別・評価・行動の四区分のそれぞれについて基礎的な機能を具備すること (CR-2a)、
  それらが機能的に退化していないこと (CR-2b)、
  およびそれらの活動を支える基盤機能を具備すること (CR-2c)。
\end{description}

CR-1 および CR-2 は、いずれも7タプルから機械的に検査できる
条件ではなく、その実現に対する意味論的条件である。
両者は対象が異なり、本論文は一方から他方を一般的な帰結として置かない。
これらの充足を有限観測から判定する手続きは Phase~1 では
与えない。

\subsubsection*{仮定一覧}

本節で導入した仮定は九件である。
導入地位の記号は、A が公理的、S が構造的、O が観察的を表す。
各仮定の内容、導入地位および強さは、それぞれの所在小節に
記載したとおりである。
「主な依拠箇所」は、当該仮定を参照する主な定義・命題・備考を示し、
参照の全件ではない。

\medskip
{\small
\begin{center}
\begin{tabular}{p{1.4cm} p{4.0cm} p{6.0cm} p{0.8cm} p{0.8cm}}
\hline
\textbf{仮定} & \textbf{所在} & \textbf{主な依拠箇所} & \textbf{地位} & \textbf{強さ} \\
\hline
\ref{ass:no_branching} & \S\ref{subsec:self_formal_definition}（構造的自己と時間的自己同一性） & 定義~\ref{def:time_identity} & A & 強 \\[2pt]
\ref{ass:compact_spaces} & \S\ref{subsec:state_space_basic_variables}（状態遷移） & 定義~\ref{def:acs}、命題~\ref{prop:At_compact}、定義~\ref{def:x_E}、定義~\ref{def:update_section5} & S & 強 \\[2pt]
\ref{ass:f_continuity} & 同上 & 定義~\ref{def:acs}、命題~\ref{prop:At_compact} & S & 中 \\[2pt]
\ref{ass:Eext_time_fixed} & 同上 & 備考~\ref{rem:fixed_environment_and_state_dependent_viability} & S & 中 \\[2pt]
\ref{ass:f_decomposition} & 同上 & 定義~\ref{def:acs}、備考~\ref{rem:oplus_neutrality}、備考~\ref{rem:blockage_auto} & S & 強 \\[2pt]
\ref{ass:env_input_nondegeneracy} & 同上 & 命題~\ref{prop:nondegeneracy_implication} & S & 弱 \\[2pt]
\ref{ass:env_viable_nondegeneracy} & \S\ref{subsec:viability_function_and_blockage}（可行性） & 命題~\ref{prop:nondegeneracy_implication}、命題~\ref{prop:feasibility_sensitivity_scope} & S & 弱 \\[2pt]
\ref{ass:viable_realizability} & 同上 & 備考~\ref{rem:viable_realizability_consequences} & S & 中 \\[2pt]
\ref{ass:M_compact} & \S\ref{subsec:internal_model_M}（内部モデル） & 定義~\ref{def:model_parameter_space}、定義~\ref{def:update_operator} & S & 強 \\
\hline
\end{tabular}
\end{center}
}

\medskip
このうち強さを強とするものは、
仮定~\ref{ass:no_branching}、
仮定~\ref{ass:compact_spaces}、
仮定~\ref{ass:f_decomposition} および
仮定~\ref{ass:M_compact} の四件である。
これらは、緩和すると本節の主要定義または後続節の主要命題の
成立範囲が変わりうる仮定である。
その弱化または代替定式化の検討は
\S\ref{sec:limitations_and_outlook} で扱う。

仮定~\ref{ass:M_compact} の強さ「強」は、
各 $\mathcal{M}_c$ の非空コンパクト性にかかる。
同仮定に含まれる距離の尺度正規化は
計量上の規約として扱い、
この強さの評価対象には含めない。

\subsubsection*{本節が主張しない事項と Phase 委譲}

本節は閉塞、機能的終端および崩壊を定義するが、単独では次のいずれも主張しない。
閉塞、機能的終端または崩壊への到達、整形式実現の終局的喪失、
終局的な構造的自己喪失、
社会的相互作用、確率的安定性、実装可能性および測定可能性で
ある。

以下は Phase~1 の枠内では形式化しない。
各項目は、後続 Phase への委譲事項、または将来の拡張・再検討の
候補として保留する。
本論文は、これらのすべてが後続 Phase において形式化されることを
主張しない。

\begin{enumerate}[label=(\arabic*)]
  \item 確率的拡張、連続時間化、および非コンパクト内部状態空間への
  一般化。

  \item 必要環境要因の経験的同定手続き、および同定誤差の扱い
  （備考~\ref{rem:necessity_identification}）。
  必要性の存立と同定可能性は区別される。

  \item 分岐を許す場合の時間的自己同一性の再定式化
  （仮定~\ref{ass:no_branching} の緩和、
  \S\ref{sec:limitations_and_outlook}）。

  \item 整形式実現指標 $\mathrm{WF}_{i,t}$
  （定義~\ref{def:wf_realization}）の値を有限観測から
  判定する操作的手続き。
  $\mathrm{WF}$ と機能的終端指標 $\mathrm{FT}_{i,t}$・崩壊指標 $\mathrm{Col}_{i,t}$ の関係は
  備考~\ref{rem:collapse_wf_relation} で述べた。

  \item 機能的終端の条件 (FT-1)・(FT-2) および崩壊の条件 (SC-1)・(SC-2) の
  判定規則の一般形式、ならびに (SC-1) が参照する、構造的個体としての同一性を
  成立させる組織関係の一般形式
  （定義~\ref{def:functional_termination}、定義~\ref{def:collapse}）。
  本論文はこれらを対象系ごとに固定されるものとして扱う。

  \item 自己--環境境界の動的再編成、および
  時間依存する外部環境状態空間への拡張。

  \item 作用差（定義~\ref{def:action_difference}）の大きさ、
  多段階にわたる作用、行動や行動応答成分の寄与の比較、
  および外部と内部に共通する資源量の収支。
  背景成分の列の生成則と、他のACS個体との結合動学は Phase~2 の課題とする
  （備考~\ref{rem:exogenous_environment}、\S\ref{subsec:limitations_list} の L2）。

  \item 実現外部環境状態 $e_t$ と内部環境推定 $\hat{e}_t$ の
  乖離の定量化（備考~\ref{rem:estimate_gap}）。

  \item 必要に応じた、内部表現の空間
  $\mathcal{R}_{\mathrm{int}}$ から内部モデル状態空間
  $\mathcal{M}$ への明示的な符号化写像の導入
  （備考~\ref{rem:M_encodes_representation}）。
  本論文は、そのような写像を後続 Phase で導入することを
  約束しない。

  \item 環境変化への横断的応答能力
  （備考~\ref{rem:env_change_response}）を新たな構成要件として
  形式化すること。
  環境変化の存続関連性、応答の余地が存在する時間尺度、
  予測されていなかった変化の判定、修繕・代替・回復の経路、
  および当該能力をACS成立条件へ昇格させるための必要十分条件は、
  Phase~1.$\alpha$ または Phase~2 の検討対象とする（\S\ref{subsec:phase1a}）。

  \item 必要環境要因の分類粒度の確定
  （備考~\ref{rem:eta_coarse_classes}）。
  標準的粗分類候補 $\mathsf{Cls}_\eta$ は分割ではなく、
  網羅性および排他性を証明していない。
\end{enumerate}

\subsubsection*{三原理への接続}

本節では、ACSを7タプル
$\Sigma=(X,U,E,f,\Pi,M,\mathrm{Update})$ として定義し、
外部環境状態空間 $E_{\mathrm{ext}}$ に相対的な記述の枠組み、
構造的自己・整形式実現と崩壊、必要環境空間、状態遷移、
持続評価と可行性、原始的な機能、内部モデル構造、
内部モデル更新、および認知的処理能力の記述方針を与えた。

\S\ref{sec:three_principles} では、これらの定義のもとで、
自己維持（SM）・環境利用（EU）・可行応答（FR）の三原理を形式的に定義し、
各原理の概念的定義と時刻別充足述語を区別したうえで、
時刻別充足述語相互の条件付き含意および非含意を扱う。


%================ Section 3 三原理：SM・EU・FR ================
%======================================================================
\section{三原理：SM・EU・FR}
\label{sec:three_principles}

本節では、自律認知システム（ACS）の長期的な可行性に関わる三つの制約軸、
すなわち自己維持（SM）・環境利用（EU）・可行応答（FR）を、
\S\ref{sec:acs_def} で定めた自己--環境構造と可行性構造の上に形式化する。
本節の全結果は Phase~1 の決定論的・離散時間枠組みに限定され、
各主張の強さは対応する定義・仮定・証明の強さを超えない。

\begin{definition}[三原理の概念的定義]
\label{def:three_principles_concept}

本論文における\emph{三原理}とは、
自己--環境境界をもつ ACS が、
内部状態と外部環境状態との相互作用のもとで時間的に存続する際に問題となる
可行性維持上の制約を、
自己維持（Self-Maintenance; SM）、
環境利用（Environmental Utilization; EU）、
可行応答（Feasible Response; FR）
の三つの相異なる制約軸として整理したものである。

\begin{description}[style=nextline]
  \item[SM]
  ACS の内部状態を可行集合 $K$ の内に保つことに関わる制約軸である。
  \item[EU]
  一段階の可行性に関して臨界となる必要環境要因が、その局面で充足されることに関わる制約軸である。
  \item[FR]
  必要な局面で可行な応答が得られること、
  およびその必要応答ポテンシャルの保持・獲得・回復・再構成に関わる制約軸である。
  ここでいう「可行」は、本論文の可行性、すなわち次状態が $K$ に属することに固定する
  （備考~\ref{rem:fr_scope}）。
\end{description}

これら三つの制約軸は、
自己--環境関係をもつ ACS を時間的な存続構造として記述するとき、
同時に意味をもつ。
ただし、本論文は、三原理の時刻別充足述語が同時に真であることを要求しない。
時刻別充足述語の偽と閉塞との関係は原理ごとに異なる。
EU の時刻別充足述語が偽である場合、または EU が真で FR の時刻別充足述語が偽である場合には、いずれも閉塞を伴う。
判定時刻の閉塞は、EU が偽である場合と、EU が真で FR が偽である場合とに排他的に分かれる
（命題~\ref{prop:blockage_dichotomy}）。
SM の時刻別充足述語の偽については、本論文は閉塞を一般的な帰結として置かない。
いずれの時刻別充足述語の偽からも、本論文は崩壊を一般的な帰結として置かず、
三原理が ACS の存続に関わる制約を網羅するとも主張しない。

三原理は、ACS に外部から課される価値規範または目的関数ではなく、
本論文の Phase~1 において、
自己--環境相互作用下の時間的存続を記述するために採用する
記述的な制約軸である。
本定義は三原理という総称の概念的整理であり、
原理そのものについて成立・不成立を述べない。
各原理について真偽をもつのは、
定義~\ref{def:SM}、定義~\ref{def:EU} および定義~\ref{def:FR} の
時刻別充足述語である。
\end{definition}

%----------------------------------------------------------------------
\subsection{記号・規約と記述水準}
\label{subsec:s3_setup}

本小節は、以降の各小節が繰り返し参照する記述上の規約と、
時刻別充足述語の適用域を一度だけ与える。
本節で新たに用いる主要記号の一覧は付録~\ref{app:symbols_s3} に示す。

\subsubsection*{三原理の構造的前提}

\S\ref{sec:acs_def} で導入された ACS は、
固定された外部環境状態空間 $E_{\mathrm{ext}}$ を背景型として、
内部状態型 $X$ と外部環境状態型 $E_{\mathrm{ext}}$ を区別し、
状態遷移 $f:X\times U\times E_{\mathrm{ext}}\to X$ によって
自己と外部環境との相互作用を記述する。
必要環境空間 $E$ と必要環境充足集合 $E_{\mathrm{req}}(x)$ は、
構造的自己の存続が外部環境条件と無関係ではないことを記述する。

本節では、この自己--環境構造を所与として、
現在の可行性、
一段階の可行性に関して臨界となる必要環境要因の充足、
必要な局面での可行な応答という
三つの問題を SM、EU、FR として形式化する。
Phase~1 は自己--環境境界の生成過程、
個体発生過程、または境界形成の物理機構を扱わない。

また、本節で導入される評価関連量については、
\S\ref{subsec:description_levels_global} の全論文共通規約に従い、
ACS 側に必要となる構造的評価機能と、
それを論文上で記述する数学的形式表現とを区別する。
本論文は、形式量の存在から、
ACS がその真値を完全に知ること、
または同一の数式を内部実装することを一般的な帰結として置かない。

\subsubsection*{本節共通の記述規約}

\begin{remark}[含意の扱いに関する規約：時刻別充足述語・評価量と失敗状態]
\label{rem:conv_nonimplication}
本節は、いずれかの時刻別充足述語が偽であること、または T-STAB の不充足だけから、
整形式実現の不成立 $\mathrm{WF}_{i,t}=0$ または崩壊 $\mathrm{Col}_{i,t}=1$ を
一般的な帰結として置かない。
閉塞 $A_t=\emptyset$ については次のとおりである。
\begin{enumerate}
  \item $\mathrm{SM}_t$ の偽、T-STAB の不充足、
  および $\chi_{\mathrm{EI},t}=0$ の一時刻の成立だけからは、
  閉塞を一般的な帰結として置かない。
  \item $\mathrm{EU}_t$ が偽ならば $A_t=\emptyset$ である（命題~\ref{prop:eu_blockage}）。
  \item $\mathrm{FR}_t$ が偽ならば $A_t=\emptyset$ である（定義~\ref{def:FR}）。
\end{enumerate}
関係品質 $Q$ については、定義される場合の正・零・負の符号と、
そもそも評価不能である場合とを区別する。
時刻域 $\mathbb{T}_{\mathrm{P},i}$ を明示した一覧は
備考~\ref{rem:principles_failure_table} に示す。
\end{remark}

\begin{remark}[記述水準規約：形式表現と内部評価の非同一性]
\label{rem:conv_description_level}
本節が導入する形式量は \S\ref{subsec:description_levels_global} の規約に従う理論記述であり、
ACS がその真値や集合全体を内部で完全に保持・計算することを要求しない。
内部評価・応答に用いられるのは、当該 ACS が実際に利用可能な情報である。
構造的評価機能に要求する非退化性は、\S\ref{subsec:self_env_principles} で述べる。
\end{remark}

\begin{remark}[非規範性規約]
\label{rem:conv_nonnormative}
三原理および本節の評価量は、長期可行性に関わる記述的な制約軸と形式表現であり、
外部から課す価値規範・単一の目的関数・最大化目標ではない。
このことは、時刻別充足述語の真を道徳的な善悪と同一視することも、
$Q$ や T-STAB を単一効用関数とみなすことも要求しない。
\end{remark}

\begin{remark}[非網羅・非最小規約]
\label{rem:conv_nonexhaustive}
本節は、三原理が存続に関わる全制約を網羅すること、三分解が唯一であること、
三原理の時刻別充足述語が互いに冗長でないこと、
および三原理の最小性を主張しない。
三原理の間の関係は、
\S\ref{subsec:principle_linkage} の連動と含意関係の限定として述べる。
\end{remark}

\begin{remark}[構造的環境結合・実現相互作用・EU・関係評価の水準区別]
\label{rem:conv_ro_layers}
本節では、次の数学的対象を区別する。
$\mathrm{EC}(x_t)$ は定義~\ref{def:env_coupling} の構造的環境結合条件、
$\chi_{\mathrm{EI},t}$ は定義~\ref{def:phi} の実現外部相互作用を表す二値指標、
$e_t\in E_{\mathrm{req}}(x_t)$ は現在の必要環境充足を表す命題、
$\mathrm{EU}_t$ は定義~\ref{def:EU} の時刻別充足述語、
$Q_{i\leftarrow j,t}$ は \S\ref{sec:rel_min} で定める関係品質であり、定義域上でのみ値をもつ。
これらを同一視せず、同値性または含意を用いる場合は個別に明示した条件による。
また、命題~\ref{prop:pi_E_sensitivity} の一般化勾配感応性は、
構造的環境結合条件とは別の局所的な可行性感応条件である。
\end{remark}

\begin{remark}[T-STAB 地位規約]
\label{rem:conv_tstab_scope}
T-STAB は三原理関連評価の有限窓時間統合条件であり、
力学系論の Lyapunov 安定性・漸近安定性・構造安定性、
三原理に追加される第四原理、外生的目標関数のいずれでもない。
その充足のみから、無限時間可行性、収束、吸引性、不変集合の存在、
または固定点への到達を一般的な帰結として置かない。
\end{remark}

\subsubsection{記述水準と時刻別充足述語の適用域}
\label{subsubsec:s3_levels}

\begin{remark}[本節における記述水準]
\label{rem:s3_description_levels}
本節では、次の水準を区別する。
\begin{description}[style=nextline]
  \item[原理]
  自己維持（SM）・環境利用（EU）・可行応答（FR）は、存続中のACSに対する制約軸である。
  本節は、原理そのものについて成立・不成立を述べない。
  \item[原理応答項目]
  各原理について、ACSの側で何が成り立っていることが問われるかを述べる項目である。
  手段は書かない。
  \item[時刻別充足述語]
  各原理について時刻ごとに定める述語 $\mathrm{SM}_t$・$\mathrm{EU}_t$・$\mathrm{FR}_t$ である。
  真偽をもつのはこの水準である。
  本節は判定を「真である」「偽である」と書き、
  「偽である」は、述語が定義されたうえで偽となることを指す。
  述語が定義されない時刻を、偽である時刻に含めない。
  \item[構造的機能と標準機能形式]
  \S\ref{sec:acs_def} の基礎的な機能と基盤機能、
  本節の三原理関連構造的機能（分類A）、
  および \S\ref{sec:tipa} の TIPA・PSL・ISL（分類B）である。
  作動・停止・低下を述べる対象はこの水準である。
  \item[診断量]
  応答多様度 $\mu_{\mathrm{eff},t}$、頑健存続可能段数、評価尺度 $s_{\cdot,t}$ その他の分類C量である。
\end{description}
本論文は、機能が作動していることから時刻別充足述語が真であることを、
時刻別充足述語が偽であることから特定の機能の停止を、いずれも一般的な帰結として置かない。
原理ごとの専用の機能は置かない。
一つの原理応答項目を複数の構造的機能が共同して担うこと、
および一つの構造的機能が複数の原理応答項目に関与することを排除しない。

これらとは別の軸として、本節の理論上の量（$\Pi$、$A_t$、$\mathrm{Crit}^{(1)}(x_t)$ など）、
ACS 内部でのそれらに関わる推定・評価、
および外部観測者による測定・診断を区別する。
時刻別充足述語の真理条件は理論上の量だけで定まり、
内部の推定・評価の結果も、外部観測者の測定結果も用いない
（備考~\ref{rem:crit_description_level}）。
\end{remark}

本節は単一のACS個体 $\Sigma_i$ を扱うため、
\S\ref{sec:acs_def} の規約に従い、
状態・環境・可行行動集合等の個体添字を省略して
$x_t$、$e_t$、$A_t$ と書く。
個体の同一性に関わる指標
$\mathrm{WF}_{i,t}$、$\mathrm{Col}_{i,t}$、$\mathrm{FT}_{i,t}$ には添字を付す。

\begin{definition}[原理判定時刻域]
\label{def:principle_time_domain}
ACS個体 $\Sigma_i$ について、
\[
  \mathbb{T}_{\mathrm{P},i}
  :=
  \bigl\{
    t\in\mathbb{N}_0
    \bigm|
    \mathrm{FT}_{i,t}=0
  \bigr\}
\]
を\emph{原理判定時刻域}と呼ぶ。
ここで $\mathrm{FT}_{i,t}$ は定義~\ref{def:functional_termination} の機能的終端指標である。
本節で定める時刻別充足述語は、
$t\in\mathbb{T}_{\mathrm{P},i}$ の時刻にだけ定義する
（\S\ref{sec:acs_def} 備考~\ref{rem:collapse_not_principle_failure}）。
$t\notin\mathbb{T}_{\mathrm{P},i}$ の時刻には、
いずれの時刻別充足述語も真とも偽ともしない。
\end{definition}

\begin{lemma}[原理判定時刻域の下方閉性]
\label{lem:principle_time_domain_down}
$t\in\mathbb{T}_{\mathrm{P},i}$ かつ $s\in\mathbb{N}_0$、$s\le t$ ならば、
$s\in\mathbb{T}_{\mathrm{P},i}$ である。
また、$t\in\mathbb{T}_{\mathrm{P},i}$ ならば $\mathrm{Col}_{i,t}=0$ である。
\end{lemma}

\begin{proof}
$s<t$ かつ $\mathrm{FT}_{i,s}=1$ とすると、
命題~\ref{prop:collapse_basic} の (C-2) により $\mathrm{FT}_{i,t}=1$ となり、
$t\in\mathbb{T}_{\mathrm{P},i}$ に反する。
よって $\mathrm{FT}_{i,s}=0$ である。
$s=t$ の場合は仮定そのものである。
後半は、$\mathrm{Col}_{i,t}=1$ ならば
命題~\ref{prop:termination_order}(ii) により $\mathrm{FT}_{i,t}=1$ となることの対偶である。
\end{proof}

\begin{remark}[原理判定時刻域と整形式実現・軌道記録]
\label{rem:principle_time_domain_scope}
原理判定時刻域には $\mathrm{WF}_{i,t}=1$ を課さない。
\S\ref{sec:acs_def} は実現状態について任意の時刻で
$x_t\in X$、$e_t\in E_{\mathrm{ext}}$ とし、
機能的終端と整形式実現の間にいずれの向きの含意も仮定しない
（備考~\ref{rem:collapse_wf_relation}）。
したがって、機能的終端に至る前に一時的に $\mathrm{WF}_{i,t}=0$ となる時刻も、
時刻別充足述語の判定の対象に含まれる。
同様に、原理判定時刻域には $x_t\in K$ も課さない。
状態が可行集合の外にありながら機能的終端に至っていない時刻、
たとえば身体の状態が閾値を割り込んだが回復の途上にある生物、
あるいは電力が閾値を下回ったが復帰の手段を残すロボットに相当する時刻も、
判定の対象に含まれる（例示の地位は備考~\ref{rem:example_status} による）。
このような時刻では、$\mathrm{SM}_t$ は偽であり、
$\mathrm{EU}_t$ と $\mathrm{FR}_t$ は一段階での $K$ への回復に関わる局面を記述する
（\S\ref{subsec:eu}、\S\ref{subsec:fr}）。
また補題~\ref{lem:principle_time_domain_down} により、
$t\in\mathbb{T}_{\mathrm{P},i}$ ならば、
$\Sigma_i$ の崩壊指標列を記録する軌道記録 $\sigma$ について
$t\in I_{\mathrm{ACS}}(\sigma)$ である（\S\ref{sec:collapse_cond} 定義~\ref{def:acs_valid_time_domain}）。
逆に $I_{\mathrm{ACS}}(\sigma)$ は機能的終端後の時刻を含みうるので、
$I_{\mathrm{ACS}}(\sigma)$ を原理判定時刻域として用いない。
\end{remark}

%----------------------------------------------------------------------
\subsection{自己--環境関係と三原理の前提}
\label{subsec:self_env_principles}

本小節では、三原理が本理論においてなぜ導入されるかを、
\S\ref{sec:acs_def} で確立された自己--環境構造から整理し、
三原理それぞれの前提と統一的な記述形式を与える。
三原理の時刻別充足述語の形式的定義は
\S\ref{subsec:sm_foundational}、\S\ref{subsec:eu}、\S\ref{subsec:fr} が担う。
本小節では新たな仮定、命題または証明を置かない。

\subsubsection{自己--環境境界と相互作用の前提}

\S\ref{sec:acs_def} は、固定された外部環境状態空間
$E_{\mathrm{ext}}$ を背景型として、ACS を
\[
  \Sigma=(X,U,E,f,\Pi,M,\mathrm{Update})
\]
と記述し、型分離規約 TSC のもとで
内部状態型 $X$ と外部環境状態型 $E_{\mathrm{ext}}$ を区別する。
この区別は、自己--環境境界
$B_{SE}=(X,E_{\mathrm{ext}})$
（定義~\ref{def:b_se}）として記述される。

ここで重要なのは、$X$ と $E_{\mathrm{ext}}$ の一方を
他方から時間的に生成することではない。
Phase~1 は境界形成の時間発展を扱わず、
境界に相対して内部と外部が区別された形式構造を所与とする。
したがって、本論文でいう自己と環境の「同時的成立」は、
両者が同じ境界関係によって相対的に区別されるという
\emph{構造的同時性}を意味する。

この境界は、自己と環境の切断を意味しない。
状態遷移写像
\[
  f:X\times U\times E_{\mathrm{ext}}\to X
\]
（定義~\ref{def:state_transition}）は、
次状態を現在状態、実効行動、実現外部環境状態の関数として記述する。
したがって、外部環境状態は状態遷移へ関与しうる形式変数として
ACS の時間発展に組み込まれる。

ただし、写像の型だけから
\[
  e\neq e'
  \;\Longrightarrow\;
  f(x,u,e)\neq f(x,u,e')
\]
のような環境引数への非退化性は仮定しない。
同一の $(x,u)$ に対して異なる外部環境状態が
同一の次状態を与えることも、形式上は許される。
本論文が必要環境依存として要求する内容は、
状態遷移写像の引数型だけではなく、
必要環境空間 $E$ と構成要件 CR-1
（\S\ref{subsec:environment_space}）によって与えられる。

必要環境空間 $E$ は、
固定された $E_{\mathrm{ext}}$ のもとで、
ACS の存続に必要な環境要因を状態依存に記述する。
必要環境充足集合
\[
  E_{\mathrm{req}}(x)
  =
  \{e\in E_{\mathrm{ext}}\mid
    \forall\eta\in E,\ \eta(x,e)=1\}
\]
（定義~\ref{def:E_req}）は、
状態 $x$ において必要環境要因を同時に充足する
外部環境状態を表す。
したがって、ACS の構造的自己は外部環境から区別される一方、
その存続条件の一部は外部環境条件に依存しうる。

また、持続評価関数 $\Pi$、可行集合 $K$、
可行行動対応 $A(x,e)$
（\S\ref{subsec:viability_function_and_blockage}）は、
現在状態と実現外部環境状態のもとで、
実効行動を介した一段階の状態遷移が
可行集合 $K$ に留まりうるかを記述する。
ここで「行動を介する」とは、
指向的介入だけを意味しない。
選択的待機を含む本論文の行動空間 $U$ の要素を通じて、
一段階の可行性維持可能性を記述するという意味である。

以上より、Phase~1 の ACS は、
内部と外部を区別しながら、
外部環境状態を引数として状態遷移し、
さらに必要環境条件に存続上依存しうる構造として記述される。
三原理は、この自己--環境関係を前提として導入される。

\subsubsection{三つの構造的問いと三原理}

自己--環境境界と必要環境依存を備えた ACS を
時間的に存続する構造として考えると、
少なくとも次の三つの問いが同時に意味をもつ。

\begin{description}[style=nextline]
  \item[問い1（構造的自己の可行性）]
  現在の構造的自己は、現在時刻において可行条件を満たしているか。
  形式的には、現在状態が可行集合 $K$ に属するか、
  すなわち $\Pi(x_t)\ge0$ であるかという問題に接続する。

  \item[問い2（臨界必要要因の充足）]
  ある必要環境要因が充足されなければ、
  どの行動によっても次状態を $K$ に置けない局面において、
  その要因が実際に充足されているか。
  このような要因を、以下では臨界必要要因と呼ぶ（定義~\ref{def:eu_fire}）。
  この局面は、現在の状態が $K$ の内にあるとき（維持）にも、外にあるとき（回復）にも生じうる。

  \item[問い3（必要な局面での可行な応答）]
  臨界必要要因がすべて充足されている局面において、
  次状態を $K$ に置く応答、すなわち可行性を維持する応答、
  または $K$ の外から一段階で回復する応答が存在するか。
  また、そのような応答の可能性が時間を通じて保持・獲得・回復・再構成されているか。
\end{description}

本論文は、これら三つの問いに対応する制約軸として、
それぞれ自己維持 SM、環境利用 EU、可行応答 FR を採用する。
各原理の時刻別充足述語は
定義~\ref{def:SM}、定義~\ref{def:EU}、定義~\ref{def:FR}
で与える。

ここで、三原理が「自己--環境関係から導出される」とは主張しない。
本論文が主張するのは、自己と環境を区別し、
その相互作用と必要環境依存のもとで時間的存続を扱うとき、
現在の可行性、臨界必要要因の充足、必要な局面での可行な応答という
三つの構造問題が同時に記述対象となり、
Phase~1 ではそれらを SM・EU・FR の三軸として形式化する、
という範囲である。

外部環境状態の差異に対して ACS が機能的に閉鎖し切っていないこと、
すなわち自己と環境の構造的な結合は、
三原理のいずれでもなく、EU の前提・背景構造として
\S\ref{subsec:env_coupling} で扱う。

また、問い1の数式 $\Pi(x_t)\ge0$ 自体は
$X$ と $\Pi$ が与えられれば形式的に記述できる。
したがって、本論文は「自己--環境境界がなければ
可行性という数式そのものが定義不能である」とは主張しない。
ここでの主張は、三つの問いを\emph{同一の ACS の存続問題として}
統合的に位置づける際に、自己--環境境界とその相互作用構造が
共通の前提となる、というものである。

\begin{remark}[三原理の構造的同時性]
\label{rem:principles_simultaneity}
本小節でいう同時性は、
SM、EU、FR が時間軸上の異なる時点で順次発生するという意味ではなく、
自己--環境関係を備えた ACS の存続を記述するとき、
三つの制約問題が論理的に並立して意味をもつという意味である。

したがって、
\[
  \text{三原理の構造的同時性}
  \;\neq\;
  \text{三原理の時刻別充足述語の常時の真}
\]
である。
ある時刻に一つまたは複数の時刻別充足述語が偽であることは、
本理論において排除されない。

また、一つまたは複数の時刻別充足述語が偽であることから
$\mathrm{Col}_{i,t}=1$
を一般的な帰結として置かない。
崩壊の定義は \S\ref{sec:acs_def}（定義~\ref{def:collapse}）にあり、
その周辺の状態と経路の整理は \S\ref{sec:collapse_cond} の責務である。
可行行動集合が空である状態
$A_t=\emptyset$
（閉塞、定義~\ref{def:blockage_sec2}）についても、
崩壊または整形式実現の不成立と同一視しない。

三原理の時刻別充足述語の間の関係は、
\S\ref{subsec:principle_linkage} で連動と含意関係の限定として述べる。
\end{remark}

\subsubsection{三原理の前提と統一記述形式}

三原理の前提は、原理（構造体への制約）ではなく、
対象とする ACS のクラスと世界の記述である。
前提を規約・構成要件・仮定のいずれに属するかとともに示す。

\begin{description}[style=nextline]
  \item[SM の前提]
  (P-SM1) 構造的自己には可行集合 $K=\{x\mid\Pi(x)\ge0\}$ がある
  （定義~\ref{def:Pi}、定義~\ref{def:K}）。
  (P-SM2) 状態は行動の選択によらずに変化しうる
  （仮定~\ref{ass:f_decomposition} の $f_{\mathrm{auto}}$）。
  \item[EU の前提]
  (P-EU1) ACS は外部環境の中にある
  （$E_{\mathrm{ext}}$、規約 TSC、$B_{SE}$、$f:X\times U\times E_{\mathrm{ext}}\to X$）。
  (P-EU2) 環境入力は退化していない
  （仮定~\ref{ass:env_input_nondegeneracy}、仮定~\ref{ass:env_viable_nondegeneracy}）。
  (P-EU3) 必要環境要因があり、それを欠けば存続できない
  （構成要件 CR-1、\S\ref{subsec:environment_space}）。
  自己と環境の構造的結合（\S\ref{subsec:env_coupling}）は EU の背景構造である。
  \item[FR の前提]
  (P-FR1) 将来の状態は、現在の行動と実現する外部環境状態の列に依存して定まる
  （定義~\ref{def:state_transition}）。
\end{description}

仮定~\ref{ass:nonstationary}（非定常環境）は、
FR の時刻別充足述語の前件ではなく、
Phase~1 の背景仮定として対象のクラスを限定する
（\S\ref{subsubsec:real_world_scope}）。
仮定~\ref{ass:theta_exterior_blockage} は前提ではなく、
\S\ref{subsec:sd_necessity} の条件付き解析において、可行性核の外への持続を閉塞へ接続する系にだけ用いる追加仮定である。

三原理を、次の統一的な形式で記述する。
「原理応答項目」は、各原理について ACS の側で何が成り立っていることが問われるかを述べる項目であり、
それを担う手段は書かない。
時間射程の欄は説明のためのものであり、形式上の分類ではない。

\begin{center}
\begingroup\small\setlength{\tabcolsep}{4pt}\renewcommand{\arraystretch}{1.18}
\begin{tabular}{p{2.6cm} p{3.9cm} p{3.9cm} p{3.9cm}}
\hline
\textbf{項目} & \textbf{SM} & \textbf{EU} & \textbf{FR} \\
\hline
対象 & 構造的自己の状態 & 自己と必要環境要因の関係 & 必要な局面での応答 \\[3pt]
適用域 & $t\in\mathbb{T}_{\mathrm{P},i}$ & $t\in\mathbb{T}_{\mathrm{P},i}$ & $t\in\mathbb{T}_{\mathrm{P},i}$ かつ $(x_t,e_t)\in\mathcal{D}_{\mathrm{FR}}$（$\mathrm{EU}_t$ が真の時刻） \\[3pt]
問われる局面 & すべての判定時刻 & すべての判定時刻（$\mathrm{Crit}^{(1)}(x_t)=\emptyset$ の時刻には空虚に真） & 適用域のすべての時刻（$K$ の内では維持、外では一段階の回復） \\[3pt]
時刻別充足述語 & $x_t\in K$（定義~\ref{def:SM}） & $\forall\eta\in\mathrm{Crit}^{(1)}(x_t):\eta(x_t,e_t)=1$（定義~\ref{def:EU}） & $A_t\neq\emptyset$（定義~\ref{def:FR}） \\[3pt]
原理応答項目 & 現在の状態が $K$ 内にあり、将来の可行性を損なわない状態が形成・保持されていること & 一段階の可行性に関して臨界となる必要環境要因について、その充足に至る自己と環境の関係が成り立っていること & $\mathcal{D}_{\mathrm{FR}}$ の局面で $K$ 側へつながる応答が存在し、その必要応答ポテンシャルが時間を通じて保持・獲得・回復・再構成されていること \\[3pt]
時間射程（説明） & 即時：現在の状態。拡張：$\mathrm{Hor}_f$、T-STAB (T1) & 即時：一段階の臨界必要要因の充足。拡張：持続的な充足（CR-1 の側で扱う） & 即時：局所 FR。拡張：\S\ref{subsec:sd_necessity}、T-STAB (T2) \\[3pt]
閉塞との関係 & 偽は閉塞の定義条件でない & 偽ならば閉塞（命題~\ref{prop:eu_blockage}） & 偽と閉塞は一致（$\mathcal{D}_{\mathrm{FR}}$ 上）。閉塞は EU の偽か、EU の真かつ FR の偽に排他的に分かれる（命題~\ref{prop:blockage_dichotomy}） \\
\hline
\end{tabular}
\endgroup
\end{center}

\subsubsection{三原理制約と内発的評価機能}

三原理は制約軸であるが、
それらに関係する差異が ACS の作動と完全に切断されてよいわけではない。
評価に関する構成要件としては、
\S\ref{sec:acs_def} の構成要件 CR-2 の (CR-2b)(iii) が、
評価の対象の違いを非等価に扱う場合が存在すること（非退化性）と、
$K\neq\emptyset$ かつ $X\setminus K\neq\emptyset$ のときに
$K$ に属する状態と属さない状態の違いに対して評価が異なる場合が存在すること
（可行性感応性）を要請する
（定義~\ref{def:three_primitives}、\S\ref{subsec:three_primitive_functions}）。
本節は構成要件の内容を変更しない。

本節はこれに加えて、三原理制約を自己--環境相互作用のなかで維持する
ACS の機能構造を記述するため、
可行性の低下、必要環境要因の非充足、可行な応答の縮退、
外部との関係の悪化等に関係する差異の少なくとも一部が、
ACS の応答形成に利用可能であることを、
構造的評価機能の機能要件として記述する。
評価対象に三原理に関係する状態差、関係差または時間変化が存在する場合、
その差異の少なくとも一部に対して評価結果または内部応答が変化しうることを、
評価機能の非退化性の最低線とする。
この機能要件は、三原理の時刻別充足述語の条件でも、構成要件でも、仮定でもなく、
三原理に関係する内発的評価機能の記述上の要件である。
すべての相異なる状態を識別すること、
形式値と内部評価が一致すること、
または一般的な誤差上界をこの要件から要求しない。
定量的精度条件が必要な場合には、該当する後続節で追加仮定として明示する。

本論文は、この役割を\emph{構造的評価機能}として扱う。
構造的評価機能は、外部の観測者が ACS に対して任意に設定する
合否判定規約ではなく、
自己--環境相互作用から得られる利用可能な情報に基づき、
当該 ACS が応答形成へ利用しうる評価機能として位置づけられる。

ただし、この内発性は、ACS が形式モデル上の真値を
完全に認識・計算することを意味しない。
たとえば、後続小節で用いられる $\Pi$、$K$、$A(x,e)$、
$E_{\mathrm{req}}(x)$、$\mathrm{Crit}^{(1)}(x)$、$\mu_{\mathrm{eff}}$ 等について、
ACS がそれらの真値または集合全体を完全に内部表象することは要求しない。
また、形式モデルから導かれる可行性核のような解析対象は、
ACS の内発的評価機能そのものとは区別する。
後続の外部関係に伴う可行性関連差異の評価および三原理の時間的統合評価でも、
評価機能自体は ACS 側に内発的に必要な機能として扱う一方、
具体的な尺度、閾値、精度、符号化および時間窓は
特定の実装として固定しない。

\begin{remark}[例示の地位]
\label{rem:example_status}
本節および後続節では、構造的評価機能の直感的理解を助けるために、
人間、他の生物または人工的な認知システムについて、
当該機能との構造的対応を説明する直感的な例を
少数、補足として述べる場合がある。

これらの例は、当該機能の存在証明、
ACS 一般に共通する具体実装、
または因果機構の実証ではない。
経験的事実として述べる場合には文献根拠を付し、
文献根拠を伴わない例は、Phase~1 の形式的帰結ではなく
理論理解を助ける解釈的補足として扱う。
\end{remark}

\subsubsection{後続小節への接続}

以上により、本節では、
自己--環境境界、状態遷移、必要環境依存および可行性構造を前提として、
SM、EU、FR が同時に意味をもつ三つの制約軸として導入される。

以下、\S\ref{subsec:sm_foundational} で SM を、
\S\ref{subsec:viab_core} で $E_{\mathrm{req}}$ ロバスト可行性核と有限段頑健可行集合を、
\S\ref{subsec:env_coupling} で構造的環境結合とその補助記述・診断量を、
\S\ref{subsec:eu} で EU を、
\S\ref{subsec:fr} で FR を扱う。
続いて \S\ref{subsec:principle_failure_relation} で三原理と閉塞・整形式実現・崩壊の関係を、
\S\ref{subsec:principle_linkage} で三原理の連動と含意関係の限定を、
\S\ref{subsec:sd_necessity} で応答多様性縮退と長期可行性の条件付き接続を整理する。
さらに \S\ref{subsec:s3_tstab} で三原理の時間的統合評価を、
\S\ref{subsec:principles_to_tipa} で構造的機能と標準機能形式への接続を扱う。
各小節では、主張を定義的等値性、追加仮定のもとでの条件付き含意、
または構造的接続として区別し、
主張の強さが証明および仮定の強さを超えないように扱う。

%----------------------------------------------------------------------
\subsection{SM}
\label{subsec:sm_foundational}

本小節では、SM の時刻別充足述語と、その判定に用いる可行性余裕度 $\Pi$ の構造を与える。
SM は、時刻 $t$ の状態 $x_t$ が \S\ref{sec:acs_def} で定義された可行集合 $K$ に属するかを問う
制約軸であり、新たな目的関数、最適化規則、
または他の二原理に先行する時間的生成条件ではない。
はじめに、$\Pi$ の一般化勾配構造と、$\Pi$ の最小値が捨てる成分別の診断量を整える。

\subsubsection{$\Pi$ の Clarke 一般化勾配構造}

定義~\ref{def:Pi} の $\Pi$ は
式 $\min_k(x_k-\theta_k)$ により
$\mathbb{R}^n$ 上でも同一に定義される。
Clarke 一般化勾配は開集合上の局所 Lipschitz 関数に対して取り扱うのが標準であるため、
まず $\mathbb{R}^n$ 上への拡張
\[
  \widetilde{\Pi}:\mathbb{R}^n\to\mathbb{R},\qquad \widetilde{\Pi}(x):=\min_k\bigl(x_k-\theta_k\bigr)
\]
を定義し、
ACS の状態空間 $X$ 上ではその制限 $\Pi:=\widetilde{\Pi}|_X$ を用いる。
Clarke 一般化勾配は $\widetilde{\Pi}$ に対して取り、
その値を $\partial^C\widetilde{\Pi}(x)$ と書く。
$x\in X$ に対しては、以下すべてにわたり略記として
\[
  \partial\Pi(x)\;:=\;\partial^C\widetilde{\Pi}(x)
\]
を用いる。
この拡張により、$X$ の境界点 $x\in\partial X$ においても
一般化勾配が $\mathbb{R}^n$ の開近傍上の局所 Lipschitz 関数に対するものとして矛盾なく定義される。
以下では、この $\Pi$ について、以降で用いる基本的性質を命題として確立する。

\begin{proposition}[$\Pi$ の Lipschitz 連続性と Clarke 一般化勾配構造]
\label{prop:pi_subdiff}
定義~\ref{def:Pi} の持続評価関数
\[
  \Pi=\widetilde{\Pi}|_X,
  \qquad
  \widetilde{\Pi}(x)
  =
  \min_k(x_k-\theta_k)
\]
のもとで、以下が成立する。

\textbf{(i) Lipschitz 連続性：}$\Pi$ は $X$ 上 Lipschitz 定数 $L=1$ で Lipschitz 連続である：
\[
  |\Pi(x)-\Pi(x')| \;\le\; \|x-x'\|_\infty \qquad \forall\, x,x'\in X.
\]

\textbf{(ii) Clarke 一般化勾配の明示形式：}任意の $x_t\in X$ において、
Clarke 一般化勾配 $\partial\Pi(x_t)=\partial^C\widetilde{\Pi}(x_t)$ は
非空コンパクト凸集合をなし、
\[
  \partial\Pi(x_t)\;=\;\operatorname{conv}\bigl\{\,e_k \;\bigm|\; k\in\arg\min_j\bigl(x_{j,t}-\theta_j\bigr)\,\bigr\}
\]
と表される。
ここで $e_k\in\mathbb{R}^n$ は第 $k$ 成分が $1$、
他が $0$ の標準基底ベクトルである。
\end{proposition}

\begin{proof}
\textbf{(i)} 各アフィン関数 $\ell_k(x)=x_k-\theta_k$ は
Lipschitz 定数 $1$ で Lipschitz 連続である。
有限個の Lipschitz 連続関数の下限（pointwise minimum）が
同一の Lipschitz 定数で Lipschitz 連続であることは直接確認できる標準的事実であり、
変分解析の一般的背景については
Rockafellar and Wets~\cite{RockafellarWets1998} を参照する。

\textbf{(ii)} $\widetilde{\Pi}$ は
有限個の $C^1$ アフィン関数の pointwise minimum であるから、
Clarke の一般化勾配の標準的性質\cite{Clarke1990}
（一般化勾配の和・最大／最小に関する計算規則）により、
Clarke 一般化勾配は最小値を達成するインデックスに対応する勾配の凸包と一致する：
\[
  \partial\Pi(x_t)=\operatorname{conv}\bigl\{\nabla\ell_k(x_t)\mid k\in\arg\min_j\ell_j(x_t)\bigr\}.
\]
$\nabla\ell_k(x)=e_k$ より所期の形を得る。
有限個のベクトルの凸包は $\mathbb{R}^n$ の非空コンパクト凸集合をなす。
\end{proof}

\begin{remark}[$L_\infty$ ノルムの採用]
\label{rem:linf}
命題~\ref{prop:pi_subdiff}(i) で採用する $L_\infty$ ノルムは、
定義~\ref{def:Pi} の $\Pi(x_t)=\min_k(x_{k,t}-\theta_k)$ が
各成分の最小値として定義されることと構造的に整合する。
各成分の最大誤差が可行性余裕度 $\Pi$ の変動上限を直接規定する。
\end{remark}

\subsubsection{閾値割れ集合と不足量}

\begin{definition}[閾値割れ集合と不足量ベクトル（分類C）]
\label{def:violation_profile}
状態 $x_t$ に対して、\emph{閾値割れ集合} $V^-(x_t)$ を
\[
  V^-(x_t)\;:=\;\bigl\{\,k\in\{1,\ldots,n\}\;\bigm|\;x_{k,t}-\theta_k<0\,\bigr\}
\]
と定義する。
また、\emph{不足量ベクトル} $d^-(x_t)\in\mathbb{R}^n$ を
\[
  d^-(x_t)\;:=\;\bigl((\theta_1-x_{1,t})_+,\ldots,(\theta_n-x_{n,t})_+\bigr),\qquad (a)_+:=\max(a,0)
\]
と定義する。

$V^-(x_t)$ は閾値を下回る成分の添字集合、
$d^-(x_t)$ は各成分の閾値不足量を表し、
\[
  V^-(x_t)
  =
  \bigl\{
    k
    \mid
    d^-_k(x_t)>0
  \bigr\}
\]
が成り立つ。
\end{definition}

\begin{proposition}[$\Pi$ と閾値割れ集合の関係]
\label{prop:pi_violation_profile}
定義~\ref{def:Pi} の $\Pi(x_t)=\min_k(x_{k,t}-\theta_k)$ および
定義~\ref{def:violation_profile} のもとで、
\[
  \Pi(x_t)\ge0\;\iff\;V^-(x_t)=\emptyset\;\iff\;d^-(x_t)=\mathbf{0},
\]
かつ
\[
  \Pi(x_t)<0\;\iff\;V^-(x_t)\neq\emptyset.
\]
\end{proposition}

\begin{proof}
まず第一の同値連鎖を示す。
\[
  \Pi(x_t)\ge0
  \iff
  \min_k(x_{k,t}-\theta_k)\ge0
  \iff
  \forall k,\;
  x_{k,t}-\theta_k\ge0
\]
である。
最右辺は、定義~\ref{def:violation_profile} より
$V^-(x_t)=\emptyset$ と同値であり、
また $\forall k,\; (\theta_k-x_{k,t})_+=0$、
すなわち $d^-(x_t)=\mathbf{0}$ とも同値である。
よって
\[
  \Pi(x_t)\ge0
  \iff
  V^-(x_t)=\emptyset
  \iff
  d^-(x_t)=\mathbf{0}
\]
を得る。

第二の同値は第一式の否定として、
\[
  \Pi(x_t)<0
  \iff
  \exists k,\;
  x_{k,t}-\theta_k<0
  \iff
  V^-(x_t)\neq\emptyset
\]
により得られる。
\end{proof}


\subsubsection{SM の時刻別充足述語と可行集合}

\begin{definition}[SM の時刻別充足述語]
\label{def:SM}
$t\in\mathbb{T}_{\mathrm{P},i}$ について、
SM の時刻別充足述語 $\mathrm{SM}_t$ は
\[
  x_t\in K
\]
のとき真、$x_t\notin K$ のとき偽とする。
\end{definition}

\begin{proposition}[SM と可行集合の同値性]
\label{prop:sm_K_equiv}
\begin{enumerate}[label=(\roman*)]
  \item 任意の $t\in\mathbb{N}_0$ について
  $\Pi(x_t)\ge0 \iff x_t\in K$ である。
  \item $t\in\mathbb{T}_{\mathrm{P},i}$ ならば、
  $\mathrm{SM}_t$ が真であること、$\Pi(x_t)\ge0$、$x_t\in K$ の三条件は同値である。
\end{enumerate}
\end{proposition}

\begin{proof}
(i) は $K$ の定義（定義~\ref{def:K}）と $x_t\in X$ による。
(ii) は定義~\ref{def:SM} と (i) による。
\end{proof}

\begin{proposition}[一段階の $K$ 保持と $A_t$]
\label{prop:sm_one_step_A}
\begin{enumerate}[label=(\roman*)]
  \item 任意の $t\in\mathbb{N}_0$ と $u\in U$ について
  \[
    u\in A_t
    \iff
    f(x_t,u,e_t)\in K
    \iff
    \Pi\bigl(f(x_t,u,e_t)\bigr)\ge0 .
  \]
  \item 任意の $t\in\mathbb{N}_0$ について、
  $x_{t+1}\in K \iff u_t\in A_t$ である。
  \item $t+1\in\mathbb{T}_{\mathrm{P},i}$ ならば、
  $\mathrm{SM}_{t+1}$ が真であることと $u_t\in A_t$ は同値である。
\end{enumerate}
\end{proposition}

\begin{proof}
(i) の前半は可行行動対応の定義（定義~\ref{def:at}）、後半は $K$ の定義による。
(ii) は状態遷移 $x_{t+1}=f(x_t,u_t,e_t)$（定義~\ref{def:state_transition}）と (i) による。
(iii) は (ii) と定義~\ref{def:SM} による。
\end{proof}

\subsubsection{SM、閉塞、崩壊、必要環境充足の区別}

\begin{remark}[SM の基底的位置]
\label{rem:sm_foundational_status}
命題~\ref{prop:sm_K_equiv} が示す意味で、
SM は三原理のうち、
現在状態が可行集合 $K$ に属するか否かを直接記述する制約軸である。

ここでいう「基底的」とは、
可行集合 $K$ が、現在状態の判定だけでなく、
可行行動集合 $A_t$、
EU の一段階臨界域 $\mathrm{Crit}^{(1)}_\eta$、
FR の定義域 $\mathcal{D}_{\mathrm{FR}}$、
可行性保持指示関数 $\varphi_u$ および応答多様度 $\mu_{\mathrm{eff},t}$ の
基準集合としても参照されるという、形式体系上の位置を指す。
$\mathrm{Crit}^{(1)}_\eta$ と $\mathcal{D}_{\mathrm{FR}}$ は $K$ を $A$ を通して参照するが、
$\mathrm{EU}_t$ と $\mathrm{FR}_t$ の定義域は $x_t\in K$ を要求しない（定義~\ref{def:EU}、定義~\ref{def:FR}）。

したがって本論文は、$\mathrm{SM}_t$ の真から
$\mathrm{EU}_t$ または $\mathrm{FR}_t$ の真を一般的な帰結として置かず、
SM が他の二原理を論理的に包含することも主張しない
（備考~\ref{rem:principle_implication_scope}）。
\end{remark}

\begin{remark}[現在の SM と次時刻の $K$ 保持は別条件である]
\label{rem:sm_current_next}
命題~\ref{prop:sm_one_step_A} は、
現在時刻の状態条件と一段階先の $K$ 保持の可能性を区別する。
現在時刻の SM は $x_t\in K$ という状態条件である。
これに対して $A_t\neq\emptyset$ は、
実現外部環境状態 $e_t$ のもとで、
一段階先を $K$ に置く行動が存在するという遷移の条件である。

したがって本論文は、定義だけから
$x_t\in K \iff A_t\neq\emptyset$
を仮定しない。
$A_t=\emptyset$ となる時刻の扱いは、$x_t\in K$ の場合も $x_t\notin K$ の場合も、
\S\ref{subsec:principle_failure_relation} の閉塞の二分
（命題~\ref{prop:blockage_dichotomy}）で整理する。
\end{remark}

\begin{remark}[SM と閉塞の論理的区別]
\label{rem:sm_blockage_distinction}
$x_t\in K$ と閉塞 $A_t=\emptyset$ は同一の条件ではない。
SM は現在状態 $x_t$ の $K$ への所属によって定まり、
$A_t$ は $f(x_t,u,e_t)\in K$ を満たす行動 $u$ の存在によって定まる。
前者は時刻 $t$ の状態条件、
後者は同じ時刻の状態と実現外部環境状態から
一段階先の可行状態へ至る行動があるかを表す遷移の条件である。

Phase~1 の一般定義は、
\[
  A_t=\emptyset
  \Longrightarrow
  x_t\notin K
\]
も、
\[
  x_t\notin K
  \Longrightarrow
  A_t=\emptyset
\]
も仮定しない。
個別の ACS で一方の含意が成り立つか否かは、
$f$、$U$、$e_t$ および $K$ の具体的構造に依存する。
\end{remark}

\begin{remark}[SM の時刻別充足述語が偽であることと崩壊の区別]
\label{rem:sm_not_collapse}
$t\in\mathbb{T}_{\mathrm{P},i}$ で $\mathrm{SM}_t$ が偽であること、
すなわち $\Pi(x_t)<0$（$x_t\notin K$）は、
時刻 $t$ の状態が可行集合の外にあることを表す。
これは \S\ref{sec:acs_def} の崩壊 $\mathrm{Col}_{i,t}=1$（定義~\ref{def:collapse}）と同一ではない。
本論文は
\[
  x_t\notin K
  \Longrightarrow
  \mathrm{Col}_{i,t}=1
\]
を一般には仮定しない。
また、$\mathrm{SM}_t$ が判定される時刻では
補題~\ref{lem:principle_time_domain_down} により $\mathrm{Col}_{i,t}=0$ である。

逆に、崩壊の判定条件を SM の定義へ組み込むこともしない。
SM は単一時刻の状態条件であり、
崩壊は同一個体としての非回復性を含む \S\ref{sec:acs_def} の概念である（定義~\ref{def:collapse}）。
整形式実現の不成立 $\mathrm{WF}_{i,t}=0$ についても、
本論文は $\mathrm{SM}_t$ の偽だけからこれを一般的な帰結として置かない
（備考~\ref{rem:principles_failure_table}）。
\end{remark}

\begin{remark}[SM と必要環境充足の区別]
\label{rem:sm_vs_Ereq}
必要環境充足集合
\[
  E_{\mathrm{req}}(x_t)
  =
  \bigl\{
    e\in E_{\mathrm{ext}}
    \bigm|
    \forall\eta\in E,\;
    \eta(x_t,e)=1
  \bigr\}
\]
は、状態 $x_t$ に相対して
すべての必要環境要因を充足する外部環境状態を表す。
これに対し SM は、現在状態 $x_t$ 自体が $K$ に属するかを問う。
したがって、$e_t\in E_{\mathrm{req}}(x_t)$ と $x_t\in K$ は別の条件であり、
本論文は両者の同値性を仮定しない。

これは、必要環境要因の充足が構造的自己の存続に無関係であることを意味しない。
構成要件 CR-1 は、構造的自己の存続が外部環境条件と関係することを要求する
（\S\ref{subsec:environment_space}）。
ただし、その関係は
「必要環境要因が現在充足されているならば現在状態は $K$ に属する」
という一時刻の同値では表されない。
内部損傷、内部資源状態、過去からの履歴その他の状態差は $x_t$ に含まれうるため、
必要環境要因が現在充足されていても $x_t\notin K$ である場合を、
一般定義から排除しない。
また、必要環境要因の一時的な非充足から、
直ちに $x_t\notin K$ または崩壊を導くことも仮定しない。
必要環境要因のうち、時刻 $t$ の臨界必要要因の充足は EU の時刻別充足述語として扱う
（\S\ref{subsec:eu}）。
\end{remark}

\begin{remark}[必要環境依存性と SM の構造的接続]
\label{rem:sm_environment_dependency}
SM は内部状態だけを引数とする $\Pi:X\to\mathbb{R}$ によって判定されるが、
そのことは SM が外部環境から因果的・構造的に独立であることを意味しない。
状態遷移 $f:X\times U\times E_{\mathrm{ext}}\to X$ は外部環境状態を引数に含み、
必要環境空間 $E$ と構成要件 CR-1 は
構造的自己の存続に必要な外部条件を記述する。
したがって、外部環境状態の差異は状態遷移を介して
将来の $x_{t+1}$、ひいては $\Pi(x_{t+1})$ に関係しうる。

ただし、$f$ の型だけから外部環境引数への非退化な依存を導くことはしない。
どの外部環境の差が将来の $K$ 保持に関与するかは、
当該 ACS の $f$、$E$、$E_{\mathrm{req}}$ および状態構造に相対して決まる。
外部環境状態の差異が環境応答成分の次状態に反映されうるという構造は、
\S\ref{subsec:env_coupling} の構造的環境結合条件として、
三原理とは別に記述する。
\end{remark}

\begin{remark}[単一時刻条件と時間的維持の区別]
\label{rem:sm_temporal_scope}
定義~\ref{def:SM} の $\mathrm{SM}_t$ は、
各時刻 $t\in\mathbb{T}_{\mathrm{P},i}$ ごとに判定される単一時刻の述語である。
本論文は、ある時刻 $t$ で $\mathrm{SM}_t$ が真であることだけから、
後続の $s>t$ について $x_s\in K$ が成り立つことを一般的な帰結として置かない。

有限時間区間 $\{t,t+1,\ldots,t+T\}$（$T\in\mathbb{N}_0$）にわたり
状態が $K$ に留まるという主張は、
\[
  x_{t+\tau}\in K
  \qquad
  \forall\,\tau\in\{0,1,\ldots,T\}
\]
という時刻列上の条件を別に要求する。
将来へ向かう頑健性の段数は
\S\ref{subsec:viab_core} の $\mathrm{Hor}_f$ で、
有限窓内の正の可行性余裕度は
\S\ref{subsec:s3_tstab} の (T1) で、それぞれ記述する。
これらはいずれも $\mathrm{SM}_t$ の真理条件ではない。
無期限の可行性保存を論じる場合には、
状態遷移、可行行動、必要環境条件、
および可行性核等の動学的構造が必要となる。
本論文は、単一時刻の $\mathrm{SM}_t$ の真から長期可行性を一般的な帰結として置かない。
\end{remark}

\subsubsection{本小節の位置づけ}

以上より、Phase~1 における SM の時刻別充足述語は、
$t\in\mathbb{T}_{\mathrm{P},i}$ の時刻について
\[
  \mathrm{SM}_t\ \text{が真}
  \iff
  \Pi(x_t)\ge0
  \iff
  x_t\in K
\]
という現在状態の可行性条件として位置づけられる。
この条件は、一段階の可行行動の存在、
必要環境要因の現在充足、閉塞、崩壊、
または将来にわたる可行性保存と同一ではない。

次小節では、$E_{\mathrm{req}}$ ロバスト可行性核と有限段頑健可行集合を導入し、
SM の原理応答項目のうち将来へ向かう側面を、
分類Cの量として記述する。

%----------------------------------------------------------------------
\subsection{$E_{\mathrm{req}}$ ロバスト可行性核と有限段頑健可行集合}
\label{subsec:viab_core}

本小節では、必要環境充足集合 $E_{\mathrm{req}}(x)$ のすべての外部環境状態に対して
一つの共通の行動で可行性を保つことができる状態の集合を導入し、
その有限段の版と段数を分類Cの量として定める。
これらは SM の原理応答項目のうち将来へ向かう側面を記述するものであり、
いずれの時刻別充足述語の真理条件でもない。
後続の \S\ref{subsec:fr} と \S\ref{subsec:sd_necessity} は本小節の結果を用いる。

\subsubsection{必要環境充足状態域と $E_{\mathrm{req}}$ ロバスト可行性核}

\begin{definition}[必要環境充足状態域]
\label{def:req_domain}
必要環境要因を充足する外部環境状態が少なくとも一つ存在する
状態の集合を
\[
  D_{\mathrm{req}}
  :=
  \bigl\{\,x\in X \bigm| E_{\mathrm{req}}(x)\neq\emptyset \,\bigr\}
  \subseteq X
\]
と書き、その可行集合との共通部分を
\[
  K_{\mathrm{req}}
  :=
  K\cap D_{\mathrm{req}}
\]
と書く。

$D_{\mathrm{req}}$ および $K_{\mathrm{req}}$ は、
必要環境空間 $E$ と可行集合 $K$ から
\S\ref{sec:acs_def} の形式体系において一意に定まる派生集合であり、
解析者が選択する対象ではない。
$E_{\mathrm{req}}(x)\neq\emptyset$ ならば、ある $e\in E_{\mathrm{ext}}$ がすべての $\eta\in E$ について $\eta(x,e)=1$ を満たすので、
\S\ref{sec:acs_def} の要因別充足状態域（定義~\ref{def:eta_satisfiable_domain}）とのあいだに
$D_{\mathrm{req}}\subseteq\bigcap_{\eta\in E}D_\eta$ が成り立つ。

\textit{注：}
\S\ref{sec:acs_def} は
$\forall x\in X:E_{\mathrm{req}}(x)\neq\emptyset$ を要求していない。
構成要件 CR-1 の (CR-1b) が要求するのは
$E_{\mathrm{req}}(x)\neq\emptyset$ なる $x$ の存在であり、
全域的な非空性ではない。
したがって $D_{\mathrm{req}}\subsetneq X$ および
$K_{\mathrm{req}}\subsetneq K$ である場合を排除しない。
\end{definition}

\begin{definition}[$E_{\mathrm{req}}$ ロバスト可行性核]
\label{def:viab_core}
状態遷移 $x_{t+1}=f(x_t,u_t,e_t)$
（定義~\ref{def:state_transition}）のもとで、
部分集合 $K'\subseteq K_{\mathrm{req}}$ が
\emph{$E_{\mathrm{req}}$ 不変}であるとは、
\[
  \forall x\in K',\;
  \exists u\in U
  \;\text{ s.t. }\;
  \forall e\in E_{\mathrm{req}}(x),\;
  f(x,u,e)\in K'
  \tag{$\ast$}
\]
が成立することをいう。

$E_{\mathrm{req}}$ 不変な $K_{\mathrm{req}}$ の部分集合全体の族を
$\mathfrak{K}$ と書き、
\emph{$E_{\mathrm{req}}$ ロバスト可行性核}を
\[
  \mathrm{Viab}_f(K)
  :=
  \bigcup_{K'\in\mathfrak{K}} K'
\]
と定義する。
この和集合が $\mathfrak{K}$ の最大元であることは
命題~\ref{prop:viab_core_max} で示す。
定義より
\[
  \mathrm{Viab}_f(K)
  \subseteq
  K_{\mathrm{req}}
  \subseteq
  K
\]
である。

\textit{注（量化順序）：}
条件 $(\ast)$ は $\exists u\,\forall e$ の順序であり、
\emph{一つの共通行動}が $E_{\mathrm{req}}(x)$ の
すべての外部環境状態に対して可行性を保持することを要求する。
これは
\[
  \forall e\in E_{\mathrm{req}}(x),\;
  \exists u\in U
  \;\text{ s.t. }\;
  f(x,u,e)\in K'
\]
という、外部環境状態ごとに異なる行動を選んでよい条件とは
異なる。本小節ではこの量化順序を入れ替えない。

\textit{注（定義域の限定）：}
$\mathfrak{K}$ の要素を $K_{\mathrm{req}}$ の部分集合に限るのは、
$E_{\mathrm{req}}(x)=\emptyset$ である状態において
$(\ast)$ の全称量化が空虚に成立し、
そのような状態が無条件に核へ含まれることを避けるためである。
$x\notin K_{\mathrm{req}}$ の状態は、本診断の対象外として扱う。
本論文は、この扱いから当該状態が将来失敗することを一般的な帰結として置かない。

\textit{注（標準 Viability Kernel との関係）：}
本定義は $E_{\mathrm{req}}(x)$ 上のロバスト条件を採用した
本論文固有の構成であり、
標準的な Viability Kernel との数学的同一性を主張しない。
\end{definition}

\begin{proposition}[最大 $E_{\mathrm{req}}$ 不変集合の存在]
\label{prop:viab_core_max}
$\mathrm{Viab}_f(K)$ は $E_{\mathrm{req}}$ 不変であり、
したがって $\mathfrak{K}$ の最大元である。
\end{proposition}

\begin{proof}
$\mathfrak{K}\ni\emptyset$ であるから
$\mathfrak{K}\neq\emptyset$ であり、和集合は定義される
（$\mathfrak{K}=\{\emptyset\}$ のときは
$\mathrm{Viab}_f(K)=\emptyset$ となる）。

$x\in\mathrm{Viab}_f(K)$ を任意にとる。
和集合の定義より、ある $K'\in\mathfrak{K}$ が存在して
$x\in K'$ である。
$K'$ の $E_{\mathrm{req}}$ 不変性より、ある $u\in U$ が存在して
\[
  \forall e\in E_{\mathrm{req}}(x),\;
  f(x,u,e)\in K'.
\]
$K'\subseteq\mathrm{Viab}_f(K)$ であるから、
同じ $u$ について
\[
  \forall e\in E_{\mathrm{req}}(x),\;
  f(x,u,e)\in\mathrm{Viab}_f(K).
\]
$x$ は任意であったから $\mathrm{Viab}_f(K)$ は
$E_{\mathrm{req}}$ 不変である。
また $\mathrm{Viab}_f(K)\subseteq K_{\mathrm{req}}$ であるから
$\mathrm{Viab}_f(K)\in\mathfrak{K}$ であり、
和集合の構成より任意の $K'\in\mathfrak{K}$ について
$K'\subseteq\mathrm{Viab}_f(K)$ である。
\end{proof}

\begin{remark}[Viability Theory との構造的関係]
\label{rem:viability_relation}
本論文の $E_{\mathrm{req}}$ に相対する
ロバスト不変集合構成は、
Viability Theory~\cite{Aubin1991,SaintPierre1994} および
集合値解析~\cite{AubinFrankowska1990} における
可行性・不変性の考え方と構造的に関係する。

本論文は状態空間のコンパクト性
（仮定~\ref{ass:compact_spaces}）および
状態遷移写像の連続性
（仮定~\ref{ass:f_continuity}）を用いる一方、
定義~\ref{def:viab_core} では、
状態依存の必要環境充足集合
$E_{\mathrm{req}}(x)$ を量化域として、
\[
  \exists u\in U\;
  \forall e\in E_{\mathrm{req}}(x)
\]
という本論文固有のロバスト条件を採用する。

したがって、
本論文の $\mathrm{Viab}_f(K)$ と
標準的な Viability Kernel との数学的同一性、
または特定の Viability Theory の
標準公理体系への一致を主張しない。
ここでいう関係は、
可行集合内に状態を保持しうる制御・不変性を扱うという
概念的・構造的対応に限られる。

また、
$E_{\mathrm{req}}:X\rightrightarrows E_{\mathrm{ext}}$ について、
閉グラフ性、上半連続性またはコンパクト値性を
\S\ref{sec:acs_def} は一般には仮定していない
（定義~\ref{def:acs}(3) は $\eta$ に正則性を課さない）。
したがって本論文は、
Viability Theory の一般論のみから
$\mathrm{Viab}_f(K)$ が $X$ の閉部分集合であることを主張しない。
閉性を要する主張は本小節では用いない。
\end{remark}

\subsubsection{可行性核の理論上の地位}

\begin{remark}[$\mathrm{Viab}_f(K)$ の存在論的・認識論的地位]
\label{rem:viab_core_status}
$\mathrm{Viab}_f(K)$ は、ACS の内部に実装される認知表象ではなく、
ACS が計算する対象でもなく、
ACS が直接の目標とする集合でもない。

本論文における $\mathrm{Viab}_f(K)$ は、
状態遷移写像 $f$、可行集合 $K$、行動空間 $U$、および
必要環境充足集合 $E_{\mathrm{req}}$ を固定したときに、
どの状態から $E_{\mathrm{req}}$ ロバストな可行性保持が
可能であるかを数学的に特徴づける、
モデル由来の派生的解析対象である。

したがって本論文では、
「ACS が可行性核を目指す」、
「ACS が可行性核へ戻ろうとする」、
「ACS が可行性核を認識している」
といった目的論的・認知論的表現を用いない。

\textit{注（モデル相対性）：}
$\mathrm{Viab}_f(K)$ は解析者が恣意的に選ぶ集合ではない。
モデルが固定されれば、それは当該モデルから数学的に定まる。
ただし、現実のACSを記述するにあたって
どの $f$、どの $K$、どの必要環境要因を採用するかには
モデル化が介在する。
したがって本論文は、現実そのものに対して観測者非依存の
既知な可行性核が与えられているとは扱わず、
$\mathrm{Viab}_f(K)$ を
指定された形式モデルに相対的に定まる解析対象として扱う。
\end{remark}

\begin{remark}[形式モデル上の核・内部予測・実現軌道の区別]
\label{rem:three_layers}
本論文では次の三者を区別し、同一視しない。

第一に、形式モデル上の可行性核 $\mathrm{Viab}_f(K)$。
これは解析者が形式体系から数学的に導く集合である。

第二に、ACS 内部の予測および評価。
これは内部モデル $M$、内部環境推定 $\hat e_t$
（定義~\ref{def:internal_env_estimate}）、および
\S\ref{sec:tipa} の機構を通じて ACS が扱う、
利用可能な内部情報である。
本論文は、それが形式モデルと一致することを仮定しない。

第三に、実際に実現した軌道 $x_0,x_1,x_2,\ldots$ である。

本論文は、現実のACSについて、
真の状態遷移法則の把握、将来の外部環境状態の把握、
必要環境要因の特定、可行性核の計算、
または将来軌道の予測が可能であることを仮定しない。
Phase~1 において $\mathrm{Viab}_f(K)$ を定義できることと、
実在ACSがその集合を認識または計算できることとは別問題である。

したがって、ACS が $\mathrm{Viab}_f(K)$ を知らない場合でも、
結果として長期にわたりその内部に留まることはありうる。
逆に、ACS 自身の内部評価が良好であっても、
内部モデルが不完全であれば
形式モデル上の可行領域から外れることはありうる。

現実のACSについて本論文が述べるのは、
局所的な評価・応答・更新の結果として
可行性核内に留まりうる、という水準までである。

\textit{注：}
モデル不確実性、部分観測、推定誤差、および
将来の外部環境状態に関する不確実性により、
形式モデル上の可行性核と ACS 内部の予測可能領域は
一致しない場合がある。
確率モデル、部分観測モデル、推定誤差の分布、信念状態、
および確率的可行性核の形式化は Phase~1 では導入せず、
Phase~2 の課題とする（\S\ref{subsec:limitations_list} の L1、within present framework）。
\end{remark}

\begin{remark}[可行性核外部性の意味の限定]
\label{rem:viab_core_exterior}
本論文は、ある状態が $\mathrm{Viab}_f(K)$ の外部にあることから、
その状態から将来必ず失敗することを一般的な帰結として置かない。

$x\notin\mathrm{Viab}_f(K)$ から従うのは、
定義~\ref{def:viab_core} の量化条件のもとで、
$E_{\mathrm{req}}$ ロバストな可行性保持を与える
共通方策が存在しないことに限られる。

実現する外部環境状態の列が
$E_{\mathrm{req}}(x_s)$ 全体を走査するとは限らず、
また ACS が $\mathrm{Update}$ によって
内部モデルおよび応答方略を変化させることもありうる。
したがって、実現軌道が可行性を保つことは排除されない。

本論文では「核外部にあること」を
「将来の破綻が確定していること」と記述しない。

\textit{注（時刻間変化との関係）：}
本小節の $\mathrm{Viab}_f(K)$ は、
$f$、$K$、$U$、$E_{\mathrm{req}}$ を固定した
Phase~1 の形式構造に対する解析対象である。
これに対し \S\ref{sec:extended_dynamics} では、
内部モデル・記憶・応答方略が $\mathrm{Update}$ によって
時刻間で変化する過程を扱う。
したがって本論文は、
一度可行性核の外部へ出たACSが回復しないとは述べない。
実効的な応答能力の時刻間変化と可行性構造との関係は
\S\ref{sec:extended_dynamics} の備考~\ref{rem:s5_learning_viability} で述べる。
\end{remark}

\subsubsection{有限段頑健可行集合と頑健存続可能段数}

\begin{definition}[有限段 $E_{\mathrm{req}}$ 頑健可行集合（分類C）]
\label{def:viab_stage}
$K_{\mathrm{req}}=K\cap D_{\mathrm{req}}$（定義~\ref{def:req_domain}）とする。
$\nu\in\mathbb{N}_0$ について、集合 $\mathrm{Viab}_f^{(\nu)}(K)\subseteq K_{\mathrm{req}}$ を
\[
  \mathrm{Viab}_f^{(0)}(K):=K_{\mathrm{req}},
\]
\[
  \mathrm{Viab}_f^{(\nu)}(K)
  :=
  \bigl\{
    x\in K_{\mathrm{req}}
    \bigm|
    \exists u\in U\;
    \forall e\in E_{\mathrm{req}}(x):\;
    f(x,u,e)\in\mathrm{Viab}_f^{(\nu-1)}(K)
  \bigr\}
  \qquad(\nu\ge1)
\]
により帰納的に定め、
\emph{$\nu$ 段 $E_{\mathrm{req}}$ 頑健可行集合}と呼ぶ。
上付きの $(\nu)$ は反復の段数を表す。
量化の順序は定義~\ref{def:viab_core} と同じく $\exists u\,\forall e$ である。
各段の元を $K_{\mathrm{req}}$ に限るので、
$E_{\mathrm{req}}(x)=\emptyset$ の状態で全称量化が空虚に成立する場合は含まれない。
\end{definition}

\begin{definition}[頑健存続可能段数（分類C）]
\label{def:viab_horizon}
$x\in K_{\mathrm{req}}$ について、
\[
  \mathrm{Hor}_f(x)
  :=
  \sup
  \bigl\{
    \nu\in\mathbb{N}_0
    \bigm|
    x\in\mathrm{Viab}_f^{(\nu)}(K)
  \bigr\}
  \in\mathbb{N}_0\cup\{\infty\}
\]
を\emph{頑健存続可能段数}と呼ぶ。
$x\in\mathrm{Viab}_f^{(0)}(K)$ であるから、上の集合は空でない。
\end{definition}

\begin{proposition}[有限段頑健可行集合の入れ子性]
\label{prop:viab_stage_nested}
\begin{enumerate}[label=(\roman*)]
  \item すべての $\nu\in\mathbb{N}_0$ について
  $\mathrm{Viab}_f^{(\nu+1)}(K)\subseteq\mathrm{Viab}_f^{(\nu)}(K)$ である。
  \item $x\in K_{\mathrm{req}}$ と $\nu\in\mathbb{N}_0$ について、
  $x\in\mathrm{Viab}_f^{(\nu)}(K)\iff\nu\le\mathrm{Hor}_f(x)$ である。
  $\mathrm{Hor}_f(x)<\infty$ のとき、
  定義~\ref{def:viab_horizon} の上限は最大値として達成される。
\end{enumerate}
\end{proposition}

\begin{proof}
(i)：$\nu$ についての帰納法による。
$\nu=0$ のとき、定義から $\mathrm{Viab}_f^{(1)}(K)\subseteq K_{\mathrm{req}}=\mathrm{Viab}_f^{(0)}(K)$ である。
$\nu\ge1$ とし、$\mathrm{Viab}_f^{(\nu)}(K)\subseteq\mathrm{Viab}_f^{(\nu-1)}(K)$ を仮定する。
$x\in\mathrm{Viab}_f^{(\nu+1)}(K)$ をとると、$x\in K_{\mathrm{req}}$ であり、
ある $u\in U$ が存在して、すべての $e\in E_{\mathrm{req}}(x)$ について
$f(x,u,e)\in\mathrm{Viab}_f^{(\nu)}(K)\subseteq\mathrm{Viab}_f^{(\nu-1)}(K)$ である。
よって $x\in\mathrm{Viab}_f^{(\nu)}(K)$ である。

(ii)：$S_{\mathrm{Hor}}(x):=\{\nu\in\mathbb{N}_0\mid x\in\mathrm{Viab}_f^{(\nu)}(K)\}$ とおく。
(i) を繰り返し用いると、$\nu\in S_{\mathrm{Hor}}(x)$ かつ $\mu\le\nu$ ならば $\mu\in S_{\mathrm{Hor}}(x)$ である。
$\nu\in S_{\mathrm{Hor}}(x)$ ならば上限の定義から $\nu\le\mathrm{Hor}_f(x)$ である。
逆に $\nu\le\mathrm{Hor}_f(x)$ とする。
$\mathrm{Hor}_f(x)<\infty$ のとき、$S_{\mathrm{Hor}}(x)$ は上に有界な $\mathbb{N}_0$ の空でない部分集合なので
最大元をもち、それは $\mathrm{Hor}_f(x)$ に等しい。
下方閉性から $\nu\in S_{\mathrm{Hor}}(x)$ である。
$\mathrm{Hor}_f(x)=\infty$ のとき、$S_{\mathrm{Hor}}(x)$ は上に有界でないので
$\nu\le\nu'$ となる $\nu'\in S_{\mathrm{Hor}}(x)$ が存在し、下方閉性から $\nu\in S_{\mathrm{Hor}}(x)$ である。
\end{proof}

\begin{proposition}[$E_{\mathrm{req}}$ ロバスト可行性核と有限段頑健可行集合]
\label{prop:viab_stage_core}
すべての $\nu\in\mathbb{N}_0$ について
$\mathrm{Viab}_f(K)\subseteq\mathrm{Viab}_f^{(\nu)}(K)$ である。
したがって、$x\in\mathrm{Viab}_f(K)$ ならば $\mathrm{Hor}_f(x)=\infty$ である。
\end{proposition}

\begin{proof}
$\nu$ についての帰納法による。
$\nu=0$ のとき、定義~\ref{def:viab_core} から $\mathrm{Viab}_f(K)\subseteq K_{\mathrm{req}}$ である。
$\nu\ge1$ とし、$\mathrm{Viab}_f(K)\subseteq\mathrm{Viab}_f^{(\nu-1)}(K)$ を仮定する。
$x\in\mathrm{Viab}_f(K)$ をとる。
命題~\ref{prop:viab_core_max} により $\mathrm{Viab}_f(K)$ は $E_{\mathrm{req}}$ 不変であるから、
ある $u\in U$ が存在して、すべての $e\in E_{\mathrm{req}}(x)$ について
$f(x,u,e)\in\mathrm{Viab}_f(K)\subseteq\mathrm{Viab}_f^{(\nu-1)}(K)$ である。
また $x\in K_{\mathrm{req}}$ である。
よって $x\in\mathrm{Viab}_f^{(\nu)}(K)$ である。
後半は定義~\ref{def:viab_horizon} による。
\end{proof}

\begin{proposition}[有限状態空間における可行性核と有限段頑健可行集合]
\label{prop:viab_stage_finite}
内部状態空間 $X$ が有限集合であるとする。
このとき、$\nu^{\ast}\le|K_{\mathrm{req}}|$ なる $\nu^{\ast}\in\mathbb{N}_0$ が存在して、
すべての $\nu\ge\nu^{\ast}$ について
$\mathrm{Viab}_f^{(\nu)}(K)=\mathrm{Viab}_f^{(\nu^{\ast})}(K)$ であり、
\[
  \mathrm{Viab}_f(K)
  =
  \mathrm{Viab}_f^{(\nu^{\ast})}(K)
  =
  \bigcap_{\nu\in\mathbb{N}_0}\mathrm{Viab}_f^{(\nu)}(K)
\]
である。
したがって、$x\in K_{\mathrm{req}}$ について
$\mathrm{Hor}_f(x)=\infty\iff x\in\mathrm{Viab}_f(K)$ である。
本命題は仮定を用いない。
\end{proposition}

\begin{proof}
各段の集合を $V^{(\nu)}:=\mathrm{Viab}_f^{(\nu)}(K)$ と書く。
$K_{\mathrm{req}}$ の部分集合 $V$ に対し
\[
  \Phi(V):=\bigl\{x\in K_{\mathrm{req}}\bigm|\exists u\in U\;\forall e\in E_{\mathrm{req}}(x):\;f(x,u,e)\in V\bigr\}
\]
とおくと、定義~\ref{def:viab_stage} より $V^{(0)}=K_{\mathrm{req}}$、$V^{(\nu)}=\Phi(V^{(\nu-1)})$（$\nu\ge1$）である。

第一に、段数の停止を示す。
命題~\ref{prop:viab_stage_nested} (i) より $V^{(\nu+1)}\subseteq V^{(\nu)}$ であり、
$|V^{(0)}|=|K_{\mathrm{req}}|$ は有限である。
$\nu=0,1,\ldots,|K_{\mathrm{req}}|$ のすべてで $V^{(\nu+1)}\subsetneq V^{(\nu)}$ であったとすると、
$|V^{(|K_{\mathrm{req}}|+1)}|\le|K_{\mathrm{req}}|-(|K_{\mathrm{req}}|+1)<0$ となり矛盾する。
よって $V^{(\nu^{\ast}+1)}=V^{(\nu^{\ast})}$ なる $\nu^{\ast}\le|K_{\mathrm{req}}|$ が存在する。
$V^{(\nu+1)}=V^{(\nu)}$ ならば $V^{(\nu+2)}=\Phi(V^{(\nu+1)})=\Phi(V^{(\nu)})=V^{(\nu+1)}$ であるから、
帰納法により、すべての $\nu\ge\nu^{\ast}$ について $V^{(\nu)}=V^{(\nu^{\ast})}$ である。
列は減少するので、$\bigcap_{\nu}V^{(\nu)}=V^{(\nu^{\ast})}$ である。

第二に、$V^{\ast}:=V^{(\nu^{\ast})}$ が $E_{\mathrm{req}}$ 不変であることを示す。
$x\in V^{\ast}=V^{(\nu^{\ast}+1)}=\Phi(V^{\ast})$ ならば、
$x\in K_{\mathrm{req}}$ であり、ある $u\in U$ が存在して、
すべての $e\in E_{\mathrm{req}}(x)$ について $f(x,u,e)\in V^{\ast}$ である。
よって $V^{\ast}\subseteq K_{\mathrm{req}}$ は定義~\ref{def:viab_core} の意味で $E_{\mathrm{req}}$ 不変であり、
$V^{\ast}\subseteq\mathrm{Viab}_f(K)$ である。
逆の包含 $\mathrm{Viab}_f(K)\subseteq\bigcap_{\nu}V^{(\nu)}=V^{\ast}$ は命題~\ref{prop:viab_stage_core} による。

第三に、$x\in K_{\mathrm{req}}$ について、
命題~\ref{prop:viab_stage_nested} (ii) より
$\mathrm{Hor}_f(x)=\infty$ はすべての $\nu$ で $x\in V^{(\nu)}$ であることと同値であり、
これは $x\in\bigcap_{\nu}V^{(\nu)}=\mathrm{Viab}_f(K)$ と同値である。
\end{proof}

\begin{remark}[頑健存続可能段数の地位と限定]
\label{rem:viab_stage_scope}
\begin{enumerate}
  \item 本論文は $\mathrm{Viab}_f(K)=\bigcap_{\nu\in\mathbb{N}_0}\mathrm{Viab}_f^{(\nu)}(K)$ を主張しない。
  右辺の元は各有限段で頑健であるが、本論文は、右辺が $E_{\mathrm{req}}$ 不変な集合であることを示していない。
  同じ理由で、本論文は $\mathrm{Hor}_f(x)=\infty\Rightarrow x\in\mathrm{Viab}_f(K)$ を主張しない。
  ただし $X$ が有限集合の場合は、命題~\ref{prop:viab_stage_finite} により両者が成立する。
  \item $\mathrm{Hor}_f(x)$ は $E_{\mathrm{req}}(x)$ のすべての外部環境状態に対する頑健性の段数であり、
  実現する外部環境状態列のもとで存続する時刻数の予測ではない。
  \item $\mathrm{Hor}_f$ は $\mathrm{SM}_t$ の真理条件ではない。
  SM の原理応答項目のうち将来へ向かう側面を記述する分類Cの量である。
  \item $\mathrm{Viab}_f(K)=\emptyset$ であっても、ある $\nu$ について
  $\mathrm{Viab}_f^{(\nu)}(K)\neq\emptyset$ となる系を記述できる。
  例えば $X=\{0,1,2\}\subseteq\mathbb R$、$\Pi(x)=x-1$（したがって $K=\{1,2\}$）、$U=\{u\}$、
  $E_{\mathrm{ext}}=\{e\}$、$E=\{\eta\}$、$\eta\equiv1$、$f(2,u,e)=1$、$f(1,u,e)=f(0,u,e)=0$ とすると、
  $\mathrm{Viab}_f^{(1)}(K)=\{2\}$、$\mathrm{Viab}_f^{(2)}(K)=\emptyset$ であり、
  命題~\ref{prop:viab_stage_finite} により $\mathrm{Viab}_f(K)=\emptyset$ である（型の水準の例）。
  本論文は、どの系でも有限段の集合が空でないとは主張しない。
  \item ACS がこれらの量を内部に保持または計算することを要求しない
  （\S\ref{subsec:description_levels_global}）。
\end{enumerate}
\end{remark}

\subsubsection{全環境に対する頑健な核と楽観的な核との関係}

定義~\ref{def:viab_core} の可行性核は、外部環境状態の量化の範囲を各状態の $E_{\mathrm{req}}(x)$ に取る。
本小節では、比較のために、量化の範囲を外部環境状態空間 $E_{\mathrm{ext}}$ 全体に取る二つの核を置き、
可行性核との包含関係を示す。

\begin{definition}[全環境頑健核と楽観核（分類C）]
\label{def:kernel_all_any}
部分集合 $K'\subseteq K$ が\emph{全環境頑健不変}であるとは、
\[
  \forall x\in K',\;\exists u\in U\;\text{ s.t. }\;\forall e\in E_{\mathrm{ext}},\;f(x,u,e)\in K'
\]
が成立することをいい、\emph{楽観不変}であるとは、
\[
  \forall x\in K',\;\exists u\in U,\;\exists e\in E_{\mathrm{ext}}\;\text{ s.t. }\;f(x,u,e)\in K'
\]
が成立することをいう。
全環境頑健不変な $K$ の部分集合全体の和集合を\emph{全環境頑健核} $\mathrm{Viab}_f^{\mathrm{all}}(K)$、
楽観不変な $K$ の部分集合全体の和集合を\emph{楽観核} $\mathrm{Viab}_f^{\mathrm{any}}(K)$ と書く。
上付きの $\mathrm{all}$・$\mathrm{any}$ は量化の種類を表すラベルであり、
有限段の段数を表す括弧付きの上付き $(\nu)$ とは異なる。

\textit{注：}
二つの核は必要環境空間 $E$ と $E_{\mathrm{req}}$ を用いず、$K$ の部分集合として定める。
したがって $\mathrm{Viab}_f(K)$ と異なり、$K_{\mathrm{req}}$ に含まれることを要件としない。
いずれも比較のための集合であり、三原理の時刻別充足述語の真理条件ではない。
\end{definition}

\begin{proposition}[可行性核と全環境頑健核・楽観核の包含関係]
\label{prop:kernel_standard_comparison}
\begin{enumerate}[label=(\roman*)]
  \item $\mathrm{Viab}_f^{\mathrm{all}}(K)$ は全環境頑健不変であり、$\mathrm{Viab}_f^{\mathrm{any}}(K)$ は楽観不変である。
  \item $\mathrm{Viab}_f(K)\subseteq\mathrm{Viab}_f^{\mathrm{any}}(K)$ である。
  \item $\mathrm{Viab}_f^{\mathrm{all}}(K)\subseteq\mathrm{Viab}_f(K)$ であることと、
  $\mathrm{Viab}_f^{\mathrm{all}}(K)\subseteq D_{\mathrm{req}}$ であることは同値である。
\end{enumerate}
本命題は仮定を用いない。
\end{proposition}

\begin{proof}
(i)：$x\in\mathrm{Viab}_f^{\mathrm{all}}(K)$ ならば、ある全環境頑健不変な $K'\subseteq K$ について $x\in K'$ であり、
ある $u\in U$ が存在して、すべての $e\in E_{\mathrm{ext}}$ について
$f(x,u,e)\in K'\subseteq\mathrm{Viab}_f^{\mathrm{all}}(K)$ である。
楽観核についても、存在量化を保ったまま同じ議論が成り立つ。

(ii)：$x\in\mathrm{Viab}_f(K)$ とする。
$\mathrm{Viab}_f(K)\subseteq K_{\mathrm{req}}$ より $E_{\mathrm{req}}(x)\neq\emptyset$ である。
命題~\ref{prop:viab_core_max} により、ある $u\in U$ が存在して、
すべての $e\in E_{\mathrm{req}}(x)$ について $f(x,u,e)\in\mathrm{Viab}_f(K)$ である。
$E_{\mathrm{req}}(x)\subseteq E_{\mathrm{ext}}$ の元 $e$ を一つとれば、
$f(x,u,e)\in\mathrm{Viab}_f(K)$ となる $u$ と $e$ が存在する。
よって $\mathrm{Viab}_f(K)\subseteq K$ は楽観不変であり、$\mathrm{Viab}_f(K)\subseteq\mathrm{Viab}_f^{\mathrm{any}}(K)$ である。

(iii)：$R:=\mathrm{Viab}_f^{\mathrm{all}}(K)$ とおく。
$R\subseteq\mathrm{Viab}_f(K)$ ならば、$\mathrm{Viab}_f(K)\subseteq K_{\mathrm{req}}\subseteq D_{\mathrm{req}}$ より $R\subseteq D_{\mathrm{req}}$ である。
逆に $R\subseteq D_{\mathrm{req}}$ とする。
$R\subseteq K$ であるから $R\subseteq K\cap D_{\mathrm{req}}=K_{\mathrm{req}}$ である。
(i) により、各 $x\in R$ についてある $u\in U$ が存在して、
すべての $e\in E_{\mathrm{ext}}$ について $f(x,u,e)\in R$ である。
$E_{\mathrm{req}}(x)\subseteq E_{\mathrm{ext}}$ であるから、
すべての $e\in E_{\mathrm{req}}(x)$ についても $f(x,u,e)\in R$ である。
よって $R$ は $K_{\mathrm{req}}$ の $E_{\mathrm{req}}$ 不変な部分集合であり、
定義~\ref{def:viab_core} により $R\subseteq\mathrm{Viab}_f(K)$ である。
\end{proof}

\begin{remark}[包含関係の読み方と限定]
\label{rem:kernel_standard_comparison_scope}
\begin{enumerate}
  \item 命題~\ref{prop:kernel_standard_comparison} (iii) は、すべての外部環境状態に対して頑健に保てる状態であっても、
  必要環境要因を充足する外部環境状態が存在しない状態を含むならば、
  可行性核に含まれないことを表す。
  可行性核は、$E_{\mathrm{ext}}$ 全体への頑健性ではなく、
  必要環境要因が充足される範囲での頑健性を表す集合である。
  \item 本論文は、(ii) および (iii) の包含が真部分集合となるか等号となるかを一般的な帰結として置かない。
  それは $f$、$K$、$E$ の具体的構造に依存する。
  \item 本命題は、差別化核その他の既存の核との定義の同一性を主張しない
  （\S\ref{subsec:related}）。
  \item ACS がこれらの集合を内部に保持または計算することを要求しない
  （\S\ref{subsec:description_levels_global}）。
\end{enumerate}
\end{remark}

$E_{\mathrm{req}}$ ロバスト可行性核と比較用の集合との包含関係を図~\ref{fig:s3_kernel_inclusion} に示す。

\begin{figure}[tbp]
\centering
\begin{tikzpicture}[
  font=\footnotesize,
  set/.style={draw, rounded corners=2pt, align=center, inner sep=3pt},
  cond/.style={font=\scriptsize, inner sep=1pt},
  note/.style={font=\scriptsize, align=left, inner sep=1pt},
  plabel/.style={font=\small, anchor=west, inner sep=0pt},
  sub/.style={-{Latex[length=1.6mm]}},
  csub/.style={-{Latex[length=1.6mm]}, dashed},
]
\node[plabel] at (0,1.6) {(a)\ 追加の条件なしに成り立つ包含};
\node[set] (VF)  at (0.9,0)  {$\mathrm{Viab}_f(K)$};
\node[set] (Vn1) at (3.6,0)  {$\mathrm{Viab}_f^{(\nu+1)}(K)$};
\node[set] (Vn)  at (6.0,0)  {$\mathrm{Viab}_f^{(\nu)}(K)$};
\node[set] (V0)  at (9.0,0)  {$\mathrm{Viab}_f^{(0)}(K)$\\$=K_{\mathrm{req}}$};
\node[set] (K)   at (11.6,0) {$K$};
\node[set] (Dr)  at (9.0,-1.4) {$D_{\mathrm{req}}$};
\node[set] (Vany) at (3.6,-2.3) {$\mathrm{Viab}_f^{\mathrm{any}}(K)$};
\node[set] (Vall) at (11.6,1.3) {$\mathrm{Viab}_f^{\mathrm{all}}(K)$};
\draw[sub] (VF) -- node[cond, above]{$\forall\nu$} (Vn1);
\draw[sub] (Vn1) -- (Vn);
\draw[sub] (Vn) -- node[cond, above=3pt]{入れ子} (V0);
\draw[sub] (V0) -- (K);
\draw[sub] (V0) -- (Dr);
\draw[sub] (VF) |- (Vany);
\draw[sub] (Vany) -| (K);
\draw[sub] (Vall) -- (K);
\node[note, anchor=north west] at (5.0,-2.55)
 {$\nu\in\mathbb{N}_0$ は任意。$\mathrm{Viab}_f(K)$ は各段に含まれる};

\begin{scope}[yshift=-5.0cm]
\node[plabel] at (0,1.0) {(b)\ 条件のもとでの関係};
\node[set] (Vall2) at (1.4,0) {$\mathrm{Viab}_f^{\mathrm{all}}(K)$};
\node[set] (VF2)   at (10.6,0) {$\mathrm{Viab}_f(K)$};
\draw[csub] (Vall2) -- node[cond, above]{$\mathrm{Viab}_f^{\mathrm{all}}(K)\subseteq D_{\mathrm{req}}$ のとき、かつそのときに限り成り立つ}
 node[cond, below]{命題~\ref{prop:kernel_standard_comparison} (iii)} (VF2);
\node[set, dashed, anchor=north west, align=left] at (0,-1.0)
 {$X$ が有限集合のとき（命題~\ref{prop:viab_stage_finite}）：
  ある $\nu^{\ast}\le|K_{\mathrm{req}}|$ について\\
  $\mathrm{Viab}_f(K)=\mathrm{Viab}_f^{(\nu^{\ast})}(K)=\bigcap_{\nu\in\mathbb{N}_0}\mathrm{Viab}_f^{(\nu)}(K)$};

\begin{scope}[shift={(0,-2.9)}]
\draw[sub] (0,0) -- (0.8,0); \node[anchor=west, font=\scriptsize] at (0.9,0) {$A\to B$：$A\subseteq B$};
\draw[csub] (6.3,0) -- (7.1,0); \node[anchor=west, font=\scriptsize] at (7.2,0) {条件のもとで成り立つ包含};
\end{scope}
\end{scope}
\end{tikzpicture}
\caption{$E_{\mathrm{req}}$ ロバスト可行性核と、有限段頑健可行集合・全環境頑健核・楽観核の包含関係。
矢印の向きは部分集合の向きを表す。
線のない組については、本図は包含関係を表示しない。
本図は、包含が真部分集合であるか等号であるかを表さない（備考~\ref{rem:kernel_standard_comparison_scope}）。
(a) の根拠は定義~\ref{def:req_domain}・定義~\ref{def:viab_core}・定義~\ref{def:viab_stage}・定義~\ref{def:kernel_all_any}・
命題~\ref{prop:viab_stage_nested}・命題~\ref{prop:viab_stage_core}・命題~\ref{prop:kernel_standard_comparison} (ii) である。
(b) の等号は、$X$ が有限集合の場合に命題~\ref{prop:viab_stage_finite} により成り立つ。
一般の場合については本論文は主張しない（備考~\ref{rem:viab_stage_scope}）。
$\mathrm{Viab}_f(K)$ が空で、ある段の $\mathrm{Viab}_f^{(\nu)}(K)$ が空でない系も記述できる。}
\label{fig:s3_kernel_inclusion}
\end{figure}

\subsubsection{行動応答成分の応答を考慮した頑健性}

定義~\ref{def:viab_core} の可行性核は、各状態 $x$ で $E_{\mathrm{req}}(x)$ のすべての元にわたって量化する。
\S\ref{sec:acs_def} の分解 $E_{\mathrm{ext}}=E^{\mathrm{bg}}\times O$（定義~\ref{def:Eext_decomposition}）のもとでは、
この量化は行動応答成分 $o$ についても、同じ時刻にありえた別の値をすべて含む。
本小節では、これとは別に、現在の行動応答成分を固定し、
その次の値を遷移写像 $g_O$（定義~\ref{def:env_realization}）で定める頑健性を置く。

\begin{definition}[応答込み頑健可行性核（分類C）]
\label{def:viab_resp}
$(x,o)\in X\times O$ について、
\[
  B_{\mathrm{req}}(x,o)
  :=
  \bigl\{\,b\in E^{\mathrm{bg}}\bigm|(b,o)\in E_{\mathrm{req}}(x)\,\bigr\}
\]
を\emph{背景充足集合}と呼び、
\[
  K_{\mathrm{req}}^{\mathrm{resp}}
  :=
  \bigl\{\,(x,o)\in K\times O\bigm| B_{\mathrm{req}}(x,o)\neq\emptyset\,\bigr\}
\]
と置く。
部分集合 $K'\subseteq K_{\mathrm{req}}^{\mathrm{resp}}$ が\emph{応答込みで $E_{\mathrm{req}}$ 不変}であるとは、
\[
  \forall (x,o)\in K',\;
  \exists u\in U
  \;\text{ s.t. }\;
  \forall b\in B_{\mathrm{req}}(x,o),\;
  \bigl(f(x,u,(b,o)),\ g_O(o,x,u,b)\bigr)\in K'
\]
が成立することをいう。
そのような部分集合全体の族を $\mathfrak{K}^{\mathrm{resp}}$ と書き、
\[
  \mathrm{Viab}^{\mathrm{resp}}_f(K)
  :=
  \bigcup_{K'\in\mathfrak{K}^{\mathrm{resp}}}K'
  \subseteq X\times O
\]
を\emph{応答込み頑健可行性核}と呼ぶ。
$B_{\mathrm{req}}(x,o)\neq\emptyset$ を要求するのは、
定義~\ref{def:viab_core} で $K_{\mathrm{req}}$ に限ったのと同じく、
全称量化が空虚に成立する組を含めないためである。
\end{definition}

\begin{proposition}[応答込み頑健可行性核の不変性]
\label{prop:viab_resp_max}
$\mathrm{Viab}^{\mathrm{resp}}_f(K)$ は応答込みで $E_{\mathrm{req}}$ 不変であり、
したがって $\mathfrak{K}^{\mathrm{resp}}$ の最大元である。
\end{proposition}

\begin{proof}
$\emptyset\in\mathfrak{K}^{\mathrm{resp}}$ であるから和集合は定義される。
$(x,o)\in\mathrm{Viab}^{\mathrm{resp}}_f(K)$ をとれば、ある $K'\in\mathfrak{K}^{\mathrm{resp}}$ について $(x,o)\in K'$ であり、
ある $u\in U$ が存在して、すべての $b\in B_{\mathrm{req}}(x,o)$ について次の組が $K'$ に属する。
$K'\subseteq\mathrm{Viab}^{\mathrm{resp}}_f(K)$ であるから、同じ $u$ について次の組は $\mathrm{Viab}^{\mathrm{resp}}_f(K)$ に属する。
また $\mathrm{Viab}^{\mathrm{resp}}_f(K)\subseteq K_{\mathrm{req}}^{\mathrm{resp}}$ である。
\end{proof}

二つの核を比べるため、$\mathrm{Viab}^{\mathrm{resp}}_f(K)$ の $X$ への射影を
\[
  \pi_X\bigl(\mathrm{Viab}^{\mathrm{resp}}_f(K)\bigr)
  :=
  \bigl\{\,x\in X\bigm|\exists o\in O:\ (x,o)\in\mathrm{Viab}^{\mathrm{resp}}_f(K)\,\bigr\}
\]
と書く。

\begin{proposition}[二つの頑健性の非包含]
\label{prop:robust_noninclusion}
本小節と定義~\ref{def:viab_core} の定義が適用できる型の条件をみたす
$(X,U,E^{\mathrm{bg}},O,f,K,E,g_O)$ について、
$\mathrm{Viab}_f(K)$ と $\pi_X\bigl(\mathrm{Viab}^{\mathrm{resp}}_f(K)\bigr)$ の間には、
どちら向きの包含も一般には成り立たない。
すなわち、次の (i)(ii) をそれぞれみたす例が存在する。
\begin{enumerate}[label=(\roman*)]
  \item $\mathrm{Viab}_f(K)\not\subseteq\pi_X\bigl(\mathrm{Viab}^{\mathrm{resp}}_f(K)\bigr)$。
  \item $\pi_X\bigl(\mathrm{Viab}^{\mathrm{resp}}_f(K)\bigr)\not\subseteq\mathrm{Viab}_f(K)$。
\end{enumerate}
\end{proposition}

\begin{proof}
いずれの例でも $X=\{0,1\}\subseteq\mathbb R$、$\Pi(x)=x-1$（したがって $K=\{1\}$）、
$U=\{u_{\emptyset}\}$ とし、枠は一つ（$\mathcal A_O=\{a\}$）、$E=\{\eta\}$ とする。
有限集合には離散位相を置く。

(i)：$E^{\mathrm{bg}}=\{b\}$、$O=O_a=\{\bot_a,\mathrm{food}\}$ とし、
$\eta(x,(b,o))=1\iff o=\mathrm{food}$ と置く。
$f(1,u_{\emptyset},(b,\mathrm{food}))=1$、それ以外の引数では $f=0$ とし、
$g_O\equiv\bot_a$ とする。
このとき $E_{\mathrm{req}}(x)=\{(b,\mathrm{food})\}$ であり、$K_{\mathrm{req}}=\{1\}$ である。
$\{1\}$ は $E_{\mathrm{req}}$ 不変であるから $1\in\mathrm{Viab}_f(K)$ である。
一方、$(1,\mathrm{food})$ の次の組は $(1,\bot_a)$ であり、$B_{\mathrm{req}}(1,\bot_a)=\emptyset$ なので
$(1,\bot_a)\notin K_{\mathrm{req}}^{\mathrm{resp}}$ である。
したがって $(1,\mathrm{food})$ を含む応答込み不変集合は存在せず、
$(1,\bot_a)$ と $(0,o)$ はそもそも $K_{\mathrm{req}}^{\mathrm{resp}}$ に属さないので、
$\mathrm{Viab}^{\mathrm{resp}}_f(K)=\emptyset$ である。

(ii)：$E^{\mathrm{bg}}=\{b_1,b_2\}$、$O=O_a=\{\bot_a,\mathrm{good},\mathrm{bad}\}$ とし、
$\eta(x,(b,o))=1\iff b=b_1$ と置く。
$f(1,u_{\emptyset},(b_1,\mathrm{good}))=1$、それ以外の引数では $f=0$ とし、
$g_O\equiv\mathrm{good}$ とする。
このとき $(b_1,\mathrm{bad})\in E_{\mathrm{req}}(1)$ に対して $f=0\notin K$ であるから、
どの $K'\subseteq K_{\mathrm{req}}$ も $1$ を含んで $E_{\mathrm{req}}$ 不変にはならず、
$\mathrm{Viab}_f(K)=\emptyset$ である。
一方、$B_{\mathrm{req}}(1,\mathrm{good})=\{b_1\}$ であり、$(1,\mathrm{good})$ の次の組はつねに $(1,\mathrm{good})$ であるから、
$\{(1,\mathrm{good})\}$ は応答込みで $E_{\mathrm{req}}$ 不変であり、
$1\in\pi_X\bigl(\mathrm{Viab}^{\mathrm{resp}}_f(K)\bigr)$ である。
\end{proof}

\begin{remark}[二つの例の読みと射程]
\label{rem:robust_noninclusion_reading}
(i) は、取り込みによって行動応答成分が失われる場合である。
既存の核は次の時刻の成分の値を固定せずに量化するので、資源の枯渇を反映しない。
(ii) は、行動応答成分が当該ACSのもとで保たれる場合である。
既存の核は、実現しない成分の値まで量化に含める。
二つの例は、本小節の定義が適用できる型の水準で構成したものであり、
\S\ref{sec:acs_def} の構成要件と仮定のすべてをみたす個体であることまでは、本命題の主張に含めない。
\end{remark}

\begin{remark}[同じ時刻にありえた外部環境状態にわたる量化]
\label{rem:same_time_alternatives}
本節で $E_{\mathrm{ext}}$ または $E_{\mathrm{req}}(x)$ にわたって量化する条件
（一段階臨界域 $\mathrm{Crit}^{(1)}_\eta$、$\mathcal D_{\mathrm{FR}}$、$\mu_{\mathrm{eff}}$、$\mathrm{EC}(x)$、
定義~\ref{def:viab_core}・\ref{def:viab_stage} の頑健性、
\S\ref{subsec:sd_necessity} の応答不能環境集合）は、
行動応答成分についても、同じ時刻にありえた別の値を含む。
これらは、履歴から当該の値に到達できるかを問わない、同じ時刻の外部環境状態の比較である。
真理条件は \S\ref{sec:acs_def} の分解の前後で変わらない。
応答込みの版を置くのは定義~\ref{def:viab_resp} だけとし、
応答不能環境集合、$\mu_{\mathrm{eff}}$ および \S\ref{subsec:sd_necessity} の条件付き解析の応答込みの版は、
Phase~1.$\alpha$ の課題とする（\S\ref{subsec:phase1a}）。
\end{remark}

%----------------------------------------------------------------------
\subsection{構造的環境結合とその補助記述・診断量}
\label{subsec:env_coupling}

本小節では、ACS が外部環境状態の差異に対して機能的に閉鎖し切っていないという
自己と環境の構造的な結合を、三原理から独立した条件として形式化し、
EU の前提・背景構造として位置づける。
あわせて、環境応答成分、一般化勾配の環境応答方向への感応性、
および実現外部相互作用を、補助的な記述・診断量として整理する。
ここでいう「補助」とは、三原理の時刻別充足述語の真理条件ではないことを意味する。

本小節では、次の三つを互いに同一視しない。
\begin{enumerate}
  \item 構造的環境結合：外部環境状態の差異が ACS の次状態に反映されうる構造があるか
  （定義~\ref{def:env_coupling}）。
  \item 実現した環境相互作用・環境感応性：
  境界横断的な外部作用が実現しているか（定義~\ref{def:phi}）、
  および可行性余裕度を定める最弱成分が環境応答座標に属するか
  （命題~\ref{prop:pi_E_sensitivity}）。
  \item EU：一段階の臨界必要要因がすべて充足されているか（定義~\ref{def:EU}）。
\end{enumerate}

\subsubsection{環境応答成分 $x_E$}

\begin{assumption}[環境応答成分の構造的決定]
\label{ass:env_response}
ACS $\Sigma$ に対して、状態空間 $X$ の座標系が所与である。
この座標系において、状態ベクトル $x_t$ の成分のうち、
外部環境状態との差異に対する応答に関係する成分インデックス集合
\[
  I_E\subseteq\{1,2,\ldots,n\},
  \qquad
  |I_E|=n_E,\quad
  1\le n_E\le n
\]
が、当該 ACS の形式記述に相対して固定される。

Phase~1 では $I_E$ を、
解析者が構造的環境結合条件の判定のため任意に選択する評価集合ではなく、
当該 ACS に固有の構造パラメータとして扱う。
ただし、この仮定は特定の神経回路、アルゴリズム、
センサ構成またはハードウェア実装を要求するものではない。

導入地位：S（構造的）。強さ：中。
依拠箇所：環境応答成分の定義（定義~\ref{def:x_E}）と命題~\ref{prop:pi_E_sensitivity}。
\end{assumption}

\begin{definition}[環境応答成分 $x_E$ と座標射影 $\pi_E$]
\label{def:x_E}
仮定~\ref{ass:env_response} のもとで、
$I_E=\{k_1<k_2<\cdots<k_{n_E}\}$ とし、
座標射影
\[
  \pi_E:\mathbb{R}^n\to\mathbb{R}^{n_E}
\]
を
\[
  \pi_E(z)
  :=
  \bigl(
    z_{k_1},
    z_{k_2},
    \ldots,
    z_{k_{n_E}}
  \bigr),
  \qquad
  z\in\mathbb{R}^n
\]
で一意に定める。

$\pi_E$ は $\mathbb{R}^n$ 上の線形写像であり、
状態ベクトル
$x_t\in X\subseteq\mathbb{R}^n$
に対する $\pi_E(x_t)$ と、
一般化勾配ベクトル
$g\in\partial\Pi(x_t)\subseteq\mathbb{R}^n$
に対する $\pi_E(g)$ は、
いずれも同一の座標射影 $\pi_E$ の適用である。
型の曖昧さを避けるため、
$\pi_E$ の定義域を状態空間 $X$ ではなく
$\mathbb{R}^n$ とする。

記号
\[
  x_E:=\pi_E(x_t)\in\mathbb{R}^{n_E}
\]
を\emph{環境応答成分}と呼ぶ。
$x_E$ は状態ベクトル $x_t$ の
$I_E$ 成分への制限を表す。
ここで $x_E$ は $x_t\in X$ の内部状態成分であり、
実現外部環境状態 $e_t\in E_{\mathrm{ext}}$、
外部環境状態空間 $E_{\mathrm{ext}}$ そのもの、
または必要環境空間 $E$ を表すものではない。

$\pi_E(X)$ は $X$ のコンパクト性
（仮定~\ref{ass:compact_spaces}）と
$\pi_E$ の連続性により、
$\mathbb{R}^{n_E}$ のコンパクト部分集合をなす。

さらに、$x_E$ 方向の一般化勾配射影集合を
\[
  \partial_{x_E}\Pi(x_t)
  :=
  \bigl\{
    \pi_E(g)
    \;\bigm|\;
    g\in\partial\Pi(x_t)
  \bigr\}
  \subseteq
  \mathbb{R}^{n_E}
\]
と定義する。

$\partial_{x_E}\Pi(x_t)$ は
$\partial\Pi(x_t)$ の線形像として
非空コンパクト凸集合をなす
（命題~\ref{prop:pi_subdiff} と
$\pi_E$ の線形性による）。

この記号は、
命題~\ref{prop:pi_E_sensitivity} において、
可行性余裕度を規定する最弱成分のうち
環境応答座標に属するものが存在するかを記述するために用いる。
条件
\[
  \partial_{x_E}\Pi(x_t)
  \neq
  \{\mathbf{0}_{n_E}\}
\]
は構造的環境結合に関連する環境感応性の補助条件であるが、
本小節で定義する構造的環境結合条件（定義~\ref{def:env_coupling}）そのものではない。
\end{definition}

\begin{remark}[環境応答成分の観察的対応]
\label{rem:env_response_observation}
現実の生物的・人工的認知システムにおいて、
異なる外部環境状態 $e,e'\in E_{\mathrm{ext}}$ に対して
内部状態成分が系統的に異なる応答を示す場合、
その成分は $I_E$ の観察的候補となりうる。
形式的には、ある
$x\in X$、$u\in U$、$e\neq e'\in E_{\mathrm{ext}}$ について
$[f(x,u,e)]_k\neq [f(x,u,e')]_k$
が確認されるとき、第 $k$ 成分はその候補となる。

ただし、本条件は $I_E$ の一般的な因果同定規則を定義するものではない。
観測ノイズ、時間遅延、内部フィードバック等を考慮した
操作的・因果的決定基準の形式化はPhase~2 の課題とする。
\end{remark}

\subsubsection{構造的環境結合条件}

\begin{definition}[構造的環境結合条件（分類C）]
\label{def:env_coupling}
状態 $x\in X$ が\emph{構造的環境結合条件}を満たすとは、
外部環境状態の差異が、少なくとも一つの行動のもとで
環境応答成分の次状態に非自明な差異を与えうること、すなわち
\[
  \exists\,u\in U,\;
  \exists\,e,e'\in E_{\mathrm{ext}}
  \quad\text{s.t.}\quad
  e\neq e',
  \qquad
  \pi_E\bigl(f(x,u,e)\bigr)
  \neq
  \pi_E\bigl(f(x,u,e')\bigr)
\]
が成り立つことをいい、これを $\mathrm{EC}(x)$ と書く。
時刻 $t$ については $\mathrm{EC}(x_t)$ を用いる。

本条件は、自己--環境境界をもつ ACS が
外部環境状態の差異に対して機能的に閉じていないことを表す
Phase~1 の最小構造条件であり、
三原理の時刻別充足述語ではない。
本条件は入力側の構造の部分的な形式化である。
外生背景成分の列の生成則などの独立の環境動学は Phase~1 では定式化せず、
行動応答成分の遷移は定義~\ref{def:env_realization} の写像 $g_O$ として与えるにとどめる
（\S\ref{subsec:limitations_list} の L2）。

本条件は、特定の外部対象との関係が現在実現していること、
当該外部状態へ実際に到達可能であること、
必要環境要因が充足されていること、
EU の時刻別充足述語が真であること、
または関係品質が正であることを、その定義条件に含まない。
\end{definition}

\begin{remark}[構造的環境結合条件と一般化勾配感応性の区別]
\label{rem:ro_two_layer_intuition}
定義~\ref{def:env_coupling} の構造的環境結合条件は、
外部環境状態の差異が環境応答成分の次状態へ
非自明に反映されうるという\emph{構造的結合可能性}を記述する。

これに対して、命題~\ref{prop:pi_E_sensitivity} の条件
\[
  \exists\,g\in\partial\Pi(x_t)
  \quad\text{s.t.}\quad
  \pi_E(g)\neq\mathbf{0}_{n_E}
\]
は、現在の可行性余裕度を規定する最弱成分の少なくとも一つが
環境応答座標に属することを表す\emph{局所的な可行性感応条件}である。

両者は異なる条件であり、本論文はいずれの向きの含意も一般的な帰結として置かない。
したがって本論文は、構造的結合可能性、局所的可行性感応性、
実現外部相互作用、EU、および関係品質を相互に区別する。
\end{remark}

\begin{remark}[出力側の構造的手がかり：外部作用行動]
\label{rem:u_ext_output_cue}
\S\ref{sec:acs_def} の備考~\ref{rem:action_space_structure} は、
外部環境へ直接作用する行動の集合 $U_{\mathrm{ext}}\subseteq U$ を置く。
$U_{\mathrm{ext}}$ は、ACS から外部環境へ向かう作用の、
出力側の構造的手がかりを与える。
ただし、$U_{\mathrm{ext}}$ が空でないことだけからは、
実現した環境相互作用、臨界必要要因の存在、EU の時刻別充足述語の真のいずれも
一般的な帰結として置かない。
$U_{\mathrm{ext}}$ は三原理の時刻別充足述語の条件ではない。
\end{remark}

\subsubsection{$\pi_E$ と $\Pi$ の関係に関する命題}

\begin{proposition}[$\pi_E$ の最小感応性：同値版]
\label{prop:pi_E_sensitivity}
仮定~\ref{ass:env_response} および命題~\ref{prop:pi_subdiff} のもとで、
\[
  \bigl[\exists\,g\in\partial\Pi(x_t)\;\text{ s.t. }\;\pi_E(g)\neq\mathbf{0}_{n_E}\bigr]
  \;\iff\;
  I_E\cap\arg\min_k\bigl(x_{k,t}-\theta_k\bigr)\neq\emptyset.
\]
\end{proposition}

\begin{proof}
命題~\ref{prop:pi_subdiff}(ii) より、
\[
  \partial\Pi(x_t)
  =
  \operatorname{conv}
  \bigl\{
    e_k
    \bigm|
    k\in\arg\min_j(x_{j,t}-\theta_j)
  \bigr\}
\]
である。

\textbf{($\Rightarrow$)}\;
任意の $g\in\partial\Pi(x_t)$ は
\[
  g
  =
  \sum_{k\in\arg\min_j(x_{j,t}-\theta_j)}
  \lambda_k e_k
  \qquad
  (\lambda_k\ge0,\quad \textstyle\sum_k\lambda_k=1)
\]
と表される。
$\pi_E$ の線形性より
$\pi_E(g)=\sum_{k\in\arg\min_j(x_{j,t}-\theta_j)}\lambda_k\,\pi_E(e_k)$
となる。
$k\notin I_E$ のとき $\pi_E(e_k)=\mathbf{0}_{n_E}$ であるため、
$I_E\cap\arg\min_j(x_{j,t}-\theta_j)=\emptyset$ と仮定すると、
すべての $k\in\arg\min_j(x_{j,t}-\theta_j)$ に対して
$\pi_E(e_k)=\mathbf{0}_{n_E}$ となり、
任意の $g\in\partial\Pi(x_t)$ について
$\pi_E(g)=\mathbf{0}_{n_E}$ を得る。
よって対偶により、
$\pi_E(g)\neq\mathbf{0}_{n_E}$ を満たす
$g\in\partial\Pi(x_t)$ が存在するならば、
$I_E\cap\arg\min_j(x_{j,t}-\theta_j)\neq\emptyset$ である。

\textbf{($\Leftarrow$)}\;
$k^*\in I_E\cap\arg\min_j(x_{j,t}-\theta_j)$ を取る。
$k^*\in\arg\min_j(x_{j,t}-\theta_j)$ より
$e_{k^*}\in\partial\Pi(x_t)$ である。
また $k^*\in I_E$ より $k^*=k_r$ となる
$r\in\{1,\ldots,n_E\}$ が存在し、
\[
  \pi_E(e_{k^*})
  =
  e_r^{(n_E)}
  \neq
  \mathbf{0}_{n_E}
\]
となる（ここで $e_r^{(n_E)}$ は
$\mathbb{R}^{n_E}$ の第 $r$ 標準基底ベクトル）。
したがって $g:=e_{k^*}\in\partial\Pi(x_t)$ は
$\pi_E(g)\neq\mathbf{0}_{n_E}$ を満たす。
\end{proof}

\begin{remark}[$x_E$ 方向の局所感応性]
\label{rem:pi_E_local_sensitivity}
命題~\ref{prop:pi_E_sensitivity} の条件
\[
  I_E\cap\arg\min_k(x_{k,t}-\theta_k)\neq\emptyset
\]
が成立し、
かつ $x_t\in\operatorname{int}X$ である場合には、
最小値を達成する環境応答成分の一つを十分小さく減少させることで、
$\Pi(x_t)$ の値も減少する。
したがって、この条件は内点において
$\Pi$ が少なくとも一つの $x_E$ 方向に
局所的に非定値であることと整合する。
境界点では許容される摂動方向が $X$ の局所形状に依存するため、
一般的な方向解析は Phase~2 の課題とする。
\end{remark}

\subsubsection{環境応答成分の閾値割れと診断指標の役割}

\begin{definition}[環境応答成分の閾値割れ集合（分類C）]
\label{def:ro_related_violation}
環境応答成分の添字集合 $I_E$（定義~\ref{def:x_E}）を用いて
\emph{環境応答成分の閾値割れ集合} $V^-_E(x_t)$ を
\[
  V^-_E(x_t)\;:=\;I_E\cap V^-(x_t)
\]
と定義する。
$V^-_E(x_t)\neq\emptyset$ は、
環境応答成分の少なくとも一つが閾値を下回っていることを意味する。
\end{definition}

\begin{remark}[$\min$ 型 $\Pi$ の診断限界と補完指標の役割]
\label{rem:diagnostic_role}
$\Pi(x_t)=\min_k(x_{k,t}-\theta_k)$ は
最悪成分のみを代表するスカラーであり、
複数の閾値割れを識別しない。
たとえば
$(x_{k,t}-\theta_k)_k=(-2,-1,+3)$ のとき
$\Pi=-2$ となり、
第二成分の閾値割れ（$-1$）は $\Pi$ 単独では診断上区別されない。
$V^-(x_t)$ は割れている成分の全体を、
$d^-(x_t)$ は各成分の不足量を明示的に保持することで、この診断限界を補完する。

各指標の役割は次のように分担される。
$\Pi$ は可行・不可行判定および最悪成分の代表スカラー、
$\partial\Pi=\partial^C\widetilde{\Pi}$ は最弱成分の方向、
$V^-$ は閾値割れ成分の全体、
$d^-$ は各成分の不足量、
$V^-_E$ は環境応答成分の閾値割れを、それぞれ理論上表す。
これらは診断上の理論量であり、ACS 内部の感知・評価が同じ量を検出することを表さない。

$V^-_E(x_t)$ は、
命題~\ref{prop:pi_E_sensitivity} の
\[
  I_E\cap
  \arg\min_k(x_{k,t}-\theta_k)
\]
とは一般に異なる集合である。
後者は最弱成分
（$\arg\min$）のうち環境応答成分であるものを表すのに対し、
$V^-_E$ は最弱であるか否かによらず、
閾値を下回る環境応答成分を表す。
したがって、環境応答成分が閾値割れしていても
最弱成分でない場合、
$\partial\Pi$ に基づく環境応答座標への局所感応性判定
（命題~\ref{prop:pi_E_sensitivity}）はこれを表さないが、
$V^-_E$ はこれを表す。

これらの補完指標は、
本節における三原理の時刻別充足述語の定義および
構造的命題には要しない診断上の精密化であり、
その動学的利用は Phase~2 の課題とする
（within present framework）。

したがって
$V^-_E(x_t)\neq\emptyset$
から、構造的環境結合条件（定義~\ref{def:env_coupling}）を一般的な帰結として置かない。
$V^-_E$ は環境応答成分の閾値割れを理論上表す診断指標であり、
命題~\ref{prop:pi_E_sensitivity} の環境応答座標への局所感応性条件、
および定義~\ref{def:env_coupling} の構造的環境結合条件とは
それぞれ別概念である。
これらの間に一般的な必要十分関係は仮定しない。
\end{remark}

\subsubsection{実現外部相互作用の記述}

\begin{definition}[外部相互作用指標 $\chi_{\mathrm{EI},t}$（分類C）]
\label{def:phi}
時刻 $t\in\mathbb{N}_0$ における ACS $\Sigma$ の
\emph{外部相互作用指標}
\[
  \chi_{\mathrm{EI},t}\in\{0,1\}
\]
を、自己--環境境界を横断する外部作用が
時刻 $t$ に実現しているか否かを表す記述量として
\[
  \chi_{\mathrm{EI},t}
  :=
  \begin{cases}
    1,
    &
    \text{境界横断的な外部作用が少なくとも一つ実現している場合},
    \\[4pt]
    0,
    &
    \text{そのような外部作用が一つも実現していない場合}
  \end{cases}
\]
と定義する。

ここでいう境界横断的な外部作用は、
外部から ACS への入力だけに限定しない。
自己--環境境界を横断する作用である限り、
外部から内部への取り込み、
内部から外部への作用・排出・放散、
および双方向的な作用を含みうる。

$\chi_{\mathrm{EI},t}=1$ は、
少なくとも一つの境界横断作用が
時刻 $t$ に実現していることのみを表す。
したがって、それだけから
必要環境要因が充足されていること、
当該相互作用が可行性維持に十分であること、
または当該相互作用の関係品質が正であることを、本論文は一般的な帰結として置かない。

一方、
\[
  \chi_{\mathrm{EI},t}=0
\]
を、本論文では時刻 $t$ における
\emph{完全孤立状態}と呼ぶ。
これは境界横断的作用の実現がないことを表す記述であり、
構造的環境結合条件が成り立たないこと、$x_t\notin K$、
EU の時刻別充足述語の偽、閉塞または崩壊の別名ではない。
\end{definition}

\begin{remark}[$\chi_{\mathrm{EI}}$ の理論上の地位]
\label{rem:chi_ei_status}
$\chi_{\mathrm{EI},t}$ は分類Cの記述・診断量であり、
7タプル $\Sigma$ の新たな構成要素ではない。
また、実在 ACS が同一の二値変数を
内部で保持または計算することを要求しない。

本指標は、
「その時刻に境界横断作用が実現しているか」
という実現状態を理論上区別するための数学的表現である。
作用の強度、流量、方向別分解、持続時間、
複数相互作用の集約、および操作的測定法は
Phase~1 では固定しない。
\end{remark}

\begin{definition}[終端的完全孤立]
\label{def:terminal_isolation}
ある時刻 $t_0\in\mathbb{N}_0$ が存在して
\[
  \forall t\ge t_0,
  \qquad
  \chi_{\mathrm{EI},t}=0
\]
が成立するとき、
その実現軌道は時刻 $t_0$ 以降
\emph{終端的完全孤立}にあるという。

終端的完全孤立は、
完全孤立状態がある一時刻に成立することより強い
時間的記述である。
ただし、本定義それ自体は
可行性喪失、閉塞または崩壊を帰結として含まない。
\end{definition}

\subsubsection{構造的環境結合・実現外部相互作用・必要環境充足・EU の区別}

\begin{remark}[構造的結合可能性と実現相互作用の区別]
\label{rem:ro_realized_interaction_distinction}
定義~\ref{def:env_coupling} の構造的環境結合条件は、
状態 $x_t$ を基準として、
異なる外部環境状態 $e,e'\in E_{\mathrm{ext}}$ が
環境応答成分の次状態に異なる結果を与えうるという
構造的結合可能性を記述する。

これに対し、
$\chi_{\mathrm{EI},t}$ は
時刻 $t$ に実際の境界横断作用が
成立しているか否かを記述する。

したがって両者は、
\[
  \text{結合構造}
  \qquad\text{と}\qquad
  \text{実現イベント}
\]
という異なる論理的対象である。
Phase~1 の一般定義は、
\[
  \mathrm{EC}(x_t)
  \Longrightarrow
  \chi_{\mathrm{EI},t}=1
\]
も、
\[
  \chi_{\mathrm{EI},t}=1
  \Longrightarrow
  \mathrm{EC}(x_t)
\]
も仮定しない。
\end{remark}

\begin{remark}[必要環境空間・構造的環境結合・EU は同一ではない]
\label{rem:ereq_not_ro}
必要環境空間 $E$ と
必要環境充足集合 $E_{\mathrm{req}}(x)$ は、
ACS の存続が外部環境条件と無関係ではないという
外部依存の背景構造を与える。

これに対し構造的環境結合条件は、
定義~\ref{def:env_coupling} により、
外部環境状態の差異が
ACS の環境応答成分の次状態へ
非自明に反映されうるという
ACS 側の構造的な結合を記述する。
また EU は、時刻ごとの臨界必要要因の充足を問う時刻別充足述語である。

したがって、
\[
  E,\quad E_{\mathrm{req}}(x),\quad \mathrm{EC}(x),\quad \mathrm{EU}_t
\]
は相互に関連するが同一ではない。
とくに本論文は、
$E_{\mathrm{req}}(x_t)\neq\emptyset$ であることのみから
$\mathrm{EC}(x_t)$ を、
$\mathrm{EC}(x_t)$ から
$E_{\mathrm{req}}(x_t)\neq\emptyset$ を、
$\mathrm{EC}(x_t)$ から $\mathrm{EU}_t$ の真を、
いずれも一般的な帰結として置かない。
\end{remark}

\begin{remark}[補助記述・診断量と EU の真理条件]
\label{rem:env_coupling_not_eu_truth}
本小節の量は、EU の時刻別充足述語の真理条件ではない。
とくに本論文は、
$\chi_{\mathrm{EI},t}=0$ から $\mathrm{EU}_t$ の偽を、
命題~\ref{prop:pi_E_sensitivity} の環境応答方向への感応性から $\mathrm{EU}_t$ の真を、
いずれも一般的な帰結として置かない。
EU の時刻別充足述語は、定義~\ref{def:EU} の臨界必要要因の充足だけで定まる。
\end{remark}

\begin{remark}[実現相互作用と必要環境充足の区別]
\label{rem:ei_vs_ereq}
\S\ref{sec:acs_def} で定義された
必要環境充足集合
\[
  E_{\mathrm{req}}(x_t)
  =
  \bigl\{
    e\in E_{\mathrm{ext}}
    \bigm|
    \forall\eta\in E,\;
    \eta(x_t,e)=1
  \bigr\}
\]
に対し、
$e_t\in E_{\mathrm{req}}(x_t)$
は、現在状態 $x_t$ に相対して
実現外部環境状態 $e_t$ が
すべての必要環境要因を充足していることを表す。

これは、
境界横断作用が現在実現していることを表す
$\chi_{\mathrm{EI},t}=1$
とは別の条件である。
本論文は一般に、
\[
  e_t\in E_{\mathrm{req}}(x_t)
  \Longrightarrow
  \chi_{\mathrm{EI},t}=1
\]
も、
\[
  \chi_{\mathrm{EI},t}=1
  \Longrightarrow
  e_t\in E_{\mathrm{req}}(x_t)
\]
も仮定しない。

同様に、
$e_t\notin E_{\mathrm{req}}(x_t)$
は、少なくとも一つの必要環境要因が
現在の $(x_t,e_t)$ において充足されていないことを表す。
本論文は、これから完全孤立
$\chi_{\mathrm{EI},t}=0$
を一般的な帰結として置かない。
逆に、完全孤立から
$e_t\notin E_{\mathrm{req}}(x_t)$ を
定義だけで導くこともしない。
\end{remark}

\subsubsection{実現相互作用に関する時間的限定}

\begin{remark}[一時的孤立と終端的完全孤立の区別]
\label{rem:temporary_terminal_isolation}
単一時刻における
$\chi_{\mathrm{EI},t}=0$
と、定義~\ref{def:terminal_isolation} の
終端的完全孤立
\[
  \exists t_0\;
  \forall t\ge t_0:\
  \chi_{\mathrm{EI},t}=0
\]
は区別される。

前者は一時的な実現相互作用の不在であり、
後者はある有限時刻以降の恒常的な実現相互作用不在である。
ただし後者についても、
CR-1 に関係する必要環境要因の持続的欠損、
$K$ 外部への持続、閉塞または崩壊との接続は
別途示されなければならない。

したがって、
時間長だけを根拠として
\[
  \chi_{\mathrm{EI},t}=0
  \Longrightarrow
  \mathrm{Col}_{i,t}=1
\]
という一般則を導入しない。
\end{remark}

\subsubsection{本小節の位置づけ}

以上の構造的環境結合条件、実現外部相互作用、必要環境充足、EU の時刻別充足述語は、
備考~\ref{rem:conv_ro_layers} に示した異なる記述水準に属し、同一視しない。
構造的環境結合条件は、必要環境要因の充足が問われる自己と環境の関係の背景構造であり、
次小節では、その充足を EU の時刻別充足述語として定める。

%----------------------------------------------------------------------
\subsection{EU}
\label{subsec:eu}

本小節では、EU の時刻別充足述語を定める。
EU が問うのは、必要環境要因のすべてが常時充足されていることではなく、
充足されなければ次状態を $K$ に置く行動が残らない局面において、
当該要因が充足されていることである。
この局面は、現在の状態が $K$ の内にあるか外にあるかによらず生じうる。
$K$ の内では、要因の充足は可行性の維持に必要であり、
$K$ の外では、一段階での $K$ への回復に必要である。
この局面を、要因ごとの一段階臨界域と一段階臨界必要要因集合によって定める。
いずれも \S\ref{sec:acs_def} の必要環境空間 $E$ と可行行動対応 $A$ だけから定まり、
新たな仮定を用いない。

\begin{definition}[一段階臨界域（分類C）]
\label{def:eu_critical_domain}
必要環境要因 $\eta\in E$ について、
次の (i)〜(iii) をすべて満たす $x\in X$ の集合を $\mathrm{Crit}^{(1)}_\eta\subseteq X$ と書き、
$\eta$ の\emph{一段階臨界域}と呼ぶ。
\begin{enumerate}[label=(\roman*)]
  \item ある $e^0\in E_{\mathrm{ext}}$ について $\eta(x,e^0)=0$。
  \item すべての $e'\in E_{\mathrm{ext}}$ について、$\eta(x,e')=0$ ならば $A(x,e')=\emptyset$。
  \item ある $e\in E_{\mathrm{ext}}$ について、$\eta(x,e)=1$ かつ $A(x,e)\neq\emptyset$。
\end{enumerate}
\end{definition}

各条件は、$\eta:X\times E_{\mathrm{ext}}\to\{0,1\}$ と
可行行動対応 $A:X\times E_{\mathrm{ext}}\rightrightarrows U$（定義~\ref{def:at}）の値だけで定まり、
$x\in X$ のすべてについて意味をもつ。
(i) は、(ii) が空虚に成り立つ場合を除くための条件である。
(ii) は、当該要因が充足されない外部環境状態ではどの行動も次状態を $K$ に置かないことを、
(iii) は、当該要因が充足される外部環境状態のなかに次状態を $K$ に置く行動がある場合があることを表す。
$x\in K$ では、$A(x,e)$ は可行性を維持する行動の集合であり、
$x\notin K$ では、一段階で $K$ へ回復する行動の集合である（定義~\ref{def:at}）。
$\mathrm{Crit}^{(1)}_\eta$ は一段階の条件であり、$K$ の位相的境界とは同一視しない。
本論文は $\mathrm{Crit}^{(1)}_\eta\neq\emptyset$ を一般には主張しない。

\begin{definition}[一段階臨界必要要因集合（分類C）]
\label{def:eu_fire}
写像 $\mathrm{Crit}^{(1)}:X\to 2^{E}$ を
\[
  \mathrm{Crit}^{(1)}(x)
  :=
  \bigl\{
    \eta\in E
    \bigm|
    x\in\mathrm{Crit}^{(1)}_\eta
  \bigr\}
\]
で定め、$\mathrm{Crit}^{(1)}(x)$ を $x$ における\emph{一段階臨界必要要因集合}と呼ぶ。
定義により $\eta\in\mathrm{Crit}^{(1)}(x)\iff x\in\mathrm{Crit}^{(1)}_\eta$ であり、
$\mathrm{Crit}^{(1)}_\eta\subseteq X$ は状態の集合、$\mathrm{Crit}^{(1)}(x)\subseteq E$ は要因の集合である。
下付きの $\eta$ との混同を避けるため、時刻については $\mathrm{Crit}^{(1)}(x_t)$ と書き、時刻添字による略記は置かない。
$\mathrm{Crit}^{(1)}(x)$ は $x$ だけで定まり、実現外部環境状態には依存しない。
$\eta\in\mathrm{Crit}^{(1)}(x)$ であることを、$\eta$ が状態 $x$ で\emph{臨界必要要因である}という。
$E$ のすべての元は必要環境要因であるが、そのうち当該状態で一段階の可行性に関して臨界となっているものが臨界必要要因である。
上付きの $(1)$ は一段階の条件であることを表す（備考~\ref{rem:eu_one_step_scope}）。
\end{definition}

\begin{remark}[一段階臨界必要要因集合の記述水準]
\label{rem:crit_description_level}
$\mathrm{Crit}^{(1)}(x)$ は、$\eta$、$f$、$K$ から理論上定まる分類Cの量であり、
ある要因が当該状態で一段階の可行性の維持または回復に関して臨界となっていることを、
ACS がそれを認識しているか否かとは独立に記述する。
本論文では次の三つを区別する。
\begin{enumerate}
  \item 理論上の臨界必要要因：$\mathrm{Crit}^{(1)}(x_t)$。
  \item ACS 内部での必要性の推定・評価：感知・識別・評価の機能と、
  \S\ref{sec:tipa} の TIPA 等の標準機能形式によって形成される内部の評価。
  \item 外部観測者による測定・診断：観測可能な情報と測定規則に依存する。
\end{enumerate}
理論上は臨界必要要因である要因が内部で必要と評価されないこと（見逃し）も、
理論上は臨界必要要因でない要因が内部で必要と評価されること（誤評価）も、
本論文の一般構造は排除しない。
EU の時刻別充足述語（定義~\ref{def:EU}）の真理条件は (1) だけで定まり、
(2) の内部評価の結果も、(3) の測定結果も用いない。
(2) の標準的な記述は \S\ref{sec:tipa} の責務であり、本節ではそのための記号を置かない。
ACS がその内部評価で $\mathrm{Crit}^{(1)}(x_t)$ を保持または計算することは要求しない
（備考~\ref{rem:s3_description_levels}）。

本論文は、三つの水準を語彙の上でも分ける。
(1) には「臨界である」「必要である」「充足される」を、
(2) には「内部で推定する」「検出する」「評価する」「発火する」「作動する」を、
(3) には「観測者が推定する」「観測する」「測定する」「診断する」を用いる。
「推定」は (2) と (3) のいずれにも現れうるので、推定の主体を明示する。
とくに「発火」は (2) の内部の作動を表す語として留保し、(1) の理論上の量には用いない。
\end{remark}

\begin{definition}[EU の時刻別充足述語]
\label{def:EU}
$t\in\mathbb{T}_{\mathrm{P},i}$ について、
EU の時刻別充足述語 $\mathrm{EU}_t$ は、
\[
  \forall\eta\in\mathrm{Crit}^{(1)}(x_t):\quad \eta(x_t,e_t)=1
\]
のとき真、そうでないとき偽とする。
$t\notin\mathbb{T}_{\mathrm{P},i}$ の時刻には、$\mathrm{EU}_t$ を定義しない。
\end{definition}

\begin{remark}[EU の真が表すことの範囲]
\label{rem:eu_vacuous}
$\mathrm{Crit}^{(1)}(x_t)=\emptyset$ のとき、定義~\ref{def:EU} の全称量化は空虚に成立し、
$\mathrm{EU}_t$ は真である。
$\mathrm{EU}_t$ が真であることの定義条件は、時刻 $t$ の臨界必要要因がすべて充足されていることだけであり、
臨界必要要因でない要因の充足を含めない。
したがって $\mathrm{EU}_t$ の真は、すべての必要環境要因が利用可能であること、
または環境を利用する機能が健全であることを表さない
（備考~\ref{rem:eu_env_utilization}）。
$\mathrm{EU}_t$ の定義域は $\mathrm{SM}_t$ の真を前提としない。
SM と EU は相異なる制約軸として、それぞれ別個に判定される。
\end{remark}

\begin{proposition}[臨界必要要因の非充足と閉塞]
\label{prop:eu_blockage}
\begin{enumerate}[label=(\roman*)]
  \item 任意の $t\in\mathbb{N}_0$ と $\eta\in E$ について、
  $x_t\in\mathrm{Crit}^{(1)}_\eta$ かつ $\eta(x_t,e_t)=0$ ならば $A_t=\emptyset$ である。
  \item $t\in\mathbb{T}_{\mathrm{P},i}$ で $\mathrm{EU}_t$ が偽ならば、$A_t=\emptyset$ である。
\end{enumerate}
\end{proposition}

\begin{proof}
(i)：定義~\ref{def:eu_critical_domain} (ii) に $e'=e_t$ を代入すると $A(x_t,e_t)=\emptyset$ である。
(ii)：$\mathrm{EU}_t$ が偽ならば、ある $\eta\in\mathrm{Crit}^{(1)}(x_t)$ について $\eta(x_t,e_t)\neq1$、
すなわち $\eta(x_t,e_t)=0$ である。
$\eta\in\mathrm{Crit}^{(1)}(x_t)$ は $x_t\in\mathrm{Crit}^{(1)}_\eta$ を意味するので、(i) を適用する。
\end{proof}

\begin{proposition}[臨界必要要因の非充足と次時刻の状態]
\label{prop:eu_sm_next}
命題~\ref{prop:eu_blockage} の (i) または (ii) の前件が成り立つならば、
$x_{t+1}\notin K$ である。
さらに $t+1\in\mathbb{T}_{\mathrm{P},i}$ ならば、$\mathrm{SM}_{t+1}$ は偽である。
\end{proposition}

\begin{proof}
命題~\ref{prop:eu_blockage} により $A_t=\emptyset$ であり、
命題~\ref{prop:sm_one_step_A} (ii) と $u_t\in U$ により $x_{t+1}\notin K$ である。
後半は定義~\ref{def:SM} による。
\end{proof}

\begin{remark}[維持の局面と回復の局面]
\label{rem:eu_maintenance_recovery}
$x_t\in K$ で $\mathrm{EU}_t$ が偽ならば、命題~\ref{prop:eu_sm_next} により
次時刻に状態は $K$ を外れる。
$x_t\notin K$ で $\mathrm{EU}_t$ が偽ならば、次時刻にも状態は $K$ の外にある。
すなわち、臨界必要要因の非充足は、$K$ の内では一段階での可行性の維持を、
$K$ の外では一段階での $K$ への回復を妨げる。
いずれの場合も、本論文は $\mathrm{EU}_t$ の偽だけから、
多段の時間幅での回復不能、$\mathrm{WF}_{i,t}=0$ または崩壊を一般的な帰結として置かない。
\end{remark}

\begin{remark}[一段階の必要性の射程]
\label{rem:eu_one_step_scope}
$\mathrm{Crit}^{(1)}_\eta$ と $\mathrm{Crit}^{(1)}$ は一段階の条件である。
したがって次の二つは、$\mathrm{Crit}^{(1)}$ と EU の時刻別充足述語の真理条件に含まれない。
\begin{enumerate}
  \item 一時刻の欠如では次状態を $K$ に置く行動が残るが、欠如の持続によって可行性が失われる要因。
  \item $K$ の外の状態から数段階を経て $K$ へ回復するために要する要因。
\end{enumerate}
(1) のうち要因の非充足が持続する場合は、
構成要件 CR-1 の (CR-1a)(ii) が整形式実現の水準で扱う（備考~\ref{rem:eu_cr1_connection}）。
$K$ への多段の回復可能性・有限時間到達可能性と必要環境要因との接続は、
本論文の範囲では定式化しない（\S\ref{sec:limitations_and_outlook}）。
\end{remark}

\begin{remark}[EU と構成要件 CR-1 の接続範囲]
\label{rem:eu_cr1_connection}
$\mathrm{EU}_t$ が偽であることは、ある臨界必要要因 $\eta$ が時刻 $t$ に充足されないことであり、
命題~\ref{prop:eu_blockage} により時刻 $t$ の閉塞を伴う。
一方、整形式実現の喪失との接続は構成要件 CR-1 の (CR-1a)(ii)
（\S\ref{subsec:environment_space}）による。
すなわち、同一の要因 $\eta$ について、ある $t_0\ge1$ で
$\mathrm{WF}_{i,t_0-1}=1$、$\eta(x_{t_0-1},e_{t_0-1})=1$ であり、
すべての $s\ge t_0$ で $\eta(x_s,e_s)=0$ となるならば、
ある有限時刻 $T\ge t_0$ 以後のすべての時刻で $\mathrm{WF}_{i,s}=0$ である。
本論文は、$\mathrm{EU}_t$ が偽である時刻があることだけから、この前件を一般的な帰結として置かない。
臨界必要要因が時刻ごとに異なる場合、および要因の非充足が持続しない場合は、
(CR-1a)(ii) の前件を満たさない。
また (CR-1a)(ii) の帰結は整形式実現の喪失であって、
機能的終端 $\mathrm{FT}_{i,t}=1$ でも崩壊 $\mathrm{Col}_{i,t}=1$ でもない。
\end{remark}

\begin{proposition}[臨界必要要因と要因別充足状態域]
\label{prop:eu_fire_domain}
$\eta\in E$ とし、$D_\eta$ を \S\ref{sec:acs_def} の要因別充足状態域
（定義~\ref{def:eta_satisfiable_domain}）とする。
\begin{enumerate}[label=(\roman*)]
  \item $\mathrm{Crit}^{(1)}_\eta\subseteq D_\eta$ である。
  したがって、$\eta\in\mathrm{Crit}^{(1)}(x)$ ならば $x\in D_\eta$ である。
  \item $x\notin D_\eta$ ならば、$\eta\notin\mathrm{Crit}^{(1)}(x)$ である。
\end{enumerate}
\end{proposition}

\begin{proof}
(i)：定義~\ref{def:eu_critical_domain} の (iii) から、$\eta(x,e)=1$ となる $e\in E_{\mathrm{ext}}$ が存在するので、
定義~\ref{def:eta_satisfiable_domain} により $x\in D_\eta$ である。
後半は定義~\ref{def:eu_fire} による。
(ii)：(i) の対偶である。
\end{proof}

\begin{remark}[現在状態で充足不能な要因と EU]
\label{rem:eu_env_utilization}
$x\notin D_\eta$ は、現在の状態 $x$ では、当該要因 $\eta$ がどの外部環境状態によっても充足されないことを表す
（定義~\ref{def:eta_satisfiable_domain}）。
必要環境充足判定 $\eta(x,e)$ は内部状態 $x$ に依存するので、
摂取・取得・接続に関わる内部状態の変化もこの関係のなかに表される。
命題~\ref{prop:eu_fire_domain} (ii) により、このとき $\eta$ は臨界必要要因にならず、
EU の時刻別充足述語は $\eta$ の充足を問わない。
このような状態で $E_{\mathrm{req}}(x)=\emptyset$ となることは、
\S\ref{sec:acs_def} 備考~\ref{rem:env_utilization_function} の (a) による。
また、$x\notin D_\eta$ だけから閉塞、$\mathrm{WF}_{i,t}=0$、機能的終端または崩壊を一般的な帰結として置かないことは、
同備考の (c) と同じ扱いである。
このような状態で閉塞が生じる場合、その閉塞は命題~\ref{prop:blockage_dichotomy} の分割において、
$\eta$ の非充足を理由とする (B1) には数えられない。
本論文は、この場合を EU の側の事態として別に数える定義を置かない。

同備考の必要環境利用機能は、内部状態を $D_\eta$ に属させる条件の成立に関与する
既存機能の協働の総称であり、EU の原理応答項目に関与する。
ただし本論文は、必要環境利用機能が作動していることから $\mathrm{EU}_t$ の真を、
$\mathrm{EU}_t$ の偽から当該機能の停止を、いずれも一般的な帰結として置かない
（備考~\ref{rem:s3_description_levels}）。
\end{remark}

\begin{remark}[CR-1 と実現外部相互作用の接続範囲]
\label{rem:cr1_ei_connection}
構成要件 CR-1（\S\ref{subsec:environment_space}）は、
$E$ の各要素が
構造的自己の存続にとって実質的な必要環境要因であることを、
欠損の実現可能性、欠損時の存続不能性、
および共同充足可能性によって規定する。

したがって、ある必要環境要因 $\eta\in E$ について、
CR-1 が定める持続的欠損の前件が
実現軌道上で成立する場合には、
その帰結は CR-1 に従って評価される。
しかし、
$\chi_{\mathrm{EI},t}=0$
という事実だけから、
特定の $\eta$ の持続的欠損条件が成立したとはいえない。

この区別は、
必要環境要因の充足判定が
\[
  \eta:X\times E_{\mathrm{ext}}\to\{0,1\}
\]
として内部状態 $x_t$ にも依存することと整合する。
内部備蓄、内部補償、状態依存の許容時間等によって、
一時的な外部作用の欠如が
直ちに必要環境要因の持続的欠損へ一致しない場合を
Phase~1 の一般構造は排除しない。

したがって本論文では、
「終端的完全孤立だから必ず可行性を失う」
という無条件命題を置かない。
終端的完全孤立と存続不能性を接続する場合には、
少なくとも、当該孤立が CR-1 における
必要環境要因の持続的欠損条件へ実際に接続することを
別途確認または仮定する必要がある。
\end{remark}

%----------------------------------------------------------------------
\subsection{FR}
\label{subsec:fr}

本小節では、FR の時刻別充足述語（局所 FR）を定める。
局所 FR は、時刻 $t$ の臨界必要要因がすべて充足されている局面に限って、
次状態を $K$ に置く行動が存在するかを問う。
現在の状態が $K$ の内にあれば、これは可行性を維持する行動の存在であり、
$K$ の外にあれば、一段階で $K$ へ回復する行動の存在である。
臨界必要要因が充足されないことに伴う閉塞は EU の偽として扱い、局所 FR の判定には含めない。
あわせて、応答の多様性を表す診断量として応答多様度を定める。

\begin{definition}[FR の定義域]
\label{def:fr_domain}
\[
  \mathcal{D}_{\mathrm{FR}}
  :=
  \bigl\{
    (x,e)\in X\times E_{\mathrm{ext}}
    \bigm|
    \forall\eta\in\mathrm{Crit}^{(1)}(x):\ \eta(x,e)=1
  \bigr\}
\]
を FR の\emph{定義域}と呼ぶ。
$\mathrm{Crit}^{(1)}$ は $X$ 上の写像であるから（定義~\ref{def:eu_fire}）、条件はすべての $(x,e)$ で定まる。
$\mathrm{Crit}^{(1)}(x)=\emptyset$ のとき、すべての $e\in E_{\mathrm{ext}}$ について $(x,e)\in\mathcal{D}_{\mathrm{FR}}$ である。
\end{definition}

\begin{definition}[FR の時刻別充足述語（局所 FR）]
\label{def:FR}
$t\in\mathbb{T}_{\mathrm{P},i}$ かつ $(x_t,e_t)\in\mathcal{D}_{\mathrm{FR}}$ の時刻について、
FR の時刻別充足述語 $\mathrm{FR}_t$ は
\[
  A_t\neq\emptyset
\]
のとき真、$A_t=\emptyset$ のとき偽とする。
これを\emph{局所 FR} とも呼ぶ。
$t\notin\mathbb{T}_{\mathrm{P},i}$ または $(x_t,e_t)\notin\mathcal{D}_{\mathrm{FR}}$ の時刻には、
$\mathrm{FR}_t$ を定義しない。
\end{definition}

\begin{remark}[FR の定義上の限定]
\label{rem:fr_scope}
\begin{enumerate}
  \item 「可行」は本論文の可行性、すなわち $f(x,u,e)\in K$ に固定する。
  行動の実行可能性一般、行動の試行の成否、または実行の成功の度合いを意味しない
  （\S\ref{sec:acs_def} 備考~\ref{rem:f_reading}）。
  \item 制約軸としての FR は、必要な局面で可行な応答が得られること、
  およびその必要応答ポテンシャルの保持・獲得・回復・再構成に関わる。
  局所 FR はその一時刻の側面であり、長期の側面は \S\ref{subsec:sd_necessity} と
  \S\ref{subsec:s3_tstab} で部分的に扱う。
  \item $\mathcal{D}_{\mathrm{FR}}$ 上で $\mathrm{FR}_t$ の偽と閉塞 $A_t=\emptyset$ は一致する。
  しかし FR は閉塞の言い換えではない。
  閉塞は時刻ごとの状態の診断であり、
  FR は長期の側面を含む応答側の制約軸である。
  \item $x_t\notin K$ の時刻の局所 FR は、一段階で $K$ へ回復する行動が存在するかを問う。
  この時刻に $\mathrm{FR}_t$ が偽であることは、一段階では $K$ へ戻れないことであり、
  数段階を経た $K$ への回復の可能性については何も述べない（備考~\ref{rem:eu_one_step_scope}）。
  この時刻に当該ACSが回復に寄与する行動を実行していることは排除されず、
  $\mathrm{FR}_t$ の偽は、回復に寄与する行動が存在しないこと、回復の過程が進んでいないこと、
  または応答の機構が損なわれていることを意味しない
  （\S\ref{sec:acs_def} 備考~\ref{rem:blockage_scope_sec2}）。
  \item 自己発展は、必要応答ポテンシャルを獲得・拡張する原理応答項目の一つの様式として位置づける。
  \item 「応答能力」は一般語としてだけ用い、形式上の用語は「必要応答ポテンシャル」とする。
\end{enumerate}
\end{remark}

\begin{remark}[定義域の役割]
\label{rem:fr_domain_role}
$\mathcal{D}_{\mathrm{FR}}$ は $A_t$ の値を変えない。
$\mathcal{D}_{\mathrm{FR}}$ が定めるのは、閉塞のうち、
どれを FR の時刻別充足述語が偽である時刻として数えるかである。
臨界必要要因が充足されないことに伴う閉塞は EU の偽として扱い、FR の判定には含めない
（命題~\ref{prop:blockage_dichotomy}）。
$\mathcal{D}_{\mathrm{FR}}$ は $x\in K$ を条件に含まないので、
$K$ の外での一段階の回復の可否も FR の判定に含まれる。
\end{remark}

\begin{definition}[応答多様度（分類C）]
\label{def:mu_eff}
$(x,e)\in X\times E_{\mathrm{ext}}$ を固定する。
各 $u\in A(x,e)$ に対し、可行性保持指示関数
$\varphi_u:E_{\mathrm{req}}(x)\to\{0,1\}$ を
\[
  \varphi_u(e'):=\mathbf{1}\bigl[f(x,u,e')\in K\bigr]
\]
で定め、$A(x,e)$ 上の同値関係 $\sim$ を
$u\sim u' :\Longleftrightarrow \varphi_u=\varphi_{u'}$（$E_{\mathrm{req}}(x)$ 上の関数として等しい）で定める。
\emph{応答多様度}を商集合の基数
\[
  \mu_{\mathrm{eff}}(x,e):=\bigl|A(x,e)/\!\sim\bigr|
\]
として定め、$\mu_{\mathrm{eff},t}:=\mu_{\mathrm{eff}}(x_t,e_t)$ と書く。
$\mu_{\mathrm{eff}}$ の有限性は仮定しない。
$\mu_{\mathrm{eff}}$ はすべての $(x,e)$ で定まる診断量であり、
いずれの原理の時刻別充足述語の定義にも用いない。
\end{definition}

\begin{proposition}[定義域と EU の時刻別充足述語]
\label{prop:fr_domain_eu}
$t\in\mathbb{T}_{\mathrm{P},i}$ ならば、
\[
  \mathrm{EU}_t\ \text{が真}
  \iff
  (x_t,e_t)\in\mathcal{D}_{\mathrm{FR}}
\]
である。
したがって、$t\in\mathbb{T}_{\mathrm{P},i}$ の時刻に $\mathrm{FR}_t$ が定義されることは、
$\mathrm{EU}_t$ が真であることと同値である。
\end{proposition}

\begin{proof}
$\mathrm{EU}_t$ が真である条件（定義~\ref{def:EU}）と
$(x_t,e_t)\in\mathcal{D}_{\mathrm{FR}}$ の条件（定義~\ref{def:fr_domain}）は、
いずれも $\forall\eta\in\mathrm{Crit}^{(1)}(x_t):\eta(x_t,e_t)=1$ である。
後半は定義~\ref{def:FR} による。
\end{proof}

\begin{proposition}[系：局所 FR と次状態の $K$ への所属]
\label{prop:fr_one_step}
$t\in\mathbb{T}_{\mathrm{P},i}$ かつ $(x_t,e_t)\in\mathcal{D}_{\mathrm{FR}}$ ならば、
\[
  \mathrm{FR}_t\ \text{が真}
  \iff
  A_t\neq\emptyset
  \iff
  \exists u\in U:\ f(x_t,u,e_t)\in K
  \iff
  \exists u\in U:\ \Pi\bigl(f(x_t,u,e_t)\bigr)\ge0
\]
である。
\end{proposition}

\begin{proof}
第一の同値は定義~\ref{def:FR}、第二・第三は命題~\ref{prop:sm_one_step_A} (i) による。
\end{proof}

\begin{proposition}[局所 FR の偽と次時刻の状態]
\label{prop:fr_next_sm}
$t\in\mathbb{T}_{\mathrm{P},i}$、$(x_t,e_t)\in\mathcal{D}_{\mathrm{FR}}$ で $\mathrm{FR}_t$ が偽ならば、
$A_t=\emptyset$ であり、$x_{t+1}\notin K$ である。
さらに $t+1\in\mathbb{T}_{\mathrm{P},i}$ ならば、$\mathrm{SM}_{t+1}$ は偽である。
\end{proposition}

\begin{proof}
定義~\ref{def:FR} により $A_t=\emptyset$ である。
命題~\ref{prop:sm_one_step_A} (ii) と $u_t\in U$ により $x_{t+1}\notin K$ である。
後半は定義~\ref{def:SM} による。
\end{proof}

\begin{proposition}[局所 FR と応答多様度]
\label{prop:fr_mu_eff}
\begin{enumerate}[label=(\roman*)]
  \item 任意の $(x,e)\in X\times E_{\mathrm{ext}}$ について、
  $A(x,e)=\emptyset\iff\mu_{\mathrm{eff}}(x,e)=0$ である。
  \item $t\in\mathbb{T}_{\mathrm{P},i}$ かつ $(x_t,e_t)\in\mathcal{D}_{\mathrm{FR}}$ ならば、
  $\mathrm{FR}_t$ が真であることと $\mu_{\mathrm{eff},t}\ge1$ は同値である。
  とくに $\mu_{\mathrm{eff},t}\ge2$ ならば $\mathrm{FR}_t$ は真である。
\end{enumerate}
\end{proposition}

\begin{proof}
(i)：$A(x,e)=\emptyset$ ならば商集合は空であり、基数は $0$ である。
$A(x,e)\neq\emptyset$ ならば、$u\in A(x,e)$ の同値類が商集合の元となるので、基数は $1$ 以上である。
(ii)：定義~\ref{def:FR} と (i) による。
\end{proof}

\begin{remark}[$E_{\mathrm{req}}(x_t)=\emptyset$ の場合の応答多様度]
\label{rem:mu_eff_empty_req}
$E_{\mathrm{req}}(x_t)=\emptyset$ のとき、すべての $u\in A_t$ について $\varphi_u$ は空集合上の唯一の関数となるので、
$A_t\neq\emptyset$ ならば $\mu_{\mathrm{eff},t}=1$ である。
これは比較の対象となる必要環境充足状態がないことの集合論上の帰結であり、
応答の可能性が一種類に縮退していることを示すものではない。
命題~\ref{prop:fr_mu_eff} はこの場合も含めて成り立つ。
\end{remark}

\begin{remark}[環境比較基準の内在性]
\label{rem:env_basis_intrinsic}
$\varphi_u$、$\sim$、$\mu_{\mathrm{eff},t}$ の比較基準には、
\S\ref{sec:acs_def} の必要環境充足集合
$E_{\mathrm{req}}(x_t)$
を用いる。
これは、固定された外部環境状態空間 $E_{\mathrm{ext}}$ のもとで、
必要環境空間 $E$ と状態 $x_t$ から一意に派生する集合であり、
解析者が応答多様度の評価のため任意に選択する評価領域ではない。

また、
$E_{\mathrm{req}}(x_t)$ と
一段階可行環境対応
$\mathcal{E}_{\mathrm{viab}}^{(1)}(x_t)$
（定義~\ref{def:one_step_viable_environment}）との間に、
いずれの向きの包含関係も仮定しない。
\end{remark}

\begin{remark}[局所 FR と長期の側面の時間的区別]
\label{rem:sd_local_scope}
定義~\ref{def:FR} の局所 FR は、
固定された現在状態 $x_t$ と実現外部環境状態 $e_t$ を基準とする
一段階の条件である。
したがって本論文は、ある時刻 $t$ で $\mathrm{FR}_t$ が真であることから、
後続時刻 $s>t$ で $\mathrm{FR}_s$ が定義されて真であることを一般的な帰結として置かない。
応答多様度についても同様に、
$\mu_{\mathrm{eff},t}\ge2$ や $\mu_{\mathrm{eff},t}=1$ は時刻 $t$ の局所量であり、
後続時刻の値を定めない。

\S\ref{subsubsec:real_world_scope} で述べたとおり、
本論文が対象とする ACS は、時間変化する環境のもとで、
有限の計算資源、有限の応答時間、有限の記憶容量、
供給・エネルギー制約、および構成要素の損傷・劣化可能性を伴って作動する。
したがって、後続時刻 $s>t$ において
\[
  x_s=x_t,\qquad
  e_s=e_t,\qquad
  A_s=A_t,\qquad
  E_{\mathrm{req}}(x_s)=E_{\mathrm{req}}(x_t)
\]
が維持されることを一般には仮定しない。
各時刻の $x_s$ と $e_s$ に応じて、
$A_s$、$E_{\mathrm{req}}(x_s)$、$\mathrm{FR}_s$ および
$\mu_{\mathrm{eff},s}$ は改めて評価される。

また、Phase~1 で実現列を所与として解析することは、
ACS 自身または外部観測者による将来系列の事前把握可能性を
仮定するものではない（備考~\ref{rem:three_layers}）。
Phase~1 は将来系列に確率測度を導入しない。

FR の長期の側面、すなわち必要応答ポテンシャルの保持・獲得・回復・再構成について、
Phase~1 は独立した無限時間の判定条件を与えない。
局所 FR の有限窓内の持続は \S\ref{subsec:s3_tstab} の (T2) で、
応答多様性の縮退の持続と長期可行性との条件付き接続は
\S\ref{subsec:sd_necessity} で、
有限段の頑健性は \S\ref{subsec:viab_core} の $\mathrm{Hor}_f$ で、
それぞれ部分的に記述する。
内部状態の時刻間変化と応答方略の更新は
\S\ref{sec:extended_dynamics} の責務である。
長期の側面の本格的な理論は Phase~1.$\alpha$ の課題とする
（\S\ref{sec:limitations_and_outlook}）。
\end{remark}

\begin{remark}[可行性余裕度の改善可能性]
\label{rem:sd_improvement_aux}
条件
\[
  \exists\,u\in A_t
  \quad\text{s.t.}\quad
  \Pi\bigl(f(x_t,u,e_t)\bigr)>\Pi(x_t)
\]
は、現在の実現外部環境状態 $e_t$ のもとで
可行性余裕度を一段階で厳密に改善できる行動が存在することを表す。

本論文では、この条件を
\emph{可行性余裕度の改善可能性}を表す補助診断条件として保持するが、
定義~\ref{def:FR} の局所 FR の条件とはしない。
現在すでに高い可行性余裕をもつ状態や、
可行な応答を保ちつつ一段階では
$\Pi$ を厳密増加させない状態を、
この条件だけを理由に局所 FR の偽とはしないためである。
\end{remark}

\begin{proposition}[有限段頑健可行集合と局所 FR]
\label{prop:viab_stage_fr}
$\nu\ge1$、$x_t\in\mathrm{Viab}_f^{(\nu)}(K)$、$e_t\in E_{\mathrm{req}}(x_t)$ ならば、
$(x_t,e_t)\in\mathcal{D}_{\mathrm{FR}}$ かつ $A_t\neq\emptyset$ である。
さらに $t\in\mathbb{T}_{\mathrm{P},i}$ ならば、$\mathrm{FR}_t$ は真である。
\end{proposition}

\begin{proof}
定義~\ref{def:viab_stage} により、ある $u\in U$ が存在して、
すべての $e\in E_{\mathrm{req}}(x_t)$ について
$f(x_t,u,e)\in\mathrm{Viab}_f^{(\nu-1)}(K)\subseteq K_{\mathrm{req}}\subseteq K$ である。
$e_t\in E_{\mathrm{req}}(x_t)$ であるから $u\in A_t$ であり、$A_t\neq\emptyset$ である。
$e_t\in E_{\mathrm{req}}(x_t)$ と定義~\ref{def:E_req} により、すべての $\eta\in E$ について $\eta(x_t,e_t)=1$ である。
$\mathrm{Crit}^{(1)}(x_t)\subseteq E$（定義~\ref{def:eu_fire}）であるから、
すべての $\eta\in\mathrm{Crit}^{(1)}(x_t)$ について $\eta(x_t,e_t)=1$ であり、
$(x_t,e_t)\in\mathcal{D}_{\mathrm{FR}}$ である。
後半は定義~\ref{def:FR} による。
\end{proof}

\begin{proposition}[選択の責務]
\label{prop:selection_responsibility}
任意の $t\in\mathbb{N}_0$ について、$u_t\notin A_t$ ならば $x_{t+1}\notin K$ である。
とくに、$t\in\mathbb{T}_{\mathrm{P},i}$ で $\mathrm{SM}_t$・$\mathrm{EU}_t$・$\mathrm{FR}_t$ がすべて真であっても、
$u_t\notin A_t$ ならば $x_{t+1}\notin K$ である。
\end{proposition}

\begin{proof}
命題~\ref{prop:sm_one_step_A} (ii) による。
\end{proof}

\begin{remark}[選択の責務と構成要件]
\label{rem:selection_responsibility}
本論文は $u_t\in A_t$ を構成要件にも仮定にも置かない。
三原理の時刻別充足述語がすべて真であれば $A_t\neq\emptyset$ である（定義~\ref{def:FR}）。
しかし、時刻別充足述語の定義には $u_t\in A_t$ を含めず、
本論文は、三原理の時刻別充足述語がすべて真であることから $u_t\in A_t$ を要求しない。
実効行動の形成は \S\ref{sec:tipa} の TIPA・PSL の責務であり、
同節は $\mathrm{Cand}_t=A_t$ を要求しない。
\end{remark}

%----------------------------------------------------------------------
\subsection{三原理と閉塞・整形式実現・崩壊の関係}
\label{subsec:principle_failure_relation}

本小節では、閉塞を臨界必要要因の充足の有無によって二つに分け、
三原理の時刻別充足述語と閉塞・整形式実現・崩壊との関係を整理する。
二分の対象は、現在の状態が $K$ の内にある時刻の閉塞（一段階で可行性を維持できない局面）と、
$K$ の外にある時刻の閉塞（一段階では $K$ へ戻れない局面）の両方である。

\begin{proposition}[閉塞の二分]
\label{prop:blockage_dichotomy}
$A_t=\emptyset$ の時刻 $t$ について、次の (B1) と (B2) のちょうど一方が成り立つ。
\begin{description}[style=nextline]
  \item[(B1) 臨界必要要因非充足閉塞]
  ある $\eta\in\mathrm{Crit}^{(1)}(x_t)$ について $\eta(x_t,e_t)=0$。
  \item[(B2) 臨界必要要因充足下閉塞]
  $(x_t,e_t)\in\mathcal{D}_{\mathrm{FR}}$。
\end{description}
また、任意の時刻 $t$ について、(B1) の条件が成り立てば $A_t=\emptyset$ である。
さらに $t\in\mathbb{T}_{\mathrm{P},i}$ の時刻について、次が成り立つ。
\begin{enumerate}[label=(\roman*)]
  \item (B1) の閉塞の時刻とは、$\mathrm{EU}_t$ が偽である時刻であり、
  このとき $(x_t,e_t)\notin\mathcal{D}_{\mathrm{FR}}$ であるから $\mathrm{FR}_t$ は定義されない。
  \item (B2) の閉塞の時刻とは、$\mathrm{EU}_t$ が真であり、かつ $\mathrm{FR}_t$ が偽である時刻である。
  \item $A_t=\emptyset$ であることは、
  「$\mathrm{EU}_t$ が偽である」か「$\mathrm{EU}_t$ が真であり、かつ $\mathrm{FR}_t$ が偽である」かの
  いずれかが成り立つことと同値であり、この二つは同時には成り立たない。
\end{enumerate}
\end{proposition}

\begin{proof}
定義~\ref{def:fr_domain} により、(B2) の条件は $\forall\eta\in\mathrm{Crit}^{(1)}(x_t):\eta(x_t,e_t)=1$ であり、
$\eta$ の値は $\{0,1\}$ に属するので、(B1) の条件はその否定である。
よって、$A_t=\emptyset$ の時刻ではちょうど一方が成り立つ。
(B1) の条件から $A_t=\emptyset$ が従うことは命題~\ref{prop:eu_blockage} (i) による。

$t\in\mathbb{T}_{\mathrm{P},i}$ とする。
(i)：定義~\ref{def:EU} により、$\mathrm{EU}_t$ が偽であることは (B1) の条件と同じであり、
そのとき $A_t=\emptyset$ である。
命題~\ref{prop:fr_domain_eu} により $(x_t,e_t)\notin\mathcal{D}_{\mathrm{FR}}$ であり、
定義~\ref{def:FR} により $\mathrm{FR}_t$ は定義されない。
(ii)：$\mathrm{EU}_t$ が真であることは $(x_t,e_t)\in\mathcal{D}_{\mathrm{FR}}$ と同値であり（命題~\ref{prop:fr_domain_eu}）、
そのもとで $\mathrm{FR}_t$ が偽であることは $A_t=\emptyset$ と同値である（定義~\ref{def:FR}）。
(iii)：$A_t=\emptyset$ ならば、前半の二分と (i)(ii) により、いずれか一方が成り立つ。
逆に、$\mathrm{EU}_t$ が偽ならば命題~\ref{prop:eu_blockage} (ii) により、
$\mathrm{EU}_t$ が真で $\mathrm{FR}_t$ が偽ならば定義~\ref{def:FR} により、$A_t=\emptyset$ である。
二つが同時に成り立たないことは、$\mathrm{EU}_t$ の真偽が異なることによる。
\end{proof}

\begin{remark}[二分の地位]
\label{rem:blockage_dichotomy}
命題~\ref{prop:blockage_dichotomy} は定義から従う診断上の分割であり、閉塞の原因の分類ではない。
(B1)・(B2) の名称は時刻別充足述語を用いずに定めており、
$t\notin\mathbb{T}_{\mathrm{P},i}$ の時刻の閉塞にも適用される。
時刻別充足述語との対応は、$t\in\mathbb{T}_{\mathrm{P},i}$ の時刻に限って述べる。
(B1) の時刻には $\mathrm{FR}_t$ は偽ではなく定義されないのであって、
定義されないことを偽に含めない（備考~\ref{rem:s3_description_levels}）。
(B1) の閉塞において、要因 $\eta$ の非充足が閉塞の唯一の原因であるとは主張しない。
$x_t\notin K$ の時刻の閉塞は一段階では $K$ へ戻れないことを表し、
数段階を経た $K$ への回復の可能性については何も述べない（備考~\ref{rem:eu_one_step_scope}）。
現在の状態でどの外部環境状態によっても充足されない要因は臨界必要要因に含まれないので、
そのような要因に関わる閉塞は (B1) には数えられない（備考~\ref{rem:eu_env_utilization}）。

(B2) は、閉塞のうち (B1) の条件を満たさないものの診断上の残余であり、
応答の側の失敗を表す原因の分類ではない。
たとえば、複数の必要環境要因がそろって充足されないことによってのみ一段階の閉塞が生じ、
各要因が単独では臨界必要要因に属さない場合、その閉塞は (B2) に属する。
また、$x_t\notin K$ で、どの外部環境状態のもとでも一段では $K$ に戻れない場合、
定義~\ref{def:eu_critical_domain} の (iii) を満たす要因はなく臨界必要要因集合は空となるので、
その時刻の閉塞は (B2) に属する。
この二分は、固定した要因の分解 $E$ と一段の判定に相対する診断上の分割であり、
環境の側と応答の側の因果的な寄与を分類するものではない。
\end{remark}

時刻別充足述語の定義域と閉塞の二分の関係を図~\ref{fig:s3_predicate_domain} に示す。

\begin{figure}[tbp]
\centering
\begin{tikzpicture}[
  font=\footnotesize,
  box/.style={draw, rounded corners=2pt, align=center, inner sep=3pt},
  cond/.style={font=\scriptsize, inner sep=1pt},
  note/.style={font=\scriptsize, align=center, inner sep=1pt},
  plabel/.style={font=\small, anchor=west, inner sep=0pt},
  arr/.style={-{Latex[length=1.6mm]}},
]
\node[plabel] at (0,0.8) {(a)\ 時刻別充足述語の定義域};
\node[box] (t) at (5.0,0) {時刻 $t$};
\node[box] (out) at (1.9,-1.3)
 {時刻別充足述語を定めない\\（真でも偽でもない）};
\node[box] (in) at (8.0,-1.3) {$t\in\mathbb{T}_{\mathrm{P},i}$};
\draw[arr] (t) -| node[cond, pos=0.3, above]{$\mathrm{FT}_{i,t}=1$} (out);
\draw[arr] (t) -| node[cond, pos=0.3, above]{$\mathrm{FT}_{i,t}=0$} (in);

\node[box] (sm) at (4.6,-3.0)
 {$\mathrm{SM}_t$：$x_t\in K$ なら真\\$x_t\notin K$ なら偽};
\node[box] (eu) at (9.0,-3.0)
 {$\mathrm{EU}_t$：すべての $\eta\in\mathrm{Crit}^{(1)}(x_t)$ で\\$\eta(x_t,e_t)=1$ なら真\\そうでなければ偽\\（$\mathrm{Crit}^{(1)}(x_t)=\emptyset$ なら真）};
\coordinate (br) at (8.0,-1.8);
\draw (in) -- (br);
\draw[arr] (br) -| (sm.north);
\draw[arr] (br) -| (eu.north);
\node[note, anchor=north] at (sm.south) {SM の真偽は EU・FR の\\判定の条件ではない};

\node[box] (frd) at (7.3,-5.4)
 {$(x_t,e_t)\in\mathcal{D}_{\mathrm{FR}}$（命題~\ref{prop:fr_domain_eu}）\\$\mathrm{FR}_t$ を定義する\\$A_t\neq\emptyset$ なら真\\$A_t=\emptyset$ なら偽};
\node[box] (fru) at (10.9,-5.4)
 {$\mathrm{FR}_t$ は定義しない\\（偽ではない）\\このとき $A_t=\emptyset$\\（命題~\ref{prop:eu_blockage}）};
\draw[arr] (frd.north |- eu.south) -- node[cond, left]{真} (frd.north);
\draw[arr] (fru.north |- eu.south) -- node[cond, right]{偽} (fru.north);

\begin{scope}[yshift=-8.1cm]
\node[plabel] at (0,0.8) {(b)\ 閉塞の二分（$A_t=\emptyset$ のすべての時刻）};
\node[box] (bl) at (5.8,0) {$A_t=\emptyset$};
\node[box] (b1) at (3.0,-1.5)
 {(B1)：ある $\eta\in\mathrm{Crit}^{(1)}(x_t)$ で\\$\eta(x_t,e_t)=0$};
\node[box] (b2) at (8.6,-1.5)
 {(B2)：$(x_t,e_t)\in\mathcal{D}_{\mathrm{FR}}$};
\draw[arr] (bl) -| (b1);
\draw[arr] (bl) -| (b2);
\node[cond] at (5.8,-0.6) {ちょうど一方が成り立つ};
\node[box, dashed] (b1p) at (3.0,-3.0)
 {$t\in\mathbb{T}_{\mathrm{P},i}$ のとき\\$\mathrm{EU}_t$ が偽（$\mathrm{FR}_t$ は定義されない）};
\node[box, dashed] (b2p) at (8.6,-3.0)
 {$t\in\mathbb{T}_{\mathrm{P},i}$ のとき\\$\mathrm{EU}_t$ が真で $\mathrm{FR}_t$ が偽};
\draw (b1) -- (b1p);
\draw (b2) -- (b2p);
\node[note, align=left, anchor=north west] at (0.3,-3.75)
 {(B1) の条件は、どの時刻でも単独で $A_t=\emptyset$ を導く（命題~\ref{prop:eu_blockage}）。\\
  点線の枠は原理判定時刻域の中での読み替えである。\\
  その外の時刻にも (B1)・(B2) の二分は適用される。};
\end{scope}
\end{tikzpicture}
\caption{時刻別充足述語の定義域と閉塞の二分（命題~\ref{prop:blockage_dichotomy}）。
矢印は定義による場合分けを表し、因果や時間の順序を表さない。
(a) SM と EU は原理判定時刻域の上でそれぞれ判定され、SM から EU・FR への線はない。
FR は EU が真の時刻にだけ定義され、定義されないことを偽に含めない。
(b) 二分は定義から従う診断上の分割であり、閉塞の原因の分類ではない（備考~\ref{rem:blockage_dichotomy}）。}
\label{fig:s3_predicate_domain}
\end{figure}

\begin{remark}[三原理の時刻別充足述語と閉塞・整形式実現・崩壊]
\label{rem:principles_failure_table}
$t\in\mathbb{T}_{\mathrm{P},i}$ の時刻について、次が成り立つ。
\begin{enumerate}
  \item $\mathrm{SM}_t$ の偽は閉塞の定義条件ではない。
  本論文は、$\mathrm{SM}_t$ の偽だけから閉塞 $A_t=\emptyset$ を一般的な帰結として置かない。
  \item $\mathrm{EU}_t$ が偽ならば $A_t=\emptyset$ である（命題~\ref{prop:eu_blockage}）。
  \item $\mathrm{FR}_t$ が偽ならば $A_t=\emptyset$ であり、
  $\mathcal{D}_{\mathrm{FR}}$ 上で $\mathrm{FR}_t$ の偽と閉塞は一致する（定義~\ref{def:FR}）。
  \item 閉塞の時刻は、$\mathrm{EU}_t$ が偽である時刻と、
  $\mathrm{EU}_t$ が真で $\mathrm{FR}_t$ が偽である時刻に排他的・網羅的に分かれる
  （命題~\ref{prop:blockage_dichotomy}）。
  \item 本論文は、いずれの時刻別充足述語が偽であることだけからも、
  $\mathrm{WF}_{i,t}=0$ または崩壊 $\mathrm{Col}_{i,s}=1$（$s\ge t$）を一般的な帰結として置かない。
  \item 時刻別充足述語を判定する時刻 $t\in\mathbb{T}_{\mathrm{P},i}$ では、
  定義~\ref{def:principle_time_domain} により $\mathrm{FT}_{i,t}=0$ であり、
  補題~\ref{lem:principle_time_domain_down} により $\mathrm{Col}_{i,t}=0$ である。
\end{enumerate}
\end{remark}

%----------------------------------------------------------------------
\subsection{三原理の連動と含意関係の限定}
\label{subsec:principle_linkage}

本小節では、三原理の時刻別充足述語の間に本節で示した条件付き含意をまとめ、
本論文が主張しない含意を明示する。

\begin{remark}[三原理の連動]
\label{rem:principle_linkage}
三原理は相互を定義するものではなく、概念的に異なる制約軸である。
ただし、明示した前提または適用域のもとで、
それぞれの時刻別充足述語の間に条件付き含意が成立することを妨げない。
本節で示した条件付き含意等は、次のとおりである。
\begin{enumerate}
  \item 命題~\ref{prop:eu_sm_next}：EU の偽から次時刻の $x_{t+1}\notin K$。
  \item 命題~\ref{prop:fr_next_sm}：局所 FR の偽から次時刻の $x_{t+1}\notin K$。
  \item 命題~\ref{prop:viab_stage_fr}：有限段頑健可行集合から局所 FR の真。
  \item 命題~\ref{prop:fr_domain_eu}：FR の定義域と EU の真。
  \item 命題~\ref{prop:blockage_dichotomy}：閉塞の二分と EU・FR の時刻別充足述語との対応。
\end{enumerate}
\end{remark}

\begin{remark}[含意関係の限定]
\label{rem:principle_implication_scope}
本論文は、$t\in\mathbb{T}_{\mathrm{P},i}$ について、
$\mathrm{SM}_t$ の真から $\mathrm{EU}_t$ の真を、
$\mathrm{EU}_t$ の真から $\mathrm{SM}_t$ の真を、
$\mathrm{SM}_t$ と $\mathrm{EU}_t$ の真から $\mathrm{FR}_t$ の真を、
いずれも一般的な帰結として置かない。
これは含意を主張しないということであり、
これらの含意が成り立たないことを示したものではない。
含意が成り立たないことを示すには、
\S\ref{sec:acs_def} の仮定と構成要件をすべて満たす系を構成し、
前件が真で後件が偽となる時刻を示す必要がある。
その構成は本論文では行わない。
\end{remark}

%----------------------------------------------------------------------
\subsection{応答多様性縮退と長期可行性：条件付き必要性}
\label{subsec:sd_necessity}

\subsubsection{本小節の目的と射程}

本小節では、FR の長期の側面のうち、
必要環境充足状態に対する応答の多様性が縮退したまま持続する場合に、
長期可行性とどのように関係するかを検討する。
扱う量は応答多様度 $\mu_{\mathrm{eff},t}$（定義~\ref{def:mu_eff}、分類C）であり、
これは三原理のいずれの時刻別充足述語の定義にも用いない診断量である。

本小節で中心的に扱う縮退モードは、
\[
  e_t\in E_{\mathrm{req}}(x_t),
  \qquad
  A_t\neq\emptyset,
  \qquad
  \mu_{\mathrm{eff},t}=1
\]
である。
これは、現在可行な行動が存在する一方で、
$E_{\mathrm{req}}(x_t)$ 上の可行性保持パターンが
単一の同値類へ縮退している場合を表す。
このとき、命題~\ref{prop:viab_stage_fr} の証明と同じ論法により
$(x_t,e_t)\in\mathcal{D}_{\mathrm{FR}}$ であり、
$t\in\mathbb{T}_{\mathrm{P},i}$ の時刻では $\mathrm{FR}_t$ は真である。
すなわち、本小節が扱う縮退は局所 FR が真である時刻にも生じうる。

以下で形式的に示すのは、
ある時刻の $\mu_{\mathrm{eff},t}=1$ から閉塞を導く無条件命題ではない。
示すのは次の二段である。
\begin{enumerate}
  \item 主命題（強仮定を用いない）：
  実現外部環境状態が必要環境要因を充足し、可行な状態で応答多様性の縮退と応答の弱点が持続するとき、
  頑健応答不能が持続し、軌道は可行性核の外に留まる（命題~\ref{prop:sd_robust_nonresponse}）。
  \item 系（強仮定のもとで）：
  可行性核の外への持続を閉塞へ接続する強仮定を加えると、
  有限時間で閉塞へ至る（命題~\ref{prop:sd_necessity}）。
\end{enumerate}
いずれも条件付き必要性の方向の主張である。
また、閉塞 $A_t=\emptyset$ から崩壊を導くこともしない。
閉塞と崩壊の関係は \S\ref{sec:collapse_cond} の責務である。

以下では、\S\ref{sec:acs_def} の必要環境充足集合 $E_{\mathrm{req}}(x)$ と、
定義~\ref{def:mu_eff} の $\varphi_u$、$\sim$、$\mu_{\mathrm{eff}}(x,e)$、
$\mu_{\mathrm{eff},t}$ を上流の定義として固定し、本小節で変更しない。
証明中で時刻 $t$ の同値関係であることを明示する場合は $\sim_t$ と書く。
本小節は、解析者が外部から選択する環境部分集合を導入しない
（備考~\ref{rem:env_basis_intrinsic}）。
環境に関する量化は、$E_{\mathrm{ext}}$ または
$E_{\mathrm{req}}(x)$ によって与える。

\subsubsection{本小節で用いる仮定と応答不能の記述}

\begin{assumption}[非定常環境]
\label{ass:nonstationary}
実現外部環境状態の列 $(e_t)_{t\ge0}$ は時間的に一定ではない。
すなわち、少なくとも二つの時刻 $t\neq t'$ が存在して
\[
  e_t\neq e_{t'}
\]
が成立する。ここで $e_t,e_{t'}\in E_{\mathrm{ext}}$ である。

導入地位：O（観察的）。強さ：弱。

\textit{注：}
本仮定は確率的分散または統計的変動を仮定するものではなく、
Phase~1 の決定論的・離散時間枠組みにおいて、
外部環境状態の列が時間的に一定でないことのみを意味する。
その変化は、外生背景成分の変化によるものでも、
当該ACS自身の行動に応答した行動応答成分の変化
（\S\ref{sec:acs_def} 定義~\ref{def:env_realization}）によるものでもよい。

本仮定は、
環境応答多様性を保持することの構造的意義を解釈する
背景条件として用いるが、
後述する可行な応答の弱点（定義~\ref{def:nonresponse_env_set}）の存在を
本仮定から導くことは主張しない。
また、命題~\ref{prop:sd_robust_nonresponse} と命題~\ref{prop:sd_necessity} の
形式的証明には使用しない。
\end{assumption}

\begin{definition}[応答不能環境集合と可行な応答の弱点（分類C）]
\label{def:nonresponse_env_set}
状態 $x\in X$ と行動 $u\in U$ について、
\[
  N(x,u)
  :=
  \bigl\{\,e\in E_{\mathrm{req}}(x)\bigm| f(x,u,e)\notin K\,\bigr\}
  \subseteq E_{\mathrm{req}}(x)
\]
を、状態 $x$ における行動 $u$ の\emph{応答不能環境集合}と呼ぶ。
$N(x,u)\neq\emptyset$ は、状態 $x$ で行動 $u$ をとるとき、
必要環境要因がすべて充足される外部環境状態のうちに、次状態を $K$ に置かないものがあることを表す。

時刻 $t$ において
\[
  \exists\,u^\star\in A_t:\quad N(x_t,u^\star)\neq\emptyset
\]
が成り立つとき、時刻 $t$ に\emph{可行な応答の弱点}があるという。
\end{definition}

\begin{remark}[応答不能は状態・行動・環境の関係である]
\label{rem:nonresponse_relational}
$N(x,u)$ は状態 $x$、行動 $u$、外部環境状態 $e$ の三つ組の関係として定まり、
外部環境状態だけの性質ではない。
本論文は、存続できない局面を外部環境の側だけに帰属させず、
状態がある外部環境状態に応答できないという関係として記述する。

可行な応答の弱点は、ACS の構成要件でも、すべての時刻に課す仮定でもない。
本小節では、\S\ref{subsubsec:sd_main} の命題の前提としてだけ用いる。
弱点がない時刻も存在しうる。
例えば $E_{\mathrm{req}}(x_t)=\{e_t\}$ であれば、
任意の $u\in A_t$ について $f(x_t,u,e_t)\in K$ であるから、
時刻 $t$ に可行な応答の弱点はない。
本論文は、弱点がない状態の存在を排除しない。

弱点の証人 $e^\star\in N(x_t,u^\star)$ は $E_{\mathrm{ext}}$ 全体ではなく
$E_{\mathrm{req}}(x_t)$ に属する。
これは、定義~\ref{def:viab_core} の可行性核が $E_{\mathrm{req}}(x)$ 上の量化で定義されることとの整合による。
\end{remark}

\begin{remark}[可行な応答の弱点が含意する環境差]
\label{rem:weakness_env_difference}
時刻 $t$ に可行な応答の弱点があり、その証人を
$u^\star\in A_t$、$e^\star\in N(x_t,u^\star)$ とする。
このとき
\[
  f(x_t,u^\star,e_t)\in K,
  \qquad
  f(x_t,u^\star,e^\star)\notin K
\]
であるから、$e^\star\neq e_t$ であり、
\[
  u^\star\in A(x_t,e_t)\setminus A(x_t,e^\star),
  \qquad
  A(x_t,e_t)\neq A(x_t,e^\star)
\]
が成立する。
したがって、弱点がある局面では、可行行動対応 $A(x,e)$ が外部環境状態の差異に対して
少なくとも一箇所で非退化である。
ただし、この帰結から
構造的環境結合条件（定義~\ref{def:env_coupling}）や
特定の環境応答座標集合 $I_E$ の選択を直接導くことはしない。
\end{remark}

\begin{definition}[応答不能領域（分類C）]
\label{def:nonresponse_region}
\[
  \mathcal N^{(1)}
  :=
  \bigl\{\,(x,e)\in X\times E_{\mathrm{ext}}\bigm| A(x,e)=\emptyset\,\bigr\}
\]
を\emph{応答不能領域}と呼ぶ。
$(x,e)\in\mathcal N^{(1)}$ は $e\notin\mathcal E_{\mathrm{viab}}^{(1)}(x)$
（\S\ref{sec:acs_def} 定義~\ref{def:one_step_viable_environment}）と同値であり、
本定義は同対応の補集合を組として並べたものである。
定義から、時刻 $t$ の閉塞 $A_t=\emptyset$ は $(x_t,e_t)\in\mathcal N^{(1)}$ と同値である。
\end{definition}

\begin{definition}[頑健応答集合と頑健応答不能（分類C）]
\label{def:robust_response}
状態 $x\in X$ について、
\[
  U^{\mathrm{rob}}(x)
  :=
  \bigl\{\,u\in U\bigm| N(x,u)=\emptyset\,\bigr\}
\]
を\emph{頑健応答集合}と呼ぶ。
$U^{\mathrm{rob}}(x)=\emptyset$ のとき、状態 $x$ は\emph{頑健応答不能}であるという。
すなわち頑健応答不能とは、どの行動をとっても、
必要環境要因がすべて充足される外部環境状態のうちに次状態を $K$ に置かないものが残ることである。
\end{definition}

\begin{proposition}[頑健応答不能と頑健可行集合]
\label{prop:robust_nonresponse_viab}
$x\in X$ について $U^{\mathrm{rob}}(x)=\emptyset$ ならば、
$x\notin\mathrm{Viab}_f(K)$ であり、
すべての $\nu\ge1$ について $x\notin\mathrm{Viab}_f^{(\nu)}(K)$ である。
\end{proposition}

\begin{proof}
対偶を示す。
$x\in\mathrm{Viab}_f(K)$ ならば、命題~\ref{prop:viab_core_max} より、ある $u\in U$ が存在して
すべての $e\in E_{\mathrm{req}}(x)$ について $f(x,u,e)\in\mathrm{Viab}_f(K)\subseteq K$ であるから、
$N(x,u)=\emptyset$、すなわち $u\in U^{\mathrm{rob}}(x)$ である。
$\nu\ge1$ について $x\in\mathrm{Viab}_f^{(\nu)}(K)$ ならば、定義~\ref{def:viab_stage} より、ある $u\in U$ が存在して
すべての $e\in E_{\mathrm{req}}(x)$ について
$f(x,u,e)\in\mathrm{Viab}_f^{(\nu-1)}(K)\subseteq K_{\mathrm{req}}\subseteq K$ であるから、
同様に $u\in U^{\mathrm{rob}}(x)$ である。
\end{proof}

本論文は、命題~\ref{prop:robust_nonresponse_viab} の逆を一般的な帰結として置かない。
$E_{\mathrm{req}}(x)\neq\emptyset$ のとき、$U^{\mathrm{rob}}(x)\neq\emptyset$ は
$E_{\mathrm{req}}(x)$ のすべての外部環境状態のもとで次状態を $K$ に置く行動の存在だけを表し、
$K_{\mathrm{req}}$ や可行性核に留まることまでは表さない。
$E_{\mathrm{req}}(x)=\emptyset$ のときは、定義から $U^{\mathrm{rob}}(x)=U$ であり、
この非空性は頑健性について何も表さない。

\subsubsection{$K$ 外部持続と閉塞の接続}

\begin{definition}[$K$ 外部持続]
\label{def:theta_exterior}
以下、補集合は状態空間 $X$ における相対補集合とし、
\[
  K^c:=X\setminus K,
  \qquad
  \bigl(\mathrm{Viab}_f(K)\bigr)^c
  :=
  X\setminus\mathrm{Viab}_f(K)
\]
と書く。

ACS $\Sigma$ が時刻 $t$ 以降
\emph{$K$ 外部に持続する}とは、
\[
  \forall s\ge t,\;
  x_s\in K^c
\]
が成立することをいう。

\textit{注：}
本定義は、一時的な離脱後に $K$ へ再進入する
オーバーシュートを排除する。
$x_t\notin K$ となった直後に $x_{t+1}\in K$ に戻る場合は、
$K$ 外部に持続していないとして扱う。

ある有限時刻に $x_{t+T}\notin K$ となることから、
\[
  \forall s\ge0,\;
  x_{t+T+s}\notin K
\]
を一般的な帰結として置かない。
一度 $K$ を離脱した後の復帰は Phase~1 では排除されていない。
$K$ 外部への持続は、離脱とは独立に述べられる条件である。

有限時間幅 $\tau$ による持続の略記は、
$\tau$ の下限を指定しない限り一時的離脱と区別できないため
採用しない。
\end{definition}

\begin{assumption}[可行性核外部持続と閉塞の有限時間接続]
\label{ass:theta_exterior_blockage}
任意の時刻 $t\in\mathbb{N}_0$ について、
ACS $\Sigma$ の状態が時刻 $t$ 以降持続的に
$\mathrm{Viab}_f(K)$（定義~\ref{def:viab_core}）の外部に留まる、
すなわち
\[
  \forall s\ge t,\;
  x_s\notin\mathrm{Viab}_f(K)
\]
が成立するとき、
\[
  \exists\,T'\in\mathbb{N}_0
  \;\text{ s.t. }\;
  A_{t+T'}=\emptyset.
\]

導入地位：S（構造的）。強さ：強（強仮定）。

\textit{注（適用条件）：}
本仮定の前件は、可行性核外部への\emph{持続}であり、
一時的な外部化ではない。
また本仮定の帰結は閉塞 $A_{t+T'}=\emptyset$ であって、
崩壊ではない。
本論文では
\[
  A_t=\emptyset
  \;\Longrightarrow\;
  \mathrm{Col}_{i,t}=1
\]
とは述べない。閉塞と崩壊の正式な接続は
\S\ref{sec:collapse_cond} の責務である。

\textit{注（前件の包含関係）：}
定義~\ref{def:viab_core} より
$\mathrm{Viab}_f(K)\subseteq K$ であるから
$K^c\subseteq\bigl(\mathrm{Viab}_f(K)\bigr)^c$ である。
したがって $K$ 外部への持続
（定義~\ref{def:theta_exterior}）は本仮定の前件を含意する。
前件を $K$ 外部持続ではなく可行性核外部持続として置くのは、
$K$ への一時的な再進入を挟みながらも
ロバスト可行性を回復できない軌道を対象に含めるためである。

\textit{注：}
本仮定を ACS の内部動学から導出するための十分条件の形式化は
Phase~2 の課題とする。
Phase~1 の本小節では具体的経路を形式仮定として課さない
（within present framework）。

\textit{注（役割）：}
本仮定は、命題~\ref{prop:exterior_blockage}、系である命題~\ref{prop:sd_necessity} および命題~\ref{prop:sd_necessity_fr} にだけ用いる。
主命題（命題~\ref{prop:sd_robust_nonresponse}）は本仮定を用いない。
可行性核の外に持続しても閉塞に至らない軌道は型の上では排除されず
（備考~\ref{rem:exterior_without_blockage}）、
本仮定はそのような軌道を条件付き解析の対象から除く強仮定である。
\S\ref{sec:acs_def} の分解のもとでは、当該ACSの行動に応答する行動応答成分が保たれることにより、
可行性核の外に留まりながら可行行動を保つ軌道も、本仮定によって除かれる。
\end{assumption}

\begin{proposition}[$K$ 外部持続と閉塞]
\label{prop:exterior_blockage}
仮定~\ref{ass:theta_exterior_blockage} のもとで、
ACS $\Sigma$ が時刻 $t$ 以降 $K$ 外部に持続する
（定義~\ref{def:theta_exterior}）ならば、
\[
  \exists\,T'\in\mathbb{N}_0
  \;\text{ s.t. }\;
  A_{t+T'}=\emptyset.
\]
\end{proposition}

\begin{proof}
定義~\ref{def:viab_core} より
$\mathrm{Viab}_f(K)\subseteq K$、したがって
$K^c\subseteq\bigl(\mathrm{Viab}_f(K)\bigr)^c$ である。
$K$ 外部持続は可行性核外部持続を含意するので、
仮定~\ref{ass:theta_exterior_blockage} を適用して結論を得る。
\end{proof}

\begin{remark}[可行性核の外への持続と閉塞の関係]
\label{rem:exterior_without_blockage}
本論文は、仮定~\ref{ass:theta_exterior_blockage} を置かない解析において、
可行性核の外への持続から閉塞を一般的な帰結として置かない。
例えば命題~\ref{prop:robust_noninclusion} (ii) の例で、
外生背景成分がつねに $b_1$ であり $o_0=\mathrm{good}$ であれば、
$x_t=1\notin\mathrm{Viab}_f(K)$ がすべての時刻で成り立つ一方、
$A_t=\{u_{\emptyset}\}\neq\emptyset$ もすべての時刻で成り立つ。
この例も型の水準で構成したものであり、備考~\ref{rem:robust_noninclusion_reading} と同じ射程をもつ。
\end{remark}

\subsubsection{応答多様性縮退に関する条件付き決定論的証明}

\begin{lemma}[応答多様性縮退と応答の弱点による頑健応答不能]
\label{lem:sd_viab}
定義~\ref{def:mu_eff}、定義~\ref{def:viab_core}、
定義~\ref{def:nonresponse_env_set} および定義~\ref{def:robust_response} のもとで、
時刻 $t$ において
\[
  x_t\in K,
  \qquad
  e_t\in E_{\mathrm{req}}(x_t),
  \qquad
  A_t\neq\emptyset,
  \qquad
  \mu_{\mathrm{eff},t}=1
\]
が成立し、かつ時刻 $t$ に可行な応答の弱点があるならば、
\[
  U^{\mathrm{rob}}(x_t)=\emptyset,
  \qquad
  x_t\notin\mathrm{Viab}_f(K).
\]
\end{lemma}

\begin{proof}
$e_t\in E_{\mathrm{req}}(x_t)$ より
$E_{\mathrm{req}}(x_t)\neq\emptyset$、したがって
$x_t\in K_{\mathrm{req}}$ である。

時刻 $t$ に可行な応答の弱点があるので（定義~\ref{def:nonresponse_env_set}）、ある
\[
  u^\star\in A_t,
  \qquad
  e^\star\in N(x_t,u^\star)\subseteq E_{\mathrm{req}}(x_t)
\]
が存在して $f(x_t,u^\star,e^\star)\notin K$ が成立する。
定義~\ref{def:mu_eff} の可行性保持指示関数
$\varphi_u:E_{\mathrm{req}}(x_t)\to\{0,1\}$ を用いれば
$\varphi_{u^\star}(e^\star)=0$ である。

\textbf{ステップ1：応答多様性縮退による弱点の共有。}
\[
  \mu_{\mathrm{eff},t}
  =
  \bigl|A_t/\!\sim_t\bigr|
  =
  1
\]
より、$A_t$ 上の全ての行動は単一の同値類に属する。
$\sim_t$ は $E_{\mathrm{req}}(x_t)$ 上の関数としての
$\varphi_u$ の一致で定義されるから
\[
  \forall u\in A_t,
  \qquad
  \varphi_u=\varphi_{u^\star}.
\]
特に $e^\star\in E_{\mathrm{req}}(x_t)$ において
\[
  \forall u\in A_t,
  \qquad
  \varphi_u(e^\star)=0,
\]
すなわち
\[
  \forall u\in A_t,
  \qquad
  f(x_t,u,e^\star)\notin K.
\]

\textbf{ステップ2：$A_t$ 外の行動。}
定義~\ref{def:at} より
$A_t=\{u\in U\mid f(x_t,u,e_t)\in K\}$ であるから、
任意の $u\in U\setminus A_t$ について
$f(x_t,u,e_t)\notin K$ である。
ここで前提より $e_t\in E_{\mathrm{req}}(x_t)$ であるから、
$e_t$ は $E_{\mathrm{req}}(x_t)$ 内の証人として用いることができる。

\textbf{ステップ3：ロバスト可行性条件の否定。}
ステップ1より $u\in A_t$ は $e^\star\in E_{\mathrm{req}}(x_t)$ で、
ステップ2より $u\in U\setminus A_t$ は
$e_t\in E_{\mathrm{req}}(x_t)$ で、それぞれ可行性を失う。
したがって
\[
  \forall u\in U,\;
  \exists e\in E_{\mathrm{req}}(x_t)
  \;\text{ s.t. }\;
  f(x_t,u,e)\notin K.
\]
すなわち、すべての $u\in U$ について $N(x_t,u)\neq\emptyset$ であり、
$U^{\mathrm{rob}}(x_t)=\emptyset$ である（定義~\ref{def:robust_response}）。
これは
\[
  \exists u\in U
  \;\text{ s.t. }\;
  \forall e\in E_{\mathrm{req}}(x_t),\;
  f(x_t,u,e)\in K
  \tag{$\dagger$}
\]
の否定である。

一方、$x_t\in\mathrm{Viab}_f(K)$ であるならば、
命題~\ref{prop:viab_core_max} より $\mathrm{Viab}_f(K)$ は
$E_{\mathrm{req}}$ 不変であるから、ある $u\in U$ が存在して
\[
  \forall e\in E_{\mathrm{req}}(x_t),\;
  f(x_t,u,e)\in\mathrm{Viab}_f(K)\subseteq K
\]
が成立し、これは $(\dagger)$ を含意する。
$(\dagger)$ の否定を得ているから
$x_t\notin\mathrm{Viab}_f(K)$ である。
\end{proof}

\begin{remark}[補題~\ref{lem:sd_viab} の論理構造]
\label{rem:sd_viab_structure}
本補題では、
可行な応答の弱点、
$\mu_{\mathrm{eff},t}=1$、および
$e_t\in E_{\mathrm{req}}(x_t)$ が
それぞれ異なる役割を担う。

可行な応答の弱点は、
現在可行な行動のうち少なくとも一つ $u^\star\in A_t$ について
$N(x_t,u^\star)\neq\emptyset$ であることだけを前提とする。

$\mu_{\mathrm{eff},t}=1$ は、$A_t$ 上の全可行行動が
$E_{\mathrm{req}}(x_t)$ 上で同一の可行性保持パターンをもつことを
意味し、$u^\star$ の弱点を $A_t$ 全体へ伝播させる。
この伝播がなければ、$A_t$ の他の行動が
$E_{\mathrm{req}}(x_t)$ 上でロバストである可能性を排除できない。
したがって $\mu_{\mathrm{eff},t}=1$ は、本補題では状態分類のラベルではなく、結論の導出に用いる前提である。
この前提を除くと結論が成り立たない例がある。
$X=\{0,1\}$、$\Pi(x)=x-1$（$K=\{1\}$）、$U=\{u_1,u_2\}$、$E_{\mathrm{ext}}=\{e^1,e^2\}$、$E=\{\eta\}$、$\eta\equiv1$ とし、
$f(1,u_1,e^1)=f(1,u_2,e^1)=f(1,u_2,e^2)=1$、それ以外の引数では $f=0$ とする。
$x_t=1$、$e_t=e^1$ では $A_t=\{u_1,u_2\}$ であり、$u_1$ に弱点 $e^2$ があるが、$\mu_{\mathrm{eff},t}=2$ である。
このとき $u_2\in U^{\mathrm{rob}}(1)$ であり、$\{1\}$ は $E_{\mathrm{req}}$ 不変であるから $1\in\mathrm{Viab}_f(K)$ である
（型の水準の例）。

$e_t\in E_{\mathrm{req}}(x_t)$ は、$A_t$ の外側にある行動を
$E_{\mathrm{req}}(x_t)$ 内の証人によって排除するために用いる。
この条件がない場合、$E_{\mathrm{req}}(x_t)$ 上でロバストな行動が
$A_t$ の外に存在しうる。
$\mu_{\mathrm{eff},t}$ は $A_t$ 上の分割の濃度であるから、
そのような行動を評価できない。
したがって本条件は、現在の可行行動集合 $A_t$ と
ロバスト候補行動とを接続する役割をもつ。
この前提を除くと結論が成り立たない例がある。
前の例で $E_{\mathrm{ext}}=\{e^1,e^2,e^3\}$、$\eta(x,e)=1\iff e\neq e^1$ とし、
$f(1,u_1,e^1)=f(1,u_1,e^2)=f(1,u_2,e^2)=f(1,u_2,e^3)=1$、それ以外の引数では $f=0$ とする。
$x_t=1$、$e_t=e^1\notin E_{\mathrm{req}}(1)$ では $A_t=\{u_1\}$、$\mu_{\mathrm{eff},t}=1$ であり、$u_1$ に弱点 $e^3$ があるが、
$u_2\in U^{\mathrm{rob}}(1)$ であり、$1\in\mathrm{Viab}_f(K)$ である（型の水準の例）。

なお、より強い条件、すなわち頑健応答不能 $U^{\mathrm{rob}}(x_t)=\emptyset$ を前提に置けば、
$\mu_{\mathrm{eff},t}=1$ を用いずに
命題~\ref{prop:robust_nonresponse_viab} から同じ非包含が得られる。
本小節はこの強い条件を前提とせず、
応答の弱点と $\mu_{\mathrm{eff},t}=1$ から頑健応答不能を導く。
\end{remark}

\begin{remark}[$\mu_{\mathrm{eff},t}=1$ の構造的意味：傾向的解釈]
\label{rem:mu_eff_tendency}
$\mu_{\mathrm{eff},t}=1$ は、
現在の可行行動集合 $A_t$ 上で
$E_{\mathrm{req}}(x_t)$ に対する応答パターンが
単一化している状態を表す。

ただし、$\mu_{\mathrm{eff},t}=1$ のみから
構造的脆弱性が一般に導かれるわけではない。
共通の応答パターンが
$\varphi(e)\equiv1$（$e\in E_{\mathrm{req}}(x_t)$）
であるならば、単一の応答パターンであっても
時刻 $t$ の状態 $x_t$ における一段階先については、
$E_{\mathrm{req}}(x_t)$ 全体に対して可行性を保つ。
ただし、この局所的一段階性質は
備考~\ref{rem:sd_local_scope} の時間的射程に従い、
本論文は、この性質から、後続時刻でも同じ応答パターンのまま長期存続可能であることを一般的な帰結として置かない。
なお、時刻 $t$ に可行な応答の弱点がある場合には、
$\varphi(e)\equiv1$ の場合は生じない。
上記は、弱点がない時刻における境界事例を記述する。

これに対し時刻 $t$ に可行な応答の弱点がある場合には、共通応答パターンは
少なくとも一つの $e^\star\in E_{\mathrm{req}}(x_t)$ に対して
$\varphi(e^\star)=0$ を含む。
このとき $\mu_{\mathrm{eff},t}=1$ は、
その弱点を回避する代替的な可行行動が
$A_t$ 内に存在しないことを意味する。

外部環境状態の列が時間的に一定でないこと
（仮定~\ref{ass:nonstationary}）は、
応答多様性の構造的重要性を理解する背景を与える。
ただし本論文は、非定常性のみから、弱点をもつ外部環境状態 $e^\star$ が
実現する外部環境状態列上で実際に出現することを一般的な帰結として置かない。

また、$E_{\mathrm{req}}(x_t)=\emptyset$ に起因する
$\mu_{\mathrm{eff},t}=1$ は本備考の対象ではない
（備考~\ref{rem:mu_eff_empty_req}）。
本補題および本備考は $e_t\in E_{\mathrm{req}}(x_t)$ を
前提としており、この場合 $E_{\mathrm{req}}(x_t)\neq\emptyset$ である。

したがって本論文は、$\mu_{\mathrm{eff},t}=1$ 単独から
閉塞または崩壊を一般的な帰結として置かない。本備考は、
可行な応答の弱点と組み合わされた場合の
構造的脆弱性を記述する傾向的解釈である。
\end{remark}

\begin{remark}[内部状態差と本小節の時間的射程]
\label{rem:sd_internal_state_scope}
定義~\ref{def:mu_eff} の $\mu_{\mathrm{eff},t}$ は、
固定された現在状態 $x_t$ における局所量である。
したがって本補題が直接扱うのは、
その時点の内部状態に条件づけられた
$E_{\mathrm{req}}(x_t)$ 上の応答多様性である。

一方、\S\ref{subsec:self_env_principles} で概念的に述べた
時間とともに生じうる内部状態の差異は、
時刻 $s$ ごとに $x_s$ が変化し、それに伴って
$E_{\mathrm{req}}(x_s)$、$A_s$、$\mu_{\mathrm{eff},s}$ が
再評価されることを通じて本小節の時間的議論へ接続する。

ただし、本小節は老化、損耗、内部資源減少等の
個別の内部変化機構を新たに定式化しない。
それらの時刻間更新は \S\ref{sec:extended_dynamics} の責務であり、
本小節は各時刻で与えられた $x_s$ に相対して
応答多様性縮退の帰結を評価する。
\end{remark}

\begin{lemma}[可行性核外部におけるロバスト可行性の喪失]
\label{lem:finite_exit}
$x_t\in K_{\mathrm{req}}$ かつ
$x_t\notin\mathrm{Viab}_f(K)$ とする。
このとき、任意の状態フィードバック方策
$\pi:K_{\mathrm{req}}\to U$ に対して、
ある $T\in\mathbb{N}$（$T\ge1$）と有限列
$(e_s)_{s=t}^{t+T-1}$ が存在して、
\[
  e_s\in E_{\mathrm{req}}(x_s),
  \qquad
  x_{s+1}=f(x_s,\pi(x_s),e_s)
  \qquad
  (s=t,\ldots,t+T-1)
\]
のもとで
\[
  x_{t+T}\notin K_{\mathrm{req}}
\]
となる。

ここで列は離脱時刻までの有限 prefix であり、
$x_{t+T}\notin K_{\mathrm{req}}$ において $\pi$ は評価されない。
前提 $x_t\in K_{\mathrm{req}}$ より $T=0$ は結論を与えないため、
$T\ge1$ とする。
\end{lemma}

\begin{proof}
背理法による。
ある状態フィードバック方策 $\pi$ が存在して、
任意の $T\in\mathbb{N}$ および上記の条件を満たす任意の有限列
$(e_s)_{s=t}^{t+T-1}$ について
$x_{t+T}\in K_{\mathrm{req}}$ が成立すると仮定する。

$R$ を、$x_t$ から $\pi$ と上記の有限列によって
到達しうる状態全体の集合とする
（$T=0$ の場合として $x_t\in R$）。
背理法の仮定より $R\subseteq K_{\mathrm{req}}$ であるから、
$R$ の各元において $\pi$ は定義される。

任意の $y\in R$ と任意の $e\in E_{\mathrm{req}}(y)$ をとる。
$y$ に至る有限列の末尾に $e$ を継ぎ足した列も
同じ条件を満たすから、$f(y,\pi(y),e)\in R$ である。
したがって $u=\pi(y)$ ととれば
\[
  \forall y\in R,\;
  \exists u\in U
  \;\text{ s.t. }\;
  \forall e\in E_{\mathrm{req}}(y),\;
  f(y,u,e)\in R
\]
が成立する。
すなわち $R$ は $K_{\mathrm{req}}$ の $E_{\mathrm{req}}$ 不変な
部分集合であり、$R\in\mathfrak{K}$ である。

命題~\ref{prop:viab_core_max} より
$R\subseteq\mathrm{Viab}_f(K)$、
したがって $x_t\in\mathrm{Viab}_f(K)$ となり、
前提に矛盾する。
\end{proof}

\begin{remark}[補題~\ref{lem:finite_exit} の射程]
\label{rem:uniform_bound}
本補題は、各状態フィードバック方策に対して
離脱を生じさせる有限の外部環境状態列と時刻 $T\in\mathbb{N}$ が
\emph{存在すること}のみを主張する。

本補題は、実際に実現した外部環境状態列のもとで
軌道が $K_{\mathrm{req}}$ を離脱することを主張しない。
可行性核の外部にあることから従うのは、
任意の状態フィードバック方策に対して、
離脱を生じさせる外部環境状態列が少なくとも一つ存在するという
ロバスト可行性の喪失に限られる。
本補題は $E_{\mathrm{req}}(x_s)$ の任意の元を許して量化しており、
外生背景成分の実現列と行動応答成分の遷移写像 $g_O$
（\S\ref{sec:acs_def} 定義~\ref{def:env_realization}）に従う実現軌道を対象としないため、
実現軌道の有限時間離脱を主張しない。

また、本補題は状態フィードバック方策を対象とする。
過去の履歴に依存する方策への一般化、
および $T$ が初期状態と方策にわたって一様に有界であること
$\sup T<\infty$
を、本補題は主張しない。
これらには追加の構造的条件が必要であり、Phase~2 の課題とする。

なお、補題~\ref{lem:finite_exit} は
命題~\ref{prop:sd_robust_nonresponse} と命題~\ref{prop:sd_necessity} の形式的証明には使用しない
（within present framework）。
\end{remark}

\subsubsection{主命題と系：頑健応答不能の持続と閉塞への接続}
\label{subsubsec:sd_main}

本小節の結果は、強仮定を用いない主命題と、
強仮定~\ref{ass:theta_exterior_blockage} を加えた系に分けて述べる。

\begin{proposition}[主命題：応答多様性縮退と応答の弱点の持続による頑健応答不能の持続]
\label{prop:sd_robust_nonresponse}
ある時刻 $t$ 以降の各 $s\ge t$ について
\begin{enumerate}[label=(\alph*)]
  \item $x_s\in K$ ならば $e_s\in E_{\mathrm{req}}(x_s)$、
  \item $x_s\in K$ かつ $A_s\neq\emptyset$ ならば、
  $\mu_{\mathrm{eff},s}=1$ であり、かつ時刻 $s$ に可行な応答の弱点がある
\end{enumerate}
が成立するとする。このとき、各 $s\ge t$ について
\[
  x_s\notin\mathrm{Viab}_f(K)
\]
であり、さらに $x_s\in K$ である時刻 $s$ では $U^{\mathrm{rob}}(x_s)=\emptyset$ である。
本命題は仮定~\ref{ass:theta_exterior_blockage} を用いない。
\end{proposition}

\begin{proof}
任意の $s\ge t$ をとる。

$x_s\notin K$ の場合、定義~\ref{def:viab_core} より
$\mathrm{Viab}_f(K)\subseteq K$ であるから
$x_s\notin\mathrm{Viab}_f(K)$ である。

$x_s\in K$ かつ $A_s\neq\emptyset$ の場合、条件 (a) より
$e_s\in E_{\mathrm{req}}(x_s)$ であり、条件 (b) より
$\mu_{\mathrm{eff},s}=1$ で、時刻 $s$ に可行な応答の弱点がある。
したがって補題~\ref{lem:sd_viab} より
$U^{\mathrm{rob}}(x_s)=\emptyset$ かつ $x_s\notin\mathrm{Viab}_f(K)$ である。

$x_s\in K$ かつ $A_s=\emptyset$ の場合、条件 (a) より
$e_s\in E_{\mathrm{req}}(x_s)$ である。
定義~\ref{def:at} より、すべての $u\in U$ について $f(x_s,u,e_s)\notin K$ であるから、
$e_s\in N(x_s,u)$ であり、$U^{\mathrm{rob}}(x_s)=\emptyset$ である。
命題~\ref{prop:robust_nonresponse_viab} より $x_s\notin\mathrm{Viab}_f(K)$ である。
\end{proof}

\begin{proposition}[系：応答多様性縮退の持続からの有限時間閉塞；強仮定のもとで]
\label{prop:sd_necessity}
仮定~\ref{ass:theta_exterior_blockage} のもとで、
ある時刻 $t$ 以降の各 $s\ge t$ について
命題~\ref{prop:sd_robust_nonresponse} の条件 (a)(b) が成立するとする。このとき、
\[
  \exists\,T\in\mathbb{N}_0
  \;\text{ s.t. }\;
  A_{t+T}=\emptyset.
\]
\end{proposition}

\begin{proof}
命題~\ref{prop:sd_robust_nonresponse} より
$\forall s\ge t,\ x_s\notin\mathrm{Viab}_f(K)$ である。
これは仮定~\ref{ass:theta_exterior_blockage} の前件そのものであるから、
同仮定より、ある $T\in\mathbb{N}_0$ について $A_{t+T}=\emptyset$ である。
\end{proof}

\begin{remark}[主命題と系の役割分担]
\label{rem:sd_role_division}
主命題（命題~\ref{prop:sd_robust_nonresponse}）と系（命題~\ref{prop:sd_necessity}）は、依拠する前提が異なる。
\begin{enumerate}
  \item 主命題は、定義・補題~\ref{lem:sd_viab}・命題~\ref{prop:robust_nonresponse_viab} だけから従い、強仮定を用いない。
  結論は、可行性核の外への持続と、可行な時刻における頑健応答不能である。
  閉塞は結論に含まない。
  \item 系は、主命題の結論を強仮定~\ref{ass:theta_exterior_blockage} によって閉塞へ接続する。
  有限時間の閉塞という結論は、この強仮定に依存する。
  \item 強仮定を置かない解析では、本論文の主張は主命題の結論までである。
  主命題の結論が成り立っていても、閉塞は一般には従わない（備考~\ref{rem:exterior_without_blockage}）。
\end{enumerate}
\end{remark}

\begin{proposition}[系：応答多様性縮退の持続と局所 FR]
\label{prop:sd_necessity_fr}
仮定~\ref{ass:theta_exterior_blockage} と命題~\ref{prop:sd_robust_nonresponse} の条件 (a)(b) のもとで、
ある $T\in\mathbb{N}_0$ が存在して $A_{t+T}=\emptyset$ であり、さらに次のいずれかが成り立つ。
\begin{enumerate}[label=(\roman*)]
  \item $x_{t+T}\notin K$ であり、時刻 $t+T$ の閉塞は (B1) または (B2) である
  （命題~\ref{prop:blockage_dichotomy}）。
  このとき $t+T\in\mathbb{T}_{\mathrm{P},i}$ ならば、$\mathrm{SM}_{t+T}$ は偽であり、
  かつ「$\mathrm{EU}_{t+T}$ が偽」か「$\mathrm{EU}_{t+T}$ が真で $\mathrm{FR}_{t+T}$ が偽」のいずれか一方が成り立つ。
  \item $x_{t+T}\in K$ であり、時刻 $t+T$ の閉塞は (B2) である。
  このとき $t+T\in\mathbb{T}_{\mathrm{P},i}$ ならば、$\mathrm{EU}_{t+T}$ は真であり、$\mathrm{FR}_{t+T}$ は偽である。
\end{enumerate}
\end{proposition}

\begin{proof}
命題~\ref{prop:sd_necessity} により、$A_{t+T}=\emptyset$ となる $T\in\mathbb{N}_0$ が存在する。
$x_{t+T}\notin K$ ならば、命題~\ref{prop:blockage_dichotomy} により (B1)・(B2) のちょうど一方が成り立つ。
後半は定義~\ref{def:SM} と同命題 (iii) による。
$x_{t+T}\in K$ ならば、前提 (a) により $e_{t+T}\in E_{\mathrm{req}}(x_{t+T})$ である。
命題~\ref{prop:viab_stage_fr} の証明と同じく、
定義~\ref{def:E_req} と $\mathrm{Crit}^{(1)}(x_{t+T})\subseteq E$ により
$(x_{t+T},e_{t+T})\in\mathcal{D}_{\mathrm{FR}}$ であり、閉塞は (B2) である。
後半は命題~\ref{prop:blockage_dichotomy} (ii) による。
\end{proof}

命題~\ref{prop:sd_necessity_fr} は、命題~\ref{prop:sd_necessity} と同じく、
強仮定~\ref{ass:theta_exterior_blockage} に依存する系である。
前提 (a) のもとでは、$s\ge t$ かつ $x_s\in K$ の時刻に (B1) の閉塞は生じず、
そのような時刻が $\mathbb{T}_{\mathrm{P},i}$ に属するならば $\mathrm{EU}_s$ は真である。
(i) の場合の閉塞は、一段階では $K$ へ戻れないことを表し、
数段階を経た $K$ への回復の可能性については何も述べない（備考~\ref{rem:eu_one_step_scope}）。

\begin{remark}[条件 (a) の役割と強さ]
\label{rem:req_condition_role}
命題~\ref{prop:sd_robust_nonresponse} の条件 (a) は、
可行状態にある各時刻において
実現外部環境状態が必要環境要因を充足していることを要求する。

これは補題~\ref{lem:sd_viab} を各時刻へ適用するために要する条件で
あり、本論文はこれを \S\ref{sec:acs_def} の構成要件 CR-1 から導出していない。
CR-1 の (CR-1b) は
$E_{\mathrm{req}}(x)\neq\emptyset$ なる $x$ の存在を要求するのみで、
実現外部環境状態が各時刻に $E_{\mathrm{req}}(x_s)$ へ属することを
要求していない。

したがって条件 (a) は命題~\ref{prop:sd_robust_nonresponse} の前件に属する限定条件であり、
本論文はこれを一般に成立する性質としては扱わない。
条件 (a) が成立しない時刻を含む軌道に対して、
命題~\ref{prop:sd_robust_nonresponse}、命題~\ref{prop:sd_necessity} および命題~\ref{prop:sd_necessity_fr} は何も述べない。
\end{remark}

\begin{remark}[本命題が扱う縮退モード]
\label{rem:sd_failure_mode_scope}
定義~\ref{def:mu_eff} により、
$A_t=\emptyset$ ならば $\mu_{\mathrm{eff},t}=0$ であり
（命題~\ref{prop:fr_mu_eff}）、これはすでに閉塞である。
また $E_{\mathrm{req}}(x_t)=\emptyset$ かつ $A_t\neq\emptyset$ ならば
$\mu_{\mathrm{eff},t}=1$ であるが、
これは比較の対象となる必要環境充足状態がないことの集合論上の帰結である
（備考~\ref{rem:mu_eff_empty_req}）。
本小節の中核補題、主命題（命題~\ref{prop:sd_robust_nonresponse}）および系（命題~\ref{prop:sd_necessity}）が扱うのは、
$e_t\in E_{\mathrm{req}}(x_t)$ かつ $A_t\neq\emptyset$ のもとで
$\mu_{\mathrm{eff},t}=1$ となる応答多様性縮退モードである。

したがって、本命題の結論を
$\mu_{\mathrm{eff},t}\le1$ となるすべての場合へ拡張しない。
また、局所 FR の偽は $\mathcal{D}_{\mathrm{FR}}$ 上の閉塞そのものであり
（定義~\ref{def:FR}）、
本命題は局所 FR の偽から何かを導くものではない。
\end{remark}

\subsubsection{核に留まる応答と結果の差}

主命題は、応答できない状態の持続を述べる。
これと対になるものとして、可行性核に留まる応答、あるいは可行性核に入る応答を
選べる場合の帰結を、仮定を用いずに述べる。

\begin{definition}[核維持応答集合（分類C）]
\label{def:kernel_response}
状態 $x\in X$ について、
\[
  U^{\mathrm{Viab}}(x)
  :=
  \bigl\{\,u\in U\bigm|\forall e\in E_{\mathrm{req}}(x):\ f(x,u,e)\in\mathrm{Viab}_f(K)\,\bigr\}
\]
を\emph{核維持応答集合}と呼ぶ。
$x\in\mathrm{Viab}_f(K)$ のとき、その元は可行性核に留まる応答であり、
$x\notin\mathrm{Viab}_f(K)$ のとき、その元は必要充足環境のいずれのもとでも可行性核に入る応答である。
\end{definition}

$x\in\mathrm{Viab}_f(K)$ ならば、命題~\ref{prop:viab_core_max} より $U^{\mathrm{Viab}}(x)\neq\emptyset$ である。
また $\mathrm{Viab}_f(K)\subseteq K$ であるから、$U^{\mathrm{Viab}}(x)\subseteq U^{\mathrm{rob}}(x)$ である。
$E_{\mathrm{req}}(x)=\emptyset$ のときは、定義から $U^{\mathrm{Viab}}(x)=U$ である。

\begin{proposition}[核に留まる応答を選び続ける軌道]
\label{prop:kernel_maintenance}
ある時刻 $T$ について $x_T\in\mathrm{Viab}_f(K)$ であり、
すべての $s\ge T$ について
\[
  e_s\in E_{\mathrm{req}}(x_s),
  \qquad
  u_s\in U^{\mathrm{Viab}}(x_s)
\]
が成立するとする。このとき、すべての $s\ge T$ について
\[
  x_s\in\mathrm{Viab}_f(K),
  \qquad
  A_s\neq\emptyset
\]
である。本命題は仮定を用いない。
\end{proposition}

\begin{proof}
$s$ に関する帰納法による。
$x_s\in\mathrm{Viab}_f(K)$ とすれば、$e_s\in E_{\mathrm{req}}(x_s)$ と $u_s\in U^{\mathrm{Viab}}(x_s)$ より
$x_{s+1}=f(x_s,u_s,e_s)\in\mathrm{Viab}_f(K)$ である。
さらに $\mathrm{Viab}_f(K)\subseteq K$ であるから $f(x_s,u_s,e_s)\in K$、
すなわち $u_s\in A_s$ であり、$A_s\neq\emptyset$ である。
\end{proof}

\begin{definition}[核維持応答の獲得（分類C）]
\label{def:kernel_response_acquisition}
軌道について、ある時刻 $T$ が存在して、すべての $s\ge T$ で
$E_{\mathrm{req}}(x_s)\neq\emptyset$ かつ $u_s\in U^{\mathrm{Viab}}(x_s)$ が成り立つとき、
その軌道は時刻 $T$ に\emph{核維持応答を獲得している}という。
条件 $E_{\mathrm{req}}(x_s)\neq\emptyset$ を含めるのは、
$E_{\mathrm{req}}(x_s)=\emptyset$ のとき $U^{\mathrm{Viab}}(x_s)=U$ となり、
$u_s\in U^{\mathrm{Viab}}(x_s)$ が空虚に成立するからである。
これは軌道の性質であり、実効行動 $u_s$ の生成機構
（\S\ref{sec:tipa}）と内部モデルの更新（\S\ref{sec:extended_dynamics}）を通じて
成立する場合も成立しない場合もある。
\end{definition}

\begin{remark}[結果の差と非主張事項]
\label{rem:kernel_response_scope}
命題~\ref{prop:kernel_maintenance} と主命題（命題~\ref{prop:sd_robust_nonresponse}）は、
同じ枠組みの中で次の二つの場合を書き分ける。
可行性核に留まる応答が存在し、それが選ばれ続ける場合には、閉塞は生じない。
可行な状態で頑健応答不能が続く場合には、可行性核に留まる応答は存在しない。
後者の場合にその後の軌道がどうなるかは、実現する外部環境状態の列と行動の選択に依存し、
閉塞への接続は強仮定のもとでの系としてだけ述べる。

本論文は次のいずれも一般的な帰結として置かない。
\begin{enumerate}
  \item 核維持応答の獲得が起こること、またはその時刻の上界。
  \item 内部モデルの更新（学習）が核維持応答の獲得を導くこと。
  \item $U^{\mathrm{Viab}}(x)\neq\emptyset$ から、実際に選ばれる行動が $U^{\mathrm{Viab}}(x)$ に属すること。
\end{enumerate}
ACS が $\mathrm{Viab}_f(K)$ や $U^{\mathrm{Viab}}(x)$ を内部に保持または計算することも要求しない
（\S\ref{subsec:description_levels_global}）。
\end{remark}

\subsubsection{結果の論理的地位}

\begin{remark}[命題~\ref{prop:sd_necessity} の必要性方向の意味]
\label{rem:sd_necessity_scope}
命題~\ref{prop:sd_necessity} は、
ある一時点で $\mu_{\mathrm{eff},t}=1$ となっただけで
直ちに閉塞へ到達することを主張しない。

同命題が示すのは、条件 (a)、条件 (b)（可行な応答の弱点を含む）、および
仮定~\ref{ass:theta_exterior_blockage} が成立する場合に、
$\mu_{\mathrm{eff},s}=1$ という応答多様性縮退を
閉塞に到達しないまま無期限に持続させることができない、
という条件付き必要性の方向である。

したがって Phase~1 で得られる帰結は、
各時刻における $\mu_{\mathrm{eff},t}\ge2$ を
無条件の普遍必要条件として証明した、という主張ではない。
本小節が示すのは、指定された追加条件のもとで
\emph{応答多様性の持続的縮退を回避すること}が
長期的な非閉塞維持に対して条件付きで要求される、
という構造的帰結である。
これは FR の長期の側面の一部を、追加仮定のもとで記述するものである。
強仮定を用いずに示せるのは、同じ条件 (a)(b) のもとでの頑健応答不能と可行性核の外への持続
（主命題、命題~\ref{prop:sd_robust_nonresponse}）までである。

また、$E_{\mathrm{req}}(x_t)=\emptyset$ となる時刻について、本命題は閉塞を帰結としない。
本論文は、その事実だけから $A_t=\emptyset$ または崩壊を一般的な帰結として置かない。
上記の追加条件のいずれかが成立しない場合、
本命題は閉塞について何も述べない。
\end{remark}

\begin{remark}[循環参照の不在]
\label{rem:no_circular}
主命題（命題~\ref{prop:sd_robust_nonresponse}）の形式的証明は、
定義~\ref{def:at}、定義~\ref{def:mu_eff}、定義~\ref{def:req_domain}、
定義~\ref{def:viab_core}、定義~\ref{def:nonresponse_env_set}、定義~\ref{def:robust_response}、
命題~\ref{prop:viab_core_max}、命題~\ref{prop:robust_nonresponse_viab}
および補題~\ref{lem:sd_viab} に依拠し、仮定を用いない。
系（命題~\ref{prop:sd_necessity}）は、これに加えて
仮定~\ref{ass:theta_exterior_blockage} に依拠する。
命題~\ref{prop:sd_necessity_fr} は、これに加えて
定義~\ref{def:E_req}、定義~\ref{def:eu_fire}、
定義~\ref{def:fr_domain}、定義~\ref{def:SM} および定義~\ref{def:FR} に依拠する。
命題~\ref{prop:kernel_maintenance} は、定義~\ref{def:at}、定義~\ref{def:viab_core} および定義~\ref{def:kernel_response} だけに依拠する。

仮定~\ref{ass:nonstationary} は
環境応答多様性の構造的意義を述べる背景条件であり、
備考~\ref{rem:viability_relation} は
Viability Theory との概念的位置づけを整理するものであって、
いずれもこれらの命題の形式的証明には使用しない。
補題~\ref{lem:finite_exit} および
命題~\ref{prop:exterior_blockage} もこれらの命題の証明に使用しない。

また、これらの命題の形式的証明の依存関係には
\S\ref{subsec:sd_necessity} より後の概念を用いない。
したがって、証明構造に後続節からの循環的依存は生じない。
\end{remark}

\begin{remark}[条件付き必要性の論理的地位]
\label{rem:sd_status}
本小節で扱う応答多様性縮退と長期可行性の接続は、
SM の定義的同値性（命題~\ref{prop:sm_K_equiv}）とも、
局所 FR と閉塞の定義上の一致（定義~\ref{def:FR}）とも異なる論理的地位をもつ。

SM と局所 FR は単一時刻の時刻別充足述語である。
これに対し応答多様度 $\mu_{\mathrm{eff},t}$ は、
必要環境要因が充足される外部環境状態の差異に対して、
一段階の可行性保持の様式がいくつに分かれているかを表す分類Cの診断量であり、
いずれの時刻別充足述語の定義にも用いない。
本論文は、単一時刻の $\mu_{\mathrm{eff},t}$ の値だけから、
長期的な応答多様性の維持に関する帰結を一般的な帰結として置かない。

本小節で条件付き決定論的に示したのは、
長期的な応答多様性維持の十分条件でも、
$\mu_{\mathrm{eff},t}\le1$ の一般からの閉塞でもない。
示したのは、$E_{\mathrm{req}}(x_s)$ が実現環境を含む状態が継続し、
可行な状態で応答の弱点と $\mu_{\mathrm{eff},s}=1$ という応答多様性縮退が持続する場合の、
頑健応答不能と可行性核の外への持続という条件付き帰結（主命題）であり、
さらに可行性核外部持続から閉塞への強い接続仮定を置いた場合の、
有限時間閉塞という条件付き帰結（系）である。
命題~\ref{prop:sd_necessity_fr} は、その閉塞時刻が
$x\notin K$ の時刻であるか局所 FR が偽となる時刻であるかを分ける系である。

\textit{注（三原理と可行性核の関係）：}
本論文は三原理を、可行性核を計算してそこへ移動する
アルゴリズムとして定義しない。
SM は現在状態の可行性を、EU は時刻ごとの臨界必要要因の充足を、
局所 FR は $\mathcal{D}_{\mathrm{FR}}$ の局面での可行な行動の存在を、
それぞれ時刻別に記述し、
これらが長期的な可行性維持と整合しうる、という関係に留める。
有限段の頑健性との関係は命題~\ref{prop:viab_stage_fr} の一方向に限る。

本論文は、
「三原理の時刻別充足述語がすべて真ならば可行性核内に留まる」とも、
「可行性核内にいれば三原理の時刻別充足述語がすべて真である」とも述べない。
また、可行性核への所属を ACS の成立条件へ格上げしない。

したがって本小節の結果は、
Phase~1 において形式的に証明できる範囲を、
応答多様性縮退の持続に対する条件付き必要性の方向として
明確化するものである
（主張の強さ $\le$ 証明の強さ）。
\end{remark}

\begin{remark}[応答多様度の「程度」命題は Phase~2 の課題]
\label{rem:sd_degree}
$\mu_{\mathrm{eff},t}$ の値と閉塞到達時間との単調関係、
例えば「応答多様度が大きいほど
可行行動集合が空になるまでの時間が長くなる」
という定量的関係は、
命題~\ref{prop:sd_necessity} は主張しない。

この種の程度命題を形式化するには、
追加の動学的仮定、比較原理、時間平均条件、
または確率的定式化等が必要となる。

Phase~1 で確立するのは、
$\mu_{\mathrm{eff},t}=1$ という応答多様性縮退が
指定された条件のもとで持続する場合の
頑健応答不能の持続（主命題）と、強仮定のもとでの有限時間閉塞（系）という条件付き帰結であり、
$\mu_{\mathrm{eff},t}$ の大小と閉塞到達時間との
定量的単調関係ではない。
また本論文は、$\mu_{\mathrm{eff},t}=1$ 単独から
閉塞を一般的な帰結として置かない。
\end{remark}

\begin{remark}[可行性余裕度の改善可能性は補助診断である]
\label{rem:sd1_unfulfilled}
備考~\ref{rem:sd_improvement_aux} で導入した
\[
  \exists u\in A_t
  \quad\text{s.t.}\quad
  \Pi\bigl(f(x_t,u,e_t)\bigr)>\Pi(x_t)
\]
という条件は、現在の実現外部環境状態のもとで
可行性余裕度を一段階改善できるかを記述する補助診断であり、
定義~\ref{def:FR} の局所 FR の条件でも、
本小節の命題の前件でもない。

本論文は、この改善可能性が成り立たないことだけから、
局所 FR の偽、$\mu_{\mathrm{eff},t}\le1$、閉塞、
または将来の可行性喪失を一般的な帰結として置かない。
例えば $\Pi(x_t)$ がモデル上の高い可行性余裕度にあり、
一段階でさらに増加させる行動が存在しない場合でも、
\[
  E_{\mathrm{req}}(x_t)\neq\emptyset,
  \qquad
  \mu_{\mathrm{eff},t}\ge2
\]
が成立することはありうる。

逆に、可行性余裕度を一段階改善できる行動が存在しても、
本論文は、そのことだけから応答多様性の非退化や長期可行性を一般的な帰結として置かない。
\end{remark}

\begin{remark}[$\mu_{\mathrm{eff}}$ の測度論的位置づけと Phase~2 への委譲]
\label{rem:mu_measure}
$\mu_{\mathrm{eff}}(x,e)$ は、
$E_{\mathrm{req}}(x)$ 上の応答同値類による商集合
$A(x,e)/\!\sim$ の基数として定義される
（定義~\ref{def:mu_eff}）。

Phase~1 では $\mu_{\mathrm{eff}}$ を、
可行性保持の様式の数を表す濃度の診断量として用いる。
$\mu_{\mathrm{eff}}$ の有限性は仮定せず、
$\mu_{\mathrm{eff}}$ に関する閾値をいずれの時刻別充足述語の条件にも置かない。

$A(x,e)$ が連続的な構造をもつ場合に、
行動多様性を基数だけでなく連続量として定量化するためには、
同値類または商空間に対する可測性等の追加構造が必要となる。
適切な測度、商測度、被覆数、有効次元、
その他の連続的多様性指標を用いた
$\mu_{\mathrm{eff}}$ の精密化は Phase~2 の課題とする
（within present framework）。
\end{remark}

%----------------------------------------------------------------------
\subsection{三原理の時間的統合評価}
\label{subsec:s3_tstab}

本小節では、
三原理 SM・EU・FR に関係する状態を、
時間方向に統合して評価するための
Phase~1 の形式表現を与える。

ここで扱うのは、三原理そのものの再定義でも、
時刻別充足述語の変更でもない。
また、\S\ref{sec:rel_min} の関係品質 $Q$ を
EU の時刻別充足述語や構造的環境結合条件の
追加条件として組み込むことでもない。

本論文が構造的に要求するのは、
ACS が自己--環境相互作用のもとで、
可行性余裕の低下、
必要環境要因の非充足、
可行な応答の縮退、
および外部関係に伴う可行性関連差異等の少なくとも一部へ感応し、
それらを応答形成または内部更新に利用しうることである。

以下で導入する評価尺度、
評価ベクトル $s_t$、
有限時間窓、
および T-STAB は、
理論上の状態列について三原理に関係する状態を時間方向に集約する
記述・比較のための条件である。
これらは、ACS 内部で時間方向の評価を担う機能の実装でも、その機能の出力でもない。
ACS の内部評価が時間方向の統合を行うことは機能水準の事柄であり、
T-STAB はその内部評価の結果ではなく、理論上の値によって判定される。
したがって、
ACS がこれらの量の真値を
同一の数式、同一の尺度、
または同一の精度で
内部計算することは要求しない。

\subsubsection{三原理評価表現の参照スケール}

\begin{definition}[三原理評価表現の参照スケール]
\label{def:s3_eval_scales}
SM に関係する評価を $[0,1]$ 値へ正規化して表現するため、
対象 ACS に相対する正の可行性余裕度参照スケール $\Pi_{\mathrm{ref}}>0$ を用いる。
$\Pi_{\mathrm{ref}}$ はすべての ACS に共通する定数ではなく、
$\max_{x\in K}\Pi(x)$ と同一視することも要求しない。
\end{definition}

\begin{remark}[参照スケールの理論上の地位]
\label{rem:s3_scale_status}
定義~\ref{def:s3_eval_scales} の $\Pi_{\mathrm{ref}}$ は、
三原理の時刻別充足述語の条件ではない。
したがって、参照スケールの選択を変更しても、
定義~\ref{def:SM}、定義~\ref{def:EU}、定義~\ref{def:FR}
の時刻別充足述語の真偽は変化しない。

この参照量は、三原理に関係する差異を程度表現へ写像するための
数学的較正パラメータであり、
ACS の7タプルに新たな構成要素を追加するものでもない。
応答多様度 $\mu_{\mathrm{eff}}$ には参照スケールを置かず、
T-STAB の条件にも用いない。

具体的な数値較正、内部信号との対応、
および対象種別ごとの尺度設計は、
Phase~1 では規定しない
（within present framework）。
\end{remark}

\subsubsection{三原理評価の形式表現}

\begin{definition}[SM 評価尺度（分類C）]
\label{def:s_sm}
すべての $t\in\mathbb{N}_0$ について
\[
  s_{\mathrm{SM},t}
  :=
  \min\Bigl\{1,\ \max\Bigl\{0,\ \frac{\Pi(x_t)}{\Pi_{\mathrm{ref}}}\Bigr\}\Bigr\}
  \in[0,1]
\]
と定める。
$\Pi(x_t)\le0$ のとき $s_{\mathrm{SM},t}=0$ である。
$s_{\mathrm{SM},t}$ は $\mathrm{SM}_t$ の真偽ではなく、正の可行性余裕度の程度を表す。
\end{definition}

\begin{definition}[EU・FR 評価尺度（分類C）]
\label{def:s_eu}
\label{def:s_fr}
$t\in\mathbb{T}_{\mathrm{P},i}$ の時刻について、
$\mathrm{EU}_t$ が真のとき $s_{\mathrm{EU},t}:=1$、偽のとき $s_{\mathrm{EU},t}:=0$ と定める。
$t\in\mathbb{T}_{\mathrm{P},i}$ かつ $(x_t,e_t)\in\mathcal{D}_{\mathrm{FR}}$ の時刻について、
$A_t\neq\emptyset$ のとき $s_{\mathrm{FR},t}:=1$、$A_t=\emptyset$ のとき $s_{\mathrm{FR},t}:=0$ と定める。
それ以外の時刻には定義しない。
\end{definition}

\begin{definition}[三原理評価ベクトル（分類C）]
\label{def:S3}
$S_3:=[0,1]^3$ とする。
$t\in\mathbb{T}_{\mathrm{P},i}$ かつ $(x_t,e_t)\in\mathcal{D}_{\mathrm{FR}}$ の時刻について、
\[
  s_t:=\bigl(s_{\mathrm{SM},t},\,s_{\mathrm{EU},t},\,s_{\mathrm{FR},t}\bigr)\in S_3
\]
と定める。
$s_{\mathrm{SM},t}$ はすべての時刻で、$s_{\mathrm{EU},t}$ は $t\in\mathbb{T}_{\mathrm{P},i}$ で定義されるので、
この時刻では三成分はすべて定義される。
それ以外の時刻には $s_t$ を定義しない。
\end{definition}

\begin{remark}[三成分は同一物理量ではない]
\label{rem:s3_noncommensurate}
$s_{\mathrm{SM},t}$、$s_{\mathrm{EU},t}$、$s_{\mathrm{FR},t}$ は
いずれも $[0,1]$ に値をとるが、
同一の物理量または同一の測定尺度を共有するとは限らない。
$s_{\mathrm{SM},t}$ は正の可行性余裕度の程度を、
$s_{\mathrm{EU},t}$ と $s_{\mathrm{FR},t}$ は
それぞれの時刻別充足述語の真偽を表す。

したがって、三成分を単純平均して単一の「総合点」を構成すること、
または一つの成分の増大が他の成分の不足を補償すると解釈することは、
本小節では行わない。
\end{remark}

\subsubsection{T-STAB：三原理時間統合評価条件}

\begin{definition}[T-STAB 評価パラメータ]
\label{def:tstab_params}
SM 評価の下限参照値 $\underline{s}_{\mathrm{SM}}\in(0,1]$、
一時刻変動上限 $\delta_{\max}\in[0,1]$、
および有限窓幅 $W_T\in\mathbb{N}$（$W_T\ge1$）を用いる。
これらは対象 ACS の評価構造に相対する参照条件であり、
時刻別充足述語を変更しない。
\end{definition}

\begin{definition}[T-STAB：三原理時間統合評価条件]
\label{def:tstab}
$t\in\mathbb{N}_0$ について、
\[
  \mathcal{W}_{T,t}:=\bigl\{r\in\mathbb{N}_0\bigm|\max(0,t-W_T+1)\le r\le t\bigr\},
  \qquad
  \Delta\mathcal{W}_{T,t}:=\bigl\{r\in\mathcal{W}_{T,t}\bigm|r-1\in\mathcal{W}_{T,t}\bigr\}
\]
とする。
T-STAB は $t\in\mathbb{T}_{\mathrm{P},i}$ の時刻についてだけ判定する。
$t\in\mathbb{T}_{\mathrm{P},i}$ の時刻において ACS が \emph{T-STAB 条件を充足する}とは、
次の四条件がすべて成り立つことをいう。
\begin{description}[style=nextline]
  \item[(T1) SM 評価の下限]
  すべての $r\in\mathcal{W}_{T,t}$ について $s_{\mathrm{SM},r}\ge\underline{s}_{\mathrm{SM}}$。
  \item[(T2) 局所 FR の窓内充足]
  すべての $r\in\mathcal{W}_{T,t}$ について、$(x_r,e_r)\in\mathcal{D}_{\mathrm{FR}}$ ならば $A_r\neq\emptyset$。
  \item[(T3) EU の窓内充足]
  すべての $r\in\mathcal{W}_{T,t}$ について、$\mathrm{EU}_r$ が真である。
  \item[(T4) SM 評価の一時刻変動]
  すべての $r\in\Delta\mathcal{W}_{T,t}$ について $|s_{\mathrm{SM},r}-s_{\mathrm{SM},r-1}|\le\delta_{\max}$。
  $\Delta\mathcal{W}_{T,t}=\emptyset$ のときは、この条件は空虚に成り立つ。
\end{description}
$t\in\mathbb{T}_{\mathrm{P},i}$ ならば、補題~\ref{lem:principle_time_domain_down} により
$\mathcal{W}_{T,t}\subseteq\mathbb{T}_{\mathrm{P},i}$ であり、
(T3) の $\mathrm{EU}_r$ は各 $r\in\mathcal{W}_{T,t}$ で定義される。
(T1)・(T2)・(T4) は、各時刻で定義される量だけを用いる。
したがって四条件はそれぞれ独立に定義される。
\end{definition}

T-STAB の (T1)〜(T3) のもとで評価ベクトルがとる範囲は
$S_3^{\mathrm T}:=\{s\in S_3\mid s_{\mathrm{SM}}\ge\underline{s}_{\mathrm{SM}},\ s_{\mathrm{EU}}=1,\ s_{\mathrm{FR}}=1\}$
で表せる。
$S_3^{\mathrm T}$ は閉区間と一点集合の直積であり、$S_3$ の閉部分集合である。

\begin{proposition}[T-STAB から有限窓内の時刻別充足述語への含意]
\label{prop:tstab_implies_principles}
$t\in\mathbb{T}_{\mathrm{P},i}$ において T-STAB 条件が充足されるとする。
このとき、すべての $r\in\mathcal{W}_{T,t}$ について、
$r\in\mathbb{T}_{\mathrm{P},i}$、$\Pi(x_r)>0$ であり、
$\mathrm{SM}_r$・$\mathrm{EU}_r$・$\mathrm{FR}_r$ はいずれも定義されて真である。
また $s_r$ が定義され、$s_r\in S_3^{\mathrm T}$ である。
\end{proposition}

\begin{proof}
$r\in\mathcal{W}_{T,t}$ をとる。
補題~\ref{lem:principle_time_domain_down} により $r\in\mathbb{T}_{\mathrm{P},i}$ である。
(T1) と $\underline{s}_{\mathrm{SM}}>0$ により $s_{\mathrm{SM},r}>0$ であり、
定義~\ref{def:s_sm} により $\Pi(x_r)>0$、したがって $x_r\in K$ で $\mathrm{SM}_r$ は真である。
(T3) により $\mathrm{EU}_r$ は真である。
命題~\ref{prop:fr_domain_eu} により $(x_r,e_r)\in\mathcal{D}_{\mathrm{FR}}$ であり、
(T2) により $A_r\neq\emptyset$、したがって $\mathrm{FR}_r$ は真である。
以上から $s_r$ が定義され、$s_{\mathrm{EU},r}=s_{\mathrm{FR},r}=1$、
$s_{\mathrm{SM},r}\ge\underline{s}_{\mathrm{SM}}$ であるから $s_r\in S_3^{\mathrm T}$ である。
\end{proof}

命題~\ref{prop:tstab_implies_principles} の含意は有限窓 $\mathcal{W}_{T,t}$ 内の各時刻に限られる。
本論文は、T-STAB の充足から、窓の外の時刻における時刻別充足述語の真、
または無限時間にわたる可行性の保持を一般的な帰結として置かない。

\begin{proposition}[時刻別充足述語の真と T-STAB]
\label{prop:tstab_zero_margin}
$t\in\mathbb{T}_{\mathrm{P},i}$ で $\mathrm{SM}_t$・$\mathrm{EU}_t$・$\mathrm{FR}_t$ がすべて真であっても、
$\Pi(x_t)=0$ ならば、時刻 $t$ において T-STAB 条件は充足されない。
\end{proposition}

\begin{proof}
$\Pi(x_t)=0$ ならば、定義~\ref{def:s_sm} により $s_{\mathrm{SM},t}=0<\underline{s}_{\mathrm{SM}}$ である。
$t\in\mathcal{W}_{T,t}$ であるから (T1) が成り立たない。
\end{proof}

本命題は、$\Pi(x_t)=0$ となる時刻が存在することを主張しない。
時刻別充足述語の真から T-STAB の充足が従わないことは、
この命題の前件が成り立つ時刻について示されている。

関係品質の有限窓評価は \S\ref{subsec:ro_quality} で扱い、T-STAB の条件には含めない（備考~\ref{rem:Q_not_tstab_core}）。

\subsubsection{本小節の記述理論上の地位}

\begin{remark}[時間集約条件の記述上の地位]
\label{rem:tstab_descriptive_status}
本小節が扱うのは、三原理に関係する理論上の状態を有限窓で集約する条件である。
ACS が三原理に関係する状態差を時間方向に統合して
応答形成または内部更新へ利用しうるという構造的評価機能は、
これとは別の機能水準の事柄である。

$s_t$、$S_3$、$\mathcal{W}_{T,t}$、T-STAB、
および \S\ref{sec:rel_min} の $\bar Q_{i\leftarrow j,W_Q,t}$ は、
理論上の値によって定まる数学的表現であり、
ACS 内部の評価機能の実装またはその出力ではない。
したがって、これらの形式量は
\[
  \text{ACS が完全に知る真値}
\]
とも、
\[
  \text{ACS が外生的に最大化すべき目標}
\]
とも同一ではない。

また本小節は、三原理の時間統合評価から
新たな第四原理、新たな失敗状態、
または力学系論上の安定性定理を導入するものではない。
T-STAB は $t\in\mathbb{T}_{\mathrm{P},i}$ の時刻にだけ判定され、
その充足から窓内の時刻別充足述語の真が従う
（命題~\ref{prop:tstab_implies_principles}）。
逆向きの含意は、命題~\ref{prop:tstab_zero_margin} の前件が成り立つ場合に成り立たない。

次小節では、この時間統合評価とは別に、
三原理から構造的機能・標準機能形式への接続を扱う。
\end{remark}

%----------------------------------------------------------------------
\subsection{三原理から構造的機能・標準機能形式への接続}
\label{subsec:principles_to_tipa}

本小節では、
\S\ref{subsec:self_env_principles}--\S\ref{subsec:s3_tstab}
で定義・整理した三原理 SM・EU・FR と、
\S\ref{sec:tipa} で導入する TIPA・PSL・ISL との
記述水準上の接続を明確にする。

本小節の目的は、
三原理を新たな機構へ再定義することではない。
また、SM・EU・FR を
$M_1$・$M_2$・$M_3$ の三機構へ一対一に対応させることでもない。

\S\ref{subsec:description_levels_global} の
構造的機能と数学的形式表現の区別、
ならびに定義~\ref{def:abc_classification} の
分類A・B・Cの区別に従い、
本小節では
\[
  \text{原理水準の構造的制約}
  \;\longrightarrow\;
  \text{構造的機能}
  \;\longrightarrow\;
  \text{標準機能形式}
\]
という記述上の接続だけを与える。

\subsubsection{原理水準と機構水準の非同一性}

三原理は概念的に異なる制約軸である（備考~\ref{rem:principle_linkage}）。
しかし、そのことから
\[
  \mathrm{SM}\leftrightarrow M_1,
  \qquad
  \mathrm{EU}\leftrightarrow M_2,
  \qquad
  \mathrm{FR}\leftrightarrow M_3
\]
のような一対一の対応を置かない。
同一の情報統合処理、同一の評価処理、同一の記憶、
同一の応答形成、または同一の物理基盤が、
複数の原理に関係する差異を同時に担いうる。
逆に、一つの原理に関係する機能が
複数の機構・基盤・実装部位へ分散して担われることも許容する
（備考~\ref{rem:s3_description_levels}）。

\subsubsection{三原理に関係する構造的機能}

\begin{definition}[三原理関連構造的機能（分類A）]
\label{def:principle_related_functions}
本論文が対象とする ACS において、
三原理に関係する差異を
自己--環境相互作用のなかで
応答形成または内部更新へ接続するために要求される
機能水準の役割を、
\emph{三原理関連構造的機能}と呼ぶ。

本論文では少なくとも、次の機能群を区別する。

\begin{enumerate}

  \item
  \textbf{情報取得・統合機能}

  外部環境状態、内部状態、
  記憶、内部モデルその他の利用可能情報を取得し、
  後続の評価および応答形成に利用可能な
  内部表現へ統合する機能。

  \item
  \textbf{可行性感応評価機能}

  SM に関係する可行性余裕の低下、閾値接近、
  内部状態の悪化その他の
  可行性関連差異の少なくとも一部へ感応し、
  応答形成に利用可能な評価へ変換する機能。

  \item
  \textbf{応答可能性評価・再構成機能}

  FR に関係する、
  可行な応答の有無とその縮退、
  内部状態差に伴う応答余地の変化、
  必要環境状態差に対する応答差その他の差異の少なくとも一部へ感応し、
  必要応答ポテンシャルの保持・獲得・回復・再構成に
  利用可能な評価へ変換する機能。

  \item
  \textbf{環境利用・外部結合評価機能}

  EU に関係する、必要環境要因の充足・非充足とその接近、
  外部環境状態差への応答可能性、
  実現外部相互作用、
  および外部関係に伴う可行性関連差異の
  少なくとも一部を識別・評価する機能。

  \item
  \textbf{応答形成・選択機能}

  上記の統合情報および評価を用いて、
  現在の状態・環境に相対して応答候補を構成し、
  実効的な行動または内部応答へ接続する機能。

  \item
  \textbf{時間横断的情報保持・更新機能}

  単一時刻に閉じず、
  過去の情報、評価、応答結果その他を
  後続時刻の評価・応答形成に
  利用可能な形で保持または更新する機能。

  \item
  \textbf{物理的感知・作用・基盤維持機能}

  外部・内部信号を物理的に取得し、
  応答指令を実世界の作用へ接続し、
  さらに作動を支える物理的基盤を維持する機能。

\end{enumerate}

これらは定義~\ref{def:abc_classification} の意味で
分類Aの機能水準に属する。
ただし、上記の列挙は七つの独立モジュールの存在を要求するものではない。
同一の実装構造が複数の機能を兼ねること、
また一つの機能が複数の実装構造へ分散することを許容する。
各機能群の名称に付した原理名は主な関係を示すにとどまり、
その機能群が当該原理だけに関係することを意味しない
（備考~\ref{rem:principle_function_overlap}）。
\end{definition}

\begin{remark}[原始的な機能・基盤機能と三原理関連構造的機能]
\label{rem:primitive_vs_principle_functions}
\S\ref{sec:acs_def} の原始的な機能（定義~\ref{def:primitive_functions}）は
感知・識別・評価・行動の四区分であり、
基盤機能（定義~\ref{def:base_functions}）は
その責務を担う機能の活動を支え、担い手を持続させる機能である。
定義~\ref{def:principle_related_functions} の七つの機能群は、
これらの区分を三原理に関係する差異の側から記述し直したものであり、
新たな原始的な区分を加えるものではない。
これらの機能群の列挙は、\S\ref{sec:acs_def} の構成要件 CR-2 に条件を加えるものではない。

両者の間に一対一の対応は置かない。
主な関係は次のとおりである。
\begin{enumerate}
  \item 第1群（情報取得・統合）は、主に感知と識別の責務に関わる。
  \item 第2群から第4群（三原理に関係する評価）は、主に評価の責務に関わる。
  ただし、評価の前提となる差異の識別にも関わる。
  \item 第5群（応答形成・選択）は、主に行動の責務に関わる。
  \item 第6群（時間横断的情報保持・更新）は、四区分のいずれの責務にも関与しうる。
  情報を保持する担い手の持続という側面では基盤機能にも関わる。
  \item 第7群（物理的感知・作用・基盤維持）のうち、
  信号の物理的取得と作用への接続は感知と行動の責務に、
  作動を支える物理的基盤の維持は基盤機能に関わる。
  同一の過程について両方の側面を区別して記述しうることは、
  \S\ref{sec:acs_def} の扱いと同じである。
\end{enumerate}
ここに挙げた関係は記述上の整理であり、
各機能群をいずれか一つの区分または基盤機能に割り当てることを求めない。
また \S\ref{sec:acs_def} 備考~\ref{rem:primitive_vs_principles} と同じく、
本論文は、いずれかの機能が作動していることから時刻別充足述語が真であることを、
いずれかの時刻別充足述語が偽であることから特定の機能の停止を、いずれも一般的な帰結として置かない。
\end{remark}

\begin{remark}[三原理関連機能と時刻別充足述語の真の非同一性]
\label{rem:functions_not_principle_truth}
定義~\ref{def:principle_related_functions} の
構造的機能が存在することから、
各時刻における $\mathrm{SM}_t$・$\mathrm{EU}_t$・$\mathrm{FR}_t$ の真を、本論文は一般的な帰結として置かない。

例えば、可行性感応評価機能を備えていても、
実現状態が $\Pi(x_t)<0$ となることはありうる。
同様に、応答可能性評価・再構成機能が存在していても、
$\mathrm{FR}_t$ が偽となる時刻、
または $\mu_{\mathrm{eff},t}\le1$ となる時刻を排除しない。
また、環境利用・外部結合評価機能が存在していても、
$\mathrm{EU}_t$ が偽となる時刻を排除しない。

したがって、
\[
  \text{原理関連機能の存在}
  \neq
  \text{時刻別充足述語の常時の真}
\]
である。
\end{remark}

\begin{remark}[理論上の量への関連と内部での計算の区別]
\label{rem:theory_quantity_vs_internal}
定義~\ref{def:principle_related_functions} の機能群は、
$\Pi$、$K$、$A_t$、$\mathrm{Crit}^{(1)}(x_t)$、$\mathcal{D}_{\mathrm{FR}}$、$\mu_{\mathrm{eff}}$、$Q$ などの理論上の量が
表す差異に構造的に関連する。
しかし、関連することと、その量の真値を内部で保持または計算することは異なる。
本論文は、いずれの機能群についても、これらの量の真値を内部で保持または計算することを要求しない。
たとえば環境利用・外部結合評価機能について、
内部で必要と評価された要因が $\mathrm{Crit}^{(1)}(x_t)$ と一致すること、
または内部で形成された関係の評価が $Q$ と一致することを要求しない
（備考~\ref{rem:crit_description_level}）。
内部の評価と理論上の量との不一致（見逃し・誤評価）は、
\S\ref{sec:tipa} の機構水準で記述される事柄である。
\end{remark}

\subsubsection{TIPA への標準的接続}

\begin{remark}[TIPA は分類Bの標準機能形式である]
\label{rem:tipa_as_standard_form}
\S\ref{sec:tipa} で導入する
TIPA
（Triadic Information-Processing Architecture）は、
定義~\ref{def:principle_related_functions} の
構造的機能その他を
本論文の形式体系内で分析可能に分解する
分類Bの標準機能形式である。

TIPA は
$M_1,M_2,M_3$
からなり、
それぞれ情報統合、適応的評価、応答形成・行動選択を標準的に担う。

ここで重要なのは、
\[
  \text{情報統合機能}
  \neq
  M_1,
\]
\[
  \text{可行性感応評価機能}
  \neq
  M_2,
\]
\[
  \text{応答形成・選択機能}
  \neq
  M_3
\]
という区別である。
左辺は機能水準の要求、
右辺は本論文が採用する標準機能形式である。
\end{remark}

\begin{definition}[三原理から TIPA への機能的接続]
\label{def:principles_tipa_connection}
三原理に関係する情報処理と
TIPA の標準的役割分担を、
次のように記述する。

\medskip
\noindent
\textbf{$M_1$：情報統合機構。}

$M_1$ は、
外部・内部から利用可能な信号、内部モデル、記憶その他を統合し、
後続の評価機構が利用可能な内部表現を形成する標準機構である。
したがって、
SM に関係する内部状態信号、
EU に関係する必要環境要因の充足に関わる情報と外部環境差・外部対象差、
FR に関係する応答可能性または状態差その他の情報は、
必要に応じて $M_1$ の統合対象となりうる。
ただし、$M_1$ が理論上の
\[
  \Pi(x_t),\quad
  E_{\mathrm{req}}(x_t),\quad
  \mathrm{Crit}^{(1)}(x_t),\quad
  \mu_{\mathrm{eff},t},\quad
  Q
\]
を再構成することは要求しない。

\medskip
\noindent
\textbf{$M_2$：適応的評価機構。}

$M_2$ は、
$M_1$ 等から利用可能な情報に基づき、
可行性脅威、必要環境要因の非充足とその接近、
応答余地の縮退、外部関係の変化、切迫度その他を
応答形成へ利用可能な形で評価する標準機構である。
したがって、SM・EU・FR のいずれに関係する差異も、
その一部が $M_2$ の評価対象となりうる。
ただし、$M_2$ が理論上の値として
\[
  \Pi,\quad
  K,\quad
  A_t,\quad
  E_{\mathrm{req}},\quad
  \mathrm{Crit}^{(1)},\quad
  \mathcal{D}_{\mathrm{FR}},\quad
  \mu_{\mathrm{eff}},\quad
  Q,\quad
  s_t,\quad
  \text{T-STAB}
\]
を直接読み出すこと、
またはそれらを数値計算することは要求しない。

\medskip
\noindent
\textbf{$M_3$：応答行動機構。}

$M_3$ は、
統合された情報および評価を用いて、
応答候補の構成、要請適合の判定、行動指令その他の応答形成を担う
標準機構である。
したがって、三原理に関係する評価差は、
必要に応じて $M_3$ の候補形成または選択へ反映されうる。
ただし、三原理は $M_3$ に外部から与えられる
単一の目的関数ではない。
また本論文は、$M_3$ の出力だけから時刻別充足述語の真を一般的な帰結として置かない。
$\mathrm{FR}_t$ が真であっても、$M_3$ の行動指令または実効行動が
$A_t$ に属することは要求しない（備考~\ref{rem:selection_responsibility}）。
\end{definition}

\begin{remark}[M1・M2・M3 は三原理専用機構ではない]
\label{rem:tipa_not_principle_modules}
定義~\ref{def:principles_tipa_connection} は、
$M_1$、$M_2$、$M_3$ を
それぞれ特定の原理専用の機構として定義するものではない。

例えば、SM に関係する危険信号の処理には
$M_1$ の情報統合、$M_2$ の評価、$M_3$ の応答形成が
すべて関与しうる。
同様に、EU に関係する必要環境要因の充足や外部差・外部関係に伴う可行性関連差異、
FR に関係する応答の縮退も、複数機構を横断して扱われうる。

したがって、TIPA における三機構の区別は
\[
  \text{入力統合}
  \rightarrow
  \text{評価}
  \rightarrow
  \text{応答形成}
\]
という機能責務上の標準分解であり、
SM・EU・FR の三分解と同型であることを主張しない。
\end{remark}

\subsubsection{PSL への接続}

\begin{definition}[原理関連機能と PSL の接続]
\label{def:principles_psl_connection}
\S\ref{subsec:psl} で導入する
PSL（Physical Substrate Layer）は、
TIPA と物理的実装部位とのあいだで、
感知、作用、基盤維持、局所的な保護応答その他を
物理的に接続する分類Bの標準機能形式である。

三原理との関係は、次のように記述する。

\begin{enumerate}

  \item
  \textbf{入力側接続。}

  SM・EU・FR に関係する外部・内部信号が
  TIPA の情報処理へ利用可能となるためには、
  それらを物理的に取得・伝達する入力側の基盤機能を要する。

  \item
  \textbf{出力側接続。}

  TIPA で形成された応答指令が実効的な作用へ接続されるためには、
  出力側の物理的実現機能を要する。

  \item
  \textbf{基盤維持。}

  情報処理・記憶・感知・作用を継続可能にするためには、
  それらを担う基盤の維持機能を要する。

  \item
  \textbf{局所急性異常への保護応答。}

  熟慮的な評価完了を待つことが可行性維持上不適切となりうる
  急性局所異常に対しては、
  局所的・低遅延の保護応答機能が三原理関連作動を支えることがある。

\end{enumerate}

これらの具体的な標準形式は \S\ref{subsec:psl} に委ねる。
\end{definition}

\begin{remark}[PSL 機能は三原理そのものではない]
\label{rem:psl_not_principles}
PSL の各標準機能形式は、SM・EU・FR のいずれかと同一ではない。

例えば、基盤維持機能が SM に強く関係しても、
\[
  \text{PSL 基盤維持}
  \iff
  \mathrm{SM}
\]
とはしない。
また、局所的保護応答が急性の可行性脅威へ対処しうるからといって、
本論文は、その作動だけから $\mathrm{EU}_t$ または $\mathrm{FR}_t$ を含む
三原理の時刻別充足述語の真を一般的な帰結として置かない。

PSL は、原理関連機能を物理的作動へ接続する標準機能形式である。
\end{remark}

\subsubsection{ISL への接続}

\begin{definition}[原理関連機能と ISL の接続]
\label{def:principles_isl_connection}
\S\ref{subsec:impl} で導入する
ISL（Implementation Substrate Layer）は、
感知、作用、情報保持その他を担う物理的実装部位を
本論文の標準形式で整理する分類Bの記述層である。

三原理関連機能との接続は、
少なくとも次の三つの物理的能力を介して記述される。
\[
  \text{感知可能な物理担体},
\]
\[
  \text{作用可能な物理担体},
\]
\[
  \text{情報保持可能な物理担体}.
\]
\S\ref{subsec:impl} では、これらを標準的に
\[
  B_{\mathrm{sens}},
  \qquad
  B_{\mathrm{eff}},
  \qquad
  B_{\mathrm{mem}}
\]
として整理する。
ただし、これら三集合が実在 ACS において
互いに完全に分離した三つの独立器官または
三つの独立ハードウェアとして存在することは要求しない。
\end{definition}

\begin{remark}[実装部位と原理の局在を同一視しない]
\label{rem:principles_not_localized}
SM・EU・FR のいずれも、
単一の $B_{\mathrm{sens}}$、$B_{\mathrm{eff}}$、$B_{\mathrm{mem}}$
または単一の物理部位へ局在するとはしない。
原理関連機能は、感知・記憶・評価・応答・基盤維持の
複数過程を横断して成立しうる。
したがって、生物実装、人工実装、その他の実装形態において、
同じ原理水準の制約が異なる物理配置によって担われることを許容する。
\end{remark}

\subsubsection{時間方向の変化との役割分担}

\begin{remark}[\S\ref{sec:tipa} と \S\ref{sec:extended_dynamics} の責務境界]
\label{rem:s4_s5_boundary_from_s3}
本小節は、三原理に関係する構造的機能を
TIPA・PSL・ISL の標準機能形式へ接続する。

\S\ref{sec:tipa} は、同一時刻における
機構・基盤・実装部位の構造と作動を扱う。
これに対し \S\ref{sec:extended_dynamics} は、
\S\ref{sec:tipa} で定義済みの構造を再定義せず、
内部モデル、機能的記憶、偏り、実効的発展自由度、予測地平、
PSL・ISL の健全度その他について、
時刻 $t$ から $t+1$ への決定論的な変化を扱う。

したがって、
\[
  \text{\S\ref{sec:tipa}}
  =
  \text{同一時刻の構造・作動}
\]
\[
  \text{\S\ref{sec:extended_dynamics}}
  =
  \text{既定義構造・内部量の時刻間変化}
\]
という責務境界を維持する。
\end{remark}

\begin{remark}[Update と三原理関連機能]
\label{rem:update_principle_connection}
7タプルの $\mathrm{Update}$ は、内部モデルの時刻間更新を表す。
三原理に関係する評価結果、記憶、内部状態差、外部環境差その他が
後続時刻の内部モデルへ反映されうることは、
三原理関連構造的機能と Update の接続として理解できる。
FR の長期の側面である必要応答ポテンシャルの保持・獲得・回復・再構成は、
この接続を通じて記述されうる。
ただし、本節は Update の具体的アルゴリズム、学習率、収束特性、
または最適化法を規定しない。
それらの時刻間更新構造は \S\ref{sec:extended_dynamics} へ委ねる。
\end{remark}

\subsubsection{三原理と標準機能形式の対応範囲}

\begin{remark}[原理別の主な機能接続は重複しうる]
\label{rem:principle_function_overlap}
三原理と機構層の関係を直感的に整理すると、
次のような主な接続を記述できる。

\begin{description}

  \item[SM 関連]
  内部状態および可行性脅威に関係する信号の取得、
  統合、評価、応答形成、物理的保護、基盤維持。

  \item[EU 関連]
  必要環境要因の充足・非充足とその接近に関係する信号の取得と評価、
  外部環境状態差の感知、外部入力の識別、
  実現外部相互作用の処理、外部関係に伴う可行性関連差異の評価、
  必要環境要因の充足へ向かう応答形成。

  \item[FR 関連]
  可行な応答の有無とその縮退の評価、
  内部状態差に伴う応答余地の変化の評価、
  必要環境状態差に対する応答差の評価、
  記憶・内部モデル更新、応答候補の再構成。

\end{description}

ただし、この整理は原理ごとの排他的な担当表ではない。
例えば、外部環境差の感知は EU に関係するだけでなく、
SM の危険検出や FR の応答多様性の維持にも同時に関係しうる。
同様に、記憶は FR だけでなく、SM・EU に関係する評価にも利用されうる。
したがって、三原理と機構機能の対応は多対多である。
\end{remark}

\subsubsection{本小節の記述理論上の地位}

\begin{remark}[原理水準から機構水準への接続の最小的主張]
\label{rem:principles_to_mechanisms_scope}
本小節が主張するのは、
三原理に関係する構造的差異が
ACS の作動と完全に切断されていてはならず、
それらの差異の少なくとも一部が、
情報取得・統合、評価、応答形成、記憶・更新、
物理的感知・作用・基盤維持を通じて
後続の作動へ接続されうるという
機能水準の要請である。

これを、本論文では分類Aの構造的機能として整理し、
\S\ref{sec:tipa} では TIPA・PSL・ISL という
分類Bの標準機能形式によって記述する。
したがって本小節は、
\[
  \text{三原理}
  \longrightarrow
  \text{構造的機能}
  \longrightarrow
  \text{標準機能形式}
\]
という記述上の橋渡しを与える。
ただし、
三原理、特定機構および特定物理部位は、
相異なる記述水準に属するため、
相互に同一視しない。

具体的な TIPA・PSL・ISL の定義、
入力・処理・出力の責務境界、高速処理、Automaticity、
物理的入出力接続、および実装部位構造は
\S\ref{sec:tipa} に委ねる。
時刻間の更新・学習・健全度変化その他は
\S\ref{sec:extended_dynamics} に委ねる。
\end{remark}


%================ Section 4 TIPA・PSL・実装部位層 ================
%======================================================================
\section{TIPA・PSL・ISL}
\label{sec:tipa}

\S\ref{sec:three_principles} では、
自己維持（SM）、環境利用（EU）、可行応答（FR）の三原理を、
存続中のACSに対する制約軸として定式化した。
本節ではそれらを再定義せず、
単一 ACS が同一時刻に情報を取得・統合し、
可行性維持に関係する差異を評価し、
応答を形成して実効行動へ接続するための標準機能形式を与える。

本節で扱う主要な標準機能形式は、
三機構情報処理アーキテクチャ・TIPA（Triadic Information-Processing Architecture）、
物理的基盤層・PSL（Physical Substrate Layer）、
実装部位層・ISL（Implementation Substrate Layer）である。
TIPA は情報統合、適応的評価、応答行動形成を担い、
PSL は TIPA と実装側の入力・出力・基盤維持・局所保護応答を接続し、
ISL は感知・作用・情報保持等を担いうる実装上の担体を整理する。

本節の責務は、
これらの同一時刻における構造、情報依存、作動形式および接続関係を定めることである。
記憶更新、学習、内部モデル更新、
統合的高速処理および Automaticity の形成・強化・弱化・解除、
PSL・ISL の健全度変化その他の
$t\to t+1$ の時刻間動学は
\S\ref{sec:extended_dynamics} に委ねる。

\begin{remark}[原理水準と機構記述の関係]
\label{rem:principle_level_locus}
SM、EU、FR は \S\ref{sec:three_principles} で定める制約軸であり、
TIPA、PSL、ISL のいずれか一つと同一視しない。
原理そのものについては成立・不成立を述べず、
真偽をもつのは各原理の時刻別充足述語である（備考~\ref{rem:s3_description_levels}）。
本論文は、本節の機構出力または作動状態だけから、
三原理の時刻別充足述語の真偽を一般的な帰結として置かない。
\end{remark}

以下、\S\ref{subsec:s4_scope}で本節の記述範囲と共通規約を、
\S\ref{subsec:tipa_architecture} で TIPA の三機構と三機能を、
\S\ref{subsec:psl} で PSL の下位責務と実効行動解決を、
\S\ref{subsec:purpose_action_result} で目的・行動・結果の記述を、
\S\ref{subsec:impl} で ISL を、
\S\ref{subsec:fast_processing} で統合的高速処理と Automaticity を、
\S\ref{subsec:functional_axis} で機能軸を扱う。
同一時刻の機能的接続の全体は、小括の図~\ref{fig:s4_same_time} に示す。

%----------------------------------------------------------------------
\subsection{本節の記述範囲と共通規約}
\label{subsec:s4_scope}

本小節では、
本節全体を貫く責務境界と共通数学装置のみを確認する。
本節で新たに用いる主要記号の一覧は付録~\ref{app:symbols_s4} に示す。
各機構または作動形態に固有の仮定・区別は、
それぞれを担う後続小節で与える。

\subsubsection{同一時刻構造と時刻間動学の責務境界}
\label{subsubsec:s4_same_time_boundary}

本節は TIPA、PSL、ISL の
同一時刻における構造、情報依存、作動および接続を扱う。
これらの形成、学習、更新、劣化、回復その他の時刻間変化は、
\S\ref{sec:extended_dynamics} に委ねる。

したがって本節では、
ある時刻にどの情報がどの機構から利用可能となり、
各機構が何を生成し、
最終的に状態遷移へ入力される行動がどのように構成されるかを定める。

\subsubsection{構成的必要機能と標準機能形式の区別}
\label{subsubsec:s4_abc_application}

本節で用いる分類A、分類B、分類Cは、
定義~\ref{def:abc_classification} の一般規則に従う。
本節では分類そのものを再定義せず、
構成的に必要な機能と、
TIPA・PSL・ISL その他の分類Bの標準機能形式とを同一視しないこと、
および分類Cの量を ACS の追加構成要素とみなさないことのみを確認する。

具体的な器官、回路、ハードウェア、ソフトウェアモジュールその他の実装は、
分類Bの標準機能形式と一対一に対応することを要求しない。

本節はまた、\S\ref{subsec:description_levels_global} の記述水準の規約を
そのまま適用する。
すなわち、ある形式量を導入することは ACS がその真値を知ることを意味せず、
ある構造的機能が必要であると述べることは
ACS が本論文の特定の数式を内部に実装することを意味しない。
本節が構造的機能の存在について要求するのは、
当該機能に関係する差異の少なくとも一部に対して
評価結果または内部応答が変化しうることに限られる。

本節は単一のACS個体 $\Sigma_i$ を扱うため、
\S\ref{sec:acs_def} の規約に従い、
状態・環境・可行行動集合等の個体添字を省略して
$x_t$、$e_t$、$A_t$ と書く。
個体の同一性に関わる崩壊指標には添字を付して $\mathrm{Col}_{i,t}$ と書く。

\subsubsection{決定論的選択射影}
\label{subsubsec:selection_projection}

\begin{definition}[決定論的選択射影（分類B）]
\label{def:sel_proj}
集合 $Y$ 上の非空候補集合族を
\[
  \mathfrak C_Y
  \subseteq
  \mathcal P(Y)\setminus\{\emptyset\}
\]
とする。
決定論的選択射影を
\[
  \mathrm{Proj}_{\mathrm{sel},Y}:
  \mathfrak C_Y
  \to Y
\]
とし、任意の $C\in\mathfrak C_Y$ について
\[
  \mathrm{Proj}_{\mathrm{sel},Y}(C)\in C
\]
を満たす固定された選択規則とする。

この写像は、複数の同等候補が存在する場合にも
Phase~1 の決定論性を維持するための共通形式装置である。
具体的なタイブレーク規則は実装依存とする。
\end{definition}

\subsubsection{状態担体と情報来歴}
\begin{definition}[同時刻作動の担体規約（分類B）]
\label{def:s4_carrier_contract}
本節の時刻依存の機構出力は、固定した規則のもとで $(M_t,I_t)$ から生成する。
$I_t=P_{\mathrm{sens}}(x_t,e_t)$ とし、$I_{\mathrm{sense},t}=I_t$ と書く。
物理的信号・PSL基盤・感度・効果器の状態は $x_t$ に、
保持内容は $M_t^{\mathrm L},M_t^{\mathrm W}$ に担わせる。
本節で用いる各信号空間・情報空間は非空集合とする。
位相を用いる結果は、既定義の $X,\mathcal M$ の位相に限り、
信号空間のコンパクト性や読み出しの連続性は要求しない。

各前時刻入力は、以下の固定された読み出しによって与える。
\[
\begin{aligned}
\rho_{\mathrm{past}}:\mathcal M^{\mathrm W}\longrightarrow
 \mathcal P_{\mathrm{prio}}\times U\times X^{\le W}
 \times\mathsf S_{M1}\times\mathcal I^{\le W_{\mathrm{short}}}
 \times\mathcal P_{\mathrm{fin}}(\mathsf{Pur})\times\mathcal G_O.
\end{aligned}
\]
右辺の型は後掲の M1・M2・M3 の定義による
（$\mathsf{Pur}$ は定義~\ref{def:purpose_formation}、$\mathcal G_O$ は定義~\ref{def:env_set_prediction}）。
各成分を、保持された前回Priority、実効行動の記録、状態推定履歴、
M1出力の記録、感知情報履歴、前時刻の目的の集合の記録、
および前時刻の実効行動についての外部側予測集合の記録として用いる。
最後の二つの記録が形成されるのは、
Update の詳細型（\S\ref{sec:acs_def} 備考~\ref{rem:update_type_two_layers}）が
同一ステップの中間値を追加引数として受け取る範囲に限られる。
目的の集合の記録は M3 の作動成分 $\mathcal M_3^{\mathrm W}$ に担わせる。
以降、機構への入力に現れる前時刻記号はこの記録値を表し、
実際の過去出力との同一性は仮定しない。
初期記録も $M_0^{\mathrm W}$ から読み出し、独立初期状態を追加しない。
記録の欠落・誤りを許すが、機構はその扱いを含め固定した全域写像として与える。
空履歴または固定した初期規約値を用いる場合、それを実際の過去値とはしない。
欠落の識別が必要な実現では、その識別情報も既存の作動成分に保持する。
不足する履歴を真の過去値で補完しない。

$\rho_{\mathrm{past}}$ の存在を任意の符号化可能性の定理とはしない。
採用する Update の成分型の下で実現される記録のみを用いる。
特に Update$_1$ は他機構の状態を直接入力としないので、
他機構由来の記録を機構1に正確に保存できるとは一般にはいえない。
読み出しの担体を $\mathcal M^{\mathrm W}$ 全体とすることは、この更新制約を解除しない。
\end{definition}

%----------------------------------------------------------------------
\subsection{TIPA：三機構情報処理アーキテクチャ}
\label{subsec:tipa_architecture}

TIPA を
\[
  \mathrm{TIPA}:=(M_1,M_2,M_3)
\]
として、情報統合、適応的評価、応答行動形成を
三つの機構へ機能分解する分類Bの標準機能形式とする。
この分解は、三つの独立した物理モジュールの存在を要求するものではない。

Priority の標準型を
\[
  \mathcal P_{\mathrm{prio}}
  :=
  \{
    \mathrm{Priority\ 1},
    \mathrm{Priority\ 2},
    \mathrm{Priority\ 3},
    \mathrm{None}
  \}
\]
とし、標準動作モード集合を
\[
  \mathsf{Mode}
  :=
  \{\mathrm{3A},\mathrm{3B},\mathrm{3C}\}
\]
とする。
また、M3 の応答様式出力の型を
\[
  \mathsf{Style}
  :=
  \{\mathrm{3B\mbox{-}own},\mathrm{3B\mbox{-}rel},\mathrm{3B\mbox{-}comb},\mathrm{none}\}
\]
とする。
これらの意味は M2 および M3 の該当箇所で与える。

TIPA 内の概念的な生成依存は
\[
  M_1\longrightarrow M_2\longrightarrow M_3
\]
とする。
これは、必要な上流出力が下流機構から利用可能となる情報依存を表す。

\begin{assumption}[TIPA三機構の機能的並行性]
\label{ass:functional_concurrency}
$M_1$、$M_2$、$M_3$ は、
上記の生成依存を保ちながら、
同一時刻に機能的に並行して作動しうるものとする。

したがって、
$M_1\to M_2\to M_3$ は情報生成依存を表し、
三機構を物理時間上またはハードウェア上の
厳密な逐次モジュールへ固定するものではない。

導入地位：S（構造的）。強さ：中。
依拠箇所：内部モデル構造の定義（\S\ref{sec:acs_def} 定義~\ref{def:internal_model}）における三機構の作動の読み方。
本節の命題の前提としては用いない。
\end{assumption}

\begin{remark}[機構間の送受信と同一時刻の記述]
\label{rem:tipa_bidirectional}
離散時間の標準形では、同一時刻の中には上記の生成依存の向き $M_1\to M_2\to M_3$ だけを置く。
同一時刻の中で出力が互いに入力となる関係は、
各時刻の値を定めるために、不動点の存在と、
その一意性または値を一つに定める決定論的な規則とを別に要する記述となるので、本節では置かない。

逆向きの影響は、時刻をまたいで次の時刻の入力として表す。
本節の標準形がもつ経路は次のとおりである。
\begin{enumerate}[label=(\roman*)]
  \item 保持記録を通じる経路（定義~\ref{def:s4_carrier_contract}）。
  M1 は、前時刻の M1 出力の記録 $\mathsf S_{M1,t-1}$ と状態推定履歴、
  前回 Priority $P_{t-1}$、最終実効行動の記録 $u_{t-1}$ を入力とし、
  目的の集合の記録を目的対象性の成分の形成に用いうる。
  M2 は自身の過去の評価の記録を読み出しうる
  （定義~\ref{def:m2_latent_instability} の $\rho_{\mathrm{eval,past}}$）。
  M3 は前時刻の目的の集合 $\mathrm{Pur}_{t-1}$ を入力とする。
  前時刻の動作モード $\mathrm{mode}_{t-1}$ は、
  標準形では独立の保持記録または明示的入力として設けない。
  ただし、前時刻のモードに由来する情報が
  $M_t$ などの既存の入力に実際に保持されている場合、
  その情報が現在のモード判断に関与することは排除しない。
  \item 状態と外部環境状態を通じる経路。
  実効行動 $u_t$ は $f$ と行動応答成分の遷移を通じて次の時刻の状態と外部環境状態を変え、
  その変化は感知を経て M1 に入る。
\end{enumerate}
いずれの経路も、記録または状態が実際にその内容を担う範囲に限られる。
M3 の出力から M2 への経路は置かない（定義~\ref{def:m2_map}、備考~\ref{rem:m2_no_return_scope}）。
また、PSL-4 の局所保護応答は三機構を経由しない経路であり（\S\ref{subsubsec:psl4}）、
$M_1\to M_2\to M_3$ は当該ACSの信号のすべてが通る経路ではない。

実装において、三機構は、それぞれ接続先から情報を受け取り、
処理した情報を出力しつづけるものとして読む。
$M_1$、$M_2$、$M_3$ の番号は記述上の順番であり、
作動中には機構の間で双方向の送受信が起こりうる。
本論文が各機構について記述するのは、
受け取った情報をその機構の枠組みの中で処理し、
型宣言に合う出力を与えるところまでである。
出力の性質のクラス分けの詳細は、本節で意味の解釈を与えたものを除いて定めない。
一つの機構の内部の反復は、その機構の時刻別の写像の中に縮約して表す
（備考~\ref{rem:internal_steps_time}）。
これに対し、実装が同一時刻の中で機構の間を往復する場合、
その往復は標準形の写像に縮約しない。
その場合は時刻の刻みを細かく取り、戻りを次の時刻の入力として表す。
時刻の刻みの取り方は記述上の選択であり、
本論文は、そのような刻みがどの実現にも存在することを主張しない。
刻みの取り方を変えると一段遷移仕様が変わり、
可行行動集合と閉塞の判定もそれに相対して定まる
（\S\ref{sec:acs_def} 備考~\ref{rem:blockage_scope_sec2} (c)）。
連続時間の記述との対応は Phase~2 の課題とする（\S\ref{subsec:limitations_list} の L1）。
\end{remark}

\begin{remark}[複数の情報の並行的な処理の範囲]
\label{rem:multi_item_concurrency}
同一時刻の識別状態は複数の内部対象トークンを、
要請信号と要求×対象の評価は複数の要求と組を、
目的の集合は複数の目的を、それぞれ含みうる。
生成依存 $M_1\to M_2\to M_3$ は、
ある情報の処理の完了を他の情報の処理の条件とする規則を含まない。
したがって本論文は、標準形の上では、
一つの情報が未処理であることそれ自体から、他の情報の処理の停止を一般的な帰結として置かない。
実装では、複数の情報が異なる処理の段階に同時にありうる。

一方、実装の処理能力や仕様によって、
個々の情報の処理が遅れること、または処理されないまま残ることは起こりうる。
標準形の上では、これは各機構の写像が当該の情報を出力に反映しないこととして表れる
（たとえば、トークンを割り当てないこと、要求×対象の評価の定義域に含めないこと）。
本論文は、すべての情報がいつか処理されること、処理の遅れに上限があること、
および処理の公平性を一般的な帰結として置かない。
トークンの未割当、要求×対象の未評価、目的の状態判断の保留は
それぞれ異なる水準の量であり、本論文はこれらを一つの未処理の状態としてまとめない。
\end{remark}

\begin{definition}[機構出力の個体への帰属（分類B）]
\label{def:s4_output_attribution}
本節の M1、M2、M3 の時刻別の出力は、
$\mathrm{WF}_{i,t}=1$（\S\ref{sec:acs_def} 定義~\ref{def:wf_realization}）の時刻 $t$ について
$\Sigma_i$ の出力として定める。
$\mathrm{WF}_{i,t}=0$ の時刻には、これらを $\Sigma_i$ の出力として定めない。
このとき出力は、特定の値（$0$、$\mathsf{und}$、空集合その他）をとるのではなく、未定義である。
崩壊の後の時刻はこの場合に含まれる（同節定義~\ref{def:collapse} の帰属の規定）。
各機構の写像は型の上で全域であり、この規定は写像の定義域ではなく、
時刻別の値を個体に帰属させるかどうかに関わる。
機能軸（定義~\ref{def:functional_axis}）その他の分類Cの記述量は本規定の対象ではない。
\end{definition}

\begin{remark}[離散時刻と内部処理工程]
\label{rem:internal_steps_time}
本節の $M_1\to M_2\to M_3$ と各時刻の標準写像は、
実装内部の個々の計算、神経処理、反復比較、再評価または確認操作を一対一に表すものではない。
一つの時刻の出力の形成には多数の内部工程が含まれうる。
また、一つの目的の形成から遂行と達成の判断までは、複数の時刻にまたがりうる。
本論文は、内部工程の数、処理時間および反復回数を規定せず、その最低値も置かない。
一つの機構の内部の反復・再比較・再統合は、その機構の時刻別の写像に縮約して表す。
同一時刻の中での機構の間の往復は縮約しない（備考~\ref{rem:tipa_bidirectional}）。
長期にわたる目的を持つかどうか、および期間にわたる持続を達成の条件とする目的を持つかどうかは、
実装依存である。
\end{remark}


\subsubsection{M1：情報統合機構}
\label{subsec:m1}

M1 は、PSL-1 を介して利用可能となる感知情報、
内部モデル、有限履歴、前時刻の評価・行動情報その他を統合し、
M2 以降から利用可能な現在状態・外部環境状態についての内部表現を構成する。

有限履歴長を表す実装依存の正整数を
\[
  W\in\mathbb N,
  \qquad
  W_{\mathrm{short}}\in\mathbb N,
  \qquad
  1\le W_{\mathrm{short}}\le W
\]
とする。
$X^{\le W}$ は長さ高々 $W$ の $X$ 値有限履歴の集合、
$\mathcal I^{\le W_{\mathrm{short}}}$ は長さ高々 $W_{\mathrm{short}}$ の
$\mathcal I$ 値有限履歴の集合を表す。

また、M1-b で成分構造を具体化する識別状態空間として、
あらかじめ非空集合 $\Xi$ を置く。
動態重みの値域として標準単体
\[
  \Delta^{n-1}
  :=
  \left\{
    w\in[0,1]^n
    \;\middle|\;
    \sum_{k=1}^{n}w_k=1
  \right\}
\]
を用いる。

\begin{definition}[M1内部状態（分類B）]
\label{def:m1_state}
M1 が時刻 $t$ に生成する標準出力を束ねた内部状態を
\[
  \mathsf S_{M1,t}
  :=
  \bigl(
    \hat x_t,
    \hat e_t,
    \xi_t,
    w_t,
    D_{M1,t}
  \bigr)
\]
とし、その型を
\[
  \mathsf S_{M1}
  :=
  X
  \times E_{\mathrm{ext}}
  \times\Xi
  \times\Delta^{n-1}
  \times\mathcal D_{M1}
\]
とする。

$\mathsf S_{M1,t}$ は M1 の既定義出力を束ねる標準内部表現であり、
独立した追加状態変数または別機構を導入するものではない。
前時刻の $\mathsf S_{M1,t-1}$ を次時刻の M1 が参照することを許すが、
学習・可塑的変化その他の追加的な時刻間更新則は
\S\ref{sec:extended_dynamics} に委ねる。
\end{definition}

\begin{definition}[情報統合機構 $M_1$（分類B）]
\label{def:m1_map}
感知情報空間を $\mathcal I$、
内部モデル状態空間を $\mathcal M$、
M1内部状態空間を $\mathsf S_{M1}$、
M2用時間特徴量空間を $\mathcal D_{M1}$ とする。

M1 を
\[
  \mathrm{M1}:
  \mathcal I
  \times\mathcal M
  \times\mathcal P_{\mathrm{prio}}
  \times U
  \times X^{\le W}
  \times\mathsf S_{M1}
  \to
  X
  \times E_{\mathrm{ext}}
  \times\Xi
  \times\Delta^{n-1}
  \times\mathcal D_{M1}
\]
なる決定論的標準機能形式とする。

\[
\begin{aligned}
(&\hat x_t,\hat e_t,\xi_t,w_t,D_{M1,t})
=
\mathrm{M1}\bigl(
I_{\mathrm{sense},t},
M_t,
P_{t-1},
u_{t-1},
\{\hat x_s\}_{s\in[\max(0,t-W),t-1]},
\mathsf S_{M1,t-1}
\bigr).
\end{aligned}
\]

ここで
\[
  \hat x_t\in X,
  \qquad
  \hat e_t\in E_{\mathrm{ext}},
  \qquad
  \xi_t\in\Xi,
  \qquad
  w_t\in\Delta^{n-1},
  \qquad
  D_{M1,t}\in\mathcal D_{M1}.
\]
前時刻入力 $u_{t-1}$ は最終実効行動について保持された記録であり、
M3 行動指令の記録とは区別する。実際の前時刻の行動との一致は別条件である。

$t=0$ の前時刻記号には、定義~\ref{def:s4_carrier_contract} に従い
$M_0^{\mathrm W}$ の読み出し値を用いる。これは負時刻の実現を仮定しない。
\end{definition}

\paragraph*{4.2.1-a\quad M1-a：情報統合}

M1-a は、現在利用可能な情報と時間的文脈を統合し、
現在状態および外部環境状態についての機構層推定を形成する。

\noindent\textbf{外部環境状態の型。}
\S\ref{sec:acs_def} の型規約に従い、
実現外部環境状態および M1 の環境推定は
$e_t,\hat e_t\in E_{\mathrm{ext}}$
とする。
必要環境要因の集合 $E$ を外部環境状態の型として用いない。

\begin{definition}[M1-aの予備状態推定（分類B）]
\label{def:std_init}
予備状態推定写像を
\[
  \Phi_{\mathrm{pre}}:
  \mathcal I\times\mathcal M\times\mathsf S_{M1}
  \to X
\]
とし、
\[
  \tilde x_t
  :=
  \Phi_{\mathrm{pre}}
  \bigl(
    I_{\mathrm{sense},t},
    M_t,
    \mathsf S_{M1,t-1}
  \bigr)
  \in X
\]
とする。
\end{definition}

外部環境状態推定については、
\[
  \Phi_{\mathrm{env}}:
  \mathcal I\times\mathcal M\times X^{\le W}
  \to E_{\mathrm{ext}}
\]
を用いて
\[
  \hat e_t
  =
  \Phi_{\mathrm{env}}
  \bigl(
    I_{\mathrm{sense},t},
    M_t,
    \{\hat x_s\}
  \bigr)
\]
としてよい。

\begin{definition}[M1-aの動態重みと最終統合（分類B）]
\label{def:m1_integration}
動態重みを $w_t\in\Delta^{n-1}$ とし、
最終統合写像を
\[
  \Phi_{\mathrm{integrate}}:
  X\times\mathcal M\times X\times\Delta^{n-1}
  \to X
\]
として、
\[
  \hat x_t
  =
  \Phi_{\mathrm{integrate}}
  \bigl(
    \tilde x_t,
    M_t,
    \bar x_{\mathrm{base},t},
    w_t
  \bigr)
  \in X
\]
とする。
\end{definition}
\begin{definition}[M1補助出力の情報源（分類B）]
\label{def:s4_m1_sources}
予備推定等から最終統合へ用いる補助量は、固定写像
$\Phi_{\mathrm{weight}}:\mathcal M\times\mathcal I\to\Delta^{n-1}$、
$\Phi_{\mathrm{baseline}}:\mathcal M\times\mathcal I\to X$ により
$w_t=\Phi_{\mathrm{weight}}(M_t,I_t)$、
$\bar x_{\mathrm{base},t}=\Phi_{\mathrm{baseline}}(M_t,I_t)$ とする。
$\Phi_{\mathrm{baseline}}$ は、$M_t$ の作動成分から $\rho_{\mathrm{past}}$ により読み出す状態推定履歴（第3成分）に、
定義~\ref{def:m1_temporal_base} の履歴集約写像 $\Phi_{\mathrm{hist}}$ を適用する合成写像とする。
一段予測を用いる場合、
$\Phi_{\mathrm{predict}}:\mathcal M\to X^{X\times U\times E_{\mathrm{ext}}}$
を固定し、$\hat f_t=\Phi_{\mathrm{predict}}(M_t)$ とする。
右辺の指数記法は、$X\times U\times E_{\mathrm{ext}}$ から $X$ への関数全体の集合を表し、連続性は要求しない。
予測に用いる前時刻推定も保持記録から得る。
\end{definition}


\begin{remark}[状態・環境推定と実現値]
\label{rem:m1_estimate_nontruth}
$\hat x_t$ および $\hat e_t$ は内部推定であり、
$x_t$ および $e_t$ と一致することを一般には要求しない。
また、Phase~1 の本節では普遍的な推定誤差上界または収束性を要求しない。
\end{remark}
\begin{definition}[状態推定誤差（分類C）]
\label{def:s4_state_estimation_error}
$\mathrm{WF}_{i,t}=1$ の時刻について、
$X\subseteq\mathbb R^n$ の制限距離 $d_X(x,x')=\|x-x'\|_\infty$ を用い、
$\varepsilon_{\mathrm{est},t}:=d_X(x_t,\hat x_t)$ と定める。
これは真の状態との照合量であり、内部の予測整合性誤差
$\varepsilon_{\mathrm{pred},t}$ とは別である。
内部計算可能性や外部測定可能性を要求しない。
\end{definition}


\paragraph*{4.2.1-b\quad M1-b：境界管理}

M1-b は、自己--環境境界その他の応答形成に必要な識別状態を、
複数の抽象的識別軸によって構成する。

\begin{definition}[識別状態と内部対象トークン（分類B）]
\label{def:id_vector}
識別軸の添字集合を
\[
  \mathcal R_\xi
  :=
  \{
    \mathrm{bnd},
    \mathrm{tmp},
    \mathrm{org},
    \mathrm{dyn},
    \mathrm{rep},
    \mathrm{mval},
    \mathrm{prox},
    \mathrm{tgt}
  \}
\]
とする。
各 $\kappa\in\mathcal R_\xi$ に未決定カテゴリを含む有限非空集合 $\Xi_\kappa$ を置き、
一つの行の値域を
\[
  \Xi^{(8)}:=\prod_{\kappa\in\mathcal R_\xi}\Xi_\kappa
\]
とする。

ただし、\S\ref{sec:acs_def} の自己--環境境界との接続に用いる
境界帰属軸については、上流の型インターフェースを保つため
\[
  \Xi_{\mathrm{bnd}}
  =
  \{\mathsf{self},\mathsf{env},\mathsf{und}\}
\]
とする。
ここで $\mathsf{self}$、$\mathsf{env}$、$\mathsf{und}$ は、
それぞれ自己側への帰属、環境側への帰属、未決定を表す標準カテゴリである。
他の七軸の具体的カテゴリ分解は Phase~1 では一般に固定しない。

内部対象トークンのラベル集合として、有限非空集合 $\mathsf{Lab}$ を固定する。
以下で全体の行の添字に用いる記号 $\ast$ と、自己を表す記号 $\mathsf{self}$ は $\mathsf{Lab}$ に属さないものとし、
$\mathsf{Lab}\cap\bigl(\{\mathsf{self}\}\times\mathsf{Lab}\bigr)=\emptyset$ とする。
その大きさは実装依存とする。
識別状態空間を
\[
  \Xi
  :=
  \Xi^{(8)}
  \times
  \bigcup_{\mathcal K\subseteq\mathsf{Lab}}
  \bigl(\Xi^{(8)}\bigr)^{\mathcal K}
\]
とし、M1-b の識別状態を
\[
  \xi_t
  =
  \bigl(
    \xi_t(\ast),\
    (\xi_t(k))_{k\in\mathcal K_t}
  \bigr)
  \in\Xi
\]
と書く。
第1成分 $\xi_t(\ast)\in\Xi^{(8)}$ を\emph{全体の行}、
第2成分の添字集合 $\mathcal K_t\subseteq\mathsf{Lab}$ を時刻 $t$ の\emph{内部対象トークン集合}、
各 $\xi_t(k)\in\Xi^{(8)}$ をトークン $k$ の\emph{行}と呼ぶ。
$\mathcal K_t=\emptyset$ を許す。
全体の行の第 $\kappa$ 成分を $\xi_{\kappa,t}$、トークン $k$ の行の第 $\kappa$ 成分を $\xi_{\kappa,t}(k)$ と書く。
とくに \S\ref{sec:acs_def} が参照する $\xi_{\mathrm{bnd},t}$ は、全体の行の境界帰属成分である。

各軸はそれぞれ、
境界帰属、時間的位置、情報由来、動的性、表象基盤、
記憶照合上の有益性・有害性、相互作用近接性、目的対象性を表す。
具体的カテゴリ、特徴抽出規則および内部符号化は実装依存とする。
\end{definition}

\begin{remark}[内部対象トークンの形成と実装依存の範囲]
\label{rem:token_generation}
$\mathcal K_t$ は M1 の出力 $\xi_t$ の一部であり、
M1 の入力（感知情報 $I_{\mathrm{sense},t}$、内部モデル $M_t$ および保持記録）から同一時刻に形成されうる。
この形成は M2・M3 の同一時刻の出力を用いないので、同一時刻の循環を生じない。
本論文は次を要求しない。
\begin{enumerate}[label=(\alph*)]
  \item 外部に存在する対象のすべてにトークンが割り当てられること。
  \item 新たに現れた対象に、それが感知された時刻にトークンが割り当てられること。
  \item トークンの行の各成分が対象の性質と一致すること。
\end{enumerate}
一時刻に形成されうるトークンの数、割当ての遅れと精度、
複数の対象を一つずつ識別するか同時に識別するかは実装依存である。
$\mathcal K_t=\emptyset$ の場合にも全体の行 $\xi_t(\ast)$ は定まり、
M2 と M3 はそれを用いて評価と応答形成を行いうる。
ラベル $k\in\mathsf{Lab}$ は内部の識別の単位であり、
\S\ref{sec:acs_def} の枠 $a\in\mathcal A_O$ や外部の個物と同一視しない。
同じラベルを時刻をまたいで再び用い、異なる対象を表してもよい。
トークンが時刻をまたいで持続することは、
前時刻の $\mathsf S_{M1,t-1}$ を入力とする M1 が同じラベルを与えることとして表し、
その持続の規則は \S\ref{sec:extended_dynamics} に委ねる。
\end{remark}

\noindent\textbf{境界帰属。}
対象または情報を自己側、環境側、未決定のいずれとして扱うかを表す。

\noindent\textbf{時間的位置。}
対象または情報が現在、過去、未決定のいずれに関係するかを表す。

\noindent\textbf{情報由来。}
情報が内部由来、外部由来、未決定のいずれであるかを表す。

\noindent\textbf{動的性。}
対象を相互作用上の動的対象、非動的対象、未決定のいずれとして識別する。

\noindent\textbf{表象基盤。}
表象が観測由来、内部生成由来、未決定のいずれであるかを表す。

\noindent\textbf{記憶照合上の有益性・有害性。}
現在情報と過去経験との記憶照合上の関係を表し、
現在の可行性評価そのものとはしない。

\begin{remark}[記憶照合上の有益性・有害性とM2評価の分離]
\label{rem:mval_identification}
本論文は、過去に有益または有害であった対象との類似性だけから、
現在の Priority、可行性脅威または行動要請を一般的な帰結として置かない。
現在の可行性関連評価は M2 の責務である。
\end{remark}

\noindent\textbf{相互作用近接性。}
対象との相互作用が現在の ACS にとって成立しうる程度を表し、
単なる物理距離とは同一視しない。

\noindent\textbf{目的対象性。}
対象または情報が、時刻 $t$ の開始時点で保持されている記録のうえで、
当該ACSの目的の向き先になっているか、なっていないか、または未決定であるかを表す。
同じ時刻に M3 が形成する目的 $\mathrm{Pur}_t$（定義~\ref{def:purpose_formation}）は参照しない。
時刻 $t$ に新たに目的の向き先となったトークンは、記録を通じて時刻 $t+1$ 以後の本軸に現れうる。
記録の正確さは一般には要求しない（定義~\ref{def:s4_carrier_contract}）。
M2 はこの軸を評価に用いない（定義~\ref{def:m2_map}）。

\begin{remark}[識別状態の役割と非真値性]
\label{rem:id_vector_role}
$\xi_t$ は M1 の標準内部表現であり、
ACS の存在論的な自己--環境境界、外部対象の客観的属性、
または三原理の時刻別充足述語の判定そのものではない。
\end{remark}

\paragraph*{4.2.1-c\quad M1-c：時間的連続性統合}

M1-c は、有限履歴と内部モデルによる予測を利用し、
現在処理に必要な時間的文脈および時間特徴量を構成する。

\begin{definition}[M1-cの時間的基底表現（分類B）]
\label{def:m1_temporal_base}
履歴集約写像を
\[
  \Phi_{\mathrm{hist}}:
  X^{\le W}\to X
\]
とし、
\[
  \bar x_{\mathrm{base},t}
  :=
  \Phi_{\mathrm{hist}}
  \bigl(\{\hat x_s\}\bigr)
  \in X
\]
とする。
\end{definition}

\noindent\textbf{一段階予測写像の型。}
\S\ref{sec:acs_def} の外部環境状態型に従い、
内部モデルが一段階予測を用いる場合には
\[
  \hat f_t:
  X\times U\times E_{\mathrm{ext}}
  \to X
\]
とする。
前時刻の最終実効行動 $u_{t-1}$ と環境推定 $\hat e_{t-1}$ を用いる場合、
\[
  \tilde x_t^{-}
  :=
  \hat f_t
  \bigl(
    \hat x_{t-1},
    u_{t-1},
    \hat e_{t-1}
  \bigr)
\]
とする。

\begin{definition}[M1-cの予測整合性と先行時間特徴量（分類C）]
\label{def:m1_temporal_features}
状態空間上の距離 $d_X$ を用いる場合、
予測整合性誤差を
\[
  \varepsilon_{\mathrm{pred},t}
  :=
  d_X(\tilde x_t,\tilde x_t^{-})
  \in\mathbb R_{\ge0}
\]
として記述してよい。

これらは分類Cの機構層記述量であり、
予測の正しさや、三原理の時刻別充足述語の真偽を表さない。
\end{definition}

\begin{definition}[外部側の集合値予測（分類B）]
\label{def:env_set_prediction}
\S\ref{sec:acs_def} の分解 $E_{\mathrm{ext}}=E^{\mathrm{bg}}\times O$
（定義~\ref{def:Eext_decomposition}）に従い、
外部環境状態とその推定を $e=(b,o)$、$\hat e_t=(\hat b_t,\hat o_t)$ と書く。
外部側の予測値の集合を
\[
  \mathcal G_O:=\bigl(\mathcal P(O)\setminus\{\emptyset\}\bigr)\cup\{\bot\}
\]
とする。$\bot$ は予測を与えないことを表し、空集合ではない。
内部モデルが行動応答成分の予測を用いる場合、固定写像
$\Phi_{\mathrm{pred},O}:\mathcal M\to\mathcal G_O^{\,X\times U\times E_{\mathrm{ext}}}$
により $\hat g_t:=\Phi_{\mathrm{pred},O}(M_t)$ とし、行動 $u\in U$ についての予測集合を
\[
  \hat G_t(u):=\hat g_t(\hat x_t,u,\hat e_t)\in\mathcal G_O
\]
とする。
$\hat G_t(u)$ は、次の時刻の行動応答成分として起こりうると内部で見込まれる値の集合であり、
確率や頻度を表さない。
$\hat g_t$ は状態側の一段予測 $\hat f_t$ とは別の写像である。
\end{definition}

\begin{definition}[予測集合の記録と二種類の予測外（分類C）]
\label{def:prediction_record_check}
時刻 $t$ の実効行動 $u_t$ が定まった後、
行動した時点の予測集合 $\hat G_t(u_t)$ が保持記録に残され、
時刻 $t+1$ にはその記録値が $\hat G^{\mathrm{rec}}_{t}\in\mathcal G_O$ として
定義~\ref{def:s4_carrier_contract} の読み出しから得られるとする。
記録値は $\hat G_t(u_t)$ との一致を仮定しない。
次の二つを定める。
\begin{enumerate}[label=(\roman*)]
  \item \emph{内部の予測外}：$\mathrm{WF}_{i,t+1}=1$ かつ $\hat G^{\mathrm{rec}}_{t}\neq\bot$ のとき、
  $\hat o_{t+1}\notin\hat G^{\mathrm{rec}}_{t}$ であること。
  これは時刻 $t+1$ の M1-c が利用可能な量だけで定まり、
  定義~\ref{def:m1_temporal_features} の $\varepsilon_{\mathrm{pred},t}$ と同じ水準の記述量である。
  \item \emph{理論上の予測外}：$\mathrm{WF}_{i,t}=1$ かつ $\hat G_t(u_t)\neq\bot$ のとき、
  $o_{t+1}\notin\hat G_t(u_t)$ であること。
  これは実現値との照合量であり、当該ACSの入力ではない
  （定義~\ref{def:s4_state_estimation_error} の $\varepsilon_{\mathrm{est},t}$ と同じ扱い）。
\end{enumerate}
予測が $\bot$ の場合はそれぞれ定めず、偽に置換しない。
両者は一致しなくてよい。
たとえば、対象の誤認により $\hat o_{t+1}$ が $o_{t+1}$ と異なる場合には、一方だけが成り立ちうる。
\end{definition}

\begin{remark}[二つの予測の照合基準の違い]
\label{rem:prediction_reference_difference}
状態側の $\tilde x_t^{-}$ は、時刻 $t$ の予測写像 $\hat f_t$ により前時刻の記録を計算し直した値である。
これに対し外部側では、行動した時点の予測集合の記録と照合する。
内部モデルが時刻 $t$ から $t+1$ の間に更新された場合、
計算し直した予測と行動した時点の予測は一致しなくてよい。
本論文は、前者を後者で置き換えることも、その逆も要求しない。
予測外にもとづく内部モデルの更新の規則は与えない（\S\ref{sec:extended_dynamics}）。
\end{remark}

\begin{definition}[M1-cの最終時間特徴量（分類B）]
\label{def:m1_temporal_output}
M1-c が M2 に提供する時間特徴量を
\[
  D_{M1,t}\in\mathcal D_{M1}
\]
とする。
$\mathcal D_{M1}$ の内部成分には、
状態変化、入力変化、予測差、変化傾向その他を実装に応じて含めてよい。
\end{definition}

\noindent\textbf{M1内部の生成依存。}
概念的には
\[
  \text{現在入力・履歴}
  \longrightarrow
  \text{予備推定・時間的文脈}
  \longrightarrow
  \text{最終統合}
  \longrightarrow
  \text{識別・時間特徴}
\]
として整理する。
これは物理的逐次実行を要求しない。

\begin{assumption}[M1-cの時間統合構造]
\label{ass:m1c_temporal}
M1-c は、有限履歴を現在処理へ接続し、
M2 から利用可能な時間特徴量を形成できるものとする。

導入地位：S（構造的）。強さ：中。
本仮定は、M1-c の標準形式（定義~\ref{def:m1_temporal_base}、定義~\ref{def:m1_temporal_output}）を適用する ACS の範囲を限定する構造的条件であり、本節の命題の前提としては用いない。
\end{assumption}

\begin{remark}[M1の責務境界]
\label{rem:m1_responsibility}
M1 は状態・環境推定、識別および時間統合を担うが、
Priority 生成、要請信号の生成、目的の形成、動作モードの判断、行動候補比較、
または三原理の時刻別充足述語の判定を担当しない。
\end{remark}


\subsubsection{M2：適応的評価機構}
\label{subsec:m2}

M2 は、M1 が形成した現在像および時間特徴を受け取り、
可行性脅威、構造的脆弱性、および潜在的不安定性に関係する差異を
機構層で代理評価し、M3 の応答形成から利用可能にする。

本論文は M2 に、理論上の
$\Pi$、$K$、$A_t$、$E_{\mathrm{req}}(x_t)$、$\mathrm{Crit}^{(1)}(x_t)$、
$\mathcal{D}_{\mathrm{FR}}$、$\mu_{\mathrm{eff},t}$、頑健存続可能段数 $\mathrm{Hor}_f$、$Q$ その他の
真値へのアクセス、内部での保持、または正確な計算を要求しない。
M2 の評価は、M1 等から利用可能な情報に依存する
（\S\ref{sec:three_principles} 備考~\ref{rem:theory_quantity_vs_internal}）。

\begin{definition}[M2内部評価情報空間（分類B）]
\label{def:m2_eval_space}
M2 が M3 へ提供する内部評価情報の型を非空集合
\[
  \mathfrak E
\]
とし、
\[
  \mathrm{ev}_t\in\mathfrak E
\]
とする。
$\mathfrak E$ の内部表現形式は Phase~1 では一意に固定しない。
\end{definition}

\begin{definition}[要求の一覧と要請信号の型（分類B）]
\label{def:request_signal_types}
要求ラベルの非空集合 $\mathsf{Req}$ と、要請の質の非空集合 $\mathsf{Qual}$ を置く。
固定した読み出し $\rho_{\mathrm{req}}:\mathcal M\to\mathcal P_{\mathrm{fin}}(\mathsf{Req})$ により、
時刻 $t$ の\emph{要求の一覧}を
\[
  \mathrm{Req}_t:=\rho_{\mathrm{req}}(M_t)
\]
とする。
要請信号の空間を
\[
  \mathsf{Sig}:=\mathcal P_{\mathrm{fin}}(\mathsf{Req}\times\mathsf{Qual})
\]
とし、M2 が出力する\emph{要請信号}を $\mathrm{sig}_t\in\mathsf{Sig}$、
標準形では $\mathrm{sig}_t\subseteq\mathrm{Req}_t\times\mathsf{Qual}$ とする。
$(q,c)\in \mathrm{sig}_t$ は、要求 $q$ が質 $c$ の要請として時刻 $t$ に立ち上がっていることを表す。
$\mathsf{Qual}$ の元の列挙と抽象的な分類は Phase~1 では与えず、Phase~1.$\alpha$ の課題とする。
記号 $\mathsf{Qual}$ は \S\ref{sec:rel_min} の関係品質 $Q_{i\leftarrow j,t}$ とは無関係である。
要請信号 $\mathrm{sig}_t$ は、\S\ref{sec:three_principles} の三原理評価ベクトル $s_t$（定義~\ref{def:S3}）とは別の量である。
\end{definition}

\begin{definition}[要求×対象の評価の型（分類B）]
\label{def:rq_eval_types}
評価値の非空集合 $\mathfrak E^{\mathrm{rq}}$ を置き、
\[
  \mathfrak E^{\mathrm{rq}}_{\mathrm{map}}
  :=
  \bigcup_{D}
  \bigl(\mathfrak E^{\mathrm{rq}}\bigr)^{D}
\]
とする。ここで $D$ は $\mathsf{Req}\times(\mathsf{Lab}\cup\{\ast\})$ の有限部分集合を動く。
M2 が出力する\emph{要求×対象の評価}を $\mathrm{ev}^{\mathrm{rq}}_t\in\mathfrak E^{\mathrm{rq}}_{\mathrm{map}}$ とし、
その定義域を $D^{\mathrm{rq}}_t$ と書く。
標準形では $D^{\mathrm{rq}}_t\subseteq\mathrm{Req}_t\times(\mathcal K_t\cup\{\ast\})$ とする。
$(q,k)\in D^{\mathrm{rq}}_t$ はトークン $k$ が要求 $q$ に照らして評価されたことを、
$(q,\ast)$ は対象を介さずに評価されたことを表す。
定義域に含まれない組は\emph{未評価}であり、低い評価値として扱わない。
\end{definition}

\begin{definition}[適応的評価機構 $M_2$（分類B）]
\label{def:m2_map}
M2 を
\[
\begin{aligned}
\mathrm{M2}:&
X\times E_{\mathrm{ext}}\times\Xi
\times\Delta^{n-1}\times\mathcal D_{M1}
\times\mathcal M\times X^{\le W}
\times\mathcal I^{\le W_{\mathrm{short}}}
\\
&\to
\mathcal P_{\mathrm{prio}}
\times\mathfrak E
\times\mathsf{Sig}
\times\mathfrak E^{\mathrm{rq}}_{\mathrm{map}}
\end{aligned}
\]
なる決定論的標準写像とする。

\[
\begin{aligned}
(&P_t,\mathrm{ev}_t,\mathrm{sig}_t,\mathrm{ev}^{\mathrm{rq}}_t)
=
\mathrm{M2}\bigl(
\hat x_t,\hat e_t,\xi_t,w_t,D_{M1,t},M_t,
\{\hat x_s\},\{I_{\mathrm{sense},s}\}
\bigr).
\end{aligned}
\]
M2 は、M3 の出力と M3 の目的基準の評価（定義~\ref{def:eval_vs_judgment}）を、
同一時刻のものも保持記録を通じたものも入力としない。
標準形では、M2 の値は識別状態の目的対象性の成分
$\xi_{\mathrm{tgt},t}$ と $\xi_{\mathrm{tgt},t}(k)$ に依存せず、
$M_t$ の成分のうち M3 の出力と目的基準の評価の保持記録を担う部分にも依存しない。
\end{definition}

\begin{remark}[M2 に戻らないことの範囲]
\label{rem:m2_no_return_scope}
定義~\ref{def:m2_map} の条件が禁じるのは、
M3 の出力、目的基準の評価、判断とそれらの保持記録を M2 の入力とすること、
および識別状態の目的対象性の成分を用いることである。
この条件は、M3 から生じた情報が後の時刻の M2 に影響しないことを意味しない。
本論文は次の二つの間接的な影響を排除しない。
\begin{enumerate}[label=(\alph*)]
  \item 実効行動が状態と外部環境状態を変え、
  その結果が感知と M1 の統合を経て後の時刻の M2 の入力に現れること。
  内的認知状態を意図的に変化させる行動 $u_t\in U_{\mathrm{int}}$
  （\S\ref{sec:acs_def} 備考~\ref{rem:action_space_structure}）の作用もこれに含まれる。
  \item M1 が保持内容 $M_t$（目的の集合の記録を含む）を用いて注意の配分その他の統合を行い、
  その出力が M2 の入力となること。
  M1 の写像は $M_t$ を入力とし（定義~\ref{def:m1_map}）、
  $\mathrm{Update}_2$ は $\mathcal M_1$ を引数とする
  （\S\ref{sec:acs_def} 定義~\ref{def:update_operator}）。
\end{enumerate}
M2 の評価は、こうして与えられた材料について存続基準によって行われ、目的を基準としない。
M3 の目的基準の評価・判断と M2 の評価との区別は評価の基準の区別であり、
情報の因果的な隔離ではない。
最終実効行動の記録 $u_{t-1}$ は M3 行動指令の記録とは区別され（定義~\ref{def:m1_map}）、
M3 の出力の記録には当たらない。
(a)・(b) の影響が実際に生じるかどうかは $f$、$\mathrm{Update}$、感知および M1 の実現に依存し、
本論文はそれを一般的な帰結として置かない。
\end{remark}

\paragraph*{4.2.2-a\quad M2-a：可行性脅威検出}

M2-a は、現在または近接した将来における急性の可行性脅威を、
利用可能な機構層情報から評価する分類Bの標準機能である。

\begin{definition}[M2-a可行性脅威検出機能（分類B）]
\label{def:m2a_threat}
M2-a が用いる機構層の可行性代理評価を、
実装依存の決定論的内部評価写像
\[
  \widehat{\Pi}_{M2}:
  X\times\mathcal M
  \to
  \mathbb R
\]
により
\[
  \Pi_{\mathrm{eval},t}
  :=
  \widehat{\Pi}_{M2}
  \bigl(
    \hat x_t,
    M_t
  \bigr)
\]
とする。
$\Pi_{\mathrm{eval},t}$ は真の $\Pi(x_t)$ と一致することを要求せず、
さらに ACS が理論上の $\Pi$ を内部で既知として
$\Pi(\hat x_t)$ を直接計算することも要求しない。
これは、\S\ref{subsec:description_levels_global} の規約に従う
機構層の代理評価である。

急性脅威判定を
\[
  \Psi_{\mathrm{acute}}:
  X\times E_{\mathrm{ext}}\times\Xi\times\mathcal D_{M1}\times\mathcal M
  \to\{0,1\}
\]
により
\[
  C_{1,t}
  :=
  \Psi_{\mathrm{acute}}
  \bigl(
    \hat x_t,\hat e_t,\xi_t,D_{M1,t},M_t
  \bigr)
\]
とする。
$C_{1,t}=1$ は M2 が急性脅威を検出したことを表し、
SM の時刻別充足述語の真偽とは同一視しない。
\end{definition}
\begin{definition}[内部可行判定と理論上の可行判定（分類C）]
\label{def:s4_internal_viability_check}
内部の判定規則を採用する場合、固定写像
$\Psi_K:\mathcal M\times\mathcal R_{\mathrm{int}}\to\{0,1,\bot\}$
で内部判定を記述する。$\bot$ は未判定である。
$\Psi_K(M_t,r_t)\in\{0,1\}$ の場合だけ、
$\mathbf1[\Pi(x_t)\ge0]$ との一致・不一致を比較する。
未判定を不一致や偽に数えず、この比較を M2 の入力として要求しない。
$\Pi_{\mathrm{eval},t}$ の符号がこの判定と一致することも仮定しない。
\end{definition}


\paragraph*{4.2.2-b\quad M2-b：構造的脆弱性評価}

M2-b は、現在の脅威または要請に対する応答構造の不足を、
対応不可能性、手段不足性、資源不足性、支援不足性、
関係品質劣化の観点から代理評価する分類Bの標準機能である。
ここでいう関係は、\S\ref{sec:rel_min} が扱う ACS 関係の最小構造を指し、
本節はその分類および関係記述子を再定義しない。
類型の名称の「関係品質劣化」は、外部関係に伴う可行性関連差異についての内部の評価の類型を指し、
\S\ref{sec:rel_min} の関係品質 $Q_{i\leftarrow j,t}$（記述者の量）の代理値ではない
（備考~\ref{rem:Q_internal_formal}）。

\begin{definition}[機構層有効行動多様度の一般形（分類C）]
\label{def:mech_mu_eff}
機構層で利用可能と推定された行動集合 $\hat A\subseteq U$ と、
その上の決定論的同値関係 $\hat\sim$ に対し、
\[
  \hat\mu_{\mathrm{eff}}
  \bigl(
    \hat x;\hat A,\hat\sim
  \bigr)
  :=
  |\hat A/\hat\sim|
\]
とする。
これは \S\ref{sec:three_principles} の応答多様度 $\mu_{\mathrm{eff},t}$（定義~\ref{def:mu_eff}）の真値とは別の分類Cの代理量である。
$\mu_{\mathrm{eff},t}$ 自体も診断量であり、いずれの原理の時刻別充足述語の定義にも用いない。
\end{definition}

\begin{definition}[M2-b構造的脆弱性評価機能（分類B）]
\label{def:m2b_vuln}
M2-b が扱う標準脆弱性を
\begin{enumerate}
  \item 対応不可能性、
  \item 手段不足性、
  \item 資源不足性、
  \item 支援不足性、
  \item 関係品質劣化
\end{enumerate}
の五類型とする。

脆弱性判定写像を
\[
  \Psi_{\mathrm{vuln}}:
  X\times E_{\mathrm{ext}}\times\Xi\times\mathcal M
  \to\{0,1\}^{5}
\]
とし、検出状態を
\[
  \mathbf C_{2,t}
  :=
  \Psi_{\mathrm{vuln}}
  \bigl(
    \hat x_t,\hat e_t,\xi_t,M_t
  \bigr)
  \in\{0,1\}^{5}
\]
とする。
各成分の具体的判定規則は実装依存とする。
\end{definition}

M2 が粗い行動可能性推定を用いる場合には、
\[
  \mathrm{Available}_{M2}:
  X\times E_{\mathrm{ext}}\times\mathcal M
  \to\mathcal P_{\mathrm{fin}}(U)
\]
により
\[
  \hat A_{M2,t}
  :=
  \mathrm{Available}_{M2}
  \bigl(
    \hat x_t,\hat e_t,M_t
  \bigr)
\]
を構成してよい。
$\hat A_{M2,t}$ を真の可行行動集合 $A_t$ と同一視しない。

M2 が予測応答同値関係 $\hat\sim_{M2,t}$ を利用する場合、
\[
  \hat\mu_{\mathrm{eff},M2,t}
  :=
  \hat\mu_{\mathrm{eff}}
  \bigl(
    \hat x_t;
    \hat A_{M2,t},
    \hat\sim_{M2,t}
  \bigr)
\]
とする。
非空集合 $U$ 上の同値関係全体を $\operatorname{Eq}(U)$ とする。
固定写像 $\Phi_{\mathrm{eq}}:\mathcal M\to\operatorname{Eq}(U)$ を与え、
$\hat\sim_{M2,t}$ は $\Phi_{\mathrm{eq}}(M_t)$ の $\hat A_{M2,t}$ への制限とする。
空の行動推定集合の場合、その商の基数は $0$ である。
具体的な同値判定規則は実装依存である。

\paragraph*{4.2.2-c\quad M2-c：潜在的不安定性の機構層評価}

M2-c は、短期的な可行性余裕の改善と
機構層で推定された応答多様性の縮小が併存するような
潜在的不安定性を評価する第三の分類B標準機能である。
その評価結果として用いる指標は分類Cとして区別する。

\begin{definition}[M2-c潜在的不安定性評価機能（分類B）]
\label{def:m2c_latent_function}
M2-c は、現在および利用可能な過去の
$\Pi_{\mathrm{eval}}$、$\hat\mu_{\mathrm{eff},M2}$ その他から、
短期改善と応答多様性縮小の併存を検出しうる標準評価機能とする。
比較窓、履歴参照方法および具体的判定規則は実装依存とする。
\end{definition}

\begin{definition}[潜在的不安定性の機構層指標（分類C）]
\label{def:m2_latent_instability}
比較間隔 $\Delta t_{\mathrm{long}}\in\mathbb N$ を固定する。
保持記録を読み出す写像を
\[
\rho_{\mathrm{eval,past}}:\mathcal M^{\mathrm W}\to
 (\mathbb R\times\mathbb N_0)\cup\{\bot\}
\]
とする。$\bot$ は当該比較に利用できる記録の欠落を表し、数値ではない。
$t\ge\Delta t_{\mathrm{long}}$ で読み出しが対 $(a,b)$ を返す場合に限り、
\[
\Delta\Pi_{\mathrm{long},t}=\Pi_{\mathrm{eval},t}-a,\qquad
\Delta\hat\mu_{\mathrm{eff,long},t}=\hat\mu_{\mathrm{eff},M2,t}-b,
\]
\[
 C_{3,t}=\mathbf1[\Delta\Pi_{\mathrm{long},t}>0\ \wedge\
 \Delta\hat\mu_{\mathrm{eff,long},t}<0]
\]
と定める。$a,b$ は比較時点の値についての保持記録であり、過去の真値とは限らない。
短い履歴・欠落の場合は $C_{3,t}$ を定義せず、偽または $0$ に置換しない。
M2 の全域写像は欠落時の応答も定めるが、それを $C_{3,t}$ の値とはしない。
本指標は T-STAB や三原理の時刻別充足述語と同一ではない。
\end{definition}

\noindent\textbf{Priority警告体系と要請信号。}

\begin{definition}[Priority警告体系（分類B）]
\label{def:m2_priority}
$\mathcal P_{\mathrm{prio}}$ の各カテゴリを次のように解釈する。
\begin{description}
  \item[Priority 1] 急性危険または可行性悪化に関する警告。
  \item[Priority 2] 手段・資源・支援・関係等の構造的脆弱性に関する警告。
  \item[Priority 3] 将来的な脆弱化または応答可能性縮退の予兆に関する警告。
  \item[None] 上記の標準警告を生成しない状態。
\end{description}
具体的な優先規則は上記の意味を保つ範囲で実装依存とする。
\end{definition}

\begin{definition}[M2による要請信号の生成（分類B）]
\label{def:m2_request_signal}
M2 は、可行性と要求に関わる評価にもとづいて要請信号 $\mathrm{sig}_t$ を生成する。
$\mathrm{sig}_t$ は応答の動作モードを指定しない。
どの要求をどの質の要請として立ち上げるかの規則は実装依存とする。
内部の必要性の評価（たとえば欠乏の内部信号や電力残量の警告）は、
$\mathrm{sig}_t$ の要素として表れうる（備考~\ref{rem:m2_need_evaluation}）。
\end{definition}

\begin{remark}[要求×対象の評価の読み]
\label{rem:rq_evaluation_reading}
同じトークンに複数の要求が向く場合、評価は要求ごとに与えられる。
たとえば同じ大型動物のトークンが、食料の要求に照らしては獲物として、
安全の要求に照らしては脅威として評価されうる。
評価は保持内容 $M_t$ に依存し、理論上の帰結と食い違ってよい。
その動物の危険を記憶していない個体は、同じ対象を獲物としてだけ評価しうる。
M2 は M3 の目的や判断を評価に用いない。
要求の切り替わりは、行動の結果として生じた状況（たとえば相手を制圧できず危険が増したこと）を
M1 と M2 が感知・評価し、次の時刻の要請信号と評価が変わることとして、時刻をまたいで表す。
どの評価をどれだけ重く扱うかに専用の記号は置かない。
\end{remark}

M2 内部の概念的生成依存は
\[
  \text{M2-a / M2-b / M2-c}
  \longrightarrow
  (P_t,\mathrm{ev}_t)
  \longrightarrow
  (\mathrm{sig}_t,\mathrm{ev}^{\mathrm{rq}}_t)
\]
として整理する。
これらは情報依存の整理であり、実行の順序を定めない（備考~\ref{rem:tipa_bidirectional}）。

\begin{remark}[M2評価の情報品質と非定量性]
\label{rem:error_bounds}
M2 の内部評価と理論上の量（三原理の時刻別充足述語の真偽を含む）との間に、
普遍的な誤差上界、完全一致または収束性を要求しない。
したがって本論文は、
Priority の発火から三原理のいずれかの時刻別充足述語が偽であること、または閉塞を、
Priority が発火しないことからそれらの不在を、
いずれも一般的な帰結として置かない。
\end{remark}

\begin{remark}[内部の必要性評価と臨界必要要因]
\label{rem:m2_need_evaluation}
M2-a の急性脅威検出と、M2-b の資源不足性・支援不足性の評価は、
\S\ref{sec:three_principles} の臨界必要要因（定義~\ref{def:eu_fire}）に関わる差異を
内部で評価する標準的な箇所となりうる。
たとえば、生物における欠乏の内部信号や、
人工システムにおける電力残量の警告は、このような内部評価にあたる。
ただし本論文は、内部で必要と評価された要因と
理論上の一段階臨界必要要因集合 $\mathrm{Crit}^{(1)}(x_t)$ との一致を要求しない。
理論上は臨界必要要因である要因が内部で必要と評価されないこと（見逃し）も、
臨界必要要因でない要因が内部で必要と評価されること（誤評価）も、
本節の標準機能形式は排除しない。
このような内部評価は、現在の状態が $K$ の内にある局面（維持）にも
外にある局面（回復）にも生じうる。
EU の時刻別充足述語の真理条件は理論上の $\mathrm{Crit}^{(1)}(x_t)$ だけで定まり、
M2 の評価結果を用いない（備考~\ref{rem:crit_description_level}）。
本節は内部の必要性評価に専用の記号を置かず、
その結果は $C_{1,t}$、$\mathbf C_{2,t}$、$P_t$、$\mathrm{ev}_t$ および要請信号 $\mathrm{sig}_t$ に含まれうるものとして扱う。
\end{remark}

\begin{remark}[TIPAと三原理の多対多対応]
\label{rem:m2_three_principles}
M2-a、M2-b、M2-c および Priority 1、2、3 を、
SM、EU、FR のいずれか一つへ固定的に一対一対応させない。
同様に、M1 を SM と、M3 の応答形成を FR と一対一に同一視しない。
原理と機構の対応は多対多である（\S\ref{sec:three_principles} 定義~\ref{def:principles_tipa_connection}）。
\end{remark}


\subsubsection{M3：応答行動機構}
\label{subsec:m3}

M3 は、M1 と M2 から利用可能となる情報を判断の材料として、目的基準の評価を行い、
動作モード、目的および目的の状態を判断し、行動候補を構成・比較して、
PSL へ渡す TIPA 側の行動指令を形成する。

\begin{definition}[二つの評価と判断の区別（分類B）]
\label{def:eval_vs_judgment}
本節では次の三つを区別する。
\begin{enumerate}[label=(\roman*)]
  \item \emph{存続基準の評価}（M2）：三原理に関わる自己の可行性と根本要求を基準として、
  状態、要求または対象に値を与えること。当該ACSの構造の発露として生じる評価であり、
  M2 の出力 $P_t$、$\mathrm{ev}_t$、$\mathrm{sig}_t$、$\mathrm{ev}^{\mathrm{rq}}_t$ がこれにあたる。
  \item \emph{目的基準の評価}（M3）：目的、好みその他の思考の中で用いる基準に照らして、
  候補、状況または目的の進み具合に値を与えること。
  M3 の内部で処理され、候補適合比較 $\mathrm{fit}_t$（定義~\ref{def:m3_fit}）はその一つである。
  \item \emph{判断}（M3）：二つの評価と統合情報を材料に、
  動作モード、目的、目的の状態および行動候補を比べて決めること。
  M3 の出力 $\mathrm{mode}_t$、$\mathrm{Pur}_t$、$\mathrm{stat}_t$、$u_{\mathrm{M3},t}$ がこれにあたる。
\end{enumerate}
本節で修飾なしに「評価」と書く場合は (i) を指し、(ii) は「目的基準の評価」と書く。
M3 は存続基準の評価を生成しない。
目的基準の評価と判断は、M2 の入力に戻らない（定義~\ref{def:m2_map}）。
思考・熟慮を介さずに現れる選択の傾向は、目的基準の評価ではなく偏向性
（\S\ref{sec:extended_dynamics} 定義~\ref{def:bias_tendency}）として記述し、
存続基準の評価とも区別する。
\end{definition}

\begin{definition}[応答行動機構 $M_3$（分類B）]
\label{def:m3_map}
M3 を
\[
\begin{aligned}
\mathrm{M3}:&
X\times E_{\mathrm{ext}}\times\Xi
\times\mathcal P_{\mathrm{prio}}\times\mathfrak E
\times\mathcal P_{\mathrm{fin}}(\mathsf{Req})\times\mathsf{Sig}\times\mathfrak E^{\mathrm{rq}}_{\mathrm{map}}
\times\mathcal P_{\mathrm{fin}}(\mathsf{Pur})\times\mathcal M
\\
&\to
U\times\mathsf{Mode}\times\mathsf{Style}\times\mathcal P_{\mathrm{fin}}(\mathsf{Pur})\times\mathsf{St}
\end{aligned}
\]
なる決定論的標準写像とする。

\[
  \bigl(
    u_{\mathrm{M3},t},
    \mathrm{mode}_t,
    \mathrm{style}_t,
    \mathrm{Pur}_t,
    \mathrm{stat}_t
  \bigr)
  =
  \mathrm{M3}
  \bigl(
    \hat x_t,\hat e_t,\xi_t,
    P_t,\mathrm{ev}_t,
    \mathrm{Req}_t,\mathrm{sig}_t,\mathrm{ev}^{\mathrm{rq}}_t,\mathrm{Pur}_{t-1},M_t
  \bigr).
\]
$\mathrm{Pur}_{t-1}$ は前時刻の目的の集合の保持記録である（定義~\ref{def:s4_carrier_contract}）。
$\mathsf{St}$ と $\mathrm{stat}_t$ は定義~\ref{def:purpose_status} で定める。
M3 の定義域は、要求の一覧・要請信号・要求×対象の評価について、
それぞれの型の任意の元を含む。
したがって、要求の一部を取り除いた入力に対しても M3 の値は定まる
（定義~\ref{def:req_involvement}）。
\end{definition}

\paragraph*{4.2.3-a\quad M3-a：応答形成用情報の構成}

M3-a は、上流機構で既に生成された情報を、
候補形成と比較から利用可能な統合入力情報として構成する。
M3-a は Priority、識別状態、要請信号および評価を再生成しない。

\begin{definition}[M3-a統合入力情報（分類B）]
\label{def:m3a_info}
M3 の統合入力空間を
\[
  \mathcal I_{M3}
  :=
  \mathcal P_{\mathrm{prio}}
  \times\mathfrak E
  \times\mathcal P_{\mathrm{fin}}(\mathsf{Req})
  \times\mathsf{Sig}
  \times\mathfrak E^{\mathrm{rq}}_{\mathrm{map}}
  \times\mathcal P_{\mathrm{fin}}(\mathsf{Pur})
  \times X
  \times E_{\mathrm{ext}}
  \times\Xi
  \times\mathcal M
\]
とする。
時刻 $t$ の統合入力を
\[
  I_{M3,t}
  :=
  \bigl(
    P_t,
    \mathrm{ev}_t,
    \mathrm{Req}_t,
    \mathrm{sig}_t,
    \mathrm{ev}^{\mathrm{rq}}_t,
    \mathrm{Pur}_{t-1},
    \hat x_t,
    \hat e_t,
    \xi_t,
    M_t
  \bigr)
  \in\mathcal I_{M3}
\]
とする。
\end{definition}

\paragraph*{4.2.3-b\quad M3-b：候補形成・比較}

M3-b は、目的を形成し、利用可能な行動候補を構成して、
目的と要請に対する内部的な適合差を比較する。

\begin{definition}[目的の型と目的の形成（分類B）]
\label{def:purpose_formation}
目的の向き先の集合を
\[
  \mathsf{Dir}:=\{\mathsf{self}\}\cup\mathsf{Lab}\cup\bigl(\{\mathsf{self}\}\times\mathsf{Lab}\bigr)
\]
とする。$\mathsf{self}$ は自己の状態を、$k\in\mathsf{Lab}$ は内部対象トークンを、
$(\mathsf{self},k)$ は自己とトークン $k$ の関係を、それぞれ向き先とすることを表す。
目的の内容の非空集合 $\mathsf{Cnt}$ を置き、目的の集合を
\[
  \mathsf{Pur}:=\mathsf{Req}\times\mathsf{Dir}\times\mathsf{Cnt}
\]
とする。目的 $p=(q,d,z)$ の $q$ を $p$ の\emph{根本要求}と呼ぶ。
M3 が時刻 $t$ に判断として形成する目的の有限集合を
$\mathrm{Pur}_t\in\mathcal P_{\mathrm{fin}}(\mathsf{Pur})$ とし（定義~\ref{def:m3_map}）、
$\mathrm{Pur}_t=\emptyset$ を許す。
目的の形成の材料は M3 の入力に限られ、
M1 の統合情報と保持記録、M2 の評価と要請信号を含む。
要請信号は目的の主な源泉の一つであるが、
保持記録に残る継続中の目的や、トークンとその識別も源泉となりうる。
本論文は源泉の網羅を主張しない。
$\mathsf{Cnt}$ の元の列挙、形成の規則、および形成と動作モードの判断の内部での順序は実装依存とする。
\end{definition}

\begin{remark}[目的の階層と入れ子]
\label{rem:purpose_hierarchy}
目的は、根本要求、向き先と内容からなる近位目的、行動、結果の順に読む。
向き先が $\mathsf{self}$ の目的は対象トークンを要さない。
同じ根本要求をもつ複数の目的が同時に、または時刻をまたいで並ぶことを、目的の入れ子として読む。
たとえば渇きにもとづく「水を探す」と、見つけた水場のトークンに向かう「汲む」は、同じ根本要求を共有する。
手段と目的の関係そのものの定義は Phase~1.$\alpha$ の課題とする（\S\ref{subsec:phase1a}）。
目的化は内部のトークンの側の変化であり、
外部の成分が目的化によって行動応答成分になるとは読まない
（\S\ref{sec:acs_def} 備考~\ref{rem:Eext_decomposition_reading}）。
\end{remark}

\begin{remark}[目的の向き先と内部対象トークン]
\label{rem:purpose_direction_token}
目的の向き先 $k\in\mathsf{Lab}$ は内部対象トークンのラベルへの参照であり、
外部の個体の時間的な同一性を表すものではない。
同じラベルは時刻をまたいで異なる対象を表しうるので（備考~\ref{rem:token_generation}）、
時刻をまたぐ目的の向き先が同じ外部対象を指しつづけることは別の条件である。
トークンの持続の規則は \S\ref{sec:extended_dynamics} に委ねる。
向き先は $\mathsf{Lab}$ 上に定まるので、
$k\notin\mathcal K_t$ の時刻（たとえば対象が感知されなくなった時刻）にも、
M3 は $k$ に向いた目的を保持し、その状態を判断しうる。
一方、標準形の要求×対象の評価の定義域は $\mathrm{Req}_t\times(\mathcal K_t\cup\{\ast\})$ に含まれるので
（定義~\ref{def:rq_eval_types}）、その時刻に M2 は組 $(q,k)$ を評価せず、
全体の行についての評価 $(q,\ast)$ を与えうるにとどまる。
\end{remark}

\begin{definition}[目的の状態判断（分類B）]
\label{def:purpose_status}
判断値の集合を $\{\mathsf{cont},\mathsf{ach},\mathsf{term},\mathsf{und}\}$ とし、
\[
  \mathsf{St}:=\bigcup_{P}\{\mathsf{cont},\mathsf{ach},\mathsf{term},\mathsf{und}\}^{P}
\]
とする。ここで $P$ は $\mathsf{Pur}$ の有限部分集合を動く。
M3 が時刻 $t$ に出力する\emph{目的の状態判断}を $\mathrm{stat}_t\in\mathsf{St}$ とし、
標準形ではその定義域を前時刻の目的の集合の記録 $\mathrm{Pur}_{t-1}$ とする。
$\mathrm{stat}_t(p)=\mathsf{cont}$、$\mathsf{ach}$、$\mathsf{term}$、$\mathsf{und}$ はそれぞれ、
目的 $p$ を継続すること、達成したと判断すること、達成せずに終えること、
判断を保留することを表す。
$\mathsf{term}$ は放棄と別の目的への切り替えを含み、その理由の区別は Phase~1.$\alpha$ の課題とする。
標準形では、$\mathrm{stat}_t(p)=\mathsf{cont}$ ならば $p\in\mathrm{Pur}_t$、
$\mathrm{stat}_t(p)\in\{\mathsf{ach},\mathsf{term}\}$ ならば $p\notin\mathrm{Pur}_t$ とする。
$\mathrm{stat}_t(p)=\mathsf{und}$ の場合、$p$ が $\mathrm{Pur}_t$ に残るか否かは定めない。
\end{definition}

\begin{remark}[目的の状態判断の読み]
\label{rem:purpose_status_reading}
目的の達成は M3 が判断として定める。
M2 の評価の変化（たとえば危険評価の低下や警告の解除）は判断の材料であり、
それ自体を達成の判断とは読まない。
M3 は、M2 の評価、保持内容、M1 の統合情報および予測を材料に、
目的に照らした目的基準の評価を経て判断する。
確認が足りないと判断すれば $\mathsf{cont}$ または $\mathsf{und}$ とし、
情報取得の行動を候補に含めうる。
判断に用いる情報量、比較の詳細度および追加の確認の有無は、
動作モードと実装に依存しうる。本論文はその手順を固定しない。
単純な ACS では、明示的な達成の判断を経ずに、
要請が消えることによって目的が形成されなくなることがありうる。
この場合は $\mathsf{und}$ のもとで目的が $\mathrm{Pur}_t$ に残らないこととして記述され、
本論文はすべての ACS に明示的な達成の判断を要求しない。
内部の判断は理論上の達成（定義~\ref{def:purpose_course}）と一致しなくてよい。
\end{remark}

\begin{definition}[M3候補集合（分類B）]
\label{def:m3_candidate_set}
時刻 $t$ に M3 が応答形成のために扱う候補集合を
\[
  \mathrm{Cand}_t\subseteq U,
  \qquad
  \mathrm{Cand}_t\neq\emptyset
\]
とする。

他の通常候補が存在しない場合にも、
\S\ref{sec:acs_def} の選択的待機 $u_\emptyset\in U$ を用いて
\[
  \mathrm{Cand}_t=\{u_\emptyset\}
\]
となりうる。

$\mathrm{Cand}_t$ を真の可行行動集合 $A_t$ と同一視せず、
一般にいずれの向きの包含関係も仮定しない。
最終実効行動が $u_t\notin A_t$ となった場合の帰結は
\S\ref{sec:three_principles} の命題~\ref{prop:selection_responsibility}（$x_{t+1}\notin K$）による。
本論文は $u_t\in A_t$ を構成要件にも仮定にも置かない（備考~\ref{rem:selection_responsibility}）。
\end{definition}

\begin{definition}[候補適合比較（分類B）]
\label{def:m3_fit}
候補間の内部的な適合差を記述する写像を
\[
  \mathrm{fit}_t:
  \mathrm{Cand}_t\to\mathbb R
\]
とする。
$\mathrm{fit}_t$ は M2 評価情報、内部モデル、状態・環境推定、
Priority、保持情報その他に依存してよい。
ただし、真の $\Pi$、$A_t$、$\mathrm{Crit}^{(1)}(x_t)$、$\mu_{\mathrm{eff},t}$、$Q$ その他と同一視しない。
$\mathrm{fit}_t$ は目的基準の評価（定義~\ref{def:eval_vs_judgment}）の一つであり、存続基準の評価ではない。
\end{definition}
\begin{definition}[候補と適合比較の情報来歴（分類B）]
\label{def:s4_candidate_sources}
$\Phi_{\mathrm{cand}}:\mathcal I_{M3}\to\mathcal P(U)\setminus\{\emptyset\}$、
$\Phi_{\mathrm{fit}}:\mathcal I_{M3}\times U\to\mathbb R$ を固定し、
$\mathrm{Cand}_t=\Phi_{\mathrm{cand}}(I_{M3,t})$、
$\mathrm{fit}_t(u)=\Phi_{\mathrm{fit}}(I_{M3,t},u)$
（$u\in\mathrm{Cand}_t$）とする。
候補のコンパクト性や最大値の達成は仮定せず、
行動確定は既定義の決定論的選択規則による。
\end{definition}

\begin{remark}[候補ごとの予測の利用]
\label{rem:candidate_prediction}
M3 は、内部モデルが与える予測写像 $\hat g_t$（定義~\ref{def:env_set_prediction}）を
候補 $u\in\mathrm{Cand}_t$ に当てはめた $\hat G_t(u)$ を、比較の材料に用いてよい。
これは予測の当てはめであり、存続基準の評価の生成ではない。
当てはめの結果に目的基準の評価を与えることはありうる。
最終の実効行動についての予測集合は、$u_t$ が定まった後に記録される
（定義~\ref{def:prediction_record_check}）。
同一時刻の比較に $u_t$ についての予測を用いる記述は置かない。
\end{remark}


\begin{remark}[要請適合性と識別の区別]
\label{rem:fit_vs_identity}
M1-b は対象または情報を内部的に識別し、
M3-b は現在の目的と要請に対する行動候補の適合差を比較する。
したがって本論文は、M1 の識別結果だけから特定の行動候補の適合を一般的な帰結として置かない。
\end{remark}

情報が不足している場合、
情報取得自体が $U$ の行動として定義される実装では、
情報取得行動を $\mathrm{Cand}_t$ に含めてよい。

\paragraph*{4.2.3-c\quad M3-c：動作モードの判断と行動指令の確定}

M3-c は、動作モードの判断と M3-b の候補比較結果に基づき、
PSL へ渡す一つの TIPA 行動指令を確定する。

\begin{definition}[動作モードの判断（分類B）]
\label{def:mode_judgment}
$\mathrm{mode}_t\in\mathsf{Mode}$ は M3 の出力であり（定義~\ref{def:m3_map}）、
M3 が要請信号 $\mathrm{sig}_t$、Priority $P_t$、評価その他を材料に判断する。
M2 は動作モードを指定しない。
判断の手続き、および各モード内部の具体的な候補比較手続きは実装依存とする。
\end{definition}

\begin{definition}[要請と応答の不一致（分類C）]
\label{def:request_response_mismatch}
要請の質と動作モードの対応として、実装が与える関係
$\mathsf{Corr}\subseteq\mathsf{Qual}\times\mathsf{Mode}$ を固定する。
時刻 $t$ の要請 $(q,c)\in \mathrm{sig}_t$ について $(c,\mathrm{mode}_t)\notin\mathsf{Corr}$ であるとき、
その要請は時刻 $t$ の応答と\emph{不一致}であるという。
不一致の要請の集合を
\[
  \mathrm{Mis}_t:=\{(q,c)\in \mathrm{sig}_t\mid (c,\mathrm{mode}_t)\notin\mathsf{Corr}\}
\]
とする。
$\mathsf{Corr}$ の中身は Phase~1 では与えない。
本論文は、不一致が M2 または M3 の不具合に由来すること、
不一致から不利な結果が生じること、およびそれらの逆を、
いずれも一般的な帰結として置かない。
\end{definition}

\begin{definition}[M3動作モード（分類B）]
\label{def:M3_modes}
$\mathsf{Mode}$ の各モードを次のように解釈する。
\begin{description}
  \item[3A：緊急応答モード]
  高い切迫性を示す評価または要請に応じて、詳細な比較よりも応答形成の速さを相対的に優先する標準モード。
  \item[3B：熟慮的応答モード]
  複数候補および複数評価差を比較的詳細に比較する標準モード。
  \item[3C：通常応答モード]
  緊急性も特別な詳細比較要求も支配的でない場合の標準モード。
\end{description}
各モードは三原理の時刻別充足述語の真偽を表すラベルではない。
\end{definition}

\begin{remark}[M3-3AとPSL-4の区別]
\label{rem:m3_3a_vs_psl4}
3A は TIPA 内で M1・M2・M3 の利用可能情報に基づいて
緊急応答形成を行う動作モードであり、
詳細な TIPA 処理の完了を一般要件としない
局所保護応答である PSL-4 とは異なる。
したがって本論文は、3A の成立から PSL-4 の作動を、
PSL-4 の作動から 3A の成立を、いずれも一般的な帰結として置かない。
\end{remark}

\begin{definition}[3B応答様式（分類B）]
\label{def:m3_3b_styles}
3B 内で利用可能な標準応答様式を
\[
  \mathsf{Style}_{3B}
  :=
  \{\mathrm{3B\mbox{-}own},\mathrm{3B\mbox{-}rel},\mathrm{3B\mbox{-}comb}\}
\]
とする。
4.2 冒頭で定めた
\[
  \mathsf{Style}
  =
  \mathsf{Style}_{3B}\cup\{\mathrm{none}\}
\]
に対し、$\mathrm{mode}_t\neq\mathrm{3B}$ の場合には
\[
  \mathrm{style}_t=\mathrm{none}
\]
とする。
$\mathrm{mode}_t=\mathrm{3B}$ の場合に限り、M3 は必要に応じて
$\mathsf{Style}_{3B}$ の要素を応答様式として利用してよい。
要素 $\mathrm{none}$ は応答様式を区別しないことを表し、
応答形成そのものの不在を意味しない。

3B-own は自己側の可行性・内部維持に関係する評価差を相対的に重視する様式、
3B-rel は外部関係・接続条件・外部関係に伴う可行性関連差異に関係する評価差を相対的に重視する様式、
3B-comb は複数の評価差を実装依存の規則で統合する様式とする。
これらは規範的善悪または倫理的優劣の順位を意味しない。
\end{definition}

\begin{definition}[TIPA行動指令の決定（分類B）]
\label{def:action_decision}
M3-c は、$\mathrm{Cand}_t$、$\mathrm{fit}_t$、
$\mathrm{mode}_t$ および必要な応答様式を用いて、
決定論的な選択規則により
\[
  u_{\mathrm{M3},t}\in\mathrm{Cand}_t\subseteq U
\]
を一意に確定する。
モードごとの比較アルゴリズム、探索量、停止条件およびタイブレーク規則は、
決定論性と上記の型整合性を保つ範囲で実装依存とする。

$u_{\mathrm{M3},t}$ は M3 が生成する TIPA 側の行動指令であり、
状態遷移へ入力される最終実効行動 $u_t$ ではない。
両者の関係は \S\ref{subsec:psl} の実効行動解決で定める。
\end{definition}

M3-a、M3-b、M3-c は M3 の役割の整理であり、
その内部の順序と相互のやりとりは固定しない（備考~\ref{rem:tipa_bidirectional}）。
動作モード、目的および行動指令は、いずれも定義~\ref{def:m3_map} の M3 の出力として定まる。

\begin{remark}[M3の成功非保証]
\label{rem:m3_output_quality}
本論文は、M3 が $u_{\mathrm{M3},t}$ を確定したことから、
$u_{\mathrm{M3},t}\in A_t$、三原理の時刻別充足述語の真、閉塞回避、崩壊回避、
または期待結果の達成を一般的な帰結として置かない。
\end{remark}

%----------------------------------------------------------------------
\subsection{PSL：物理的基盤層}
\label{subsec:psl}

PSL は、TIPA と ISL の間で、
入力・出力、記憶I/O、基盤維持、局所保護応答、感度制御および実効行動形成を
機能的に接続する分類Bの標準機能形式である。
PSL は第四の認知機構ではなく、M1・M2・M3 の認知的処理を独自に再実行しない。

\begin{definition}[利用可能性（分類C）]
\label{def:psl_availability}
情報、信号、資源または作用能力がある機構から利用可能であるとは、
その機構の許容された入力または作用対象としてアクセス可能な状態にあることをいう。
本論文は、利用可能であることから、実際の利用、正確な解釈または後続処理の成功を一般的な帰結として置かない。
\end{definition}

\begin{definition}[PSLの形式的定義（分類B）]
\label{def:psl}
PSL を、入力取得、記憶I/O、行動実現、基盤維持、局所保護応答、
感度制御および実効行動解決を担う標準機能構成として定める。
具体的な内部配置、物理媒体、通信方式および実装トポロジーは固定しない。
\end{definition}

\subsubsection{PSL-1：入力・記憶出力インターフェース}
\label{subsubsec:psl1}

PSL-1 は、実装側に存在する外部・内部信号および記憶出力を、
TIPA から利用可能な情報へ接続する。

\begin{definition}[PSL-1の形式的定義（分類B）]
\label{def:psl1}
実装側入力信号空間を $\mathcal S_{\mathrm{in}}$ とし、
\[
  \Phi_{\mathrm{sense}}:
  \mathcal S_{\mathrm{in}}\to\mathcal I
\]
を入力インターフェースとする。
\[
  I_{\mathrm{raw},t}
  :=
  \Phi_{\mathrm{sense}}\bigl(s_{\mathrm{in},t}\bigr)
  \in\mathcal I
\]
とし、入力感度制御を用いる場合には
\[
  I_{\mathrm{sense},t}
  :=
  \Gamma_{\mathrm{in},t}\bigl(I_{\mathrm{raw},t}\bigr)
  \in\mathcal I
\]
とする。

さらに、記憶実装側の出力信号空間を
$\mathcal S_{\mathrm{mem},\mathrm{out}}$、
TIPA 側から利用可能な記憶出力情報空間を
$\mathcal I_{\mathrm{mem},\mathrm{out}}$ とし、
\[
  \Phi_{\mathrm{mem\_out}}:
  \mathcal S_{\mathrm{mem},\mathrm{out}}
  \to
  \mathcal I_{\mathrm{mem},\mathrm{out}}
\]
を記憶出力インターフェースとする。
\end{definition}

\begin{remark}[PSL-1とM1の役割分離]
\label{rem:psl1_m1}
PSL-1 は情報を利用可能にする接続を担い、
状態・環境推定、識別または可行性関連評価を行わない。
\end{remark}

\subsubsection{PSL-2：出力・記憶入力インターフェース}
\label{subsubsec:psl2}

PSL-2 は、PSL により確定された実効行動を作用側実装へ接続し、
必要に応じて保持対象情報を記憶実装側へ入力可能にする。

\begin{definition}[PSL-2の形式的定義（分類B）]
\label{def:psl2}
非空な駆動信号空間 $\mathcal S_{\mathrm{drive}}$ を置く。実効行動 $u_t\in U$ を実装側駆動信号へ接続する写像を
\[
  \Phi_{\mathrm{action}}:U\to \mathcal S_{\mathrm{drive}}
\]
とし、
\[
  x_{\mathrm{drive},\mathrm{raw},t}
  :=
  \Phi_{\mathrm{action}}(u_t)
\]
とする。
出力感度制御を用いる場合、
\[
  x_{\mathrm{drive},t}
  :=
  \Gamma_{\mathrm{out},t}
  \bigl(x_{\mathrm{drive},\mathrm{raw},t}\bigr)
\]
とする。

また、TIPA 側から記憶実装側へ渡される情報空間を
$\mathcal I_{\mathrm{mem},\mathrm{in}}$、
記憶実装側の入力信号空間を
$\mathcal S_{\mathrm{mem},\mathrm{in}}$ とし、
\[
  \Phi_{\mathrm{mem\_in}}:
  \mathcal I_{\mathrm{mem},\mathrm{in}}
  \to
  \mathcal S_{\mathrm{mem},\mathrm{in}}
\]
を記憶入力インターフェースとする。
\end{definition}

\begin{remark}[記憶入力とUpdateの境界]
\label{rem:mem_input_update_boundary}
PSL-2 は情報を記憶実装側へ接続しうるが、
保持内容そのものの $t\to t+1$ 更新則は定めない。
その動学は \S\ref{sec:extended_dynamics} の責務である。
\end{remark}

\subsubsection{PSL-3：基盤維持と基盤状態信号}
\label{subsubsec:psl3}

PSL-3 は、TIPA、PSL、ISL が作動可能であるための
基盤条件および資源条件を維持し、その状態を必要な機構から利用可能にする。

\begin{definition}[PSL-3の形式的定義（分類B）]
\label{def:psl3}
PSL基盤状態を
\[
  s_{\mathrm{PSL},t}\in\mathcal S_{\mathrm{PSL}}
\]
とする。
利用可能資源空間を $\mathcal R_{\mathrm{PSL}}$ とし、
基盤維持写像を
\[
  \Phi_{\mathrm{maint}}:
  \mathcal S_{\mathrm{PSL}}\times\mathcal R_{\mathrm{PSL}}
  \to\mathcal S_{\mathrm{PSL}}
\]
とする。
また、TIPA 側から利用可能となりうる基盤情報を
\[
  \Phi_{\mathrm{base}}:
  \mathcal S_{\mathrm{PSL}}\to\mathcal I_{\mathrm{base}}
\]
により抽出してよい。
\end{definition}

\begin{remark}[PSL-3とSMの時刻別充足述語の区別]
\label{rem:psl3_sm_base}
PSL-3 は自己維持に関わる基盤条件に関係するが、
本論文は、PSL-3 の存在または作動から SM の時刻別充足述語の真を一般的な帰結として置かない。
\end{remark}

\subsubsection{PSL-4：局所急性異常への保護応答}
\label{subsubsec:psl4}

PSL-4 は、局所的に取得可能な急性異常情報に対し、
M1・M2・M3 の詳細な処理完了を一般要件とせず、
保護応答へ接続しうる分類Bの標準機能形式である。
ここでいう急性は、普遍的な時間閾値、締切時刻または
最大許容遅延を意味しない。

\begin{definition}[PSL-4：局所急性異常への保護応答（分類B）]
\label{def:psl4}
局所信号空間を
\[
  \mathcal S_{\mathrm{local}}
\]
とし、局所保護応答を要する条件集合を
\[
  \mathcal C_{\mathrm{prot}}
  \subseteq
  \mathcal S_{\mathrm{local}}
\]
とする。
局所保護応答生成写像を
\[
  \Phi_{\mathrm{prot}}:
  \mathcal S_{\mathrm{local}}
  \to
  U\cup\{\bot\}
\]
とし、
\[
  r_{\mathrm{local},t}
  :=
  \Phi_{\mathrm{prot}}
  \bigl(s_{\mathrm{local},t}\bigr)
  \in U\cup\{\bot\}
\]
とする。
$\bot$ は局所保護応答を生成しないことを表す。

標準作動条件として
\[
  s_{\mathrm{local},t}\notin\mathcal C_{\mathrm{prot}}
  \Longrightarrow
  \Phi_{\mathrm{prot}}\bigl(s_{\mathrm{local},t}\bigr)=\bot
\]
とする。
$\mathcal C_{\mathrm{prot}}$ の具体的内容、局所信号構成、応答種類および時間尺度は実装依存とする。
\end{definition}

\begin{remark}[PSL-4の急性・局所・非熟慮性]
\label{rem:psl4_no_deliberation}
「急性」は普遍的な時間閾値を意味せず、
詳細な TIPA 処理の完了を一般要件とせず保護応答へ接続しうるという機能的区分を表す。
また「局所」は特定器官の物理的位置ではなく、
局所異常情報と保護応答との機能的接続を表す。
\end{remark}

\begin{remark}[PSL-4の構造と作動成功の区別]
\label{rem:psl4_required}
PSL-4 の標準構造を備えることと、各時刻に適切な保護応答が生成されることは同一ではない。
本論文は、$r_{\mathrm{local},t}\neq\bot$ から、SM の時刻別充足述語の真、閉塞回避または崩壊回避を一般的な帰結として置かない。
仮定~\ref{ass:psl_minimal} の (4) は、局所保護応答を生成する局所信号の存在だけを要求し、
各時刻に保護応答が生成されることを要求しない。
\end{remark}

\subsubsection{M3行動指令と局所保護応答の解決}
\label{subsubsec:effective_action_resolution}

M3 が形成する $u_{\mathrm{M3},t}$ は TIPA 側の行動指令であり、
状態遷移へ入力される最終実効行動 $u_t$ とは一般に異なりうる。
PSL は局所保護応答その他との関係を決定論的に解決して $u_t$ を確定する。

\begin{definition}[実効行動解決写像（分類B）]
\label{def:effective_action_resolve}
解決補助情報空間を $\mathcal R$ とする。
実効行動解決写像を
\[
  \Phi_{\mathrm{resolve}}:
  U\times(U\cup\{\bot\})\times\mathcal R
  \to U
\]
とする。
最終実効行動を
\[
  u_t
  =
  \Phi_{\mathrm{resolve}}
  \bigl(
    u_{\mathrm{M3},t},
    r_{\mathrm{local},t},
    R_t
  \bigr)
  \in U
\]
とする。

局所保護応答が存在しない場合には
\[
  \Phi_{\mathrm{resolve}}
  \bigl(u_{\mathrm{M3},t},\bot,R_t\bigr)
  =u_{\mathrm{M3},t}
\]
を標準退化条件とする。

局所保護応答が担う行動の成分を読み出す固定写像
\[
  \mathrm{pr}_{\mathrm{loc}}:U\to U_{\mathrm{loc}}
\]
を置く（$U_{\mathrm{loc}}$ は非空集合）。
$\mathrm{pr}_{\mathrm{loc}}(u)$ を $u$ の\emph{局所応答成分}と呼ぶ。
標準形では次の二つを置く。
\begin{enumerate}[label=(\alph*)]
  \item 局所保護応答の非退化：$\Phi_{\mathrm{prot}}(s)\neq\bot$ となるすべての $s\in\mathcal S_{\mathrm{local}}$ について
  $\mathrm{pr}_{\mathrm{loc}}\bigl(\Phi_{\mathrm{prot}}(s)\bigr)\neq\mathrm{pr}_{\mathrm{loc}}(u_{\emptyset})$。
  \item 発動の保存：すべての $u\in U$、$r\in U$、$R\in\mathcal R$ について
  \[
    \mathrm{pr}_{\mathrm{loc}}\bigl(\Phi_{\mathrm{resolve}}(u,r,R)\bigr)
    =
    \mathrm{pr}_{\mathrm{loc}}(r).
  \]
\end{enumerate}
(b) は、$\Phi_{\mathrm{resolve}}(u,r,R)=r$ とする解決によって満たされるので、
空でない条件である。
(a) により、局所保護応答の局所応答成分は選択的待機のそれと異なり、
(b) はそれを M3 行動指令によって待機その他の値に置き換えないことを要求する。
$U_{\mathrm{loc}}$ と $\mathrm{pr}_{\mathrm{loc}}$ の具体的な内容は実装依存とする。
局所応答成分は、局所保護応答の発動を実効行動の上で識別する抽象的な成分であり、
身体の特定の自由度、運動の波形またはその強さを表すものではない。
(b) が要求するのは、M3 行動指令によって発動が実効行動から消去されないことである。
発動した応答の強さや形が他の行動との干渉によって変わることは、
局所応答成分以外の差として表しうるので、(b) はこれを排除しない。
\end{definition}

局所保護応答が生じた時刻の実効行動について、
局所応答成分は局所保護応答の値をとり（発動の保存）、
それ以外の成分を M3 行動指令に従わせるか、
上書き、修飾、合成その他のどの方式で解決するかは実装依存とする。
本論文は、局所保護応答が実効行動の全体を置き換えることを仮定しない。

\begin{remark}[局所保護応答と TIPA の並行性]
\label{rem:psl4_concurrency}
局所保護応答 $r_{\mathrm{local},t}=\Phi_{\mathrm{prot}}(s_{\mathrm{local},t})$ は、
局所信号 $s_{\mathrm{local},t}=\rho_{\mathrm{local}}(I_t)$ だけから生成され
（定義~\ref{def:s4_psl_sources}）、TIPA の出力を用いない。
したがって局所保護応答は、TIPA の処理とは別に、同じ時刻に並行して生じうる。
実効行動解決写像が同じ時刻の $u_{\mathrm{M3},t}$ と $r_{\mathrm{local},t}$ をともに引数とすることは、
二つの値を同じ時刻の実効行動へ形式的にまとめることを表し、
M3 の処理の完了を局所保護応答の発生の条件とするものではない。
本論文は、同一時刻の中での物理的な先行関係を定めない。

局所保護応答の発動は、M3 の行動指令によって取り消されない
（定義~\ref{def:effective_action_resolve} の発動の保存）。
行動の実行中に局所保護応答が生じた場合、局所応答成分は局所保護応答に従い、
それ以外の成分は M3 行動指令に従いうる。
したがって局所保護応答は、行動の実行中にも並行して生じうる。

局所保護応答が生じた後の影響は、最終実効行動の記録 $u_{t-1}$ と、
状態および感知の変化を通じて、次の時刻以後の M1 で統合される
（備考~\ref{rem:tipa_bidirectional}）。
M3 は同じ時刻の局所保護応答を入力としないので、
これに応じる M3 の行動は次の時刻以後に選ばれる。
M3 が局所保護応答を止める応答は、発動を取り消すことではなく、発動の後の応答である。
それは、行動の作用が状態を通じて後の時刻の局所信号に及び、
局所信号が条件集合 $\mathcal C_{\mathrm{prot}}$ に入らなくなることにより、
後の時刻に局所保護応答が生成されなくなることとして表す。
後の時刻に再び局所保護応答が生成されれば、その発動も同様に保存される。
局所信号が $\mathcal C_{\mathrm{prot}}$ に入らない時刻には局所保護応答は発動しておらず、
発動の保存の対象ではない。
前の時刻の行動によって局所信号が $\mathcal C_{\mathrm{prot}}$ に入らないようにすること
（たとえば刺激を前もって避けること）は、発動を止めることではなく、
発動の条件を成立させないことである。
本論文は、発動の条件を成立させないこと、発動の後に続きを止めること、
および再び発動しないようにすることを区別する。
たとえば、咳やまばたきのような局所的な保護応答を PSL-4 の形式で表す場合、
M3 が選びうるのは発動そのものを止めることではなく、発動の後に続きを止める応答である。
\end{remark}

\subsubsection{PSL感度制御}
\label{subsubsec:psl_sensitivity}

PSL 感度制御は、入力側または出力側の接続特性を、
既に利用可能な制御情報に応じて調整する。

\begin{definition}[PSL感度制御写像（分類B）]
\label{def:psl_sensitivity}
入力側および出力側の許容感度写像族を
\[
  \mathfrak G_{\mathrm{in}},
  \qquad
  \mathfrak G_{\mathrm{out}}
\]
とする。
時刻 $t$ の感度写像を
\[
  \Gamma_{\mathrm{in},t}\in\mathfrak G_{\mathrm{in}},
  \qquad
  \Gamma_{\mathrm{out},t}\in\mathfrak G_{\mathrm{out}}
\]
とする。
各写像は少なくとも
\[
  \Gamma_{\mathrm{in},t}:\mathcal I\to\mathcal I,
  \qquad
  \Gamma_{\mathrm{out},t}:\mathcal S_{\mathrm{drive}}\to \mathcal S_{\mathrm{drive}}
\]
として型整合的に作用するものとする。
具体的なパラメータ化、選択規則、連続性および時刻間更新則は固定しない。
ただし、入力感度写像 $\Gamma_{\mathrm{in},t}$ は、
$I_{\mathrm{sense},t}$ を用いる当該時刻の TIPA 処理に先立って確定されるものとする。
$\Gamma_{\mathrm{in},t}$ の選択に TIPA 由来の内部情報を用いる場合には、
時刻 $t$ の当該 TIPA 処理より前に既に利用可能な情報
（標準的には前時刻までの TIPA 出力）にのみ依存させ、
同一ステップ内の M1・M2・M3 の出力から
$\Gamma_{\mathrm{in},t}$ を再決定する同時刻循環を許さない。
\end{definition}
\begin{definition}[PSLの信号・基盤担体と標準写像への接続（分類B）]
\label{def:s4_psl_sources}
次の読み出しを固定する。
\[
\begin{aligned}
\rho_{\mathrm{in}}&:X\times E_{\mathrm{ext}}\to\mathcal S_{\mathrm{in}},&
\rho_{\mathrm{PSL}}&:X\to\mathcal S_{\mathrm{PSL}},\\
\rho_{\Gamma,\mathrm{in}}&:X\to\mathfrak G_{\mathrm{in}},&
\rho_{\Gamma,\mathrm{out}}&:X\to\mathfrak G_{\mathrm{out}},\\
\rho_{\mathrm{mem,out}}&:X\to\mathcal S_{\mathrm{mem},\mathrm{out}},&
\rho_{\mathrm{local}}&:\mathcal I\to\mathcal S_{\mathrm{local}}.
\end{aligned}
\]
$s_{\mathrm{in},t}=\rho_{\mathrm{in}}(x_t,e_t)$、
$s_{\mathrm{PSL},t}=\rho_{\mathrm{PSL}}(x_t)$、
$\Gamma_{\mathrm{in},t}=\rho_{\Gamma,\mathrm{in}}(x_t)$、
$\Gamma_{\mathrm{out},t}=\rho_{\Gamma,\mathrm{out}}(x_t)$ とする。
PSL-1 と標準感知の接続は
\[
P_{\mathrm{sens}}(x,e)=
\rho_{\Gamma,\mathrm{in}}(x)
 \bigl(\Phi_{\mathrm{sense}}(\rho_{\mathrm{in}}(x,e))\bigr)
\]
を満たすように指定する。この等式は任意の既存写像から導かれる定理ではなく、
本節の標準形式を採用する際の接続条件である。
記憶実装出力 $\Phi_{\mathrm{mem\_out}}(\rho_{\mathrm{mem,out}}(x_t))$ と
基盤情報 $\Phi_{\mathrm{base}}(\rho_{\mathrm{PSL}}(x_t))$ を感知情報の成分として
用いる実現では、それぞれへの固定読み出し
$\mathcal I\to\mathcal I_{\mathrm{mem},\mathrm{out}}$、
$\mathcal I\to\mathcal I_{\mathrm{base}}$ がこの値を返すよう接続する。
これは物理担体の信号についての条件であり、任意の $M_t$ を $x_t$ から
復元する要請ではない。機能的な保持内容は $M_t$ から別途利用可能である。
局所信号は $s_{\mathrm{local},t}=\rho_{\mathrm{local}}(I_t)$ とする。
解決補助情報には固定写像 $\Phi_R:\mathcal M\times\mathcal I\to\mathcal R$
を用い、$R_t=\Phi_R(M_t,I_t)$ とする。
したがって同じ $M,I$ のもとで未観測の状態差だけを使って実効行動を変えない。

基盤維持への資源入力を読む写像
$\rho_{\mathrm{res}}:X\times U\times E_{\mathrm{ext}}\to\mathcal R_{\mathrm{PSL}}$
を固定する。PSL-3を独立した状態遷移とせず、採用する実現では
\[
\rho_{\mathrm{PSL}}(f(x,u,e))=
\Phi_{\mathrm{maint}}(\rho_{\mathrm{PSL}}(x),\rho_{\mathrm{res}}(x,u,e))
\]
として $f$ と整合させる。
PSL-2 の駆動信号と効果器の状態に依存する実際の作用も $f$ に含め、
駆動信号に独立した時刻間状態を置かない。
感度の後続値は $\rho_{\Gamma,\cdot}(f(x_t,u_t,e_t))$ であり、
指向的内部行動を通じた変更を含みうるが、それだけに限定しない。

記憶入力を利用する場合、固定写像
$\Phi_{\mathrm{write}}:\mathcal M\times\mathcal I\to\mathcal I_{\mathrm{mem},\mathrm{in}}$
により書き込み情報を与え、$\Phi_{\mathrm{mem\_in}}$ で実装信号へ接続する。
保持内容の次状態は Update が定め、別の記憶更新則を並置しない。
書き込み信号と Update の実装上の対応は、採用する各成分更新の型の範囲で指定する。
\end{definition}


\begin{remark}[感度制御の非認知性]
\label{rem:sensitivity_noncognitive}
PSL 感度制御は、認知評価、Priority生成、動作モード判定、PSL-4、
Automaticity または統合的高速処理そのものではない。
\end{remark}

\subsubsection{PSLの最小構造条件}
\label{subsubsec:psl_minimal_conditions}

\begin{assumption}[PSLの最小構造条件と非認知性]
\label{ass:psl_minimal}
本節で扱う ACS の PSL は、少なくとも次を満たすものとする。
\begin{enumerate}
  \item PSL-1 により必要な入力および記憶出力を TIPA から利用可能にしうる。
  \item PSL-2 により最終実効行動および必要な記憶入力を実装側へ接続しうる。
  \item PSL-3 により必要な基盤・資源維持機能を担いうる。
  \item PSL-4 により局所急性異常への保護応答を担持しうる。
  すなわち、ある $s\in\mathcal C_{\mathrm{prot}}$ について $\Phi_{\mathrm{prot}}(s)\neq\bot$ である。
  \item M1・M2・M3 の認知的責務を PSL が独自に代替しない。
\end{enumerate}
具体的な物理部位、回路、通信方式またはソフトウェア構成は固定しない。

導入地位：S（構造的）。強さ：中。
本仮定は、\S\ref{subsec:psl} の PSL の標準形式を適用する ACS の範囲を限定する構造的条件であり、本節の命題の前提としては用いない。
\end{assumption}

%----------------------------------------------------------------------
\subsection{目的・行動・結果の記述}
\label{subsec:purpose_action_result}

本小節では、\S\ref{subsec:tipa_architecture} と \S\ref{subsec:psl} の標準形式のもとで、
目的の帰属、要求の関与、行動の区分、および一つの行動の結果の記述を与える。
いずれも新たな機構や状態変数を導入しない。

\begin{proposition}[目的の自己帰属性]
\label{prop:purpose_self_attribution}
本節の標準形式のもとで、各時刻 $t$ の
$\mathrm{Req}_t$、$\mathrm{sig}_t$、$\mathrm{ev}^{\mathrm{rq}}_t$、$\mathrm{Pur}_t$、$\mathrm{stat}_t$ は、
$(M_t,I_t)$ を引数とする固定した写像の値として定まる。
したがって、当該ACSの外部にある主体の要求や評価の基準がこれらに入るのは、
感知情報 $I_t$ を通じる場合と、保持内容 $M_t$ に保持されている場合に限られる。
\end{proposition}

\begin{proof}
定義~\ref{def:s4_carrier_contract} により、
本節の時刻依存の機構出力は固定した規則のもとで $(M_t,I_t)$ から生成され、
前時刻記号は $M_t$ の成分 $M_t^{\mathrm W}$ の固定した読み出しの値である。
$\mathrm{Req}_t=\rho_{\mathrm{req}}(M_t)$ である（定義~\ref{def:request_signal_types}）。
M1 の値は $I_{\mathrm{sense},t}=I_t$、$M_t$ および前時刻記号の値であり（定義~\ref{def:m1_map}）、
M2 の値は M1 の出力、$M_t$ および前時刻記号の値である（定義~\ref{def:m2_map}）。
よって $\mathrm{sig}_t$ と $\mathrm{ev}^{\mathrm{rq}}_t$ は $(M_t,I_t)$ の固定した写像の値である。
M3 の入力はこれらと $M_t$ および前時刻記号からなるから（定義~\ref{def:m3_map}）、
$\mathrm{Pur}_t$ と $\mathrm{stat}_t$ も同様である。
後半は、写像の引数が $(M_t,I_t)$ に限られることから従う。
\end{proof}

命題~\ref{prop:purpose_self_attribution} は、固定した写像の値についての主張である。
これらの時刻別の値を $\Sigma_i$ の出力として帰属させるかどうかは、
定義~\ref{def:s4_output_attribution} に従う。

\begin{remark}[自己帰属性の射程と異常域]
\label{rem:self_attribution_scope}
命題~\ref{prop:purpose_self_attribution} は、
目的の基準が当該ACSの内部の担体から読み出されるという構造上の性質であり、
目的が当該ACSの利益に向かうことや $\Pi$ の維持に向かうことを意味しない。
他者の保護や自己の犠牲に向かう目的も、当該ACSが保持する要求から形成される限りでこの性質をもつ。
外部の主体が行動を直接に定めることは、型の上では生じない。
最終の実効行動 $u_t$ は $\Phi_{\mathrm{resolve}}$ の値としてだけ定まるからである
（定義~\ref{def:effective_action_resolve}）。
一方、外部からの作用によって $M_t$ が書き換えられる場合は、
型の上では学習による変化と区別できない。
機能不全、被操作および外生的な支配による評価機構の迂回・置換は本備考の記述にとどめ、
その由来の追跡は Phase~1.$\alpha$ 以後の課題とする（\S\ref{subsec:phase1a}）。
\end{remark}

\begin{definition}[要求の関与（分類C）]
\label{def:req_involvement}
時刻 $t$ について、M3 の入力のうち
$\hat x_t$、$\hat e_t$、$\xi_t$、$P_t$、$\mathrm{ev}_t$、$\mathrm{Pur}_{t-1}$、$M_t$ と、
実効行動解決の $r_{\mathrm{local},t}$、$R_t$ を固定する。
要求の一覧・要請信号・要求×対象の評価の組
$(R,\gamma,v)\in\mathcal P_{\mathrm{fin}}(\mathsf{Req})\times\mathsf{Sig}\times\mathfrak E^{\mathrm{rq}}_{\mathrm{map}}$ に対し、
これらを入力とする M3 の行動指令を $\Phi_{\mathrm{resolve}}$ で解決した値を
$\mathcal U_t(R,\gamma,v)\in U$ と書く。
要求 $q\in\mathrm{Req}_t$ について
\[
  R^{\setminus q}_t:=\mathrm{Req}_t\setminus\{q\},\quad
  \mathrm{sig}^{\setminus q}_t:=\{(q',c)\in \mathrm{sig}_t\mid q'\neq q\},\quad
  v^{\setminus q}_t:=\mathrm{ev}^{\mathrm{rq}}_t\big|_{D^{\mathrm{rq}}_t\setminus(\{q\}\times(\mathsf{Lab}\cup\{\ast\}))}
\]
とおく。
\[
  \mathcal U_t\bigl(\mathrm{Req}_t,\mathrm{sig}_t,\mathrm{ev}^{\mathrm{rq}}_t\bigr)
  \neq
  \mathcal U_t\bigl(R^{\setminus q}_t,\mathrm{sig}^{\setminus q}_t,v^{\setminus q}_t\bigr)
\]
が成り立つとき、要求 $q$ は時刻 $t$ の実効行動に\emph{関与した}という。
成り立たないとき、$q$ は時刻 $t$ に存在したが関与しなかったという。
定義~\ref{def:m3_map} により右辺も定まる。
\end{definition}

\begin{remark}[関与の判定の射程]
\label{rem:involvement_scope}
定義~\ref{def:req_involvement} は、要求の一覧、要請信号および要求×対象の評価という
明示的な経路で関与を判定する。
保持内容 $M_t$ は固定するので、$M_t$ を通じた影響は判定に含まれない。
したがって、関与したという判定はこの写像の上で結果を左右したことを表すが、
関与しなかったという判定は明示的な経路で結果を左右しなかったことにとどまり、
$M_t$ を通じた影響がないことを意味しない。
たとえば、要求に由来する応答が保持内容の習慣として蓄えられ、
要請信号から当該要求を除いても同じ行動が生じる場合、
その行動は明示的な経路では当該要求の関与をもたない。
関与の判定は最終の実効行動 $u_t$ について行い、
$u_{\mathrm{M3},t}$ についての同様の判定は補助の記述として扱う。
\end{remark}

\begin{definition}[行動の区分（分類C）]
\label{def:action_classes}
時刻 $t$ の実効行動 $u_t$ を次のように区分する。
\begin{enumerate}[label=(\roman*)]
  \item ある $q\in\mathrm{Req}_t$ が $u_t$ に関与したとき、$u_t$ を\emph{目的行動}と呼ぶ。
  \item どの $q\in\mathrm{Req}_t$ も関与しないとき、
  $u_t$ を\emph{目的を介さない自己生成の行動}と呼ぶ。
\end{enumerate}
局所保護応答（\S\ref{subsec:psl}）が M3 の行動指令を上書きした場合や、
Automaticity（\S\ref{subsec:fast_processing}）による応答は、(ii) に区分されうる。
外部による強制は当該ACSの行動ではなく、
外部環境状態 $e_t$ を通じて $f$ に入る状態変化として表す
（\S\ref{sec:acs_def} 備考~\ref{rem:action_space_structure}）。
同じ時刻に、目的行動と外部による強制は共存しうる。
\end{definition}

\begin{definition}[目的に関わる作用差と一段の条件充足（分類C）]
\label{def:purpose_action_difference}
記述者が目的 $p\in\mathsf{Pur}$ に割り当てる達成条件を $G_p\subseteq X\times O$ とする。
$G_p$ は当該ACSの内部量ではない。
$e=(b,o)$ について、状態と行動応答成分の合同遷移を
\[
  \bar f(x,u,e):=\bigl(f(x,u,e),\ g_O(o,x,u,b)\bigr)\in X\times O
\]
とする（\S\ref{sec:acs_def} 定義~\ref{def:env_realization}）。
\[
  \mathbf 1\bigl[\bar f(x,u,e)\in G_p\bigr]
  \neq
  \mathbf 1\bigl[\bar f(x,u_{\emptyset},e)\in G_p\bigr]
\]
が成り立つとき、行動 $u$ は $(x,e)$ において目的 $p$ に関わる\emph{作用差}をもつという。
また、$(x_{t+1},o_{t+1})\in G_p$ であることを、時刻 $t$ の行動の後の\emph{一段の条件充足}と呼ぶ。
これは目的の達成一般ではない（定義~\ref{def:purpose_course}）。
\end{definition}

\begin{proposition}[外部側の区分と枠の作用差]
\label{prop:purpose_ad_slot}
ある $G_O\subseteq O$ について $G_p=X\times G_O$ とする。
行動 $u$ が $(x,e)$ において目的 $p$ に関わる作用差をもつならば、
ある枠 $a\in\mathcal A_O$ について、$u$ は $(o,x,b)$ において枠 $a$ に作用差をもつ
（\S\ref{sec:acs_def} 定義~\ref{def:action_difference}）。
\end{proposition}

\begin{proof}
$G_p=X\times G_O$ であるから
$\mathbf 1[\bar f(x,u,e)\in G_p]=\mathbf 1[g_O(o,x,u,b)\in G_O]$ であり、
$u_{\emptyset}$ についても同様である。
仮定よりこの二つの値は異なるから、$g_O(o,x,u,b)\neq g_O(o,x,u_{\emptyset},b)$ である。
$O=\prod_{a\in\mathcal A_O}O_a$ であるから、ある $a\in\mathcal A_O$ について第 $a$ 成分が異なる。
$\mathcal A_O=\emptyset$ の場合は $O$ が一点集合であり、仮定が成り立たない。
\end{proof}

本論文は命題~\ref{prop:purpose_ad_slot} の逆を一般的な帰結として置かない。

\begin{remark}[外部による強制と作用差]
\label{rem:forcing_action_difference}
外部による強制がある時刻の当該ACS自身の寄与は、強制を含む $e_t$ を固定した作用差で記述する。
どの行動をとっても次の状態と行動応答成分が同じになる場合は、
すべての行動について目的に関わる作用差がないこととして表し、新たな概念を置かない。
実際に生じた強制は軌道の上で $e_t$ として記述し、
許容外的作用 $\Upsilon_{\mathrm{adm},i}$ は回復可能性の判定にだけ用いる
（\S\ref{sec:acs_def} 定義~\ref{def:collapse}）。
作用差は一段階の量であり、何段にもわたる寄与の比較、
および行動をとらなかった場合の軌道の全体は本論文では定めない。
\end{remark}

\begin{definition}[一つの行動の結果の記述（分類C）]
\label{def:result_tuple}
$\mathrm{WF}_{i,t}=1$ である時刻 $t$ と $p\in\mathrm{Pur}_t$ について、
目的 $p$ に対する時刻 $t$ の実効行動 $u_t$ の一段の結果を、次の三つ組で記述する。
\begin{enumerate}[label=(\roman*)]
  \item 目的に関わる作用差：$u_t$ が $(x_t,e_t)$ において $p$ に関わる作用差をもつか否か。
  \item 一段の条件充足：$(x_{t+1},o_{t+1})\in G_p$ か否か。
  $G_p=X\times G_O$ の場合は $o_{t+1}$ だけで定まる。
  \item 予測集合への帰属：内部の予測外と理論上の予測外のそれぞれの成否
  （定義~\ref{def:prediction_record_check}）。
  予測が $\bot$ の場合、および内部の予測外について $\mathrm{WF}_{i,t+1}=0$ の場合は未定義とする。
\end{enumerate}
未定義の成分は偽として扱わない。
内部の達成の判断は一段の結果には含めず、目的の経過の記述（定義~\ref{def:purpose_course}）で扱う。
\end{definition}

\begin{definition}[目的の経過の記述（分類C）]
\label{def:purpose_course}
時刻 $\tau$ に目的 $p$ が新たに目的の集合に入ったとする
（$p\in\mathrm{Pur}_{\tau}$ であり、$\tau\geq1$ のときは $p\notin\mathrm{Pur}_{\tau-1}$）。
ここでの $\mathrm{Pur}_{\tau-1}$ は時刻 $\tau-1$ の M3 の出力であり、保持記録の値ではない。
$\tau$ を起点とする $p$ の経過について、次の時刻を定める。
該当する時刻が存在しない場合は定めない。
\begin{enumerate}[label=(\roman*)]
  \item \emph{理論上の達成時刻}：$s^{\mathrm{th}}:=\min\{s>\tau\mid (x_s,o_s)\in G_p\}$。
  \item \emph{内部の達成判断時刻}：$s^{\mathrm{in}}:=\min\{s>\tau\mid \mathrm{stat}_s(p)=\mathsf{ach}\}$。
  ここで $s$ は $\mathrm{WF}_{i,s}=1$ であり $\mathrm{stat}_s(p)$ が定まる時刻を動く。
  \item \emph{終了時刻}：$s^{\mathrm{end}}:=\min\{s>\tau\mid \mathrm{WF}_{i,s}=1,\ p\notin\mathrm{Pur}_s\}$ と、
  その時刻の判断による\emph{終了の種別}。
  種別は、$\mathrm{stat}_{s^{\mathrm{end}}}(p)=\mathsf{ach}$ ならば達成の判断による終了、
  $\mathsf{term}$ ならば達成せずに終えた終了、
  $\mathsf{und}$ であるか定まらない場合は判断を経ない消失とする。
\end{enumerate}
$s^{\mathrm{th}}$ は記述者の水準の量、$s^{\mathrm{in}}$ は当該ACSの内部の水準の量であり、
一致しなくてよい。
たとえば $s^{\mathrm{th}}<s^{\mathrm{in}}$ は達成に遅れて気づく場合を、
$s^{\mathrm{in}}$ が定まり $s^{\mathrm{th}}$ が定まらない場合は達成を誤って判断した場合を表しうる。
一段で終わる目的は、$s^{\mathrm{th}}=\tau+1$ となる特別な場合である。
\end{definition}

\begin{remark}[存続の終わりと結果の記述]
\label{rem:result_after_collapse}
たとえば、最後の食料を子に与えた親が後に崩壊する場合を考える。
子の状態を親から見た行動応答成分として記述すれば
（\S\ref{sec:acs_def} 備考~\ref{rem:exogenous_environment}）、
与える行動の作用差は親の行動の時刻の値で定まり、
子の側の条件充足と理論上の達成時刻は子の成分の値で定まる。
親の目的の状態判断は、$\mathrm{WF}_{i,t}=0$ の時刻には定まらず（定義~\ref{def:s4_output_attribution}）、
「達成せずに終えた」とも「判断を保留した」とも読まない。
崩壊ではないが $\mathrm{WF}_{i,t}=0$ である時刻
（たとえば機能が停止し本節の機構が実現されていない時刻）も同様に扱う。
本論文は、与えなかった場合に親が存続したことを主張しない。
目的の自己帰属性（命題~\ref{prop:purpose_self_attribution}）は、
このような自己の犠牲に向かう目的にも成り立つ。
\end{remark}

%----------------------------------------------------------------------
\subsection{ISL：実装部位層}
\label{subsec:impl}

ISL は、感知、作用、情報保持その他を担いうる ACS の実装要素を、
機能上の役割によって整理する分類Bの標準機能形式である。
特定の器官、電子部品、ソフトウェアモジュールまたは配置を一般要件としない。

\begin{definition}[実装部位層 ISL（分類B）]
\label{def:impl_layer}
ACS の実装要素全体を
\[
  \mathcal B_\Sigma
\]
とする。
主要な標準部分集合として
\[
  B_{\mathrm{sens}},
  \qquad
  B_{\mathrm{eff}},
  \qquad
  B_{\mathrm{mem}}
  \subseteq\mathcal B_\Sigma
\]
を置き、それぞれ感知、作用、情報保持を担いうる実装要素集合とする。
相互排反性も三集合による網羅性も一般には要求しない。
\end{definition}

\subsubsection{感覚実装役割}
\label{subsubsec:sensor}

\begin{definition}[感覚実装役割（分類B）]
\label{def:sensor}
$B_{\mathrm{sens}}$ は、外部または内部の変化を
PSL-1 から利用可能な信号として担持しうる実装要素集合である。
\end{definition}

\begin{remark}[感覚役割とM1の区別]
\label{rem:sensor_distinction}
$B_{\mathrm{sens}}$ は情報取得を担持しうる実装側の役割であり、
M1 の情報統合・識別そのものではない。
\end{remark}

\subsubsection{効果実装役割}
\label{subsubsec:effector}

\begin{definition}[効果実装役割（分類B）]
\label{def:effector}
$B_{\mathrm{eff}}$ は、PSL により確定された実効行動を
外部または内部への作用として担持しうる実装要素集合である。
\end{definition}

\begin{remark}[効果役割とM3・PSLの区別]
\label{rem:effector_distinction}
$B_{\mathrm{eff}}$ は作用の実装担体であり、
行動指令を形成する M3 または実効行動を解決する PSL と同一ではない。
\end{remark}

\subsubsection{記憶実装役割}
\label{subsubsec:mem_substrate}

\begin{definition}[記憶実装役割（分類B）]
\label{def:mem_substrate}
$B_{\mathrm{mem}}$ は、過去情報を後続時刻で再利用可能な状態として
保持しうる実装要素集合である。
具体的な媒体、符号化法、分散配置および保持期間は実装依存とする。
\end{definition}

\begin{remark}[記憶実装役割と内部モデルの区別]
\label{rem:mem_model_distinction}
$B_{\mathrm{mem}}$ は実装担体であり、
内部モデル $M_t$、状態履歴または五機能的記憶集合そのものではない。
\end{remark}

\subsubsection{ISLとPSLの接続}
\label{subsubsec:impl_connection}

\begin{definition}[ISLとPSLの接続規約（分類B）]
\label{def:impl_psl_connection}
ISL と PSL の接続は、固定された解剖学的配線またはハードウェアトポロジーではなく、
定義~\ref{def:psl_availability} の利用可能性に基づく機能的関係として定める。
概念的には
\[
  B_{\mathrm{sens}}
  \longrightarrow
  \mathrm{PSL\mbox{-}1}
  \longrightarrow
  \mathrm{TIPA}
\]
および
\[
  \mathrm{TIPA}
  \longrightarrow
  \mathrm{PSL}
  \longrightarrow
  B_{\mathrm{eff}}
\]
という方向を許し、記憶については
\[
  B_{\mathrm{mem}}
  \longleftrightarrow
  \mathrm{PSL}
  \longleftrightarrow
  \mathrm{TIPA}
\]
という利用関係を許す。
\end{definition}

PSL-4 については、局所信号取得から局所保護応答の実現までが、
同一時刻の詳細な TIPA 処理完了を必要条件としない局所接続として
ISL 上に担持されうる。

以上より、M3、PSL、ISL の責務は
\[
  \text{M3行動指令}
  \longrightarrow
  \text{PSLによる実効行動解決}
  \longrightarrow
  \text{ISL上の作用}
\]
として区別される。

\subsubsection{存在・作動・健全度の区別}
\label{subsubsec:impl_state_distinction}

\begin{remark}[実装要素の存在・作動・健全度]
\label{rem:impl_exist_activity}
実装要素が存在すること、当該時刻に作動していること、
および健全であることは相異なる。
したがって本論文は、ISL の特定要素が停止または不健全であることだけから
$A_t=\emptyset$ または $\mathrm{Col}_{i,t}=1$ を一般的な帰結として置かない。
\end{remark}

\subsubsection{ISLの最小構造条件}
\label{subsubsec:impl_minimal_conditions}

\begin{assumption}[実装部位層の構造的仮定]
\label{ass:impl_structure}
本節で扱う ACS の ISL は、
感知、作用、および時刻をまたぐ情報保持等の必要な実装役割を
担持可能な構造をもつものとする。
同一要素が複数役割を担うことを許し、
役割ごとの独立器官・部品・モジュールの存在を要求しない。

導入地位：S（構造的）。強さ：中。
本仮定は、\S\ref{subsec:impl} の ISL の標準形式を適用する ACS の範囲を限定する構造的条件であり、本節の命題の前提としては用いない。
\end{assumption}

%----------------------------------------------------------------------
\subsection{統合的高速処理と Automaticity}
\label{subsec:fast_processing}

本小節では、TIPA、PSL、ISL の既存構造が取りうる
低遅延の横断的作動形態と、詳細な再評価・再比較への依存が低い作動形態を区別する。
新たな認知機構または機構層は導入しない。

\subsubsection{統合的高速処理の位置づけ}
\label{subsubsec:fast_position}

\begin{definition}[統合的高速処理（分類B）]
\label{def:integrated_fast_processing}
統合的高速処理とは、TIPA を中心とする既存機能構造において、
\begin{enumerate}
  \item 対応する詳細処理より詳細な再評価・再比較への依存が低く、
  \item 当該実装における対応する詳細処理より行動指令形成までの実効遅延が低い
\end{enumerate}
という二条件を満たして応答形成が進行しうる作動形態をいう。
Phase~1 では普遍的な時間閾値、演算量閾値または比較回数閾値を置かない。
\end{definition}

\begin{remark}[高速処理は独立構成要素ではない]
\label{rem:fast_not_component}
統合的高速処理は、TIPA、PSL、ISL に追加される独立機構、独立層、
または固定された専用経路として要求しない。
\end{remark}

\subsubsection{三層横断性と責務境界}
\label{subsubsec:fast_cross_layer}

\begin{definition}[高速処理の三層横断性（分類C）]
\label{def:fast_cross_layer}
統合的高速処理は、TIPA、PSL、ISL の一つまたは複数の既存機能が
低遅延化へ寄与することにより成立しうる。
\end{definition}

\begin{remark}[高速処理はTIPAの責務を代替しない]
\label{rem:fast_no_tipa_substitution}
低遅延化によって M1、M2、M3 の責務が PSL または ISL へ移ることを意味しない。
\end{remark}

\subsubsection{Automaticity：自動性と熟慮依存}
\label{subsubsec:automaticity}

\begin{remark}[Automaticityと熟慮依存度]
\label{rem:automaticity_deliberation}
Automaticity とは、当該時刻の情報処理または行動形成が、
対応する詳細な再評価・再比較への依存を比較的低くしたまま成立する作動形態をいう。
本節では分類Bの作動形態として扱い、
独立した機構、独立経路、新たな M3 動作モード、
または独立スカラー状態変数として導入しない。
また Automaticity と熟慮的作動の間に普遍的な二値境界を置かない。
\end{remark}

\begin{remark}[高速性・Automaticity・PSL-4の非同一性]
\label{rem:automaticity_fast_psl4}
統合的高速処理、Automaticity、PSL-4 は相異なる機能概念である。
低遅延性と詳細再評価への依存低下は定義~\ref{def:integrated_fast_processing} の別個の条件であり、
本論文は一方から他方を一般的な帰結として置かない。
また PSL-4 は局所急性異常への保護応答の標準形式であり、
Automaticity 一般を代表しない。
したがって、
統合的高速処理、Automaticity および PSL-4 は、
扱う概念の範囲を区別し、
相互に同一視しない。
\end{remark}

\begin{remark}[実装上の付随的事象と機能概念の区別]
\label{rem:structural_operation_implementation_event}
本論文は、筋収縮、神経・電子信号の発火、混線、ノイズ、誤作動、ハードウェア割込み、
アクチュエータの意図されない作動その他が、高速、局所的、非随意的または非熟慮的であることだけから、
統合的高速処理、Automaticity または PSL-4 の成立を一般的な帰結として置かない。
各概念として分類するためには、それぞれ固有の入力--処理--出力の機能関係が成立している必要がある。
\end{remark}

\subsubsection{高速処理を支える五機能的記憶集合}
\label{subsubsec:fast_memory}

\begin{definition}[五機能的記憶集合（分類B）]
\label{def:fast_memory_sets}
記憶役割の添字集合を
\[
  \mathcal L
  :=
  \{\mathrm{rout},\mathrm{sal},\mathrm{err},\mathrm{eval},\mathrm{resp}\}
\]
とする。
各 $\ell\in\mathcal L$ に対して非空な母集合 $\mathfrak T_\ell$ を置き、
時刻 $t$ に保持され、機能的役割によって分類される内容を
\[
  \mathcal T_{\ell,t}\subseteq\mathfrak T_\ell
\]
とする。

$\mathcal T_{\mathrm{rout},t}$ は反復・定型関連、
$\mathcal T_{\mathrm{sal},t}$ は顕著性関連、
$\mathcal T_{\mathrm{err},t}$ は誤差・異常関連、
$\mathcal T_{\mathrm{eval},t}$ は評価関連、
$\mathcal T_{\mathrm{resp},t}$ は応答関連の保持内容を表す。
保持内容の担体を明示するため、固定写像
$\rho_\ell:\mathcal M^{\mathrm L}\to\mathcal P(\mathfrak T_\ell)$ を与え、
$\mathcal T_{\ell,t}=\rho_\ell(M_t^{\mathrm L})$ とする。
ここでの保持内容と、当該時刻に読み出し・利用に成功することを区別する。
空集合値を許し、冪集合の位相や読み出しの連続性は要求しない。
\end{definition}

\begin{remark}[五記憶集合の非排他性・非網羅性]
\label{rem:fast_memory_nonexclusive}
五集合は相互排反である必要はなく、同一の保持内容が複数集合に属してよい。
また五集合によって ACS の全保持内容を網羅することも要求しない。
本論文は、各集合が空であることだけから対応機能の欠如、閉塞または崩壊を一般的な帰結として置かない。
\end{remark}

\begin{remark}[五記憶集合とTIPA機構の非固定対応]
\label{rem:fast_memory_no_fixed_mapping}
五つの記憶役割を M1、M2、M3 のいずれかへ一対一に固定配置しない。
\end{remark}

\begin{remark}[五記憶集合とISL記憶実装役割]
\label{rem:fast_memory_isl}
$\mathcal T_{\ell,t}$ は保持内容の機能分類であり、
$B_{\mathrm{mem}}$ の物理的・計算的部分集合と一対一対応することを要求しない。
\end{remark}

\subsubsection{高速処理の成立条件と非保証事項}
\label{subsubsec:fast_conditions}

\begin{assumption}[高速処理時の構造整合条件]
\label{ass:fast_structural_consistency}
統合的高速処理が成立する場合にも、
TIPA の型整合性、必要な生成依存、TIPA・PSL・ISL の責務境界、
および単一実効行動の形成を維持するものとする。

導入地位：S（構造的）。強さ：中。
本仮定は、\S\ref{subsec:fast_processing} の統合的高速処理の標準形式を適用する ACS の範囲を限定する構造的条件であり、本節の命題の前提としては用いない。
\end{assumption}

\begin{remark}[高速性と正確性の区別]
\label{rem:fast_not_accuracy}
本論文は、処理が高速であることから、推定・評価・Priority・応答選択の正確性、
最適性、安全性または三原理の時刻別充足述語の真を一般的な帰結として置かない。
\end{remark}

\begin{remark}[高速性の境界水準を一意に固定しないこと]
\label{rem:fast_boundary_levels}
低遅延性は当該 ACS の実装、課題および比較対象となる詳細処理に相対的である。
したがって統合的高速処理と詳細処理の境界は実装依存の記述的区分とし、
普遍的な時間閾値、演算量閾値または比較回数閾値を与えない。
\end{remark}

\subsubsection{時刻間動学への委譲}
\label{subsubsec:fast_boundary}

本節では、統合的高速処理および Automaticity の同一時刻上の位置づけを定める。
五機能的記憶集合の更新、統合的高速処理の形成・強化・弱化、
Automaticity の増減・解除その他の時刻間変化は、
\S\ref{sec:extended_dynamics} に委ねる。
Automaticity の形成を一方向的または不可逆な過程とは仮定しない。

%----------------------------------------------------------------------
\subsection{機能軸：横断的記述・診断量}
\label{subsec:functional_axis}

\subsubsection{機能軸（分類C）}
\label{subsubsec:functional_axis}

\begin{definition}[機能軸（分類C）]
\label{def:functional_axis}
TIPA、PSL および ISL の主要作動水準を横断的に記述する量を
\[
  F_t
  :=
  \bigl(
    F_{\mathrm{M1},t},
    F_{\mathrm{M2},t},
    F_{\mathrm{M3},t},
    F_{\mathrm{PSL},t},
    F_{\mathrm{ISL}_{\mathrm{sens}},t},
    F_{\mathrm{ISL}_{\mathrm{mem}},t},
    F_{\mathrm{ISL}_{\mathrm{eff}},t}
  \bigr)
  \in[0,1]^7
\]
と定め、機能軸と呼ぶ。
各成分は規格化された分類Cの記述・診断尺度であり、
普遍的な測定写像、各成分間の共通単位、異なる ACS 間の直接比較可能性を要求しない。
対象ごとに固定した読み出し
$\mathrm{rd}_F:X\times\mathcal M\times E_{\mathrm{ext}}\to[0,1]^7$ により
$F_t=\mathrm{rd}_F(x_t,M_t,e_t)$ と表す場合、独立した状態・更新則は追加しない。
現在の担体だけでは定まらない診断には、軌道上の読み出しを別途指定する。
この規則は内部診断が理論上の値を知ることを要求しない。
\end{definition}

\begin{remark}[機能軸と健全度]
\label{rem:psl_health_bridge_s5}
$F_t$ は同一時刻の作動水準を記述する量であり、
\S\ref{sec:extended_dynamics} で扱う PSL・ISL 健全度量とは別の分類C量である。
本節は両者の間に一般的な関数関係を置かない。
\end{remark}

\begin{remark}[機能軸と閉塞・崩壊の非同一性]
\label{rem:functional_axis_noncollapse}
本論文は、機能軸の特定値だけから閉塞または崩壊を一般的な帰結として置かない。
とくに、$F_t=(0,\ldots,0)$ から
$A_t=\emptyset$ も $\mathrm{Col}_{i,t}=1$ も一般的な帰結として置かない。
閉塞は \S\ref{sec:acs_def} の可行行動集合に関する定義により、
崩壊は \S\ref{sec:acs_def} の定義~\ref{def:collapse} により判定する。
本節は $\mathrm{Col}_{i,t}$ を定義せず、機能軸から崩壊判定を行わない。
\end{remark}


\begin{definition}[標準作動形式との合成対応（分類B）]
\label{def:s4_primitive_composition}
本節の標準形式では $\mathcal R_{\mathrm{int}}=\mathcal I\times\mathsf S_{M1}$、
$\mathcal V=\mathcal P_{\mathrm{prio}}\times\mathfrak E\times\mathsf{Sig}\times\mathfrak E^{\mathrm{rq}}_{\mathrm{map}}$ と具体化してよい。
$P_{\mathrm{disc}}(M,I)$ は $I$ と、保持記録を読み出して形成した M1 出力との対、
$P_{\mathrm{val}}(M,r)$ はその M1 出力と保持履歴に対する M2 出力とする。
$P_{\mathrm{act}}(M,r,v)$ は M3 指令、$r$ に保持された局所信号、
$\Phi_R(M,I)$ を PSL で解決した実効行動とする。
局所信号を M1 の推定値に置換することは要求しない。
この具体化は機能区分と機構の一対一対応ではなく、
固定した全域写像の合成を標準写像の型へ収めるための接続である。
第2節の非退化性はこの合成について別途確認する。
本論文は、本節の各機構が型を満たすことだけから、その非退化性を一般的な帰結として置かない。
\end{definition}

%----------------------------------------------------------------------
\subsection{小括}
\label{subsec:s4_summary}

本節では、単一 ACS の同一時刻における機能構造を、
TIPA、PSL、ISL の三つの機能水準から定式化した。

TIPA については
\[
  M_1\longrightarrow M_2\longrightarrow M_3
\]
という情報生成依存を置き、各機構を三つの機能へ分解した。

\[
\begin{aligned}
M_1 &: \quad
\text{M1-a 情報統合}
+
\text{M1-b 境界管理}
+
\text{M1-c 時間的連続性統合},
\\
M_2 &: \quad
\text{M2-a 可行性脅威検出}
+
\text{M2-b 構造的脆弱性評価}
+
\text{M2-c 潜在的不安定性評価},
\\
M_3 &: \quad
\text{M3-a 応答形成用情報構成}
+
\text{M3-b 候補形成・比較}
+
\text{M3-c 行動指令確定}.
\end{aligned}
\]

M2 は存続基準の評価として Priority、内部評価情報、要請信号および要求×対象の評価を与え、
M3 は目的基準の評価を経た判断として、動作モード、目的、目的の状態および行動指令を与える。
同一時刻の中には上記の情報依存の向きだけを置き、
逆向きの送受信は保持記録を通じて次の時刻に表した。

PSL については、入力取得、出力実現、基盤維持、
局所急性異常への保護応答、感度制御、および実効行動解決を定めた。
さらに、目的の自己帰属性、要求の関与、行動の区分、
および一つの行動の結果の記述を与えた（\S\ref{subsec:purpose_action_result}）。
ISL については、感知、作用、情報保持を担いうる実装要素を、
特定の生物学的器官または人工部品へ固定せず整理した。

したがって、単一 ACS の標準的な同一時刻機能関係は概念的に
\[
  \begin{aligned}
  \text{ISL上の情報取得}
  &\longrightarrow
  \text{PSLによる利用可能化}
  \longrightarrow
  M_1
  \longrightarrow
  M_2
  \longrightarrow
  M_3
  \\
  &\longrightarrow
  \text{PSLによる実効行動解決}
  \longrightarrow
  \text{ISL上の作用}
  \longrightarrow
  \text{状態遷移}
  \end{aligned}
\]
として整理される。
写像の型と、TIPA を経由しない PSL-4 の経路を含めた接続を図~\ref{fig:s4_same_time} に示す。

\begin{figure}[tbp]
\centering
\begin{tikzpicture}[
  font=\footnotesize,
  fn/.style={draw, rounded corners=2pt, align=center, inner sep=2.5pt, minimum height=6mm},
  val/.style={align=center, inner sep=1.5pt},
  note/.style={font=\scriptsize, align=center, inner sep=1pt},
  isl/.style={draw, dotted, thick, align=center, inner sep=2pt},
  grp/.style={draw, gray, rounded corners=3pt, inner sep=4pt},
  glabel/.style={gray, font=\scriptsize, anchor=south west, inner sep=1pt},
  plabel/.style={font=\small, anchor=west, inner sep=0pt},
  inp/.style={-{Latex[length=1.6mm]}},
  dep/.style={-{Latex[length=2.2mm]}, line width=1.2pt},
  nxt/.style={-{Latex[length=1.6mm]}, dashed},
  car/.style={dotted, thick},
  over/.style={preaction={draw, white, line width=3pt}},
]
\node[plabel] at (-0.3,1.2) {(a)\ 入力と応答の形成（時刻 $t$）};

\node[val] (xt) at (0,0.35) {$x_t$};
\node[val] (et) at (0,-0.35) {$e_t$};
\node[fn] (rin)  at (1.1,0) {$\rho_{\mathrm{in}}$};
\node[fn] (sens) at (2.7,0) {$\Phi_{\mathrm{sense}}$};
\node[fn] (gin)  at (4.5,0) {$\Gamma_{\mathrm{in},t}$};
\node[val] (It)  at (5.9,0) {$I_t$};
\draw[inp] (xt) -- (rin);
\draw[inp] (et) -- (rin);
\draw[inp] (rin) -- node[above]{$s_{\mathrm{in},t}$} (sens);
\draw[inp] (sens) -- node[above]{$I_{\mathrm{raw},t}$} (gin);
\draw[inp] (gin) -- (It);
\begin{scope}[on background layer]
\node[grp, fit=(rin)(sens)(gin)] (gPSL1) {};
\end{scope}
\node[glabel] at (gPSL1.north west) {PSL-1};
\node[note] at (4.5,-0.75) {$\Gamma_{\mathrm{in},t}=\rho_{\Gamma,\mathrm{in}}(x_t)$};
\node[isl] (Bsens) at (2.7,-1.1) {$B_{\mathrm{sens}}$};
\draw[car] (Bsens) -- (gPSL1.south -| Bsens);

\node[fn] (Mt) at (1.4,-2.9) {保持内容 $M_t$};
\node[isl] (Bmem) at (1.4,-4.1) {$B_{\mathrm{mem}}$};
\draw[car] (Bmem) -- (Mt);

\node[fn] (M1) at (7.0,-2.0) {$M_1$};
\node[fn] (M2) at (8.5,-2.0) {$M_2$};
\node[fn] (M3) at (10.0,-2.0) {$M_3$};
\draw[dep] (M1) -- (M2);
\draw[dep] (M2) -- (M3);
\draw[dep] (M1.north east) to[out=35,in=145] (M3.north west);
\begin{scope}[on background layer]
\node[grp, fit=(M1)(M2)(M3), inner ysep=8pt] (gTIPA) {};
\end{scope}
\node[glabel] at (gTIPA.north west) {TIPA};

\node[fn] (rloc) at (7.0,-4.1) {$\rho_{\mathrm{local}}$};
\node[fn] (prot) at (8.5,-4.1) {$\Phi_{\mathrm{prot}}$};
\draw[inp] (rloc) -- (prot);
\begin{scope}[on background layer]
\node[grp, fit=(rloc)(prot)] (gPSL4) {};
\end{scope}
\node[glabel] at (gPSL4.north west) {PSL-4（TIPA の出力を用いない）};

\node[fn] (PhiR) at (8.5,-5.3) {$\Phi_R$};

\draw (It) -- (5.9,-0.4);
\draw[over] (5.9,-0.4) -- (5.9,-5.3);
\draw[inp] (5.9,-2.0) -- (M1.west);
\draw[inp] (5.9,-4.1) -- (rloc.west);
\draw[inp] (5.9,-5.3) -- (PhiR.west);

\draw (Mt.east) -- (5.65,-2.9);
\draw[over] (5.65,-2.9) to[out=50,in=130] (6.15,-2.9);
\draw (6.15,-2.9) -- (10.0,-2.9);
\draw[inp] (7.0,-2.9) -- (M1.south);
\draw[inp] (8.5,-2.9) -- (M2.south);
\draw[inp] (10.0,-2.9) -- (M3.south);
\draw (Mt.south) ++(0.6,0) coordinate (ms) -- (ms |- 0,-5.9) coordinate (ms2);
\draw[inp] (ms2) -| (PhiR.south);

\node[val, anchor=west] (oM3) at (10.9,-2.0) {$u_{\mathrm{M3},t}$};
\node[val, anchor=west] (oR)  at (10.9,-4.1) {$r_{\mathrm{local},t}$};
\node[val, anchor=west] (oRt) at (10.9,-5.3) {$R_t$};
\draw[inp] (M3) -- (oM3);
\draw[inp] (prot) -- (oR);
\draw[inp] (PhiR) -- (oRt);
\node[note, anchor=north] at (11.3,-5.65) {(b) へ};

\begin{scope}[yshift=-8.0cm]
\node[plabel] at (-0.3,1.6) {(b)\ 実効行動の解決と時刻 $t+1$ への写像};

\node[val] (iM3) at (0.2,0.6)  {$u_{\mathrm{M3},t}$};
\node[val] (iR)  at (0.2,0)    {$r_{\mathrm{local},t}$};
\node[val] (iRt) at (0.2,-0.6) {$R_t$};
\node[note, anchor=south] at (0.2,0.85) {(a) から};

\node[fn] (res) at (2.6,0) {$\Phi_{\mathrm{resolve}}$};
\draw[inp] (iM3) -| ([xshift=-3mm]res.north);
\draw[inp] (iR) -- (res.west);
\draw[inp] (iRt) -| ([xshift=-3mm]res.south);
\begin{scope}[on background layer]
\node[grp, fit=(res)] (gRES) {};
\end{scope}
\node[glabel, anchor=north west] at (gRES.south east) {PSL（解決）};
\node[val] (ut) at (4.4,0) {$u_t$};
\draw[inp] (res) -- (ut);
\node[note, anchor=north west, align=left] at (-0.25,-1.0)
 {$r_{\mathrm{local},t}\in U\Rightarrow\mathrm{pr}_{\mathrm{loc}}(u_t)=\mathrm{pr}_{\mathrm{loc}}(r_{\mathrm{local},t})$\\
  $r_{\mathrm{local},t}=\bot\Rightarrow u_t=u_{\mathrm{M3},t}$};

\node[fn] (act)  at (6.2,0) {$\Phi_{\mathrm{action}}$};
\node[fn] (gout) at (8.6,0) {$\Gamma_{\mathrm{out},t}$};
\node[val] (xd)  at (10.6,0) {$x_{\mathrm{drive},t}$};
\draw[inp] (ut) -- (act);
\draw[inp] (act) -- node[above]{$x_{\mathrm{drive},\mathrm{raw},t}$} (gout);
\draw[inp] (gout) -- (xd);
\begin{scope}[on background layer]
\node[grp, fit=(act)(gout)] (gPSL2) {};
\end{scope}
\node[glabel] at (gPSL2.north west) {PSL-2};
\node[note] at (8.6,-0.75) {$\Gamma_{\mathrm{out},t}=\rho_{\Gamma,\mathrm{out}}(x_t)$};
\node[isl] (Beff) at (10.6,-1.1) {$B_{\mathrm{eff}}$};
\draw[car] (Beff) -- (xd);

\node[val] (iMt) at (0.2,-2.5) {$M_t$};
\node[fn] (Up) at (2.6,-2.5) {$\mathrm{Update}(M_t,x_t,e_t)$};
\node[fn] (f)  at (6.2,-2.5) {$f(x_t,u_t,e_t)$};
\node[fn] (gO) at (9.4,-2.5) {$g_O(o_t,x_t,u_t,b_t)$};
\node[val] (M1n) at (2.6,-3.45) {$M_{t+1}$};
\node[val] (x1)  at (6.2,-3.45) {$x_{t+1}$};
\node[val] (o1)  at (9.4,-3.45) {$o_{t+1}$};
\draw[inp] (iMt) -- (Up);
\coordinate (ud) at (5.15,-1.85);
\draw (5.15,0) -- (ud);
\draw[inp] (ud) -| (f.north);
\draw[inp] (ud) -| (gO.north);
\draw[nxt] (Up) -- (M1n);
\draw[nxt] (f) -- (x1);
\draw[nxt] (gO) -- (o1);

\begin{scope}[shift={(0,-4.2)}]
\draw[inp] (0,0) -- (0.8,0); \node[anchor=west, font=\scriptsize] at (0.9,0) {写像の入力（時刻 $t$）};
\draw[dep] (5.0,0) -- (5.8,0); \node[anchor=west, font=\scriptsize] at (5.9,0) {情報生成依存};
\draw[nxt] (0,-0.5) -- (0.8,-0.5); \node[anchor=west, font=\scriptsize] at (0.9,-0.5) {出力が時刻 $t+1$ に属する};
\draw[car] (5.0,-0.5) -- (5.8,-0.5); \node[anchor=west, font=\scriptsize] at (5.9,-0.5) {実装担体の役割との対応（写像ではない）};
\end{scope}
\end{scope}
\end{tikzpicture}
\caption{単一 ACS の時刻 $t$ における TIPA・PSL・ISL の機能的接続。
(a) と (b) は同じ時刻 $t$ の値でつながり、(a) の右端の値が (b) の左端の値と同一である。
太線は情報生成依存（M1 の出力が M2・M3 の入力、M2 の出力が M3 の入力）で、逐次実行の順序ではない。
箱は機能形式・写像・信号または実装役割を図示したものであり、物理的な部位との一対一対応を表さない。
$B_{\mathrm{sens}},B_{\mathrm{eff}},B_{\mathrm{mem}}\subseteq\mathcal B_\Sigma$ には排反性も網羅性も要求しない。
$M_t$ から三機構への入力は保持記録の読み出し（定義~\ref{def:s4_carrier_contract}）である。
時刻をまたぐ影響は、保持内容と、状態・外部環境状態の変化を通じて次の時刻の感知に入る経路に限られ、
M3 の出力を M2 の入力とする経路は置かない（備考~\ref{rem:m2_no_return_scope}）。
局所保護応答は TIPA の出力を用いずに生成され（備考~\ref{rem:psl4_concurrency}）、
実効行動の局所応答成分は局所保護応答に従う（定義~\ref{def:effective_action_resolve}）。
PSL-3 の基盤情報と記憶実装の出力を感知情報の成分に含める接続は、採用する実現に依存するため省略した。}
\label{fig:s4_same_time}
\end{figure}

また、統合的高速処理、Automaticity、PSL-4 を相異なる機能概念として区別し、
高速処理に利用されうる保持内容を五機能的記憶集合として整理した。
さらに、TIPA、PSL、ISL の主要作動水準を機能軸 $F_t$ により横断的に記述した。

次節では、本節で定めた構造を再定義せず、
記憶、内部量、健全度、統合的高速処理、Automaticity その他が
$t\to t+1$ でどのように変化しうるかを扱う。


%================ Section 5 標準作動形式と時刻間の接続 ================
%======================================================================
\section{標準作動形式と時刻間の接続}
\label{sec:extended_dynamics}

本節では、\S\ref{sec:acs_def} の型と更新、\S\ref{sec:three_principles} の記述水準、
\S\ref{sec:tipa} の同時刻構造を接続する。
対象は決定論的・離散時間の単一ACSである。原始的な機能の標準写像および
読み出し規則を固定し、環境実現（定義~\ref{def:env_realization}）を与えた場合に、
一段の形式的記録が定まり、次時刻へ閉じることを示す。
数学的記録の存在と、実在個体の機能の活動・整形式実現・存続は区別する。
本節で新たに用いる主要記号の一覧は付録~\ref{app:symbols_s5} に示す。

%----------------------------------------------------------------------
\subsection{責務 R1〜R4}
\label{subsec:update_dependency}
\label{subsubsec:update_role}
\begin{definition}[本節の責務]
\label{def:s5_responsibilities}
本節の責務を次の四項目とする。
\begin{description}
\item[R1] 行動応答成分を含めて、型整合的な一段の作動記録を与える。
\item[R2] 状態・内部モデル・行動応答成分の次時刻への接続を、
既存の $f,\mathrm{Update}$ と環境実現の $g_O$ で閉じる。
\item[R3] 固定した背景成分の列・初期値・標準写像のもとで、数学的記録の存在と一意性を示す。
\item[R4] 後続節へ可行行動集合、作動記録、分類C記述量を受け渡す。
\end{description}
R3 は可行性維持、非閉塞、非崩壊、実装上の活動の継続を意味しない。
\end{definition}

\begin{remark}[他節との責務境界]
\label{rem:s5_scope_boundary}
\S\ref{sec:tipa} の機構・基盤・実装の責務を再定義しない。
関係記述は \S\ref{sec:rel_min}、実現の終端・崩壊とその記録の解釈は
\S\ref{sec:acs_def}・\S\ref{sec:collapse_cond} に従う。
本節の全域写像の評価は、機能的終端後にも実装が作動するとの主張ではない。
\end{remark}

%----------------------------------------------------------------------
\subsection{一段の作動形式}
\label{subsec:s5_operation}
\begin{definition}[標準作動の一段]
\label{def:s5_standard_step}
定義~\ref{def:three_primitives} の全域一価写像と、
環境実現 $\bigl(E^{\mathrm{bg}},O,g_O,(b_t)_{t\ge0}\bigr)$（定義~\ref{def:env_realization}）を固定し、
$(x_t,M_t,o_t)\in X\times\mathcal M\times O$ と $b_t\in E^{\mathrm{bg}}$ に対して
\[
\begin{aligned}
e_t&=(b_t,o_t),\\
I_t&=P_{\mathrm{sens}}(x_t,e_t),\\
r_t&=P_{\mathrm{disc}}(M_t,I_t),\\
v_t&=P_{\mathrm{val}}(M_t,r_t),\\
u_t&=P_{\mathrm{act}}(M_t,r_t,v_t),\\
x_{t+1}&=f(x_t,u_t,e_t),\\
M_{t+1}&=\mathrm{Update}(M_t,x_t,e_t),\\
o_{t+1}&=g_O(o_t,x_t,u_t,b_t)
\end{aligned}
\]
とする。$u_t$ は PSL の解決を含む実効行動であり、M3 指令と同一視しない。
具体化する場合は定義~\ref{def:s4_primitive_composition} による。
$x_t,M_t,o_t$ はこの段の開始時の値であり、同時刻の出力を入力へ循環させない。
\S\ref{sec:tipa} の機構出力（目的の集合 $\mathrm{Pur}_t$、目的の状態判断 $\mathrm{stat}_t$ など）は、
固定した規則のもとで $(M_t,I_t)$ から生成される（定義~\ref{def:s4_carrier_contract}）。
したがってこの段の値から一意に定まり、作動の一段に別の成分として加えない。
\end{definition}

\begin{definition}[Updateの参照と一段遷移写像]
\label{def:update_section5}
\label{def:s5_step_map}
\S\ref{sec:acs_def} の $\mathrm{Update}:\mathcal M\times X\times E_{\mathrm{ext}}\to\mathcal M$ と
定義~\ref{def:env_realization} の $g_O$ を用い、一段遷移写像を
\[
\begin{aligned}
\Phi_{\mathrm{step}} &:X\times\mathcal M\times O\times E^{\mathrm{bg}}\to X\times\mathcal M\times O,\\
\Phi_{\mathrm{step}}(x,M,o,b)&:=
\bigl(f(x,u,e),\ \mathrm{Update}(M,x,e),\ g_O(o,x,u,b)\bigr),\\
e&=(b,o),\qquad r=P_{\mathrm{disc}}(M,P_{\mathrm{sens}}(x,e)),\\
v&=P_{\mathrm{val}}(M,r),\qquad u=P_{\mathrm{act}}(M,r,v)
\end{aligned}
\]
と定める。$X\times\mathcal M\times O$ は既存の状態と行動応答成分の積であり、第八のタプル成分ではない。
$g_O$ も7タプルの成分ではなく、環境実現の側の構造である。
$X$ には既定の部分空間位相、$\mathcal M$ には max 距離の位相、積には積位相を入れる。
仮定~\ref{ass:compact_spaces}・\ref{ass:M_compact} のもとで $X\times\mathcal M$ はコンパクトである。
$O$ のコンパクト性は仮定しないので（定義~\ref{def:Eext_decomposition}）、
$X\times\mathcal M\times O$ のコンパクト性は述べない。
本節は $\Phi_{\mathrm{step}}$ の連続性を仮定も結論もしない。
$\mathcal A_O=\emptyset$ のとき $O$ は一点集合であり、第三成分は定値となる。
\end{definition}

%----------------------------------------------------------------------
\subsection{学習成分・作動成分と保持記録}
\label{subsec:memory_dyn}
\label{subsubsec:s5_memory_update}
\begin{definition}[保持内容の時刻間読み出し]
\label{def:memory_intertime_update}
定義~\ref{def:model_learning_working} の $M_t^{\mathrm L},M_t^{\mathrm W}$ を用いる。
定義~\ref{def:fast_memory_sets} の固定読み出しによって
\[
\mathcal T_{\ell,t}=\rho_\ell(M_t^{\mathrm L}),\qquad
\mathcal T_{\ell,t+1}=\rho_\ell(\operatorname{pr}_L\mathrm{Update}(M_t,x_t,e_t))
\]
とする。空集合値を許し、保持内容に独立更新写像を置かない。
$\rho_\ell$ は全域写像であるが、保持内容が実際に利用できるかは別の問題である。
\end{definition}

\begin{remark}[M1内部状態と履歴の担体]
\label{rem:s5_m1_state_ref}
$\mathsf S_{M1,t}$ は同時刻出力の束である。
その後の参照には定義~\ref{def:s4_carrier_contract} の保持記録を用いる。
本論文は、二成分の直積や包括的Updateの型だけから、任意の履歴を正確に保持すること、
各機構成分間で任意の情報を書き込めることを一般的な帰結として置かない。
成分更新の依存順序と、同時刻の M1・M2・M3 の生成依存は区別する。
\end{remark}

%----------------------------------------------------------------------
\subsection{一段遷移と作動記録の継続}
\label{subsec:s5_continuation}
\begin{proposition}[P1：一段の整形式性]
\label{prop:s5_one_step}
定義~\ref{def:three_primitives}、定義~\ref{def:state_transition}、
定義~\ref{def:update_operator} の全域一価写像と、
定義~\ref{def:env_realization} の全域一価写像 $g_O$ を固定すると、
各 $(x,M,o,b)\in X\times\mathcal M\times O\times E^{\mathrm{bg}}$ に対し、
定義~\ref{def:s5_standard_step} の各値と $\Phi_{\mathrm{step}}(x,M,o,b)$ は
一意に定まり、後者は $X\times\mathcal M\times O$ に属する。
\end{proposition}
\begin{proof}
定義~\ref{def:Eext_decomposition} より $e=(b,o)\in E_{\mathrm{ext}}$ である。
$P_{\mathrm{sens}}$ の値は $\mathcal I$ に、次に $P_{\mathrm{disc}}$ の値は
$\mathcal R_{\mathrm{int}}$ に、$P_{\mathrm{val}}$ の値は $\mathcal V$ に属する。
したがって $P_{\mathrm{act}}$ の値は $U$ に属し、$f$ の型から次状態は $X$ に属する。
Updateの型から次内部モデルは $\mathcal M$ に、$g_O$ の型から次の行動応答成分は $O$ に属する。
各写像は一価であり、入力の決定に同時刻の下流出力を要求しないので一意である。
\end{proof}

\begin{definition}[標準作動記録]
\label{def:s5_operation_record}
環境実現 $\bigl(E^{\mathrm{bg}},O,g_O,(b_t)_{t\ge0}\bigr)$ と
初期値 $(x_0,M_0,o_0)\in X\times\mathcal M\times O$ を与え、
定義~\ref{def:s5_standard_step} に従う列
$(x_t,M_t,o_t,b_t,e_t,I_t,r_t,v_t,u_t)_{t\ge0}$ を標準作動記録と呼ぶ。
外生的に与えるのは背景成分の列だけであり、
行動応答成分 $o_t$ は当該ACSの過去の行動に応答しうる（備考~\ref{rem:exogenous_environment}）。
その状態・環境部分列 $(x_t,M_t,e_t)_{t\ge0}$ の属する集合を
$\mathcal Z_{\mathrm{op}}\subseteq(X\times\mathcal M\times E_{\mathrm{ext}})^{\mathbb N_0}$
と書く。これは固定した標準写像と $g_O$ に相対する集合である。
閉集合性・コンパクト性は要求しない。集合としての存在だけを以下で用いる。
作動記録の $(x_t,o_t)$ の部分は、記録の $(u_t)$ を実効行動の列として
定義~\ref{def:env_realization} に続く式により定まる列に一致する。
$\mathcal A_O=\emptyset$ のとき、本定義は背景成分の列を外部環境状態の列とみなして与えた記録と同じである。
\end{definition}

\begin{proposition}[P2：数学的作動記録の継続]
\label{prop:s5_continuation}
P1の条件のもとで、任意の背景成分の列と初期値に対して、全 $t\in\mathbb N_0$ 上で
標準作動記録がただ一つ存在する。
\end{proposition}
\begin{proof}
$t=0$ の三つ組 $(x_0,M_0,o_0)$ は与えられている。
$t$ の三つ組が $X\times\mathcal M\times O$ に属するとき、
P1により $b_t$ とあわせて当該段の中間値と次の三つ組が一意に定まる。
自然数上の再帰により全時刻の列が得られる。
同じ初期値と背景成分の列をもつ二列を比較すれば、帰納法により各時刻で一致する。
\end{proof}

\begin{proposition}[P3：作動記録の因果性]
\label{prop:s5_causality}
固定した標準写像と $g_O$ のもとで、$(x_t,M_t,o_t)$ は $(x_0,M_0,o_0)$ と
$b_0,\ldots,b_{t-1}$ のみで定まる。同時刻の $e_t,I_t,r_t,v_t,u_t$ は、さらに $b_t$ を用いる。
\end{proposition}
\begin{proof}
$t=0$ は初期値そのものである。$t$ までの主張を仮定すると、
定義~\ref{def:s5_step_map} は $t$ の三つ組と $b_t$ のみから次の三つ組を返す。
したがって $t+1$ にも主張が成り立つ。中間値の主張は $e_t=(b_t,o_t)$ と各合成の型から従う。
$o_t$ は過去の行動に依存しうるので、因果性は背景成分の列に対して述べる。
\end{proof}

\begin{remark}[形式的継続と実現の継続]
\label{rem:s5_formal_not_survival}
P1〜P3 は $K,A_t,\mathrm{WF}_{i,t},\mathrm{FT}_{i,t},\mathrm{Col}_{i,t}$ を前提に用いない。
$P_{\mathrm{act}}$ は全域写像なので $A_t=\emptyset$ の時刻にも形式的行動値 $u_t$ は定まるが、
この値を可行行動とはしない。$x_t\notin K$ からも記録は続けられる。
$\mathrm{WF}_{i,t}=0$ の時刻の機構出力の扱いは定義~\ref{def:s4_output_attribution} による。
機能的終端後や崩壊後の形式的記録を、同一個体の活動の実現と読み替えない。
同一ACSの記述域は \S\ref{sec:collapse_cond} の $I_{\mathrm{ACS}}(\sigma)$ に従い、
原理判定はさらに \S\ref{sec:three_principles} の $\mathbb T_{\mathrm P,i}$ と FR の定義域に従う。
\end{remark}

%----------------------------------------------------------------------
\subsection{保持記録の受け渡し}
\label{subsec:s5_record_handover}
\begin{definition}[保持記録の一致（分類C）]
\label{def:s5_record_agreement}
標準作動記録と定義~\ref{def:s4_carrier_contract} の読み出し $\rho_{\mathrm{past}}$ を固定する。
$k=1,\ldots,7$ について、$\rho_{\mathrm{past}}$ の値の第 $k$ 成分への射影を $\operatorname{pr}_k$ とし、
時刻 $t$ の対応する値 $y_{k,t}$ を次のとおり定める。
\begin{enumerate}
\item Priority $P_t$。
\item 実効行動 $u_t$。
\item 時刻 $\max\{0,t-W+1\}$ から $t$ までの状態推定 $\hat x_s$ の列。
\item M1 出力の束 $\mathsf S_{M1,t}$。
\item 時刻 $\max\{0,t-W_{\mathrm{short}}+1\}$ から $t$ までの感知情報 $I_s$ の列。
\item 目的の集合 $\mathrm{Pur}_t$。
\item 実効行動についての外部側の予測集合 $\hat G_t(u_t)$（定義~\ref{def:env_set_prediction}）。
\end{enumerate}
$\operatorname{pr}_k\bigl(\rho_{\mathrm{past}}(M^{\mathrm W}_{t+1})\bigr)=y_{k,t}$ が成り立つとき、
第 $k$ 成分の記録は時刻 $t$ で\emph{一致する}という。
定義~\ref{def:m2_latent_instability} の $\rho_{\mathrm{eval,past}}$ など他の保持記録の読み出しについても、
対応する過去の値を指定すれば同じ形で定める。
\end{definition}

記録が一致するかどうかは、Update の詳細型と作動成分の内容に依存する。
本論文はその一致を一般的な帰結として置かず、仮定としても置かない（定義~\ref{def:s4_carrier_contract}）。
命題~\ref{prop:s5_continuation} は記録の一致を前提としない。
\S\ref{sec:tipa} の機構は記録の値によらず値を返す全域写像として与えられるので、
記録が一致しない時刻を含めて、作動記録はただ一つ定まる。

\begin{remark}[目的の受け渡し]
\label{rem:s5_purpose_handover}
時刻 $t+1$ の M3 の入力に現れる前時刻の目的の集合は、第 $6$ 成分の記録値である。
目的の状態判断 $\mathrm{stat}_{t+1}$ の定義域は、標準形ではこの記録値であり（定義~\ref{def:purpose_status}）、
第 $6$ 成分の記録が時刻 $t$ で一致する場合に限り、その定義域は $\mathrm{Pur}_t$ に等しい。
一致しない場合の目的の経過は、定義~\ref{def:purpose_course} の範囲で記述される。
記録から欠けた目的は判断を経ずに目的の集合から外れうる。
記録にだけ現れる目的を M3 が継続とした場合、その目的は時刻 $t$ の出力 $\mathrm{Pur}_t$ に含まれないので、
定義~\ref{def:purpose_course} では時刻 $t+1$ に新たに目的の集合へ入ったものとして扱われる。
\end{remark}

\begin{remark}[内部対象トークンの持続]
\label{rem:s5_token_persistence}
トークンの時刻をまたぐ持続（備考~\ref{rem:token_generation}）は、
時刻 $t+1$ の M1 が第 $4$ 成分の記録値を入力の一つとして、時刻 $t$ と同じラベルを与えることとして表される。
同じラベルを与えるかどうかは M1 の写像の値として定まり、本節はこれとは別の持続の規則を置かない。
したがって、ここでいうトークンの持続は内部ラベルの時刻間継続である。
ラベルの継続 $k\in\mathcal K_t\cap\mathcal K_{t+1}$ は、二時刻の M1 出力の比較として作動記録の上で定まる。
同じラベルが時刻をまたいで同じ外部の枠 $a\in\mathcal A_O$ または同じ個物を指すかどうかは、
記述者がラベルと外部との対応を与えたときに問える別の条件であり、本論文はその対応を定義しない。
\end{remark}

%----------------------------------------------------------------------
\subsection{Updateの観測経由性}
\label{subsec:s5_observation}
\begin{assumption}[Updateの観測経由性]
\label{ass:s5_observation_factorization}
全域写像 $\Phi_{\mathrm{upd}}:\mathcal M\times\mathcal I\to\mathcal M$ が存在し、
\[
\mathrm{Update}(M,x,e)=\Phi_{\mathrm{upd}}(M,P_{\mathrm{sens}}(x,e))
\qquad\forall(M,x,e)\in\mathcal M\times X\times E_{\mathrm{ext}}
\]
が成立するとする。導入地位：S（構造的）。強さ：中。
本仮定は観測以外の状態差・環境差が更新へ直接影響する系を除く追加条件であり、
\S\ref{sec:acs_def} のACSの構成要件とはしない。$\Phi_{\mathrm{upd}}$ の連続性は仮定しない。
依拠箇所は命題~\ref{prop:s5_observation_causality} であり、P1〜P3は本仮定を要しない。
\end{assumption}

\begin{proposition}[観測経由性の検査条件と内部モデルの因果性]
\label{prop:s5_observation_causality}
仮定~\ref{ass:s5_observation_factorization} は、各固定 $M$ について $P_{\mathrm{sens}}(x,e)=P_{\mathrm{sens}}(x',e')$ ならば
$\mathrm{Update}(M,x,e)=\mathrm{Update}(M,x',e')$ となることと同値である。
そのもとで $M_t$ は $M_0$ と $I_0,\ldots,I_{t-1}$ のみで定まる。
\end{proposition}
\begin{proof}
因数分解があれば同じ入力 $M,I$ は同じ更新値を返す。
逆に同じ感知情報を与える引数上で更新が一致するとする。
$I\in\operatorname{Im}P_{\mathrm{sens}}$ では、対応する更新値を $\Phi_{\mathrm{upd}}(M,I)$ と定めると一意である。
像外では $\mathcal M\ne\emptyset$ により固定した一元を値として延長する。
これで全域 $\Phi_{\mathrm{upd}}$ を得る。後半は $M_{t+1}=\Phi_{\mathrm{upd}}(M_t,I_t)$ の帰納法による。
$t=0$ では入力履歴は空であり、$M_0$ 自体が必要である。
\end{proof}

\begin{remark}[因数分解の射程]
\label{rem:s5_observation_scope}
観測経由性は真値を知ることと同値ではない。
真値を知らなくても、未観測の基盤変化が内部モデルへ作用する実現はこの仮定を満たさない場合がある。
また $\Phi_{\mathrm{upd}}$ の存在だけから \S\ref{sec:acs_def} の Update$_1$〜Update$_3$ の詳細型を実現できるとは限らない。
本節では既存の成分分解と因数分解の双方を満たす更新を対象とする。
$\Phi_{\mathrm{upd}}$ は更新成分の表現であり、更新の構造的役割を感知から一意に導くものではない。
7タプルの非冗長性は備考~\ref{rem:tuple_nonredundancy} の役割分離の意味で維持される。
\end{remark}

\begin{remark}[予測外と内部モデルの更新]
\label{rem:s5_prediction_learning}
定義~\ref{def:prediction_record_check} の理論上の予測外は、実現値 $o_{t+1}$ との照合量である。
仮定~\ref{ass:s5_observation_factorization} のもとでは
命題~\ref{prop:s5_observation_causality} により $M_s$ は $M_0$ と $I_0,\ldots,I_{s-1}$ のみで定まるので、
$o_{t+1}$ が内部モデルに関与するのは感知情報の列を経る範囲に限られる。
本節は予測外にもとづく更新の規則を与えない（備考~\ref{rem:prediction_reference_difference}）。
また、予測精度係数 $\pi_{\tau,t}$（定義~\ref{def:s5_prediction_coefficient}）は
記述者が指定する全域の係数であり、予測集合が $\bot$ の時刻にも値をとる。
二種類の予測外は $\bot$ の時刻に定まらないので、本論文は両者を同一視せず、一般的な関係も置かない。
\end{remark}

%----------------------------------------------------------------------
\subsection{時刻間変化}
\label{subsec:s5_intertime}
\label{subsubsec:intertime_update}
\begin{definition}[担体から誘導される時刻間変化]
\label{def:intertime_update}
非空集合 $\mathsf D_{\mathrm{rd}}$ と固定読み出し
$\mathrm{rd}:X\times\mathcal M\times E_{\mathrm{ext}}\to\mathsf D_{\mathrm{rd}}$ を指定した場合、
$\mathrm{rd}_t=\mathrm{rd}(x_t,M_t,e_t)$ とし、次時刻値は
\[
\mathrm{rd}_{t+1}=\mathrm{rd}\bigl(f(x_t,u_t,e_t),\mathrm{Update}(M_t,x_t,e_t),e_{t+1}\bigr),
\qquad e_{t+1}=\bigl(b_{t+1},g_O(o_t,x_t,u_t,b_t)\bigr)
\]
で定める。この型保存に独立更新写像は不要である。
$e_{t+1}$ を要する場合、$t$ 時点での内部予測とはしない。
軌道を引数とする量は、当該軌道と指定した読み出し規則から評価する。
\end{definition}

\begin{definition}[学習・忘却の時刻間的位置づけ（分類B）]
\label{def:learning_intertime_change}
学習・忘却に関わる持続的内容の変化は $M_t^{\mathrm L}$ の変化として記述する。
その変化を改善・悪化と同一視しない。作動成分だけの変更は、この定義では学習と呼ばない。
系列上の学習停止は定義~\ref{def:update_stop_on_trajectory} に従う。
$\delta_{M,c,t},\delta_t,\lambda_{T,t}$ は作動成分も含む変化量のままとする。
\end{definition}

\begin{remark}[高速処理・Automaticityと状態の変化]
\label{rem:s5_fast_change}
\S\ref{sec:tipa} の作動形態は、学習成分の変化に伴って形成・変容しうる。
その成立を学習成分だけの関数とはしない。作動成分、現在入力、基盤・感度の状態も
各時刻の作動形態に関与しうる。
本論文は、学習停止から高速処理・Automaticityの不変を一般的な帰結として置かず、独立した形成速度や学習則も置かない。
感度・基盤・効果器・健全度の時刻変化は $x_t$ の遷移を通じて記述する。
\end{remark}

\begin{remark}[学習と可行性構造]
\label{rem:s5_learning_viability}
学習（定義~\ref{def:learning_intertime_change}）は $M_t^{\mathrm L}$ の変化であり、
7タプルは時刻によらない（\S\ref{sec:acs_def} の記述水準の規約）。
したがって $K$、可行行動対応 $A$、$\mathrm{Viab}_f(K)$ および
核維持応答集合 $U^{\mathrm{Viab}}$（定義~\ref{def:kernel_response}）は学習によって変わらない。
学習が関与しうるのは、作動記録の上でどの実効行動が選ばれるか、
とくに $u_t\in U^{\mathrm{Viab}}(x_t)$ が成り立つか（定義~\ref{def:kernel_response_acquisition}）と、
それを通じた軌道である。
身体の変化のように状態の変化として表されるものは、$x_t$ の遷移として $f$ を通じて記述され、
これらの集合を変えない。
本論文は、学習が核維持応答の獲得を導くことを主張しない（備考~\ref{rem:kernel_response_scope}）。
\end{remark}

%----------------------------------------------------------------------
\subsection{分類C記述量}
\label{subsec:s5_internal_quantities}

\subsubsection{偏向性と偏向強度の位置づけ}
\label{subsec:beta_dyn}

\begin{definition}[偏向性（分類C）]
\label{def:bias_tendency}
原始的な機能（\S\ref{sec:acs_def}、定義~\ref{def:primitive_functions}）の
いずれかの区分の責務の遂行に、指定した参照基準に相対して現れる一つの偏りの傾向を、\emph{偏向性}と呼ぶ。
ここでいう責務の遂行には、そのための処理資源・注意その他の配分または利用を含める。
本節の抽象的構造記述では、記述の対象とする期間を通じて考慮する偏向性を
添字づける非空集合 $\mathcal J_{\beta}$ を置き、
各 $\iota\in\mathcal J_{\beta}$ を一つの偏向性の添字とする。
$\mathcal J_{\beta}$ は時刻によらず固定し、その有限性は仮定しない。
$\mathcal J_{\beta}$ による偏向性の切り分けは記述の解像度に相対的であり、
一つの偏向性をより細かい複数の偏向性に分けて記述することも、その逆も排除しない。
本論文は、複数の偏向性が併存すること、
およびそれらが互いに異なる方向をもつことを排除しない。
添字 $\iota$ は記述上の識別子であり、偏向性の方向および内容を表さない。
各偏向性には、その偏りを定める参照基準を記述の対象ごとに指定する。
参照基準は、対象とする責務の遂行において偏りがないとする状態の取り決めであり、
規範的な理想、正しい遂行または最適な遂行を表さない。
本論文はすべての対象に共通する参照基準を固定しない。
同じ添字について時刻間の強度を比較する記述では、参照基準と尺度を対象期間を通じて固定する。
\end{definition}

\begin{definition}[偏向性の帰属区分（分類C）]
\label{def:bias_attribution}
原始的な機能の四区分を表す添字の集合を
\[
  \mathsf{Prim}:=\{\mathrm{sens},\mathrm{disc},\mathrm{val},\mathrm{act}\}
\]
とする。各添字は感知・識別・評価・行動の区分を表し、
定義~\ref{def:three_primitives} の標準写像形式の添字と対応させて用いる。
各偏向性に、その偏りが現れる区分の空でない集合を対応させる写像
\[
  \mathrm{loc}_{\beta}:\mathcal J_{\beta}\to\mathcal P(\mathsf{Prim})\setminus\{\emptyset\}
\]
を\emph{帰属写像}と呼び、$\mathrm{loc}_{\beta}(\iota)$ を偏向性 $\iota$ の\emph{帰属区分}と呼ぶ。
$|\mathrm{loc}_{\beta}(\iota)|=1$ のとき偏向性 $\iota$ は単一の区分に帰属し、
$|\mathrm{loc}_{\beta}(\iota)|\ge2$ のとき複数の区分に複合的に帰属する。
帰属区分は時刻によらず定める。
定義~\ref{def:bias_tendency} の参照基準は、この帰属区分の責務に対応させて指定する。


\end{definition}

\begin{remark}[帰属区分の地位]
\label{rem:bias_attribution_status}
帰属区分は、偏向性を原始的な機能のどの区分の偏りとして記述するかを
理論の側で指定するものであり、
ACS が自らの偏向性を内部で表現または把握することを要求しない
（\S\ref{subsec:description_levels_global}）。
帰属区分は記述の対象ごとに指定するものであり、
本論文は、観測された振る舞いから帰属区分が一意に定まることを主張しない。
同じ振る舞いの偏りが、異なる区分に帰属する記述と両立しうるためである。
知覚的意思決定の研究では、観測される選択の偏りが、偏りの種類や実験条件に応じて、
証拠蓄積の開始点の変化として説明される場合と証拠蓄積の過程の変化として説明される場合があり、
いずれに帰するかは拡散モデル等の当てはめと比較によって判別されている~\cite{Mulder2012,Urai2019}。
こうした判別について、観測者モデルを構築し検証することに重点を置くべきだとする議論がある~\cite{RahnevDenison2018}。
帰属区分の空でない部分集合は $2^{4}-1=15$ 通りあるが、
本論文はこれを記述のための分類として用いるにとどめ、
特定の組み合わせが現れやすいといった規則性を主張しない。
\end{remark}

\begin{definition}[偏向強度（分類C）]
\label{def:beta_accumulated}
各 $\iota\in\mathcal J_{\beta}$ と時刻 $t$ について、
偏向性 $\iota$ の時刻 $t$ における参照基準からの偏りの程度を表す非負の量を $\beta_{\iota,t}$ と書き、
偏向性 $\iota$ の\emph{偏向強度}と呼ぶ。
$\beta_{\iota,t}=0$ は、時刻 $t$ において偏向性 $\iota$ に参照基準からの偏りがないことを表す。
偏りの程度を非負の実数で表す尺度は偏向性ごとに指定し、
本論文は異なる偏向性の偏向強度の値どうしが比較可能であることを仮定しない。
偏向性ごとの偏向強度の族を
\[
  \boldsymbol{\beta}_t:=\bigl(\beta_{\iota,t}\bigr)_{\iota\in\mathcal J_{\beta}}
\]
と書く。
ここでは非負実数として定め、本節で採用する有界な読み出しの型は
後述の定義~\ref{def:s5_readout_types} に従う。
参照基準または尺度が未指定の場合、あるいは強度が未測定・未同定の場合に、
そのことだけを理由として $\beta_{\iota,t}=0$ を割り当てない。
$\beta_{\iota,t}$ は偏向性 $\iota$ の程度を表し、その方向および内容を表さない。
$\boldsymbol{\beta}_t$ は偏向強度の族であり、
それらを一つの実数へ集約した量ではない。
\end{definition}

偏向強度は、能力そのものではなく、能力・資源の使い方に現れる偏向の記述量である。
\S\ref{sec:acs_def} の認知的処理能力の記述方針
（備考~\ref{rem:cognitive_capacity_policy}）のとおり、
本論文は偏向強度を認知的処理能力に含めない。
偏向性の方向および内容の形式化は Phase~2 の課題とする（\S\ref{subsec:limitations_list} の L7）。

\begin{remark}[偏向性の抽象的構造記述と実装記述の区別]
\label{rem:beta_multiple_biases}
定義~\ref{def:bias_tendency} と定義~\ref{def:beta_accumulated} は、
偏向性と偏向強度を抽象的構造記述の水準で定める。
この水準では、複数の偏向性が存在しえ、それぞれが偏向強度をもつ。

これに対し実装記述の水準では、一つの偏向性とその偏向強度は、
更新の内容によって時刻とともに変容しえ、また他の偏向性と統合されうる。
個体ACSが学習によって形成する偏向性についても、
実装・設計によって与えられる偏向性についても同様である。
本論文は、実装上のどの内容が抽象的構造記述のどの添字 $\iota$ に対応するかを
時刻を通じて固定しない。
抽象的構造記述でこの変容・統合を表す場合には、
定義~\ref{def:beta_intertime_update} の読み出しを担う状態・内部モデルの変化として記述する。
偏向強度の族に独立した更新写像は置かない。

本論文は、複数の偏向強度を一つの実数へ集約する一般式を置かず、
方向の異なる偏向性が互いに打ち消し合うとも強め合うとも仮定しない。
実装上の偏向性と添字との対応づけ、各偏向性の方向と内容の表現、
および偏向性どうしの関係の形式化は Phase~2 の課題とする（\S\ref{subsec:limitations_list} の L7）。
\end{remark}

\begin{remark}[偏向性の形成と帰属の区別]
\label{rem:bias_formation_vs_attribution}
偏向性は、\S\ref{sec:tipa} の TIPA の三機構 $M_1$・$M_2$・$M_3$、PSL および ISL の
いずれか、またはそれらの複合的な作用によって形成されうる。
本論文は、偏向性がどの機構の作動によって形成されるかと、
それがどの区分に帰属するか（定義~\ref{def:bias_attribution}）とを区別し、
両者の間に一対一の対応を置かない。
これは、TIPA の三機構と原始的な機能の区分の間に
一対一の対応を置かない \S\ref{sec:acs_def} の扱いと同じである。
基盤機能（定義~\ref{def:base_functions}）に関わる PSL・ISL の状態に由来する偏りも、
それが現れる区分に帰属させて記述する。
偏向性の形成を機構ごとに記述すること、形成の組み合わせと帰属区分の対応、
および帰属区分の同定は、Phase~1.$\alpha$ または Phase~2 の課題とする
（\S\ref{subsec:limitations_list} の L7、\S\ref{subsec:phase1a}）。
\end{remark}

\begin{remark}[偏向強度の値と良否の区別]
\label{rem:beta_value_not_quality}
本論文は、各 $\iota\in\mathcal J_{\beta}$ について、
$\beta_{\iota,t}$ の値の大小を、改善、悪化、誤りまたは非合理性の
いずれとしても読まない。
偏向性が当該ACSの可行性その他にどのように関わるかは、
偏向性の内容と、それが働く環境および内部状態に依存しうるため、
$\beta_{\iota,t}$ の値だけからは定めない。
また、$\beta_{\iota,t}$ の増加、維持、減少をいずれも許し、
$\beta_{\iota,t}$ の値を定めることと、偏向が一方向へ累積し続けることとを同一視しない。
本論文は、$\beta_{\iota,t}$ の増減から予測精度の変化方向を一般的な帰結として置かない。
\end{remark}

\begin{remark}[偏向強度と状態推定誤差の区別]
\label{rem:beta_vs_estimation_error}
偏向強度は、定義~\ref{def:bias_tendency} の参照基準に相対した
原始的な機能の責務の遂行上の偏りの程度である。
状態推定誤差 $\varepsilon_{\mathrm{est},t}$
（定義~\ref{def:s4_state_estimation_error}）は、内部推定と理論上の状態との隔たりであり、
偏向強度とは参照する対象と尺度が異なる。
本論文は、偏向強度が正であることと状態推定誤差が正であることの間に、
いずれの向きの一般的含意も置かない。
この区別は、それらの間に関係がないことや統計的独立性を主張するものではない。
\end{remark}

\subsubsection{分類C読み出しの型と偏向強度の時刻間変化}
\begin{definition}[読み出しの型と解析範囲]
\label{def:s5_readout_types}
有限な $\beta_{\max}>0,D_{\max}>0$ と予測段数 $\tau\in\mathbb N$、$\tau\ge1$ を固定する。
$\mathcal H=\{\mathrm{psl},\mathrm{mem},\mathrm{sens},\mathrm{eff}\}$ とし、
本節で採用する記述規則を次の全域写像として指定する。
\[
\begin{aligned}
\mathrm{rd}_{\beta,\iota} &:X\times\mathcal M\times E_{\mathrm{ext}}\to[0,\beta_{\max}]
 &&(\iota\in\mathcal J_\beta),\\
\mathrm{rd}_{h,q} &:X\to[0,1]&&(q\in\mathcal H),\\
\mathrm{rd}_d &:X\times\mathcal M\to[0,D_{\max}],\\
\mathrm{rd}_{\pi,\tau} &: (X\times\mathcal M\times E_{\mathrm{ext}})^{\tau+1}\to[0,1].
\end{aligned}
\]
$\mathcal J_\beta$ は定義~\ref{def:bias_tendency} による。
$\mathrm{rd}_{\beta,\iota}$ は、同定義の参照基準と定義~\ref{def:beta_accumulated} の尺度に
対応する偏向強度を返す規則として指定する。
各偏向性の尺度は、零点の意味を保ったまま $[0,\beta_{\max}]$ に収まるよう正規化する。
本論文は、共通の値域を用いることから、異なる偏向性どうしの比較可能性を一般的な帰結として置かない。
これらは具体的規則を一意に導く主張ではなく、対象ごとに指定する記述規則の型である。
各値域には通常の部分空間位相を入れるが、読み出しの連続性は課さない。
有限上界への正規化が可能な記述を本節の解析範囲とし、
任意の現象がこの範囲に収まると主張しない。
現在の担体だけで定まらない量には、この瞬時読み出しを適用しない。
履歴を要する別の記述は、独立状態を加えず作動記録上の読み出しとして別途指定する。
\end{definition}

\begin{definition}[偏向強度の時刻間読み出し]
\label{def:beta_intertime_update}
定義~\ref{def:s5_readout_types} の規則を固定し、
$\beta_{\iota,t}=\mathrm{rd}_{\beta,\iota}(x_t,M_t,e_t)$ とする。
族 $[0,\beta_{\max}]^{\mathcal J_\beta}$ には積位相を入れる。
各因子がコンパクトなので積はコンパクトであるが、
これを独立状態空間や更新写像の定義域として用いない。
\end{definition}

\subsubsection{予測精度係数と実効的時間地平}
\label{subsubsec:s5_prediction_coefficient}
\label{subsec:tau_hat_section5}
\begin{definition}[予測精度係数（分類C）]
\label{def:s5_prediction_coefficient}
標準作動記録の状態・環境部分列 $z_s=(x_s,M_s,e_s)$ について
$\pi_{\tau,t}=\mathrm{rd}_{\pi,\tau}(z_t,\ldots,z_{t+\tau})\in[0,1]$
を、時刻 $t$ の予測の有効性を事後的に記述する係数とする。
採用する規則は予測対象・照合方法を対象ごとに固定する。
その数値化・測定可能性・予測成功を本節では保証しない。
未来の実現値を参照する分類C記述であり、ACSの時刻 $t$ の内部入力にはしない。
\end{definition}

\begin{definition}[実効的時間地平（分類C）]
\label{def:tau_hat_section5}
$\widehat\tau_t:=\tau\pi_{\tau,t}\in[0,\tau]$ と定める。
したがって
$\widehat\tau_{t+1}-\widehat\tau_t=\tau(\pi_{\tau,t+1}-\pi_{\tau,t})$ である。
$\tau$ は離散的な予測段数、$\widehat\tau_t$ は実数値の尺度であり、
$x_{t+\widehat\tau_t}$ のような非整数時刻での状態を導入しない。
\end{definition}

\subsubsection{実効的発展自由度}
\label{subsec:D_eff_section5}
\begin{definition}[実効的発展自由度（分類C）]
\label{def:deff_section5}
$d_{\mathrm{eff},t}:=\mathrm{rd}_d(x_t,M_t)\in[0,D_{\max}]$ を、
発展上の実効的自由度の記述量とする。
値域と担体を指定するが、算出規則の具体化は Phase~1.$\alpha$ に委ねる。
これは $|A_t|$、M2の行動多様度または関係品質（\S\ref{subsec:ro_quality}）とは異なる。
本論文は、その数値から可行行動の存在・増減を一般的な帰結として置かない。
\end{definition}

\subsubsection{PSL・ISL健全度}
\label{subsec:health_dyn}
\begin{definition}[PSL・ISL健全度量（分類C）]
\label{def:s5_health_quantities}
$h_{q,t}:=\mathrm{rd}_{h,q}(x_t)\in[0,1]$ と定める。
$q=\mathrm{psl},\mathrm{mem},\mathrm{sens},\mathrm{eff}$ はそれぞれ
基盤、記憶実装、感知、作用の健全度を指定する。
学習内容そのものの評価は記憶実装の健全度と区別する。
共通単位、相互独立性、直接比較可能性を要求しない。
\end{definition}

\begin{remark}[健全度・機能軸・診断の区別]
\label{rem:s5_health_noncollapse}
健全度と \S\ref{sec:tipa} の機能軸は別の記述量であり、一般的な相互決定関係を置かない。
ACS内部の診断は利用可能な情報とTIPA機能の状態に依存し、
外部診断は観測情報と測定規則に依存する。
本論文は、$h_{q,t}=0$、$d_{\mathrm{eff},t}=0$、$\pi_{\tau,t}=0$ だけから、
閉塞、FR の時刻別充足述語が偽であること、機能的終端、崩壊を一般的な帰結として置かない。
\end{remark}

\begin{definition}[時刻依存量の整理用タプル（分類C）]
\label{def:extended_variable_tuple}
\label{subsubsec:desc_tuple}
\[
\mathrm{Desc}_t=(\boldsymbol\beta_t,d_{\mathrm{eff},t},\pi_{\tau,t},\widehat\tau_t,
 h_{\mathrm{psl},t},h_{\mathrm{mem},t},h_{\mathrm{sens},t},h_{\mathrm{eff},t})
\]
とする。これは読み出した量をまとめる八つ組であり、ACSの7タプルを変更しない。
独立状態でも、同時刻の内部利用可能情報の束でもない。
予測係数を含むため、記録の事後的評価としての性格をもつ。
\end{definition}

%----------------------------------------------------------------------
\subsection{後続節への受け渡し}
\label{subsec:s5_summary}
\label{subsubsec:s5_deff_interface}
\begin{definition}[後続節への記録インターフェース]
\label{def:s5_downstream_interface}
\S\ref{sec:rel_min} へは $A_t,d_{\mathrm{eff},t},\widehat\tau_t$ の列を渡す。
個体ごとの添字と有限窓差分の定義は \S\ref{sec:rel_min} に従う。
予測係数を含む窓の評価には窓端より先の実現記録も必要であり、
同時刻のオンライン診断可能性を含意しない。

\S\ref{sec:collapse_cond} へは標準作動記録と $\mathrm{Desc}_t$ を渡す。
\S\ref{sec:collapse_cond} の軌道記録に必要な $\mathrm{WF},\mathrm{FT},\mathrm{Col}$ その他の
実現水準の記録が別途与えられ、同節が軌道記録に前提とする動学との整合（備考~\ref{rem:trajectory_space_scope}）を満たす場合に、
標準作動記録をその軌道記録へ対応づける。
P2だけからこれらの指標や許容外的作用下の回復性が一意に定まるとはしない。
同一ACSとしての解釈は $I_{\mathrm{ACS}}(\sigma)$ に限定する。
崩壊をまたぐ予測窓を同一ACSの予測評価として用いる場合には、
その窓全体が同一ACSの記述域に含まれることを確認する。
\end{definition}

本節の仮定は観測経由性の一件である。
P1〜P3 は固定した全域写像の合成と自然数上の再帰による。
有限時間到達・回復・不変性・安定性・実装可能性は別の証明義務であり、
本節の型保存や分類C量の有界性からは主張しない。


%================ Section 6 ACS関係性の最小構造 ================
%======================================================================
\section{ACS関係性の最小構造}
\label{sec:rel_min}

本節では、
単一ACS $\Sigma_i$ と外部対象 $Y_j$ の間に成立しうる関係の最小構造を、
対象分類、
単方向関係記述、
双方向合成、
関係品質、
代表カテゴリーの順に定める。
関係を倫理的善悪へ分類するのではなく、
関係継続中に個体側で記録・評価しうる可行性関連量の変化を記述するための枠組みを与える。

ここでいう記録は分類Cの理論上の記述を含み、
ACS自身が各値を同じ尺度で観測・算出できることを要求しない。
前節 \S\ref{sec:extended_dynamics} の $d_{\mathrm{eff},i,t}$ と $\widehat\tau_{i,t}$ は再定義せず、
関係記録に付随する有限窓差分へ接続する。
本節の局所的サイズ指標 $\mathrm{sz}_i$ と差分記法 $\Delta_L$ も、
$\mathrm{Desc}_t$ の八成分を変更しない。
これらの差分は外部対象の単独因果効果を表さず、
\S\ref{sec:acs_def} が定める閉塞・崩壊の条件の代替にもならない。

以下、
\S\ref{subsec:rel_objects} で外部対象、
\S\ref{subsec:rel_unidirect} で単方向記述と有限窓差分、
\S\ref{subsec:rel_bidirect} で双方向合成、
\S\ref{subsec:ro_quality} で関係品質、
\S\ref{subsec:rel_categories} で18カテゴリー・六系分類を扱い、
\S\ref{subsec:rel_phase1a} で限界と後続節への接続を整理する。
本節で新たに用いる主要記号の一覧は付録~\ref{app:symbols_s6} に示す。

ここで単方向・双方向は関係を記述する水準の違いであり、
対象の価値や
関係がもたらす結果の優劣を区別するものではない。
同一ラベルの関係でも、
関係継続中のコスト、
環境条件および当該ACSの状態によって、
可行性関連量の異なる変化が記録されうる。
本節のカテゴリー名は、
その変化の符号や原因をあらかじめ決める規則ではない。

\begin{remark}[本節の射程]
\label{rem:rel_scope}
本節が対象とするのは、
外部対象の三分類、
単方向関係記述子、
ACS同士の双方向合成、
関係品質、
代表カテゴリー18種と六系分類である。
単方向記述子の有限窓差分は、
必要な時刻域と記録の条件が成立するときに限り定義される。
関係が存在すれば、
任意の対象・時刻で必ず六成分が得られるという全域性は要求しない。
外部対象の分類可能性、
関係カテゴリーの判定可能性、
有限窓差分の整形式性はそれぞれ別の問題である。
本節では、
実際に記述子を採用する場合にだけ必要な成分と時刻域を確認する。

カテゴリーの成立条件・閾値・強度・持続性・解除可能性・代替可能性・認識との対応は定めない。
決定論的な補助的詳細化はPhase~1.$\alpha$へ、
確率的成立・遷移・終了、
連続時間化および社会的ネットワーク動学はPhase~2へ委ねる
（\S\ref{subsec:limitations_list} の L8、\S\ref{subsec:phase1a}）。
\S\ref{sec:tipa} の識別・候補形成と、\S\ref{sec:acs_def}・\S\ref{sec:collapse_cond} の閉塞・崩壊の判定を本節で再定義しない。

外部対象の分類と各時刻の関係記録は区別するため、
分類可能な対象に対して
カテゴリーや差分が決まらないことは、
本節の整形式性を直ちに損なわない。
関係の時間発展や関係性ネットワークの生成則も本節の記述子には含めない。
\end{remark}

\begin{remark}[閉塞・崩壊の節との関係]
\label{rem:rel_to_sec7_intro}
\S\ref{sec:collapse_cond} は、$A_{i,t}=\emptyset$ の閉塞と
$\mathrm{Col}_{i,t}=1$ の崩壊を別個の条件で扱う。
本節の関係記録に付随する差分は、
閉塞経路集合 $\mathfrak{P}_{\mathrm{blk}}$ や
崩壊到達経路集合 $\mathfrak{P}_{\mathrm{col}}$ を分析する際の記述素材になりうるが、
差分の符号、
関係カテゴリー、
記述子の存在だけでそれらの経路集合への帰属を決定しない。
崩壊後の通常ACS状態量を差分計算のために要求することもない。

関係記録を閉塞・崩壊の経路記述に併記する場合にも、
経路への帰属を判定するのは後続節が定める状態・軌道の条件である。
記録がないことと、
関係が存在しないことも同値ではない。
\end{remark}

%----------------------------------------------------------------------
\subsection{外部対象集合と分類写像}
\label{subsec:rel_objects}

本小節では、
固定した $\Sigma_i$ から見た外部対象集合と三分類を導入する。

\begin{definition}[外部対象集合]
\label{def:external_object_set}
固定したACS個体 $\Sigma_i$ に対し、
自身を除き、
関係の対象として採用する外部対象の集合を
\[
  \mathcal Y_{-i}
\]
と書き、
個別の対象を $Y_j\in\mathcal Y_{-i}$ と表す。
$(-i)$ は自身を除く標識であり冪指数ではない。
この集合は記述する対象の領域であり、
各 $Y_j$ の内部状態空間、
意思決定機構、
自己維持機構、
または時刻を通じた確実な観測可能性を仮定するものではない。
$\mathcal Y_{-i}$ の非空性も、
全対象と関係が成立することも要求しない。

対象集合への帰属と、
特定の時刻にその対象との六成分記述子が定義されることを区別する。
「観測」「参照」「作用」「依存」「支援」等は関係対象になりうる様式を
言葉で示したものであり、
列挙した様式の成立条件を本定義で与えるものではない。

外部対象の添字 $j$ は、
関係記録中に同一対象を追跡するための標識であり、
対象にACSとしての構造を追加するものではない。
\end{definition}

\begin{remark}[外部対象集合の役割]
\label{rem:external_object_role}
$\mathcal Y_{-i}$ は $\Sigma_i$ 自身の状態空間 $X_i$ と異なる対象領域である。
外部対象 $Y_j$ と、
\S\ref{sec:acs_def} の外部環境状態 $e_t\in E_{\mathrm{ext}}$ も別の種類の対象であり、
$Y_j$ を $e_t$ と同一視したり、
$Y_j\in E_{\mathrm{ext}}$ と解釈したりしない。

非ACS対象には $X_j$、
$A_{j,t}$ または通常ACS状態量を自動的には割り当てない。
また、
必要環境要因の集合 $E$ と外部対象集合 $\mathcal Y_{-i}$ は別の記述対象であり、
ある $Y_j$ が必要環境の特定要因に一対一に対応するとは仮定しない。

$Y_j$ を \S\ref{sec:acs_def} の行動応答成分の枠 $a\in\mathcal A_O$ または背景成分、
あるいは \S\ref{sec:tipa} の内部対象トークン $k$ に対応させるかどうかは記述者が与えるものであり、
本論文はその対応を定義しない。
本節の記述子は $\Sigma_i$ 側の値だけを用いるので、この対応がなくても定義される。

本節は外部対象の内部動学を展開せず、
採用した関係記録に付随する $\Sigma_i$ 側の量を記述する。
対象が同じであっても、
異なるACSからは別の関係記録の対象となりうる。
その場合も、
記録の視点を変更しただけで対象側の状態方程式を得たとはみなさない。
\end{remark}

\begin{remark}[外部対象の識別と行動候補との区別]
\label{rem:external_object_cand}
外部対象の自己--環境境界への帰属に関わる内部識別は、
\S\ref{sec:tipa} のM1-bの識別状態（定義~\ref{def:id_vector}）に接続する。
境界帰属軸の値域は $\{\mathsf{self},\mathsf{env},\mathsf{und}\}$ であるが、
この値域指定は個々の $Y_j$ の同一性が常に内部で確定するという主張ではない。
本節の $\mathrm{type}_i$ は関係性の記述上の分類写像であり、
M1-b の識別状態そのものでも、
ACSが同じ分類を内部に保持するという仮定でもない。

また、
M3の $\mathrm{Cand}_t\subseteq U$ は行動候補集合
（定義~\ref{def:m3_candidate_set}）であり、
外部対象集合ではない。
適合評価 $\mathrm{fit}_t$（定義~\ref{def:m3_fit}）と対象同定も異なる操作である
（備考~\ref{rem:fit_vs_identity}）。
外部対象の知覚・識別がどの行動候補に接続されるかを、
本節では固定しない。
特に本論文は、
$Y_j\in\mathcal Y_{-i}$ から $Y_j\in\mathrm{Cand}_t$ を一般的な帰結として置かない。

対象を同定できたこと、
対象が当該時刻の要請に適合すること、
対象に向けた行動が候補として生成されることは、
いずれも異なる条件である。
本論文は、本節の分類ラベル $\mathrm{type}_i(Y_j)$ だけから、
これらの機構上の判定を一般的な帰結として置かない。

対象同定を本節で前提とする記録上の対応付けと、
M1-b が実際に同定できたかという機構上の判定は区別する。
\end{remark}

\begin{definition}[外部対象の三分類]
\label{def:object_type}
固定した $\Sigma_i$ に相対する全域的・単値の記述上の分類写像を
\[
 \mathrm{type}_i:\mathcal Y_{-i}\longrightarrow
 \{\mathrm{ACS},\mathrm{agent},\mathrm{obj}\}
\]
と定め、
その逆像として
\[
\begin{aligned}
\mathcal Y_{\mathrm{ACS},-i}&:=\mathrm{type}_i^{-1}(\{\mathrm{ACS}\}),\\
\mathcal Y_{\mathrm{agent},-i}&:=\mathrm{type}_i^{-1}(\{\mathrm{agent}\}),\\
\mathcal Y_{\mathrm{obj},-i}&:=\mathrm{type}_i^{-1}(\{\mathrm{obj}\})
\end{aligned}
\]
を定める。
三つの値は、
それぞれACSとして扱う外部対象、
ACSとしては扱わない行為主体、
非エージェント的対象という関係記述上の分類ラベルである。
本写像は対象の規範的地位や実際の内部認識の正確さを定めない。

これは三つの記述上のラベルへの割当てであり、
いずれかの逆像が空であることも許す。
各対象が現実には複数の性質を併せもつ可能性と、
記述に際して単一のラベルを割り当てることは矛盾しない。
\end{definition}

\begin{assumption}[外部対象分類の固定]
\label{ass:obj_stable}
本節で同一対象を追跡する期間中、
分類写像 $\mathrm{type}_i$ は時刻 $t$ によって変更されないものとする。
これは外部対象の現実の属性変化を禁じる仮定ではなく、
本節で用いる対象分類の時間的固定だけを求める。

導入地位：S（構造的・記述域の制限）。
強さ：弱（対象分類の変化を追跡しない範囲に限定）。
依拠箇所：\S\ref{subsec:rel_objects} および \S\ref{subsec:rel_bidirect}。
\end{assumption}

\begin{remark}[分類固定仮定の意味]
\label{rem:obj_classification}
仮定~\ref{ass:obj_stable} は、
採用した記述期間中の分類を固定するという制限であり、
対象の物理的性質、
行為主体性またはACS構造が現実に不変であるという主張ではない。
分類が途中で変わる対象を扱うには、
本節で固定した写像とは別の時間依存分類が必要になる。
その詳細な記述は本節では扱わない。

例えば、
同じ外部対象について記述途中で「非エージェント的対象」から
「エージェント」として扱う分類へ変更するには、
時刻依存の分類規則を別途定める必要がある。
本節はその規則を置かず、
分類を固定している期間内の関係記述に限る。
\end{remark}

\begin{remark}[三分類の排他性（定義からの帰結）]
\label{rem:object_partition}
定義~\ref{def:object_type} の $\mathrm{type}_i$ が全域的な単値写像なので、
異なる値の逆像は交わらず、
三つの逆像の和集合は $\mathcal Y_{-i}$ に一致する。
すなわち
\[
\mathcal Y_{\mathrm{ACS},-i}\mathbin{\dot\cup}
\mathcal Y_{\mathrm{agent},-i}\mathbin{\dot\cup}
\mathcal Y_{\mathrm{obj},-i}=\mathcal Y_{-i}.
\]
ここで $\dot\cup$ は互いに素な和集合を表す。
これは新たな独立仮定ではなく、
分類写像の定義から従う帰結である。

とくに各逆像が非空であることは従わず、
外部対象が存在しない場合にも分割の等式は成り立つ。
\end{remark}

\begin{remark}[三分類の位置づけ]
\label{rem:object_partition_role}
三分類の排他性は写像の値の排他性であり、
現実の対象が複数の性質を同時にもつことを否定しない。
分類は $\Sigma_i$ に相対し、
別のACS $\Sigma_k$ は異なる記述上の分類写像を採用しうる。
M1-b の未決定カテゴリ $\mathsf{und}$ は内部識別の値であり、
本節の三分類に第四の外部対象型を追加するものではない。

ある $Y_j$ に対して $\mathrm{type}_i(Y_j)=\mathrm{ACS}$ と記述することは、
ACS自身がその対象を $\mathsf{env}$ と認識済みであることも、
他のACSが同じ型を付すことも意味しない。
\end{remark}

\begin{remark}[外部対象の型と閉塞・崩壊]
\label{rem:objects_to_sec7}
本論文は、外部対象の型のみから $A_{i,t}=\emptyset$ または
$\mathrm{Col}_{i,t}=1$ を一般的な帰結として置かない。
$\mathcal Y_{-i}$ と $\mathrm{type}_i$ は関係記述の対象領域を定めるだけであり、
閉塞経路・崩壊到達経路の判定には \S\ref{sec:collapse_cond} の条件を別途用いる。
\end{remark}

%----------------------------------------------------------------------
\subsection{単方向関係記述子と有限観測窓}
\label{subsec:rel_unidirect}

本小節では、
局所サイズ指標と有限観測窓を先に定め、
その適用条件を満たす場合に六成分の単方向関係記述子を構成する。

\begin{definition}[可行行動集合の局所的サイズ指標（分類C）]
\label{def:rel_size_functional}
固定した $\Sigma_i$ の行動空間を $U_i$ とし、
評価対象とする非空の集合族
\[
 \emptyset\neq\mathcal D_{\mathrm{sz},i}\subseteq\mathcal P(U_i)
\]
および非負実数値の記述写像
\[
 \mathrm{sz}_i:\mathcal D_{\mathrm{sz},i}\longrightarrow\mathbb R_{\geq0}
\]
を置く場合、
$\mathrm{sz}_i$ を本節における\emph{局所的サイズ指標}と呼ぶ。
$\mathcal P(U_i)$ は $U_i$ の冪集合である。

「サイズ」は評価する集合に付した記述名であって、
型指定のみからは
集合包含に関する単調性、
非退化性、
有限集合の濃度との一致、
$\mathrm{sz}_i(\emptyset)=0$ などの性質を与えない。
したがって、
指標値の大小・差分を可行な行動の増減へ一般に読み替えず、
$\mathrm{sz}_i(A)=0\Rightarrow A=\emptyset$ も性質として置かない。

この指標を採用するか否か、
および評価域・具体形は本節の一般形式では固定しない。
評価域に $\emptyset$ が含まれない場合には $\mathrm{sz}_i(\emptyset)$ 自体を要求しない。

また、
同一個体の時刻間差分を扱うときは、
一つの評価域と一つの写像を
当該観測窓の両端に共通して用いることを条件とする。
写像を時刻間で取り替えた値の差を、
同じ指標の時間変化として読み替えない。

\end{definition}

本指標は、
\S\ref{sec:extended_dynamics} の $d_{\mathrm{eff},t}$、
\S\ref{sec:three_principles} で扱う可行性制約、または応答多様度と同一視しない。
本節で指標を実際に用いる場合に限り、
評価する可行行動集合が
$\mathcal D_{\mathrm{sz},i}$ に属することを確認する。

本論文は具体化を与えないので、
非自明なサイズ性を必要とする追加の比較や推論を一般的な帰結として置かない。
そのため、
$\Delta_L\mathrm{sz}_i(A_{i,t})$ の符号は、
同一指標を用いた数値差の符号として読む。
集合の包含関係、
濃度の大小、
応答多様度の変化について主張するには、
別途対応する条件が必要である。

例えば有限集合に対する濃度や、
一般の集合に対する固定精度の被覆数は、
将来採用しうる具体化の候補であるが、
本論文ではいずれも
$\mathrm{sz}_i$ の定義そのものとしては採用しない。

したがって $\mathrm{sz}_i$ と応答多様度 $\mu_{\mathrm{eff},t}$ を
同じ量とみなすことも、
その大小関係を仮定することもしない。

\begin{definition}[有限観測窓における差分記法（分類C）]
\label{def:finite_window_diff}
正の整数 $L\in\mathbb N_{\geq1}$ と非負整数時刻 $t\in\mathbb N_0$ をとる。
以下、
一つの関係記録または双方向合成を評価する間は、
窓長 $L$ を共通に固定する。
異なる窓長を採用した差分は一般に異なりうるため、
同じ $t,i,j$ でも
$L$ が異なる記録を同一の $\mathrm{Rel}_{i\leftarrow j,t}$ として混用しない。
記号の上付き・下付きに $L$ を重ねて表記しないのは、
固定窓長を文脈で指定するためである。

$t\geq L$ のとき、
有限観測窓を
\[
 W_{L,t}:=\{t-L,t-L+1,\ldots,t\}
\]
と定める。

個体 $\Sigma_i$ に対する同一の実数値量の時系列 $z_{i,s}$ が
両端 $s=t-L,t$ で定義されている場合、
その後向き差分を
\[
 \Delta_L z_{i,t}:=z_{i,t}-z_{i,t-L}
\]
と定義する。
$t<L$、
または両端のどちらかの値がない場合、
この差分は定義しない。

差分そのものの定義には、
窓内の中間時刻すべてで当該実数値が得られることまでは要求しない。

一方、
関係記録の成立を述べるときは、
以下で定める窓全体の別条件を課す。

本節で用いる量は、
$\mathrm{sz}_i(A_{i,s})$、
$ d_{\mathrm{eff},i,s}$、
$\widehat\tau_{i,s}$ の三種類である。

第一の量については
\[
 A_{i,t},\ A_{i,t-L}\in\mathcal D_{\mathrm{sz},i}
\]
を要求し、
かつ両端で同じ $\mathrm{sz}_i$ を用いる。

後二者は \S\ref{sec:extended_dynamics} の量を、
\S\ref{sec:acs_def} の記述水準の規約に従って個体 $\Sigma_i$ について
$d_{\mathrm{eff},i,s}$、$\widehat\tau_{i,s}$ と書いたものであり、
両端の値が定義される場合にのみ差を取る。
$\widehat\tau_{i,s}$ は時刻 $s$ から $s+\tau$ までの記録を用いるので
（定義~\ref{def:s5_prediction_coefficient}）、
その差分は、当該の値が定義される時刻についてだけとる。
各差分の実数値性はこの適用域の下での帰結であり、
実測可能性や内部計算可能性の一般的な主張ではない。

差分を外部対象 $Y_j$ との\emph{関係記録に付随する変化}として
単方向記述子に組み込む場合には、
さらに窓内の全時刻で
関係が継続し、
添字 $j$ が同一対象を指すことを要求する。

軌道記録 $\sigma$ に対して \S\ref{sec:collapse_cond} の
ACS記録有効時刻域を用いる場合には、
定義~\ref{def:acs_valid_time_domain} に従って
\[
 W_{L,t}\subseteq I_{\mathrm{ACS}}(\sigma)
\]
を要求する。
ここで $\sigma$ は記録対象である同一個体 $\Sigma_i$ の軌道記録を表す。
この条件はACSとしての通常記録の有効時刻域を指定するものであり、
全時刻で三つの差分対象量が実際に取得できることを意味しない。
条件を満たさない窓で未定義の値をゼロに置換しない。

$t=L$ の場合には左端が時刻 $0$ となり、
負時刻の履歴を参照しない。

また、
両端の差が計算可能でも、
窓内の関係継続や対象同定が確認できなければ、
当該対象との関係記録にその差を組み込む条件は満たされない。
\end{definition}

\begin{remark}[有限観測窓差分の制限とサイズ汎関数の委譲]
\label{rem:finite_window_diff_limit}
$\Delta_L$ は有限個の離散時刻の後向き差分であり、
連続時間微分でも、
対象 $Y_j$ が存在しなかった場合との差でもない。
$L$ は観測窓の時刻差を表し、
$W_{L,t}$ に含まれる観測時刻数は $L+1$ である。
$t<L$ では負の添字へ履歴を延長しない。

また、
$I_{\mathrm{ACS}}(\sigma)$ の外では通常ACS状態量を要求しない
（\S\ref{sec:collapse_cond}）。
崩壊が初めて成立する時刻をまたぐ窓に対して、
失われた量を補間して差分を作らない。
一方、
閉塞 $A_{i,s}=\emptyset$ は崩壊と同義ではないため、
閉塞時刻が窓に含まれることだけを理由に全ての差分を排除しない。

ただし、
閉塞が差分の端点 $s=t-L$ または $s=t$ で生じ、
$A_{i,s}=\emptyset\notin\mathcal D_{\mathrm{sz},i}$ であれば、
$\mathrm{sz}_i(A_{i,s})$ は定義域外であり、
$\Delta_L\mathrm{sz}_i(A_{i,t})$ を定義できない。
この場合、
$d_{\mathrm{eff},i,t}$ と $\widehat\tau_{i,t}$ の差分が
それぞれ整形式であっても、
定義~\ref{def:rel_unidirect} の六成分記述子は構成しない。

一方、
空集合が評価域に含まれる場合でも、
その指標値がゼロとなることは要求しない。
差分が定義できるか否かは、
各量の値と評価域、
および採用する軌道記録の時刻域で判定する。

指標 $\mathrm{sz}_i$ が導入されても、
本論文は、その値の減少から直ちに集合の包含関係や
$A_{i,t}=\emptyset$ を一般的な帰結として置かない。

指標の具体化、
有限近似、
固定精度の被覆数等はPhase~1.$\alpha$（\S\ref{subsec:phase1a}）、
測度論的な精密化はPhase~2で扱う。

ここでの「委譲」は、
任意のコンパクトな行動空間上に
非退化なサイズ指標が自動的に得られるとの存在主張ではない。
特定の集合族と指標を選ぶ場合に必要な性質は、
その具体化の側で別途検討する。
\end{remark}

\begin{definition}[単方向関係記述子（分類C）]
\label{def:rel_unidirect}
固定した $\Sigma_i$、
$Y_j\in\mathcal Y_{-i}$、
正の窓長 $L$ と時刻 $t$ をとる。
以下の六成分を記録する期間には定義~\ref{def:finite_window_diff} の窓長 $L$ を固定し、
別の窓長による記述と区別して用いる。

定義~\ref{def:finite_window_diff} の三差分が関係記録用の条件を満たして定義され、
さらにカテゴリー $\mathrm{cat}_{i\leftarrow j,t}\in\mathcal C_{\rightarrow}$、
コスト $\mathrm{cost}_{i\leftarrow j,t}\in\mathbb R_{\geq0}$、
摂動 $p_{\mathrm{pert},i\leftarrow j,t}\in P_{\mathrm{pert}}$ が記録されるとき、
$Y_j$ から $\Sigma_i$ に向かう単方向関係記述子を
\[
\begin{aligned}
 \mathrm{Rel}_{i\leftarrow j,t}:=\bigl(&
 \mathrm{cat}_{i\leftarrow j,t},\ \mathrm{cost}_{i\leftarrow j,t},\
 p_{\mathrm{pert},i\leftarrow j,t},\\
 &\Delta_L\mathrm{sz}_i(A_{i,t}),\
 \Delta_Ld_{\mathrm{eff},i,t},\
 \Delta_L\widehat\tau_{i,t}\bigr)
\end{aligned}
\]
と定める。
$\mathcal C_{\rightarrow}$ は定義~\ref{def:category_set} の15ラベル、
$P_{\mathrm{pert}}$ は定義~\ref{def:rel_unidirect_type} の非空な摂動値域である。

前三成分は関係カテゴリー・非負コスト・摂動についての記録、
後三成分は関係継続中に記録・評価された $\Sigma_i$ 側の値の差分である。
前三成分の記録が成立しても、
後三成分が未定義なら六成分の記述子は置かない。
逆に三差分が定義できても、
関係カテゴリー等の記録を採用しなければ
特定の $Y_j$ に紐づく六成分記述子を構成したことにはならない。

下付きの $i\leftarrow j$ は記述方向であり、
後三成分の式自体には $j$ が現れない。
したがって、
記述子は $Y_j$ の一意の因果的寄与を同定するものではない。
同一の $\Sigma_i$ と観測窓に複数の外部対象との関係が記録される場合、
後三成分に同じ数値差が複数の記述子で現れうる。
その一致を対象間の同等の作用や寄与量の配分と解釈しない。
本定義は、
上記条件が満たされない対象・時刻にも六成分を割り当てるものではない。
関係のコストはその関係に付随する資源負荷の非負記録であり、
$\mathrm{cost}_{i\leftarrow j,t}=0$ は関係が成立しないことを意味しない。
摂動の記録は $P_{\mathrm{pert}}$ の型に従うが、
記号 $p$ は確率を表さず、
摂動値の正負や大きさによるカテゴリー判定も本節では行わない。

後三成分の差が同じであっても、
前三成分のカテゴリー・コスト・摂動は異なりうる。
\end{definition}

\begin{remark}[対象添字の時間的持続と差分の非因果性]
\label{rem:object_index_persistence}
添字 $j$ は $W_{L,t}$ の全時刻にわたって同一の外部対象を指し、
当該対象との関係が継続していることを、
関係記録を採用する条件とする。
対象同定は \S\ref{sec:tipa} のM1-bの識別（定義~\ref{def:id_vector}）に関係するが、
識別状態の値から個体同一性を確定する判定規則は本節では定めない。
特に、
境界帰属軸が $\mathsf{env}$ を示すことと、
異なる時刻の記録が
同じ個体 $Y_j$ に属すると判定されることは別の条件である。
本節の添字 $j$ の持続は、
当該窓の記録を採用する際に置く対応付け上の要請であって、
M1-b の内部同定手続きの存在や正確さを追加で仮定するものではない。
M3の要請適合評価は対象同定とは別の機能である（備考~\ref{rem:fit_vs_identity}）。

三差分は $\Sigma_i$ 側の値に依拠し、
右辺に $j$ を含まない。
同時に存在する他の対象、
行動、
環境条件等の寄与を差し引く規則も置かない。
よって差分は「関係継続中に記録された変化」であり、
反実仮想的な因果効果、
関係カテゴリーの成立証明または対象別の寄与配分ではない。
例えば二つの対象との関係が同一の窓で継続していた場合、
それぞれの記述子には同じ $\Sigma_i$ 側の三差分が併記されうる。
両記録の差分値を足し合わせて個体側の変化を二重に計上したり、
各対象がその差分の全量を生じさせたと解釈したりしない。
\end{remark}

\begin{remark}[記号上の注意]
\label{rem:rel_unidirect_notation}
$\mathrm{Rel}_{i\leftarrow j,t}$ の向きは記述視点を示すものであり、
一方的な決定・支配を表す記号ではない。
$p_{\mathrm{pert},i\leftarrow j,t}$ の $p$ は perturbation（摂動）の略記であり、
probability（確率）ではない。
この区別により、
本論文の $p_{\mathrm{pert},i\leftarrow j,t}$ には
分布、
期待値または確率的な推移規則を割り当てない。
また、
矢印は情報の記述方向であり、
対象が意図をもつことまで要求しない。
\end{remark}

\begin{definition}[単方向関係記述子の型]
\label{def:rel_unidirect_type}
単方向記述子が定義される時刻での六成分の型を
\[
\begin{gathered}
 \mathrm{cat}_{i\leftarrow j,t}\in\mathcal C_{\rightarrow},\qquad
 \mathrm{cost}_{i\leftarrow j,t}\in\mathbb R_{\ge0},\\
 p_{\mathrm{pert},i\leftarrow j,t}\in P_{\mathrm{pert}},\\
 \Delta_L\mathrm{sz}_i(A_{i,t})\in\mathbb R,\qquad
 \Delta_L d_{\mathrm{eff},i,t}\in\mathbb R,\\
 \Delta_L\widehat\tau_{i,t}\in\mathbb R
\end{gathered}
\]
とする。
$\mathcal C_{\rightarrow}$ は定義~\ref{def:category_set} に従う。
$P_{\mathrm{pert}}$ は摂動的影響を表す非空の値域であり、
値の符号・順序・具体的尺度は本節では固定しない。
$P_{\mathrm{pert}}$ が順序を備えるとも、
摂動量をコストや三差分と
算術的に合算できるとも仮定しない。
三差分が実数であることは、
差分の定義域と \S\ref{sec:extended_dynamics} の
各量の値域が成立する場合の整形式性を表す。
$P_{\mathrm{pert}}$ の具体化は有限ラベル集合、
順序集合、
実数値区間等として
検討しうるが、
本論文はどれか一つを採用しない。

したがって摂動の大小を用いた比較は、
追加の型を定めない限り一般には行えない。
\end{definition}

\begin{remark}[型指定の役割]
\label{rem:rel_unidirect_type_role}
型の指定は各成分の種類を与えるだけで、
計測法、
推定誤差、
カテゴリー成立閾値、
摂動の加法性または強度尺度を与えない。
本論文は、$\mathrm{sz}_i$ の差分が正であることから行動集合の包含関係での拡大を、
負であることから閉塞を、いずれも一般的な帰結として置かない。

分類Cの記述子全体をACS内部が同じ形式で計算・保持することも要求しない。
\end{remark}

\begin{remark}[関係性記述子と可行性関連量]
\label{rem:rel_descriptor_viability}
$\mathrm{sz}_i(A_{i,t})$ は本節で採用する局所的指標であり、
$ d_{\mathrm{eff},i,t}$ と $\widehat\tau_{i,t}$ は
\S\ref{sec:extended_dynamics} の個体別量を参照する。
本論文は、三差分が正・零・負のいずれでも、
その符号から直ちに
$A_{i,t}=\emptyset$、
$\mathrm{Col}_{i,t}=1$、
回復可能性、
三原理の時刻別充足述語の真偽または特定対象の因果的寄与を一般的な帰結として置かない。

これらは異なる定義・条件によって扱われる量である。
\end{remark}

\begin{assumption}[単方向関係記述の対象範囲]
\label{ass:unidirect_scope}
単方向記述子を実際に採用する対象 $Y_j$、
窓長 $L$、
時刻 $t$ について、
カテゴリー・コスト・摂動が定義~\ref{def:rel_unidirect_type} の型で記録され、
三差分が定義~\ref{def:finite_window_diff} の条件を満たすことを要求する。
本仮定は、
任意の外部対象・任意の時刻に関係が成立すること、
六成分を記録できること、
および記録を採用する組の存在や非空性を要求しない。

この要請は採用済みの記録の整形式性を条件付きで指定するものであり、
外部対象とACS状態との関係が一般に計測可能であるとの存在仮定ではない。

導入地位：S（構造的・採用する記録の整形式性）。
強さ：弱（記述子を採用する場合の成分取得条件に限る）。
依拠箇所：\S\ref{subsec:rel_unidirect} および \S\ref{subsec:rel_bidirect}。
\end{assumption}

\begin{remark}[非ACS対象への適用]
\label{rem:unidirect_nonacs}
非ACSエージェントと非エージェント的対象も、
$\Sigma_i$ 側の六成分が
記録条件を満たす場合には単方向記述の対象となりうる。
その際、
$Y_j$ 側の $X_j$、
$A_{j,t}$、
$d_{\mathrm{eff},j,t}$ 等を仮定しない。
記録が採用されない対象や時刻について、
未知の値をゼロで埋めた記述子を作らない。
非エージェント的対象との関係であっても、
対象が個体側の選択を制約する場合や
資源負荷を伴う場合に、
個体側の差分を併記することは形式上可能である。
ただし、
これを対象自身による意思決定や自己維持の存在の証拠とはしない。
\end{remark}

\begin{remark}[閉塞・崩壊の節への接続]
\label{rem:to_section7}
関係記録は \S\ref{sec:collapse_cond} の経路記述への入力となりうるが、
$\mathrm{Rel}_{i\leftarrow j,t}$ の存在や差分の符号だけで
$\sigma\in\mathfrak{P}_{\mathrm{blk}}$ または
$\sigma\in\mathfrak{P}_{\mathrm{col}}$ を決定しない。
関係記録が経路の分析に併記される場合、
記述対象の軌道における
閉塞・崩壊の成立条件は \S\ref{sec:acs_def} の定義（定義~\ref{def:blockage_sec2}、定義~\ref{def:collapse}）に従う。
関係性のカテゴリー名は、
その条件の代わりになる論理式ではない。
\end{remark}

%----------------------------------------------------------------------
\subsection{双方向合成関係}
\label{subsec:rel_bidirect}

本小節では、
双方をACSとして扱える対象間で、
両方向の記述がともに整形式な時刻に二つの記述子を組にする。

ただし、
双方向合成は二つの単方向記録の順序対であり、
両個体の内部状態を結合した新しいACSや
共通の状態空間を構成する操作ではない。

\begin{assumption}[双方向合成関係の標準対象]
\label{ass:bidirect_scope}
本節の双方向合成の標準対象は、
$Y_j\in\mathcal Y_{\mathrm{ACS},-i}$ を
ACS個体 $\Sigma_j$ として扱い、
かつ $\Sigma_i\in\mathcal Y_{\mathrm{ACS},-j}$ として
逆向きの記述対象も認められる組に限る。

両方向の六成分記述子が、
共通の正の窓長 $L$ と同じ時刻 $t$ でそれぞれ定義できる場合にのみ合成する。
この条件は、
両個体が同じ局所サイズ指標 $\mathrm{sz}_i=\mathrm{sz}_j$ や
同じカテゴリー・コスト・摂動値を採用することを要求しない。
一方向の記述子があることだけから、
逆方向の記述子の存在を要求しない。
双方の対象同定と関係継続は、
各個体側の記録としてそれぞれ扱う。
両者が同じ対象を指していることが記述上確認されても、
両者の内部モデルが相互に同じ内容を保持するとまでは仮定しない。

導入地位：S（構造的・対象域の制限）。
強さ：弱（標準形式の適用対象をACS同士に限定する）。
依拠箇所：\S\ref{subsec:rel_bidirect} および \S\ref{subsec:rel_categories}。
\end{assumption}

\begin{definition}[双方向合成関係（分類C）]
\label{def:bidirect_relation}
$Y_j\in\mathcal Y_{\mathrm{ACS},-i}$ をACS個体 $\Sigma_j$ として扱い、
逆方向でも $\Sigma_i\in\mathcal Y_{\mathrm{ACS},-j}$ とする。
各個体の局所サイズ指標の評価域と通常ACS状態量の定義域をそれぞれ確認したうえで、
共通の正の窓長 $L$ と時刻 $t$ において
$\mathrm{Rel}_{i\leftarrow j,t}$ と $\mathrm{Rel}_{j\leftarrow i,t}$ の
両方が整形式である場合、
双方向合成関係を
\[
 \mathrm{Rel}_{ij,t}:=
 \bigl(\mathrm{Rel}_{i\leftarrow j,t},\mathrm{Rel}_{j\leftarrow i,t}\bigr)
\]
と定める。
両方向とも $W_{L,t}$ の同一の時刻範囲を用いるが、
同一対象の追跡、
関係継続、
および各個体の通常ACS状態量の定義域は
それぞれ別に満たされなければならない。
\S\ref{sec:collapse_cond} の軌道記録を双方について採用する場合、
$\Sigma_i$ の記録を $\sigma_i$、
$\Sigma_j$ の記録を $\sigma_j$ と呼べば、
\[
 W_{L,t}\subseteq I_{\mathrm{ACS}}(\sigma_i),\qquad
 W_{L,t}\subseteq I_{\mathrm{ACS}}(\sigma_j)
\]
を各々に要求する。
これは一つの軌道が両個体の通常状態を兼ねるという主張ではない。

窓全体が両者の有効時刻域に含まれていても、
三差分と前三成分の取得条件は別途必要である。
本定義は外部対象の内部動学を新たに構成せず、
二つのACS側の記述を組にするだけである。
順序対の第一成分と第二成分はそれぞれ異なる個体側の観測変化を含むため、
双方向合成の定義は、両者の可行性関連量の増加、
相互利益、
または
組織水準の統合構造の成立を要件として含まない。

$ij$ は対象の組を区別する標識であり、
成分の積や数値演算を表すものではない。
\end{definition}

\begin{remark}[双方向合成関係の非対称性]
\label{rem:bidirect_asymmetry}
$\mathrm{Rel}_{ij,t}$ は二つの単方向記述子の順序対であり、
$\mathrm{Rel}_{i\leftarrow j,t}=\mathrm{Rel}_{j\leftarrow i,t}$ を仮定しない。
カテゴリーが同じでもコスト・摂動・三差分は別個体側の値として異なりうる。
両方向が記録されることは、
双方に同じ作用や相互利益が存在することを意味しない。

各方向の三差分は各個体に帰属する時系列から計算されるため、
二つの数値差を無条件に加算したり、
双方に共通の可行性指標へ読み替えたりしない。
一方で記述子が定義できても、
逆方向に必要な記録が欠ければ標準的な双方向合成は置かない。
例えば同じ関係カテゴリーを両方向で記録した場合でも、
一方の資源負荷が他方より大きいことや、
実効的時間地平の差分が
逆符号になることは形式上排除されない。
双方向性を対称性の仮定へ読み替えない。
\end{remark}

\begin{definition}[方向性指標（分類C）]
\label{def:relation_direction}
双方向合成関係が定義されるとき、
解釈上の主要な記述方向を示す補助指標を
\[
 \mathrm{dir}_{ij,t}\in
 \{i\leftarrow j,\ i\rightarrow j,\ i\leftrightarrow j,\ \mathrm{undetermined}\}
\]
と書く。
順に、
$j$ から $i$ の記述が主要、
逆方向が主要、
両方向が主要、
主要方向を本節の規則では判定しない、
を表す。
「主要」の判定基準・閾値・重み付けを定めないため、
本指標が二つの単方向記述子から一意に計算されるとは主張しない。
判定基準を採用しない場合には $\mathrm{undetermined}$ と記録してよい。
これは二方向とも値が零であるという数値判定でも、
当該二者に関係がないという分類でもない。
方向性指標の値は六成分記述子そのものには含めない。
二方向のうち何を「主要」とするかは、
記録対象や採用する比較基準にも依存しうる。

本節で値域を示したことは、
主要方向の一意な算出規則を与えたことではない。
\end{definition}

\begin{remark}[方向性指標の補助的地位]
\label{rem:dir_auxiliary}
方向性指標は記述補助であり、
二つの単方向記述子を置き換えず、
閉塞・崩壊や可行性の改善方向を規定しない。
\end{remark}

\begin{remark}[方向性判定の未定義部分]
\label{rem:dir_undefined}
コスト・摂動・三差分のどれをどの程度重視するかという
主要方向の定量的判定規則はPhase~1.$\alpha$に委ねる（\S\ref{subsec:phase1a}）。
$\mathrm{undetermined}$ は関係の不存在、
両方向の対称性、
数値ゼロのいずれも意味しない。
また、
$\mathrm{dir}_{ij,t}$ は対象間の因果の向きを特定する量ではない。
\end{remark}

\begin{definition}[同種相互関係と異種相互関係]
\label{def:homo_hetero_relation}
$\mathrm{Rel}_{ij,t}$ の両単方向記述子に含まれる
$\mathrm{cat}_{i\leftarrow j,t}, \mathrm{cat}_{j\leftarrow i,t}\in\mathcal C_{\rightarrow}$ について、
\[
 \mathrm{cat}_{i\leftarrow j,t}=\mathrm{cat}_{j\leftarrow i,t}
\]
のときを同種相互関係、
\[
 \mathrm{cat}_{i\leftarrow j,t}\neq \mathrm{cat}_{j\leftarrow i,t}
\]
のときを異種相互関係と呼ぶ。

この区別は、
各方向のカテゴリーが既に記録されている場合の
ラベル一致・不一致のみを調べるものであり、
双方向成立カテゴリーの判定と同一の操作ではない。
\end{definition}

\begin{remark}[同種・異種の分類の制限]
\label{rem:homo_hetero_rel}
同種性はラベルの一致だけを意味し、
コスト・強度・持続時間・三差分の一致を意味しない。
応答側カテゴリー（被支配・被管理等）の体系化はPhase~1.$\alpha$に委ねる。
同種の場合であっても双方のコスト・摂動・三差分は異なりうる。
異種の場合も双方の差分が同じ符号を示す可能性を排除しない。
「被操作」「被搾取」「被利用」などを追加ラベルとして用いるには、
応答側の成立条件を別途定める必要がある。
\end{remark}

\begin{remark}[同種相互形を持たないカテゴリー]
\label{rem:no_mutual_form}
支配や寄生が両方向で記録されても、
それを「相互支配」や「相互寄生」の
独立した成立カテゴリーとして直ちに定義しない。
必要な複合カテゴリーと成立規則はPhase~1.$\alpha$の対象とする。
例えば競合的制御、
相互制約、
相互侵食的依存といった記述への整理は
考えられるが、
本論文はそれらを第19以降の正式カテゴリーとして定義しない。
同じ単方向ラベルが二つ現れたという形式と、
自然言語上の
「相互」という名称に固有の成立条件とは区別する。
\end{remark}

\begin{definition}[双方向成立カテゴリー（分類C）]
\label{def:bidirectional_category}
双方向合成関係 $\mathrm{Rel}_{ij,t}$ が定義され、
その双方向的な成立様式についてラベルを付す場合の値域を
\[
 \mathrm{cat}_{\leftrightarrow,ij,t}\in\mathcal C_{\leftrightarrow},\qquad
 \mathcal C_{\leftrightarrow}=
 \{\mathrm{coop},\mathrm{sym},\mathrm{conf}\}
\]
と書く。
この集合の正式な定義は定義~\ref{def:category_set} で与える。

これらは協力・共生・対立の双方向成立ラベルであり、
単方向記述子のカテゴリー成分 $\mathrm{cat}_{i\leftarrow j,t}\in\mathcal C_{\rightarrow}$ とは区別する。

本節は双方向ラベルの成立条件や、
二つの単方向ラベルからの導出規則を与えない。
そのため、
この補助ラベルを採用するか否かは、
双方向合成の整形式性とは別の記述上の選択である。
同種・異種の分類が定まっても、
協力・共生・対立のいずれかを自動的に決定できるわけではない。
したがって、
$\mathrm{Rel}_{ij,t}$ が定義されたすべての組に
$\mathrm{cat}_{\leftrightarrow,ij,t}$ を必ず割り当てるものではない。
未判定を新たな第19の関係カテゴリーとして追加しない。
\end{definition}

\begin{remark}[協力・共生・対立の双方向性]
\label{rem:cooperation_bidirectional}
協力・共生・対立の三ラベルは単方向カテゴリー集合に含めず、
二方向の関係記録を照合する水準に置く。
双方向合成関係の存在だけでは、
三ラベルのいずれが成立するかは定まらない。

寄与・相互性・対立性の成立条件と閾値を与えることはPhase~1.$\alpha$の課題とする。
例えば片方向に支援が記録されたという事実は、
逆方向にも支援があることを示さない。

両方向に支援が記録された場合にも、
それだけで協力の必要十分条件を得たことにはならない。
相互にどのような作用が成立したとき協力・共生・対立と呼ぶかは、
双方向記述を照合する規則が具体化されて初めて判断できる。
\end{remark}

\begin{remark}[競争の方向性に関する特殊性]
\label{rem:competition_directionality}
競争 $\mathrm{comp}$ は単方向カテゴリーであり、
一方の競争的作用だけを記録できる。
両方向で $\mathrm{cat}_{i\leftarrow j,t}=\mathrm{cat}_{j\leftarrow i,t}=\mathrm{comp}$ なら
同種相互関係として記録できるが、
それを別の双方向成立カテゴリーとはしない。
一方の対象が競争的な作用を記録し、
相手側に同種の作用が記録されない場合も、
この単方向ラベルを双方向成立ラベルへ移し替える必要はない。
競争が双方で現れた場合の強度や成立閾値の定量的区別はPhase~1.$\alpha$へ委ねる。
\end{remark}

\begin{remark}[双方向成立カテゴリーと同種相互関係の関係]
\label{rem:bidirect_vs_homo}
二つの単方向ラベルの一致・不一致を述べる同種／異種の区別
（定義~\ref{def:homo_hetero_relation}）と、
双方向成立ラベル（定義~\ref{def:bidirectional_category}）は別の記述軸である。
同種の支援ラベルが二つあるだけで協力の成立を導く規則は本節にない。
一方、
双方向成立ラベルを付す場合も、
両方向の単方向記録が不要になるわけではない。
本節は $\mathcal C_{\rightarrow}\times\mathcal C_{\rightarrow}$ から
$\mathcal C_{\leftrightarrow}$ への全域写像を定義しない。
二つの単方向ラベルが同種か異種かという判定は、
双方向成立ラベルがどれかを決定する十分条件にはならない。
また、
ある双方向成立ラベルを記録したとしても、
そのことだけから両方向の単方向カテゴリーの組が一意に復元できるわけではない。
\end{remark}

\begin{remark}[非ACS対象との双方向関係の制限]
\label{rem:nACS_limit}
非ACSエージェントや非エージェント的対象には、
本節で $\Sigma_j$ 側の状態空間、
可行行動集合または実効的自由度を定義しない。
そのため $\mathrm{Rel}_{j\leftarrow i,t}$ をACS型の六成分形式で自動生成しない。

$\Sigma_i$ 側の単方向記述は引き続き許される。
非ACS側の内部構造を明示する双方向形式は本論文の対象外とする。
非ACSエージェントに何らかの行為主体性があるとしても、
本論文はその対象へACSの七タプルや三機構を割り当てない。
物体との作用についても個体側の記録は行えるが、
物体側の可行行動集合を仮定して逆方向の六成分を作らない。
\end{remark}

\begin{remark}[双方向合成関係と閉塞・崩壊の節]
\label{rem:bidirect_to_sec7}
双方向合成に含まれる二つの単方向記述のいずれかは
\S\ref{sec:collapse_cond} の関係記録に参照されうる。
ただし本論文は、双方向合成の存在やカテゴリの組から、
いずれかのACSの失敗状態を一般的な帰結として置かない。
二つの単方向記述子に併記される差分を検討する際も、
閉塞と崩壊の成立は各個体の軌道ごとに別途判定する。
一方の失敗状態から、
他方の失敗状態を一般に導く規則は本節にはない。
\end{remark}

\begin{remark}[組織ACSの節への接続]
\label{rem:bidirect_to_sec8}
$\mathrm{Rel}_{ij,t}$ は、
\S\ref{sec:org_acs} で扱う
組織水準の結合構造を記述する素材になりうる。
ただし本論文は、
二つの関係記述子の組から組織自身の自己--環境境界、
機能の分担・調整、
または組織水準のACS構造が自動的に成立することを一般的な帰結として置かない。
本節は組織ACSの成立判定やネットワークの時間発展を定義しない。
双方向合成を多数並べることと、
組織自身について組織水準のACS構造を
成立させることは異なる。
組織ACSを単なるACS集団から区別する条件は後続節が定める。
関係の組を列挙するだけで、
新しい組織水準の状態空間や組織水準の可行性の
存在を示したことにはならない。
\end{remark}

%----------------------------------------------------------------------
\subsection{関係品質}
\label{subsec:ro_quality}

本小節では、
\S\ref{sec:three_principles} の構造的環境結合条件（\S\ref{subsec:env_coupling}）
および EU の時刻別充足述語（\S\ref{subsec:eu}）とは区別して、
ACS $\Sigma_i$ と外部対象 $Y_j\in\mathcal Y_{-i}$ との評価可能な関係が
$\Sigma_i$ の可行性維持に対してどの向きにあるかを記述する
関係品質 $Q_{i\leftarrow j,t}$ を形式化する。

\S\ref{subsec:env_coupling} で整理したとおり、
本論文は
\[
  \text{構造的環境結合 } \mathrm{EC}(x_t),
  \qquad
  \text{実現外部相互作用 } \chi_{\mathrm{EI},t},
  \qquad
  \text{EU の時刻別充足述語 } \mathrm{EU}_t,
  \qquad
  \text{関係品質評価 } Q
\]
を相互に異なる対象として扱う。
したがって、本小節で導入する $Q$ は、
構造的環境結合条件の条件でも、
EU の時刻別充足述語の追加条件でもない。
$Q$ は三原理の時刻別充足述語の定義に用いない。


本小節の $Q_{i\leftarrow j,t}$ は、記述者が指定する評価規則のもとで定める分類Cの量である。
理論上の $Q_{i\leftarrow j,t}$ の値または後述する各成分を
ACS が同一の数値表現で保持・計算することは要求しない。

\subsubsection{関係品質の評価可能域}

外部対象 $Y_j\in\mathcal{Y}_{-i}$（定義~\ref{def:external_object_set}）は、
他の ACS または行為主体に限定されない。
非主体的対象、資源、物質、エネルギー源、
その他の外部対象も含みうる。

\begin{definition}[関係品質の評価規則（分類C）]
\label{def:Q_rule}
記述者は、評価の対象とする関係 $(i,j)$ ごとに
\emph{評価規則} $\mathrm{Rule}_{Q,i\leftarrow j}$ を指定する。
評価規則は、参照する作動記録の区間またはモデル、比較の基準、評価の時間幅、
および定義~\ref{def:Q} の三成分と重みの基準を定め、
比較を行う期間を通じて固定する。
評価規則は規範的な善悪の基準を表さない。
これは、偏向性の参照基準（定義~\ref{def:bias_tendency}）と同じ扱いである。
評価の時間幅が時刻 $t$ より先の記録を含む場合、
時刻 $t$ について与える値は事後的な記述であり、
時刻 $t$ における当該 ACS の入力ではない
（\S\ref{sec:extended_dynamics} の予測精度係数と同じ扱い）。
\end{definition}

\begin{definition}[関係品質評価可能域]
\label{def:Q_domain}
時刻 $t$ において、
評価規則 $\mathrm{Rule}_{Q,i\leftarrow j}$ が値を与える外部対象の集合を
\[
  \mathcal D_{Q,i,t}
  :=
  \Bigl\{
    Y_j\in\mathcal{Y}_{-i}
    \ \Bigm|\
    \mathrm{Rule}_{Q,i\leftarrow j}\text{ が時刻 $t$ に値を与える}
  \Bigr\}
\]
と定義し、関係品質の\emph{評価可能域}と呼ぶ。
$\mathcal D_{Q,i,t}$ は理論上の定義域であり、
ACS がその内部で当該関係の品質を評価できる対象の集合ではない。
$Q$ が理論上定義されることと、ACS が内部で $Q$ に相当する評価を形成・保持することは区別する
（\S\ref{subsec:description_levels_global}）。

添字 $i\leftarrow j$ の向きは記述の視点を表し（備考~\ref{rem:rel_unidirect_notation}）、
$Y_j$ が行為主体であることや、時刻 $t$ に因果作用が実現していることを意味しない。

Phase~1 は
\[
  \mathcal D_{Q,i,t}\neq\emptyset
\]
を一般には要求しない。
\end{definition}

\begin{remark}[評価不能と中立評価の区別]
\label{rem:Q_undefined}
定義~\ref{def:Q_domain} により、
$Y_j\notin\mathcal D_{Q,i,t}$
である場合、
本小節の $Q_{i\leftarrow j,t}$ は
\emph{定義されない}。

したがって本論文は、
評価不能な関係に対して規約的に
$Q=0$
を割り当てない。

$Q_{i\leftarrow j,t}=0$
という記述は、
$Y_j\in\mathcal D_{Q,i,t}$ であり、
品質評価が実際に定義されたうえで
その総合評価が中立境界にある場合に限って用いる。

この区別により、
\[
  \text{評価不能}
  \neq
  \text{中立品質}
\]
を明示的に維持する。
\end{remark}

\subsubsection{関係品質評価 $Q$ の形式化}

\begin{definition}[関係品質評価 $Q$]
\label{def:Q}
時刻 $t$ において
\[
  Y_j\in\mathcal D_{Q,i,t}
\]
とする。

評価規則 $\mathrm{Rule}_{Q,i\leftarrow j}$ のもとで、
$Y_j$ との関係が
ACS $\Sigma_i$ の可行性維持に関係する
情報形成、資源収支、および可行性方向に対してもつ
総合的な関係品質を
\[
  Q_{i\leftarrow j,t}
  \in[-1,1]
\]
とする。

Phase~1 では、
この総合評価を
\[
  Q_{i\leftarrow j,t}
  :=
  \alpha^{\mathrm{inf}}_{i\leftarrow j} Q^{\mathrm{inf}}_{i\leftarrow j,t}
  +
  \alpha^{\mathrm{res}}_{i\leftarrow j} Q^{\mathrm{res}}_{i\leftarrow j,t}
  +
  \alpha^{\mathrm{via}}_{i\leftarrow j} Q^{\mathrm{via}}_{i\leftarrow j,t}
\]
として表す。

各成分
\[
  Q^{\mathrm{inf}}_{i\leftarrow j,t},
  \qquad
  Q^{\mathrm{res}}_{i\leftarrow j,t},
  \qquad
  Q^{\mathrm{via}}_{i\leftarrow j,t}
  \in[-1,1]
\]
は、次の関係評価軸を表す。

\begin{itemize}
  \item
  \textbf{情報品質成分}
  \[
    Q^{\mathrm{inf}}_{i\leftarrow j,t}\in[-1,1].
  \]
  $Y_j$ に関係する情報的作用または情報内容が、
  ACS $\Sigma_i$ の内部モデル、評価、または応答形成における
  可行性適合的な情報形成に対してもつ
  正負の寄与方向を正規化して表す。

  \[
    Q^{\mathrm{inf}}_{i\leftarrow j,t}>0
  \]
  は可行性維持に関する情報形成を支える側、
  \[
    Q^{\mathrm{inf}}_{i\leftarrow j,t}<0
  \]
  はこれを阻害する側の評価を表す。

  \item
  \textbf{資源・物質・エネルギー収支成分}
  \[
    Q^{\mathrm{res}}_{i\leftarrow j,t}\in[-1,1].
  \]
  $Y_j$ との関係に伴う
  ACS $\Sigma_i$ 側の資源、物質、エネルギー等の
  可行性維持上の正味寄与方向を正規化して表す。

  \[
    Q^{\mathrm{res}}_{i\leftarrow j,t}>0
  \]
  は補給、獲得、負荷軽減等の正側、
  \[
    Q^{\mathrm{res}}_{i\leftarrow j,t}<0
  \]
  は損失、枯渇、負荷増大等の負側の評価を表す。

  \item
  \textbf{可行性方向評価成分}
  \[
    Q^{\mathrm{via}}_{i\leftarrow j,t}\in[-1,1].
  \]
  $Y_j$ との関係が、
  ACS $\Sigma_i$ の可行性余裕度または
  可行性維持に関係する状態変化方向に対してもつ
  正負の寄与を評価した形式量である。

  \[
    Q^{\mathrm{via}}_{i\leftarrow j,t}>0
  \]
  は可行性を支える方向、
  \[
    Q^{\mathrm{via}}_{i\leftarrow j,t}<0
  \]
  は可行性を損なう方向の評価を表す。
\end{itemize}

重み
\[
  \alpha^{\mathrm{inf}}_{i\leftarrow j},
  \qquad
  \alpha^{\mathrm{res}}_{i\leftarrow j},
  \qquad
  \alpha^{\mathrm{via}}_{i\leftarrow j}
  \ge0
\]
は
\[
  \alpha^{\mathrm{inf}}_{i\leftarrow j}
  +
  \alpha^{\mathrm{res}}_{i\leftarrow j}
  +
  \alpha^{\mathrm{via}}_{i\leftarrow j}
  =1
\]
を満たすものとする。

各成分の値と重みは評価規則 $\mathrm{Rule}_{Q,i\leftarrow j}$ が定める。
重みは、評価対象 $Y_j$、その対象種別、
および当該 ACS の構造に相対する形式パラメータであり、
すべての ACS に共通する普遍定数とはしない。

特定の評価対象について
三成分のうち一部が
当該関係の評価軸として意味をもたない場合には、
その成分に対応する重みを $0$ とすることを許す。
ただし、重みの総和は $1$ とする。

各成分が $[-1,1]$ に属し、
重みが非負で総和 $1$ であるから、
$Q$ は三成分の凸結合となり、
\[
  Q_{i\leftarrow j,t}\in[-1,1]
\]
が成立する。
\end{definition}

\begin{remark}[関係品質と内部の評価の区別]
\label{rem:Q_internal_formal}
$Q_{i\leftarrow j,t}$ は記述者が評価規則のもとで定める量であり、
当該 ACS の内部の評価ではない。
\S\ref{sec:three_principles} の環境利用・外部結合評価機能（定義~\ref{def:principle_related_functions}）は、
外部関係に伴う可行性関連差異の少なくとも一部を識別・評価する分類Aの機能であり、
\S\ref{sec:tipa} の要求×対象の評価（定義~\ref{def:rq_eval_types}）は、
内部対象トークンについての内部の評価である。
本論文は、これらの内部の評価が $Q_{i\leftarrow j,t}$ と一致することを要求しない
（備考~\ref{rem:theory_quantity_vs_internal}）。
したがって、両者の食い違い（たとえば内部では好ましいと評価される関係について
$Q_{i\leftarrow j,t}<0$ であること）は、それだけでは機能の不全を意味しない。
\end{remark}

\begin{remark}[関係品質と対象の性質の区別]
\label{rem:Q_not_property}
$Q_{i\leftarrow j,t}$ は $Y_j$ の性質だけからは定まらず、
当該 ACS、時刻および評価規則に相対する。
同じ $Y_j$ について、異なる ACS の間で、また同じ ACS の異なる時刻の間で、符号は異なりうる。
\end{remark}

\begin{remark}[$Q$ の成分形式に関する Phase~1 の地位]
\label{rem:Q_form}
三成分
\[
  Q^{\mathrm{inf}}_{i\leftarrow j,t},
  \qquad
  Q^{\mathrm{res}}_{i\leftarrow j,t},
  \qquad
  Q^{\mathrm{via}}_{i\leftarrow j,t}
\]
は、
関係品質評価を
情報、資源収支、可行性方向という
三つの機能的観点に分解する
Phase~1 の形式表現である。

これは、
あらゆる生物的・人工的 ACS が
必ず三つの独立した内部変数を保持すること、
または関係品質評価の実装が
この三成分分解に一意であることを意味しない。

各成分の操作的測定方法、
対象種別ごとの正規化規則、
各重みの経験的同定、
およびより高次元の関係品質表現は、
Phase~1 では規定しない
（within present framework）。
\end{remark}

\subsubsection{$Q$ の符号分類}

\begin{definition}[$Q$ の符号による関係品質分類]
\label{def:Q_sign}
時刻 $t$ において
\[
  Y_j\in\mathcal D_{Q,i,t}
\]
である評価可能な関係について、
$Q$ の符号により次のように記述する。

\begin{itemize}
  \item
  \[
    Q_{i\leftarrow j,t}>0
  \]
  のとき、
  \emph{正側関係品質}と呼ぶ。

  \item
  \[
    Q_{i\leftarrow j,t}=0
  \]
  のとき、
  \emph{中立関係品質}と呼ぶ。

  \item
  \[
    Q_{i\leftarrow j,t}<0
  \]
  のとき、
  \emph{負側関係品質}と呼ぶ。
\end{itemize}

この分類は、
当該 ACS の可行性維持に関係する
構造的評価方向を表す記述的分類であり、
外部対象または関係の
倫理的・規範的な善悪を判定するものではない。
\end{definition}

\begin{remark}[$Q=0$ は定義された評価の中立境界である]
\label{rem:Q_threshold}
定義~\ref{def:Q_sign} における
$Q=0$
は、
別個の外生的な普遍閾値
$Q_{\mathrm{threshold}}$
を導入して設定したものではなく、
定義~\ref{def:Q} の符号付き成分表現における
中立境界である。

したがって、
$Q>0$、$Q=0$、$Q<0$
は、それぞれ
本品質評価上の正側、中立、負側を表す。

ただし、
この $Q=0$ は
$Y_j\in\mathcal D_{Q,i,t}$
である場合にのみ意味をもつ。

評価不能
$Y_j\notin\mathcal D_{Q,i,t}$
を
$Q=0$ と同一視してはならない。
\end{remark}

\subsubsection{$Q$ と構造的環境結合・実現外部相互作用の区別}

\begin{proposition}[$Q$ と構造的環境結合条件の定義式上の相互非参照]
\label{prop:Q_not_RO_condition}
\begin{enumerate}[label=(\roman*)]
  \item 構造的環境結合条件 $\mathrm{EC}(x_t)$（定義~\ref{def:env_coupling}）が
  成り立つか否かは、$x_t$、$f$、$U$、$E_{\mathrm{ext}}$ および $\pi_E$ だけで定まり、
  評価可能域 $\mathcal D_{Q,i,t}$ と $Q$ の値をその定義に含まない。
  \item 関係品質評価 $Q_{i\leftarrow j,t}$（定義~\ref{def:Q}）は、
  $Y_j\in\mathcal D_{Q,i,t}$ である評価可能な関係について定まり、
  $\mathrm{EC}(x_t)$ をその定義に含まない。
\end{enumerate}
\end{proposition}

\begin{proof}
(i)：定義~\ref{def:env_coupling} の条件は、
$u\in U$ と $e,e'\in E_{\mathrm{ext}}$ についての量化と
$\pi_E(f(x_t,u,e))\neq\pi_E(f(x_t,u,e'))$ だけからなる。
(ii)：定義~\ref{def:Q_rule}、定義~\ref{def:Q_domain} および定義~\ref{def:Q} の条件に
定義~\ref{def:env_coupling} は現れない。
\end{proof}

本命題は、両者の定義が互いを含まないことを述べるものであり、
両者の間に含意が成り立たないことを示したものではない。
本論文は、評価可能な関係について
$Q_{i\leftarrow j,t}>0$ または $Q_{i\leftarrow j,t}<0$ であることから
$\mathrm{EC}(x_t)$ が成り立つか否かを、
また $\mathrm{EC}(x_t)$ から特定の $Y_j$ について
$Y_j\in\mathcal D_{Q,i,t}$ または $Q_{i\leftarrow j,t}>0$ を、
いずれも一般的な帰結として置かない。
同じく、$Q$ の符号から $\mathrm{EU}_t$ の真偽を一般的な帰結として置かない。

\begin{remark}[$Q$ と実現外部相互作用の非同一性]
\label{rem:Q_chi_distinction}
\S\ref{subsec:env_coupling} の
実現外部相互作用指標
$\chi_{\mathrm{EI},t}$
と $Q$ も同一ではない。

$\chi_{\mathrm{EI},t}=1$
は、
時刻 $t$ に境界横断的な外部相互作用が
実現していることを表すが、
本論文はそれだけから
特定の $Y_j$ との関係について
$Y_j\in\mathcal D_{Q,i,t}$
または
$Q_{i\leftarrow j,t}>0$
を一般的な帰結として置かない。

逆に、
関係品質評価が定義されていることは、
過去の相互作用、内部モデル、記憶、
その他の関係情報に基づく場合を
Phase~1 は排除しないため、
$Y_j\in\mathcal D_{Q,i,t}$
のみから
$\chi_{\mathrm{EI},t}=1$
を一般的な帰結として置かない。

したがって、
\[
  \mathrm{EC}(x_t),
  \qquad
  \chi_{\mathrm{EI},t},
  \qquad
  \mathrm{EU}_t,
  \qquad
  Q
\]
は相互に区別して扱う。
\end{remark}

\subsubsection{負の関係品質を伴う環境結合}

\begin{definition}[負の関係品質を伴う環境結合（分類C）]
\label{def:contaminated_ro}
時刻 $t$ において、
構造的環境結合条件 $\mathrm{EC}(x_t)$（定義~\ref{def:env_coupling}）が成り立ち、
かつ、ある評価可能な外部対象
\[
  Y_j\in\mathcal D_{Q,i,t}
\]
について
\[
  Q_{i\leftarrow j,t}<0
\]
が成立するとき、
この配置を当該関係 $Y_j$ に関する
\emph{負の関係品質を伴う環境結合}と呼ぶ。
すなわち、
\[
  \mathrm{EC}(x_t)
  \quad\wedge\quad
  \exists\,Y_j\in\mathcal D_{Q,i,t}
  :
  Q_{i\leftarrow j,t}<0
\]
という記述的配置であり、
外部対象または関係の善悪を判定するものではない。
\end{definition}

\begin{remark}[負の関係品質を伴う環境結合の適用範囲]
\label{rem:contaminated_ro_scope}
定義~\ref{def:contaminated_ro} は、
構造的環境結合条件が成り立っているにもかかわらず、
少なくとも一つの評価可能な関係が
負側の品質評価を受けるという
構造的配置を区別するための記述概念である。

これは、
構造的環境結合自体が「悪い」こと、
外部との関係を閉じるべきこと、
または当該外部対象が
倫理的に有害であることを意味しない。

また本論文は、
$Q_{i\leftarrow j,t}<0$
単独から
$\Pi(x_{t+1})<\Pi(x_t)$、
$A_t=\emptyset$、
または
$\mathrm{Col}_{i,t}=1$
を一般的な帰結として置かない。

$Q$ は関係品質の評価方向を表すが、
その評価と実際の状態遷移との定量的接続には、
$f$、実現行動、実現外部環境状態、
内部状態、および必要に応じた追加条件が関与する。
\end{remark}

\begin{remark}[負側関係品質と負の関係品質を伴う環境結合は同一ではない]
\label{rem:contaminated_status}
$Q_{i\leftarrow j,t}<0$
という負側関係品質は、
評価可能な一つの関係についての記述である。
これに対し、
負の関係品質を伴う環境結合は
$\mathrm{EC}(x_t)$ が成り立つことを追加的に要求する。
したがって、
$Q<0$
と
\[
  \text{負の関係品質を伴う環境結合}
\]
は同一概念ではない。

また、負の関係品質を伴う環境結合は、
三原理に追加される第四の原理、
構造的環境結合条件の別種の形、
EU の時刻別充足述語の別種の条件、
または新たな失敗状態ではない。
\end{remark}


\subsubsection{$Q$ と可行性変化の関係}

\begin{remark}[$Q$ の符号と $\Delta\Pi$ の接続を一般則として置かない]
\label{rem:Q_connection}
両時刻の状態 $x_{t-1}$ と $x_t$ が定義される $t\ge1$ に対して、時刻差分
\[
  \Delta\Pi_t
  :=
  \Pi(x_t)-\Pi(x_{t-1})
\]
を考える場合でも、本論文は
\[
  Q_{i\leftarrow j,t}<0
  \quad\Longrightarrow\quad
  \Delta\Pi_t<0
\]
および
\[
  Q_{i\leftarrow j,t}>0
  \quad\Longrightarrow\quad
  \Delta\Pi_t>0
\]
を一般的な帰結として置かない。
これらの含意が成り立たないことを示したものではない。

理由は、
$Q$ が特定の評価対象関係についての
総合的な関係品質評価であるのに対し、
$\Delta\Pi_t$ は
ACS の実現状態遷移全体の結果として定まるからである。

したがって、
$Q$ の符号と
$\Pi$ の時間変化方向との間に
一般的な必要十分関係を置かない。

特定の ACS または特定の関係型について
両者を定量的に接続するためには、
追加の構造仮定または動学モデルを要する。
\end{remark}

\subsubsection{時間集約との役割分担}

\begin{remark}[$Q$ の時間集約は後続小節で扱う]
\label{rem:Q_temporal_deferred}
本小節の $Q$ は、
時刻 $t$ における
評価可能な個別関係の品質評価を定義する。

有限窓平均は本小節の後半で、
三原理関連評価の時間統合（T-STAB）は
\S\ref{subsec:s3_tstab} で扱う。

したがって本小節では、
評価可能な $Q$ が存在しない場合に
時間平均へ規約値 $0$ を代入すること、
単一の評価対象関係を
普遍的に「代表関係」として固定すること、
または $Q$ を EU の時刻別充足述語や構造的環境結合条件の程度へ
直接変換することは行わない。
\end{remark}



\subsubsection{関係品質 $Q$ の有限窓集約}

\begin{definition}[関係品質の有限窓評価]
\label{def:Q_window_average}
任意の有限窓幅
\[
  W_Q\in\mathbb{N},
  \qquad
  W_Q\ge1
\]
を固定する。

時刻 $t\in\mathbb{N}_0$ における
$Q$ 用有限窓を
\[
  \mathcal{W}_t^{Q}
  :=
  \Bigl\{
    r\in\mathbb{N}_0
    \Bigm|
    \max(0,t-W_Q+1)\le r\le t
  \Bigr\}
\]
とする。

外部対象
$Y_j\in\mathcal{Y}_{-i}$
について、
この窓内で
$Q_{i\leftarrow j,r}$ が定義される
時刻集合を
\[
  \mathbb T^{Q}_{i\leftarrow j,t}
  :=
  \Bigl\{
    r\in\mathcal{W}_t^{Q}
    \Bigm|
    Y_j\in\mathcal D_{Q,i,r}
  \Bigr\}
\]
とする。

\[
  \mathbb T^{Q}_{i\leftarrow j,t}\neq\emptyset
\]
の場合に限り、
当該関係の有限窓平均を
\[
  \bar Q_{i\leftarrow j,W_Q,t}
  :=
  \frac{1}{|\mathbb T^{Q}_{i\leftarrow j,t}|}
  \sum_{r\in\mathbb T^{Q}_{i\leftarrow j,t}}
  Q_{i\leftarrow j,r}
  \in[-1,1]
\]
と定義する。

一方、
\[
  \mathbb T^{Q}_{i\leftarrow j,t}=\emptyset
\]
の場合、
$\bar Q_{i\leftarrow j,W_Q,t}$ は
\emph{未定義}とする。
\end{definition}

\begin{remark}[$Q$ 未評価窓に規約値 $0$ を与えない]
\label{rem:Q_window_undefined}
定義~\ref{def:Q_window_average} は、
本小節の
\[
  \text{評価不能}
  \neq
  \text{中立品質}
\]
という区別を時間集約後にも維持する。

したがって、
有限窓内に
評価可能な $Q$ が一つも存在しない場合に
\[
  \bar Q_{i\leftarrow j,W_Q,t}:=0
\]
とは置かない。

また、
Phase~1 では
複数の外部対象に対する
$\bar Q_{i\leftarrow j,W_Q,t}$ を
単一の総合 $\bar Q_{i,t}$ へ集約する
普遍的規則を置かない。

どの関係を評価対象として優先するか、
複数関係をどの作用素で統合するか、
および対象依存の重みを
どのように更新するかは、
実装および後続の動学に依存しうる
（within present framework）。
\end{remark}

\begin{remark}[$Q$ の時間集約は EU・構造的環境結合の再定義ではない]
\label{rem:Q_window_not_RO}
\[
  \bar Q_{i\leftarrow j,W_Q,t}>0,
  \qquad
  \bar Q_{i\leftarrow j,W_Q,t}=0,
  \qquad
  \bar Q_{i\leftarrow j,W_Q,t}<0
\]
は、評価可能な一つの関係についての時間集約された品質記述である。
本論文は、これらの符号から、
窓内の時刻における $\mathrm{EU}_r$ の真偽、
または構造的環境結合条件 $\mathrm{EC}(x_r)$（定義~\ref{def:env_coupling}）が
成り立つか否かを、一般的な帰結として置かない。
したがって、$Q$ の有限窓評価は EU の条件でも構造的環境結合条件でもなく、
関係品質の時間方向への集約として扱う。
\end{remark}

\subsubsection{T-STAB との関係}

\begin{remark}[$Q$ を T-STAB の普遍条件へ組み込まない]
\label{rem:Q_not_tstab_core}
本論文は、定義~\ref{def:tstab} の T-STAB に
\[
  \bar Q_{i\leftarrow j,W_Q,t}\ge Q^{\min}
\]
のような関係品質下限を普遍的な必須条件として組み込まない。

理由は、本小節で定めた $Q$ が
特定の評価可能な外部関係に相対する量であり、
時刻または対象によって評価可能域 $\mathcal D_{Q,i,t}$ が
空でありうるためである。
また、複数の評価対象関係を単一の総合品質へ集約する普遍的規則も
Phase~1 では定めていない。
したがって、$Q$ の時間集約は
T-STAB と並行する関係品質の時間評価チャネルとして保持する。

特定の ACS 実装において、
ある評価対象関係または複数関係の集約量を
T-STAB 型の内部参照条件へ組み込むことは排除しないが、
その場合の対象選択、集約作用素、閾値、および欠測処理は
明示的に追加定義されなければならない。
\end{remark}

\begin{remark}[負の関係品質を伴う環境結合の時間評価]
\label{rem:contaminated_ro_temporal}
負の関係品質を伴う環境結合は、
時刻ごとの配置として
\[
  \mathrm{EC}(x_t)
  \quad\wedge\quad
  \exists Y_j\in\mathcal D_{Q,i,t}:
  Q_{i\leftarrow j,t}<0
\]
で記述される（定義~\ref{def:contaminated_ro}）。

定義~\ref{def:Q_window_average} により、
特定の $Y_j$ について $\bar Q_{i\leftarrow j,W_Q,t}$ が定義される場合には、
その符号を記述に用いることができる。
ただし窓平均が負であることは、窓内の各時刻で $Q_{i\leftarrow j,r}$ が定義され負であることを含意しない。
負側関係品質の時間的持続を述べるには、窓内の各時刻での定義と符号を別に確かめる。

ただし本論文は、
$\bar Q_{i\leftarrow j,W_Q,t}<0$
から、窓内の時刻においていずれかの時刻別充足述語が偽であること、
構造的環境結合条件が成り立たないこと、
T-STAB 条件の不充足、閉塞、または崩壊を、
一般的な帰結として置かない。
この有限窓平均は関係品質評価の時間方向への集約であり、
状態遷移の結果そのものではない。
\end{remark}

%----------------------------------------------------------------------
\subsection{代表カテゴリー集合と六系分類}
\label{subsec:rel_categories}

本小節では、
関係カテゴリーを成立水準と内容上の系分類という二つの軸で整理する。

\begin{definition}[代表カテゴリーラベル集合（分類C）]
\label{def:category_set}
単方向関係記述子のカテゴリー成分の値域を
\[
\begin{aligned}
\mathcal C_{\rightarrow}:=\{&\mathrm{obs},\mathrm{ref},\mathrm{aff},
\mathrm{dep},\mathrm{sup},\mathrm{par},\mathrm{dom},\mathrm{mgmt},\mathrm{manip},\\
&\mathrm{exploit},\mathrm{use},\mathrm{comp},\mathrm{alert},\mathrm{avoid},\mathrm{excl}\}
\end{aligned}
\]
とし、
ACS同士の双方向成立カテゴリーの値域を
\[
\mathcal C_{\leftrightarrow}:=
\{\mathrm{coop},\mathrm{sym},\mathrm{conf}\}
\]
とする。
全体の代表ラベル集合は
\[
\mathcal C:=\mathcal C_{\rightarrow}\cup\mathcal C_{\leftrightarrow}
\]
であり、
単方向15種と双方向成立3種の計18種である。
ラベルは関係形式を記述する名称であり、
各カテゴリーの成立条件・閾値・規範的価値は定めない。
$\mathcal C_{\rightarrow}$ と $\mathcal C_{\leftrightarrow}$ はラベルとして交わらず、
全体の18種類は「単方向記述子のカテゴリー成分の15種類の値」と「3個の双方向補助ラベル」の総数を表す。
18種類すべてが、
同じ単方向記述子の一つの成分値になるわけではない。
\end{definition}

各ラベルの読解上の名称と成立水準は次のとおりである。
\[
\begin{array}{c|l|c@{\qquad}c|l|c}
\text{ラベル}&\text{名称}&\text{水準}&\text{ラベル}&\text{名称}&\text{水準}\\\hline
\mathrm{obs}&\text{観測}&\text{単方向}&\mathrm{ref}&\text{参照}&\text{単方向}\\
\mathrm{aff}&\text{親和}&\text{単方向}&\mathrm{dep}&\text{依存}&\text{単方向}\\
\mathrm{sup}&\text{支援}&\text{単方向}&\mathrm{par}&\text{寄生}&\text{単方向}\\
\mathrm{dom}&\text{支配}&\text{単方向}&\mathrm{mgmt}&\text{管理}&\text{単方向}\\
\mathrm{manip}&\text{操作}&\text{単方向}&\mathrm{exploit}&\text{搾取}&\text{単方向}\\
\mathrm{use}&\text{利用}&\text{単方向}&\mathrm{comp}&\text{競争}&\text{単方向}\\
\mathrm{alert}&\text{警戒}&\text{単方向}&\mathrm{avoid}&\text{回避}&\text{単方向}\\
\mathrm{excl}&\text{排除}&\text{単方向}&\mathrm{coop}&\text{協力}&\text{双方向}\\
\mathrm{sym}&\text{共生}&\text{双方向}&\mathrm{conf}&\text{対立}&\text{双方向}
\end{array}
\]
競争は片方向に記録可能であり、
相互競争は同種相互関係として扱う（備考~\ref{rem:competition_directionality}）。
表中の「単方向」は $\mathrm{cat}_{i\leftarrow j,t}$ の値域、
「双方向」は $\mathrm{cat}_{\leftrightarrow,ij,t}$ の値域を指す。
ラベルの一覧は読解用の対応であり、
個別ラベルの成立を判定する基準表ではない。

\begin{remark}[カテゴリー集合の役割]
\label{rem:category_set_role}
$\mathcal C$ は代表的な18ラベルであって、
全ての関係を網羅する分類体系ではない。
単方向ラベルの値域を与えることは、
あらゆる対象・時刻に
$\mathrm{cat}_{i\leftarrow j,t}$ が必ず割り当てられることを意味しない。
各ラベルの成立条件、
カテゴリー間の包含関係、
遷移規則、
強度尺度は本節では定めない。

一つの関係が複数の自然言語上の性質を帯びうることと、
本節で一時刻の記録に一つのカテゴリーラベルを付すことは区別する。
複合カテゴリーや多重ラベルの形式的な規則は後続の詳細化に委ねる。
\end{remark}

\begin{remark}[中立カテゴリーの除外]
\label{rem:neutral_exclusion}
本節は「中立」または「接触」を第19の代表ラベルとして追加しない。
記述子が定義されない場合、
関係ラベルを判定できない場合、
記述子が存在して三差分がすべてゼロの場合は互いに異なる。
未分類であることから関係の不存在・無影響・中立を推論しない。
$\mathrm{undetermined}$ は方向性指標にだけ用いる値であり、
$\mathcal C$ の外に関係カテゴリーを増設する名称としては用いない。
したがって「接触した」という記録だけから、
関係に固有のコストや
可行性関連量の変化があると結論することもしない。
カテゴリーを付けないという処置と、
値がゼロの六成分記述子を作る処置を混同しない。
\end{remark}

\begin{definition}[六系分類（分類C）]
\label{def:six_system_classification}
$\mathcal C$ の18ラベルを記述上の便宜により六つの部分集合へ分ける：
\[
\begin{aligned}
\mathcal C_{\mathrm{recog}}&:=\{\mathrm{obs},\mathrm{ref}\},&
\mathcal C_{\mathrm{aff}}&:=\{\mathrm{aff}\},\\
\mathcal C_{\mathrm{bind}}&:=\{\mathrm{dep},\mathrm{sup},\mathrm{coop},\mathrm{sym}\},&
\mathcal C_{\mathrm{ctrl}}&:=\{\mathrm{dom},\mathrm{mgmt},\mathrm{manip}\},\\
\mathcal C_{\mathrm{extract}}&:=\{\mathrm{par},\mathrm{exploit},\mathrm{use}\},&
\mathcal C_{\mathrm{comp\mbox{-}iso}}&:=\{\mathrm{comp},\mathrm{alert},\mathrm{avoid},\mathrm{excl},\mathrm{conf}\}.
\end{aligned}
\]
順に認識系、
親和系、
結合系、
制御系、
収奪系、
競合隔離系と呼ぶ。
列挙した六集合はラベルとして互いに交わらず、
その和集合は $\mathcal C$ に一致する。

この分割はラベルの列挙から従うもので、
カテゴリーの成立条件や価値序列を追加するものではない。
各系に対応する名称とラベルを、
読解用に次のとおり整理する。
\[
\begin{array}{c|l|l}
\text{系の記号}&\text{系名}&\text{含まれるカテゴリー}\\\hline
\mathcal C_{\mathrm{recog}}&\text{認識系}&\text{観測、参照}\\
\mathcal C_{\mathrm{aff}}&\text{親和系}&\text{親和}\\
\mathcal C_{\mathrm{bind}}&\text{結合系}&\text{依存、支援、協力、共生}\\
\mathcal C_{\mathrm{ctrl}}&\text{制御系}&\text{支配、管理、操作}\\
\mathcal C_{\mathrm{extract}}&\text{収奪系}&\text{寄生、搾取、利用}\\
\mathcal C_{\mathrm{comp\mbox{-}iso}}&\text{競合隔離系}&\text{競争、警戒、回避、排除、対立}
\end{array}
\]
\end{definition}

\begin{remark}[六系分類の位置づけ]
\label{rem:six_system_role}
六系は内容上の整理であり、
優劣や善悪の序列ではない。
結合系にも制約方向の変化が、
競合隔離系にも可行性保存方向の変化が記録されうる。
また、
同じカテゴリーに属する二つの記録で三差分が一致するとは限らない。
六系の所属は、
対象別効果の算出規則や可行性の変化方向の関数を与えない。
例えば依存や支援は、
条件によって個体側の選択を広げる記録にも、
資源負荷や制約に伴う記録にも現れうる。
同様に本論文は、
警戒や回避という名称から、
可行性関連量の低下を一般的な帰結として置かない。
いずれの場合も、
個別の差分の値と関係の因果的寄与は別々の問題である。
\end{remark}

\begin{remark}[六系分類と成立水準の独立性]
\label{rem:six_system_directionality}
六系はカテゴリーの内容上の整理であり、
成立水準を一律には指定しない。
結合系には単方向の依存・支援と双方向成立の協力・共生を含み、
競合隔離系には単方向の競争・警戒・回避・排除と双方向成立の対立を含む。
成立水準は定義~\ref{def:category_set} の
$\mathcal C_{\rightarrow}$ と $\mathcal C_{\leftrightarrow}$ の区別に従う。
したがって六系の一つを選んだだけでは、
その内部の単方向・双方向成立のどちらを記録すべきかは定まらない。
成立水準を特定する際は、
各カテゴリーがどちらのラベル集合に属するかを確認する。
\end{remark}

\begin{remark}[カテゴリーと規範評価の分離]
\label{rem:category_neutral}
協力、
寄生、
支配、
回避などは構造的な関係記述ラベルであり、
その名称だけで望ましさ、
有害性、
可行性の変化方向を定めない。

本論文は、カテゴリー名や六系への所属から、
\S\ref{sec:three_principles} の SM・EU・FR の時刻別充足述語の真偽が定まることも
一般的な帰結として置かない。
関係の解釈とACS自身の内部認識については備考~\ref{rem:recognized_relation_category} の区別を用いる。
規範的評価を与えないことは、
関係記録に付随する差分が
正・負のいずれであるかを問わないという意味ではない。
差分の符号を記録することと、
カテゴリー名に倫理的善悪を割り当てることを分離する。
\end{remark}

\begin{remark}[カテゴリー成立条件の未定義性]
\label{rem:category_phase1a}
代表18ラベルを列挙したことと、
各ラベルの成立条件を与えたことは異なる。
本節は必要十分条件、
判定閾値、
持続時間、
強度尺度、
時間的下位類型およびカテゴリー間の遷移規則を定義しない。
これらの決定論的な詳細化はPhase~1.$\alpha$へ委ねる。
例えば「どの程度の資源負荷を伴えば搾取と呼ぶか」、
「どの程度の相互性があれば協力と呼ぶか」は、
本節のラベル集合から導かれる判定規則ではない。

長期的・累積的な支援のような時間的下位類型も、
本節の一時刻ラベルからは決まらない。
\end{remark}

\begin{remark}[派生カテゴリーと複合カテゴリー]
\label{rem:derived_composite_categories}
委任、
取引、
所属、
親交などは複数の関係カテゴリーの組合せとして
記述されうるが、
本節はその構成規則を定めない。
あるカテゴリーから別のカテゴリーへの時間的変化も、
本節の固定した記述形式だけで自動的に生成されるものではない。

派生・複合カテゴリーの定義と遷移規則はPhase~1.$\alpha$へ委ねる。
例えば参照・依存・支援・管理・利用等の記録を組み合わせることで
委任や取引といった関係を記述する可能性はあるが、
本論文はその組合せを正式なカテゴリー定義とはしない。
支援から依存へ、
警戒から対立へといった分類変化も、
単一時刻の関係記述子だけからは演繹されない。
\end{remark}

\begin{remark}[カテゴリーと認識された関係性]
\label{rem:recognized_relation_category}
$\mathrm{cat}_{i\leftarrow j,t}$ は理論上の記述ラベルであり、
$\Sigma_i$ がその関係を同じラベルで認識しているとは仮定しない。
その区別は、
\S\ref{sec:tipa} の内部識別と
本節の記述上の分類の役割差とも整合する。

認識された関係の指標 $\mathrm{recog}_{i\leftarrow j,t}$ 等は本節では記号候補として挙げるにとどめ、
その値域・算出規則・カテゴリーとの対応は定義しない。
したがって、
本節の $\mathrm{cat}_{i\leftarrow j,t}$ からACS自身の認識内容や正誤を推論せず、
詳細化をPhase~1.$\alpha$へ委ねる。
例えば個体側が支援と認識している関係を理論上どのように記録するかは、
ACS内部の認識内容と外部記述の対応を別途定義しなければ決まらない。
本論文の $\mathrm{recog}_{i\leftarrow j,t}$ はその詳細化を予告する名称にとどまる。
\end{remark}

\begin{remark}[カテゴリーと閉塞・崩壊の節]
\label{rem:categories_to_sec7}
18カテゴリーや六系への所属は、
\S\ref{sec:collapse_cond} の
$\mathfrak{P}_{\mathrm{blk}}$ または $\mathfrak{P}_{\mathrm{col}}$ への帰属条件ではない。
関係記録に付随して可行性関連量が変化しても、
閉塞・崩壊の判定には状態量と軌道記録の条件を別途要する。
局所サイズ指標がゼロとなることと、
$A_{i,t}=\emptyset$ も同一視しない。
結合系や親和系であっても、
個体側の条件によって制約方向の差分が
記録される可能性は排除されない。
他方、
競合隔離系のラベルをもつ記録から
直ちに閉塞・崩壊への到達可能性が導かれるわけではない。
経路集合との接続は、
後続節で規定する時刻域と軌道の条件に依存する。
\end{remark}

%----------------------------------------------------------------------
\subsection{射程外事項・後続節への接続と小括}
\label{subsec:rel_phase1a}

本小節では、
確立した記述骨格とPhaseごとの委譲事項を区別し、
閉塞・崩壊の節と組織ACSの節への接続を整理する。

\begin{remark}[本節で確立した最小構造]
\label{rem:sec6_established}
本節が与えるのは、
外部対象の記述上の三分類、
条件付きで整形式となる六成分の単方向関係記述子、
ACS同士の双方向合成と補助方向性、
評価規則に相対する関係品質、
代表カテゴリー18種と六系分類である。
局所指標・有限窓差分・六成分記述子は、
適用域と値が与えられた場合に
記述を構成するための形式であり、
任意の関係を必ず数値化できるという存在定理ではない。
記述子の値が各時刻に実在すること、
カテゴリーの成立条件、
または関係の固有の因果効果を確立したものではない。
\end{remark}

\begin{remark}[Phase~1.$\alpha$ への委譲事項]
\label{rem:phase1a_delegation}
Phase~1.$\alpha$には、
決定論的・離散時間の枠内で、
次の詳細化を委ねる（\S\ref{subsec:phase1a}）。
ここに挙げる名称は検討課題であり、
現段階の定義や成立定理ではない。
\begin{enumerate}
 \item 18カテゴリーそれぞれの成立条件・成立閾値
 （定義~\ref{def:category_set}、
備考~\ref{rem:category_phase1a}）。
 \item 関係の強度・持続時間と、長期的・累積的な支援等の時間的下位類型。
 \item 関係の解除可能性・代替可能性・媒介性・終了条件。
 \item 被支配・被管理・被操作・被搾取・被利用等の応答側カテゴリー
 （備考~\ref{rem:homo_hetero_rel}）。
 \item 支配や寄生等が双方向に現れた場合の複合的な記述規則
 （備考~\ref{rem:no_mutual_form}）。
 \item 委任・取引・所属・親交等の派生・複合カテゴリーとカテゴリー間遷移
 （備考~\ref{rem:derived_composite_categories}）。
 \item ACS内部が認識した関係と理論上のカテゴリーの対応
 （備考~\ref{rem:recognized_relation_category}）。
 \item 非ACSエージェント等との双方向的相互作用を記述する拡張形式
 （備考~\ref{rem:nACS_limit}）。
 \item 方向性指標 $\mathrm{dir}_{ij,t}$ の「主要」の定量的基準
 （備考~\ref{rem:dir_undefined}）。
 \item 局所サイズ指標 $\mathrm{sz}_i$ の評価域・具体形、
 有限近似や固定精度の被覆数による具体化、
および必要に応じたサイズ性の追加条件
 （定義~\ref{def:rel_size_functional}）。
 \item 協力・共生・対立の双方向成立規則、単方向カテゴリーとの対応、
 ならびに競争の単方向形と相互形を区別する定量的基準
 （定義~\ref{def:bidirectional_category}、
備考~\ref{rem:competition_directionality}）。
\end{enumerate}

対象同定の具体的手続きと、
有限窓での関係記録を生成・更新する規則も、
必要に応じてこの詳細化の側で扱う。
これらの詳細化を、
既に本節で定義されたものとして用いない。
\end{remark}

\begin{remark}[Phase 2への委譲事項]
\label{rem:phase2_delegation}
Phase~2には、
確率的な成立・遷移・終了、
連続時間化、
外部対象の内部動学を含む相互作用、
および社会的ネットワーク動学を委ねる（\S\ref{subsec:phase_outlook}）。
分類写像が時刻によって変わる対象の記述も、
本節の固定分類を超える課題である。
外部対象の内部状態と内部動学を伴う相互作用や、
複数の関係がネットワーク上で変化する過程は、
本節の単方向・双方向の記述子だけでは規定されない。

本論文で時刻添字を付したことを、
確率的遷移則やネットワーク動学の導入と読み替えない。
この委譲は、
個別の記録窓における差分を確率量とみなすことを意味しない。
\end{remark}

\begin{remark}[閉塞・崩壊の節と組織ACSの節への接続]
\label{rem:sec6_to_sec78}
\S\ref{sec:collapse_cond} は閉塞・崩壊と各到達経路を、
\S\ref{sec:org_acs} は組織自身についてACS構造が成立する組織ACSの条件を扱う。
本節の関係記述子は両節への参照素材であって、
失敗状態の判定や組織ACS成立の十分条件ではない。

\S\ref{sec:collapse_cond} への接続では、関係記録に付随する三差分と
個体の軌道上の閉塞・崩壊条件を区別する。
\S\ref{sec:org_acs} への接続では、関係の順序対が記録されることと、
組織自身についての組織水準のACS構造の成立とを区別する。
本論文は、これらの後続節を読み込まない単独の本節から、
両節の成立判定を一般的な帰結として置かない。
\end{remark}

\begin{remark}[本節の記述方針の総括]
\label{rem:sec6_summary}
本節の関係カテゴリー、
観測窓差分、
局所サイズ指標は
価値中立的な構造記述であり、
規範判断、
対象ごとの因果寄与、
閉塞・崩壊への到達または組織ACSの成立を証明するものではない。
さらに、
指標値の大小を可行行動集合の包含関係へ読み替えることや、
全ての外部対象・全時刻で関係記述子が定義されることも主張しない。
差分を比べるときは窓長と指標の尺度を固定し、
双方向の記録では
それぞれの個体の記録有効時刻域と関係継続条件を区別する。
後続節との接続は、
各節が別途定める定義域・条件のもとでのみ行う。
\end{remark}


%================ Section 7 崩壊・閉塞・崩壊経路集合 ================
%======================================================================
\section{崩壊・閉塞・崩壊経路集合}
\label{sec:collapse_cond}

\S\ref{sec:acs_def} で与えたACSの形式的定義、
\S\ref{sec:three_principles} の可行性構造、
\S\ref{sec:extended_dynamics} の時刻依存量、
および
\S\ref{sec:rel_min} の関係性記述を受けて、
本節では単一ACS個体 $\Sigma_i$ における
閉塞、記述的崩壊前状態帯、外的回復可能性、崩壊、
およびそれらの軌道記録上の表現を整理する。

本節の主眼は、
これらを一つの失敗概念へまとめることではなく、
互いに異なる水準の状態として区別することにある。
とくに、
閉塞、整形式実現の不成立および崩壊は、
それぞれ異なる判定条件に基づく概念であり、
相互に同一視しない。
本節はこの区別を基本とする。
閉塞は、ある時刻の実現環境のもとで
可行行動が失われることに関する局所的な状態である。
これに対して崩壊は、
単一時刻の機能低下や閉塞そのものではなく、
本枠組みにおいて終端的に扱われる状態である（\S\ref{sec:acs_def} 定義~\ref{def:collapse}）。

本節ではさらに、
崩壊周辺の状態を記述するための状態帯と、
閉塞または崩壊への到達を記録する経路集合を導入する。
本節は、崩壊の一般原因論、確率的到達理論、
あるいは崩壊までの距離理論を構成するものではない。

Phase~1 の本節が直接対象とするのは単一ACS個体である。

本節で新たに用いる主要記号の一覧は付録~\ref{app:symbols_s7} に示す。
なお、
$A_{i,t}$、$\mathrm{WF}_{i,t}$、
$\mathrm{Desc}_t$、$Q$ 等は
前節までに導入された記号であり、
本節では必要に応じて参照するが再定義しない。

\begin{remark}[本節における三原理と機能の記述水準]
\label{rem:s7_principle_levels}
本節で三原理と機能に言及するときは、次の水準を区別する。
\begin{description}[style=nextline]
  \item[原理]
  自己維持（SM）・環境利用（EU）・可行応答（FR）の三原理は、
  \S\ref{sec:three_principles} で定める、存続中のACSに対する制約軸である。
  本節は、原理そのものについて成立・不成立を述べない。
  \item[時刻別充足述語]
  各原理について \S\ref{sec:three_principles} で定める時刻別の条件であり、
  真偽をもつのはこの水準である。
  SM の時刻別充足述語は、\S\ref{sec:three_principles} で定める適用域において
  $x_t\in K$ のとき真となる。
  EU の時刻別充足述語は同節の適用域で定義され、
  FR の時刻別充足述語（局所 FR）はさらに定義域 $\mathcal{D}_{\mathrm{FR}}$ 上でだけ定義される。
  本節は各時刻の判定を「SM の時刻別充足述語が真である」「偽である」のように記述し、
  「偽である」は、時刻別充足述語が定義されたうえで偽となることを指す。
  定義されない時刻を、偽である時刻に含めない。
  \item[機能]
  原始的な機能の各区分に関して個体が有する基礎的な機能と基盤機能
  （\S\ref{sec:acs_def}）、
  \S\ref{sec:three_principles} の構造的機能、
  および \S\ref{sec:tipa} の TIPA・PSL・ISL の標準機能形式は、機能の水準に属する。
  作動・停止・低下を述べる対象はこの水準であり、
  時刻別充足述語の真偽とは区別する。
  \item[診断量]
  機能軸、健全度、記述的状態帯、$\mu_{\mathrm{eff},t}$ などの分類C量である。
\end{description}
本節は、機能が作動していることから時刻別充足述語が真であることを一般的な帰結として置かず、
時刻別充足述語が偽であることから特定の機能の停止を一般的な帰結として置かない。
また、原理ごとの専用の機能を置かない。
\end{remark}

以下では、
\S\ref{subsec:blockage} で閉塞と諸概念との関係を整理し、
\S\ref{subsec:precollapse_states} で崩壊周辺の記述的状態帯を整理する。
続いて、
\S\ref{subsec:external_recovery} で外的回復可能性と同一個体性を扱い、
\S\ref{subsec:collapse_def} で崩壊（\S\ref{sec:acs_def} で定義）を整理する。
その後、
\S\ref{subsec:blk_path_set} と
\S\ref{subsec:col_path_set} で
閉塞・崩壊を軌道記録上の経路集合として表現し、
\S\ref{subsec:inflow_paths} で
内部動学および関係性との記述的接続を整理する。
最後に、
\S\ref{subsec:to_section8} で
個体レベルの定義を組織ACSの議論へ接続する際の境界を示す。

%----------------------------------------------------------------------
\subsection{閉塞}
\label{subsec:blockage}

本小節では、ACS個体 $\Sigma_i$ の閉塞（\S\ref{sec:acs_def} 定義~\ref{def:blockage_sec2}）と、
現在状態の可行性、必要環境充足、
$E_{\mathrm{req}}$ ロバスト可行性核、
整形式実現、および崩壊との関係を整理する。

\subsubsection{閉塞の定義の参照}

閉塞は \S\ref{sec:acs_def} の定義~\ref{def:blockage_sec2}（可行行動集合が空であること）による。
本節はこれを再定義せず、時刻 $t$ の内部状態 $x_t\in X$ と
実現外部環境状態 $e_t\in E_{\mathrm{ext}}$ に対する
\[
  A_{i,t}
  =
  A(x_t,e_t)
  =
  \bigl\{
    u\in U
    \bigm|
    f(x_t,u,e_t)\in K
  \bigr\}
  =\emptyset
\]
の形で用いる。

\begin{remark}[閉塞の持続評価関数による書き換え]
\label{rem:blockage_pi_form}
可行集合 $K=\{x\in X\mid\Pi(x)\ge0\}$ の定義より、
\[
  A_{i,t}=\emptyset
  \iff
  \forall u\in U,\;
  \Pi\!\bigl(f(x_t,u,e_t)\bigr)<0
\]
である。
\end{remark}

\begin{remark}[閉塞定義における環境状態の型]
\label{rem:blockage_environment_type}
定義~\ref{def:blockage_sec2} に現れる $e_t$ は、
7タプルの第三成分 $E$ の要素ではなく、
外部環境状態空間
$e_t\in E_{\mathrm{ext}}$
の要素である。

\[
  E
  \subseteq
  \{0,1\}^{X\times E_{\mathrm{ext}}}
\]
は必要環境要因の集合であり、
環境状態そのものの集合ではない。
したがって本節でも、
実現外部環境状態は
$E_{\mathrm{ext}}$ の要素として扱い、
必要環境要因集合 $E$ とは型を区別する。

また、状態遷移写像の型は
\[
  f:
  X\times U\times E_{\mathrm{ext}}
  \longrightarrow X
\]
であり、
閉塞はこの実現外部環境状態 $e_t$ に相対して定まる。
\end{remark}

\begin{remark}[閉塞は「遷移不能」を意味しない]
\label{rem:blockage_not_no_transition}
閉塞
$A_{i,t}=\emptyset$
が意味するのは、
実現外部環境状態 $e_t$ のもとで
\emph{次状態を $K$ 内に保つ行動が存在しない}ことである。

これは、
遷移写像 $f(x_t,u,e_t)$ 自体が未定義であること、
物理的変化が停止すること、
または次時刻が存在しないことを意味しない。

したがって閉塞は、
「遷移そのものの不存在」ではなく、
「一ステップ先の可行性を保持する行動の不存在」
として理解される。
\end{remark}

\subsubsection{閉塞と SM・必要環境充足・整形式実現の区別}

\begin{remark}[閉塞と SM の時刻別充足述語は異なる時点条件である]
\label{rem:blockage_vs_sm}
\S\ref{sec:three_principles} の SM の時刻別充足述語は、
時刻 $t$ の現在状態について
$\Pi(x_t)\ge0 \iff x_t\in K$
のとき真となる、単一時刻の可行状態条件である。

これに対し閉塞は、
同じ時刻 $t$ の $(x_t,e_t)$ のもとで
一ステップ後を $K$ 内に保つ行動が存在するかを問う条件である。

したがって、
\[
  \text{SM の時刻別充足述語}
  =
  \text{現在状態の可行性}
\]
と
\[
  \text{閉塞}
  =
  \text{実現環境下での一ステップ可行行動の消失}
\]
は異なる。

本論文は一般に、
\[
  x_t\in K
  \Longrightarrow
  A_{i,t}\neq\emptyset
\]
も、
\[
  A_{i,t}\neq\emptyset
  \Longrightarrow
  x_t\in K
\]
も仮定しない。

したがって、
SM の時刻別充足述語の真偽と
閉塞の成立・不成立を
同一視しない。
\end{remark}

\begin{remark}[閉塞と必要環境充足は異なる]
\label{rem:blockage_vs_required_environment}
必要環境充足集合は
\[
  E_{\mathrm{req}}(x_t)
  =
  \bigl\{
    e\in E_{\mathrm{ext}}
    \bigm|
    \forall\eta\in E,\;
    \eta(x_t,e)=1
  \bigr\}
\]
であり、
状態 $x_t$ において
すべての必要環境要因を充足する
外部環境状態の集合を表す。

一方、閉塞は
実際に実現した $e_t$ に相対する
$A(x_t,e_t)=\emptyset$
という行動可行性条件である。

したがって、
$e_t\in E_{\mathrm{req}}(x_t)$
は閉塞の定義条件ではなく、
また
\[
  E_{\mathrm{req}}(x_t)=\emptyset
\]
も閉塞の定義と同一ではない。

必要環境充足と閉塞のあいだに
追加の含意を主張する場合には、
状態遷移および可行性に関する
追加条件を要する。
\end{remark}

閉塞と同じ時刻の整形式実現の関係は、
\S\ref{sec:acs_def} の備考~\ref{rem:blockage_not_collapse_sec2} に従う。

\subsubsection{閉塞の時刻局所性}

\begin{remark}[閉塞は吸収状態として定義されない]
\label{rem:blockage_time_dependence}
$A_{i,t}$ は
$x_t$ および $e_t$ に相対して定まるため、
閉塞は時刻 $t$ ごとの条件として判定される。

ある時刻 $t$ で
$A_{i,t}=\emptyset$
が成立しても、
後続時刻 $s>t$ における
$A_{i,s}$
の値は、
その時刻の状態 $x_s$、
実現外部環境状態 $e_s$、
および必要に応じて
後続節で扱う外的作用等に依存して
別途定まる。

したがって、
閉塞それ自体を
吸収状態として定義しない。

特に、本論文は、
$A_{i,t}=\emptyset$
から
\[
  A_{i,s}=\emptyset
  \qquad
  \forall s>t
\]
を一般的な帰結として置かない。
\end{remark}

\subsubsection{閉塞と $E_{\mathrm{req}}$ ロバスト可行性核}

\S\ref{sec:three_principles} では、
状態依存の必要環境充足集合
$E_{\mathrm{req}}(x)$ に対して
一つの共通行動で可行性を保持できる状態を捉えるため、
$E_{\mathrm{req}}$ ロバスト可行性核
$\mathrm{Viab}_f(K)$
を導入した。

ここで重要なのは、
閉塞
$A(x_t,e_t)=\emptyset$
が\emph{実現した一つの環境状態 $e_t$}
に相対するのに対し、
$\mathrm{Viab}_f(K)$ は
\emph{$E_{\mathrm{req}}(x)$ のすべての環境状態}
に対するロバスト不変性を要求する点である。

したがって、
両者を接続するには、
少なくとも実現環境が
必要環境充足集合に属するという条件が必要となる。

\begin{proposition}[閉塞と $E_{\mathrm{req}}$ ロバスト可行性核の関係]
\label{prop:blockage_viability_kernel}
時刻 $t\in\mathbb{N}_0$ において
\[
  e_t\in E_{\mathrm{req}}(x_t)
\]
かつ
\[
  A_{i,t}=\emptyset
\]
が成立するならば、
\[
  x_t\notin\mathrm{Viab}_f(K)
\]
が成立する。
\end{proposition}

\begin{proof}
背理法による。

\[
  e_t\in E_{\mathrm{req}}(x_t),
  \qquad
  A_{i,t}=\emptyset
\]
を仮定し、さらに
\[
  x_t\in\mathrm{Viab}_f(K)
\]
と仮定する。

定義~\ref{def:viab_core} より、
\[
  \mathrm{Viab}_f(K)
  =
  \bigcup_{K'\in\mathfrak{K}} K',
\]
ここで $\mathfrak{K}$ は
$E_{\mathrm{req}}$ 不変な
$K_{\mathrm{req}}$ の部分集合の族である。

したがって
$x_t\in\mathrm{Viab}_f(K)$ ならば、
ある
\[
  K'\in\mathfrak{K}
\]
が存在して
\[
  x_t\in K'
\]
である。

$K'$ の $E_{\mathrm{req}}$ 不変性より、
ある一つの行動
\[
  u^\star\in U
\]
が存在して、
\[
  \forall e\in E_{\mathrm{req}}(x_t),\qquad
  f(x_t,u^\star,e)\in K'
\]
が成立する。

仮定より
\[
  e_t\in E_{\mathrm{req}}(x_t)
\]
であるから、
特に
\[
  f(x_t,u^\star,e_t)\in K'.
\]

また、
定義~\ref{def:viab_core} より
\[
  K'
  \subseteq
  K_{\mathrm{req}}
  \subseteq
  K
\]
であるため、
\[
  f(x_t,u^\star,e_t)\in K.
\]

したがって、
定義~\ref{def:at} より
\[
  u^\star\in A_{i,t},
\]
ゆえに
\[
  A_{i,t}\neq\emptyset
\]
となる。

これは
\[
  A_{i,t}=\emptyset
\]
に矛盾する。

したがって
\[
  x_t\notin\mathrm{Viab}_f(K).
\]
\end{proof}

\begin{remark}[実現環境が $E_{\mathrm{req}}(x_t)$ 外にある場合]
\label{rem:blockage_viab_requires_ereq}
命題~\ref{prop:blockage_viability_kernel} における
$e_t\in E_{\mathrm{req}}(x_t)$
は本質的な前件である。

実現環境が
$e_t\notin E_{\mathrm{req}}(x_t)$
である場合、
閉塞
$A(x_t,e_t)=\emptyset$
が成立しても、
それだけから
$x_t\notin\mathrm{Viab}_f(K)$
を一般には導けない。

なぜなら、
$\mathrm{Viab}_f(K)$ の不変性条件が量化するのは
$E_{\mathrm{req}}(x_t)$ 内の外部環境状態であり、
その外部にある実現環境 $e_t$ に対する
一ステップ可行性は
ロバスト可行性核の定義条件に含まれないからである。
例えば $X=\{0,1\}$、$\Pi(x)=x-1$（したがって $K=\{1\}$）、$U=\{u\}$、
$E_{\mathrm{ext}}=\{e^1,e^2\}$、$E=\{\eta\}$、$\eta(x,e)=1\iff e=e^1$ とし、
$f(1,u,e^1)=1$、それ以外の引数では $f=0$ とする。
このとき $\{1\}$ は $E_{\mathrm{req}}$ 不変であるから $1\in\mathrm{Viab}_f(K)$ であるが、
$x_t=1$、$e_t=e^2$ では $A_{i,t}=\emptyset$ である（型の水準の例）。

したがって、
\[
  \text{閉塞}
  +
  \text{実現必要環境充足}
  \Longrightarrow
  \text{ロバスト可行性核外}
\]
という条件付き含意としてのみ用いる。
\end{remark}

\begin{remark}[逆方向の非採用]
\label{rem:viability_kernel_not_blockage}
命題~\ref{prop:blockage_viability_kernel} の逆向き、
すなわち
\[
  x_t\notin\mathrm{Viab}_f(K)
  \Longrightarrow
  A_{i,t}=\emptyset
\]
は、
Phase~1 の一般的主張として採用しない。

$x_t\notin\mathrm{Viab}_f(K)$ は、
$E_{\mathrm{req}}(x)$ に対する
ロバスト不変性を満たす状態集合の外部にあることを意味する。

しかし本論文は、
そのことから
実現環境 $e_t$ のもとで
一ステップだけ $K$ 内へ遷移できる行動が
存在しないことを一般的な帰結として置かない。

したがって、
ロバスト可行性核の外部にある状態でも、
$A_{i,t}\neq\emptyset$
でありうる。

本論文は、
核外性それ自体を
「閉塞へ向かうこと」
または
「閉塞経路集合への流入」
と同一視しない。
そのような動学的接続を述べるには、
追加の条件を要する。
\end{remark}

\subsubsection{閉塞と崩壊の区別}

閉塞と崩壊の区別は、
\S\ref{sec:acs_def} の備考~\ref{rem:blockage_not_collapse_sec2} に従う。
閉塞から後続時刻に可行行動が存在する状態へ移行しうることは、
備考~\ref{rem:blockage_time_dependence} に従い、
その外的回復可能性は \S\ref{subsec:external_recovery} で扱う。

\begin{remark}[閉塞・SM・整形式実現・崩壊の記述水準]
\label{rem:blockage_fourway_distinction}
本節では、少なくとも次の四条件を区別する。

\[
  \Pi(x_t)\ge0
  \qquad
  \text{（SM の時刻別充足述語／現在状態の可行性）},
\]
\[
  A_{i,t}=\emptyset
  \qquad
  \text{（閉塞／一ステップ可行行動の消失）},
\]
\[
  \mathrm{WF}_{i,t}=0
  \qquad
  \text{（整形式実現の不成立）},
\]
\[
  \mathrm{Col}_{i,t}=1
  \qquad
  \text{（崩壊）}.
\]

これらは異なる形式対象に関する条件であり、
本論文は追加仮定なしに
相互の同値性を置かない。

とりわけ、本論文は、
閉塞 $A_{i,t}=\emptyset$ から
$\mathrm{WF}_{i,t}=0$ および $\mathrm{Col}_{i,t}=1$ を
いずれも一般的な帰結として置かない。

閉塞と崩壊の詳細な区別、
および崩壊側から見た非回復性（\S\ref{sec:acs_def} 定義~\ref{def:collapse} の (SC-2)）の読みは
\S\ref{subsec:collapse_def}
に委ねる。
\end{remark}

\begin{remark}[「閉塞」という用語の使用範囲]
\label{rem:blockage_terminology}
本論文では、
$A_{i,t}=\emptyset$
が成立する状態だけを
「閉塞」と呼ぶ。

したがって、
SM の時刻別充足述語が偽であること、
必要環境の非充足、
ロバスト可行性核の外部化、
整形式実現の不成立、
全機能停止、
または崩壊を、
それだけを理由として
「閉塞」と呼ばない。

この用語規約により、
可行行動の消失と、
それ以外の失敗関連状態との
概念的混同を避ける。
閉塞を維持の局面と回復の局面に分けて読むこと、
閉塞が一段遷移仕様に相対すること、
および日常語の行き詰まりとの関係は、
\S\ref{sec:acs_def} の備考~\ref{rem:blockage_scope_sec2} による。
\end{remark}

\begin{remark}[閉塞の定義に含まれない事項]
\label{rem:blockage_in_precollapse}
閉塞から崩壊への決定論的遷移、
閉塞の持続、
閉塞からの回復、
または閉塞経路集合
$\mathfrak{P}_{\mathrm{blk}}$
への動学的流入条件は、
定義~\ref{def:blockage_sec2} に含まれず、
本論文はそれらを同定義だけからの一般的な帰結として置かない。
\end{remark}

%----------------------------------------------------------------------
\subsection{崩壊周辺の記述的状態帯}
\label{subsec:precollapse_states}


本小節では、
閉塞と崩壊のあいだに存在しうる
機能低下・作動停止等を記述的に整理するため、
補助的な状態帯
\[
  \mathsf{Z}_0,\ldots,\mathsf{Z}_5
\]
を導入する。

ここで重要なのは、
これらの状態帯が
Phase~1 における正式な状態空間 $X$ の分割ではなく、
閉塞・機能低下・崩壊を相互に混同しないための
\emph{分類Cの記述用ラベル}であるという点である。

本小節では、
$\mathsf{Z}_0$
を通常作動の\emph{基準帯}、
\[
  \mathsf{Z}_1,\ldots,\mathsf{Z}_4
\]
を\emph{記述的崩壊前状態帯}、
$\mathsf{Z}_5$
を崩壊状態に対応する終端側の記述帯として扱う。

したがって、
本節で「記述的崩壊前状態帯」と呼ぶ範囲は
$\mathsf{Z}_1,\ldots,\mathsf{Z}_4$
とし、
$\mathsf{Z}_0$ は含めない。

\subsubsection{機能作動・健全度量との関係}

\begin{remark}[機能作動記述の参照]
\label{rem:functional_axis_reference}
本小節でいう
「部分的機能不全」および
「全機能停止」は、
ACS個体 $\Sigma_i$ の
TIPA、PSL、ISL によって表現される
標準機能形式の作動状態に関する
\emph{記述的な呼称}である。

これらの機構・基盤・実装部位の正式な構造と
同一時刻における作動記述は
\S\ref{sec:tipa}
に委ねる。

本小節は、
本小節独自の新しい機能軸、
作動指標、
閾値、
または全機能停止判定写像を導入しない。

したがって、
「全機能停止」という呼称は、
TIPA、PSL、ISL に関係する
標準機能が総体として
機能的に作動していないと記述される状態を指すが、
その判定の具体的な操作化は
\S\ref{sec:tipa}
の機能軸（\S\ref{subsec:functional_axis}）に
従う。

また、
この記述は
特定の生物器官、
神経回路、
ハードウェアモジュール、
または特定アルゴリズムが
一律に停止していることを要求するものではない。
\end{remark}

\begin{remark}[健全度量と状態帯を同一視しない]
\label{rem:health_not_zone_identity}
\S\ref{sec:extended_dynamics} で扱う
PSL・ISL 関連の健全度量
\[
  h_{\mathrm{psl},t},
  \qquad
  h_{\mathrm{mem},t},
  \qquad
  h_{\mathrm{sens},t},
  \qquad
  h_{\mathrm{eff},t}
\]
は、
時刻間変化を記述・診断・比較するための
分類C量である。

これらの健全度成分 $h_{q,t}$ が
所定の閾値以下であることと、
特定の状態帯 $\mathsf{Z}_j$ への分類とのあいだに、
Phase~1 で普遍的な同値条件を置かない。

健全度の値だけから閉塞・崩壊を一般的な帰結として置かないことは、
\S\ref{sec:extended_dynamics} の備考~\ref{rem:s5_health_noncollapse} に従う。

逆に、
閉塞または崩壊が成立したことだけから
特定の健全度成分が
必ず $0$ になるとも仮定しない。

したがって、
健全度量と状態帯は関連しうるが、
別個の記述対象である。
\end{remark}

\subsubsection{記述的状態帯}

\begin{remark}[閉塞・崩壊周辺の記述的状態帯]
\label{rem:descriptive_zones}
閉塞、機能低下、および崩壊の位置づけを
記述的に区別するため、
次の補助ラベルを用いる：
\[
  \mathsf{Z}_0,\;
  \mathsf{Z}_1,\;
  \mathsf{Z}_2,\;
  \mathsf{Z}_3,\;
  \mathsf{Z}_4,\;
  \mathsf{Z}_5.
\]

これらは
分類Cの記述用ラベルであり、
ACS自身が
$\mathsf{Z}_k$ という内部変数を
保持または計算することを要求しない。

また、
$\mathsf{Z}_k$ の添字は
時間順序、
因果順序、
重症度の厳密な尺度、
距離、
または遷移確率を表さない。

\medskip
\noindent
\textbf{(a) $\mathsf{Z}_0$：通常作動・基準帯。}

$\mathsf{Z}_0$ は、
閉塞・崩壊周辺の状態を比較するための
通常作動側の参照帯である。

典型的には、
主要な機能作動が維持され、
閉塞
$A_{i,t}=\emptyset$
や顕著な機能低下が
前景化していない状態を表す。

ただし、
$\mathsf{Z}_0$ に属するための
普遍的な健全度閾値、
T-STAB 閾値、
または三原理の時刻別充足述語に関する閾値を
本小節では定めない。

$\mathsf{Z}_0$ は
障害状態を意味しないため、
本論文では
\emph{崩壊前状態帯には含めない}。

\medskip
\noindent
\textbf{(b) $\mathsf{Z}_1$：軽度障害顕在化帯。}

$\mathsf{Z}_1$ は、
通常作動との差異として、
軽度の機能低下、
慢性的負荷、
維持コスト増大、
または健全度低下に相当する兆候が
記述的に顕在化している状態を表す。

この呼称は、
特定の健全度量が
所定の閾値を下回ったことを
普遍的に意味しない。

また、本論文は、
$\mathsf{Z}_1$ であることから
SM・EU・FR のいずれかの時刻別充足述語が偽であること、
閉塞、
整形式実現の不成立、
または崩壊を
一般的な帰結として置かない。

\medskip
\noindent
\textbf{(c) $\mathsf{Z}_2$：部分的機能不全帯。}

$\mathsf{Z}_2$ は、
TIPA、PSL、ISL によって表現される
標準機能の一部について、
機能的な低下または停止が
通常作動と区別される程度に
記述される状態を表す。

ただし、
部分的機能不全が存在しても、
同一ACS個体としての構造的自己が
直ちに失われるとはしない。

また本論文は、
$\mathsf{Z}_2$ であることだけから
$A_{i,t}=\emptyset$
または
$\mathrm{WF}_{i,t}=0$
または
$\mathrm{Col}_{i,t}=1$
を一般的な帰結として置かない。

\medskip
\noindent
\textbf{(d) $\mathsf{Z}_3$：全機能停止帯。}

$\mathsf{Z}_3$ は、
\S\ref{sec:tipa} において
標準機能形式として記述される
TIPA、PSL、ISL に関係する作動機能が、
総体として機能的に停止していると
記述される状態を表す。

本小節では、
この状態を独立した
Phase~1 の形式的 Definition として
新設しない。

また、
全機能停止は
閉塞と同一ではない。

すなわち、
全機能停止帯 $\mathsf{Z}_3$ と記述されることと、
$A_{i,t}=\emptyset$ であることは
同義ではない。

さらに、本論文は、
全機能停止だけから
整形式実現の不成立または崩壊を
一般的な帰結として置かない。

\medskip
\noindent
\textbf{(e) $\mathsf{Z}_4$：閉塞帯。}

$\mathsf{Z}_4$ は、
定義~\ref{def:blockage_sec2} の
閉塞
$A_{i,t}=\emptyset$
が成立する状態に対応する。

$\mathsf{Z}_4$ は、
$\mathsf{Z}$ 帯のうち、
本小節以前に
厳密な形式条件が確定している状態帯である。

ただし\S\ref{subsec:blockage}で示したとおり、
閉塞は
SM の時刻別充足述語が偽であること、
必要環境非充足、
ロバスト可行性核外、
整形式実現の不成立、
または崩壊と同一ではない。
$\mathsf{Z}_4$ に記述される時刻は、
\S\ref{sec:three_principles} の閉塞の二分
（命題~\ref{prop:blockage_dichotomy}）によって
(B1) と (B2) に記述的に分けうる。
この二分は診断上の分割であり、原因の分類ではない（備考~\ref{rem:blockage_dichotomy}）。

\medskip
\noindent
\textbf{(f) $\mathsf{Z}_5$：崩壊帯。}

$\mathsf{Z}_5$ は、
\S\ref{sec:acs_def} 定義~\ref{def:collapse} の
崩壊指標
$\mathrm{Col}_{i,t}=1$
が成立する終端側の状態に対応する。

崩壊の正式な定義と吸収性は
\S\ref{sec:acs_def}（定義~\ref{def:collapse}、命題~\ref{prop:collapse_basic}）で与えられ、
本枠組みにおける非回復性の射程は
\S\ref{subsec:collapse_def}
で扱う。

したがって本小節では、
$\mathsf{Z}_5$ の定義条件を
先取りして追加しない。
\end{remark}

\subsubsection{状態帯の論理的関係}

\begin{remark}[$\mathsf{Z}$ 帯は状態空間分割ではない]
\label{rem:Z_not_partition}
$\mathsf{Z}_0,\ldots,\mathsf{Z}_5$ は、
状態空間 $X$ の部分集合族として
形式的に定義されたものではない。

したがって、
本論文は、これらのラベルに基づく分類が
$X$ の被覆や排反な分割を与えるとは仮定しない。

特に、
$\mathsf{Z}_1$、
$\mathsf{Z}_2$、
$\mathsf{Z}_3$ は
記述的な機能状態ラベルであるため、
閉塞
$\mathsf{Z}_4$
と同一時刻に重なって記述される場合を
排除しない。

一方、
$\mathsf{Z}_5$ は
\S\ref{sec:acs_def} 定義~\ref{def:collapse} の
崩壊終端に対応するため、
通常ACS状態量との同時的な扱いは
同定義の帰属の規定と軌道記録の規約（\S\ref{subsec:blk_path_set}）に従う。

したがって、
$\mathsf{Z}$ 帯は
正式な状態空間分割ではなく、
異なる失敗関連現象を
比較可能に記述するための
補助的整理である。
\end{remark}

\begin{remark}[$\mathsf{Z}$ 帯は全順序ではない]
\label{rem:Z_not_total_order}
添字
$0,1,2,3,4,5$
は、
記述上の整理番号であり、
厳密な重症度尺度を定義しない。

したがって、
状態帯ラベルの添字に沿った重症度の全順序を置かず、
$\mathsf{Z}_k$ と記述される状態から
$\mathsf{Z}_{k+1}$ と記述される状態への
決定論的遷移も仮定しない。

例えば、
部分的機能不全
$\mathsf{Z}_2$
から閉塞
$\mathsf{Z}_4$
へ移ること、
全機能停止
$\mathsf{Z}_3$
から通常作動側へ回復すること、
あるいは
閉塞や全機能停止を事前に経由せず
崩壊へ到達することを、
状態帯の番号規約だけでは
肯定も否定もしない。

この点は、
後続の経路集合
$\mathfrak{P}_{\mathrm{blk}}$、
$\mathfrak{P}_{\mathrm{col}}$
の定義とは別問題である。
\end{remark}

\begin{remark}[全機能停止と閉塞の非同一性]
\label{rem:full_stop_not_blockage}
全機能停止
$\mathsf{Z}_3$
と閉塞
$\mathsf{Z}_4$
は、
異なる記述対象に関する条件である。

全機能停止は、
TIPA・PSL・ISL に関係する
機能作動の停止についての
記述的状態である。

これに対し閉塞は、
実現外部環境状態 $e_t$ のもとで
$A_{i,t}=\emptyset$
が成立するという
厳密な行動可行性条件である。

したがって本論文は、
$\mathsf{Z}_3$ と記述されることから
$\mathsf{Z}_4$ と記述されることを、
またその逆を、一般的な帰結として置かない。

両者が同時に記述される場合はありうるが、
その同時成立も
本小節の定義からは要求されない。
\end{remark}

\subsubsection{三原理の時刻別充足述語・整形式実現・崩壊との区別}

\begin{remark}[状態帯と三原理の時刻別充足述語が偽であることを同一視しない]
\label{rem:zones_not_principle_failure}
$\mathsf{Z}_1$、
$\mathsf{Z}_2$、
$\mathsf{Z}_3$
は、
機能の水準における低下または作動停止に関する
記述的ラベルである（備考~\ref{rem:s7_principle_levels}）。

これらのいずれかが記述されることから、
SM・EU・FR のいずれかの時刻別充足述語が
偽であるとはしない。

逆に、
SM・EU・FR のいずれかの時刻別充足述語が偽であることから、
特定の
$\mathsf{Z}_1$、
$\mathsf{Z}_2$、
$\mathsf{Z}_3$
へ分類されるとも仮定しない。

また、
閉塞
$\mathsf{Z}_4$
についても、
\S\ref{subsec:blockage}で明示したとおり、
SM の時刻別充足述語が偽であることとは同一ではない。

したがって、
\[
  \text{三原理の時刻別充足述語の真偽}
  \neq
  \text{$\mathsf{Z}$帯の記述}
\]
である。
\end{remark}

\begin{remark}[崩壊前状態帯と整形式実現の不成立を同一視しない]
\label{rem:precollapse_not_wf_failure}
本論文は、$\mathsf{Z}_1$ から
$\mathsf{Z}_4$ のいずれかに
記述されることだけから、
$\mathrm{WF}_{i,t}=0$
を一般的な帰結として置かない。

特に、
部分的機能不全、
全機能停止、
または閉塞が存在しても、
同一ACS個体としての構造的自己の
形式的実現がその時刻に
失われたとは限らない。

逆に本論文は、
$\mathrm{WF}_{i,t}=0$
であることから、
本小節のどの記述帯へ
分類されるかを一般的な帰結として置かない。

整形式実現と崩壊の正式な関係は、
後続の
\S\ref{subsec:collapse_def}
で射程を限定して扱う。
\end{remark}

\subsubsection{「崩壊前状態帯」という呼称}

\begin{remark}[崩壊前状態帯という呼称の用法]
\label{rem:precollapse_terminology}
本論文では、
\[
  \mathsf{Z}_1,\;
  \mathsf{Z}_2,\;
  \mathsf{Z}_3,\;
  \mathsf{Z}_4
\]
を総称して、
\emph{記述的崩壊前状態帯}
と呼ぶ。
以下では略して
\emph{崩壊前状態帯}
ともいう。
帯を付けない「崩壊前状態」は、
特定の状態帯への分類を意図しない
一般的な説明語として用いる。

この呼称に
$\mathsf{Z}_0$
は含めない。
$\mathsf{Z}_0$ は
通常作動側の比較基準であり、
障害または崩壊前状態を意味しない。

また、
「崩壊前」という語は、
これらの状態が
時間的・因果的に
崩壊の前駆段階として
必ず出現することを意味しない。

すなわち、本論文は、
$\mathsf{Z}_j$（$j=1,2,3,4$）
にあることから
\[
  \exists s>t:
  \mathrm{Col}_{i,s}=1
\]
を一般的な帰結として置かない。

同様に、
崩壊が成立するために
$\mathsf{Z}_1$ から
$\mathsf{Z}_4$ のいずれかを
事前に経由することも
必要条件とはしない。

したがって、
この呼称は
\[
  \text{崩壊とは区別されるが、
  崩壊周辺の障害・閉塞状態を記述する帯}
\]
という分類上の意味に限定される。
\end{remark}

\begin{remark}[状態帯と回復・崩壊・経路集合の関係]
\label{rem:zones_no_path_implication}
本論文は、特定の
$\mathsf{Z}_1$ から
$\mathsf{Z}_4$
に記述されることだけから、
外的回復可能性、
崩壊到達経路集合への帰属、
または崩壊への到達を
一般的な帰結として置かない。

閉塞経路集合への帰属については、
$\mathsf{Z}_4$ に記述される時刻が
軌道記録 $\sigma$（定義~\ref{def:trajectory_space}）の
ACS記録有効時刻域 $I_{\mathrm{ACS}}(\sigma)$（定義~\ref{def:acs_valid_time_domain}）に属するとき、
定義~\ref{def:blk_path_set} により
$\sigma\in\mathfrak{P}_{\mathrm{blk}}$ である。
$\mathsf{Z}_1$ から $\mathsf{Z}_3$ に記述されることだけからは、
この帰属を一般的な帰結として置かない。

外的回復可能性、閉塞経路集合および崩壊到達経路集合の正式な条件は、
それぞれ
\S\ref{subsec:external_recovery}、
\S\ref{subsec:blk_path_set}、
\S\ref{subsec:col_path_set}
で定める。
\end{remark}

%----------------------------------------------------------------------
\subsection{外的回復可能性と同一個体性}
\label{subsec:external_recovery}

本小節では、
$\mathsf{Z}_1$ から $\mathsf{Z}_4$ に記述される状態について、
ACS個体 $\Sigma_i$ の外部から与えられる作用により、
\emph{同一個体としての継続が回復可能であるか}
を記述するための補助的形式を導入する。

ここで扱う外的回復可能性は、
ACS内部の行動 $u\in U$ による自律的応答とは区別される。
また、
本小節は外的作用そのものの物理過程、
介入後の詳細な状態遷移、
同一個体性を有限観測から判定する一般アルゴリズム、
または社会的介入の一般理論を構成するものではない。

本小節の目的は、
\S\ref{subsec:collapse_def} で崩壊を本節の諸概念と関係づけるのに先立って、
\[
  \text{外的作用によって同一個体として回復しうる状態}
\]
と
\[
  \text{同一個体としての回復が本枠組み内で成立しない終端状態}
\]
を混同しないための記述上の境界を与えることである。

\subsubsection{許容外的作用集合}

\begin{definition}[許容外的作用集合（分類C）]
\label{def:admissible_external_action}
ACS個体 $\Sigma_i$ に対して、
Phase~1 の本節で
外的回復可能性の評価対象として考慮する
外的作用の集合を
\[
  \Upsilon_{\mathrm{adm},i}
\]
と表し、
これを
\emph{許容外的作用集合}
と呼ぶ。

各
\[
  \upsilon\in\Upsilon_{\mathrm{adm},i}
\]
は、
$\Sigma_i$ の外部から与えられる作用を表す。

$\Upsilon_{\mathrm{adm},i}$ は、
物理的・生物学的・社会的・工学的に
可能なすべての外的作用の集合ではなく、
本論文の Phase~1 において
回復可能性を記述するために
考慮対象へ含めた作用の集合である。

本節では、
$\Upsilon_{\mathrm{adm},i}$ の
位相構造、
測度構造、
代数構造、
作用時間、
強度、
複合作用規則、
または具体的実装を
一般形式としては定めない。
個々の回復可能性の判定に用いる $\Upsilon_{\mathrm{adm},i}$ は、
当該判定に先立って固定し、判定対象の経過または結果に応じて事後的に変更しない。
$\Upsilon_{\mathrm{adm},i}$ は、\S\ref{sec:tipa} の M1-b の識別状態空間 $\Xi$ とは
記号・型・役割を異にする。
\end{definition}

\begin{remark}[内部行動 $u$ と外的作用 $\upsilon$ の役割の区別]
\label{rem:external_action_type}
許容外的作用
\[
  \upsilon\in\Upsilon_{\mathrm{adm},i}
\]
と、
ACS個体 $\Sigma_i$ 自身が
行動空間 $U$ から選択する
$u\in U$
とは、
本節では異なる役割の変数として区別する。
両集合の非同一性を一般条件としては仮定せず、
$\upsilon$ と $u$ を暗黙に同一視しない。

特に、
外的作用 $\upsilon$ を
7タプルの遷移写像
$f(x_t,u,e_t)$
の第三引数または行動引数へ
暗黙に代入しない。

Phase~1 では、
$\upsilon$ が通常の
$f$ または $\mathrm{Update}$
へどのように作用するかを表す
普遍的な外的介入作用素は導入しない。

実際に起きた救助や介入は、許容外的作用 $\upsilon$ としてではなく、
実現した外部環境状態 $e_t$（背景成分または行動応答成分）を通じて
$f$ に入る状態変化として、軌道の上で記述する。
実現した軌道の上で同一個体としての再成立が判定される場合の扱いは、
\S\ref{sec:acs_def} の命題~\ref{prop:realized_reestablishment} による。
許容外的作用集合は、回復可能性の判定にだけ用いる。
\end{remark}

\begin{remark}[$\Upsilon_{\mathrm{adm}}$ はモデル境界である]
\label{rem:Upsilon_adm_scope}
$\Upsilon_{\mathrm{adm},i}$ は、
「宇宙に存在しうるすべての回復手段」
を意味しない。

したがって、
ある状態について
\[
  \forall\upsilon\in\Upsilon_{\mathrm{adm},i},
  \quad
  \text{回復が成立しない}
\]
と記述されたとしても、
それだけから
考えうるあらゆる外的作用に対する
\emph{絶対的な物理的不可能性}
を導いてはならない。

$\Upsilon_{\mathrm{adm},i}$ は、
Phase~1 の本枠組みで
回復可能性を論じるための
\emph{モデル上の許容作用域}
を定める。

この限定は、
\S\ref{subsec:collapse_def}
で崩壊の非回復性を読むときにも維持する。
\end{remark}

\begin{remark}[$\Upsilon_{\mathrm{adm}}$ の空集合可能性]
\label{rem:Upsilon_adm_empty}
本論文は一般に、
\[
  \Upsilon_{\mathrm{adm},i}\neq\emptyset
\]
を仮定しない。

したがって、
対象系または記述範囲によっては
\[
  \Upsilon_{\mathrm{adm},i}=\emptyset
\]
でありうる。

ただしその場合、
「本枠組み内で外的回復作用が指定されていない」ことは表しても、
本論文は、それだけから
$\mathrm{Col}_{i,t}=1$
を一般的な帰結として置かない。

すなわち、
許容外的作用集合の空性と
崩壊を同一視しない。
\end{remark}

\subsubsection{同一個体性の保存}

\begin{remark}[回復と別個体生成の区別]
\label{rem:same_identity}
本小節でいう
\emph{回復}は、
外的作用後も
同一ACS個体 $\Sigma_i$ が
その後続の主体として継続する場合に限る。

したがって、
外的作用によって
ACSとして作動可能な構造が形成されたとしても、
その主体が
元の $\Sigma_i$ ではなく
別個の ACS個体
$\Sigma_{i'}$（$i'\neq i$）
として扱われるべき場合、
それを
$\Sigma_i$ の回復とは数えない。

本論文の構造的自己
$S_{\mathrm{struct}}(\Sigma_i)$
は、
同一ACS個体としての形式的構造を
識別するための基準を与える。

ただしPhase~1では、
同一個体性を
単一の距離、
単一の連続量、
単一の記憶内容、
または単一の物理部位の保存へ
還元しない。
\end{remark}

\begin{remark}[同一個体性の判定を過剰形式化しない]
\label{rem:identity_scope}
同一個体性の保存には、
対象実装に応じて、
構造的自己、
内部モデル、
Update 構造、
自己--環境境界、
記憶の連続性その他が
関係しうる。

しかし本論文は、
それらのすべてを
同一個体性の普遍的必要十分条件として
Phase~1 内で固定しない。

したがって、
本小節は
「どの程度の部品交換まで同一個体か」、
「内部モデルの何パーセントが保持されれば同一か」、
「記憶の連続性がどの程度必要か」
といった実装依存の同一性判定問題へ
普遍的な閾値を与えない。

Phase~1 で必要なのは、
回復概念が
\[
  \text{同一個体の再継続}
\]
と
\[
  \text{別個体の生成}
\]
を概念上区別することである。
\end{remark}

\subsubsection{外的回復述語}

\begin{definition}[外的回復述語と外的回復可能性（分類C）]
\label{def:external_recoverability}
時刻
\[
  t\in\mathbb{N}_0
\]
において、
ACS個体 $\Sigma_i$ が
記述的崩壊前状態帯
\[
  \mathsf{Z}_1,\ldots,\mathsf{Z}_4
\]
のいずれかに記述される状態にあるとする。

このとき、
許容外的作用集合
$\Upsilon_{\mathrm{adm},i}$ 上の
決定論的二値補助写像
\[
  \mathrm{Recover}_{i,t}:
  \Upsilon_{\mathrm{adm},i}
  \longrightarrow
  \{0,1\}
\]
を
\emph{外的回復述語}
と呼ぶ。

各
\[
  \upsilon\in\Upsilon_{\mathrm{adm},i}
\]
について
\[
  \mathrm{Recover}_{i,t}(\upsilon)=1
\]
であるとは、
対象系ごとに当該判定に先立って固定された
Phase~1 の回復判定規則のもとで、
外的作用 $\upsilon$ を考慮したときに、
別個体の生成ではなく、
同一個体 $\Sigma_i$ としての
組織・結合関係（\S\ref{sec:acs_def} の定義~\ref{def:collapse} の (SC-1)）を保存または再成立させつつ、
通常ACSとしての後続作動へ
復帰しうると判定されることを表す。

一方、
\[
  \mathrm{Recover}_{i,t}(\upsilon)=0
\]
は、
その外的作用 $\upsilon$ について
上記の意味での回復可能性が
本枠組み内で成立すると
判定されないことを表す。

この回復判定規則は、
同一個体と別個体を区別する判定基準を含み、
判定対象の経過または結果に応じて
事後的に変更しない。
その同一個体の判定基準は、
\S\ref{sec:acs_def} の定義~\ref{def:collapse} および定義~\ref{def:functional_termination} の
判定規則に含まれる基準と同じものとする。
また、$\upsilon$ を考慮して判定する継続は、
(SC-2)・(FT-2) が量化する継続のうち $\upsilon$ を加えたものとし、
後続作動への復帰の判定は、(FT-1) の基礎的な機能のいずれかの
機能的活動が同一個体として再び成立することを含むものとする。

ACS個体 $\Sigma_i$ が
時刻 $t$ において
\emph{外的回復可能}であるとは、
\[
  \exists\upsilon\in\Upsilon_{\mathrm{adm},i}
  \quad\text{s.t.}\quad
  \mathrm{Recover}_{i,t}(\upsilon)=1
\]
が成立することをいう。
\end{definition}

\begin{remark}[外的回復述語は外的介入動学そのものではない]
\label{rem:recover_not_intervention_dynamics}
定義~\ref{def:external_recoverability} の
$\mathrm{Recover}_{i,t}$ は、
外的作用後の物理過程を
明示的に生成する
状態遷移写像ではない。

本論文は、
$(x_t,\upsilon) \longmapsto x'$
のような
普遍的な外的介入遷移写像を
Phase~1 では導入しない。

したがって、
$\mathrm{Recover}_{i,t}(\upsilon)$ は
回復作用の力学そのものではなく、
対象系ごとに定められた
外的回復可能性を記述するための
分類Cの補助述語である。

介入後の詳細な状態軌道、
介入強度、
複数作用の合成、
介入コスト、
および介入の到達可能性は、
本小節の共通形式には含めない。
\end{remark}

\begin{remark}[$\mathrm{Recover}=1$ は実際の回復発生を意味しない]
\label{rem:recoverability_not_actual_recovery}
\[
  \mathrm{Recover}_{i,t}(\upsilon)=1
\]
は、
作用 $\upsilon$ のもとで
同一個体としての回復が
\emph{可能であると判定される}ことを表す。

これは、
時刻 $t$ 以後に
実際に作用 $\upsilon$ が与えられること、
または実際の軌道上で
回復が発生することを意味しない。

したがって、
回復可能性、回復作用の実施および実際の回復到達を、
それぞれ可能性・実施・結果として区別する。

この区別により、
外的回復可能性を
将来の実現事象と同一視しない。
\end{remark}

\begin{remark}[$\mathrm{Recover}=0$ の解釈]
\label{rem:recover_zero_scope}
\[
  \mathrm{Recover}_{i,t}(\upsilon)=0
\]
は、
作用 $\upsilon$ による回復が
本枠組みの判定規則のもとで
成立すると認められないことを表す。

これは、
作用 $\upsilon$ による回復が
宇宙的・物理的に絶対不可能であることを
意味しない。

さらに、
\[
  \forall\upsilon\in\Upsilon_{\mathrm{adm},i},
  \quad
  \mathrm{Recover}_{i,t}(\upsilon)=0
\]
であっても、
それだけから
崩壊を定義しない。

崩壊には、
\S\ref{sec:acs_def} 定義~\ref{def:collapse} の
(SC-1)・(SC-2) の成立を要する。
\end{remark}

\subsubsection{記述的崩壊前状態帯との関係}

\begin{remark}[外的回復可能性の対象範囲]
\label{rem:recover_predicate_scope}
本小節の外的回復可能性は、
記述的崩壊前状態帯
\[
  \mathsf{Z}_1,\ldots,\mathsf{Z}_4
\]
にある状態について用いる。

したがって、
通常作動・基準帯
$\mathsf{Z}_0$
を
「回復対象」として扱わない。

$\mathsf{Z}_0$ の維持に
外部資源または外的支援が
必要となる場合はありうるが、
それは
障害状態からの
\emph{回復可能性}
とは別の問題である。

また、
崩壊帯
$\mathsf{Z}_5$
については、
本小節の回復可能性定義を
そのまま適用して
崩壊を否定するのではなく、
\S\ref{sec:acs_def} の崩壊の定義（定義~\ref{def:collapse}）に委ねる。
\end{remark}

\begin{remark}[状態帯番号と回復可能性の関係]
\label{rem:zone_not_recoverability_order}
本論文は、$\mathsf{Z}_1$、
$\mathsf{Z}_2$、
$\mathsf{Z}_3$、
$\mathsf{Z}_4$
という添字の大小から、
回復可能性の大小または
回復確率の大小を一般的な帰結として置かない。

例えば、
$\mathsf{Z}_4$
の閉塞状態が
ある許容外的作用によって
回復可能である一方、
ある
$\mathsf{Z}_2$
の部分的機能不全状態が
本枠組み内では
回復可能と判定されない場合を
状態帯の番号規約は排除しない。

したがって、
$\mathsf{Z}$ 帯は
回復難易度の尺度ではない。
\end{remark}

\begin{remark}[閉塞と外的回復可能性]
\label{rem:blockage_external_recovery}
閉塞
$A_{i,t}=\emptyset$
は、
$\mathsf{Z}_4$ に対応する
時刻局所的な一ステップ条件である。

閉塞の成立は、
外的回復可能性を否定しない。

すなわち、
ある時刻 $t$ に
$A_{i,t}=\emptyset$
であっても、
\[
  \exists\upsilon\in\Upsilon_{\mathrm{adm},i}
  \quad\text{s.t.}\quad
  \mathrm{Recover}_{i,t}(\upsilon)=1
\]
である場合を許容する。

この区別は、
\S\ref{subsec:blockage}
で述べた、
閉塞 $A_{i,t}=\emptyset$ から $\mathrm{Col}_{i,t}=1$ を
一般的な帰結として置かないという扱いの、
回復側の記述を与える。
\end{remark}

\subsubsection{WF・三原理の時刻別充足述語との関係}

\begin{remark}[外的回復可能性と整形式実現]
\label{rem:recovery_and_wf}
外的回復可能性は、
$\mathrm{WF}_{i,t}=0$
と同義ではない。

記述的崩壊前状態帯にある
$\Sigma_i$ が
時刻 $t$ において
整形式実現を保っている場合でも、
外的作用による機能回復が
有意味である場合はありうる。

逆に本論文は、
$\mathrm{WF}_{i,t}=0$
であることだけから、
外的回復可能性の有無を
一般的な帰結として置かない。

したがって、
$\mathrm{WF}_{i,t}$
と
$\mathrm{Recover}_{i,t}$
は異なる形式対象である。

崩壊の定義（\S\ref{sec:acs_def} 定義~\ref{def:collapse}）も、
整形式実現の不成立と
本枠組みに相対する非回復性を
区別している。
\end{remark}

\begin{remark}[外的回復可能性と三原理の時刻別充足述語]
\label{rem:recovery_not_principle_truth}
外的回復可能性は、
SM・EU・FR の
いずれの時刻別充足述語の条件でもない。

また、
\[
  \exists\upsilon\in\Upsilon_{\mathrm{adm},i}:
  \mathrm{Recover}_{i,t}(\upsilon)=1
\]
であることから、
SM・EU・FR の時刻別充足述語が
時刻 $t$ に真であるとは限らない。

逆に、
三原理の時刻別充足述語が時刻 $t$ に真であっても、
外的回復作用を必要とする状態でない場合や、
$\Upsilon_{\mathrm{adm},i}$ が
空である場合がありうる。

したがって、
\[
  \text{外的回復可能性}
  \neq
  \text{三原理の時刻別充足述語の判定}
\]
である。
\end{remark}

\subsubsection{崩壊定義への接続}

\begin{remark}[外的回復可能性と崩壊の依存方向]
\label{rem:recovery_collapse_distinction}
本小節は、
外的回復可能性を先に定義し、
崩壊そのものは定義しない。

したがって、
論理的な依存方向は
\[
  \text{許容外的作用域}
  \longrightarrow
  \text{外的回復可能性}
  \longrightarrow
  \text{崩壊定義における非回復性の参照}
\]
である。

崩壊の定義の条件 (SC-2)（\S\ref{sec:acs_def} 定義~\ref{def:collapse}）は
本小節の
$\Upsilon_{\mathrm{adm},i}$
を参照し、
\S\ref{subsec:collapse_def}
はその非回復性を外的回復可能性の射程とあわせて読む。

一方、
本小節の
$\mathrm{Recover}_{i,t}$
の定義に
$\mathrm{Col}_{i,t}$
の真理値を用いない。

これにより、
\[
  \text{回復可能性}
  \leftrightarrow
  \text{崩壊でないこと}
\]
という循環定義を避ける。
\end{remark}

\begin{remark}[回復可能なら崩壊ではない]
\label{rem:recoverable_excludes_collapse}
記述的崩壊前状態帯
$\mathsf{Z}_1$ から
$\mathsf{Z}_4$ にある時刻 $t$ について、
\[
  \exists\upsilon\in\Upsilon_{\mathrm{adm},i}
  \quad\text{s.t.}\quad
  \mathrm{Recover}_{i,t}(\upsilon)=1
\]
が成立するならば、
その時刻の状態は
崩壊状態ではない
（\S\ref{subsec:collapse_def} 命題~\ref{prop:recover_excludes_termination}）。

これは、
外的回復可能性から
新たな崩壊定義を導出するものではなく、
回復判定規則に課した (SC-2) との整合の要求
（定義~\ref{def:external_recoverability}）からの帰結である。

すなわち、
\[
  \text{本枠組みで同一個体として回復可能}
  \quad\Rightarrow\quad
  \text{崩壊として扱わない}
\]
という区別を維持する。
\end{remark}

\subsubsection{本小節の限定}

\begin{remark}[Phase~1 における外的回復可能性の限定]
\label{rem:external_recovery_phase1_scope}
本小節の形式化には、
以下の限定がある。

第一に、
$\Upsilon_{\mathrm{adm},i}$ は
全外的作用を網羅しない。

第二に、
$\mathrm{Recover}_{i,t}$ は
対象系ごとの判定規則を抽象化した
分類Cの補助述語であり、
介入後の状態遷移を
生成する写像ではない。

第三に、
同一個体性の普遍的な
必要十分条件を
Phase~1 では与えない。

第四に、
回復可能性の存在は
実際の回復作用の実施や
実際の回復到達を意味しない。

第五に、本論文は、
$\mathrm{Recover}=0$ または
$\Upsilon_{\mathrm{adm}}=\emptyset$
だけから崩壊を一般的な帰結として置かない。

したがって、
本小節は
崩壊と非崩壊を
有限観測から完全判定する理論ではなく、
崩壊の定義の (SC-2) が参照する
\emph{回復可能性側の概念境界}
を与えるものである。
\end{remark}

%----------------------------------------------------------------------
\subsection{崩壊：本枠組みにおける吸収的終端状態}
\label{subsec:collapse_def}

本小節では、
\S\ref{sec:acs_def} の定義~\ref{def:functional_termination} および
定義~\ref{def:collapse} で定義した
ACS個体 $\Sigma_i$ の機能的終端と崩壊を、
本節の諸概念と関係づけて整理する。

\subsubsection{崩壊指標}

機能的終端判定時刻域 $\mathbb{T}_{\mathrm{FT},i}$、
機能的終端状態指標 $\mathrm{FTS}_{i,t}$、その条件 (FT-1)・(FT-2)、
機能的終端指標 $\mathrm{FT}_{i,t}$、
崩壊状態指標 $\mathrm{CS}_{i,t}$、その条件 (SC-1)・(SC-2)、および崩壊指標
$\mathrm{Col}_{i,t}\in\{0,1\}$ は、
\S\ref{sec:acs_def} の定義~\ref{def:functional_termination} および
定義~\ref{def:collapse} で与えた。
吸収性 (C-2) と崩壊状態成立時刻の存在 (C-3) は同節の命題~\ref{prop:collapse_basic}、
崩壊が成立している時刻までに機能的終端状態が成立済みであることは同節の命題~\ref{prop:termination_order} である。
本節でいう崩壊の非回復性は同定義の条件 (SC-2) を指し、
機能的終端の非回復性は条件 (FT-2) を指す。
本小節ではこれらを再定義せず、
条件 (FT-2)・(SC-2) が参照する許容外的作用集合 $\Upsilon_{\mathrm{adm},i}$ の内容を
\S\ref{subsec:external_recovery}（定義~\ref{def:admissible_external_action}）に従って読む。

\begin{remark}[崩壊指標の値 $0$ の意味]
\label{rem:collapse_zero_meaning}
$\mathrm{Col}_{i,t}=0$
は、
時刻 $t$ までのいずれの時刻においても、
定義~\ref{def:collapse} の意味での
崩壊状態が成立していないことを表す。

本論文は、これから
$\mathrm{WF}_{i,t}=1$、
$A_{i,t}\neq\emptyset$、
$x_t\in K$、
または
SM・EU・FR の時刻別充足述語が真であることを
一般的な帰結として置かない。

したがって、
$\mathrm{Col}_{i,t}=0$
を「健全」「可行」「回復済み」
などの一般的な正常性指標として
使用しない。
\end{remark}

\subsubsection{整形式実現との関係}

\begin{remark}[整形式実現の不成立と崩壊の関係]
\label{rem:wf_failure_not_collapse}
機能的終端指標・崩壊指標と整形式実現指標の関係は、
\S\ref{sec:acs_def} の備考~\ref{rem:collapse_wf_relation} に従う。
\end{remark}

\begin{remark}[崩壊と構造的自己]
\label{rem:collapse_structural_self}
本論文の構造的自己
$S_{\mathrm{struct}}(\Sigma_i)$
は、
同一ACS個体としての
形式的構造を識別する基準である。

崩壊が成立した時刻以後の
$\mathrm{WF}_{i,t}=0$
（命題~\ref{prop:termination_order}）は、
その構造的自己が
時刻 $t$ において
整形式に実現されていないことを表す。
その逆、すなわち $\mathrm{WF}_{i,t}=0$ から崩壊を、本論文は一般的な帰結として置かない
（備考~\ref{rem:collapse_wf_relation}）。

ただし、
構造的自己という
形式的記述対象それ自体を
数学的対象として消去するわけではない。

したがって、
崩壊は
「7タプルの記述が論文上存在しなくなること」
ではなく、
その同一ACS個体としての
通常の実現・継続が
本枠組み内で終端化したことを表す。
\end{remark}

\subsubsection{非回復性の射程}

\begin{remark}[崩壊の非回復性は本枠組みに相対する]
\label{rem:collapse_relative_irrecoverability}
条件 (SC-2) の非回復性、および機能的終端の条件 (FT-2) の非回復性は、
$\Upsilon_{\mathrm{adm},i}$
が定める
許容外的作用域と、
実現した外生背景成分の列のもとでの
通常の自律的更新および $g_O$ による継続の記述範囲に
相対する（\S\ref{sec:acs_def} 定義~\ref{def:collapse} の (SC-2)）。

したがって、
$\mathrm{Col}_{i,t}=1$
から、
考えうるすべての物理作用、
未知の技術、
Phase~1 の外側にある操作、
または本論文がモデル化していない
介入によっても
絶対に同一個体として
回復しえない、
という主張を本論文は置かない。

本論文が主張するのは、
\[
  \text{本枠組みの許容継続の範囲内で非回復}
\]
という限定された終端性である。
\end{remark}

\begin{remark}[$\Upsilon_{\mathrm{adm}}=\emptyset$ と崩壊の関係]
\label{rem:empty_Xi_not_collapse}
\S\ref{subsec:external_recovery} で述べたとおり、
本論文は
\[
  \Upsilon_{\mathrm{adm},i}=\emptyset
\]
を許容する。

しかし本論文は、
許容外的作用が
一つも指定されていないことだけから、
$\mathrm{Col}_{i,t}=1$
を一般的な帰結として置かない。

崩壊には、
少なくとも
(SC-1) 構造的同一性を成立させる組織関係の喪失と、
(SC-2) 本枠組みに相対する
同一個体としての非回復性が
ある時刻に併せて成立することを要する。

したがって、
外的作用域の空性と
崩壊は別の条件である。
\end{remark}

\begin{remark}[$\mathrm{Recover}=0$ の全称と崩壊の関係]
\label{rem:all_recover_zero_not_collapse}
記述的崩壊前状態帯において、
仮に
\[
  \forall\upsilon\in\Upsilon_{\mathrm{adm},i},
  \quad
  \mathrm{Recover}_{i,t}(\upsilon)=0
\]
であったとしても、
本論文は、その事実だけから
崩壊を一般的な帰結として置かない。

第一に、
$\mathrm{Recover}_{i,t}$
は
\S\ref{subsec:external_recovery}
で定義された
回復可能性側の分類C補助述語であり、
崩壊そのものの判定指標ではない。

第二に、
崩壊には
(SC-1) の組織関係の喪失が必要である。

第三に、
(SC-2) は
同一個体としての
許容継続全体に関する
終端条件であり、
単一時刻における
$\mathrm{Recover}_{i,t}$
の値だけへ還元しない。

したがって、
\[
  \text{外的回復可能性が確認されない}
  \neq
  \text{崩壊}
\]
である。
\end{remark}

\subsubsection{外的回復可能性との整合}

\begin{proposition}[外的回復可能性と終端状態指標]
\label{prop:recover_excludes_termination}
時刻 $t$ について $\mathrm{Recover}_{i,t}$ が定義され、
ある $\upsilon\in\Upsilon_{\mathrm{adm},i}$ について
$\mathrm{Recover}_{i,t}(\upsilon)=1$ ならば、
$\mathrm{CS}_{i,t}=0$ である。
さらに $t\in\mathbb{T}_{\mathrm{FT},i}$ ならば $\mathrm{FTS}_{i,t}=0$ である。
\end{proposition}

\begin{proof}
定義~\ref{def:external_recoverability} により、$\upsilon$ を加えた継続のうちの一つの上で、
(SC-2)・(FT-2) と同じ同一個体の判定基準のもとで、
(SC-1) の組織・結合関係が保存されるか再び成立し、
かつ (FT-1) の基礎的な機能のいずれかの機能的活動が同一個体として再び成立すると判定される。
保存される場合は時刻 $t$ に (SC-1) が成立せず、
再び成立する場合は、(SC-2) が量化する継続の上で再成立が認められるので (SC-2) が成立しない。
いずれの場合も $\mathrm{CS}_{i,t}=0$ である。
同様に、その継続は (FT-2) が量化する継続でもあり、その上で基礎的な機能の活動が再び成立するから、
$t\in\mathbb{T}_{\mathrm{FT},i}$ ならば (FT-2) が成立せず、$\mathrm{FTS}_{i,t}=0$ である。
\end{proof}

\begin{remark}[回復可能な状態は崩壊として扱わない]
\label{rem:C2_and_recoverability}
\S\ref{subsec:precollapse_states} の
記述的崩壊前状態帯
$\mathsf{Z}_1$ から
$\mathsf{Z}_4$
にある時刻 $t$ について、
\[
  \exists\upsilon\in\Upsilon_{\mathrm{adm},i}
  \quad\text{s.t.}\quad
  \mathrm{Recover}_{i,t}(\upsilon)=1
\]
が成立するならば、
その状態を
時刻 $t$ の崩壊状態として扱わない。

すなわち、
本枠組みの範囲で
同一個体としての回復が
可能であると判定される状態は、
命題~\ref{prop:recover_excludes_termination} により崩壊状態ではなく、
時刻が機能的終端判定時刻域に属するならば機能的終端状態でもない。

この関係は、
\S\ref{subsec:external_recovery}
で確定した依存方向
\[
  \text{許容外的作用域}
  \longrightarrow
  \text{外的回復可能性}
  \longrightarrow
  \text{崩壊定義の非回復性}
\]
と整合する。

一方、
外的回復可能性の定義に
$\mathrm{Col}_{i,t}$
を用いていないため、
循環定義は生じない。
\end{remark}

\subsubsection{閉塞・状態帯・三原理の時刻別充足述語との区別}

\begin{remark}[閉塞と崩壊の帰結関係の限定]
\label{rem:collapse_non_implication}
閉塞
$A_{i,t}=\emptyset$
は、
実現外部環境状態 $e_t$ のもとで
一ステップ先の可行性を保持する
行動が存在しないことを表す。

これは、
機能的終端の条件
(FT-1)・(FT-2)
とも、崩壊状態の条件
(SC-1)・(SC-2)
とも異なる。

したがって本論文は、閉塞から崩壊を一般的な帰結として置かない
（\S\ref{sec:acs_def} の備考~\ref{rem:blockage_not_collapse_sec2}）。

特に、
閉塞状態
$\mathsf{Z}_4$
にあっても、
外的回復可能性が存在する場合、
その状態は崩壊ではない。

また、
崩壊の成立に
閉塞の事前成立を
必要条件として課さない。
\end{remark}

\begin{remark}[全機能停止と崩壊の区別]
\label{rem:full_stop_not_collapse}
\S\ref{subsec:precollapse_states} の
全機能停止帯
$\mathsf{Z}_3$
は、
TIPA・PSL・ISL に関係する
標準機能の作動停止を
記述する補助的状態帯である。

全機能停止だけから、
(FT-1)・(FT-2) が
ともに成立するとも、
(SC-1)・(SC-2) が
ともに成立するとも限らない。
本論文は、全機能停止帯 $\mathsf{Z}_3$ と記述されることだけから、
機能的終端 $\mathrm{FT}_{i,t}=1$（定義~\ref{def:functional_termination}）も
一般的な帰結として置かない。
$\mathsf{Z}_3$ は標準機能形式（分類B）の作動に関する記述であり、
(FT-1) は原始的な機能の各区分に関して個体が有する基礎的な機能と
基盤機能の停止に関する条件であって、両者は機能の記述の水準が異なる。
本節は、両者の間にいずれの向きの含意も置かない。

したがって本論文は、
全機能停止帯 $\mathsf{Z}_3$ と記述されることだけから、
$\mathrm{Col}_{i,t}=1$ を一般的な帰結として置かない。

特に、
全機能停止状態であっても、
同一個体として
回復可能である場合を
本論文は排除しない。
\end{remark}

\begin{remark}[三原理の時刻別充足述語が偽であることと崩壊の関係]
\label{rem:principle_failure_not_collapse}
本論文は、SM・EU・FR の時刻別充足述語の
いずれか、
または複数が
時刻 $t$ に偽であっても、
それだけから
$\mathrm{Col}_{i,t}=1$
を一般的な帰結として置かない。

三原理は
\S\ref{sec:three_principles} で定める
制約軸であり、
その時刻別充足述語は存続中のACSの状態と環境についての条件である。
これに対し崩壊は、
\S\ref{sec:acs_def} で定義した構造的個体の終端である。

したがって、
\[
  \text{三原理の時刻別充足述語が偽}
  \neq
  \text{崩壊判定}
\]
である。

また、
T-STAB の不成立、
関係品質 $Q<0$、
または
$\chi_{\mathrm{EI},t}=0$
だけからも、
本論文は崩壊を一般的な帰結として置かない。
\end{remark}

\subsubsection{崩壊前状態帯を経由しない崩壊}

\begin{remark}[崩壊前状態帯の事前経由を必要条件としない]
\label{rem:collapse_without_precollapse}
定義~\ref{def:collapse} は、
崩壊
$\mathrm{Col}_{i,t}=1$
の成立以前に、
\[
  \mathsf{Z}_1,
  \mathsf{Z}_2,
  \mathsf{Z}_3,
  \mathsf{Z}_4
\]
のいずれかが
観測または記述されていることを
必要条件としない。

したがって、
閉塞、
部分的機能不全、
全機能停止等を
明示的に経由せず、
本節の意味で
崩壊終端へ到達する軌道を
Phase~1 の一般論から排除しない。

逆に、
$\mathsf{Z}_1$ から
$\mathsf{Z}_4$
にある軌道が
必ず崩壊へ到達するとも
主張しない。

この区別は、
後続の
\S\ref{subsec:col_path_set}
および
\S\ref{subsec:inflow_paths}
で維持する。
\end{remark}

\subsubsection{崩壊の定義と判定可能性}

\begin{remark}[崩壊の定義と有限観測による判定可能性は異なる]
\label{rem:collapse_definition_observability}
命題~\ref{prop:collapse_basic} の
(C-2) 吸収性および
定義~\ref{def:collapse} の
(SC-2) 非回復性は、
現在時刻だけではなく
将来の許容継続に関係する。

したがって、
ある有限時刻の観測だけから
一般的・決定論的に
$\mathrm{Col}_{i,t}=1$
を判定する
普遍的手続きを、
Phase~1 は与えない。

これは、
崩壊の形式的定義が
不明確であることを意味しない。

本論文は、
\[
  \text{理論上の終端状態の定義}
\]
と
\[
  \text{有限観測からそれを認識する判定法}
\]
を区別する。

後者に必要となる
実装依存の十分条件、
診断指標、
観測窓、
誤判定率、
確率的推定等は、
Phase~2 以降の課題とする（\S\ref{subsec:limitations_list} の L9、問題~\ref{prob:collapse_observability}）。
\end{remark}

\begin{remark}[崩壊指標はACS内部変数である必要はない]
\label{rem:collapse_indicator_description_level}
$\mathrm{Col}_{i,t}$ は、
本論文が
崩壊状態を記述するために導入する
理論上の指標である。

したがって、
ACS個体 $\Sigma_i$ 自身が
内部に
$\mathrm{Col}_{i,t}$
という同一の二値変数を
保持・計算・認識することを
要求しない。

また、
M2 その他の機構が
真の
$\mathrm{Col}_{i,t}$
を直接読み取ることも
Phase~1 の要件ではない。

この意味で、
崩壊指標は
記述・診断側の形式表現であり、
特定の実装変数と
一対一対応させない。
\end{remark}

\subsubsection{崩壊後の通常ACS状態量の扱い}

\begin{remark}[崩壊後に通常ACS動学を要求しない]
\label{rem:collapse_internal_state}
$\mathrm{Col}_{i,t}=1$
が成立した後は、
条件 (C-2) により、
終端記録として
崩壊指標を保持する。

一方、
崩壊後の時刻について、
$\Sigma_i$ を
通常ACSとして
\[
  x_s\in X,
  \qquad
  e_s\in E_{\mathrm{ext}},
  \qquad
  A_{i,s},
  \qquad
  \mathrm{Desc}_s
\]
によって継続記述することを
一般には要求しない。

これは、
崩壊後の物理的対象が
文字どおり消滅することを
意味しない。

意味するのは、
同一ACS個体 $\Sigma_i$ の
通常の内部動学として
その後を記述する責務を、
Phase~1 のACS軌道記述が
負わないということである。

この終端規約は、
\S\ref{subsec:blk_path_set}
で導入する軌道記録空間において
正式に反映する。
\end{remark}

%----------------------------------------------------------------------
\subsection{軌道記録空間と閉塞経路集合}
\label{subsec:blk_path_set}

\S\ref{subsec:collapse_def} では、
崩壊
$\mathrm{Col}_{i,t}=1$
を、本枠組みにおける
吸収的・非回復的終端分類として整理し、
最初の崩壊成立後には
通常ACS状態量
\[
  x_t,\qquad
  e_t,\qquad
  A_{i,t},\qquad
  \mathrm{Desc}_t
\]
を継続して要求しないことを確認した。

本小節では、
この終端規約と整合する
\emph{軌道記録空間}
$\mathfrak{P}$ を導入し、
その部分集合として
有限時刻内に閉塞へ到達する軌道記録の集合
$\mathfrak{P}_{\mathrm{blk}}$
を定義する。

ここでの軌道記録は、
崩壊前の通常ACS状態の時間発展と、
全離散時刻上の崩壊終端記録とを
明示的に分離する。

\subsubsection{ACS記録有効時刻域}

\begin{definition}[ACS記録有効時刻域（分類C）]
\label{def:acs_valid_time_domain}
軌道記録 $\sigma$ に付随する
崩壊指標列
\[
  \bigl(
    \mathrm{Col}_{i,t}
  \bigr)_{t\in\mathbb{N}_0}
\]
が、
命題~\ref{prop:collapse_basic} の
吸収性条件 (C-2) を満たすとする。

このとき、
$\sigma$ において
通常ACS状態量を記録する時刻域を
\[
  I_{\mathrm{ACS}}(\sigma)
  :=
  \bigl\{
    t\in\mathbb{N}_0
    \bigm|
    \mathrm{Col}_{i,t}=0
  \bigr\}
\]
と定義し、
これを
\emph{ACS記録有効時刻域}
と呼ぶ。

条件 (C-2) により、
$I_{\mathrm{ACS}}(\sigma)$ は
次のいずれかの形をとる。

\begin{enumerate}
  \item
  崩壊が一度も成立しない場合：
  \[
    I_{\mathrm{ACS}}(\sigma)
    =
    \mathbb{N}_0.
  \]

  \item
  ある最小時刻 $T\in\mathbb{N}_0$、$T\ge1$ で
  \[
    \mathrm{Col}_{i,T}=1
  \]
  が初めて成立する場合：
  \[
    I_{\mathrm{ACS}}(\sigma)
    =
    \{0,1,\ldots,T-1\}.
  \]

  \item
  初期時刻ですでに
  \[
    \mathrm{Col}_{i,0}=1
  \]
  である場合：
  \[
    I_{\mathrm{ACS}}(\sigma)
    =
    \emptyset.
  \]
\end{enumerate}

したがって、
最初の崩壊成立時刻そのものは
ACS記録有効時刻域に含めない。
\end{definition}

\begin{remark}[有効時刻域は崩壊前の通常ACS記述域である]
\label{rem:acs_valid_time_domain_scope}
$I_{\mathrm{ACS}}(\sigma)$ は、
同一個体 $\Sigma_i$ について
\S\ref{sec:acs_def} および \S\ref{sec:extended_dynamics} の状態量を
軌道記録の対象として扱う時刻域である。

$I_{\mathrm{ACS}}(\sigma)$ は
$\mathrm{Col}_{i,t}=0$ によって定まるから、
機能的終端後の時刻
（$\mathrm{FT}_{i,t}=1$ かつ $\mathrm{Col}_{i,t}=0$、
定義~\ref{def:functional_termination}）を含みうる。
その時刻の記録は、
構造的個体としての追跡の記録であり、
$\Sigma_i$ を現に作動するACSの実現として扱うものではない
（備考~\ref{rem:post_collapse_record}）。
現に作動するACSとしての機能の記述は、
$\mathrm{FT}_{i,t}=0$ の時刻に限る。

ここで重要なのは、
本論文が
$\mathrm{Col}_{i,t}=0$
から
$\mathrm{WF}_{i,t}=1$
を一般的な帰結として置かないことである。
したがって、
$I_{\mathrm{ACS}}(\sigma)$ は、
崩壊には至っていないが
一時的に
$\mathrm{WF}_{i,t}=0$
である時刻を含みうる。

それでも、
$t\in I_{\mathrm{ACS}}(\sigma)$
に対しては、
崩壊終端後とは区別して、
少なくとも
\[
  x_t\in X,
  \qquad
  e_t\in E_{\mathrm{ext}}
\]
という通常の状態・外部環境状態を記録できる。

一方、
$t\notin I_{\mathrm{ACS}}(\sigma)$
に対しては、
本節は
通常ACS状態量としての
$x_t$、
$e_t$、
$A_{i,t}$、
または
$\mathrm{Desc}_t$
を要求しない。

これは、
崩壊後の物理対象の存在を否定するものではなく、
同一ACS個体としての
通常動学記述を終了するという
Phase~1 の終端規約である。
\end{remark}

\begin{remark}[崩壊成立時刻に閉塞を自動付与しない]
\label{rem:no_blockage_at_collapse_time_by_default}
最初の崩壊成立時刻 $T$ は、
定義~\ref{def:acs_valid_time_domain} により
\[
  T\notin I_{\mathrm{ACS}}(\sigma)
\]
である。

したがって本節は、
崩壊成立時刻 $T$ において
通常ACS状態量としての
$A_{i,T}$
を要求しない。

このため、
$\mathrm{Col}_{i,T}=1$
から、
同じ時刻に
$A_{i,T}=\emptyset$
を自動的に付与しない。

崩壊以前のACS記録有効時刻域内で
閉塞が成立していたか否かは、
後述の
$\mathfrak{P}_{\mathrm{blk}}$
によって別途記録する。
\end{remark}

\subsubsection{軌道記録空間}

\begin{definition}[軌道記録空間（分類C）]
\label{def:trajectory_space}
本節では、
\emph{軌道記録空間}
$\mathfrak{P}$ を、
単一ACS個体 $\Sigma_i$ について、
ACS記録有効時刻域上の状態記録と
全離散時刻上の崩壊指標列を
組として保持する記録の集合として導入する。

各
\[
  \sigma\in\mathfrak{P}
\]
は、
少なくとも次の二部分を含む。

第一に、
全離散時刻上の崩壊指標列
\[
  \bigl(
    \mathrm{Col}_{i,t}
  \bigr)_{t\in\mathbb{N}_0},
\]
ただしこの列は
命題~\ref{prop:collapse_basic} の
(C-2) 吸収性を満たす。

第二に、
その崩壊指標列から定まる
ACS記録有効時刻域
$I_{\mathrm{ACS}}(\sigma)$
上の記録
\[
  \bigl(
    x_t,\,
    e_t,\,
    \mathrm{Desc}_t
  \bigr)_{t\in I_{\mathrm{ACS}}(\sigma)},
\]
ここで
\[
  x_t\in X,
  \qquad
  e_t\in E_{\mathrm{ext}}.
\]

したがって、
軌道記録を
\[
  \sigma
  =
  \left(
    \bigl(
      x_t,e_t,\mathrm{Desc}_t
    \bigr)_{t\in I_{\mathrm{ACS}}(\sigma)},
    \bigl(
      \mathrm{Col}_{i,t}
    \bigr)_{t\in\mathbb{N}_0}
  \right)
\]
と表す。

ここで
$\mathrm{Desc}_t$ は、
\S\ref{sec:extended_dynamics}
で扱われる時刻依存量を整理して参照するための
\S\ref{sec:extended_dynamics} 由来の組である。
その成分のうち予測精度係数 $\pi_{\tau,t}$ と実効的時間地平 $\widehat\tau_t$ は
時刻 $t$ から $t+\tau$ までの記録を用いるので、
$t+\tau\in I_{\mathrm{ACS}}(\sigma)$ のときに限り記録し、
それ以外の時刻では定義しない（値 $0$ を割り当てない）。
したがって、$t+\tau\notin I_{\mathrm{ACS}}(\sigma)$ である $t\in I_{\mathrm{ACS}}(\sigma)$ の
$\mathrm{Desc}_t$ は、この二成分を未定義とする部分的な組として記録する。
崩壊後の形式的な作動記録（\S\ref{sec:extended_dynamics} 命題~\ref{prop:s5_continuation}）を
これらの評価に用いない。
また、本節は
$\sigma$ を
$\mathrm{Desc}_t$ の時系列だけに
同一視しない。

各
$t\in I_{\mathrm{ACS}}(\sigma)$
について、
可行行動集合
\[
  A_{i,t}
  =
  A(x_t,e_t)
  =
  \bigl\{
    u\in U
    \bigm|
    f(x_t,u,e_t)\in K
  \bigr\}
\]
が、
定義~\ref{def:blockage_sec2}
および
\S\ref{sec:acs_def}
の可行行動集合の定義に従って定まる。

一方、
$t\notin I_{\mathrm{ACS}}(\sigma)$
については、
通常ACS状態量としての
$A_{i,t}$ を要求しない。

記号
$\mathfrak{P}$ は、
持続評価関数
$\Pi$
とは無関係の対象である。
\end{definition}

\begin{remark}[軌道記録空間の力学的意味]
\label{rem:trajectory_space_scope}
$\mathfrak{P}$ は、
任意の数列を無条件に
「ACS軌道」と呼ぶための集合ではない。

本論文で
$\sigma\in\mathfrak{P}$
をACS個体 $\Sigma_i$ の軌道記録として扱うとき、
ACS記録有効時刻域上の記録は、
\S\ref{sec:acs_def}
および
\S\ref{sec:extended_dynamics}
で定める
対象系の決定論的動学と
整合する実現記録または
決定論的シミュレーション記録であることを前提とする。

ただし本小節では、
その整合性を判定するための
新しい軌道生成作用素や
新しい状態方程式を導入しない。

$\mathfrak{P}$ は、
後続の経路集合を記述するための
分類Cの台集合である。
\end{remark}

\begin{remark}[軌道記録と実現記録・シミュレーション記録]
\label{rem:trajectory_record_origin}
$\sigma\in\mathfrak{P}$
は、
実際に観測された
実現記録だけに限定しない。

対象系の決定論的モデルに基づいて
構成された
決定論的シミュレーション記録、
すなわち初期値と環境実現（背景成分の列と $g_O$）を固定して
\S\ref{sec:extended_dynamics} の標準作動記録として生成された記録も、
当該モデルの動学と整合する限り、
軌道記録として扱いうる。
行動の列を記述者が与えて生成した記録は、
$\Sigma_i$ 自身の作動による記録ではないので、
本節では $\Sigma_i$ の軌道記録として扱わない。

ただし、
この記録由来の違いから
$\mathfrak{P}$ 上に
確率測度や尤度を
Phase~1 で導入しない。
\end{remark}

\begin{remark}[$\mathfrak{P}$ に位相・距離・確率を要求しない]
\label{rem:trajectory_space_no_extra_structure}
本論文の Phase~1 では、
$\mathfrak{P}$ に
位相、
距離、
測度、
確率測度、
または順序構造を
一般には付与しない。

したがって、
後続節において
「経路集合への距離」、
「崩壊軌道への近さ」、
「閉塞確率」
などを、
追加定義なしに用いない。

本節が必要とするのは、
軌道記録の集合と
その集合論的部分集合だけである。
\end{remark}

\subsubsection{閉塞経路集合}

\begin{definition}[閉塞経路集合（分類C）]
\label{def:blk_path_set}
\emph{閉塞経路集合}
$\mathfrak{P}_{\mathrm{blk}}$
を、
軌道記録空間
$\mathfrak{P}$
の部分集合として
\[
  \mathfrak{P}_{\mathrm{blk}}
  :=
  \left\{
    \sigma\in\mathfrak{P}
    \;\middle|\;
    \exists T\in I_{\mathrm{ACS}}(\sigma)
    \text{ s.t. }
    A_{i,T}=\emptyset
  \right\}
\]
と定義する。

すなわち、
$\mathfrak{P}_{\mathrm{blk}}$
は、
通常ACS状態量が定義されている
有効時刻域内の
ある有限時刻において、
閉塞
\[
  A_{i,T}=\emptyset
\]
が成立する軌道記録の全体である。

したがって、
\[
  \mathfrak{P}_{\mathrm{blk}}
  \subseteq
  \mathfrak{P}.
\]

下付き
$\mathrm{blk}$
は
blockage
を表す。
\end{definition}

\begin{remark}[閉塞経路集合は閉塞の持続を要求しない]
\label{rem:blk_path_no_persistence}
\[
  \sigma\in\mathfrak{P}_{\mathrm{blk}}
\]
であるために必要なのは、
ACS記録有効時刻域内の
少なくとも一つの時刻で
$A_{i,T}=\emptyset$
が成立することである。

したがって、
その後の時刻にも
閉塞が持続することを
定義上要求しない。

特に、
閉塞後に
同一個体として
$A_{i,s}\neq\emptyset$
となる状態へ回復する軌道も、
一度閉塞へ到達している限り
$\mathfrak{P}_{\mathrm{blk}}$
に属する。
\end{remark}

\begin{remark}[崩壊と閉塞経路集合への帰属の関係]
\label{rem:collapse_not_blk_membership}
本論文は、崩壊指標
$\mathrm{Col}_{i,T}=1$
がある有限時刻に成立したことだけから、
\[
  \sigma\in\mathfrak{P}_{\mathrm{blk}}
\]
を一般的な帰結として置かない。

なぜなら、
$\mathfrak{P}_{\mathrm{blk}}$
への帰属には、
崩壊以前の
ACS記録有効時刻域内で
実際に
$A_{i,t}=\emptyset$
が成立したことを要するからである。

これは、
\S\ref{subsec:collapse_def}において
崩壊に閉塞の事前経由を必要条件として課していないことと
整合する。
\end{remark}

\subsubsection{閉塞到達時刻}

\begin{definition}[閉塞到達時刻（分類C）]
\label{def:blk_reaching_time}
軌道記録
\[
  \sigma\in\mathfrak{P}_{\mathrm{blk}}
\]
に対して、
\[
  T_{\mathrm{blk}}(\sigma)
  :=
  \min
  \left\{
    T\in I_{\mathrm{ACS}}(\sigma)
    \;\middle|\;
    A_{i,T}=\emptyset
  \right\}
\]
を
$\sigma$ の
\emph{閉塞到達時刻}
と呼ぶ。

$\sigma\in\mathfrak{P}_{\mathrm{blk}}$
であれば、
定義~\ref{def:blk_path_set} により
上記集合は非空である。

また、
\[
  I_{\mathrm{ACS}}(\sigma)
  \subseteq
  \mathbb{N}_0
\]
であるため、
自然数の整列性により
最小元が存在する。

したがって、
$T_{\mathrm{blk}}(\sigma)$
は有限の自然数として定まる。
\end{definition}

\begin{remark}[閉塞到達時刻は最初の閉塞時刻にすぎない]
\label{rem:blk_reaching_time}
$T_{\mathrm{blk}}(\sigma)$ は、
軌道記録 $\sigma$ において
閉塞が最初に成立する時刻を表す。

これは、
$t>T_{\mathrm{blk}}(\sigma)$
においても
閉塞が持続すること、
崩壊へ到達すること、
または外的回復が生じることを
意味しない。

したがって、
$T_{\mathrm{blk}}$
は
終端時刻ではなく、
最初の閉塞成立時刻を記録する
分類C量である。
\end{remark}

\subsubsection{閉塞経路集合と他概念の区別}

\begin{remark}[核外性と閉塞経路集合への帰属の関係]
\label{rem:viab_outside_not_blk_path}
\S\ref{subsec:blockage}で述べたとおり、本論文は、
$x_t\notin\mathrm{Viab}_f(K)$
だけから
$A_{i,t}=\emptyset$
を一般的な帰結として置かない。

したがって本論文は、
ある時刻で
$x_t\notin\mathrm{Viab}_f(K)$
であることだけから、
\[
  \sigma\in\mathfrak{P}_{\mathrm{blk}}
\]
を一般的な帰結として置かない。

同様に、
ロバスト可行性核の外部にあることを
「閉塞経路集合への流入」
と呼ぶことも、
追加条件なしには行わない。
\end{remark}

\begin{remark}[SM の時刻別充足述語が偽であること・整形式実現の不成立と閉塞経路集合への帰属の関係]
\label{rem:sm_wf_not_blk_path}
本論文は、時刻 $t$ における
SM の時刻別充足述語が偽であること
$\Pi(x_t)<0$
または
整形式実現の不成立
$\mathrm{WF}_{i,t}=0$
だけから、
\[
  \sigma\in\mathfrak{P}_{\mathrm{blk}}
\]
を一般的な帰結として置かない。

$\mathfrak{P}_{\mathrm{blk}}$
への帰属は、
ACS記録有効時刻域内で
実際に
$A_{i,T}=\emptyset$
となる時刻が存在することによってのみ
定義される。
\end{remark}

\begin{remark}[閉塞経路集合と崩壊到達経路集合は別対象である]
\label{rem:blk_vs_col}
本小節の
$\mathfrak{P}_{\mathrm{blk}}$
は、
有限時間内に
閉塞へ到達した軌道記録の集合である。

後続の
\S\ref{subsec:col_path_set}
で定義する
崩壊到達経路集合
$\mathfrak{P}_{\mathrm{col}}$
は、
有限時間内に
$\mathrm{Col}_{i,T}=1$
となる軌道記録の集合である。

両者は異なる条件によって定まるため、
Phase~1 では一般に
\[
  \mathfrak{P}_{\mathrm{blk}}
  \subseteq
  \mathfrak{P}_{\mathrm{col}},
\]
\[
  \mathfrak{P}_{\mathrm{col}}
  \subseteq
  \mathfrak{P}_{\mathrm{blk}},
\]
または
\[
  \mathfrak{P}_{\mathrm{blk}}
  =
  \mathfrak{P}_{\mathrm{col}}
\]
のいずれも仮定しない。
\end{remark}

\subsubsection{補集合とPhase~1の限定}

\begin{remark}[$\mathfrak{P}\setminus\mathfrak{P}_{\mathrm{blk}}$ の意味]
\label{rem:blk_in_trajectory_space}
補集合
\[
  \mathfrak{P}
  \setminus
  \mathfrak{P}_{\mathrm{blk}}
\]
は、
各軌道記録 $\sigma$ の
ACS記録有効時刻域
$I_{\mathrm{ACS}}(\sigma)$
のすべての時刻において
$A_{i,t}\neq\emptyset$
である軌道記録の集合である。

ただし本論文は、
この補集合に属することから、
SM の時刻別充足述語が全時刻で真であること、
ロバスト可行性核内に
留まり続けること、
整形式実現が常に成立すること、
または崩壊を回避することを
一般的な帰結として置かない。

特に、
閉塞を経由せずに
崩壊へ到達する軌道を
本論文は排除しない。

なお、
この補集合に属する軌道記録では、
$I_{\mathrm{ACS}}(\sigma)\cap\mathbb{T}_{\mathrm{P},i}$（\S\ref{sec:three_principles} 定義~\ref{def:principle_time_domain}）の時刻のうち
$(x_t,e_t)\in\mathcal{D}_{\mathrm{FR}}$ となる時刻において、
FR の時刻別充足述語は真である
（備考~\ref{rem:internal_quantities_not_principles}）。
\end{remark}

\begin{remark}[閉塞経路集合のPhase~1における限定]
\label{rem:blk_phase1_scope}
Phase~1 では、
$\mathfrak{P}_{\mathrm{blk}}$
を
$\mathfrak{P}$ の集合論的部分集合として
定義するに留める。

本論文は一般に、
$\mathfrak{P}_{\mathrm{blk}}$ の
測度、
確率、
位相的境界、
密度、
距離、
吸引域、
または
「閉塞への近さ」
を定義しない。

また、
有限時刻までの観測だけから
将来
$\mathfrak{P}_{\mathrm{blk}}$
へ属することを
完全に予測する一般手続きも
Phase~1 では与えない。

後続の
\S\ref{subsec:inflow_paths}
で扱う内部動学・関係性の変化も、
それだけから
$\mathfrak{P}_{\mathrm{blk}}$
への帰属を導くものとはしない。
\end{remark}

%----------------------------------------------------------------------
\subsection{崩壊到達経路集合}
\label{subsec:col_path_set}

本小節では、
同じ軌道記録空間
$\mathfrak{P}$
を用いて、
有限時刻内に崩壊終端へ到達する軌道記録の集合
$\mathfrak{P}_{\mathrm{col}}$
を定義する。

ここで、
閉塞は
通常ACS記録有効時刻域
$I_{\mathrm{ACS}}(\sigma)$
内で成立する一時刻・一ステップ条件であるのに対し、
崩壊は
全離散時刻上に保持される
崩壊指標
$\mathrm{Col}_{i,t}$
によって記録される終端分類である。

したがって、
$\mathfrak{P}_{\mathrm{blk}}$
と
$\mathfrak{P}_{\mathrm{col}}$
は、
同じ軌道記録空間の部分集合であるが、
異なる条件によって定まる。

\subsubsection{崩壊到達経路集合}

\begin{definition}[崩壊到達経路集合（分類C）]
\label{def:col_path_set}
\emph{崩壊到達経路集合}
$\mathfrak{P}_{\mathrm{col}}$
を、
軌道記録空間
$\mathfrak{P}$
の部分集合として
\[
  \mathfrak{P}_{\mathrm{col}}
  :=
  \left\{
    \sigma\in\mathfrak{P}
    \;\middle|\;
    \exists T\in\mathbb{N}_0
    \text{ s.t. }
    \mathrm{Col}_{i,T}=1
  \right\}
\]
と定義する。

すなわち、
$\mathfrak{P}_{\mathrm{col}}$
は、
ある有限時刻 $T$ において
崩壊指標
\[
  \mathrm{Col}_{i,T}=1
\]
が成立する軌道記録の全体である。

したがって、
\[
  \mathfrak{P}_{\mathrm{col}}
  \subseteq
  \mathfrak{P}.
\]

下付き
$\mathrm{col}$
は
collapse
を表す。
\end{definition}

\begin{remark}[崩壊到達経路集合は閉塞の成立を要求しない]
\label{rem:col_path_no_blockage_requirement}
\[
  \sigma\in\mathfrak{P}_{\mathrm{col}}
\]
であるために必要なのは、
ある有限時刻に
$\mathrm{Col}_{i,T}=1$
が成立することである。

したがって、
その以前に
$A_{i,t}=\emptyset$
が成立していることを
定義上要求しない。

これは、
\S\ref{subsec:collapse_def}
で確定した
「崩壊は閉塞の事前経由を必要条件としない」
という方針と整合する。
\end{remark}

\subsubsection{崩壊到達時刻}

\begin{definition}[崩壊到達時刻（分類C）]
\label{def:col_reaching_time}
軌道記録
\[
  \sigma\in\mathfrak{P}_{\mathrm{col}}
\]
に対して、
\[
  T_{\mathrm{col}}(\sigma)
  :=
  \min
  \left\{
    T\in\mathbb{N}_0
    \;\middle|\;
    \mathrm{Col}_{i,T}=1
  \right\}
\]
を
$\sigma$ の
\emph{崩壊到達時刻}
と呼ぶ。

$\sigma\in\mathfrak{P}_{\mathrm{col}}$
であれば、
定義~\ref{def:col_path_set} により
上記集合は非空である。

また、
$\mathbb{N}_0$
は整列集合であるため、
最小元が存在する。

したがって、
$T_{\mathrm{col}}(\sigma)$
は有限の自然数として定まる。
\end{definition}

\begin{remark}[崩壊到達時刻以降の指標]
\label{rem:col_reaching_time}
命題~\ref{prop:collapse_basic} の
吸収性条件 (C-2) により、
$t\ge T_{\mathrm{col}}(\sigma)$
ならば、
$\mathrm{Col}_{i,t}=1$
である。

したがって、
$T_{\mathrm{col}}(\sigma)$ は、
崩壊終端記録が最初に成立する時刻である。

一方、
\S\ref{subsec:blk_path_set}
の閉塞到達時刻
$T_{\mathrm{blk}}(\sigma)$
は、
閉塞が最初に成立する時刻にすぎず、
それ以後の閉塞持続を保証しない。

この意味で、
$T_{\mathrm{blk}}$
と
$T_{\mathrm{col}}$
は異なる型の到達時刻である。
\end{remark}

\begin{remark}[崩壊到達時刻とACS記録有効時刻域]
\label{rem:col_time_and_valid_domain}
\[
  \sigma\in\mathfrak{P}_{\mathrm{col}}
\]
であり、
$T_{\mathrm{col}}(\sigma)=T$
とする。

\S\ref{subsec:blk_path_set}
の定義~\ref{def:acs_valid_time_domain} より、
$T\ge1$ の場合は
\[
  I_{\mathrm{ACS}}(\sigma)
  =
  \{0,1,\ldots,T-1\}
\]
であり、$T=0$ の場合は
\[
  I_{\mathrm{ACS}}(\sigma)=\emptyset
\]
である。

したがって、
\[
  T_{\mathrm{col}}(\sigma)
  \notin
  I_{\mathrm{ACS}}(\sigma).
\]

すなわち、
崩壊到達時刻そのものに
通常ACS状態量
\[
  x_T,\qquad
  e_T,\qquad
  A_{i,T},\qquad
  \mathrm{Desc}_T
\]
を要求しない。

この終端規約により、
崩壊時刻へ
人工的な通常ACS状態を
付与する必要がない。
\end{remark}

\subsubsection{閉塞経路集合との関係}

\begin{remark}[閉塞経路集合と崩壊到達経路集合は一般に別対象である]
\label{rem:col_vs_blk}
閉塞経路集合
$\mathfrak{P}_{\mathrm{blk}}$
は、
通常ACS記録有効時刻域内で
少なくとも一度
$A_{i,t}=\emptyset$
が成立する軌道記録の集合である。

一方、
崩壊到達経路集合
$\mathfrak{P}_{\mathrm{col}}$
は、
ある有限時刻に
$\mathrm{Col}_{i,t}=1$
が成立する軌道記録の集合である。

両者は
異なる形式条件によって定まるため、
Phase~1 では一般に
\[
  \mathfrak{P}_{\mathrm{blk}}
  \subseteq
  \mathfrak{P}_{\mathrm{col}},
\]
\[
  \mathfrak{P}_{\mathrm{col}}
  \subseteq
  \mathfrak{P}_{\mathrm{blk}},
\]
または
\[
  \mathfrak{P}_{\mathrm{blk}}
  =
  \mathfrak{P}_{\mathrm{col}}
\]
のいずれも仮定しない。

したがって本論文は、
本節の定義と仮定だけから、
両集合のいずれの方向の包含関係も、
両集合の一致も、一般的な帰結として置かない。
\end{remark}

\begin{remark}[閉塞しても崩壊しない軌道を許容する]
\label{rem:blk_without_collapse_path}
\S\ref{subsec:blockage}および\S\ref{subsec:collapse_def}で述べたとおり、
本論文は、閉塞 $A_{i,t}=\emptyset$ から $\mathrm{Col}_{i,t}=1$ を
一般的な帰結として置かない。
これに対応して、
\[
  \mathfrak{P}_{\mathrm{blk}}
  \setminus
  \mathfrak{P}_{\mathrm{col}}
\]
に属する軌道記録を
Phase~1 は排除しない。

すなわち、
有限時刻に閉塞へ到達しても、
その後に回復し、
いかなる有限時刻にも
崩壊指標が $1$ とならない軌道を
一般論から排除しない。

ただし、
この差集合が
必ず非空であるとまでは
主張しない。

その非空性は、
対象系の動学、
環境実現（背景成分の列と $g_O$）、
外的回復可能性等に依存する。
\end{remark}

\begin{remark}[閉塞を経由せず崩壊する軌道を許容する]
\label{rem:col_without_blk_path}
\S\ref{subsec:collapse_def} では、
崩壊の成立に
閉塞の事前経由を
必要条件として課さなかった。

したがって、
軌道集合の水準でも、
\[
  \mathfrak{P}_{\mathrm{col}}
  \setminus
  \mathfrak{P}_{\mathrm{blk}}
\]
に属する軌道記録を
Phase~1 は排除しない。

すなわち、
崩壊到達時刻
$T_{\mathrm{col}}(\sigma)$
より前の
すべての
$t\in I_{\mathrm{ACS}}(\sigma)$
について
$A_{i,t}\neq\emptyset$
であるにもかかわらず、
有限時刻で崩壊終端へ到達する軌道を
一般論から排除しない。

ただし、
この差集合が
必ず非空であるとまでは
主張しない。
\end{remark}

\begin{remark}[閉塞経路集合と崩壊到達経路集合の交差]
\label{rem:blk_col_intersection}
ある軌道記録が
\[
  \mathfrak{P}_{\mathrm{blk}}
  \cap
  \mathfrak{P}_{\mathrm{col}}
\]
に属する場合、
その軌道は
崩壊以前の
通常ACS記録有効時刻域内で
少なくとも一度閉塞を経験し、
かつ有限時刻内に崩壊終端へ到達する。

しかし、
その場合でも、
閉塞が崩壊の原因であること、
閉塞が崩壊の十分条件であること、
または
最初の閉塞後に
連続して閉塞が持続したことを、
本論文は定義からの帰結として置かない。

したがって、
集合の交差は
共起を表すだけであり、
因果関係を表さない。
\end{remark}

\subsubsection{到達時刻どうしの関係}

\begin{remark}[$T_{\mathrm{blk}}$ と $T_{\mathrm{col}}$ の大小関係]
\label{rem:blk_col_time_relation}
軌道記録
\[
  \sigma
  \in
  \mathfrak{P}_{\mathrm{blk}}
  \cap
  \mathfrak{P}_{\mathrm{col}}
\]
を考える。

閉塞到達時刻
$T_{\mathrm{blk}}(\sigma)$
は
$I_{\mathrm{ACS}}(\sigma)$
内に属する。

一方、
崩壊到達時刻
$T_{\mathrm{col}}(\sigma)$
は
$I_{\mathrm{ACS}}(\sigma)$
に属さない。

したがって、
\S\ref{subsec:blk_path_set}の終端規約のもとでは、
\[
  T_{\mathrm{blk}}(\sigma)
  <
  T_{\mathrm{col}}(\sigma)
\]
が成立する。

これは、
閉塞が崩壊の原因または必要条件であることを
意味しない。

意味するのは、
同一軌道が
両集合に属する場合、
閉塞は通常ACS記録が有効な時刻に記録され、
崩壊はその通常ACS記録域を終了させる
後続の終端時刻として記録されるという
型上の順序である。
\end{remark}

\subsubsection{崩壊到達と他の失敗関連概念}

\begin{remark}[整形式実現の不成立と崩壊到達経路集合への帰属の関係]
\label{rem:wf_not_col_path}
本論文は、ある時刻 $t$ に
$\mathrm{WF}_{i,t}=0$
であっても、
それだけから
\[
  \sigma\in\mathfrak{P}_{\mathrm{col}}
\]
を一般的な帰結として置かない。

\S\ref{subsec:collapse_def}で述べたとおり、
本論文は、$\mathrm{WF}_{i,t}=0$ から $\mathrm{Col}_{i,t}=1$ を
一般的な帰結として置かない。

したがって、
整形式実現の不成立は
崩壊到達経路集合への帰属条件ではない。

$\mathfrak{P}_{\mathrm{col}}$
への帰属には、
有限時刻で実際に
$\mathrm{Col}_{i,T}=1$
が成立することを要する。
\end{remark}

\begin{remark}[三原理の時刻別充足述語・T-STAB・関係品質と崩壊到達の関係]
\label{rem:principles_not_col_path}
本論文は、ある時刻または有限窓において、
SM・EU・FR のいずれかの時刻別充足述語が偽であること、
T-STAB が不成立であること、
関係品質
$Q<0$
であること、
または
$\chi_{\mathrm{EI},t}=0$
であることだけから、
\[
  \sigma\in\mathfrak{P}_{\mathrm{col}}
\]
を一般的な帰結として置かない。

これらは、
崩壊到達経路集合への
定義上の帰属条件ではない。

それらと崩壊前状態への流入との
記述的接続は、
\S\ref{subsec:inflow_paths}
で扱う。
\end{remark}

\begin{remark}[ロバスト可行性核外と崩壊到達の関係]
\label{rem:viab_outside_not_col_path}
本論文は、ある時刻で
$x_t\notin\mathrm{Viab}_f(K)$
であっても、
それだけから
\[
  \sigma\in\mathfrak{P}_{\mathrm{col}}
\]
を一般的な帰結として置かない。

ロバスト可行性核外であることは、
$E_{\mathrm{req}}$ に対する
長期ロバスト不変性が
保証されないことを意味するが、
本論文は、そこから崩壊終端への有限時間到達を
一般的な帰結として置かない。

したがって本論文は、
$x_t\notin\mathrm{Viab}_f(K)$ から
$\sigma\in\mathfrak{P}_{\mathrm{col}}$ を
一般的な帰結として置かない。
\end{remark}

\subsubsection{補集合とPhase~1の限定}

\begin{remark}[$\mathfrak{P}\setminus\mathfrak{P}_{\mathrm{col}}$ の意味]
\label{rem:col_avoidance}
補集合
\[
  \mathfrak{P}
  \setminus
  \mathfrak{P}_{\mathrm{col}}
\]
は、
いかなる有限時刻においても
$\mathrm{Col}_{i,t}=1$
が成立しない軌道記録の集合である。

定義~\ref{def:collapse} の
崩壊指標は二値であるため、
これは
\[
  \forall t\in\mathbb{N}_0,
  \quad
  \mathrm{Col}_{i,t}=0
\]
である軌道記録の集合に一致する。

ただし本論文は、
この補集合に属することから、
閉塞を経験しないこと、
SM・EU・FR の時刻別充足述語が定義される全時刻で真であること、
WF が常に成立すること、
またはロバスト可行性核内に
留まり続けることを一般的な帰結として置かない。
\end{remark}

\begin{remark}[崩壊到達経路集合のPhase~1における限定]
\label{rem:col_phase1_scope}
Phase~1 では、
$\mathfrak{P}_{\mathrm{col}}$
を
$\mathfrak{P}$ の
集合論的部分集合として定義するに留める。

本論文は一般に、
$\mathfrak{P}_{\mathrm{col}}$ の
測度、
確率、
位相的境界、
密度、
距離、
吸引域、
または
「崩壊への近さ」
を定義しない。

また、
有限時刻までの観測から
将来
$\mathfrak{P}_{\mathrm{col}}$
へ属することを
完全に予測する一般手続きも
Phase~1 では与えない。

\S\ref{subsec:collapse_def}
で述べたとおり、
崩壊の理論上の定義と
有限観測による判定可能性は
別問題である。
\end{remark}

%----------------------------------------------------------------------
\subsection{内部動学および関係性に由来する崩壊前状態への流入}
\label{subsec:inflow_paths}

本小節では、
ACS個体 $\Sigma_i$ の
内部動学上の変化および
外部対象との関係に由来する変化が、
\S\ref{subsec:precollapse_states}で導入した
記述的崩壊前状態帯
\[
  \mathsf{Z}_1,\ldots,\mathsf{Z}_4
\]
に関連しうることを、
\emph{流入傾向}
という記述レベルで整理する。

ここでいう流入傾向は、
特定の時刻量または関係状態から
$\mathsf{Z}_j$、
閉塞、
または崩壊への
決定論的遷移を表すものではない。

また、
流入傾向は
\[
  \sigma\in\mathfrak{P}_{\mathrm{blk}}
\]
または
\[
  \sigma\in\mathfrak{P}_{\mathrm{col}}
\]
という経路集合への帰属とも区別される。

したがって本小節は、
\[
  \text{時刻依存量・関係状態の変化}
\]
と
\[
  \text{記述的崩壊前状態帯}
\]
と
\[
  \text{閉塞・崩壊経路集合への帰属}
\]
の三つを同一視しない。

\subsubsection{流入傾向の記述上の地位}

\begin{remark}[流入傾向は形式的な状態遷移関係ではない]
\label{rem:inflow_is_descriptive}
本小節で用いる
「流入」または
「流入傾向」という語は、
状態空間 $X$ 上の
新しい遷移関係、
ベクトル場、
距離関数、
勾配流、
または確率遷移を
導入するものではない。

すなわち、
状態帯 $\mathsf{Z}_j$ と記述される状態から
$\mathsf{Z}_k$ と記述される状態への
形式的遷移規則を本小節では定義しない。

また、
軌道記録空間
$\mathfrak{P}$
には
Phase~1 で
距離または位相を一般には付与していないため、
「$\mathfrak{P}_{\mathrm{blk}}$ への距離」
または
「$\mathfrak{P}_{\mathrm{col}}$ への接近」
という形式量も導入しない。

本小節の流入傾向は、
崩壊前状態に関連しうる
時刻変化や関係変化を
記述的に整理するための概念である。
\end{remark}

\begin{remark}[流入傾向は分類Cの記述である]
\label{rem:inflow_description_level}
本小節で参照する
\S\ref{sec:extended_dynamics} 由来の時刻依存量、
\S\ref{sec:rel_min} 由来の関係記述、
および
$\mathsf{Z}$ 帯は、
いずれも理論記述側の対象である。

したがって、
ACS個体 $\Sigma_i$ 自身が
「自分は崩壊前状態へ流入している」
という同一の内部変数や
離散ラベルを
保持・計算することを要求しない。

同様に、
TIPA の M2 その他の機構が
本小節の
流入傾向、
$\mathsf{Z}_j$、
$\mathfrak{P}_{\mathrm{blk}}$、
または
$\mathfrak{P}_{\mathrm{col}}$
への帰属を
真値として直接認識することも
Phase~1 の要件ではない。
\end{remark}

\subsubsection{時刻依存量との接続}

\begin{remark}[\S\ref{sec:extended_dynamics} の時刻依存量の参照]
\label{rem:internal_inflow_types}
\S\ref{sec:extended_dynamics} で定義される
時刻依存量の整理用タプル
\[
  \mathrm{Desc}_t
  =
  \bigl(
    \boldsymbol{\beta}_t,\,
    d_{\mathrm{eff},t},\,
    \pi_{\tau,t},\,
    \widehat{\tau}_{t},\,
    h_{\mathrm{psl},t},\,
    h_{\mathrm{mem},t},\,
    h_{\mathrm{sens},t},\,
    h_{\mathrm{eff},t}
  \bigr)
\]
を参照する。
このタプルは分類Cの整理用表記であり、
独立した状態ベクトルではない。

本小節では、
これらの量が
有限時刻区間にわたり変化し、
その変化が
通常作動・基準帯
$\mathsf{Z}_0$
との差異を大きくする方向に
記述される場合、
記述的崩壊前状態帯
$\mathsf{Z}_1$ から
$\mathsf{Z}_4$
との関連を
検討しうるとする。

ただし、本論文は、
特定の値、
特定の符号、
特定の閾値、
または特定の有限窓差分から
特定の
$\mathsf{Z}_j$
への帰属を
一般的な帰結として置かない。
\end{remark}

健全度量と閉塞・崩壊・状態帯の関係は、
\S\ref{sec:extended_dynamics} の備考~\ref{rem:s5_health_noncollapse} と
備考~\ref{rem:health_not_zone_identity} に従う。

\begin{remark}[$d_{\mathrm{eff},t}$ と $\mu_{\mathrm{eff},t}$ は異なる]
\label{rem:d_eff_not_mu_eff}
\S\ref{sec:extended_dynamics} の
$d_{\mathrm{eff},t}$
と、
\S\ref{sec:three_principles} で
分類Cの応答多様度として扱う
$\mu_{\mathrm{eff},t}$
は、
異なる形式対象である。

$\mu_{\mathrm{eff},t}$ は、
時刻 $t$ の
$(x_t,e_t)$
における
可行行動の
$E_{\mathrm{req}}(x_t)$ 上の
可行性プロファイル同値類数として定義される
局所的な一ステップ応答多様性である。

これに対し、
$d_{\mathrm{eff},t}$
は
\S\ref{sec:extended_dynamics}
で扱う
時刻依存の発展自由度・有効多様性側の
記述量である。

したがって、
\[
  d_{\mathrm{eff},t}
  =
  \mu_{\mathrm{eff},t}
\]
を仮定しない。

また、
\[
  d_{\mathrm{eff},t}\downarrow
  \Longrightarrow
  \mu_{\mathrm{eff},t}\downarrow
\]
や、
\[
  d_{\mathrm{eff},t}=0
  \Longrightarrow
  \mu_{\mathrm{eff},t}<2
\]
も、
追加条件なしには採用しない。

さらに、
局所 FR の定義条件は $\mu_{\mathrm{eff},t}$ を用いない。
局所 FR は
\S\ref{sec:three_principles} の定義域 $\mathcal{D}_{\mathrm{FR}}$ 上で
可行行動集合の非空性によって判定されるため、
本論文は、$\mu_{\mathrm{eff},t}<2$ であることだけから
局所 FR が偽であることを一般的な帰結として置かない。
同様に、
$d_{\mathrm{eff},t}$ の低下を
局所 FR 条件の不充足そのものとして扱わない。
\end{remark}

\begin{remark}[時間地平量と崩壊を直接結びつけない]
\label{rem:time_horizon_not_collapse}
\S\ref{sec:extended_dynamics} の
$\pi_{\tau,t}$
および
$\widehat{\tau}_{t}$
は、
時間地平に関係する
時刻依存量である。

これらの低下が
対象系によっては
予測・評価・応答余裕の縮小と
関連しうるとしても、
本論文は一般に
\[
  \widehat{\tau}_{t}=0
  \Longrightarrow
  A_{i,t}=\emptyset
\]
または
\[
  \widehat{\tau}_{t}=0
  \Longrightarrow
  \mathrm{Col}_{i,t}=1
\]
を採用しない。

同様に、本論文は、
$\pi_{\tau,t}$ の特定値から
特定の
$\mathsf{Z}_j$
への帰属も
一般的な帰結として置かない。
\end{remark}

\begin{remark}[偏向強度の変化を崩壊方向と固定しない]
\label{rem:beta_not_fixed_collapse_direction}
\S\ref{sec:extended_dynamics} の
\[
  \beta_{\iota,t}\quad(\iota\in\mathcal J_{\beta})
\]
は、
参照基準に相対して定める偏向性 $\iota$（定義~\ref{def:bias_tendency}）について、
その時刻 $t$ における程度を表す偏向強度である。
$\beta_{\iota,t}$ の値からは、偏向性の帰属区分も方向も読まない。

本小節では、
$\beta_{\iota,t}$ の増大または減少の
一方を
普遍的に
「崩壊方向」として固定しない。

対象系における
偏向性 $\iota$ の参照基準と尺度、
読み出し規則と担体の遷移、
他の時刻依存量との組合せにより、
崩壊前状態との関連は異なりうる。

したがって、
\[
  \beta_{\iota,t}\uparrow
  \Longrightarrow
  \sigma\in\mathfrak{P}_{\mathrm{blk}}
\]
または
\[
  \beta_{\iota,t}\uparrow
  \Longrightarrow
  \sigma\in\mathfrak{P}_{\mathrm{col}}
\]
のような一般的含意を
Phase~1 では置かない。
\end{remark}

\subsubsection{三原理の時刻別充足述語との接続}

\begin{remark}[内部時刻量と三原理の時刻別充足述語を同一視しない]
\label{rem:internal_quantities_not_principles}
\S\ref{sec:extended_dynamics} の
時刻依存量の変化は、
対象系によっては
SM・EU・FR の時刻別充足述語が偽となる状態変化と
併存または関連しうる。

しかし、
本論文は
\[
  h_{q,t},
  \quad
  d_{\mathrm{eff},t},
  \quad
  \boldsymbol{\beta}_t,
  \quad
  \widehat{\tau}_{t}
\]
のいずれかを
SM・EU・FR の時刻別充足述語の真理値そのものとはしない。

特に、
Phase~1 における時刻 $t$ の SM の時刻別充足述語は
$\Pi(x_t)\ge0$
によって判定される。
EU の時刻別充足述語は
\S\ref{sec:three_principles} の適用域で定義され、
同節の一段階臨界必要要因集合 $\mathrm{Crit}^{(1)}(x_t)$ に属する必要環境要因の充足によって判定される。
FR の時刻別充足述語（局所 FR）は同節の定義域
$\mathcal{D}_{\mathrm{FR}}$
上で定義され、
\[
  \mathrm{FR}_t
  \Longleftrightarrow
  A_{i,t}\neq\emptyset
\]
によって判定される。
$\mu_{\mathrm{eff},t}$ は応答多様度を表す分類C量であり、
局所 FR の真理値そのものではない。

したがって、
\S\ref{sec:extended_dynamics} の時刻依存量の変化から
三原理の時刻別充足述語の真偽を
直接置換的に判定しない。
\end{remark}

\begin{remark}[三原理の時刻別充足述語が偽であることと経路集合への帰属]
\label{rem:principle_failure_not_path_membership}
本論文は、SM の時刻別充足述語が偽であることだけからは、
\[
  \sigma\in\mathfrak{P}_{\mathrm{blk}}
\]
も
\[
  \sigma\in\mathfrak{P}_{\mathrm{col}}
\]
も一般的な帰結として置かない（備考~\ref{rem:sm_wf_not_blk_path}、備考~\ref{rem:principles_not_col_path}）。

FR の時刻別充足述語は
定義域 $\mathcal{D}_{\mathrm{FR}}$ 上で
$A_{i,t}\neq\emptyset$ と同値に定まる
（備考~\ref{rem:internal_quantities_not_principles}）。
したがって、
FR の時刻別充足述語が時刻 $t$ に偽であれば
$A_{i,t}=\emptyset$ であり、
$t\in I_{\mathrm{ACS}}(\sigma)$ のとき、
定義~\ref{def:blk_path_set} により
$\sigma\in\mathfrak{P}_{\mathrm{blk}}$ が従う。
EU の時刻別充足述語が偽であることと閉塞との関係は
\S\ref{sec:three_principles} で扱い、
本備考はその関係を前提としない。

一方、本論文は、
三原理のいずれの時刻別充足述語が偽であることからも、
それだけから
$\sigma\in\mathfrak{P}_{\mathrm{col}}$
を一般的な帰結として置かない（備考~\ref{rem:principles_not_col_path}）。

したがって、
本小節で
内部時刻量の変化と
三原理の時刻別充足述語が偽であることとの
記述的接続を述べる場合でも、
それを
閉塞または崩壊への
決定論的経路としては扱わない。
\end{remark}

\subsubsection{関係性に由来する変化との接続}

\begin{remark}[関係記述子と流入傾向]
\label{rem:relational_inflow}
\S\ref{sec:rel_min} で扱う
外部対象との関係記述は、
ACS個体 $\Sigma_i$ と
外部対象との間で実現している
関係状態を記述する。

関係状態の時間変化は、
対象系によっては、
内部状態 $x_t$、
\S\ref{sec:extended_dynamics} の時刻依存量、
必要環境充足、
または
可行行動集合の変化と
関連しうる。

本小節では、
このような関係性由来の変化が
記述的崩壊前状態帯との
関連を形成しうることを許容する。

ただし、本論文は、
特定の関係カテゴリ、
特定の関係記述子成分、
または特定の関係方向から、
特定の
$\mathsf{Z}_j$、
閉塞、
または崩壊を
一般的な帰結として置かない。
\end{remark}

\begin{remark}[関係カテゴリは崩壊原因ラベルではない]
\label{rem:relation_category_not_collapse_cause}
\S\ref{sec:rel_min} において
関係カテゴリが導入されている場合でも、
本節は
あるカテゴリを
「崩壊を生じさせるカテゴリ」
として固定しない。

同じカテゴリ名で記述される関係でも、
関係の具体的内容、
持続時間、
方向性、
対象状態、
環境状態、
内部状態、
および他の関係との組合せにより、
ACSへの影響は異なりうる。

したがって、
関係カテゴリから
\[
  \sigma\in\mathfrak{P}_{\mathrm{blk}}
\]
または
\[
  \sigma\in\mathfrak{P}_{\mathrm{col}}
\]
への一般的含意を
採用しない。
\end{remark}

\begin{remark}[関係品質 $Q$ の負値と崩壊の関係]
\label{rem:negative_Q_not_collapse}
\S\ref{subsec:ro_quality} の
関係品質
$Q_{i\leftarrow j,t}$
は、
評価可能な関係の品質を
記述する分類C量である。

$Q_{i\leftarrow j,t}<0$
であることは、
当該品質指標上で
負側の評価であることを表すが、
本論文は、それだけから
$\Pi(x_t)$
の低下方向、
閉塞、
整形式実現の不成立、
または崩壊を
一般的な帰結として置かない。

したがって本論文は、
$Q<0$ から
$\sigma\in\mathfrak{P}_{\mathrm{blk}}$ も
$\sigma\in\mathfrak{P}_{\mathrm{col}}$ も
一般的な帰結として置かない。

関係品質の時間的蓄積や
複数関係の統合評価を
崩壊経路へ接続する追加モデルは、
Phase~1 の本小節では導入しない。
\end{remark}

\begin{remark}[構造的環境結合と実現関係変化を区別する]
\label{rem:env_coupling_not_relational_inflow}
\S\ref{sec:three_principles} で定義される
構造的環境結合条件は、
ACS と外部環境との構造的な結合可能性に関する背景条件である。
これは環境利用 EU の前提・背景構造として位置づけられるが、
EU の時刻別充足述語そのものではない。

これに対して、
本小節で扱う
関係性由来の変化は、
実現した外部対象との関係状態や
その時間変化に関する記述である。

したがって、
構造的環境結合条件、実現関係状態および関係性由来の流入傾向は、
異なる記述対象として扱い、
相互に同一視しない。

構造的環境結合条件が成立していても
不利な関係状態が実現する場合を
排除しないし、
同条件が不成立であることから
直ちに閉塞・崩壊を導くこともない。
\end{remark}

\subsubsection{流入傾向と状態帯の関係}

\begin{remark}[流入傾向から特定の $\mathsf{Z}$ 帯を一意に定めない]
\label{rem:inflow_not_unique_zone}
内部時刻量または関係状態に
変化が観察されたとしても、
その変化から
\[
  \mathsf{Z}_1,
  \mathsf{Z}_2,
  \mathsf{Z}_3,
  \mathsf{Z}_4
\]
のいずれか一つを
一意に決定する
普遍的写像を
Phase~1 では定義しない。

同一の変化が、
ある対象系では
軽度障害顕在化に関連し、
別の対象系では
通常作動範囲内に留まり、
また別の対象系では
複数の記述帯と
同時に関連する場合を
排除しない。

これは、
\S\ref{subsec:precollapse_states}で
$\mathsf{Z}$ 帯を
状態空間分割または全順序として
定義しなかったことと整合する。
\end{remark}

\begin{remark}[流入傾向は基準帯からの単調離脱を意味しない]
\label{rem:inflow_not_monotone}
流入傾向という語は、
$\mathsf{Z}_0$ から $\mathsf{Z}_5$ まで
添字順に状態帯を経由する
単調な進行系列を意味しない。

内部量や関係状態の変化後に
通常作動・基準帯側へ戻ること、
異なる記述帯へ移ること、
または
明示的な崩壊前状態帯を経ずに
崩壊終端へ到達することを
Phase~1 は排除しない。
\end{remark}

\subsubsection{流入傾向と経路集合への帰属}

\begin{remark}[流入と経路集合への帰属の区別]
\label{rem:inflow_vs_membership}
本小節でいう
流入傾向と、
軌道記録が
$\mathfrak{P}_{\mathrm{blk}}$
または
$\mathfrak{P}_{\mathrm{col}}$
に属することは
別の条件である。

ある軌道記録において、
健全度低下、
$d_{\mathrm{eff},t}$ の変化、
時間地平量の変化、
偏向強度の変化、
または関係状態の変化が
崩壊前状態との関連を示していても、
本論文は、それだけから
有限時刻内の閉塞または崩壊を
一般的な帰結として置かない。

したがって本論文は、
流入傾向から
$\sigma\in\mathfrak{P}_{\mathrm{blk}}$ も
$\sigma\in\mathfrak{P}_{\mathrm{col}}$ も
一般的な帰結として置かない。
\end{remark}

\begin{remark}[経路集合への帰属から特定の流入機構を逆算しない]
\label{rem:membership_not_unique_inflow}
逆に、本論文は、
\[
  \sigma\in\mathfrak{P}_{\mathrm{blk}}
\]
または
\[
  \sigma\in\mathfrak{P}_{\mathrm{col}}
\]
であることから、
その軌道が
特定の健全度低下、
特定の $d_{\mathrm{eff},t}$ 変化、
特定の関係カテゴリ、
または特定の $Q$ 変化を
事前に経由したことを
一般的な帰結として置かない。

閉塞または崩壊へ至る
具体的機構は、
対象系ごとの
内部動学、
環境実現（背景成分の列と $g_O$）、
関係状態、
外的作用、
その他の条件に依存しうる。

したがって、
経路集合への帰属は
原因機構の一意な同定を
意味しない。
\end{remark}

\begin{remark}[崩壊前状態帯を経由しない崩壊との整合]
\label{rem:inflow_and_direct_collapse}
\S\ref{subsec:collapse_def}および
\S\ref{subsec:col_path_set}で確定したとおり、
崩壊
$\mathrm{Col}_{i,t}=1$
は、
記述的崩壊前状態帯
$\mathsf{Z}_1$ から
$\mathsf{Z}_4$
の事前経由を
必要条件としない。

したがって、
本小節で整理する
「崩壊前状態への流入」は、
崩壊に至りうる
全経路を網羅するものではない。

崩壊前状態帯を
明示的に経由しない
崩壊到達経路は、
本小節の流入整理の対象外であるが、
Phase~1 の軌道記録空間からは
排除されない。
\end{remark}

\subsubsection{M2・TIPAとの接続}

\begin{remark}[経路集合・流入傾向はM2の直接入力ではない]
\label{rem:path_sets_m2_connection}
本節で導入した
\[
  \mathfrak{P}_{\mathrm{blk}},
  \qquad
  \mathfrak{P}_{\mathrm{col}},
  \qquad
  \mathsf{Z}_1,\ldots,\mathsf{Z}_4,
\]
および本小節の
流入傾向は、
理論記述側の対象である。

したがって、
TIPA の M2 が
これらの集合、
集合への帰属、
状態帯ラベル、
または
「流入の程度」
をそのまま内部入力として
保持・計算することを
要求しない。

M2 の正式な入力型、
評価写像、
Priority、
その他の標準機能形式は
\S\ref{sec:tipa}
の責務であり、
本小節では再定義しない。

経路集合に属する軌道上で
可行性関連の状態変化が生じ、
それがM2に利用可能な
感知・統合・内部モデル情報へ
反映される実装は許容される。

しかし、本論文は、
\[
  \sigma\in\mathfrak{P}_{\mathrm{blk}}
\]
または
\[
  \sigma\in\mathfrak{P}_{\mathrm{col}}
\]
であることから、
特定のM2出力や
特定のPriority値を
一般的な帰結として置かない。
\end{remark}

\subsubsection{Phase~1 における限定}

\begin{remark}[流入傾向のPhase~1における限定]
\label{rem:inflow_phase1_scope}
Phase~1 の本小節は、
単一ACS個体 $\Sigma_i$ の
軌道記録について、
内部時刻量および
外部対象との関係状態の変化が
崩壊前状態と
関連しうることを
記述的に整理するに留める。

本小節は、
次のものを一般には導入しない。

\begin{itemize}
  \item
  崩壊前状態への距離またはポテンシャル、
  \item
  流入速度または流入率、
  \item
  閉塞・崩壊確率、
  \item
  特定量から $\mathsf{Z}$ 帯への普遍的閾値写像、
  \item
  特定の関係カテゴリから崩壊への因果法則、
  \item
  流入傾向から経路集合への帰属を保証する十分条件。
\end{itemize}

これらを導入するには、
追加の対象系依存仮定、
位相・測度・確率構造、
または実証的根拠を要する。

その精密な形式化は
Phase~2 以降の課題とする（\S\ref{subsec:limitations_list} の L9）。
\end{remark}

\subsubsection{本小節の要約}

\begin{remark}[流入記述の要約]
\label{rem:inflow_summary}
本小節で維持する関係は、
次のように要約できる。

\[
  \text{時刻依存量の変化}
  \;\dashrightarrow\;
  \text{崩壊前状態との関連可能性}
\]

\[
  \text{関係状態の変化}
  \;\dashrightarrow\;
  \text{崩壊前状態との関連可能性}
\]

ここで
$\dashrightarrow$
は
形式的含意ではなく、
記述的接続を表す。

したがって本論文は、
時刻量または関係状態の変化だけから、
特定の状態帯 $\mathsf{Z}_j$ への分類を一般的な帰結として置かない。
また本論文は、流入傾向だけから、
軌道記録 $\sigma$ が閉塞経路集合
$\mathfrak{P}_{\mathrm{blk}}$ または崩壊到達経路集合
$\mathfrak{P}_{\mathrm{col}}$ に属することを一般的な帰結として置かない。
\end{remark}

%----------------------------------------------------------------------
\subsection{組織ACSへの接続}
\label{subsec:to_section8}

本小節では、
本節の個体レベルの概念を
\S\ref{sec:org_acs}
で扱う組織ACS、
すなわち組織自身についてACS構造が成立する場合の議論へ接続する。

ただし、
\S\ref{sec:collapse_cond} の個体レベル概念を
組織水準へそのまま持ち上げることはしない。
ACS集団と組織ACSの区別は次の備考で述べる。

\subsubsection{ACS集団と組織ACSの区別}

\begin{remark}[ACS集団はそれだけでは組織ACSではない]
\label{rem:acs_group_vs_org_acs}
複数のACS個体
$\Sigma_1,\ldots,\Sigma_n$
と、
それらの間の相互作用または関係構造が存在する場合、
それらは
\S\ref{sec:org_acs}
で扱う
ACS集団の基体を構成しうる。

しかし本論文は、
複数個体が存在すること、
または関係ネットワークが存在することだけから、
組織水準の構造的自己、
組織水準の可行性、
あるいは組織ACSの成立を
一般的な帰結として置かない。

したがって、
\[
  \text{ACS集団}
  \neq
  \text{組織ACS}
\]
である。

組織ACSとしての成立には、
\S\ref{sec:org_acs}
で別途定める
組織水準のACS構造の成立を要する。

本節は、
それらの条件を先取りして
再定義しない。
\end{remark}

\begin{remark}[本節で扱う対象は個体添字 $i$ に属する]
\label{rem:individual_level_failure_objects}
本節で扱った
\[
  A_{i,t},
  \qquad
  \mathrm{WF}_{i,t},
  \qquad
  \mathrm{Col}_{i,t},
\]
\[
  \Upsilon_{\mathrm{adm},i},
  \qquad
  \mathrm{Recover}_{i,t},
\]
および
個体 $\Sigma_i$ の軌道記録に関する
\[
  \mathfrak{P},
  \qquad
  \mathfrak{P}_{\mathrm{blk}},
  \qquad
  \mathfrak{P}_{\mathrm{col}}
\]
は、
\S\ref{sec:collapse_cond} の範囲では
単一ACS個体に対する対象である。

したがって、
これらに含まれる個体添字 $i$ を除去して、
組織水準の閉塞・崩壊・回復へ
自動的に読み替えない。

複数個体を含む系についても、
各個体の閉塞または崩壊は、
まず各
$\Sigma_i$
について個別に判定される。
\end{remark}

\subsubsection{個体崩壊と組織水準状態の非同一性}

\begin{remark}[個体の崩壊と組織ACSの崩壊を同一視しない]
\label{rem:individual_vs_org_collapse}
ある構成個体について
$\mathrm{Col}_{i,t}=1$
が成立したとしても、
本論文は、それだけから
組織ACS全体の崩壊を一般的な帰結として置かない。

逆に、
\S\ref{sec:org_acs}
または後続Phaseで
組織水準の崩壊概念を扱う場合でも、
組織水準の崩壊から
すべての構成個体について
$\mathrm{Col}_{i,t}=1$
を導けるとは限らない。

したがって、
\[
  \text{個体ACSの崩壊}
  \neq
  \text{組織水準の崩壊}
\]
である。

組織水準の崩壊を形式化する場合には、
組織自身の構造的自己、
組織水準の可行性、
組織水準の結合構造その他に対する
条件を要する。
その際には、
\S\ref{sec:org_acs} の組織水準のACS構造のもとで、
\S\ref{sec:collapse_cond} の崩壊条件を
組織水準で満たすことを確認する。
\end{remark}

\begin{remark}[個体閉塞と組織水準可行性を同一視しない]
\label{rem:individual_blockage_org_viability}
ある個体
$\Sigma_i$
について
$A_{i,t}=\emptyset$
が成立しても、
本論文は、それだけから
他の個体
$\Sigma_j$
の閉塞、
ACS集団全体の機能停止、
または組織ACSの可行性喪失を
一般的な帰結として置かない。

同様に、
複数個体が同時に閉塞していることだけから、
組織ACSの崩壊を定義しない。

個体レベルの失敗状態を
組織水準の状態へ接続するには、
\S\ref{sec:org_acs}
で導入される
組織水準のACS構造を
別途用いる必要がある。
\end{remark}

\subsubsection{関係性を介した波及の位置づけ}

\begin{remark}[関係性を介した波及は可能性としてのみ接続する]
\label{rem:multi_agent_collapse_deferral}
\S\ref{subsec:inflow_paths} では、
外部対象との関係状態の変化が、
単一ACS個体の内部状態や
時刻依存量の変化と
関連しうることを記述した。

外部対象が別のACS個体
$\Sigma_j$
である場合、
$\Sigma_j$ の行動または状態変化が
関係を介して
$\Sigma_i$ の状態変化に関与する場合を
Phase~1 の一般論は排除しない。

しかし、
このことから
「崩壊が個体間を伝播する」
という形式法則を、
本節では一般的な帰結として置かない。

特に、
\[
  \mathrm{Col}_{j,t}=1
  \Longrightarrow
  \mathrm{Col}_{i,t+k}=1
\]
のような
個体間の崩壊伝播則を
追加仮定なしには置かない。

多主体的な波及を
正式な動学として扱うには、
複数個体の状態・関係構造の
同時時間発展を記述する
追加形式が必要である。

その精密化は、
\S\ref{sec:org_acs}
での構造的接続を経たうえで、
主として Phase~2 以降の課題とする（問題~\ref{prob:social_dynamics}）。
\end{remark}

\begin{remark}[関係カテゴリと多主体的崩壊の関係]
\label{rem:relation_category_collapse_scope}
\S\ref{sec:rel_min}
で用いる関係カテゴリまたは関係記述子は、
個体間関係を記述するための対象である。

本論文は、特定の関係カテゴリが存在すること、
特定の関係品質
$Q$
が負であること、
または
関係状態に大きな変化があることだけから、
多主体的な閉塞、
崩壊伝播、
または組織ACSの崩壊を
一般的な帰結として置かない。

したがって、
\S\ref{sec:collapse_cond} から
組織水準または多主体間の因果法則を
先取りして構成しない。
\end{remark}

\subsubsection{外的回復可能性の多主体文脈}

\begin{remark}[他ACSからの作用と $\Upsilon_{\mathrm{adm}}$ を自動同一視しない]
\label{rem:multi_agent_recovery_scope}
あるACS個体
$\Sigma_j$
から
$\Sigma_i$
への作用が存在するとき、
その作用が
$\Sigma_i$ の回復に関与する場合はありうる。

しかし、
$\Sigma_j$ からの任意の作用を
自動的に
\[
  \upsilon\in\Upsilon_{\mathrm{adm},i}
\]
とみなさない。

ある外部作用が
$\Upsilon_{\mathrm{adm},i}$
の要素として扱われるかは、
\S\ref{subsec:external_recovery}
で定めた
$\Sigma_i$ 側の許容作用域に依存する。

また、
複数個体が協調して行う回復作用、
組織構造そのものによる回復作用、
または組織ACSが構成個体へ与える
トップダウン作用を
Phase~1 の本節で
一般形式化しない。

これらを扱う場合には、
\S\ref{sec:org_acs}
の組織構造と整合する
追加定義を要する。
\end{remark}

\subsubsection{経路集合の多主体拡張}

\begin{remark}[経路集合の多主体的拡張は本節で導入しない]
\label{rem:path_set_multi_agent}
本節の
\[
  \mathfrak{P}_{\mathrm{blk}}
  \quad\text{および}\quad
  \mathfrak{P}_{\mathrm{col}}
\]
は、
単一ACS個体の
軌道記録空間
$\mathfrak{P}$
の部分集合として定義された。

複数ACS個体を同時に扱うために、
例えば
$\prod_{i=1}^{n}\mathfrak{P}_i$
のような積軌道空間や、
関係ネットワークの時間発展を含む
結合軌道空間を導入することは可能である。

しかし、
Phase~1 の本節では
そのような多主体軌道空間を
正式には導入しない。

したがって、
多主体的な
閉塞経路集合、
崩壊到達経路集合、
同時到達時刻、
伝播速度、
連鎖崩壊集合等も
本節では定義しない。

これらは、
\S\ref{sec:org_acs}
の組織構造を前提にした
後続の拡張課題である。
\end{remark}

\subsubsection{組織ACSの節への引き渡しインターフェース}

\begin{remark}[組織ACSの節へ渡す個体レベルの確定事項]
\label{rem:sec7_to_sec8_interface}
\S\ref{sec:org_acs}
が個体レベルの失敗状態を参照するとき、
本節から引き渡す確定事項は
少なくとも次である。

\begin{enumerate}
  \item
  個体閉塞：
  $A_{i,t}=\emptyset$。

  \item
  個体崩壊：
  $\mathrm{Col}_{i,t}=1$、
  ただし
  本枠組みにおける
  (SC-1) 構造的同一性を成立させる組織関係の喪失、
  (SC-2) 同一個体としての相対的非回復性、
  および (C-2) 吸収性を伴う。

  \item
  個体レベルでは、
  本論文は閉塞 $A_{i,t}=\emptyset$ から
  $\mathrm{Col}_{i,t}=1$ を一般的な帰結として置かない。

  \item
  個体崩壊は
  組織水準の崩壊と同一ではない。

  \item
  個体の
  $\mathfrak{P}_{\mathrm{blk}}$
  と
  $\mathfrak{P}_{\mathrm{col}}$
  には、
  一般的な包含関係を置かない。

  \item
  内部量・関係状態の変化は、
  閉塞または崩壊の
  一般的十分条件ではない。
\end{enumerate}

\S\ref{sec:org_acs} は、
これらの個体レベル定義を変更せず参照し、
組織水準の構造を必要に応じて
別途定義する。
\end{remark}

\subsubsection{本節の責務境界}

\begin{remark}[本節の責務境界]
\label{rem:sec7_scope_boundary}
\S\ref{sec:collapse_cond} の責務は、
Phase~1 の単一ACS個体について、
閉塞と崩壊を区別し、
回復可能性と終端性を整理し、
それらを軌道記録集合として記述するところまでである。

本節は、
次のものを一般理論としては導入しない。

\begin{itemize}
  \item
  組織ACSそのものの成立条件、
  \item
  組織水準の構造的自己、
  \item
  組織水準の崩壊指標、
  \item
  個体間の崩壊伝播則、
  \item
  多主体軌道空間、
  \item
  多主体的な閉塞・崩壊経路集合、
  \item
  組織水準の回復作用素。
\end{itemize}

このうち、
組織自身をACSとして扱う組織ACSの成立条件および
組織水準の可行性・閉塞・崩壊は
\S\ref{sec:org_acs} の責務であり、
個体間の崩壊伝播則、
多主体軌道空間、
および多主体的な経路集合は、
Phase~2 以降の責務である（問題~\ref{prob:social_dynamics}、\S\ref{subsec:phase_outlook}）。

したがって、
\S\ref{sec:collapse_cond} から \S\ref{sec:org_acs} への接続は、
\[
  \text{個体レベルの確定定義を渡す}
\]
ことに限定し、
\[
  \text{個体定義を組織水準へ自動拡張しない}
\]
という境界を維持する。
\end{remark}


%================ Section 8 組織ACS ================
%======================================================================
\section{組織ACS：社会的組織への構造的適用}
\label{sec:org_acs}

本節では、複数の参加主体からなる社会的組織を対象とし、
組織自身についてACSの形式的構造が成立する場合を記述する。
組織にACSが参加することと、組織そのものがACSであることは区別される。
後者を扱うため、組織水準の構造を指定し、
先行節の定義を適用できる条件を明らかにする。

本節は既に識別される組織の抽象的構造を扱い、
その発生起源、組織化の過程または形成の必然性を扱わない。
組織固有の目的、参加主体の共通意思および特定の統治形態も、
組織ACSの一般的な成立要件とはしない。
本節の区別は個体ACSの関係対象の記述と、
組織自身の可行性の記述に用いられる。
本節で新たに用いる主要記号の一覧は付録~\ref{app:symbols_s8} に示す。

%----------------------------------------------------------------------
\subsection{集団・組織・社会の区分と組織ACS}
\label{subsec:org_terms}

\begin{definition}[集団・組織・社会の記述上の区分]
\label{def:org_terms}
\emph{集団}とは、識別規則の下で複数の個体を
一つの対象集合として扱う基体をいう。
集団の定義は、集団自身の機能的統合、
自己--環境境界またはACS構造の成立を要件として含まない。

\emph{組織}とは、区別される組織的境界の下で、
複数の参加主体と必要に応じた非主体的要素の寄与が
組織水準の機能を実現するように結び付けられ、
分担・調整される系をいう。
ここでいう組織は参加主体からなる社会的組織を指し、
生物の組織または物質の組織構造を指さない。
機能の分解、分担および調整の方式は実装依存とする。
組織の記述には、組織自身の整形式なACS構造を前提としない。

\emph{社会}とは、個体ACSその他の主体、集団および組織が
帰属・所属・活動する場をいい、
これらの主体・集団・組織、その他の要素、
およびそれらの間の関係を含みうる。
\end{definition}

\begin{definition}[ACS集団・ACS組織・組織ACS]
\label{def:org_class_terms}
二体以上の個体ACSからなる集団を\emph{ACS集団}と呼ぶ。
ACS集団の定義は、
定義~\ref{def:org_terms}の組織としての構造の成立を要件として含まない。

個体ACSが参加する組織を\emph{ACS組織}と呼ぶ。
組織のうち、その組織自身について定義~\ref{def:org_acs}の
適用条件が成立するものを\emph{組織ACS}と呼ぶ。
ACS組織の定義は、組織ACSであるための
定義~\ref{def:org_acs}の適用条件の成立を要件として含まず、
組織ACSの定義は、個体ACSの参加を要件として含まない。

時刻 $t$ に組織ACSとして扱われる組織 $\mathcal O$ のうち、
参加ACS数（式~\eqref{eq:org_participants}）が
$n_{\mathrm{org},t}(\mathcal O)\geq1$ であるものを\emph{個体ACS参加型組織ACS}、
$n_{\mathrm{org},t}(\mathcal O)=0$ であるものを\emph{個体ACS非参加型組織ACS}と呼ぶ。
この区分は参加の記録による区分であり、組織ACSであるための条件を加えない。
組織自身は、その組織の参加ACSに数えない。
\end{definition}

\begin{remark}[組織ACSの記述単位]
\label{rem:org_acs_unit}
本論文の組織ACSは、
複数の参加主体からなる組織（定義~\ref{def:org_terms}）を対象とする
組織水準の最小の記述単位とする。
参加主体が個体ACSであることは要件としない。
組織または組織ACSを構成要素とする組織は
これと区別し、Phase~1では扱わない。
\end{remark}

\begin{remark}[組織と関係対象分類]
\label{rem:org_terms_nonimplication}
\S\ref{sec:rel_min}の
$\mathrm{type}_i:\mathcal Y_{-i}\to
\{\mathrm{ACS},\mathrm{agent},\mathrm{obj}\}$ は、
個体ACSに相対する関係記述上の対象分類である。
外部対象集合 $\mathcal Y_{-i}$ は組織および組織ACSを除外しないため、
これらは行為主体として $\mathrm{ACS}$ または $\mathrm{agent}$ に
分類されうる。
ただし $\mathrm{type}_i$ の値は、
定義~\ref{def:org_acs}の時刻別の適用条件の成立・不成立を表さない。
組織ACSを関係対象とする記述の詳細、
とくにACS個体としての扱いを前提とする双方向合成
（仮定~\ref{ass:bidirect_scope}）の組織ACSへの適用は、
本論文（Phase~1）では扱わず、
Phase~1.$\alpha$ または Phase~2 に委ねる（\S\ref{subsec:phase1a}）。
\end{remark}

%----------------------------------------------------------------------
\subsection{組織水準の記述と参加ACS}
\label{subsec:org_structure}

組織水準の状態と機能は、当該組織の境界に相対して指定される。
個体ACSの状態を列挙したこと、個体間の関係記録が存在すること、
または組織内に役職名があることのみから、
本論文は組織自身の状態・認知機構・更新構造を構成しない。

\begin{definition}[参加ACSの時刻別記録]
\label{def:org_participation}
考察対象とする個体ACSの識別子の集合を $\mathbb{I}$、
考察対象とする組織の集合を $\mathfrak{O}$ とする。
$\mathbb{I}$ の有限性は仮定しない。
組織 $\mathcal O\in\mathfrak{O}$ と時刻 $t$ に対し、
その時刻に当該組織へACSとして参加する相異なる個体の
識別子集合を $\mathbb{I}_{\mathrm{org},t}(\mathcal O)\subseteq\mathbb{I}$ とする。
この集合を時刻 $t$ の参加ACS集合と呼び、
参加ACS数をその濃度
\begin{equation}
  n_{\mathrm{org},t}(\mathcal O)
  :=\bigl|\mathbb{I}_{\mathrm{org},t}(\mathcal O)\bigr|
  \label{eq:org_participants}
\end{equation}
で表す。
$\mathbb{I}_{\mathrm{org},t}(\mathcal O)=\emptyset$ も許す。
これは個体の区別と当該時刻の参加が指定される場合の記録であり、
すべての候補個体の通常ACS状態が全時刻に存在するとは仮定しない。

固定した有限の候補個体集合を用いるモデルでは、
追加の記述として
\begin{equation}
  P\subseteq\mathbb{I},\quad |P|<\infty,
  \qquad
  \mathbb{I}_{\mathrm{org},t}(\mathcal O)\subseteq P
  \label{eq:org_candidate_pool}
\end{equation}
と置いてよい。
$P$ の要素数に下限は置かない。
式~\eqref{eq:org_participants}の参加数の定義に、
固定候補集合 $P$ の存在は必要としない。
\end{definition}

\begin{remark}[構成員と機能担い手]
\label{rem:org_participants_functions}
組織の機能は、参加ACSのほか、非ACSエージェント、
装置、保存された情報その他の要素によって担われてもよい。
参加ACS数と、役割の割当ておよび機能の実効的充足は異なる。
全参加ACSの共通目的・共通意思、固定役職および
特定の階層構造は要求しない。
機能充足を参加ACSの人数閾値で表す方式も、一般要件とはしない。

組織の境界外にある供給要因と、境界内の機能担い手とは区別する。
前者が組織の存続に必要ならば、
\S\ref{sec:acs_def}の必要環境要因として記述する余地がある。

参加主体と非主体的要素の区別、
および各参加主体が個体ACSに該当するか否かは、対象モデルにおいて指定する。
本論文は、その区別を一般的に判定する手続きを与えない。
\end{remark}

\begin{remark}[関係記録と組織の機能]
\label{rem:org_relation_scope}
\S\ref{sec:rel_min}の関係記述子は、
参加ACSと他の対象との関係を記述しうる。
しかし本論文は、関係記録が非空であることだけから、
組織自身の制御機構、因果的統合またはACS性を一般的な帰結として置かない。
関係記録の変化も、他者の一意な因果的寄与とは解釈しない。
\end{remark}

%----------------------------------------------------------------------
\subsection{組織水準のACS構造と組織ACSの適用規約}
\label{subsec:org_formal}

\begin{definition}[組織水準のACS構造]
\label{def:org_structural_form}
組織 $\mathcal O$ 自身を記述対象とし、
状態空間、行動空間、必要環境要因、遷移、持続評価、
機構および更新を
\begin{equation}
  \Sigma_{\mathrm{org}}
  :=\bigl(
    X_{\mathrm{org}},U_{\mathrm{org}},E_{\mathrm{org}},
    f_{\mathrm{org}},\Pi_{\mathrm{org}},M_{\mathrm{org}},
    \mathrm{Update}_{\mathrm{org}}
  \bigr)
  \label{eq:org_7tuple}
\end{equation}
として与える。
当該7タプルが\S\ref{sec:acs_def}のACS定義、
構成要件および整形式性を満たし、
組織自身の自己--環境境界と機能の実現を記述する場合、
これを\emph{組織水準のACS構造}と呼ぶ。

特に、$X_{\mathrm{org}}\subseteq\mathbb R^d$ は
空でないコンパクトな組織状態空間とし、
外部環境状態空間 $E_{\mathrm{ext},\mathrm{org}}$ と
必要環境要因集合 $E_{\mathrm{org}}$ を区別する。
両者および遷移の型は
\begin{equation}
  E_{\mathrm{org}}
  \subseteq
  \{0,1\}^{X_{\mathrm{org}}\times E_{\mathrm{ext},\mathrm{org}}},
  \qquad E_{\mathrm{org}}\neq\emptyset
  \label{eq:org_environment_type}
\end{equation}
および
\begin{equation}
  f_{\mathrm{org}}:
  X_{\mathrm{org}}\times U_{\mathrm{org}}
  \times E_{\mathrm{ext},\mathrm{org}}
  \longrightarrow X_{\mathrm{org}}
  \label{eq:org_transition_type}
\end{equation}
とする。
$\Pi_{\mathrm{org}}:X_{\mathrm{org}}\to\mathbb R$、
$U_{\mathrm{org}}\neq\emptyset$ とし、
その他の型、構成要件 CR-1（必要環境整合性）・CR-2（原始的な機能と基盤機能の具備）および
整形式実現に関する条件は\S\ref{sec:acs_def}に従う。
これらの条件を記号の割当てだけで満たしたとは扱わない。
\end{definition}

\begin{remark}[組織状態と遷移の一価性]
\label{rem:org_state_closure}
組織状態 $x_{\mathrm{org}}\in X_{\mathrm{org}}$ は、
個体ACSの状態の和、平均または直積として
一般的に定まるものではない。
集約写像を用いるモデルでは、
同じ組織状態・組織行動・外部環境状態に対応する
記述が異なる組織次状態を与えないことを確認する必要がある。
そうでなければ式~\eqref{eq:org_transition_type}の
一価の遷移は定まらない。
組織状態を直接与えるモデルに集約写像を要求しない。
\end{remark}

\begin{definition}[組織ACSの記述的適用規約]
\label{def:org_acs}
組織 $\mathcal O$（定義~\ref{def:org_terms}）を
時刻 $t$ に\emph{組織ACS}として扱うのは、
当該組織自身を対象として、
定義~\ref{def:org_structural_form}の
整形式なACS構造が時刻 $t$ に実現し、
境界・機能・更新の対象が当該組織として識別される場合に限る。
参加ACS数 $n_{\mathrm{org},t}(\mathcal O)$ はこの適用条件に含まれない。
組織水準の構造の指定は、
一つの参加主体（個体ACSを含む）の活動を
組織と呼び換えるだけでは満たされない。
他方、個体ACSからの非還元性を判定する
モデル非依存の形式述語は、本規約には導入しない。
本規約は、組織ACSの一般的存在または
個別の実在組織への該当性を主張するものではない。
\end{definition}

\begin{remark}[成立と観察による判定]
\label{rem:org_acs_observability}
境界、継続性、環境との相互作用および
内部機能の変更が観察されることは、
組織水準のACS構造を検討する根拠となりうる。
本論文は、これらの観察事実のみから、
定義~\ref{def:org_acs}の全条件を一般的な帰結として置かない。
内部情報が観察できない場合、
ACS性を確認できないこととACS性が存在しないことを区別する。
\end{remark}

\begin{remark}[個体水準と組織水準の判定の区別]
\label{rem:org_individual_nonimplication}
参加個体の可行性・三原理の時刻別充足述語の真偽と、
組織自身のACS性・可行性・三原理の時刻別充足述語の真偽は、
異なる状態空間と評価対象に関する判定である。
個体量の単純な集計を組織量とする一般則は本節で与えない。
個体水準と組織水準の間の一般的な含意関係も、
追加の接続条件なしには主張しない。
\end{remark}

%----------------------------------------------------------------------
\subsection{組織水準の可行性と三原理}
\label{subsec:org_viability}

\begin{definition}[組織水準の可行集合と可行行動集合]
\label{def:org_viable_action}
定義~\ref{def:org_structural_form}の整形式な
組織水準のACS構造が与えられる場合、
\begin{equation}
  K_{\mathrm{org}}
  :=\{x\in X_{\mathrm{org}}\mid
  \Pi_{\mathrm{org}}(x)\geq0\}
  \label{eq:org_viable_set}
\end{equation}
とする。
時刻 $t$ に組織水準の通常ACS状態
$x_{\mathrm{org},t}\in X_{\mathrm{org}}$ が実現し、
$e_{\mathrm{org},t}\in E_{\mathrm{ext},\mathrm{org}}$ が
与えられる場合の可行行動集合を
\begin{equation}
  A_{\mathrm{org},t}
  :=\bigl\{u\in U_{\mathrm{org}}\mid
  f_{\mathrm{org}}\bigl(x_{\mathrm{org},t},u,
  e_{\mathrm{org},t}\bigr)\in K_{\mathrm{org}}\bigr\}
  \label{eq:org_viable_action}
\end{equation}
とする。
この集合は一時刻先の可行状態への遷移可能性を表し、
無期限の存続を保証しない。
\end{definition}

\begin{proposition}[組織水準の閉塞の特徴づけ]
\label{prop:org_blockage_characterization}
定義~\ref{def:org_viable_action}の量が
時刻 $t$ に整形式に定義されているとする。
このとき、
\begin{equation}
  A_{\mathrm{org},t}=\emptyset
  \quad\Longleftrightarrow\quad
  \forall u\in U_{\mathrm{org}},\;
  \Pi_{\mathrm{org}}\!\left(
  f_{\mathrm{org}}\bigl(x_{\mathrm{org},t},u,
  e_{\mathrm{org},t}\bigr)\right)<0
  \label{eq:org_blockage_characterization}
\end{equation}
が成り立つ。
\end{proposition}

\begin{proof}
式~\eqref{eq:org_viable_set}および
式~\eqref{eq:org_viable_action}より、
$u\in A_{\mathrm{org},t}$ は
$\Pi_{\mathrm{org}}(f_{\mathrm{org}}(
 x_{\mathrm{org},t},u,e_{\mathrm{org},t}))\geq0$
と同値である。
可行行動が存在しないことは、
いずれの $u\in U_{\mathrm{org}}$ に対しても
この非負条件が成立しないことと同値である。
$\Pi_{\mathrm{org}}$ は実数値なので、
その否定は式~\eqref{eq:org_blockage_characterization}の
厳密な負値条件である。
\end{proof}

\begin{remark}[可行性と組織ACS資格]
\label{rem:org_viability_vs_qualification}
$K_{\mathrm{org}}$ への帰属と、
定義~\ref{def:org_acs}の時刻別の組織ACS資格は
異なる述語である。
可行集合は $\Pi_{\mathrm{org}}$ から定義され、
組織ACS資格は時刻別の整形式な実現によって定まる。
両者の含意を主張する場合には、
組織状態・遷移と整形式な実現との接続条件が必要となる。
$A_{\mathrm{org},t}\neq\emptyset$ のみから
遷移後の組織ACS資格を保証しない。
\end{remark}

\begin{remark}[先行節の定義の組織水準への適用]
\label{rem:org_principles_conditional}
組織水準のACS構造 $\Sigma_{\mathrm{org}}$ は、
定義~\ref{def:org_structural_form} により
\S\ref{sec:acs_def}のACS定義、構成要件および整形式性を満たすので、
抽象構造型の一つの実現である（\S\ref{subsec:acs_7tuple} の記述水準の規約）。
したがって、
\S\ref{sec:acs_def}の閉塞・機能的終端・崩壊、
\S\ref{sec:three_principles}の三原理SM・EU・FRの時刻別充足述語とT-STAB、
および\S\ref{sec:collapse_cond}の状態帯・外的回復可能性・軌道記録と経路集合は、
個体添字を組織 $\mathcal O$ に置き換えて $\Sigma_{\mathrm{org}}$ にそのまま適用される。
たとえば、SMの現在状態における条件は
\begin{equation}
  \Pi_{\mathrm{org}}\bigl(x_{\mathrm{org},t}\bigr)\geq0
  \label{eq:org_SM}
\end{equation}
と書ける。
条件 (SC-1)・(SC-2)・(FT-1)・(FT-2) と外的回復の判定規則は、
個体の場合と同じく対象系ごとに固定する。
組織については、組織としての同一性を成立させる組織関係と、
同一組織の判定基準がその判定規則に含まれる。

\S\ref{sec:acs_def}でいう「ACS個体」は抽象構造型の実現を指し、
本節で参加者を指す「個体ACS」とは区別される。
組織ACSは、前者の意味ではACS個体の一つである。

一方、参加する個体ACSの量
（一段階臨界必要要因の充足、局所FRの真偽、閉塞、崩壊等）を
組織水準の量へ転写または集約する規則は与えない。
組織にこれらの定義を適用できることと、
各時刻・各経路で原理を実際に充足することは異なる。
組織水準の時刻別充足述語が偽であることから導かれるのは、
\S\ref{sec:three_principles}の命題が示す帰結
（たとえばEUの偽からの閉塞、命題~\ref{prop:eu_blockage}）に限られ、
それだけから組織ACS資格の喪失または崩壊は結論しない。
\end{remark}

%----------------------------------------------------------------------
\subsection{資格・時間的同一性と失敗状態の適用範囲}
\label{subsec:org_identity}

\begin{remark}[時刻別資格と組織の継続]
\label{rem:org_qualification_loss}
参加ACS数の変化
（たとえば $n_{\mathrm{org},t}(\mathcal O)$ の減少や $0$ への到達）のみから、
定義~\ref{def:org_acs}の組織ACS資格の喪失、
組織の構造的自己の喪失、または崩壊を結論しない。
他方、参加主体の離脱によって組織自身のACS構造が実際に失われる場合は排除しない。
その判定は、組織水準の整形式な実現について行う。
組織ACS資格の喪失と組織の崩壊は、別の判定である（備考~\ref{rem:org_failure_scope}）。

組織ACS資格がいったん成立しなくなった後に再び成立しても、
それだけで同一組織ACSが復帰したとはいえない。
同一組織としての復帰を扱う場合には、
対象モデルで許容する境界・機能・更新の継続条件に照らし、
旧組織からの継承を別途確認する。
一律の復帰期限は定めない。
名称の一致や内部参加者・外部観測者の認識のみを
構造的同一性の十分条件とはしない。
\end{remark}

\begin{remark}[閉塞・崩壊の適用域]
\label{rem:org_failure_scope}
組織自身について通常のACS状態および
式~\eqref{eq:org_viable_action}が定義される時刻には、
命題~\ref{prop:org_blockage_characterization}によって
組織水準の可行行動集合の空性を判定できる。
これは\S\ref{sec:acs_def}の閉塞（定義~\ref{def:blockage_sec2}）を
$\Sigma_{\mathrm{org}}$ に適用したものである。
閉塞、可行集合からの逸脱、整形式実現の不成立、
組織ACS資格の喪失を同一視しない。

組織水準の崩壊は、
\S\ref{sec:acs_def}の定義~\ref{def:collapse}の条件 (SC-1)・(SC-2) を、
備考~\ref{rem:org_principles_conditional}の判定規則と、
組織について定めた許容外的作用の集合（\S\ref{sec:collapse_cond}）のもとで判定する。
組織ACS資格と参加ACS数は崩壊の定義に含まれないので、
参加ACS一体の離脱・崩壊や参加ACS数の減少のみから、
組織の崩壊は結論しない。
また、組織の分割のように、
同一組織の継続として対応づけられる後続の実現が複数になる場合は、
\S\ref{sec:acs_def}の仮定~\ref{ass:no_branching}の対象外である。
\end{remark}

%----------------------------------------------------------------------
\subsection{適用限界}
\label{subsec:org_scope}

本節は、既存の社会的組織について
ACSの構造的定義を適用できる条件を記述する。
組織ACSの非空な具体モデルの構成、
実在する任意の組織についての全条件の実証、
または組織化からACS構造が生じる必然性は主張しない。
特定の実装によらない非還元性の形式判定規則、
組織の同一性を判定する普遍的な手続、
および構成員の生成・離脱を含む一般的社会動態も扱わない。

組織ACSとしての成立は、
目的達成、参加個体すべての可行性、
三原理の実際の充足または倫理的責任の成立を意味しない。
ACSその他の主体のあいだの持続可能な協調の成立条件は、
ECITとの正式接続としてPhase~3で検討する（\S\ref{subsec:phase_outlook}、問題~\ref{prob:ecit_full}）。


%================ Section 9 限界・未解決問題・Phase 2・3への展望 ================
%======================================================================
\section{限界・未解決問題・Phase~2・3への展望}
\label{sec:limitations_and_outlook}

%----------------------------------------------------------------------
\subsection*{導入}
\label{subsec:s9_intro}

本節では、Phase~1（決定論的定式化；本論文）で構成した
ACS理論の理論的限界を体系的に開示し、未解決問題を整理し、
Phase~2およびPhase~3への研究方向を記述したうえで結論を述べる。
本節は新たなACS構成要件、三原理、TIPA機構、関係カテゴリー、
閉塞・崩壊条件または組織ACSの適用条件を追加する節ではない。
既節で与えた定義・命題・仮定の射程を変更せず、
それらが現在の枠組みのもとで何を主張していないかを明示する。

構成は、\S\ref{subsec:limitations_list}（理論的限界）、
\S\ref{subsec:phase1a}（Phase~1.$\alpha$ への委譲事項）、
\S\ref{subsec:open_problems}（未解決問題）、
\S\ref{subsec:phase_outlook}（Phase~2・3への展望）、
\S\ref{subsec:conclusion}（結論）である。

%----------------------------------------------------------------------
\subsection{理論的限界}
\label{subsec:limitations_list}

以下では、\S\ref{sec:acs_def}から\S\ref{sec:org_acs}までに
明示されたPhase~1の境界と、現在なお残る未確定部分を
L1--L18の連番で整理する。

\paragraph{L1：決定論的・離散時間・コンパクト空間への限定}
\label{lim:L1}
Phase~1は、決定論的・離散時間の状態遷移と、
コンパクトな状態・行動・内部モデル等の空間を中心とする枠組みである。
したがって、確率的遷移、確率分布、期待値、連続時間極限、
ジャンプ過程、非有界空間および可変次元状態空間を
一般理論としては扱わない。
この限定は、存在・閉性・コンパクト性等を扱う際の技術的明確性を与える一方、
現実の確率的・連続時間的・非有界な系への直接適用を保証しない。

\medskip
\noindent\textit{起源：\S\ref{sec:acs_def}および論文全体のPhase~1規約。
Phase~2：確率的遷移、連続時間化、局所コンパクト化その他の弱化を検討する。}

\paragraph{L2：外部環境状態の遷移・到達可能性（環境動学）の限定的な扱い}
\label{lim:L2}
\S\ref{sec:acs_def}では、ACSの存続に不可欠な環境要因を
必要環境空間$E$として与え、背景型として外部環境状態空間
$E_{\mathrm{ext}}=E^{\mathrm{bg}}\times O$を置く。
Phase~1では、当該ACSの行動に応答しうる行動応答成分の遷移を、
環境実現の側の写像$g_O$として与えるにとどめる
（定義~\ref{def:env_realization}）。
外生背景成分の列の生成則、環境状態間の到達可能性、
環境の確率的変動、他のACS個体との結合動学、
およびACSと環境の共進化は、
独立の環境動学としては定式化しない。
したがって、異なる環境状態の存在可能性と、
その状態へ実際に到達する動学は区別される。
自己--環境境界の再編成と外部環境状態空間の時間依存はL16で扱う。

\medskip
\noindent\textit{起源：\S\ref{sec:acs_def}の必要環境空間の構造的解釈と外部環境状態空間の分解。
Phase~2：背景成分の生成則、環境遷移の可到達性、確率的環境、他個体との結合動学を導入する。}

\paragraph{L3：認知的処理能力の未形式化}
\label{lim:L3}
Phase~1は認知的処理能力を形式的に定義せず、
単一の量としても導入しない。
\S\ref{sec:acs_def}の備考~\ref{rem:cognitive_capacity_policy}は、
量的容量、時間応答性、統合の質を区別すべきこと、
容量と実現量を区別すべきこと、
および能力の使い方に現れる偏向性とその程度（偏向強度）を能力に含めないことを
記述方針として述べるにとどまる。
したがって、認知的処理能力の高低を比較する指標、
三つの区分の測定、およびそれらと可行性・閉塞・崩壊との関係を、
Phase~1の結果からは論じられない。

\medskip
\noindent\textit{起源：\S\ref{sec:acs_def}、
備考~\ref{rem:cognitive_capacity_policy}。
Phase~1.$\alpha$／Phase~2：三つの区分の形式化
（時間応答性については所定の基準に到達するまでの段数を候補とする）と、
その測定を検討する。}

\paragraph{L4：同一個体性における非分岐性}
\label{lim:L4}
仮定~\ref{ass:no_branching}は、
時刻 $t$ のACS個体実現に対して、
同一個体の継続として対応づけられる時刻 $t+1$ の
ACS個体実現を高々一つとし、
コピー個体または複数の後続実現が生じる分岐を
Phase~1の対象から除外する。

したがって、完全複製、分岐的自己同一性、
複数の後続実現への同一性継承等を含む系については、
時間的自己同一性関係をそのまま適用することはできない。
組織ACSでは、組織の分割・分社のように、
同一組織の継続として対応づけられる後続の実現が複数になる場合がこれに当たる。
なお、この非分岐性は状態遷移写像 $f$ の一価性とは
別個の個体同一性上の仮定である。

\medskip
\noindent\textit{起源：\S\ref{sec:acs_def}、
仮定~\ref{ass:no_branching}、
定義~\ref{def:time_identity}。
Phase~2：個体実現間の時間的継続関係、
集合値対応または系譜構造等を候補として、
分岐を許す時間的自己同一性を再定式化する。}
\paragraph{L5：T-STABの閾値較正と推論射程}
\label{lim:L5}
\S\ref{sec:three_principles}のT-STAB
（定義~\ref{def:tstab}）は、
SM評価の下限、局所FRとEUの窓内充足、およびSM評価の一時刻変動を
有限時間窓上で統合する三原理時間統合評価条件である。
これは力学系論におけるLyapunov安定性、漸近安定性、
または構造安定性を意味するものではない。
下限参照値 $\underline{s}_{\mathrm{SM}}$、
変動上限 $\delta_{\max}$、
および窓幅 $W_T$（定義~\ref{def:tstab_params}）は
対象ACSの評価構造に相対する参照パラメータであり、
Phase~1では所与として扱われる。
その定量的較正則は一般には与えられていない。

また本論文は、T-STAB条件の違反を、
閉塞 $A_{i,t}=\emptyset$、崩壊 $\mathrm{Col}_{i,t}=1$、
負の関係品質を伴う環境結合、完全孤立、
またはPriorityの発火の十分条件として置かない。
これらへの接続を述べるには追加条件を要する。
したがってT-STABは、
Phase~1では三原理に関係する理論上の状態を有限窓で
記述・診断する条件であり、
一般的な収束、漸近安定性または長期予測を
単独で与えるものではない。

また、Phase~1は、FRの長期の側面について、
独立した無限時間判定条件を与えない。
定義~\ref{def:FR}は各時刻の局所FRを与え、
T-STABは有限窓内におけるその充足を評価する。
したがって本論文は、T-STABの成立から、
有限窓外における局所FRの持続、
または無限時間にわたる応答多様性の維持を一般的な帰結として置かない。
長期の側面の本格的な理論は、\S\ref{subsec:phase1a}に挙げる。

\medskip
\noindent\textit{起源：\S\ref{sec:three_principles}、
定義~\ref{def:FR}、備考~\ref{rem:sd_local_scope}、
定義~\ref{def:tstab_params}、定義~\ref{def:tstab}、
命題~\ref{prop:tstab_implies_principles}、
備考~\ref{rem:conv_tstab_scope}および備考~\ref{rem:tstab_descriptive_status}。
Phase~2：閾値・窓幅の較正則、パラメータ感度、
ロバスト性、およびT-STAB違反から各失敗状態への
追加条件付き接続を実装別に検討する。}

\paragraph{L6：TIPAの具体的実装写像と推定誤差の定量化}
\label{lim:L6}
\S\ref{sec:tipa}は、M1・M2・M3、PSLおよびISLの
機能的構造と最小要件を与えるが、各実装における
信号変換、評価、Priority生成、応答選択等の
具体的関数形を一意には固定しない。
また、M1の状態・環境推定 $\hat x_t$、$\hat e_t$ と
実現値 $x_t$、$e_t$ との誤差、
およびM2の代理評価と原理水準の真値との誤差について、
Phase~1は普遍的な誤差上界、収束性または一致を要求しない
（備考~\ref{rem:m1_estimate_nontruth}、備考~\ref{rem:error_bounds}）。
したがって、推定・評価の誤差が各実装でどの程度に収まるか、
その時間変化、および誤差が可行行動集合や閉塞到達へ与える影響は、
Phase~1の結果からは定まらない。
TIPAはPhase~1では構成的・機能的形式を与えるが、
特定の生物学的機構またはAI実装を一意に同定する理論ではない。
また、構造的分離 $B_{SE}$ とM1-bの自己・環境分類の不整合を検証する機構
（備考~\ref{rem:xi_se_interpretation}）、
同一の $f$ に対応しうる内部モデル構造 $M$ の一意性
（備考~\ref{rem:M_independence}）、
内部表現の空間 $\mathcal{R}_{\mathrm{int}}$ から $\mathcal{M}$ への符号化写像
（備考~\ref{rem:M_encodes_representation}）、
および $f_{\mathrm{auto}}$、$f_{\mathrm{dir}}$、$\oplus$ の具体的実装形式
（仮定~\ref{ass:f_decomposition}）も定めていない。

\medskip
\noindent\textit{起源：\S\ref{sec:tipa}、
備考~\ref{rem:m1_estimate_nontruth}、備考~\ref{rem:error_bounds}。
Phase~2：実装別写像の具体化、誤差上界を導入するための条件、
推定誤差の動学、ロバスト性および実証的較正を検討する。
整合性検証機構、$M$ の一意性および $f$ の分解の実装形式も
Phase~2以降の課題とする。
符号化写像は、後続Phaseで導入することを約束しない。}

\paragraph{L7：\S\ref{sec:extended_dynamics}更新写像の抽象性と変数間結合の未確定}
\label{lim:L7}
\S\ref{sec:extended_dynamics}は、
$\boldsymbol{\beta}_t$、$d_{\mathrm{eff},t}$、$\pi_{\tau,t}$、
$\widehat{\tau}_{t}$および四健全度等について、
固定した読み出し規則により、状態・内部モデル・外部環境状態、
または作動記録の有限区間から定まる値として扱い、
その時刻$t$から$t+1$への変化は、$f$、Update、および外部環境の実現（背景成分の列と $g_O$）を通じて定まる。
独立した更新写像は置かない。
一方、具体的学習率、回復率、収束速度、変数間の相互依存関数、
長期安定性および一般的な閉形式解は固定しない。
また、$\widehat{\tau}_{t}=\pi_{\tau,t}\tau$は派生量であり、
独立した更新則をもたない。
したがって、\S\ref{sec:extended_dynamics}は時間方向の最小構造を与えるが、
特定実装の詳細な学習・劣化・回復モデルを与えるものではない。

\medskip
\noindent\textit{起源：\S\ref{sec:extended_dynamics}。
Phase~2：具体的更新則（予測と実現の差にもとづく内部モデルの更新の規則を含む）、結合構造、収束・安定性・回復条件をモデル別に検討する。
偏向性については、各偏向性の参照基準と尺度の具体化、偏向強度の読み出し方式、TIPA・PSL・ISL による形成の記述、形成の組み合わせと帰属区分の対応、観測者モデルに基づく帰属区分の同定、実装上の偏向性の変容・統合と抽象的構造記述の添字との対応づけ、および各偏向性の方向・内容の表現を検討する
（\S\ref{sec:extended_dynamics}、定義~\ref{def:bias_attribution}、備考~\ref{rem:beta_multiple_biases}）。}

\paragraph{L8：関係記述子の非因果性と測定・分類の実装依存性}
\label{lim:L8}
\S\ref{sec:rel_min}の関係記述子は、
対象間の関係構造と、関係継続窓において観測される
可行性関連量の変化を記述する分類Cの分析形式である。
有限窓差分
$\Delta_L\mathrm{sz}_i(A_{i,t})$、
$\Delta_Ld_{\mathrm{eff},i,t}$、
$\Delta_L\widehat{\tau}_{i,t}$は、
対象$Y_j$の一意な因果的寄与を意味しない。
また、局所的サイズ指標$\mathrm{sz}_i$の具体的構成、
カテゴリー判定閾値、持続時間、複合カテゴリー、因果帰属等は
一般理論として固定していない。
局所的サイズ指標の測度論的な精密化
（備考~\ref{rem:finite_window_diff_limit}）も行わない。

\medskip
\noindent\textit{起源：\S\ref{sec:rel_min}。
Phase~1.$\alpha$：分類・測定規則の詳細化。
Phase~2：介入・反実仮想等を含む因果的拡張を必要に応じて検討し、
局所的サイズ指標の測度論的な精密化を扱う。}

\paragraph{L9：閉塞・崩壊の観測可能性、回復判定および経路動学}
\label{lim:L9}
\S\ref{sec:collapse_cond}では、
閉塞$A_{i,t}=\emptyset$と
崩壊$\mathrm{Col}_{i,t}=1$を明確に分離し、
閉塞経路集合・崩壊到達経路集合・到達時刻を定義した。
しかし、有限観測だけから崩壊を一般に判定できるとは主張せず、
外的回復可能性についても実装依存の決定論的判定規則を抽象化している。
さらに、経路集合への流入率、到達確率、平均到達時間、
回復確率等はPhase~1では定義しない。
整形式実現指標 $\mathrm{WF}_{i,t}$ の値を有限観測から判定する操作的手続き
（定義~\ref{def:wf_realization}）と、
条件 (FT-1)・(FT-2)・(SC-1)・(SC-2) の判定規則の一般形式
（定義~\ref{def:functional_termination}、定義~\ref{def:collapse}）も与えていない。

\medskip
\noindent\textit{起源：\S\ref{sec:collapse_cond}。
Phase~2：観測可能性、判定アルゴリズム、確率的到達・回復解析、回復余裕の定量化を検討する。
$\mathrm{WF}_{i,t}$ の操作的判定手続きもPhase~2の課題とする。
(FT-1)・(FT-2)・(SC-1)・(SC-2) の判定規則は、
Phase~1では対象系ごとに固定されるものとして扱う。}

\paragraph{L10：参加ACS集合の時刻別記録と組織同一性動学の未形式化}
\label{lim:L10}
\S\ref{sec:org_acs}では、
組織 $\mathcal O$ の参加ACS集合
$\mathbb{I}_{\mathrm{org},t}(\mathcal O)$ と参加ACS数
$n_{\mathrm{org},t}(\mathcal O)$ を時刻ごとに記録する。
固定した有限候補集合 $P$ は、採用するモデルで必要な場合に限る
追加的な記述であり、Phase~1の一般要件ではない。
したがって、Phase~1は構成員数を時間的に固定する理論ではない。

一方、参加ACS集合そのものの一般的な遷移則、
加入、離脱、生成、消滅等の動学、
およびそれらの変化を通じた組織同一性の保持条件は定式化しない。
また本論文は、参加ACS数の変化のみから、
組織ACS資格の成立・不成立、組織の崩壊、
または同一組織ACSの復帰を一般的な帰結として置かない。

\medskip
\noindent\textit{起源：\S\ref{sec:org_acs}、
定義~\ref{def:org_participation}、
備考~\ref{rem:org_qualification_loss}および\S\ref{subsec:org_scope}。
Phase~2：参加集合の遷移、加入・離脱、
組織同一性および復帰条件の動学を対象モデルごとに検討する。}

\paragraph{L11：組織水準7タプルの存在・構成・一価遷移は一般定理ではない}
\label{lim:L11}
\S\ref{sec:org_acs}の
組織水準のACS構造
$\Sigma_{\mathrm{org}}$ は、
組織自身について
\S\ref{sec:acs_def}のACS定義、構成要件および整形式性が
満たされる場合に用いる構造形式である。
任意の社会的組織について、
$X_{\mathrm{org}}$、$U_{\mathrm{org}}$、
$E_{\mathrm{org}}$、$f_{\mathrm{org}}$、
$\Pi_{\mathrm{org}}$、$M_{\mathrm{org}}$、
$\mathrm{Update}_{\mathrm{org}}$ が
自動的に整形式に構成できることを証明するものではない。

また、組織状態は個体状態の和、平均または直積として
一般的に定まるものではなく、
集約写像を用いる場合には
同一の組織状態・組織行動・外部環境状態から
一価の組織次状態が定まることを別途確認する必要がある。
組織ACSの適用規約も、一般的存在または
任意の実在組織への該当性を主張しない。

\medskip
\noindent\textit{起源：\S\ref{sec:org_acs}、
定義~\ref{def:org_structural_form}、
備考~\ref{rem:org_state_closure}および定義~\ref{def:org_acs}。
Phase~2：具体的組織モデルに対する存在条件、
整形式性、一価遷移、識別可能性および実証的判定を検討する。}

\paragraph{L12：組織水準の判定規則と個体--組織集約則}
\label{lim:L12}
\S\ref{sec:org_acs}では、
組織自身に整形式なACS構造が成立する場合に限って、
先行節の三原理、閉塞・崩壊、時間統合評価および経路集合の定義を
組織水準へ適用する（備考~\ref{rem:org_principles_conditional}）。
その適用には、条件 (SC-1)・(SC-2)・(FT-1)・(FT-2) の判定規則を
組織について固定することを要する。
その一般形式は、個体の場合と同じく与えていない。

個体の可行性、三原理の時刻別充足述語の真偽、関係記述、健全度等を
平均・最小・加重平均その他の作用素によって
組織量へ変換する普遍的な集約則は与えない。
また、組織ACS資格、可行性、閉塞、
整形式実現の不成立および崩壊は
相互に同一視されない。
組織の分割・分社のように後続の実現が複数になる場合は、
非分岐性の仮定（L4）のもとでは扱えない。

\medskip
\noindent\textit{起源：\S\ref{sec:org_acs}、
備考~\ref{rem:org_individual_nonimplication}、
備考~\ref{rem:org_principles_conditional}および
備考~\ref{rem:org_failure_scope}。
Phase~2：個体--組織間の集約則、
組織水準の三原理評価、時間統合、
閉塞・崩壊経路および組織動態を個別モデルとして検討する。}

\paragraph{L13：ドメイン横断的対応と実証的妥当性の未検証}
\label{lim:L13}
本論文は、人間・生物、AI・ロボット、組織等へ
適用可能な抽象構造を意図しているが、
個別ドメインにおける7タプル各成分、三原理、TIPA、PSL・ISL、
関係記述子、閉塞・崩壊条件等の測定可能性と実証的妥当性を
本論文内で検証していない。
したがって、形式的適用可能性と経験的妥当性を区別する必要がある。
既存理論との関係についても、本論文が主張するのは
明示した範囲での整合性であり、包含・優越・一般化を意味しない。
TIPA成分の更新の質を表す量（定義~\ref{def:update_quality}）の
実証的な測定・推定方法も与えていない。

\medskip
\noindent\textit{起源：\S\ref{sec:intro}および全節の適用範囲。
Phase~2以降：ドメイン別操作化、測定、比較研究および反証可能性の検討を行う。}

\paragraph{L14：ECITフルモデルおよび規範理論はPhase~1の外部にある}
\label{lim:L14}
本論文は、倫理を善悪判断そのものとして定義するのではなく、
可行性保存に関する構造的制約として記述する。
したがって、愛・智慧・責任等のECIT諸要素、
倫理的協調応答の完全な発動条件、社会的規範の生成・変化、
具体的価値選択および政策的最適化をPhase~1の形式的帰結としては主張しない。
これらはACSの記述理論とは論理的地位を分けて扱う必要がある。

\medskip
\noindent\textit{起源：\S\ref{sec:intro}および\S\ref{sec:org_acs}の
記述理論・適用範囲としての限定。
Phase~3：ECITとの正式接続を、
Phase~1・2の結果と区別して検討する。}

\paragraph{L15：行動の実行の成否の区別}
\label{lim:L15}
\S\ref{sec:acs_def}の状態遷移写像の読みの規約
（備考~\ref{rem:f_reading}）のもとでは、
行動の試行、行動を実行できなかった場合、
および実行の成功の度合いを、$f$ の記述の中で区別しない。
行動入力の実行可能性を $U$ とは別に表す構造も置かないため、
$U$ の元をある時点で実行できないことは
$f$ が与える次状態に吸収され、
その行動が当該ACSの行動空間に含まれないこととは
別の形式対象として表されない。
また、作用差（定義~\ref{def:action_difference}）の大きさ、
多段階にわたる作用、行動や行動応答成分の寄与の比較、
および外部と内部に共通する資源量の収支も定量化しない。

\medskip
\noindent\textit{起源：\S\ref{sec:acs_def}、備考~\ref{rem:f_reading}。
Phase~2：実行の成否または程度を表す変数、
および時点ごとに実行可能な行動を $U$ とは別に表す構造を検討する。
作用差の大きさ等の定量化は、後続Phaseへの委譲事項または
将来の拡張・再検討の候補として保留する（\S\ref{sec:acs_def} の末尾の列挙）。}

\paragraph{L16：個体の構造の時間変化}
\label{lim:L16}
Phase~1は、個体の7タプル $\Sigma_i$ を時刻に依存しない構造体として扱い、
時刻とともに変わるのは状態、内部モデル状態その他の時刻別の値に限る
（\S\ref{sec:acs_def}、\S\ref{subsec:acs_7tuple} の記述水準の規約、定義~\ref{def:s_struct}）。
したがって、行動空間 $U$ そのものの拡大（新たな行動の獲得）または縮小、
内部モデルの機構構成の変化、
および状態空間の変化を、
同一個体の時間発展として記述しない。
可行集合を定める閾値の時間依存（定義~\ref{def:Pi}）、
行動空間 $U$ のラベル構造が位相構造に与える影響
（備考~\ref{rem:action_space_structure}）も扱わない。
また、7タプルの成分ではないが、
自己--環境境界の動的再編成（備考~\ref{rem:TSC_character}）、
時間依存する外部環境状態空間（仮定~\ref{ass:Eext_time_fixed}）、
および個体別の外部環境状態空間を相互に対応づける構造
（備考~\ref{rem:Eext_relative}）も扱わない。
他のACS個体との結合動学はL2に含める。

\medskip
\noindent\textit{起源：\S\ref{sec:acs_def}、\S\ref{subsec:acs_7tuple} の記述水準の規約、
備考~\ref{rem:TSC_character}、備考~\ref{rem:Eext_relative}、
備考~\ref{rem:action_space_structure}、仮定~\ref{ass:Eext_time_fixed}、定義~\ref{def:Pi}。
Phase~2：時刻に依存する個体の構造と、
その変化のもとでの時間的自己同一性を検討する。
動的閾値への拡張はPhase~2以降の課題とする。}

\paragraph{L17：必要環境要因の操作的同定・分解・判定}
\label{lim:L17}
\S\ref{sec:acs_def}は、必要環境要因を充足判定写像 $\eta\in E$ として与え、
構成要件CR-1（\S\ref{subsec:environment_space}）によってその数学的地位を定める。
一方、実際の対象から要因を同定し、その充足を判定する手続きは与えない。
具体的には、
必要環境要因の経験的同定手続きと同定誤差の扱い
（備考~\ref{rem:necessity_identification}）、
外部依存が一部の内部状態においてのみ現れる要因を要因と認める範囲とその程度、
CR-1の充足を有限観測から判定する手続き、
要因の同定が外部背景の取り方に依存する程度
（備考~\ref{rem:cr1_relativity}）、
要因の分解の一意性を与える追加条件
（備考~\ref{rem:factor_nonuniqueness}）、
および要因の分類粒度（備考~\ref{rem:eta_coarse_classes}）を定めていない。
必要性の存立と同定可能性は区別される。

\medskip
\noindent\textit{起源：\S\ref{sec:acs_def}、\S\ref{subsec:environment_space}、
備考~\ref{rem:necessity_identification}、備考~\ref{rem:cr1_relativity}、
備考~\ref{rem:factor_nonuniqueness}、備考~\ref{rem:eta_coarse_classes}。
Phase~2：経験的同定手続き、同定誤差、外部依存の範囲と程度、
有限観測によるCR-1の判定、外部背景への依存、分解の一意性の条件を検討する。
分類粒度の確定は、後続Phaseへの委譲事項または
将来の拡張・再検討の候補として保留する。}

\paragraph{L18：三原理に関わる補助量と解析条件の精密化}
\label{lim:L18}
\S\ref{sec:three_principles}は、三原理の時刻別充足述語と
その条件付き解析に必要な範囲で補助量を定める。
次の精密化は行わない。
環境応答成分 $I_E$ の同定に関する、観測ノイズ・時間遅延・内部フィードバック等を考慮した
操作的・因果的な決定基準（備考~\ref{rem:env_response_observation}）、
境界点における $x_E$ 方向の局所感応性の方向解析（備考~\ref{rem:pi_E_local_sensitivity}）、
$\Pi$ を補完する診断指標の動学的利用（備考~\ref{rem:diagnostic_role}）、
補題~\ref{lem:finite_exit} の履歴依存方策への一般化と離脱時刻の一様な有界性の条件
（備考~\ref{rem:uniform_bound}）、
応答多様度の値と閉塞到達時間との定量的関係（備考~\ref{rem:sd_degree}）、
$\mu_{\mathrm{eff}}$ の測度論的な精密化（備考~\ref{rem:mu_measure}）、
および仮定~\ref{ass:theta_exterior_blockage} を内部動学から導くための十分条件である。

\medskip
\noindent\textit{起源：\S\ref{sec:three_principles}、
備考~\ref{rem:env_response_observation}、備考~\ref{rem:pi_E_local_sensitivity}、
備考~\ref{rem:diagnostic_role}、備考~\ref{rem:uniform_bound}、
備考~\ref{rem:sd_degree}、備考~\ref{rem:mu_measure}、
仮定~\ref{ass:theta_exterior_blockage}。
Phase~2：上記の各項目を検討する。}

%----------------------------------------------------------------------
\subsection{Phase~1.$\alpha$ への委譲事項}
\label{subsec:phase1a}

本論文でいうPhase~1.$\alpha$は、
Phase~1の前提（決定論的・離散時間・コンパクト空間・非分岐性・単一ACS中心）と
各対象の型を変えずに、
Phase~1で値域・型・担体だけを与えた対象について、
規則・閾値・操作的基準を具体化する補助的詳細化の段階をいう。
Phase~1の仮定を弱める拡張はPhase~2に、
ECITとの接続はPhase~3に属する
（\S\ref{subsec:phase_outlook}）。
本文で「Phase~1.$\alpha$またはPhase~2」とした事項は、
決定論的に扱える部分をPhase~1.$\alpha$で、
それ以外の部分をPhase~2で検討する。
以下は検討課題の所在を示すものであり、
各事項がPhase~1.$\alpha$で形式化されることを主張しない。

\begin{enumerate}
  \item \textbf{構成的必要機能の候補：}
  必要環境利用機能（備考~\ref{rem:env_utilization_function}）、
  環境変化への横断的応答能力（備考~\ref{rem:env_change_response}）。
  \item \textbf{機構の構成：}
  TIPAの機構統合を許す再定式化（観察~\ref{obs:tipa_integration}）、
  保護的な働きの定まらない不随意の局所応答を含む、PSL層の局所自動応答一般
  （\S\ref{subsubsec:psl4} のPSL-4の範囲の外にあるもの）。
  \item \textbf{認知的処理能力：}
  量的容量・時間応答性・統合の質の形式化（備考~\ref{rem:cognitive_capacity_policy}）。
  \item \textbf{三原理：}
  行動応答成分の応答を考慮した応答不能環境集合・応答多様度・条件付き解析の版
  （備考~\ref{rem:same_time_alternatives}）、
  FRの長期の側面（備考~\ref{rem:sd_local_scope}）、
  \S\ref{sec:acs_def}の仮定と構成要件を満たすモデルによる
  時刻別充足述語の非含意の構成（備考~\ref{rem:principle_implication_scope}）。
  \item \textbf{TIPAの内容：}
  要請の質の列挙と分類（定義~\ref{def:request_signal_types}）、
  手段と目的の関係（備考~\ref{rem:purpose_hierarchy}）、
  目的を達成せずに終える理由の区別（定義~\ref{def:purpose_status}）、
  評価機構の迂回・置換の由来の追跡（備考~\ref{rem:self_attribution_scope}）、
  内部の対象トークンと外部の枠・個物の対応（備考~\ref{rem:Eext_decomposition_reading}）。
  \item \textbf{時刻間の量と記録：}
  実効的発展自由度の算出規則（定義~\ref{def:deff_section5}）、
  偏向性の形成の機構別の記述と帰属区分の同定（備考~\ref{rem:bias_formation_vs_attribution}）、
  保持記録が一致するための $\mathrm{Update}$ の詳細型の条件
  （定義~\ref{def:s5_record_agreement}、備考~\ref{rem:update_type_two_layers}）。
  \item \textbf{関係性：}
  \S\ref{sec:rel_min}の備考~\ref{rem:phase1a_delegation}の各項、
  局所的サイズ指標の具体化と有限近似（備考~\ref{rem:finite_window_diff_limit}）、
  主要方向の定量的判定規則（備考~\ref{rem:dir_undefined}）。
  \item \textbf{組織ACS：}
  組織ACSを関係対象とする記述と双方向合成の適用（備考~\ref{rem:org_terms_nonimplication}）。
\end{enumerate}

%----------------------------------------------------------------------
\subsection{未解決問題}
\label{subsec:open_problems}

以下はPhase~1では解決していない研究問題である。
各問題は将来の研究方向を示すものであり、
解決可能性、唯一の解法または必然的な進展を主張しない。

\begin{problem}[分岐を許す場合の同一個体性の再定式化]
\label{prob:branching_identity}

\textbf{背景：}
\S\ref{sec:acs_def} の仮定~\ref{ass:no_branching}
（非分岐性；強仮定）は、
時刻 $t$ のACS個体実現に対して、
同一個体の継続として対応づけられる時刻 $t+1$ の
個体実現を高々一つとし、
コピー個体または複数の後続実現が生じる分岐を
Phase~1 の対象から除外する。
構造的自己 $S_{\mathrm{struct}}(\Sigma_i)$
（定義~\ref{def:s_struct}）および
時間的自己同一性の予備条件
（定義~\ref{def:time_identity}）は、
この非分岐性を前提として記述されている。
組織ACSの分割・分社も、同一組織の後続が複数になる場合として、この問題に含まれる。

\textbf{問題の定式化：}
分岐（コピー個体・複数の後続実現）を許す場合について、
以下を問う。

\begin{enumerate}
  \item[(i)] \textbf{同一個体性関係の再定式化：}
  時刻 $t$ の個体実現に対して複数の後続実現が生じる場合に、
  いずれを同一個体 $\Sigma_i$ の継続とみなすか、
  複数の実現を同一系譜上の継続として扱うか、
  または分岐後の各実現を別個体として扱うかを定める
  同一個体性関係の形式化。

  \item[(ii)] \textbf{構造的自己と時間的自己同一性の拡張：}
  分岐を許す枠組みにおける構造的自己
  $S_{\mathrm{struct}}$ と時間的自己同一性の予備条件
  （定義~\ref{def:time_identity}）の再定式化、
  および崩壊 $\mathrm{Col}_t$
  （\S\ref{sec:collapse_cond}）の分岐下での扱い。

  \item[(iii)] \textbf{分岐と閉塞・崩壊の関係：}
  分岐の成立と閉塞
  $A_t=\emptyset$ および崩壊
  $\mathrm{Col}_t=1$ との論理的・動学的関係、
  ならびに分岐後の複数実現を含む個体群または系譜としての
  存続をどのように記述するかの整理。
\end{enumerate}

\textbf{現状の障壁：}
Phase~1 は非分岐性（仮定~\ref{ass:no_branching}）のもとで
単一の後続実現を前提としており、
分岐下の同一個体性は、個体の存続に関する新たな
形式的同一性関係を要する。
この関係は状態遷移写像 $f$ の一価性とは別の個体同一性上の条件であり、
Phase~1 でも非分岐性は $f$ とは別の仮定として置いている
（仮定~\ref{ass:no_branching}）。

\textbf{Phase~2以降での取り組み方針：}
個体実現間の時間的継続関係、集合値対応、または系譜構造等を
候補として、複数の後続実現を許す同一個体性の形式体系を導入し、
仮定~\ref{ass:no_branching} のもとでの非分岐的時間同一性を
その特別な場合として回収できる枠組みを検討する。
\end{problem}

\begin{problem}[T-STABパラメータの較正と予測的解釈]
\label{prob:stability}
T-STABで所与とされる
$\underline{s}_{\mathrm{SM}}$、
$\delta_{\max}$、
および$W_T$について、
対象ACSの実装種別と整合する定量的較正条件を明らかにする。
さらに、
\begin{enumerate}
  \item SM評価の参照スケールおよび下限参照値の選択に対する感度；
  \item 窓幅$W_T$と変動上限$\delta_{\max}$に対するロバスト性；
  \item T-STAB違反と閉塞・崩壊・警告条件との
  追加条件付き接続の識別可能性；
  \item 実装データからのパラメータ推定および妥当性検証
\end{enumerate}
を検討する。
Phase~1のT-STAB定義は、
これらの一般的な較正則または予測性能を与えない。
\end{problem}

\begin{problem}[推定誤差と内部更新則の定量的同定]
\label{prob:estimation_error}
TIPAのM1推定誤差、M2評価誤差、内部モデル更新、
偏向性ごとの偏向強度 $\boldsymbol{\beta}_t$（各偏向性の参照基準と尺度の指定を含む）、
$d_{\mathrm{eff},t}$、$\pi_{\tau,t}$、
四健全度等について、実装ごとの読み出し規則、担体の遷移と誤差伝播を同定し、
T-STAB、可行行動集合および閉塞到達への影響を評価する。
Phase~1の抽象的な型保存条件から、一般的な収束率や
誤差分布を導くことはできないため、追加モデルとデータを要する。
\end{problem}

\begin{problem}[核維持応答の獲得条件]
\label{prob:kernel_response_acquisition_conditions}
\textbf{背景：}
命題~\ref{prop:kernel_maintenance} は、
$x_T\in\mathrm{Viab}_f(K)$ のもとで、
すべての $s\ge T$ について $e_s\in E_{\mathrm{req}}(x_s)$ と
$u_s\in U^{\mathrm{Viab}}(x_s)$ が成り立つならば、
すべての $s\ge T$ について $x_s\in\mathrm{Viab}_f(K)$ かつ $A_s\neq\emptyset$ であることを示す。
このうち後者の前件は、実効行動の生成機構（\S\ref{sec:tipa}）と
内部モデルの更新（\S\ref{sec:extended_dynamics}）を通じて成立する場合も成立しない場合もある
（定義~\ref{def:kernel_response_acquisition}）。
本論文は、核維持応答の獲得が起こること、内部モデルの更新が獲得を導くこと、
および $U^{\mathrm{Viab}}(x)\neq\emptyset$ から実際に選ばれる行動が $U^{\mathrm{Viab}}(x)$ に属することの
いずれも一般的な帰結として置かない（備考~\ref{rem:kernel_response_scope}）。
また ACS が $\mathrm{Viab}_f(K)$ や $U^{\mathrm{Viab}}(x)$ を内部に保持または計算することも要求しない
（\S\ref{subsec:description_levels_global}）。

\textbf{問題の定式化：}
\begin{enumerate}
  \item[(i)] 観測（$P_{\mathrm{sens}}$）と内部推定が、
  核維持応答とそれ以外の応答を区別するのに要する情報を保つための条件を特定すること。
  \item[(ii)] 内部モデルの更新と応答の形成・選択の機構に関するどの追加条件のもとで、
  ある時刻以降に核維持応答が選ばれ続ける（獲得が成立する）か、
  およびその時刻の上界を与えうるかを検討すること。
  \item[(iii)] 必要環境の充足 $e_s\in E_{\mathrm{req}}(x_s)$ が途切れる外部環境状態の列のもとで、
  核維持応答の選択と、実現する軌道が $\mathrm{Viab}_f(K)$ にとどまることとの関係を整理すること。
\end{enumerate}

\textbf{現状の障壁：}
Phase~1 は可行性核と核維持応答を形式量として定義するが、
内部機構がそれらを内部情報だけからどのように近似し選択するかについては、
一般的な条件を与えない。

\textbf{取り組み方針：}
決定論的な有限の実現における条件はPhase~1.$\alpha$の詳細化として、
推定の不確実性や外部環境の確率的変動を含む一般の動学への拡張はPhase~2として検討する。
\end{problem}

\begin{problem}[関係記述から因果モデルへの拡張]
\label{prob:relational_causality}
\S\ref{sec:rel_min}の記述的関係構造を保ちながら、
必要な応用に限って、対象別の因果寄与を定義する条件を明らかにする。
具体的には、介入規則、比較軌道、反実仮想、交絡する複数関係、
時間窓の選択およびサイズ汎関数の選択が
因果的解釈へ与える影響を検討する。
記述的関係カテゴリーそのものを因果カテゴリーへ読み替えないことが前提となる。
\end{problem}

\begin{problem}[閉塞・崩壊・回復の観測可能性]
\label{prob:collapse_observability}
有限観測から閉塞・崩壊・回復可能性を判定するために、
どの観測量、履歴長、内部状態へのアクセスおよび外的作用モデルが必要かを明らかにする。
あわせて、閉塞到達時刻・崩壊到達時刻の予測、
経路集合への流入条件、回復判定のロバスト性、
誤判定の上界等を定式化する。
\end{problem}

\begin{problem}[$K$ への多段の回復可能性と必要環境要因]
\label{prob:multistep_recovery}
\textbf{背景：}
\S\ref{sec:three_principles}の一段階臨界必要要因、EUおよび局所FRは
一段の条件である（備考~\ref{rem:eu_one_step_scope}）。
$x_t\notin K$ の時刻の閉塞は、一段では $K$ へ戻れないことだけを表す。
本論文は、そこから有限段で戻れないことを一般的な帰結として置かない
（\S\ref{sec:acs_def}、備考~\ref{rem:blockage_scope_sec2}）。
また $\mathrm{Viab}_f(K)\subseteq K$ であるから、
$x\notin\mathrm{Viab}_f(K)$ は多段の回復不能の根拠にならない。

\textbf{問題の定式化：}
\begin{enumerate}
  \item[(i)] $K$ の外の状態から有限段で $K$ へ到達できることを、
  行動の列と外生背景成分の列の量化順序
  （行動の列を背景成分に先立って固定するか、各段の状態と外部環境状態に応じて選ぶか）
  を明示して定義すること。
  \item[(ii)] その到達に要する必要環境要因を、
  一段階臨界必要要因との関係を含めて特徴づけること。
  \item[(iii)] 多段の回復の途中にある時刻の閉塞と、
  有限段での回復不能とを区別して記述すること。
\end{enumerate}

\textbf{現状の障壁：}
Phase~1は多段の到達可能性を定義せず、
必要環境要因との接続も一段の条件に限っている。

\textbf{取り組み方針：}
決定論的な状態遷移と環境実現（外生背景成分の列と $g_O$）のもとで扱える部分は
Phase~1.$\alpha$の詳細化として、
外部環境の確率的変動を伴う到達可能性はPhase~2として検討する。
\end{problem}

\begin{problem}[組織ACSの動態・同一性・集約則]
\label{prob:social_dynamics}
\S\ref{sec:org_acs}の時刻別参加ACS集合と
組織ACSの記述的適用規約を基礎として、
Phase~2では、参加ACSの加入・離脱・生成・消滅、
組織境界・機能分担・関係構造の時間変化、
組織同一性の保持・復帰条件、
組織水準のEU・FRの扱いとT-STABの具体化、
および組織水準の閉塞・崩壊経路を扱う拡張を構成する。
個体量から組織量への集約則はモデル依存であり、
単一の平均・最小・加重平均を一般則として前提としない。
\end{problem}

\begin{problem}[能力の向上と有効な選択肢の縮小（能力逆説）]
\label{prob:capability_paradox}

\textbf{背景：}
本問題では、認知的処理能力（\S\ref{sec:acs_def}、備考~\ref{rem:cognitive_capacity_policy}）の向上が
内部モデルの表現可能性を広げ、
それが行動の候補（\S\ref{sec:tipa} の $\mathrm{Cand}_t$ 等）を生成しうる幅の拡大に
接続しうる、という仮説を考える。
一方、各時刻に状態遷移写像へ与えられる行動入力は $u_t\in U$ の一つであり
（定義~\ref{def:state_transition}、備考~\ref{rem:f_reading}）、
この点は能力の高低によらない。
本論文では、能力の向上に伴って、自己を維持するための管理負荷と、
他のACSその他の外部対象との関係から生じる制約とが増大し、
長い時間地平で評価したときに有効な選択肢が少数に絞られうる、
という見方を\emph{能力逆説}と呼ぶ。
さらに、時間地平の見通しが長い系ほど、
ACS同士の可行性を維持する協調的な行動が
有効な選択肢として残りうるという見方を、
能力逆説から導かれうる帰結の候補として記録する。
これらは本論文の仮説であり、Phase~1 の定義・命題から得られた一般的な帰結としては主張しない。
ここでいう管理負荷と有効性は、Phase~1 では形式的に定義していない。

\textbf{問題の定式化：}
以下を問う。

\begin{enumerate}
  \item[(i)] \textbf{潜在的な選択肢と有効な選択肢の区別：}
  内部モデルが生成しうる行動の候補（\S\ref{sec:tipa} の $\mathrm{Cand}_t$ 等）と、
  時間地平を指定したときに有効とされる行動の集合とを区別して定義すること。
  有効性の基準には、一段階の可行行動集合 $A_t$（定義~\ref{def:at}）を
  多段の時間地平へ拡張したものが候補となる。
  ただし、決定論的な状態遷移のもとで初期値、外生背景成分の列および $g_O$ を固定すると、
  $T+1$ 段にわたって可行性を保つ行動列の最初の $T$ 段は
  $T$ 段にわたって可行性を保つ行動列であるから、
  そのような行動列の最初の行動がとりうる集合は、
  時間地平を長くしても大きくならない。
  この縮小は能力の高低によらず定義から生じる。
  したがって能力逆説の内容は、この縮小の有無ではなく、
  能力の向上が縮小の程度をどのように変えるかにある。

  \item[(ii)] \textbf{能力と制約の結合：}
  能力の向上が、量的容量・時間応答性・統合の質
  （備考~\ref{rem:cognitive_capacity_policy}）の変化、
  外部対象との関係（\S\ref{sec:rel_min}）その他を通じて、
  有効な選択肢の集合に与える影響を形式化すること。
  能力の向上が可行行動集合を広げる場合や、
  制約そのものを変える場合のように、
  逆向きに働く効果も含めて扱う。

  \item[(iii)] \textbf{協調的行動との関係：}
  複数のACSが互いの外部環境の一部となる状況で、
  時間地平の見通しの長さと、
  ACS同士の可行性を維持する行動が有効な選択肢として残ることとの関係を定式化すること。
  協調的な行動が有効な選択肢に含まれうることは、それだけでは、
  協調的でない行動を有効な選択肢から除く条件を定めない。
  したがって、有効な選択肢が協調的な行動に絞られるための条件そのものを問う。
\end{enumerate}

\textbf{現状の障壁：}
Phase~1 は単一ACSの決定論的記述を中心とし、
複数のACSの同時的な時間発展、
時間地平にわたる有効性の基準、
およびACS同士の可行性の相互依存を形式化していない。
\S\ref{sec:rel_min} の関係記述は因果的寄与を含まず、
\S\ref{sec:org_acs} は組織水準への限定的な構造適用にとどまる。

\textbf{Phase~2以降での取り組み方針：}
Phase~2 では、問題~\ref{prob:social_dynamics}の拡張と合わせて、
複数のACSの相互作用と、時間地平にわたる有効な選択肢の集合を形式化する。
Phase~3 では、問題~\ref{prob:ecit_full}の中で、
有効な選択肢が協調的な行動に絞られる条件と
ECITの協調概念との関係を検討する。
いずれの段階でも能力逆説を前提として採用せず、
成り立つ条件と成り立たない条件の双方を検討の対象とする。
\end{problem}

\begin{problem}[横断的対応付けと実証的検証]
\label{prob:cross_domain}
人間・生物、AI・ロボット、組織等について、
ACSの7タプル、三原理、TIPA、PSL・ISL、
\S\ref{sec:extended_dynamics}変数、関係記述子、閉塞・崩壊の操作的対応を定め、
測定可能性、再現性、予測性能および反証可能性を評価する。
形式的な対応例と、経験的に支持された対応を明確に区別する必要がある。
\end{problem}

\begin{problem}[ECITとの正式接続]
\label{prob:ecit_full}
ACS理論で与えた可行性保存の構造と、
ECITにおける倫理的協調、愛・智慧・責任等の概念との
正式な対応関係を構成する。
その際、記述的な可行性条件から規範的価値判断が
自動的に導かれるとは仮定しない。
また、問題~\ref{prob:social_dynamics}で扱う
Phase~2の組織動態の結果については、
記述的な結果としての位置づけを維持したまま、
ECITとの接続を検討する。
\end{problem}

\begin{remark}[未解決問題の依存関係]
\label{rem:problem_dependencies}
上記の問題には実務上の依存関係があるが、
厳密な論理的必須順序を主張しない。

単一ACSについては、
時間的自己同一性、T-STABの較正、推定誤差、内部更新則、
核維持応答の獲得条件、
および多段の回復可能性の検討を基礎的課題とする。
そのうえで、一つの自然な研究順序は、
\[
  \text{単一ACS}
  \;\longrightarrow\;
  \text{関係・閉塞}
  \;\longrightarrow\;
  \text{組織・社会動態}
  \;\longrightarrow\;
  \text{ECIT接続}
\]
である。

問題~\ref{prob:capability_paradox}は、
組織・社会動態と、ECIT接続の双方にまたがる。
横断的実証検証は各段階と並行して進めうる。
\end{remark}

%----------------------------------------------------------------------
\subsection{Phase~2・3への展望}
\label{subsec:phase_outlook}

本節では、Phase~1の定義・命題を比較基準として保持しつつ、
その強い境界条件を段階的に弱化する研究方向を示す。
以下は研究計画であり、達成を保証するものではない。

\subsubsection*{Phase~2：確率的・動態的拡張}
\label{subsubsec:phase2}

\noindent\textbf{Stage~1（単一ACSのロバスト化）：}
問題~\ref{prob:branching_identity}、
問題~\ref{prob:stability}、
問題~\ref{prob:estimation_error}および
問題~\ref{prob:multistep_recovery}を中心に、
分岐を許す時間的自己同一性、
T-STABパラメータの較正、
推定誤差、更新則、到達性および安定性を
具体モデル上で検討する。
同時に、Phase~1で採用した
決定論・離散時間・コンパクト性・非分岐性等の仮定を、
どの順序で弱化できるかを評価する。

\medskip
\noindent\textbf{Stage~2（関係・閉塞・回復の動態化）：}
問題~\ref{prob:relational_causality}および
問題~\ref{prob:collapse_observability}を中心に、
関係記録と因果推定を分離したまま、
確率的・連続時間的な関係変化、閉塞到達、崩壊到達、
回復可能性および観測誤差を扱う。

\medskip
\noindent\textbf{Stage~3（組織・社会動態の拡張）：}
問題~\ref{prob:social_dynamics}のPhase~2部分を中心に、
参加ACS集合の時間変化、加入・離脱、
組織同一性、機能分担、関係構造、
組織水準のEU・FRの扱いとT-STAB、
および組織水準の閉塞・崩壊経路を
モデル依存の明示的規則として定式化する。
個体水準から組織水準への集約則は、
適用対象ごとに型と根拠を明示し、
一般則として無根拠に固定しない。
あわせて、問題~\ref{prob:capability_paradox}のうち、
複数のACSの相互作用と時間地平にわたる有効な選択肢の集合の形式化を扱う。

\medskip
\noindent\textbf{Stage~4（実証的評価）：}
問題~\ref{prob:cross_domain}を通じて、
各ドメインにおける観測可能量と理論変数の対応を検証する。
Phase~2で導入する追加構造が予測力・説明力・反証可能性を
実際に改善するかを評価し、必要に応じてPhase~1の主張範囲を再限定する。
あわせて、本論文が目指す共有可能な記述言語としての性質、
すなわち、異なる観察者が同じ主体の同じ行動を
ACSの語彙で記述したとき、
その記述がどの範囲で一致し、どこでモデル化の選択によって分かれるかを、
具体的な対象について検討する。

\subsubsection*{Phase~3：ECITとの正式接続}
\label{subsubsec:phase3}

Phase~3では、問題~\ref{prob:ecit_full}を中心に、
ACS理論の記述的な可行性保存構造と
ECITの倫理的協調概念との正式な接続を検討する。
問題~\ref{prob:capability_paradox}のうち協調的な行動に関わる部分も、
この段階で扱う。
この接続においてACS理論が担う役割は、
協調の当事者となる主体それぞれの行動を、
特定の社会や文化の規範を前提とせずに、
共通の構造のもとで記述する層を与えることである
（\S\ref{subsec:background}）。
ECITが扱う協調は、
異なる規範のもとにある主体どうし、
あるいは人間と自律的AIシステムのように成り立ちの異なる主体どうしの間で
問題となる。
そうした主体の間では、
ある主体がなぜそう行動したのかを、
いずれか一方の規範によって判定する前に、
どの可行性上の要請が関わり、
どの内部機構の作動を経てその行動が成立したかを
共有可能な形式で表現できることが、
協調を論じるための前提の一つになると考えられる。
Phase~3では、この記述層の上に
ECITの協調概念をどのように構成できるかを検討する。

ここでは、ACS理論が与える構造的記述、
ECITが扱う協調の条件、
および行動の善し悪しに関する価値判断を、
互いに異なる論理的地位をもつものとして区別する。
Phase~1で定義したACS構造から、
特定の価値判断または社会規範が論理的に自動導出されるとはしない。
共通の記述層を与えることと、
その記述にもとづいて行動の善し悪しを定めることとは、
別の論理的地位をもつ。
規範的要素を導入する場合には、その追加公理・評価基準・
社会的前提をACSの構造的定義と明確に分離する。
また、社会規範そのものを判定の基準としてではなく、
主体が置かれた環境、関係、または内部モデルの内容として
記述の対象に含める可能性も、
この段階で検討の対象となる。
Phase~1は、規範をこのように記述する形式を与えていない。

\begin{remark}[Phase区分の研究戦略的意義]
\label{rem:phase_strategy}
Phase~1--3の区分は、研究の時系列だけでなく、
異なる証明強度とモデル依存性を分離するための区分である。
Phase~1は決定論的・離散時間の最小構造、
Phase~2は確率・連続時間・因果・組織／社会動態等の拡張、
Phase~3はECITとの正式接続を担当する。
後続Phaseの結果によりPhase~1の適用範囲が
修正または限定される可能性をあらかじめ排除しない。
\end{remark}

%----------------------------------------------------------------------
\subsection{結論}
\label{subsec:conclusion}

本論文は、自律認知システム理論（ACS理論）のPhase~1として、
自律的認知システムを
$\Sigma=(X,U,E,f,\Pi,M,\mathrm{Update})$という7タプルで記述し、
可行性保存を中心とする構造的枠組みを構成した。
以下では、本論文が確立した範囲と、確立していない範囲を分けて確認する。

\subsubsection*{Phase~1で確立した主要構造}
\label{subsubsec:core_results}

\begin{itemize}
  \item \textbf{ACSの形式的定義：}
  \S\ref{sec:acs_def}で、状態・行動・必要環境・状態遷移・持続評価・内部モデル・更新作用素からなる
  7タプルと、可行集合$K$、可行行動集合$A_t$その他の基本構造を定めた。
  外部環境状態空間を外生背景成分と行動応答成分に分け、
  構造的自己と整形式実現の区別のもとで、機能的終端と崩壊を状態条件として定義した。

  \item \textbf{三原理：}
  \S\ref{sec:three_principles}で、SM・EU・FRを
  同一の論理地位へ平板化せず、それぞれの時刻別充足述語・定義域・証明強度を区別して形式化し、
  閉塞を臨界必要要因の充足の有無によって二分した。
  また、SM評価・EU・局所FRを有限時間窓上で統合する
  三原理時間統合評価条件（T-STAB）を定義した。

  \item \textbf{TIPA・PSL・ISL：}
  \S\ref{sec:tipa}で、M1情報統合、M2適応的評価、M3応答行動からなるTIPAと、
  その物理的基盤層PSL、実装部位層ISLの機能的構造を整理した。

  \item \textbf{時刻間動学：}
  \S\ref{sec:extended_dynamics}で、原始的な機能の標準写像と読み出し規則を固定し
  環境実現を与えたとき、一段の作動記録がただ一つ定まり次の時刻へ閉じることを示し、
  記憶、偏向性（参照基準に相対する偏りの傾向）とその程度である偏向強度、
  実効的発展自由度、予測精度、実効的時間地平、健全度等を、
  状態・内部モデル・外部環境状態または作動記録の有限区間からの読み出しとして時刻$t$から$t+1$へ接続した。

  \item \textbf{関係性の最小構造：}
  \S\ref{sec:rel_min}で、外部対象との単方向・双方向関係、方向性、評価規則に相対する関係品質、代表カテゴリーおよび六系分類を、
  因果帰属を自動的には含まない記述・診断・比較構造として整理した。

  \item \textbf{閉塞・崩壊：}
  \S\ref{sec:collapse_cond}で、閉塞$A_t=\emptyset$と
  崩壊$\mathrm{Col}_t=1$を明確に分離し、
  外的回復可能性、軌道記録、閉塞経路集合、崩壊到達経路集合および到達時刻を整理した。

  \item \textbf{組織ACSへの限定的構造適用：}
  \S\ref{sec:org_acs}で、ACS集団、ACS組織および組織ACSを区別し、
  組織自身について整形式な7タプル構造が成立する場合に限って、
  個体ACSの参加の有無によらず組織ACSとして扱う記述的適用規約を与え、
  参加ACS数により個体ACS参加型と個体ACS非参加型を区別した。
  組織水準のACS構造が抽象構造型の実現であることにより、
  先行節の三原理、閉塞・崩壊および経路集合の定義を組織水準へ適用する形を与え、
  組織水準の可行集合・可行行動集合と閉塞の特徴づけを示した。
  個体量から組織量への普遍的集約則と、
  組織についての判定規則の一般形式は与えなかった。
\end{itemize}

\subsubsection*{理論の位置づけ}
\label{subsubsec:positioning}

ACS理論は、価値判断そのものを与える規範理論ではなく、
自律認知システムの可行性保存に関わる条件を記述する構造理論として位置づけられる。
本論文が既存のViability Theory、凸解析、動的システム論、制御理論、
認知アーキテクチャその他の理論を包含または一般化したとは主張しない。
本論文で述べる既存理論との関係は、明示した部分における整合性・接続可能性に限定される。

本論文が構造理論として記述を目指したのは、
自律的に振る舞う主体の行動を、
善悪の判定に先立って、
可行性上の要請と内部機構の作動との関係として記述するための
共通の形式である。
Phase~1で与えたのはその決定論的・構造的な基礎であり、
この形式が異なる領域や文化のもとにある主体の行動を
実際にどの程度共有可能な形で記述できるかは、
問題~\ref{prob:cross_domain}および
問題~\ref{prob:ecit_full}として残る。

また、Phase~1の形式的結果が人間、生物、AI、ロボット、組織に
経験的に妥当であることは本論文だけからは導かれない。
その評価は問題~\ref{prob:cross_domain}として残る。

\subsubsection*{主張強度の原則}
\label{subsubsec:principle_check}

本論文では、各結果の論理的地位を
定義的帰結、仮定付き命題、傾向的記述、
未解決問題等に分け、次の原則を維持する：
\[
  \text{主張の強さ}\;\leq\;\text{証明の強さ}
\]
証明または定義から導かれない内容については、
一般的含意、因果効果、安定性、最適性、実証的妥当性または規範的正当性を
自動的に付与しない。

本論文で提示したPhase~1の枠組みが、
自律認知システムと倫理的協調知性の将来的な数理研究に
どの程度有効な基盤となるかは、Phase~2以降の形式的・実証的検証によって評価される。

\medskip
本論文およびACS理論は平和利用を志向して公開される。
ACSによる協調的な社会構造、ならびに人間と他の知性体の
持続的な共生のあり方の探求に資することは、
Phase~1から証明された帰結ではなく、研究上の理念的な希望として位置づける。

\begin{remark}[AI支援の位置づけ]
\label{rem:ai_assistance}
本研究では、言語モデルを、原稿案・数式・証明候補・検証用コードの作成の補助、
文献候補の整理、および論理的整合性・記号整合性・参照整合性の確認に利用した。
理論上の採否、修正判断および最終的な研究責任は著者に帰属する。
著者以外の独立した専門家による本文と検証用コードの全面的な検証は行われていない。
\end{remark}


%================ 付録 A 記号添字規約・B 主要記号一覧・C 定義・命題等の一覧 ================
\appendix
%======================================================================
\section{記号添字規約}
\label{app:notation}

本付録では、本論文の記号に付す添字の位置と順序を、
添字が担う役割の種類ごとに一つに定め、
同じ種類の添字はどの節でも同じ位置に置く。
本付録の規約は記法上の約束であり、
仮定一覧には含めない。
各記号の数学的内容は、それぞれの初出節の定義による。

%----------------------------------------------------------------------
\subsection*{記号添字規約 NIC（添字の位置・順序・書体）}
\label{conv:NIC}

\begin{description}

  \item[(N1) 下付きの識別列]
    下付きには識別列を置く。
    対象を識別する添字は、すべて下付きの識別列に置き、
    \[
      \text{種類}\;\to\;\text{個体}\;\to\;\text{成分}\;\to\;\text{窓}\;\to\;\text{時刻}
    \]
    の順に並べる。
    種類の位置には、説明的ラベル、
    ラベル集合上を走る種類変数、
    および写像・評価関数の依存先を示すパラメータを置き、
    二つ以上が並ぶ場合は、量の直接の種別を先に、
    水準等の修飾を後に置く。
    個体の位置には、ACS個体および外部対象の識別子と、その順序対を置く。
    成分の位置には、TIPAの成分番号、状態座標の番号および組の成分の番号を置く。
    窓の位置には観測窓の長さを置き、
    時刻は下付きの列の最後に置く。
    固定の説明的ラベルを、下付きの種類の位置に置くか、
    (N2) の上付きの変種ラベルとするかは、
    その記号の初出の定義で定め、以後変えない。
    下付きの種類の位置に置いたラベルは識別列の一部であり、
    個体・成分・窓・時刻の添字はその後に続く。
    例えば $d_{\mathrm{eff},i,t}$ は個体 $i$ の時刻 $t$ における
    $d_{\mathrm{eff}}$ を、
    $\bar Q_{i\leftarrow j,W_Q,t}$ は順序対 $i\leftarrow j$ についての
    窓幅 $W_Q$ の時刻 $t$ における値を、
    $E_{\mathrm{ext},\mathrm{org}}$ は組織水準の外部環境状態空間を表す。

  \item[(N2) 上付き]
    上付きに置くのは、次のいずれかに限る。
    \begin{enumerate}[label=(\alph*)]
      \item 演算：冪、逆写像、補集合、参照点、転置、関数空間の指数、
      および括弧付きで表す反復の段数（$\mathrm{Crit}^{(1)}$、$\mathrm{Viab}_f^{(\nu)}(K)$ 等）。
      括弧付きの上付きを段数以外の意味で用いる場合は、使用箇所で宣言する。
      \item 変種ラベル：同じ基本記号で表される対象の変種を区別する、
      固定した説明的ラベル（$E^{\mathrm{bg}}$、$U^{\mathrm{rob}}$、$Q^{\mathrm{inf}}$ 等）。
      変種ラベルは集合上を走る変数ではなく、
      識別列に参与しない。
      変種ラベルを付けた記号には、(N1) の識別列を下付きに続けて置く
      （$Q^{\mathrm{inf}}_{i\leftarrow j,t}$ 等）。
    \end{enumerate}
    個体・成分・窓・時刻の識別は上付きに置かない。
    総和、積、合併等の大型演算子の範囲指定は通常の位置に置き、
    本項の対象としない。

  \item[(N3) 括弧]
    括弧内には、数学的対象を入力とする写像、対応または演算の入力を置く。
    括弧を用いるのは、値がその入力によって定まる場合に限る。
    入力は、宣言された集合の元であるか、
    または (N9) の構造型を満たす任意の構造に対して定まる演算における
    構造そのものである。
    例えば $f(x,u,e)$、$A(x,e)$、$E_{\mathrm{req}}(x)$、$\mathrm{Viab}_f(K)$、
    $S_{\mathrm{struct}}(\Sigma_i)$ の括弧内はこの意味での入力である。

  \item[(N4) 識別と入力]
    値が、識別された個体の実現、すなわちその個体の状態列や環境実現に依存する量は、
    構造を入力とする関数形では書かず、
    識別子と時刻を下付きに置く。
    このため、写像を定義したうえで、その時刻ごとの評価値を
    \[
      A_{i,t}:=A_i(x_{i,t},e_{i,t})
    \]
    のように下付きで表す。
    時刻評価値の記号は各節で定義を与えたものに限り、
    それ以外は $E_{\mathrm{req}}(x_t)$ のように
    写像に状態を入力した形で書く。

  \item[(N5) 個体と水準]
    個体ACSの識別子 $i,i'$ は、\S\ref{sec:org_acs} の識別子集合 $\mathbb I$ の元とする。
    外部対象は $Y_j\in\mathcal Y_{-i}$ と書き、
    添字 $j$ は外部対象を区別する識別子とする。
    $-i$ は $i$ を除く他者を表す派生記法であり、
    $-i$ 自身は識別子ではない。
    個体 $i$ の構造を $\Sigma_i$、
    その成分を $X_i,U_i,E_i,f_i,\Pi_i,M_i,\mathrm{Update}_i$、
    時刻 $t$ の状態を $x_{i,t}$、$M_{i,t}$、
    個体 $i$ に相対的な外部環境状態を
    $e_{i,t}\in E_{\mathrm{ext},i}$ と書く。
    組織水準等の水準名は種類のラベルとして扱い、
    $\Sigma_{\mathrm{org}}$ と書く。
    具体例では、宣言したうえで個体を数字で識別してよい。
    複数個体を扱う文脈では、個体の位置を省略しない。

  \item[(N6) 時刻]
    本論文では、時刻は列または過程の添字として下付きの最後に置き、
    写像の入力としては扱わない。

  \item[(N7) 書体]
    値が走る添字は斜体で書き、
    説明的ラベル、変種ラベルおよび構造部位名は立体で書く。
    ただし、他の記号を名指すラベル（$\mathbb T^{Q}$ の $Q$ 等）は、
    名指す記号の書体に従う。
    空間および集合は大文字で、その元および状態は小文字で書くことを原則とし、
    内部モデルの状態 $M_t$ はその例外とする（(N9)）。
    上付きと下付きの区別には、
    テンソル解析における反変・共変のような
    変換上の意味をもたせない。

  \item[(N8) 省略と式の成分]
    単一個体のみを扱う文脈では個体の添字を省略する。
    下付きが五つ以上並ぶ場合は、
    固定する添字を段落の冒頭で宣言したうえで、
    その添字を省略する。
    式の成分は $[f(x,u,e)]_k$ のように角括弧を用いて表す。

  \item[(N9) 構造と状態]
    定義~\ref{def:acs} が定める7タプルの構造型を抽象構造型と呼び、
    個体の構造 $\Sigma$、$\Sigma_i$ をその実現とする。
    抽象構造型を集合として扱うことは要求せず、
    個体 $\Sigma_i$ を抽象構造型の元とは書かない。
    内部モデルについては、
    構造 $M$（第 $c$ 機構 $M_c$）、
    状態空間 $\mathcal M$（成分 $\mathcal M_c$）、
    状態 $M_t\in\mathcal M$（成分 $M_{c,t}\in\mathcal M_c$）を区別する。
    構造に時刻添字を付けた $\Sigma_{i,t}$ は構造の時間変化を記述するための予約記号とし、
    本論文では用いない（\S\ref{subsec:limitations_list} の L16）。

  \item[(N10) 文字の予約]
    添字に用いる文字は、種類ごとに次の表のとおり予約する。
    予約は下付きの添字として用いる文字に適用する。
    集合の元を表す束縛変数を使用箇所で宣言して用いる場合
    （$q\in\mathrm{Req}_t$ 等）は、本項の対象としない。

\end{description}

\begin{center}
\begin{tabular}{ll}
  \hline
  文字 & 種類 \\
  \hline
  $i,\ i'$ & 個体ACSの識別子（$\mathbb I$ の元） \\
  $j$ & 外部対象の識別子 \\
  $c$ & TIPAの成分番号 \\
  $k$ & 状態座標の番号および組の成分の番号 \\
  $a$ & 行動応答成分の枠（$\mathcal A_O$ の元） \\
  $\kappa$ & 識別の行のラベル（$\mathcal R_\xi$ の元） \\
  $\iota$ & 偏向性の添字（$\mathcal J_\beta$ の元） \\
  $\ell$ & 機能的記憶の種類 \\
  $q$ & 健全度の種類 \\
  $T,\ L$ & 観測窓の長さ \\
  $t,\ s,\ r$ & 時刻 \\
  $\nu$ & 段数（上付きの括弧内） \\
  \hline
\end{tabular}
\end{center}

数字の下付きは、番号で区別して定義した特定の対象
（$M_1$、$P_2$、$\mathsf Z_3$、$t_0$、$\varepsilon_0$ 等）、
または (N5) により宣言した具体例の個体を表す。
軌道 $\sigma$ は軌道記録空間の元として括弧内に置く。
$n$、$m$ は自然数を表す記号として使用箇所で意味を宣言し、
添字には用いない（$\prod_{i=1}^{n}$ のような範囲指定には用いる）。

\begin{description}
  \item[(N11) 書体の使用規則]
    ひげ文字（フラクトゥール体）は、
    原則として母集合・台集合および集合族・写像族などの集合的対象に用いる
    （$\mathfrak K$、$\mathfrak P$ 等）。
    個々の写像、演算子および評価規則には、
    名称型（$\mathrm{rd}$、$\mathrm{Rule}$ 等）または通常の数式記号を用いる。
    書体の違いは記述対象を区別するための表記上の規約であり、
    書体のみから追加的な数学的性質を導かない。
\end{description}

%======================================================================
\section{主要記号一覧}
\label{app:symbols}

本付録は、各節で新たに用いる主要記号を節ごとにまとめた参照用の一覧である。
各記号の定義、型、値域、適用域、未定義の扱いおよび地位は、
本文の初出の定義・備考による。
本付録の記述と本文とが異なって読める場合は、本文に従う。
「定義・初出」の列は、その記号を定める本文の箇所を示す。
「名称・型」の列の「分類A・B・C」は \S\ref{sec:intro} の三分類に基づく
理論上・記述上の位置づけであり、集合・写像・実数値などの数学的な型とは異なる。
7タプルを構成する基礎集合・写像そのものには分類表記を付さない。

%----------------------------------------------------------------------
\subsection{自律認知システムの形式的定義（\S\ref{sec:acs_def}）}
\label{app:symbols_s2}

\begingroup
\small
\setlength{\tabcolsep}{4pt}
\renewcommand{\arraystretch}{1.18}
\begin{longtable}{p{3.0cm} p{3.0cm} p{6.6cm} p{2.0cm}}
\hline
\textbf{記号} & \textbf{名称・型} & \textbf{役割・注記} & \textbf{定義・初出} \\
\hline
\endhead
\hline
\endfoot
\multicolumn{4}{l}{\textbf{基礎構造と自己--環境分離}} \\[2pt]
$\Sigma,\ \Sigma_i$ & ACSの7タプル・構造体 &
単一ACSの形式構造 $\Sigma=(X,U,E,f,\Pi,M,\mathrm{Update})$ を表す。 &
定義~\ref{def:acs} \\[3pt]
$X,\ x_t$ & 内部状態空間・状態 $x_t\in X\subseteq\mathbb R^n$ &
単一ACSの内部状態を表す。 &
定義~\ref{def:acs} \\[3pt]
$U,\ u_t$ & 行動空間・実効行動 $u_t\in U$ &
状態遷移へ入力される実効行動。 &
定義~\ref{def:acs} \\[3pt]
$E_{\mathrm{ext}},\ e_t$ & 外部環境状態空間・実現状態 $e_t\in E_{\mathrm{ext}}$ &
7タプルに対して外生的な背景型と実現外部環境状態。 &
\S\ref{subsec:wellformedness}、定義~\ref{def:Eext_decomposition} \\[3pt]
$E^{\mathrm{bg}},\ b_t$ & 外生背景成分の空間・実現値 &
$E_{\mathrm{ext}}=E^{\mathrm{bg}}\times O$ の第一成分。当該ACSの行動に応答しない外部側の条件。必要環境要因の集合 $E$ とは無関係。 &
定義~\ref{def:Eext_decomposition} \\[3pt]
$O,\ o_t,\ \mathcal A_O,\ \bot_a$ & 行動応答成分の空間・実現値・枠の添字集合・不在状態 &
$O=\prod_{a\in\mathcal A_O}O_a$。$\mathcal A_O$ は有限・固定。枠は個物ではなく位置であり、個物の同一性は追跡しない。 &
定義~\ref{def:Eext_decomposition} \\[3pt]
$g_O$ & 行動応答成分の遷移写像 &
$g_O:O\times X\times U\times E^{\mathrm{bg}}\to O$、$o_{t+1}=g_O(o_t,x_t,u_t,b_t)$。7タプルの成分ではない。 &
定義~\ref{def:env_realization} \\[3pt]
$E,\ \eta$ & 必要環境要因集合・充足判定写像 &
$\eta:X\times E_{\mathrm{ext}}\to\{0,1\}$、$\emptyset\neq E\subseteq\{0,1\}^{X\times E_{\mathrm{ext}}}$。 &
定義~\ref{def:acs}、定義~\ref{def:environment_space} \\[3pt]
$f$ & 状態遷移写像 &
$f:X\times U\times E_{\mathrm{ext}}\to X$、$x_{t+1}=f(x_t,u_t,e_t)$。 &
定義~\ref{def:acs}、定義~\ref{def:state_transition} \\[3pt]
$S_{\mathrm{struct}}(\Sigma)$ & 構造的自己 &
個体として追跡する7タプル構造体の同一性。 &
定義~\ref{def:s_struct} \\[3pt]
$B_{SE}$ & 自己--環境境界 &
内部状態型 $X$ と外部環境状態型 $E_{\mathrm{ext}}$ の構造的分離。 &
定義~\ref{def:b_se} \\[3pt]
TSC,\ CR-1,\ CR-2 & 規約・構成要件 &
型分離規約と成立条件。 &
\S\ref{subsec:wellformedness}、\S\ref{subsec:environment_space}、\S\ref{subsec:three_primitive_functions} \\[3pt]
\multicolumn{4}{l}{\textbf{可行性・整形式実現・機能的終端と崩壊}} \\[2pt]
$\Pi,\ K$ & 持続評価関数・可行集合 &
$\Pi:X\to\mathbb R$、$K=\{x\in X\mid\Pi(x)\ge0\}$。 &
定義~\ref{def:acs}、定義~\ref{def:Pi}、定義~\ref{def:K} \\[3pt]
$\theta$ & 閾値ベクトル &
$\Pi$ の $\min$ 型表現における可行判定の基準（備考~\ref{rem:min_type_pi}）。 &
定義~\ref{def:Pi}、備考~\ref{rem:min_type_pi} \\[3pt]
$A(x,e),\ A_t$ & 可行行動対応・可行行動集合 &
$A(x,e)=\{u\in U\mid f(x,u,e)\in K\}$、$A_t=A(x_t,e_t)$。個体を明示する場合は $A_{i,t}=A_i(x_{i,t},e_{i,t})$。 &
定義~\ref{def:at} \\[3pt]
$\mathrm{WF}_{i,t}$ & 整形式実現指標・分類C &
時刻 $t$ に構造的自己が形式条件を満たして実現するかを記述。 &
定義~\ref{def:wf_realization} \\[3pt]
$S_{\mathrm{dyn},t}$ & 動的自己状態 &
時刻 $t$ の内部状態と可行行動集合の対 $(x_t,A_t)$。 &
定義~\ref{def:s_dyn} \\[3pt]
$D_\eta$ & 要因別充足状態域・分類C &
要因 $\eta$ を充足する外部環境状態が少なくとも一つ存在する内部状態の集合。 &
定義~\ref{def:eta_satisfiable_domain} \\[3pt]
$E_{\mathrm{req}}(x)$ & 必要環境充足集合・分類C &
全必要環境要因を同時に充足する外部環境状態の集合。 &
定義~\ref{def:E_req} \\[3pt]
$\mathcal E_{\mathrm{viab}}^{(1)}(x)$ & 一段階可行環境対応・分類C &
少なくとも一つの可行行動を許す外部環境状態の集合。 &
定義~\ref{def:one_step_viable_environment} \\[3pt]
$\mathbb{T}_{\mathrm{FT},i}$ & 機能的終端判定時刻域 &
崩壊状態が一度も成立していない時刻と、崩壊状態が初めて成立する時刻からなる時刻域。 &
定義~\ref{def:functional_termination} \\[3pt]
$\mathrm{FTS}_{i,t}$ & 機能的終端状態指標 &
$t\in\mathbb{T}_{\mathrm{FT},i}$ について、時刻 $t$ に (FT-1)・(FT-2) がともに成立するか否かを表す。 &
定義~\ref{def:functional_termination} \\[3pt]
$\mathrm{FT}_{i,t}$ & 機能的終端指標 &
時刻 $t$ までに機能的終端状態が成立したことを保持する履歴指標。時刻 $t$ に機能的終端状態にあることは表さない。 &
定義~\ref{def:functional_termination} \\[3pt]
$\mathrm{CS}_{i,t}$ & 崩壊状態指標 &
時刻 $t$ に (SC-1)・(SC-2) がともに成立するか否かを表す。 &
定義~\ref{def:collapse} \\[3pt]
$\mathrm{Col}_{i,t}$ & 崩壊指標 &
時刻 $t$ までに崩壊状態が成立したことを保持する履歴指標。$\mathrm{WF}_{i,t}=0$ とは区別する。 &
定義~\ref{def:collapse} \\[3pt]
$\Upsilon_{\mathrm{adm},i}$ & 許容外的作用集合 &
非回復性の条件 (FT-2)・(SC-2) が参照する外的作用の集合。\S\ref{sec:acs_def} では記号と型のみを宣言し、内容は \S\ref{sec:collapse_cond} で定める。 &
\S\ref{subsubsec:collapse_sec2}、定義~\ref{def:admissible_external_action} \\[3pt]
\multicolumn{4}{l}{\textbf{内部モデル・更新と標準機能表現}} \\[2pt]
$M,\ \mathcal M,\ M_t$ & 内部モデル構造・状態空間・状態 &
機構構造 $M$、状態空間 $\mathcal M$、時刻別状態 $M_t\in\mathcal M$ を区別する。 &
定義~\ref{def:acs}、定義~\ref{def:internal_model}、定義~\ref{def:model_parameter_space} \\[3pt]
$\mathrm{Update}$ & 内部モデル更新作用素 &
$\mathrm{Update}:\mathcal M\times X\times E_{\mathrm{ext}}\to\mathcal M$。 &
定義~\ref{def:acs}、定義~\ref{def:update_operator} \\[3pt]
$\hat e_t$ & 内部環境推定・分類B &
実現外部環境状態 $e_t$ と区別される内部表現。 &
定義~\ref{def:internal_env_estimate} \\[3pt]
$\xi_{\mathrm{bnd},t}$ & 境界帰属識別・分類B &
\S\ref{sec:tipa} で定義される境界帰属識別。構造的自己と独立（命題~\ref{prop:xi_se_independence}）。 &
\S\ref{sec:tipa} \\[3pt]
$\mathcal M_c^{\mathrm L}, \mathcal M_c^{\mathrm W}, M_t^{\mathrm L}, M_t^{\mathrm W}$ & 学習成分・作動成分 &
各機構の内部モデル状態の直積因子と座標射影。独立したタプル成分ではない。 &
定義~\ref{def:model_learning_working} \\[3pt]
$\delta^{\mathrm L}_{c,t},\ \delta^{\mathrm L}_t$ & 学習成分変化量・分類C &
学習停止の判定量。作動成分を含む $\delta_t$ と区別する。 &
定義~\ref{def:update_stop_on_trajectory} \\[3pt]
$\delta_{M,c,t},\ \delta_t$ & TIPA成分別変化量・最大集約変化量・分類C &
内部モデル状態の成分別変化と最大集約距離。 &
定義~\ref{def:tipa_component_change} \\[3pt]
$\lambda_{T,t},\ \boldsymbol{\lambda}_{M,T,t}$ & 時間平均内部モデル変化量・分類C &
時間窓において実際に生じた内部モデル変化の時間平均。認知的処理能力を表す量ではない。 &
定義~\ref{def:update_quality} \\[3pt]
$P_{\mathrm{sens}},P_{\mathrm{disc}},P_{\mathrm{val}},P_{\mathrm{act}}$ & 原始的な機能の標準写像形式・分類B &
感知・識別・評価・行動に対応する標準写像形式。原始的な機能そのものではない。 &
定義~\ref{def:three_primitives} \\[3pt]
$\mathcal{V}$ & 評価値の集合 &
評価の標準写像形式 $P_{\mathrm{val}}$ の値域。順序その他の構造を仮定しない。応答の候補ごとの評価その他の評価情報を表しうる。持続評価関数 $\Pi$ とは異なる。 &
定義~\ref{def:three_primitives} \\[3pt]
\end{longtable}
\endgroup

%----------------------------------------------------------------------
\subsection{三原理：SM・EU・FR（\S\ref{sec:three_principles}）}
\label{app:symbols_s3}

\begingroup
\small
\setlength{\tabcolsep}{4pt}
\renewcommand{\arraystretch}{1.18}
\begin{longtable}{p{3.0cm} p{3.0cm} p{6.6cm} p{2.0cm}}
\hline
\textbf{記号} & \textbf{名称・型} & \textbf{役割・注記} & \textbf{定義・初出} \\
\hline
\endhead
\hline
\endfoot
\multicolumn{4}{l}{\textbf{三原理と時刻別充足述語}} \\[2pt]
SM,\ EU,\ FR & 三原理の制約軸名 &
自己維持・環境利用・可行応答。制約軸であり、規範・目的関数ではない。原理そのものについて成立・不成立を述べない。 &
定義~\ref{def:three_principles_concept} \\[3pt]
$\mathbb{T}_{\mathrm{P},i}$ & 原理判定時刻域・$\subseteq\mathbb{N}_0$ &
$\{t\mid\mathrm{FT}_{i,t}=0\}$。時刻別充足述語はこの時刻域でだけ定義する。 &
定義~\ref{def:principle_time_domain} \\[3pt]
$\mathrm{SM}_t,\ \mathrm{EU}_t,\ \mathrm{FR}_t$ & 時刻別充足述語・真偽値 &
$\mathrm{SM}_t$・$\mathrm{EU}_t$ は $t\in\mathbb{T}_{\mathrm{P},i}$ で、$\mathrm{FR}_t$ はさらに $(x_t,e_t)\in\mathcal{D}_{\mathrm{FR}}$ の時刻でだけ定義する。定義されない時刻は偽に含めない。 &
定義~\ref{def:SM}、定義~\ref{def:EU}、定義~\ref{def:FR} \\[3pt]
$\mathrm{Crit}^{(1)}_\eta,\ \mathrm{Crit}^{(1)}(x)$ & 一段階臨界域・$\subseteq X$・分類C、一段階臨界必要要因集合・$\mathrm{Crit}^{(1)}:X\to2^E$・分類C &
$\eta\in\mathrm{Crit}^{(1)}(x)\iff x\in\mathrm{Crit}^{(1)}_\eta$。時刻添字による略記は置かない。理論上の量であり、ACS 内部の必要性評価とは別（備考~\ref{rem:crit_description_level}）。 &
定義~\ref{def:eu_critical_domain}、定義~\ref{def:eu_fire} \\[3pt]
$\mathcal{D}_{\mathrm{FR}}$ & FR の定義域・$\subseteq X\times E_{\mathrm{ext}}$ &
$t\in\mathbb{T}_{\mathrm{P},i}$ では $(x_t,e_t)\in\mathcal{D}_{\mathrm{FR}}$ と「$\mathrm{EU}_t$ が真」は同値（命題~\ref{prop:fr_domain_eu}）。 &
定義~\ref{def:fr_domain} \\[3pt]
$\varphi_u,\ \sim,\ \mu_{\mathrm{eff}}$ & 可行性保持指示・行動同値・応答多様度・分類C &
$E_{\mathrm{req}}(x)$ 上の可行性保持パターンで $A(x,e)$ を分類。有限性は仮定しない。いずれの時刻別充足述語の定義にも用いない。 &
定義~\ref{def:mu_eff} \\[3pt]
$\widetilde\Pi,\ \partial\Pi(x)$ & $\mathbb R^n$ 拡張と Clarke 一般化勾配・$\subseteq\mathbb R^n$ &
$\widetilde\Pi(x)=\min_k(x_k-\theta_k)$、$\partial\Pi(x)$ は非空コンパクト凸（命題~\ref{prop:pi_subdiff}）。 &
命題~\ref{prop:pi_subdiff} \\[3pt]
$V^-(x)$\newline $d^-(x)$ & 閾値割れ集合・不足量・分類C &
$\Pi$ の最小値が捨てる成分別診断。 &
定義~\ref{def:violation_profile} \\[3pt]
\multicolumn{4}{l}{\textbf{可行性核と構造的環境結合}} \\[2pt]
$D_{\mathrm{req}},\ K_{\mathrm{req}}$ & 必要環境充足状態域・その可行部分 &
$D_{\mathrm{req}}=\{x\mid E_{\mathrm{req}}(x)\neq\emptyset\}$、$K_{\mathrm{req}}=K\cap D_{\mathrm{req}}$。全域非空は仮定しない。 &
定義~\ref{def:req_domain} \\[3pt]
$\mathrm{Viab}_f(K)$ & $E_{\mathrm{req}}$ ロバスト可行性核・$\subseteq K_{\mathrm{req}}$ &
一つの共通行動で $E_{\mathrm{req}}(x)$ 全体に集合内可行性を保つ最大不変集合。標準 Viability Kernel と同一とは主張しない。 &
定義~\ref{def:viab_core} \\[3pt]
$\mathrm{Viab}_f^{(\nu)}(K)$\newline $\mathrm{Hor}_f(x)$ & 有限段頑健可行集合・頑健存続可能段数・分類C &
$\mathrm{Hor}_f$ は予測量でも $\mathrm{SM}_t$ の真理条件でもない。 &
定義~\ref{def:viab_stage}、定義~\ref{def:viab_horizon} \\[3pt]
$\mathrm{Viab}_f^{\mathrm{all}}(K)$\newline $\mathrm{Viab}_f^{\mathrm{any}}(K)$ & 全環境頑健核・楽観核・分類C &
量化の範囲を $E_{\mathrm{ext}}$ 全体に取る比較用の核。$\mathrm{Viab}_f(K)\subseteq\mathrm{Viab}_f^{\mathrm{any}}(K)$。前者が $\mathrm{Viab}_f(K)$ に含まれるのは $D_{\mathrm{req}}$ に含まれるときに限る（命題~\ref{prop:kernel_standard_comparison}）。 &
定義~\ref{def:kernel_all_any} \\[3pt]
$B_{\mathrm{req}}(x,o)$\newline $\mathrm{Viab}^{\mathrm{resp}}_f(K)$ & 背景充足集合・応答込み頑健可行性核・分類C &
行動応答成分 $o$ を固定したときに必要環境要因がすべて充足される外生背景成分の集合と、$g_O$ による応答を考慮した最大不変集合。$\mathrm{Viab}_f(K)$ との包含は一般には成り立たない（命題~\ref{prop:robust_noninclusion}）。 &
定義~\ref{def:viab_resp} \\[3pt]
$N(x,u)$\newline $\mathcal N^{(1)}$ & 応答不能環境集合・応答不能領域・分類C &
$N(x,u)\subseteq E_{\mathrm{req}}(x)$ は行動 $u$ で次状態を $K$ に置けない必要充足環境の集合、$\mathcal N^{(1)}\subseteq X\times E_{\mathrm{ext}}$ はどの行動でも可行性を保てない組の集合。外部環境状態だけの性質ではない。 &
定義~\ref{def:nonresponse_env_set}、定義~\ref{def:nonresponse_region} \\[3pt]
$U^{\mathrm{rob}}(x)$\newline $U^{\mathrm{Viab}}(x)$ & 頑健応答集合・核維持応答集合・分類C &
$N(x,u)=\emptyset$ となる行動の集合と、$E_{\mathrm{req}}(x)$ 全体で次状態を $\mathrm{Viab}_f(K)$ に置く行動の集合。 &
定義~\ref{def:robust_response}、定義~\ref{def:kernel_response} \\[3pt]
$\mathrm{EC}(x)$ & 構造的環境結合条件・分類C &
外部環境状態の差異が環境応答成分の次状態に反映されうること。三原理の時刻別充足述語ではない。 &
定義~\ref{def:env_coupling} \\[3pt]
$I_E,\ \pi_E,\ x_E$ & 環境応答座標・射影・成分 &
$I_E\subseteq\{1,\dots,n\}$、$x_E=\pi_E(x_t)\in\mathbb R^{n_E}$。$x_E$ は内部状態成分で $e_t$・$E$ ではない。 &
仮定~\ref{ass:env_response}、定義~\ref{def:x_E} \\[3pt]
$\partial_{x_E}\Pi(x)$\newline $V^-_E(x)$ & 環境応答方向の勾配射影集合・環境応答成分の閾値割れ &
構造的環境結合条件そのものではない。 &
命題~\ref{prop:pi_E_sensitivity}、定義~\ref{def:ro_related_violation} \\[3pt]
$\chi_{\mathrm{EI},t}$ & 外部相互作用指標・$\{0,1\}$・分類C &
境界横断的外部作用の実現有無。構造的環境結合条件・EU・崩壊とは別概念。 &
定義~\ref{def:phi} \\[3pt]
\multicolumn{4}{l}{\textbf{時間統合評価}} \\[2pt]
$s_{\mathrm{SM},t}$\newline $s_{\mathrm{EU},t}$\newline $s_{\mathrm{FR},t}$ & 三原理評価尺度・分類C &
$s_{\mathrm{SM},t}\in[0,1]$ は正の可行性余裕度、$s_{\mathrm{EU},t},s_{\mathrm{FR},t}\in\{0,1\}$ は時刻別充足述語の真偽で、それぞれ $\mathrm{EU}_t$・$\mathrm{FR}_t$ が定義される時刻でだけ定める。同一物理量ではない。 &
定義~\ref{def:s_sm}、定義~\ref{def:s_eu} \\[3pt]
$s_t,\ S_3$ & 三原理評価ベクトル・$S_3=[0,1]^3$ &
$t\in\mathbb{T}_{\mathrm{P},i}$ かつ $(x_t,e_t)\in\mathcal{D}_{\mathrm{FR}}$ の時刻に定義。 &
定義~\ref{def:S3} \\[3pt]
T-STAB & 三原理時間統合評価条件・有限窓派生条件 &
$t\in\mathbb{T}_{\mathrm{P},i}$ で判定する (T1)--(T4)。力学系の安定性・第四原理・目的関数ではない。 &
定義~\ref{def:tstab} \\[3pt]
\end{longtable}
\endgroup

%----------------------------------------------------------------------
\subsection{TIPA・PSL・ISL（\S\ref{sec:tipa}）}
\label{app:symbols_s4}

\begingroup
\small
\setlength{\tabcolsep}{4pt}
\renewcommand{\arraystretch}{1.18}
\begin{longtable}{p{3.0cm} p{3.0cm} p{6.6cm} p{2.0cm}}
\hline
\textbf{記号} & \textbf{名称・型} & \textbf{役割・注記} & \textbf{定義・初出} \\
\hline
\endhead
\hline
\endfoot
TIPA & 三機構情報処理アーキテクチャ・分類B &
$M_1,M_2,M_3$ により情報統合、適応的評価、応答形成を整理する。 &
\S\ref{subsec:tipa_architecture} \\[3pt]
$M_1$ & 情報統合機構・分類B &
状態・環境推定、識別、時間特徴を形成する。 &
定義~\ref{def:acs}、定義~\ref{def:m1_map} \\[3pt]
$\mathsf S_{M1}$ & M1内部状態空間・分類B &
M1 の当該時刻出力を束ね、次時刻の M1 から参照可能にする内部状態の型。 &
定義~\ref{def:m1_state} \\[3pt]
$M_2$ & 適応的評価機構・分類B &
可行性関連差を代理評価し、Priority 等を形成する。 &
定義~\ref{def:acs}、定義~\ref{def:m2_map} \\[3pt]
$M_3$ & 応答行動機構・分類B &
候補を形成・比較し、TIPA 行動指令を形成する。 &
定義~\ref{def:acs}、定義~\ref{def:m3_map} \\[3pt]
$\xi_t$ & 識別状態・分類B &
M1-b が形成する八軸の標準内部表現。全体の行と内部対象トークンごとの行からなる。 &
定義~\ref{def:id_vector} \\[3pt]
$\mathcal K_t,\ \mathsf{Lab}$ & 内部対象トークン集合・ラベル集合・分類B &
時刻 $t$ の識別状態がもつトークンの有限集合（空を許す）とラベルの有限集合。外部の枠とは同一視しない。 &
定義~\ref{def:id_vector} \\[3pt]
$P_t$ & Priority・分類B &
M2 が形成する標準警告カテゴリ。 &
定義~\ref{def:m2_priority} \\[3pt]
$\mathrm{ev}_t$ & 内部評価情報・分類B &
M2 が M3 へ提供する内部評価情報。 &
定義~\ref{def:m2_eval_space} \\[3pt]
$\mathrm{Req}_t$ & 要求の一覧・分類B &
保持内容から読み出す当該ACSの要求の有限集合。 &
定義~\ref{def:request_signal_types} \\[3pt]
$\mathrm{sig}_t$ & 要請信号・分類B &
M2 が出力する、立ち上がった要求と要請の質の組の有限集合。質の分類は与えない。\S\ref{sec:three_principles} の三原理評価ベクトル $s_t$ とは別。 &
定義~\ref{def:request_signal_types} \\[3pt]
$\mathrm{ev}^{\mathrm{rq}}_t$ & 要求×対象の評価・分類B &
M2 が要求と内部対象トークン（または全体）の組の一部に与える評価。 &
定義~\ref{def:rq_eval_types} \\[3pt]
$\mathrm{stat}_t$ & 目的の状態判断・分類B &
M3 が前時刻の目的に与える判断（継続・達成・未達終了・保留）。 &
定義~\ref{def:purpose_status} \\[3pt]
$\mathrm{Pur}_t$ & 目的の集合・分類B &
M3 が判断として形成する目的の有限集合。 &
定義~\ref{def:purpose_formation} \\[3pt]
$\mathrm{Cand}_t$ & M3候補集合・分類B &
M3 が応答形成の比較対象として扱う行動候補集合。 &
定義~\ref{def:m3_candidate_set} \\[3pt]
$\mathrm{fit}_t$ & 候補適合比較・分類B &
M3 が候補間の内部的適合差を比較する値。目的基準の評価の一つ。 &
定義~\ref{def:m3_fit} \\[3pt]
$\hat g_t,\ \hat G_t(u)$ & 外部側の集合値予測・分類B &
内部モデルが与える行動応答成分の予測写像と、行動 $u$ についての予測集合。 &
定義~\ref{def:env_set_prediction} \\[3pt]
$u_{\mathrm{M3},t}$ & M3行動指令 &
M3 が PSL へ渡す TIPA 側の行動指令。 &
定義~\ref{def:action_decision} \\[3pt]
PSL & 物理的基盤層・分類B &
TIPA と ISL の入出力、基盤維持、局所保護等を接続する。 &
定義~\ref{def:psl} \\[3pt]
$r_{\mathrm{local},t}$ & 局所保護応答 &
PSL-4 により生成されうる局所応答。 &
定義~\ref{def:psl4} \\[3pt]
$\Phi_{\mathrm{resolve}}$ & 実効行動解決写像・分類B &
M3指令と局所保護応答等を実効行動へ接続する。 &
定義~\ref{def:effective_action_resolve} \\[3pt]
$\mathcal B_\Sigma$ & 実装要素全体 &
ACS の実装要素の集合。 &
定義~\ref{def:impl_layer} \\[3pt]
$B_{\mathrm{sens}}$\newline $B_{\mathrm{eff}}$\newline $B_{\mathrm{mem}}$ & ISL標準部分集合・分類B &
感知、作用、情報保持を担いうる実装要素集合。 &
定義~\ref{def:impl_layer} \\[3pt]
$\mathcal T_{\ell,t}$ & 機能的記憶集合・分類B &
保持内容を五つの機能的役割から整理する。 &
定義~\ref{def:fast_memory_sets} \\[3pt]
$\rho_{\mathrm{past}},\ \rho_{\mathrm{eval,past}}$ & 保持記録の読み出し・分類B &
作動成分を担体とする。記録値を実際の過去値とは同一視しない。 &
定義~\ref{def:s4_carrier_contract} \\[3pt]
$\rho_{\mathrm{PSL}},\ \rho_{\Gamma,\cdot},\ \Phi_R$ & 基盤・感度・解決情報の来歴・分類B &
基盤と感度は $x$、解決情報は $(M,I)$ から読む。 &
定義~\ref{def:s4_psl_sources} \\[3pt]
$F_t$ & 機能軸・分類C &
TIPA、PSL、ISL の主要作動水準を横断的に記述する。 &
定義~\ref{def:functional_axis} \\[3pt]
\end{longtable}
\endgroup

%----------------------------------------------------------------------
\subsection{標準作動形式と時刻間の接続（\S\ref{sec:extended_dynamics}）}
\label{app:symbols_s5}

\begingroup
\small
\setlength{\tabcolsep}{4pt}
\renewcommand{\arraystretch}{1.18}
\begin{longtable}{p{3.0cm} p{3.0cm} p{6.6cm} p{2.0cm}}
\hline
\textbf{記号} & \textbf{名称・型} & \textbf{役割・注記} & \textbf{定義・初出} \\
\hline
\endhead
\hline
\endfoot
\multicolumn{4}{l}{\textbf{一段の作動形式と作動記録}} \\[2pt]
R1〜R4 & 本節の責務 &
型整合的な一段の作動記録、次時刻への接続、数学的記録の存在と一意性、後続節への受け渡し。R3 は可行性維持、非閉塞、非崩壊、実装上の活動の継続を意味しない。 &
定義~\ref{def:s5_responsibilities} \\[3pt]
$\Phi_{\mathrm{step}}$ & 一段遷移写像 &
$X\times\mathcal M\times O\times E^{\mathrm{bg}}\to X\times\mathcal M\times O$。$f$、$\mathrm{Update}$、$g_O$ により次の三つ組を返す。$X\times\mathcal M\times O$ は第八のタプル成分ではない。連続性は仮定も結論もしない。 &
定義~\ref{def:s5_step_map} \\[3pt]
$(x_t,\ldots,u_t)_{t\ge0}$ & 標準作動記録・列 &
列 $(x_t,M_t,o_t,b_t,e_t,I_t,r_t,v_t,u_t)_{t\ge0}$。環境実現と初期値を与え、標準作動の一段（定義~\ref{def:s5_standard_step}）に従う列。外生的に与えるのは背景成分の列だけである。 &
定義~\ref{def:s5_operation_record} \\[3pt]
$\mathcal Z_{\mathrm{op}}$ & 状態・環境部分列の集合 &
$\mathcal Z_{\mathrm{op}}\subseteq(X\times\mathcal M\times E_{\mathrm{ext}})^{\mathbb N_0}$。固定した標準写像と $g_O$ に相対する集合。閉集合性・コンパクト性は要求しない。 &
定義~\ref{def:s5_operation_record} \\[3pt]
\multicolumn{4}{l}{\textbf{保持記録と観測経由性}} \\[2pt]
$\operatorname{pr}_k,\ y_{k,t}$ & 保持記録の第 $k$ 成分への射影・対応する時刻 $t$ の値（$k=1,\ldots,7$） &
$\operatorname{pr}_k(\rho_{\mathrm{past}}(M^{\mathrm W}_{t+1}))=y_{k,t}$ のとき、第 $k$ 成分の記録は時刻 $t$ で一致するという。本論文は一致を一般的な帰結として置かず、仮定としても置かない。 &
定義~\ref{def:s5_record_agreement} \\[3pt]
$\Phi_{\mathrm{upd}}$ & 観測経由の更新写像・$\mathcal M\times\mathcal I\to\mathcal M$ &
$\mathrm{Update}(M,x,e)=\Phi_{\mathrm{upd}}(M,P_{\mathrm{sens}}(x,e))$ を満たす全域写像。その存在を仮定とし、連続性は仮定しない。 &
仮定~\ref{ass:s5_observation_factorization} \\[3pt]
\multicolumn{4}{l}{\textbf{時刻間変化}} \\[2pt]
$\mathrm{rd},\ \mathsf D_{\mathrm{rd}},\ \mathrm{rd}_t$ & 固定読み出し・その値域（非空集合）・時刻 $t$ の値 &
$\mathrm{rd}_t=\mathrm{rd}(x_t,M_t,e_t)$。次時刻値は $f$、$\mathrm{Update}$ と環境実現から定め、独立更新写像を置かない。 &
定義~\ref{def:intertime_update} \\[3pt]
\multicolumn{4}{l}{\textbf{分類C記述量}} \\[2pt]
$\mathcal J_{\beta},\ \iota$ & 偏向性の添字集合（非空集合）・その元 &
時刻によらず固定し、有限性は仮定しない。$\iota$ は記述上の識別子であり、偏向性の方向および内容を表さない。 &
定義~\ref{def:bias_tendency} \\[3pt]
$\mathsf{Prim}$ & 原始的な機能の四区分の添字集合 &
$\{\mathrm{sens},\mathrm{disc},\mathrm{val},\mathrm{act}\}$。標準写像形式の添字と対応させて用いる。 &
定義~\ref{def:bias_attribution} \\[3pt]
$\mathrm{loc}_{\beta}$ & 帰属写像・$\mathcal J_{\beta}\to\mathcal P(\mathsf{Prim})\setminus\{\emptyset\}$ &
各偏向性に、その偏りが現れる区分の空でない集合を対応させる。帰属区分は時刻によらず定める。 &
定義~\ref{def:bias_attribution} \\[3pt]
$\beta_{\iota,t},\ \boldsymbol{\beta}_t$ & 偏向強度・非負の量、その族・分類C &
参照基準からの偏りの程度。異なる偏向性の値どうしの比較可能性は仮定しない。未指定・未測定・未同定を理由として $0$ を割り当てない。族は一つの実数へ集約した量ではない。 &
定義~\ref{def:beta_accumulated}、定義~\ref{def:beta_intertime_update} \\[3pt]
$\beta_{\max},\ D_{\max},\ \tau$ & 読み出しの有限上界・予測段数 &
有限な $\beta_{\max}>0$、$D_{\max}>0$ と $\tau\in\mathbb N$、$\tau\ge1$ を固定する。 &
定義~\ref{def:s5_readout_types} \\[3pt]
$\mathcal H$ & 健全度の種類の集合 &
$\{\mathrm{psl},\mathrm{mem},\mathrm{sens},\mathrm{eff}\}$。 &
定義~\ref{def:s5_readout_types} \\[3pt]
$\mathrm{rd}_{\beta,\iota}, \mathrm{rd}_{h,q},\allowbreak \mathrm{rd}_d, \mathrm{rd}_{\pi,\tau}$ & 分類C読み出しの記述規則・全域写像 &
値域はそれぞれ $[0,\beta_{\max}]$、$[0,1]$、$[0,D_{\max}]$、$[0,1]$。具体的規則を一意に導く主張ではなく、対象ごとに指定する記述規則の型である。読み出しの連続性は課さない。 &
定義~\ref{def:s5_readout_types} \\[3pt]
$\pi_{\tau,t}$ & 予測精度係数・$[0,1]$・分類C &
時刻 $t$ の予測の有効性を事後的に記述する係数。未来の実現値を参照し、ACSの時刻 $t$ の内部入力にはしない。 &
定義~\ref{def:s5_prediction_coefficient} \\[3pt]
$\widehat\tau_t$ & 実効的時間地平・$[0,\tau]$・分類C &
$\widehat\tau_t:=\tau\pi_{\tau,t}$。非整数時刻での状態を導入しない。 &
定義~\ref{def:tau_hat_section5} \\[3pt]
$d_{\mathrm{eff},t}$ & 実効的発展自由度・$[0,D_{\max}]$・分類C &
$\mathrm{rd}_d(x_t,M_t)$。$|A_t|$、M2の行動多様度、関係品質とは異なる。数値から可行行動の存在・増減を一般的な帰結として置かない。 &
定義~\ref{def:deff_section5} \\[3pt]
$h_{q,t}$ & PSL・ISL健全度量・$[0,1]$・分類C &
$\mathrm{rd}_{h,q}(x_t)$。$q=\mathrm{psl},\mathrm{mem},\mathrm{sens},\mathrm{eff}$ は基盤、記憶実装、感知、作用の健全度を指定する。共通単位、相互独立性、直接比較可能性を要求しない。 &
定義~\ref{def:s5_health_quantities} \\[3pt]
$\mathrm{Desc}_t$ & 時刻依存量の整理用タプル・分類C &
読み出した量をまとめる八つ組。ACSの7タプルを変更せず、独立状態でも同時刻の内部利用可能情報の束でもない。 &
定義~\ref{def:extended_variable_tuple} \\[3pt]
\end{longtable}
\endgroup

%----------------------------------------------------------------------
\subsection{ACS関係性の最小構造（\S\ref{sec:rel_min}）}
\label{app:symbols_s6}

\begingroup
\small
\setlength{\tabcolsep}{4pt}
\renewcommand{\arraystretch}{1.18}
\begin{longtable}{p{3.0cm} p{3.0cm} p{6.6cm} p{2.0cm}}
\hline
\textbf{記号} & \textbf{名称・型} & \textbf{役割・注記} & \textbf{定義・初出} \\
\hline
\endhead
\hline
\endfoot
$\mathcal Y_{-i}$ & 外部対象集合・集合 &
固定した $\Sigma_i$ の関係記述対象を定める。 &
定義~\ref{def:external_object_set} \\[3pt]
$\mathrm{type}_i$ & 外部対象の三分類・写像 &
外部対象へ $\mathrm{ACS},\mathrm{agent},\mathrm{obj}$ の記述ラベルを割り当てる。 &
定義~\ref{def:object_type} \\[3pt]
$L$ & 観測窓長・正の整数 &
同じ記録を比較する間に固定する窓長。 &
定義~\ref{def:finite_window_diff} \\[3pt]
$W_{L,t}$ & 有限観測窓・時刻集合 &
時刻 $t-L$ から $t$ までの記録の対象時刻を定める。 &
定義~\ref{def:finite_window_diff} \\[3pt]
$\mathrm{sz}_i$ & 局所的サイズ指標・写像・分類C &
採用する場合、$\mathcal D_{\mathrm{sz},i}$ 上の非負実数値写像として用いる。単調性は仮定しない。 &
定義~\ref{def:rel_size_functional} \\[3pt]
$\Delta_L$ & 後向き差分演算子・分類C &
両端の値が定義される場合に $z_t-z_{t-L}$ を表す。 &
定義~\ref{def:finite_window_diff} \\[3pt]
$\mathcal C_{\rightarrow}$ & 単方向ラベル集合・有限集合・分類C &
単方向記述子のカテゴリー成分の値域（15ラベル）。 &
定義~\ref{def:rel_unidirect}、定義~\ref{def:category_set} \\[3pt]
$\mathrm{cat}_{i\leftarrow j,t}$ & 単方向カテゴリー成分・ラベル &
採用した関係記録の単方向カテゴリーを表す。 &
定義~\ref{def:rel_unidirect} \\[3pt]
$\mathrm{cost}_{i\leftarrow j,t}$ & 関係コスト・非負実数 &
当該関係に付随する資源負荷を記録する。 &
定義~\ref{def:rel_unidirect} \\[3pt]
$p_{\mathrm{pert},i\leftarrow j,t}$ & 摂動記録・$P_{\mathrm{pert}}$ の元 &
当該関係の摂動的影響を記録する。確率量ではない。 &
定義~\ref{def:rel_unidirect} \\[3pt]
$\mathrm{Rel}_{i\leftarrow j,t}$ & 単方向関係記述子・六つ組・分類C &
前三成分と $\Sigma_i$ 側の有限窓差分三成分を束ねる六成分の記述形式。特定の外部対象による因果効果を同定しない。 &
定義~\ref{def:rel_unidirect} \\[3pt]
$\mathrm{Rel}_{ij,t}$ & 双方向合成関係・順序対・分類C &
二個体側の単方向関係記述子を順序対として束ねる。 &
定義~\ref{def:bidirect_relation} \\[3pt]
$\mathrm{dir}_{ij,t}$ & 方向性指標・ラベル・分類C &
双方向合成を読むための補助的な主要方向ラベル。 &
定義~\ref{def:relation_direction} \\[3pt]
$\mathcal C_{\leftrightarrow}$ & 双方向成立ラベル集合・有限集合・分類C &
協力・共生・対立の三ラベルの値域。 &
定義~\ref{def:bidirectional_category} \\[3pt]
$\mathrm{cat}_{\leftrightarrow,ij,t}$ & 双方向成立カテゴリー・ラベル・分類C &
採用する場合に双方向成立様式を記述する補助ラベル。 &
定義~\ref{def:bidirectional_category} \\[3pt]
$\mathrm{Rule}_{Q,i\leftarrow j}$ & 関係品質の評価規則・分類C &
記述者が関係ごとに指定し、参照する記録またはモデル、比較の基準、時間幅、三成分と重みの基準を定める。規範的な善悪の基準ではない。 &
定義~\ref{def:Q_rule} \\[3pt]
$\mathcal D_{Q,i,t}$ & 関係品質の評価可能域・$\mathcal Y_{-i}$ の部分集合 &
時刻 $t$ に評価規則が値を与える外部対象の集合。空でありうる。 &
定義~\ref{def:Q_domain} \\[3pt]
$Q_{i\leftarrow j,t}$ & 関係品質・$[-1,1]$・分類C &
$Y_j\in\mathcal D_{Q,i,t}$ のときだけ定義される。成分は $Q^{\mathrm{inf}}_{i\leftarrow j,t}$・$Q^{\mathrm{res}}_{i\leftarrow j,t}$・$Q^{\mathrm{via}}_{i\leftarrow j,t}$。評価不能を $0$ と置かない。$Y_j$ の性質でも内部の評価でもない。 &
定義~\ref{def:Q} \\[3pt]
$\bar Q_{i\leftarrow j,W_Q,t}$ & 関係品質の有限窓平均・定義時 $[-1,1]$ &
窓内で $Q$ が定義される時刻 $\mathbb T^{Q}_{i\leftarrow j,t}$ の上の平均。空なら未定義。 &
定義~\ref{def:Q_window_average} \\[3pt]
$\mathcal C$ & 代表カテゴリー集合・有限集合・分類C &
単方向15種と双方向成立3種からなる18ラベルをまとめる。六系分類は定義~\ref{def:six_system_classification} に示す。 &
定義~\ref{def:category_set}、定義~\ref{def:six_system_classification} \\[3pt]
\end{longtable}
\endgroup

%----------------------------------------------------------------------
\subsection{崩壊・閉塞・崩壊経路集合（\S\ref{sec:collapse_cond}）}
\label{app:symbols_s7}

\begingroup
\small
\setlength{\tabcolsep}{4pt}
\renewcommand{\arraystretch}{1.18}
\begin{longtable}{p{3.0cm} p{3.0cm} p{6.6cm} p{2.0cm}}
\hline
\textbf{記号} & \textbf{名称・型} & \textbf{役割・注記} & \textbf{定義・初出} \\
\hline
\endhead
\hline
\endfoot
$\mathrm{Col}_{i,t}$ & 崩壊指標\newline $\{0,1\}$ 値指標 &
時刻 $t$ までに崩壊状態が成立したことを保持する履歴指標。時刻 $t$ の崩壊状態そのものは $\mathrm{CS}_{i,t}$ で表す。閉塞 $A_{i,t}=\emptyset$、整形式実現の不成立 $\mathrm{WF}_{i,t}=0$、全機能停止帯 $\mathsf{Z}_3$ のいずれとも同一視しない。 &
定義~\ref{def:collapse} \\[3pt]
$\mathsf{Z}_0,\ldots,\mathsf{Z}_5$ & 記述的状態帯\newline 分類ラベル系列 &
閉塞・崩壊周辺の状態を記述的に整理する補助分類。状態空間 $X$ の分割ではなく、添字は全順序または必然的な遷移順序を意味しない。 &
\S\ref{subsec:precollapse_states} \\[3pt]
$\Upsilon_{\mathrm{adm},i}$ & 許容外的作用集合\newline 外的作用の集合 &
外的回復可能性の評価対象として考慮する許容作用域。可能な全外的作用を網羅せず、M1-b の識別状態の空間 $\Xi$ とは記号・型・役割を峻別する。 &
定義~\ref{def:admissible_external_action} \\[3pt]
$\upsilon$ & 許容外的作用\newline $\Upsilon_{\mathrm{adm},i}$ の元 &
$\Sigma_i$ の外部から与えられる作用。ACS自身の選択行動 $u\in U$ とは異なる役割の変数として区別し、遷移写像 $f$ の引数へ暗黙に代入しない。 &
定義~\ref{def:admissible_external_action} \\[3pt]
$\mathrm{Recover}_{i,t}$ & 外的回復述語\newline 決定論的二値補助写像\newline 定義域：$\Upsilon_{\mathrm{adm},i}$\newline 値域：$\{0,1\}$ &
同一個体 $\Sigma_i$ として(SC-1) の組織・結合関係を保存または再成立させ、通常ACSとして後続作動へ復帰しうるかを記述する。作用の実施、外的介入動学、実際の回復到達そのものを表さない。 &
定義~\ref{def:external_recoverability} \\[3pt]
$I_{\mathrm{ACS}}(\sigma)$ & ACS記録有効時刻域\newline $\mathbb{N}_0$ の部分集合 &
軌道記録 $\sigma$ において通常ACS状態量を記録する時刻域。$\{t\mid \mathrm{WF}_{i,t}=1\}$ とは同一でなく、崩壊前の一時的な整形式実現の不成立時刻を含みうる。 &
定義~\ref{def:acs_valid_time_domain} \\[3pt]
$\mathfrak{P}$ & 軌道記録空間\newline 軌道記録の台集合 &
ACS記録有効時刻域上の状態記録と全離散時刻上の崩壊指標列を保持する記録集合。持続評価関数 $\Pi$ とは無関係であり、Phase~1 では一般の位相・距離・確率測度を要求しない。 &
定義~\ref{def:trajectory_space} \\[3pt]
$\sigma$ & 軌道記録\newline $\mathfrak{P}$ の元 &
対象系の決定論的動学と整合する実現記録または決定論的シミュレーション記録。ACS個体 $\Sigma_i$ そのものではなく、その時間発展と崩壊終端の記録を表す。 &
定義~\ref{def:trajectory_space} \\[3pt]
$\mathfrak{P}_{\mathrm{blk}}$ & 閉塞経路集合\newline $\mathfrak{P}$ の部分集合 &
ACS記録有効時刻域内の有限時刻で $A_{i,t}=\emptyset$ が成立する軌道記録の集合。閉塞の持続や崩壊到達を要求せず、確率事象・確率分布族を表さない。 &
定義~\ref{def:blk_path_set} \\[3pt]
$T_{\mathrm{blk}}(\sigma)$ & 閉塞到達時刻\newline 写像\newline 定義域：$\mathfrak{P}_{\mathrm{blk}}$\newline 値域：$\mathbb{N}_0$ &
$\sigma\in\mathfrak{P}_{\mathrm{blk}}$ に対して、閉塞が最初に成立する有限時刻を与える。閉塞非到達軌道には値を割り当てず、$T_{\mathrm{col}}(\sigma)$ と一致するとは限らない。 &
定義~\ref{def:blk_reaching_time} \\[3pt]
$\mathfrak{P}_{\mathrm{col}}$ & 崩壊到達経路集合\newline $\mathfrak{P}$ の部分集合 &
有限時刻で $\mathrm{Col}_{i,t}=1$ が成立する軌道記録の集合。$\mathfrak{P}_{\mathrm{blk}}$ とは独立に定義し、一般的な包含関係または一致を仮定しない。 &
定義~\ref{def:col_path_set} \\[3pt]
$T_{\mathrm{col}}(\sigma)$ & 崩壊到達時刻\newline 写像\newline 定義域：$\mathfrak{P}_{\mathrm{col}}$\newline 値域：$\mathbb{N}_0$ &
$\sigma\in\mathfrak{P}_{\mathrm{col}}$ に対して、崩壊指標が最初に $1$ となる有限時刻を与える。吸収性条件 (C-2) により、以後の全離散時刻で崩壊指標は $1$ に保持される。 &
定義~\ref{def:col_reaching_time} \\[3pt]
\end{longtable}
\endgroup

%----------------------------------------------------------------------
\subsection{組織ACS：社会的組織への構造的適用（\S\ref{sec:org_acs}）}
\label{app:symbols_s8}

\begingroup
\small
\setlength{\tabcolsep}{4pt}
\renewcommand{\arraystretch}{1.18}
\begin{longtable}{p{3.0cm} p{3.0cm} p{6.6cm} p{2.0cm}}
\hline
\textbf{記号} & \textbf{名称・型} & \textbf{役割・注記} & \textbf{定義・初出} \\
\hline
\endhead
\hline
\endfoot
$\mathbb{I}$ & 個体識別子の集合 &
考察対象とする個体ACSの識別子の全体。有限性は仮定しない。 &
定義~\ref{def:org_participation} \\[3pt]
$\mathfrak{O}$ & 組織の集合 &
考察対象とする組織の全体。 &
定義~\ref{def:org_participation} \\[3pt]
$\mathcal O$ & 組織・$\mathfrak{O}$ の元 &
ACS性に先立って識別される社会的組織。 &
定義~\ref{def:org_participation} \\[3pt]
$\mathbb{I}_{\mathrm{org},t}(\mathcal O)$ & 参加個体識別子・$\mathbb{I}$ の部分集合 &
当該時刻の相異なる参加ACSを記録する。 &
定義~\ref{def:org_participation} \\[3pt]
$n_{\mathrm{org},t}(\mathcal O)$ & 参加ACS数・濃度 &
参加する個体ACSの数。組織ACSの区分（個体ACS参加型・非参加型）に用いる。 &
定義~\ref{def:org_participation} \\[3pt]
$\Sigma_{\mathrm{org}}$ & 組織水準の7タプル &
組織自身についてACS構造が成立する場合の記述。 &
定義~\ref{def:org_structural_form} \\[3pt]
$X_{\mathrm{org}}$ & 状態空間・コンパクト集合 &
組織自身の状態の定義域。 &
定義~\ref{def:org_structural_form} \\[3pt]
$E_{\mathrm{org}}$ & 必要環境要因・関数集合 &
組織の必要環境要因を与える第三成分。 &
定義~\ref{def:org_structural_form} \\[3pt]
$E_{\mathrm{ext},\mathrm{org}}$ & 外部環境状態空間・集合 &
組織遷移の外部環境引数。 &
定義~\ref{def:org_structural_form} \\[3pt]
$\Pi_{\mathrm{org}}$ & 持続評価・実数値関数 &
組織自身の可行集合を定める。 &
定義~\ref{def:org_structural_form} \\[3pt]
$K_{\mathrm{org}}$ & 可行集合・状態部分集合 &
組織持続評価の非負領域。 &
定義~\ref{def:org_viable_action} \\[3pt]
$A_{\mathrm{org},t}$ & 可行行動集合・集合 &
一時刻先の組織状態の可行性に関する行動。 &
定義~\ref{def:org_viable_action} \\[3pt]
\end{longtable}
\endgroup

%======================================================================
\section{定義・命題等の一覧}
\label{app:index}

本付録は、本文の定義、仮定、命題、補題、観察記録および問題のすべてと、
本文中で参照される備考を、節ごとに本文での出現順に並べた参照用の一覧である。
各項目の内容は本文による。

%----------------------------------------------------------------------
\subsection*{導入（\S\ref{sec:intro}）}

\begingroup
\small
\setlength{\tabcolsep}{4pt}
\renewcommand{\arraystretch}{1.12}
\begin{longtable}{p{2.4cm} p{10.4cm} r}
定義~\ref{def:abc_classification} & 構造・機能概念の三分類 & \pageref{def:abc_classification} \\
備考~\ref{rem:abc_classification_rules} & 三分類の適用規則 & \pageref{rem:abc_classification_rules} \\
\end{longtable}
\endgroup

%----------------------------------------------------------------------
\subsection*{自律認知システムの形式的定義（\S\ref{sec:acs_def}）}

\begingroup
\small
\setlength{\tabcolsep}{4pt}
\renewcommand{\arraystretch}{1.12}
\begin{longtable}{p{2.4cm} p{10.4cm} r}
備考~\ref{rem:TSC_character} & TSCの性格 & \pageref{rem:TSC_character} \\
備考~\ref{rem:Eext_relative} & 外部環境状態空間の個体相対性 & \pageref{rem:Eext_relative} \\
定義~\ref{def:Eext_decomposition} & 外部環境状態空間の分解 & \pageref{def:Eext_decomposition} \\
備考~\ref{rem:Eext_decomposition_reading} & 枠と行動応答成分の読み & \pageref{rem:Eext_decomposition_reading} \\
定義~\ref{def:primitive_functions} & 原始的な機能（分類A） & \pageref{def:primitive_functions} \\
定義~\ref{def:base_functions} & 基盤機能（分類A） & \pageref{def:base_functions} \\
備考~\ref{rem:primitive_vs_principles} & 原始的な機能と三原理の記述水準の区別 & \pageref{rem:primitive_vs_principles} \\
定義~\ref{def:acs} & 自律認知システム & \pageref{def:acs} \\
備考~\ref{rem:action_space_structure} & 行動空間 $U$ の構造 & \pageref{rem:action_space_structure} \\
備考~\ref{rem:update_type_two_layers} & Updateの型宣言の二層構造 & \pageref{rem:update_type_two_layers} \\
備考~\ref{rem:tuple_nonredundancy} & 7タプルの非冗長性 & \pageref{rem:tuple_nonredundancy} \\
定義~\ref{def:s_struct} & 構造的自己 & \pageref{def:s_struct} \\
定義~\ref{def:wf_realization} & 整形式実現指標（分類C） & \pageref{def:wf_realization} \\
備考~\ref{rem:math_vs_realization} & 数学的水準と実現水準の分離 & \pageref{rem:math_vs_realization} \\
定義~\ref{def:s_dyn} & 動的自己状態 & \pageref{def:s_dyn} \\
定義~\ref{def:b_se} & 自己--環境境界 & \pageref{def:b_se} \\
命題~\ref{prop:xi_se_independence} & $\xi_{\mathrm{bnd},t}$ と $S_{\mathrm{struct}}$ の独立性 & \pageref{prop:xi_se_independence} \\
備考~\ref{rem:xi_se_interpretation} & 認知的識別の失敗と構造の独立性 & \pageref{rem:xi_se_interpretation} \\
仮定~\ref{ass:no_branching} & 個体同一性関係の非分岐性 & \pageref{ass:no_branching} \\
定義~\ref{def:time_identity} & 時間的自己同一性の予備条件 & \pageref{def:time_identity} \\
定義~\ref{def:collapse} & 崩壊状態指標と崩壊指標 & \pageref{def:collapse} \\
定義~\ref{def:functional_termination} & 機能的終端状態指標と機能的終端指標 & \pageref{def:functional_termination} \\
命題~\ref{prop:realized_reestablishment} & 実現した軌道上の再成立 & \pageref{prop:realized_reestablishment} \\
命題~\ref{prop:collapse_basic} & 機能的終端指標と崩壊指標の基本性質 & \pageref{prop:collapse_basic} \\
命題~\ref{prop:termination_order} & 機能的終端と崩壊の順序および崩壊後の整形式実現 & \pageref{prop:termination_order} \\
備考~\ref{rem:termination_hierarchy} & 機能的終端・崩壊と他の状態の区別 & \pageref{rem:termination_hierarchy} \\
備考~\ref{rem:primitive_activity_cessation} & 機能的活動の全面的不成立と待機・静止の区別 & \pageref{rem:primitive_activity_cessation} \\
備考~\ref{rem:collapse_not_principle_failure} & 機能的終端・崩壊と三原理 & \pageref{rem:collapse_not_principle_failure} \\
備考~\ref{rem:collapse_wf_relation} & 機能的終端・崩壊と整形式実現指標の関係 & \pageref{rem:collapse_wf_relation} \\
備考~\ref{rem:post_collapse_record} & 機能的終端後と崩壊後の記録 & \pageref{rem:post_collapse_record} \\
定義~\ref{def:environment_space} & 必要環境要因と必要環境空間 & \pageref{def:environment_space} \\
備考~\ref{rem:factor_functional_class} & 必要環境要因は特定の供給手段ではない & \pageref{rem:factor_functional_class} \\
備考~\ref{rem:eta_coarse_classes} & 必要環境要因の標準的粗分類候補 & \pageref{rem:eta_coarse_classes} \\
定義~\ref{def:eta_satisfiable_domain} & 要因別充足状態域（分類C） & \pageref{def:eta_satisfiable_domain} \\
備考~\ref{rem:env_utilization_function} & 必要環境利用機能 & \pageref{rem:env_utilization_function} \\
備考~\ref{rem:cr1_relativity} & (CR-1a)(i) の外部背景相対性 & \pageref{rem:cr1_relativity} \\
定義~\ref{def:E_req} & 必要環境充足集合（分類C） & \pageref{def:E_req} \\
備考~\ref{rem:factor_nonuniqueness} & 要因分解の非一意性 & \pageref{rem:factor_nonuniqueness} \\
備考~\ref{rem:necessity_identification} & 必要性の存立と同定可能性 & \pageref{rem:necessity_identification} \\
仮定~\ref{ass:compact_spaces} & $X$・$U$ のコンパクト性 & \pageref{ass:compact_spaces} \\
仮定~\ref{ass:f_continuity} & $f$ の連続性 & \pageref{ass:f_continuity} \\
仮定~\ref{ass:Eext_time_fixed} & $E_{\mathrm{ext}}$ の時刻固定性 & \pageref{ass:Eext_time_fixed} \\
定義~\ref{def:state_transition} & 状態遷移 & \pageref{def:state_transition} \\
備考~\ref{rem:f_reading} & 状態遷移写像の読みの規約 & \pageref{rem:f_reading} \\
定義~\ref{def:env_realization} & 行動応答成分の遷移写像と環境実現 & \pageref{def:env_realization} \\
備考~\ref{rem:exogenous_environment} & 実現外部環境状態の生成について & \pageref{rem:exogenous_environment} \\
定義~\ref{def:action_difference} & 作用差（分類C） & \pageref{def:action_difference} \\
仮定~\ref{ass:f_decomposition} & 状態遷移の二成分分解 & \pageref{ass:f_decomposition} \\
備考~\ref{rem:oplus_neutrality} & $\oplus$ と選択的待機の代数的地位 & \pageref{rem:oplus_neutrality} \\
仮定~\ref{ass:env_input_nondegeneracy} & 外部環境入力に対する動学的非退化性 & \pageref{ass:env_input_nondegeneracy} \\
備考~\ref{rem:nondegeneracy_levels} & 非退化条件の水準の区別 & \pageref{rem:nondegeneracy_levels} \\
定義~\ref{def:Pi} & 持続評価関数 & \pageref{def:Pi} \\
備考~\ref{rem:min_type_pi} & $\min$ 型 $\Pi$ の採用理由 & \pageref{rem:min_type_pi} \\
定義~\ref{def:K} & 可行集合 & \pageref{def:K} \\
命題~\ref{prop:theta_from_pi} & 可行性の等価性と $K$ の導出可能性 & \pageref{prop:theta_from_pi} \\
定義~\ref{def:at} & 可行行動対応と可行行動集合 & \pageref{def:at} \\
命題~\ref{prop:At_compact} & 可行行動集合の閉性とコンパクト性 & \pageref{prop:At_compact} \\
定義~\ref{def:blockage_sec2} & 閉塞 & \pageref{def:blockage_sec2} \\
備考~\ref{rem:blockage_not_collapse_sec2} & 閉塞と整形式実現の不成立・崩壊の区別 & \pageref{rem:blockage_not_collapse_sec2} \\
備考~\ref{rem:time_identity_collapse} & 時間的自己同一性と機能的終端・崩壊・閉塞の区別 & \pageref{rem:time_identity_collapse} \\
備考~\ref{rem:blockage_auto} & 閉塞と状態遷移の記述可能性 & \pageref{rem:blockage_auto} \\
備考~\ref{rem:blockage_scope_sec2} & 閉塞の意味の範囲 & \pageref{rem:blockage_scope_sec2} \\
定義~\ref{def:one_step_viable_environment} & 一段階可行環境対応（分類C） & \pageref{def:one_step_viable_environment} \\
備考~\ref{rem:fixed_environment_and_state_dependent_viability} & 一段階可行性と必要環境充足の区別 & \pageref{rem:fixed_environment_and_state_dependent_viability} \\
仮定~\ref{ass:env_viable_nondegeneracy} & 可行行動集合に対する環境入力の非退化性 & \pageref{ass:env_viable_nondegeneracy} \\
命題~\ref{prop:nondegeneracy_implication} & 非退化性の水準間の含意 & \pageref{prop:nondegeneracy_implication} \\
備考~\ref{rem:nondegeneracy_converse} & 逆向きの含意について & \pageref{rem:nondegeneracy_converse} \\
仮定~\ref{ass:viable_realizability} & 可行実現可能性 & \pageref{ass:viable_realizability} \\
備考~\ref{rem:viable_realizability_consequences} & 可行実現可能性から従う事項 & \pageref{rem:viable_realizability_consequences} \\
定義~\ref{def:basic_functions} & 基礎的な機能 & \pageref{def:basic_functions} \\
定義~\ref{def:three_primitives} & 原始的な機能の標準写像形式（分類B） & \pageref{def:three_primitives} \\
命題~\ref{prop:feasibility_sensitivity_scope} & 評価の可行性感応性が要求を課す範囲 & \pageref{prop:feasibility_sensitivity_scope} \\
備考~\ref{rem:P_dependency} & 原始的な機能の標準写像形式の作動依存 & \pageref{rem:P_dependency} \\
仮定~\ref{ass:M_compact} & $\mathcal{M}$ のコンパクト性と成分距離の尺度規約 & \pageref{ass:M_compact} \\
定義~\ref{def:model_learning_working} & 内部モデル状態の学習成分と作動成分 & \pageref{def:model_learning_working} \\
定義~\ref{def:internal_model} & 内部モデル構造 & \pageref{def:internal_model} \\
定義~\ref{def:model_parameter_space} & 内部モデル状態空間 & \pageref{def:model_parameter_space} \\
備考~\ref{rem:M_independence} & $M$ の独立成分としての保持 & \pageref{rem:M_independence} \\
定義~\ref{def:internal_env_estimate} & 内部環境推定（分類B） & \pageref{def:internal_env_estimate} \\
備考~\ref{rem:estimate_gap} & 実現外部環境状態と内部環境推定の乖離 & \pageref{rem:estimate_gap} \\
備考~\ref{rem:M_encodes_representation} & 内部モデル状態と内部表現の保持 & \pageref{rem:M_encodes_representation} \\
定義~\ref{def:tipa_component_change} & TIPA成分別変化量と最大集約距離（分類C） & \pageref{def:tipa_component_change} \\
観察記録~\ref{obs:tipa_integration} & $M_1/M_2/M_3$ の機能的分離方針 & \pageref{obs:tipa_integration} \\
定義~\ref{def:update_operator} & Update作用素の構造 & \pageref{def:update_operator} \\
定義~\ref{def:update_consistent_sequence} & Update整合系列 & \pageref{def:update_consistent_sequence} \\
定義~\ref{def:update_stop_on_trajectory} & 系列上の学習停止 & \pageref{def:update_stop_on_trajectory} \\
定義~\ref{def:update_quality} & 時間平均内部モデル変化量（分類C） & \pageref{def:update_quality} \\
備考~\ref{rem:env_change_response} & 環境変化応答機能；構成的必要機能候補 & \pageref{rem:env_change_response} \\
備考~\ref{rem:cognitive_capacity_policy} & 認知的処理能力の記述方針 & \pageref{rem:cognitive_capacity_policy} \\
\end{longtable}
\endgroup

%----------------------------------------------------------------------
\subsection*{三原理：SM・EU・FR（\S\ref{sec:three_principles}）}

\begingroup
\small
\setlength{\tabcolsep}{4pt}
\renewcommand{\arraystretch}{1.12}
\begin{longtable}{p{2.4cm} p{10.4cm} r}
定義~\ref{def:three_principles_concept} & 三原理の概念的定義 & \pageref{def:three_principles_concept} \\
備考~\ref{rem:conv_ro_layers} & 構造的環境結合・実現相互作用・EU・関係評価の水準区別 & \pageref{rem:conv_ro_layers} \\
備考~\ref{rem:conv_tstab_scope} & T-STAB 地位規約 & \pageref{rem:conv_tstab_scope} \\
備考~\ref{rem:s3_description_levels} & 本節における記述水準 & \pageref{rem:s3_description_levels} \\
定義~\ref{def:principle_time_domain} & 原理判定時刻域 & \pageref{def:principle_time_domain} \\
補題~\ref{lem:principle_time_domain_down} & 原理判定時刻域の下方閉性 & \pageref{lem:principle_time_domain_down} \\
備考~\ref{rem:principle_time_domain_scope} & 原理判定時刻域と整形式実現・軌道記録 & \pageref{rem:principle_time_domain_scope} \\
備考~\ref{rem:example_status} & 例示の地位 & \pageref{rem:example_status} \\
命題~\ref{prop:pi_subdiff} & $\Pi$ の Lipschitz 連続性と Clarke 一般化勾配構造 & \pageref{prop:pi_subdiff} \\
定義~\ref{def:violation_profile} & 閾値割れ集合と不足量ベクトル（分類C） & \pageref{def:violation_profile} \\
命題~\ref{prop:pi_violation_profile} & $\Pi$ と閾値割れ集合の関係 & \pageref{prop:pi_violation_profile} \\
定義~\ref{def:SM} & SM の時刻別充足述語 & \pageref{def:SM} \\
命題~\ref{prop:sm_K_equiv} & SM と可行集合の同値性 & \pageref{prop:sm_K_equiv} \\
命題~\ref{prop:sm_one_step_A} & 一段階の $K$ 保持と $A_t$ & \pageref{prop:sm_one_step_A} \\
備考~\ref{rem:sm_not_collapse} & SM の時刻別充足述語が偽であることと崩壊の区別 & \pageref{rem:sm_not_collapse} \\
定義~\ref{def:req_domain} & 必要環境充足状態域 & \pageref{def:req_domain} \\
定義~\ref{def:viab_core} & $E_{\mathrm{req}}$ ロバスト可行性核 & \pageref{def:viab_core} \\
命題~\ref{prop:viab_core_max} & 最大 $E_{\mathrm{req}}$ 不変集合の存在 & \pageref{prop:viab_core_max} \\
備考~\ref{rem:viability_relation} & Viability Theory との構造的関係 & \pageref{rem:viability_relation} \\
備考~\ref{rem:three_layers} & 形式モデル上の核・内部予測・実現軌道の区別 & \pageref{rem:three_layers} \\
定義~\ref{def:viab_stage} & 有限段 $E_{\mathrm{req}}$ 頑健可行集合（分類C） & \pageref{def:viab_stage} \\
定義~\ref{def:viab_horizon} & 頑健存続可能段数（分類C） & \pageref{def:viab_horizon} \\
命題~\ref{prop:viab_stage_nested} & 有限段頑健可行集合の入れ子性 & \pageref{prop:viab_stage_nested} \\
命題~\ref{prop:viab_stage_core} & $E_{\mathrm{req}}$ ロバスト可行性核と有限段頑健可行集合 & \pageref{prop:viab_stage_core} \\
命題~\ref{prop:viab_stage_finite} & 有限状態空間における可行性核と有限段頑健可行集合 & \pageref{prop:viab_stage_finite} \\
備考~\ref{rem:viab_stage_scope} & 頑健存続可能段数の地位と限定 & \pageref{rem:viab_stage_scope} \\
定義~\ref{def:kernel_all_any} & 全環境頑健核と楽観核（分類C） & \pageref{def:kernel_all_any} \\
命題~\ref{prop:kernel_standard_comparison} & 可行性核と全環境頑健核・楽観核の包含関係 & \pageref{prop:kernel_standard_comparison} \\
備考~\ref{rem:kernel_standard_comparison_scope} & 包含関係の読み方と限定 & \pageref{rem:kernel_standard_comparison_scope} \\
定義~\ref{def:viab_resp} & 応答込み頑健可行性核（分類C） & \pageref{def:viab_resp} \\
命題~\ref{prop:viab_resp_max} & 応答込み頑健可行性核の不変性 & \pageref{prop:viab_resp_max} \\
命題~\ref{prop:robust_noninclusion} & 二つの頑健性の非包含 & \pageref{prop:robust_noninclusion} \\
備考~\ref{rem:robust_noninclusion_reading} & 二つの例の読みと射程 & \pageref{rem:robust_noninclusion_reading} \\
備考~\ref{rem:same_time_alternatives} & 同じ時刻にありえた外部環境状態にわたる量化 & \pageref{rem:same_time_alternatives} \\
仮定~\ref{ass:env_response} & 環境応答成分の構造的決定 & \pageref{ass:env_response} \\
定義~\ref{def:x_E} & 環境応答成分 $x_E$ と座標射影 $\pi_E$ & \pageref{def:x_E} \\
備考~\ref{rem:env_response_observation} & 環境応答成分の観察的対応 & \pageref{rem:env_response_observation} \\
定義~\ref{def:env_coupling} & 構造的環境結合条件（分類C） & \pageref{def:env_coupling} \\
命題~\ref{prop:pi_E_sensitivity} & $\pi_E$ の最小感応性：同値版 & \pageref{prop:pi_E_sensitivity} \\
備考~\ref{rem:pi_E_local_sensitivity} & $x_E$ 方向の局所感応性 & \pageref{rem:pi_E_local_sensitivity} \\
定義~\ref{def:ro_related_violation} & 環境応答成分の閾値割れ集合（分類C） & \pageref{def:ro_related_violation} \\
備考~\ref{rem:diagnostic_role} & $\min$ 型 $\Pi$ の診断限界と補完指標の役割 & \pageref{rem:diagnostic_role} \\
定義~\ref{def:phi} & 外部相互作用指標 $\chi_{\mathrm{EI},t}$（分類C） & \pageref{def:phi} \\
定義~\ref{def:terminal_isolation} & 終端的完全孤立 & \pageref{def:terminal_isolation} \\
定義~\ref{def:eu_critical_domain} & 一段階臨界域（分類C） & \pageref{def:eu_critical_domain} \\
定義~\ref{def:eu_fire} & 一段階臨界必要要因集合（分類C） & \pageref{def:eu_fire} \\
備考~\ref{rem:crit_description_level} & 一段階臨界必要要因集合の記述水準 & \pageref{rem:crit_description_level} \\
定義~\ref{def:EU} & EU の時刻別充足述語 & \pageref{def:EU} \\
命題~\ref{prop:eu_blockage} & 臨界必要要因の非充足と閉塞 & \pageref{prop:eu_blockage} \\
命題~\ref{prop:eu_sm_next} & 臨界必要要因の非充足と次時刻の状態 & \pageref{prop:eu_sm_next} \\
備考~\ref{rem:eu_one_step_scope} & 一段階の必要性の射程 & \pageref{rem:eu_one_step_scope} \\
備考~\ref{rem:eu_cr1_connection} & EU と構成要件 CR-1 の接続範囲 & \pageref{rem:eu_cr1_connection} \\
命題~\ref{prop:eu_fire_domain} & 臨界必要要因と要因別充足状態域 & \pageref{prop:eu_fire_domain} \\
備考~\ref{rem:eu_env_utilization} & 現在状態で充足不能な要因と EU & \pageref{rem:eu_env_utilization} \\
定義~\ref{def:fr_domain} & FR の定義域 & \pageref{def:fr_domain} \\
定義~\ref{def:FR} & FR の時刻別充足述語（局所 FR） & \pageref{def:FR} \\
備考~\ref{rem:fr_scope} & FR の定義上の限定 & \pageref{rem:fr_scope} \\
定義~\ref{def:mu_eff} & 応答多様度（分類C） & \pageref{def:mu_eff} \\
命題~\ref{prop:fr_domain_eu} & 定義域と EU の時刻別充足述語 & \pageref{prop:fr_domain_eu} \\
命題~\ref{prop:fr_one_step} & 系：局所 FR と次状態の $K$ への所属 & \pageref{prop:fr_one_step} \\
命題~\ref{prop:fr_next_sm} & 局所 FR の偽と次時刻の状態 & \pageref{prop:fr_next_sm} \\
命題~\ref{prop:fr_mu_eff} & 局所 FR と応答多様度 & \pageref{prop:fr_mu_eff} \\
備考~\ref{rem:mu_eff_empty_req} & $E_{\mathrm{req}}(x_t)=\emptyset$ の場合の応答多様度 & \pageref{rem:mu_eff_empty_req} \\
備考~\ref{rem:env_basis_intrinsic} & 環境比較基準の内在性 & \pageref{rem:env_basis_intrinsic} \\
備考~\ref{rem:sd_local_scope} & 局所 FR と長期の側面の時間的区別 & \pageref{rem:sd_local_scope} \\
備考~\ref{rem:sd_improvement_aux} & 可行性余裕度の改善可能性 & \pageref{rem:sd_improvement_aux} \\
命題~\ref{prop:viab_stage_fr} & 有限段頑健可行集合と局所 FR & \pageref{prop:viab_stage_fr} \\
命題~\ref{prop:selection_responsibility} & 選択の責務 & \pageref{prop:selection_responsibility} \\
備考~\ref{rem:selection_responsibility} & 選択の責務と構成要件 & \pageref{rem:selection_responsibility} \\
命題~\ref{prop:blockage_dichotomy} & 閉塞の二分 & \pageref{prop:blockage_dichotomy} \\
備考~\ref{rem:blockage_dichotomy} & 二分の地位 & \pageref{rem:blockage_dichotomy} \\
備考~\ref{rem:principles_failure_table} & 三原理の時刻別充足述語と閉塞・整形式実現・崩壊 & \pageref{rem:principles_failure_table} \\
備考~\ref{rem:principle_linkage} & 三原理の連動 & \pageref{rem:principle_linkage} \\
備考~\ref{rem:principle_implication_scope} & 含意関係の限定 & \pageref{rem:principle_implication_scope} \\
仮定~\ref{ass:nonstationary} & 非定常環境 & \pageref{ass:nonstationary} \\
定義~\ref{def:nonresponse_env_set} & 応答不能環境集合と可行な応答の弱点（分類C） & \pageref{def:nonresponse_env_set} \\
定義~\ref{def:nonresponse_region} & 応答不能領域（分類C） & \pageref{def:nonresponse_region} \\
定義~\ref{def:robust_response} & 頑健応答集合と頑健応答不能（分類C） & \pageref{def:robust_response} \\
命題~\ref{prop:robust_nonresponse_viab} & 頑健応答不能と頑健可行集合 & \pageref{prop:robust_nonresponse_viab} \\
定義~\ref{def:theta_exterior} & $K$ 外部持続 & \pageref{def:theta_exterior} \\
仮定~\ref{ass:theta_exterior_blockage} & 可行性核外部持続と閉塞の有限時間接続 & \pageref{ass:theta_exterior_blockage} \\
命題~\ref{prop:exterior_blockage} & $K$ 外部持続と閉塞 & \pageref{prop:exterior_blockage} \\
備考~\ref{rem:exterior_without_blockage} & 可行性核の外への持続と閉塞の関係 & \pageref{rem:exterior_without_blockage} \\
補題~\ref{lem:sd_viab} & 応答多様性縮退と応答の弱点による頑健応答不能 & \pageref{lem:sd_viab} \\
補題~\ref{lem:finite_exit} & 可行性核外部におけるロバスト可行性の喪失 & \pageref{lem:finite_exit} \\
備考~\ref{rem:uniform_bound} & 補題~\ref{lem:finite_exit} の射程 & \pageref{rem:uniform_bound} \\
命題~\ref{prop:sd_robust_nonresponse} & 主命題：応答多様性縮退と応答の弱点の持続による頑健応答不能の持続 & \pageref{prop:sd_robust_nonresponse} \\
命題~\ref{prop:sd_necessity} & 系：応答多様性縮退の持続からの有限時間閉塞；強仮定のもとで & \pageref{prop:sd_necessity} \\
命題~\ref{prop:sd_necessity_fr} & 系：応答多様性縮退の持続と局所 FR & \pageref{prop:sd_necessity_fr} \\
定義~\ref{def:kernel_response} & 核維持応答集合（分類C） & \pageref{def:kernel_response} \\
命題~\ref{prop:kernel_maintenance} & 核に留まる応答を選び続ける軌道 & \pageref{prop:kernel_maintenance} \\
定義~\ref{def:kernel_response_acquisition} & 核維持応答の獲得（分類C） & \pageref{def:kernel_response_acquisition} \\
備考~\ref{rem:kernel_response_scope} & 結果の差と非主張事項 & \pageref{rem:kernel_response_scope} \\
備考~\ref{rem:sd_degree} & 応答多様度の「程度」命題は Phase~2 の課題 & \pageref{rem:sd_degree} \\
備考~\ref{rem:mu_measure} & $\mu_{\mathrm{eff}}$ の測度論的位置づけと Phase~2 への委譲 & \pageref{rem:mu_measure} \\
定義~\ref{def:s3_eval_scales} & 三原理評価表現の参照スケール & \pageref{def:s3_eval_scales} \\
定義~\ref{def:s_sm} & SM 評価尺度（分類C） & \pageref{def:s_sm} \\
定義~\ref{def:s_eu} & EU・FR 評価尺度（分類C） & \pageref{def:s_eu} \\
定義~\ref{def:S3} & 三原理評価ベクトル（分類C） & \pageref{def:S3} \\
定義~\ref{def:tstab_params} & T-STAB 評価パラメータ & \pageref{def:tstab_params} \\
定義~\ref{def:tstab} & T-STAB：三原理時間統合評価条件 & \pageref{def:tstab} \\
命題~\ref{prop:tstab_implies_principles} & T-STAB から有限窓内の時刻別充足述語への含意 & \pageref{prop:tstab_implies_principles} \\
命題~\ref{prop:tstab_zero_margin} & 時刻別充足述語の真と T-STAB & \pageref{prop:tstab_zero_margin} \\
備考~\ref{rem:tstab_descriptive_status} & 時間集約条件の記述上の地位 & \pageref{rem:tstab_descriptive_status} \\
定義~\ref{def:principle_related_functions} & 三原理関連構造的機能（分類A） & \pageref{def:principle_related_functions} \\
備考~\ref{rem:theory_quantity_vs_internal} & 理論上の量への関連と内部での計算の区別 & \pageref{rem:theory_quantity_vs_internal} \\
定義~\ref{def:principles_tipa_connection} & 三原理から TIPA への機能的接続 & \pageref{def:principles_tipa_connection} \\
定義~\ref{def:principles_psl_connection} & 原理関連機能と PSL の接続 & \pageref{def:principles_psl_connection} \\
定義~\ref{def:principles_isl_connection} & 原理関連機能と ISL の接続 & \pageref{def:principles_isl_connection} \\
備考~\ref{rem:principle_function_overlap} & 原理別の主な機能接続は重複しうる & \pageref{rem:principle_function_overlap} \\
\end{longtable}
\endgroup

%----------------------------------------------------------------------
\subsection*{TIPA・PSL・ISL（\S\ref{sec:tipa}）}

\begingroup
\small
\setlength{\tabcolsep}{4pt}
\renewcommand{\arraystretch}{1.12}
\begin{longtable}{p{2.4cm} p{10.4cm} r}
定義~\ref{def:sel_proj} & 決定論的選択射影（分類B） & \pageref{def:sel_proj} \\
定義~\ref{def:s4_carrier_contract} & 同時刻作動の担体規約（分類B） & \pageref{def:s4_carrier_contract} \\
仮定~\ref{ass:functional_concurrency} & TIPA三機構の機能的並行性 & \pageref{ass:functional_concurrency} \\
備考~\ref{rem:tipa_bidirectional} & 機構間の送受信と同一時刻の記述 & \pageref{rem:tipa_bidirectional} \\
定義~\ref{def:s4_output_attribution} & 機構出力の個体への帰属（分類B） & \pageref{def:s4_output_attribution} \\
備考~\ref{rem:internal_steps_time} & 離散時刻と内部処理工程 & \pageref{rem:internal_steps_time} \\
定義~\ref{def:m1_state} & M1内部状態（分類B） & \pageref{def:m1_state} \\
定義~\ref{def:m1_map} & 情報統合機構 $M_1$（分類B） & \pageref{def:m1_map} \\
定義~\ref{def:std_init} & M1-aの予備状態推定（分類B） & \pageref{def:std_init} \\
定義~\ref{def:m1_integration} & M1-aの動態重みと最終統合（分類B） & \pageref{def:m1_integration} \\
定義~\ref{def:s4_m1_sources} & M1補助出力の情報源（分類B） & \pageref{def:s4_m1_sources} \\
備考~\ref{rem:m1_estimate_nontruth} & 状態・環境推定と実現値 & \pageref{rem:m1_estimate_nontruth} \\
定義~\ref{def:s4_state_estimation_error} & 状態推定誤差（分類C） & \pageref{def:s4_state_estimation_error} \\
定義~\ref{def:id_vector} & 識別状態と内部対象トークン（分類B） & \pageref{def:id_vector} \\
備考~\ref{rem:token_generation} & 内部対象トークンの形成と実装依存の範囲 & \pageref{rem:token_generation} \\
定義~\ref{def:m1_temporal_base} & M1-cの時間的基底表現（分類B） & \pageref{def:m1_temporal_base} \\
定義~\ref{def:m1_temporal_features} & M1-cの予測整合性と先行時間特徴量（分類C） & \pageref{def:m1_temporal_features} \\
定義~\ref{def:env_set_prediction} & 外部側の集合値予測（分類B） & \pageref{def:env_set_prediction} \\
定義~\ref{def:prediction_record_check} & 予測集合の記録と二種類の予測外（分類C） & \pageref{def:prediction_record_check} \\
備考~\ref{rem:prediction_reference_difference} & 二つの予測の照合基準の違い & \pageref{rem:prediction_reference_difference} \\
定義~\ref{def:m1_temporal_output} & M1-cの最終時間特徴量（分類B） & \pageref{def:m1_temporal_output} \\
仮定~\ref{ass:m1c_temporal} & M1-cの時間統合構造 & \pageref{ass:m1c_temporal} \\
定義~\ref{def:m2_eval_space} & M2内部評価情報空間（分類B） & \pageref{def:m2_eval_space} \\
定義~\ref{def:request_signal_types} & 要求の一覧と要請信号の型（分類B） & \pageref{def:request_signal_types} \\
定義~\ref{def:rq_eval_types} & 要求×対象の評価の型（分類B） & \pageref{def:rq_eval_types} \\
定義~\ref{def:m2_map} & 適応的評価機構 $M_2$（分類B） & \pageref{def:m2_map} \\
備考~\ref{rem:m2_no_return_scope} & M2 に戻らないことの範囲 & \pageref{rem:m2_no_return_scope} \\
定義~\ref{def:m2a_threat} & M2-a可行性脅威検出機能（分類B） & \pageref{def:m2a_threat} \\
定義~\ref{def:s4_internal_viability_check} & 内部可行判定と理論上の可行判定（分類C） & \pageref{def:s4_internal_viability_check} \\
定義~\ref{def:mech_mu_eff} & 機構層有効行動多様度の一般形（分類C） & \pageref{def:mech_mu_eff} \\
定義~\ref{def:m2b_vuln} & M2-b構造的脆弱性評価機能（分類B） & \pageref{def:m2b_vuln} \\
定義~\ref{def:m2c_latent_function} & M2-c潜在的不安定性評価機能（分類B） & \pageref{def:m2c_latent_function} \\
定義~\ref{def:m2_latent_instability} & 潜在的不安定性の機構層指標（分類C） & \pageref{def:m2_latent_instability} \\
定義~\ref{def:m2_priority} & Priority警告体系（分類B） & \pageref{def:m2_priority} \\
定義~\ref{def:m2_request_signal} & M2による要請信号の生成（分類B） & \pageref{def:m2_request_signal} \\
備考~\ref{rem:error_bounds} & M2評価の情報品質と非定量性 & \pageref{rem:error_bounds} \\
備考~\ref{rem:m2_need_evaluation} & 内部の必要性評価と臨界必要要因 & \pageref{rem:m2_need_evaluation} \\
定義~\ref{def:eval_vs_judgment} & 二つの評価と判断の区別（分類B） & \pageref{def:eval_vs_judgment} \\
定義~\ref{def:m3_map} & 応答行動機構 $M_3$（分類B） & \pageref{def:m3_map} \\
定義~\ref{def:m3a_info} & M3-a統合入力情報（分類B） & \pageref{def:m3a_info} \\
定義~\ref{def:purpose_formation} & 目的の型と目的の形成（分類B） & \pageref{def:purpose_formation} \\
備考~\ref{rem:purpose_hierarchy} & 目的の階層と入れ子 & \pageref{rem:purpose_hierarchy} \\
定義~\ref{def:purpose_status} & 目的の状態判断（分類B） & \pageref{def:purpose_status} \\
定義~\ref{def:m3_candidate_set} & M3候補集合（分類B） & \pageref{def:m3_candidate_set} \\
定義~\ref{def:m3_fit} & 候補適合比較（分類B） & \pageref{def:m3_fit} \\
定義~\ref{def:s4_candidate_sources} & 候補と適合比較の情報来歴（分類B） & \pageref{def:s4_candidate_sources} \\
備考~\ref{rem:fit_vs_identity} & 要請適合性と識別の区別 & \pageref{rem:fit_vs_identity} \\
定義~\ref{def:mode_judgment} & 動作モードの判断（分類B） & \pageref{def:mode_judgment} \\
定義~\ref{def:request_response_mismatch} & 要請と応答の不一致（分類C） & \pageref{def:request_response_mismatch} \\
定義~\ref{def:M3_modes} & M3動作モード（分類B） & \pageref{def:M3_modes} \\
定義~\ref{def:m3_3b_styles} & 3B応答様式（分類B） & \pageref{def:m3_3b_styles} \\
定義~\ref{def:action_decision} & TIPA行動指令の決定（分類B） & \pageref{def:action_decision} \\
定義~\ref{def:psl_availability} & 利用可能性（分類C） & \pageref{def:psl_availability} \\
定義~\ref{def:psl} & PSLの形式的定義（分類B） & \pageref{def:psl} \\
定義~\ref{def:psl1} & PSL-1の形式的定義（分類B） & \pageref{def:psl1} \\
定義~\ref{def:psl2} & PSL-2の形式的定義（分類B） & \pageref{def:psl2} \\
定義~\ref{def:psl3} & PSL-3の形式的定義（分類B） & \pageref{def:psl3} \\
定義~\ref{def:psl4} & PSL-4：局所急性異常への保護応答（分類B） & \pageref{def:psl4} \\
定義~\ref{def:effective_action_resolve} & 実効行動解決写像（分類B） & \pageref{def:effective_action_resolve} \\
備考~\ref{rem:psl4_concurrency} & 局所保護応答と TIPA の並行性 & \pageref{rem:psl4_concurrency} \\
定義~\ref{def:psl_sensitivity} & PSL感度制御写像（分類B） & \pageref{def:psl_sensitivity} \\
定義~\ref{def:s4_psl_sources} & PSLの信号・基盤担体と標準写像への接続（分類B） & \pageref{def:s4_psl_sources} \\
仮定~\ref{ass:psl_minimal} & PSLの最小構造条件と非認知性 & \pageref{ass:psl_minimal} \\
命題~\ref{prop:purpose_self_attribution} & 目的の自己帰属性 & \pageref{prop:purpose_self_attribution} \\
備考~\ref{rem:self_attribution_scope} & 自己帰属性の射程と異常域 & \pageref{rem:self_attribution_scope} \\
定義~\ref{def:req_involvement} & 要求の関与（分類C） & \pageref{def:req_involvement} \\
定義~\ref{def:action_classes} & 行動の区分（分類C） & \pageref{def:action_classes} \\
定義~\ref{def:purpose_action_difference} & 目的に関わる作用差と一段の条件充足（分類C） & \pageref{def:purpose_action_difference} \\
命題~\ref{prop:purpose_ad_slot} & 外部側の区分と枠の作用差 & \pageref{prop:purpose_ad_slot} \\
定義~\ref{def:result_tuple} & 一つの行動の結果の記述（分類C） & \pageref{def:result_tuple} \\
定義~\ref{def:purpose_course} & 目的の経過の記述（分類C） & \pageref{def:purpose_course} \\
定義~\ref{def:impl_layer} & 実装部位層 ISL（分類B） & \pageref{def:impl_layer} \\
定義~\ref{def:sensor} & 感覚実装役割（分類B） & \pageref{def:sensor} \\
定義~\ref{def:effector} & 効果実装役割（分類B） & \pageref{def:effector} \\
定義~\ref{def:mem_substrate} & 記憶実装役割（分類B） & \pageref{def:mem_substrate} \\
定義~\ref{def:impl_psl_connection} & ISLとPSLの接続規約（分類B） & \pageref{def:impl_psl_connection} \\
仮定~\ref{ass:impl_structure} & 実装部位層の構造的仮定 & \pageref{ass:impl_structure} \\
定義~\ref{def:integrated_fast_processing} & 統合的高速処理（分類B） & \pageref{def:integrated_fast_processing} \\
定義~\ref{def:fast_cross_layer} & 高速処理の三層横断性（分類C） & \pageref{def:fast_cross_layer} \\
備考~\ref{rem:automaticity_deliberation} & Automaticityと熟慮依存度 & \pageref{rem:automaticity_deliberation} \\
定義~\ref{def:fast_memory_sets} & 五機能的記憶集合（分類B） & \pageref{def:fast_memory_sets} \\
仮定~\ref{ass:fast_structural_consistency} & 高速処理時の構造整合条件 & \pageref{ass:fast_structural_consistency} \\
定義~\ref{def:functional_axis} & 機能軸（分類C） & \pageref{def:functional_axis} \\
定義~\ref{def:s4_primitive_composition} & 標準作動形式との合成対応（分類B） & \pageref{def:s4_primitive_composition} \\
\end{longtable}
\endgroup

%----------------------------------------------------------------------
\subsection*{標準作動形式と時刻間の接続（\S\ref{sec:extended_dynamics}）}

\begingroup
\small
\setlength{\tabcolsep}{4pt}
\renewcommand{\arraystretch}{1.12}
\begin{longtable}{p{2.4cm} p{10.4cm} r}
定義~\ref{def:s5_responsibilities} & 本節の責務 & \pageref{def:s5_responsibilities} \\
定義~\ref{def:s5_standard_step} & 標準作動の一段 & \pageref{def:s5_standard_step} \\
定義~\ref{def:update_section5} & Updateの参照と一段遷移写像 & \pageref{def:update_section5} \\
定義~\ref{def:memory_intertime_update} & 保持内容の時刻間読み出し & \pageref{def:memory_intertime_update} \\
命題~\ref{prop:s5_one_step} & P1：一段の整形式性 & \pageref{prop:s5_one_step} \\
定義~\ref{def:s5_operation_record} & 標準作動記録 & \pageref{def:s5_operation_record} \\
命題~\ref{prop:s5_continuation} & P2：数学的作動記録の継続 & \pageref{prop:s5_continuation} \\
命題~\ref{prop:s5_causality} & P3：作動記録の因果性 & \pageref{prop:s5_causality} \\
定義~\ref{def:s5_record_agreement} & 保持記録の一致（分類C） & \pageref{def:s5_record_agreement} \\
仮定~\ref{ass:s5_observation_factorization} & Updateの観測経由性 & \pageref{ass:s5_observation_factorization} \\
命題~\ref{prop:s5_observation_causality} & 観測経由性の検査条件と内部モデルの因果性 & \pageref{prop:s5_observation_causality} \\
定義~\ref{def:intertime_update} & 担体から誘導される時刻間変化 & \pageref{def:intertime_update} \\
定義~\ref{def:learning_intertime_change} & 学習・忘却の時刻間的位置づけ（分類B） & \pageref{def:learning_intertime_change} \\
備考~\ref{rem:s5_fast_change} & 高速処理・Automaticityと状態の変化 & \pageref{rem:s5_fast_change} \\
備考~\ref{rem:s5_learning_viability} & 学習と可行性構造 & \pageref{rem:s5_learning_viability} \\
定義~\ref{def:bias_tendency} & 偏向性（分類C） & \pageref{def:bias_tendency} \\
定義~\ref{def:bias_attribution} & 偏向性の帰属区分（分類C） & \pageref{def:bias_attribution} \\
定義~\ref{def:beta_accumulated} & 偏向強度（分類C） & \pageref{def:beta_accumulated} \\
備考~\ref{rem:beta_multiple_biases} & 偏向性の抽象的構造記述と実装記述の区別 & \pageref{rem:beta_multiple_biases} \\
備考~\ref{rem:bias_formation_vs_attribution} & 偏向性の形成と帰属の区別 & \pageref{rem:bias_formation_vs_attribution} \\
定義~\ref{def:s5_readout_types} & 読み出しの型と解析範囲 & \pageref{def:s5_readout_types} \\
定義~\ref{def:beta_intertime_update} & 偏向強度の時刻間読み出し & \pageref{def:beta_intertime_update} \\
定義~\ref{def:s5_prediction_coefficient} & 予測精度係数（分類C） & \pageref{def:s5_prediction_coefficient} \\
定義~\ref{def:tau_hat_section5} & 実効的時間地平（分類C） & \pageref{def:tau_hat_section5} \\
定義~\ref{def:deff_section5} & 実効的発展自由度（分類C） & \pageref{def:deff_section5} \\
定義~\ref{def:s5_health_quantities} & PSL・ISL健全度量（分類C） & \pageref{def:s5_health_quantities} \\
備考~\ref{rem:s5_health_noncollapse} & 健全度・機能軸・診断の区別 & \pageref{rem:s5_health_noncollapse} \\
定義~\ref{def:extended_variable_tuple} & 時刻依存量の整理用タプル（分類C） & \pageref{def:extended_variable_tuple} \\
定義~\ref{def:s5_downstream_interface} & 後続節への記録インターフェース & \pageref{def:s5_downstream_interface} \\
\end{longtable}
\endgroup

%----------------------------------------------------------------------
\subsection*{ACS関係性の最小構造（\S\ref{sec:rel_min}）}

\begingroup
\small
\setlength{\tabcolsep}{4pt}
\renewcommand{\arraystretch}{1.12}
\begin{longtable}{p{2.4cm} p{10.4cm} r}
定義~\ref{def:external_object_set} & 外部対象集合 & \pageref{def:external_object_set} \\
定義~\ref{def:object_type} & 外部対象の三分類 & \pageref{def:object_type} \\
仮定~\ref{ass:obj_stable} & 外部対象分類の固定 & \pageref{ass:obj_stable} \\
定義~\ref{def:rel_size_functional} & 可行行動集合の局所的サイズ指標（分類C） & \pageref{def:rel_size_functional} \\
定義~\ref{def:finite_window_diff} & 有限観測窓における差分記法（分類C） & \pageref{def:finite_window_diff} \\
備考~\ref{rem:finite_window_diff_limit} & 有限観測窓差分の制限とサイズ汎関数の委譲 & \pageref{rem:finite_window_diff_limit} \\
定義~\ref{def:rel_unidirect} & 単方向関係記述子（分類C） & \pageref{def:rel_unidirect} \\
備考~\ref{rem:rel_unidirect_notation} & 記号上の注意 & \pageref{rem:rel_unidirect_notation} \\
定義~\ref{def:rel_unidirect_type} & 単方向関係記述子の型 & \pageref{def:rel_unidirect_type} \\
仮定~\ref{ass:unidirect_scope} & 単方向関係記述の対象範囲 & \pageref{ass:unidirect_scope} \\
仮定~\ref{ass:bidirect_scope} & 双方向合成関係の標準対象 & \pageref{ass:bidirect_scope} \\
定義~\ref{def:bidirect_relation} & 双方向合成関係（分類C） & \pageref{def:bidirect_relation} \\
定義~\ref{def:relation_direction} & 方向性指標（分類C） & \pageref{def:relation_direction} \\
備考~\ref{rem:dir_undefined} & 方向性判定の未定義部分 & \pageref{rem:dir_undefined} \\
定義~\ref{def:homo_hetero_relation} & 同種相互関係と異種相互関係 & \pageref{def:homo_hetero_relation} \\
備考~\ref{rem:homo_hetero_rel} & 同種・異種の分類の制限 & \pageref{rem:homo_hetero_rel} \\
備考~\ref{rem:no_mutual_form} & 同種相互形を持たないカテゴリー & \pageref{rem:no_mutual_form} \\
定義~\ref{def:bidirectional_category} & 双方向成立カテゴリー（分類C） & \pageref{def:bidirectional_category} \\
備考~\ref{rem:competition_directionality} & 競争の方向性に関する特殊性 & \pageref{rem:competition_directionality} \\
備考~\ref{rem:nACS_limit} & 非ACS対象との双方向関係の制限 & \pageref{rem:nACS_limit} \\
定義~\ref{def:Q_rule} & 関係品質の評価規則（分類C） & \pageref{def:Q_rule} \\
定義~\ref{def:Q_domain} & 関係品質評価可能域 & \pageref{def:Q_domain} \\
定義~\ref{def:Q} & 関係品質評価 $Q$ & \pageref{def:Q} \\
備考~\ref{rem:Q_internal_formal} & 関係品質と内部の評価の区別 & \pageref{rem:Q_internal_formal} \\
定義~\ref{def:Q_sign} & $Q$ の符号による関係品質分類 & \pageref{def:Q_sign} \\
命題~\ref{prop:Q_not_RO_condition} & $Q$ と構造的環境結合条件の定義式上の相互非参照 & \pageref{prop:Q_not_RO_condition} \\
定義~\ref{def:contaminated_ro} & 負の関係品質を伴う環境結合（分類C） & \pageref{def:contaminated_ro} \\
定義~\ref{def:Q_window_average} & 関係品質の有限窓評価 & \pageref{def:Q_window_average} \\
備考~\ref{rem:Q_not_tstab_core} & $Q$ を T-STAB の普遍条件へ組み込まない & \pageref{rem:Q_not_tstab_core} \\
定義~\ref{def:category_set} & 代表カテゴリーラベル集合（分類C） & \pageref{def:category_set} \\
定義~\ref{def:six_system_classification} & 六系分類（分類C） & \pageref{def:six_system_classification} \\
備考~\ref{rem:category_phase1a} & カテゴリー成立条件の未定義性 & \pageref{rem:category_phase1a} \\
備考~\ref{rem:derived_composite_categories} & 派生カテゴリーと複合カテゴリー & \pageref{rem:derived_composite_categories} \\
備考~\ref{rem:recognized_relation_category} & カテゴリーと認識された関係性 & \pageref{rem:recognized_relation_category} \\
備考~\ref{rem:phase1a_delegation} & Phase~1.$\alpha$ への委譲事項 & \pageref{rem:phase1a_delegation} \\
\end{longtable}
\endgroup

%----------------------------------------------------------------------
\subsection*{崩壊・閉塞・崩壊経路集合（\S\ref{sec:collapse_cond}）}

\begingroup
\small
\setlength{\tabcolsep}{4pt}
\renewcommand{\arraystretch}{1.12}
\begin{longtable}{p{2.4cm} p{10.4cm} r}
備考~\ref{rem:s7_principle_levels} & 本節における三原理と機能の記述水準 & \pageref{rem:s7_principle_levels} \\
備考~\ref{rem:blockage_time_dependence} & 閉塞は吸収状態として定義されない & \pageref{rem:blockage_time_dependence} \\
命題~\ref{prop:blockage_viability_kernel} & 閉塞と $E_{\mathrm{req}}$ ロバスト可行性核の関係 & \pageref{prop:blockage_viability_kernel} \\
備考~\ref{rem:health_not_zone_identity} & 健全度量と状態帯を同一視しない & \pageref{rem:health_not_zone_identity} \\
定義~\ref{def:admissible_external_action} & 許容外的作用集合（分類C） & \pageref{def:admissible_external_action} \\
定義~\ref{def:external_recoverability} & 外的回復述語と外的回復可能性（分類C） & \pageref{def:external_recoverability} \\
命題~\ref{prop:recover_excludes_termination} & 外的回復可能性と終端状態指標 & \pageref{prop:recover_excludes_termination} \\
定義~\ref{def:acs_valid_time_domain} & ACS記録有効時刻域（分類C） & \pageref{def:acs_valid_time_domain} \\
定義~\ref{def:trajectory_space} & 軌道記録空間（分類C） & \pageref{def:trajectory_space} \\
備考~\ref{rem:trajectory_space_scope} & 軌道記録空間の力学的意味 & \pageref{rem:trajectory_space_scope} \\
定義~\ref{def:blk_path_set} & 閉塞経路集合（分類C） & \pageref{def:blk_path_set} \\
定義~\ref{def:blk_reaching_time} & 閉塞到達時刻（分類C） & \pageref{def:blk_reaching_time} \\
備考~\ref{rem:sm_wf_not_blk_path} & SM の時刻別充足述語が偽であること・整形式実現の不成立と閉塞経路集合への帰属の関係 & \pageref{rem:sm_wf_not_blk_path} \\
定義~\ref{def:col_path_set} & 崩壊到達経路集合（分類C） & \pageref{def:col_path_set} \\
定義~\ref{def:col_reaching_time} & 崩壊到達時刻（分類C） & \pageref{def:col_reaching_time} \\
備考~\ref{rem:principles_not_col_path} & 三原理の時刻別充足述語・T-STAB・関係品質と崩壊到達の関係 & \pageref{rem:principles_not_col_path} \\
備考~\ref{rem:internal_quantities_not_principles} & 内部時刻量と三原理の時刻別充足述語を同一視しない & \pageref{rem:internal_quantities_not_principles} \\
\end{longtable}
\endgroup

%----------------------------------------------------------------------
\subsection*{組織ACS：社会的組織への構造的適用（\S\ref{sec:org_acs}）}

\begingroup
\small
\setlength{\tabcolsep}{4pt}
\renewcommand{\arraystretch}{1.12}
\begin{longtable}{p{2.4cm} p{10.4cm} r}
定義~\ref{def:org_terms} & 集団・組織・社会の記述上の区分 & \pageref{def:org_terms} \\
定義~\ref{def:org_class_terms} & ACS集団・ACS組織・組織ACS & \pageref{def:org_class_terms} \\
備考~\ref{rem:org_terms_nonimplication} & 組織と関係対象分類 & \pageref{rem:org_terms_nonimplication} \\
定義~\ref{def:org_participation} & 参加ACSの時刻別記録 & \pageref{def:org_participation} \\
定義~\ref{def:org_structural_form} & 組織水準のACS構造 & \pageref{def:org_structural_form} \\
備考~\ref{rem:org_state_closure} & 組織状態と遷移の一価性 & \pageref{rem:org_state_closure} \\
定義~\ref{def:org_acs} & 組織ACSの記述的適用規約 & \pageref{def:org_acs} \\
備考~\ref{rem:org_individual_nonimplication} & 個体水準と組織水準の判定の区別 & \pageref{rem:org_individual_nonimplication} \\
定義~\ref{def:org_viable_action} & 組織水準の可行集合と可行行動集合 & \pageref{def:org_viable_action} \\
命題~\ref{prop:org_blockage_characterization} & 組織水準の閉塞の特徴づけ & \pageref{prop:org_blockage_characterization} \\
備考~\ref{rem:org_principles_conditional} & 先行節の定義の組織水準への適用 & \pageref{rem:org_principles_conditional} \\
備考~\ref{rem:org_qualification_loss} & 時刻別資格と組織の継続 & \pageref{rem:org_qualification_loss} \\
備考~\ref{rem:org_failure_scope} & 閉塞・崩壊の適用域 & \pageref{rem:org_failure_scope} \\
\end{longtable}
\endgroup

%----------------------------------------------------------------------
\subsection*{限界・未解決問題・Phase~2・3への展望（\S\ref{sec:limitations_and_outlook}）}

\begingroup
\small
\setlength{\tabcolsep}{4pt}
\renewcommand{\arraystretch}{1.12}
\begin{longtable}{p{2.4cm} p{10.4cm} r}
問題~\ref{prob:branching_identity} & 分岐を許す場合の同一個体性の再定式化 & \pageref{prob:branching_identity} \\
問題~\ref{prob:stability} & T-STABパラメータの較正と予測的解釈 & \pageref{prob:stability} \\
問題~\ref{prob:estimation_error} & 推定誤差と内部更新則の定量的同定 & \pageref{prob:estimation_error} \\
問題~\ref{prob:kernel_response_acquisition_conditions} & 核維持応答の獲得条件 & \pageref{prob:kernel_response_acquisition_conditions} \\
問題~\ref{prob:relational_causality} & 関係記述から因果モデルへの拡張 & \pageref{prob:relational_causality} \\
問題~\ref{prob:collapse_observability} & 閉塞・崩壊・回復の観測可能性 & \pageref{prob:collapse_observability} \\
問題~\ref{prob:multistep_recovery} & $K$ への多段の回復可能性と必要環境要因 & \pageref{prob:multistep_recovery} \\
問題~\ref{prob:social_dynamics} & 組織ACSの動態・同一性・集約則 & \pageref{prob:social_dynamics} \\
問題~\ref{prob:capability_paradox} & 能力の向上と有効な選択肢の縮小（能力逆説） & \pageref{prob:capability_paradox} \\
問題~\ref{prob:cross_domain} & 横断的対応付けと実証的検証 & \pageref{prob:cross_domain} \\
問題~\ref{prob:ecit_full} & ECITとの正式接続 & \pageref{prob:ecit_full} \\
\end{longtable}
\endgroup


%================ 参考文献 ================
\printbibliography

\end{document}
